\documentclass[10pt,a4paper]{revtex4-2}
\usepackage{xcolor,amsmath,csquotes,graphicx,enumitem,amsthm,amsfonts,amssymb,centernot,comment,marginnote,tikz,fancyhdr,bm}
\usepackage[a4paper,left=1.0cm,top=1.5cm,bottom=1.5cm,right=1.0cm]{geometry}
\usepackage{multirow}
\usepackage[many]{tcolorbox}
\usepackage{hyperref}

\tcbset{
  before skip=10pt,
  after skip=10pt
}

\newtcolorbox{verbbox-modal}
{
%breakable,
  enhanced,
  colframe = black,
  colback  = yellow!20,
  %coltitle = #2!20!black,  
  %title    = {#1},
}

\newtcolorbox{verbbox}
{
%breakable,
  enhanced,
  colframe = black,
  %colback  = red!20,
  %coltitle = #2!20!black,  
  %title    = {#3},
}

\newtcolorbox{verbbox-nom}
{
%breakable,
  enhanced,
  colframe = black,
  colback  = red!20,
  %coltitle = #2!20!black,  
  %title    = {#3},
}

\newtcolorbox{verbbox-acc}
{
%breakable,
  enhanced,
  colframe = black,
  colback  = blue!20,
  %coltitle = #2!20!black,  
  %title    = {#1},
}

\newtcolorbox{verbbox-dat}
{
%breakable,
  enhanced,
  colframe = black,
  colback  = green!30,
  %coltitle = #2!20!black,  
  %title    = {#3},
}

\newtcolorbox{verbbox-sep}
{
%breakable,
  enhanced,
  colframe = black,
  colback  = orange!20,
  %coltitle = #2!20!black,  
  %title    = {#3},
}




%Title-Page
\title{German Verb Table}



\begin{document}

\tableofcontents


\section{Tenses}

\subsection{Präsens (Present)}
The Präsens tense is the most frequently used tense in German and covers a wide range of meanings. It is used to describe actions happening now, general truths, habitual actions, and very often future events when a time expression makes the meaning clear. For example, \textit{Ich lerne Deutsch} can mean “I learn German” or “I am learning German,” and \textit{Morgen gehe ich nach Hause} uses the present tense to refer to the future. German does not have a separate continuous tense, so Präsens fulfills both roles.

\subsection{Perfekt (Present Perfect)}
The Perfekt tense is the main tense for talking about completed actions in the past, especially in spoken German. It is formed using the present tense of \textit{haben} or \textit{sein} plus the past participle of the main verb. For instance, \textit{Ich habe gegessen} means “I ate” or “I have eaten.” In everyday conversation, the Perfekt almost completely replaces the simple past for most verbs, making it essential for understanding and speaking modern German.

\subsection{Präteritum (Simple Past)}
The Präteritum tense is primarily used in written German, such as novels, newspapers, and formal narratives, to describe past events. In spoken German, it is largely limited to a small group of very common verbs, especially \textit{sein}, \textit{haben}, and the modal verbs, as in \textit{Ich war müde} or \textit{Ich musste gehen}. Although learners often encounter this tense in textbooks, its practical spoken use is quite limited.

\subsection{Plusquamperfekt (Past Perfect)}
The Plusquamperfekt tense is used to describe an action that had already been completed before another action in the past. It is formed with the simple past of \textit{haben} or \textit{sein} plus the past participle. An example is \textit{Ich hatte gegessen, bevor er kam}, meaning “I had eaten before he arrived.” This tense helps clarify the sequence of past events and is used when the order of actions is important.

\subsection{Futur I (Future)}
The Futur I tense is formed with \textit{werden} plus the infinitive of the main verb and is used to talk about future actions or make assumptions about the present or future. For example, \textit{Ich werde lernen} means “I will study,” while \textit{Er wird krank sein} expresses a guess, such as “He is probably sick.” However, Germans often prefer the present tense with a time expression instead of Futur I, making this tense less common than learners might expect.

\newpage

\section{Auxillary Verbs}

%sein
\begin{verbbox-nom}\addcontentsline{toc}{subsubsection}{sein (to be)}
    {\Large \bf sein} (to be) \hfill Type: \textbf{irregular, + Nominativ} \hfill Auxiliary Verb: \textbf{sein} 
     \vspace{5pt}
    \hrule
    \vspace{5pt}
  \begin{tabular}{rl}
     \multicolumn{2}{l}{\textbf{PRÄSENS}}\\
    ich & \textbf{bin}\\
    du & \textbf{bist}\\
    er/sie/es & \textbf{ist}\\
    wir & \textbf{sind}\\
    ihr & \textbf{seid}\\
    sie/Sie & \textbf{sind}\\
  \end{tabular}
  \hfill
     \begin{tabular}{rl}
    \multicolumn{2}{l}{\textbf{PRÄTERITUM}}\\
    ich & \textbf{war}\\
    du & \textbf{warst}\\
    er/sie/es & \textbf{war}\\
    wir & \textbf{waren}\\
    ihr & \textbf{wart}\\
    sie/Sie & \textbf{waren}\\
  \end{tabular}
    \hfill
     \begin{tabular}{rl}
    \multicolumn{2}{l}{\textbf{KONJUNKTIV II}}\\
    ich & \textbf{wäre}\\
    du & \textbf{wärest}\\
    er/sie/es & \textbf{wär}\\
    wir & \textbf{wären}\\
    ihr & \textbf{wäret}\\
    sie/Sie & \textbf{wären}\\
  \end{tabular}
  \hfill
  \begin{tabular}{rl}
     \multicolumn{2}{l}{\textbf{IMPERATIV}}\\
   & \textbf{sei (du)!}\\
   & \textbf{seien wir!}\\
   & \textbf{seid (ihr)!}\\
   & \textbf{seiren Sie!}\\
   & \vspace{10pt}
  \end{tabular}

    \vspace{10pt}
 
    \begin{tabular}{ll}
     \multicolumn{2}{l}{\textbf{PERFEKT}}\\
   &  sein + \textbf{gewesen}\\
  \end{tabular}
\hfill
   \begin{tabular}{ll}
   
    \multicolumn{2}{l}{\textbf{FUTURE I}}\\
   & werden + \textbf{sein}\\
 
  \end{tabular}
\hfill
   \begin{tabular}{ll}
      \multicolumn{2}{l}{\textbf{PLUSQUAMPERFEKT}}\\
   & waren + \textbf{gewesen}\\
  \end{tabular}
  
  \vspace{10pt}

   \begin{tabular}{lll}
      \multicolumn{3}{l}{\textbf{MIT PRÄPOSITIONEN}}\\
& \textbf{sein in} + Dat & (to be located in)\\
& \textbf{sein auf} + Dat & (to be on / at)\\
&\textbf{sein bei} + Dat & (to be with /at)\\
&\textbf{sein von} + Dat & (to be from / made of)\\
&\textbf{sein aus} + Dat & (to be from / made of)\\
&\textbf{sein für} + Akk & (to be for / in favour of)\\
&\textbf{sein mit} + Dat & (to be with / equipped with)
  \end{tabular}
  \hfill
  % \begin{tabular}{lll}
  %    \multicolumn{3}{l}{\textbf{TRENNBARE VERBEN}}\\
% \vspace{70pt}
 % \end{tabular}



  \vspace{-5pt}

\end{verbbox-nom}

%haben
\begin{verbbox-acc}\addcontentsline{toc}{subsubsection}{haben (to have)}
    {\Large \bf haben} (to have) \hfill Type: \textbf{irregular, + Akkusativ} \hfill Auxiliary Verb: \textbf{haben} 
     \vspace{5pt}
    \hrule
    \vspace{5pt}
  \begin{tabular}{rl}
     \multicolumn{2}{l}{\textbf{PRÄSENS} }\\
    ich & \textbf{habe}\\
    du & \textbf{hast}\\
    er/sie/es & \textbf{hat}\\
    wir & \textbf{haben}\\
    ihr & \textbf{habt}\\
    sie/Sie & \textbf{haben}\\
  \end{tabular}
  \hfill
     \begin{tabular}{rl}
    \multicolumn{2}{l}{\textbf{PRÄTERITUM} }\\
    ich & \textbf{hatte}\\
    du & \textbf{hattest}\\
    er/sie/es & \textbf{hatte}\\
    wir & \textbf{hatten}\\
    ihr & \textbf{hattet}\\
    sie/Sie & \textbf{hatten}\\
  \end{tabular}
  \hfill
       \begin{tabular}{rl}
    \multicolumn{2}{l}{\textbf{KONJUNKTIV II} }\\
    ich & \textbf{hätte}\\
    du & \textbf{hättest}\\
    er/sie/es & \textbf{hätte}\\
    wir & \textbf{hätten}\\
    ihr & \textbf{hättet}\\
    sie/Sie & \textbf{hätten}\\
  \end{tabular}
  \hfill
  \begin{tabular}{rl}
     \multicolumn{2}{l}{\textbf{IMPERATIV} }\\
   & \textbf{hab(e) (du)!}\\
   & \textbf{haben wir!}\\
   & \textbf{habt (ihr)!}\\
   & \textbf{haben Sie!}\\
   & \vspace{10pt}
  \end{tabular}

    \vspace{10pt}
 
    \begin{tabular}{rl}
     \multicolumn{2}{l}{\textbf{PERFEKT}}\\
    &haben + \textbf{gehabt}\\
  \end{tabular}
\hfill
   \begin{tabular}{rl}
   
    \multicolumn{2}{l}{\textbf{FUTURE I} }\\
    &werden + \textbf{haben}\\
 
  \end{tabular}
\hfill
   \begin{tabular}{rl}
      \multicolumn{2}{l}{\textbf{PLUSQUAMPERFEKT}}\\
    &hatten + \textbf{gehabt}\\
  \end{tabular}
  


  \vspace{10pt}

   \begin{tabular}{lll}
      \multicolumn{3}{l}{\textbf{MIT PRÄPOSITIONEN}}\\
& \textbf{haben an} + Dat & (to have a problem / interest /fault)\\
& \textbf{haben für} + Akk & (to have for / feel towards)\\
&\textbf{haben von} + Dat & (to have heard / get from)\\
&\textbf{haben auf} + Akk & (to be in the mood for)\\
&\textbf{haben vor} + Dat & (to plan / intend)
  \end{tabular}
 % \hfill
  % \begin{tabular}{lll}
   %   \multicolumn{3}{l}{\textbf{TRENNBARE VERBEN}}\\
%\vspace{50pt}
 % \end{tabular}




  \vspace{-5pt}


\end{verbbox-acc}

%werden
\begin{verbbox-acc}\addcontentsline{toc}{subsubsection}{werden (to become)}
    {\Large \bf werden} (to become) \hfill Type: \textbf{irregular} \hfill Auxiliary Verb: \textbf{sein} 
     \vspace{5pt}
    \hrule
    \vspace{5pt}
  \begin{tabular}{rl}
     \multicolumn{2}{l}{\textbf{PRÄSENS} }\\
    ich & \textbf{werde}\\
    du & \textbf{wirst}\\
    er/sie/es & \textbf{wird}\\
    wir & \textbf{werden}\\
    ihr & \textbf{werdet}\\
    sie/Sie & \textbf{werden}\\
  \end{tabular}
  \hfill
     \begin{tabular}{rl}
    \multicolumn{2}{l}{\textbf{PRÄTERITUM} }\\
    ich & \textbf{wurde}\\
    du & \textbf{wurdest}\\
    er/sie/es & \textbf{wurde}\\
    wir & \textbf{wurden}\\
    ihr & \textbf{wurdet}\\
    sie/Sie & \textbf{wurden}\\
  \end{tabular}
    \hfill
     \begin{tabular}{rl}
    \multicolumn{2}{l}{\textbf{KONJUNKTIV II} }\\
    ich & \textbf{würde}\\
    du & \textbf{würdest}\\
    er/sie/es & \textbf{würde}\\
    wir & \textbf{würden}\\
    ihr & \textbf{würdet}\\
    sie/Sie & \textbf{würden}\\
  \end{tabular}
  \hfill
  \begin{tabular}{rl}
     \multicolumn{2}{l}{\textbf{IMPERATIV} }\\
   & \textbf{wurde (du)!}\\
   & \textbf{wurden wir!}\\
   & \textbf{werdet (ihr)!}\\
   & \textbf{wurden Sie!}\\
   & \vspace{10pt}
  \end{tabular}

    \vspace{10pt}
 
    \begin{tabular}{rl}
     \multicolumn{2}{l}{\textbf{PERFEKT}}\\
   & sein + \textbf{geworden}\\
  \end{tabular}
\hfill
   \begin{tabular}{rl}
   
    \multicolumn{2}{l}{\textbf{FUTURE I} }\\
    &werden + \textbf{werden}\\
 
  \end{tabular}
\hfill
   \begin{tabular}{rl}
      \multicolumn{2}{l}{\textbf{PLUSQUAMPERFEKT}}\\
   & waren + \textbf{geworden}\\
  \end{tabular}
  
\vspace{10pt}

\begin{tabular}{lll}
\multicolumn{3}{l}{\textbf{MIT PRÄPOSITIONEN}}\\
& ---\\
\end{tabular}
\hfill
\begin{tabular}{lll}
\multicolumn{3}{l}{\textbf{TRENNBARE VERBEN}}\\
& ---\\
\end{tabular}

\vspace{-5pt}

\end{verbbox-acc}

%lassen
\begin{verbbox}\addcontentsline{toc}{subsubsection}{lassen (to let)}
    {\Large \bf lassen} (to let) \hfill Type: \textbf{irregular} \hfill Auxiliary Verb: \textbf{haben} 
     \vspace{5pt}
    \hrule
    \vspace{5pt}
  \begin{tabular}{rl}
     \multicolumn{2}{l}{\textbf{PRÄSENS} }\\
    ich & \textbf{lasse}\\
    du & \textbf{lässt}\\
    er/sie/es & \textbf{lässt}\\
    wir & \textbf{lassen}\\
    ihr & \textbf{lasst}\\
    sie/Sie & \textbf{lassen}\\
  \end{tabular}
  \hfill
     \begin{tabular}{rl}
    \multicolumn{2}{l}{\textbf{PRÄTERITUM} }\\
    ich & \textbf{ließ}\\
    du & \textbf{ließt}\\
    er/sie/es & \textbf{ließ}\\
    wir & \textbf{ließen}\\
    ihr & \textbf{ließt}\\
    sie/Sie & \textbf{ließen}\\
  \end{tabular}
     \hfill
     \begin{tabular}{rl}
    \multicolumn{2}{l}{\textbf{KONJUNKTIV II} }\\
    ich & \textbf{ließe}\\
    du & \textbf{ließest}\\
    er/sie/es & \textbf{ließe}\\
    wir & \textbf{ließen}\\
    ihr & \textbf{ließet}\\
    sie/Sie & \textbf{ließen}\\
  \end{tabular}
  \hfill
  \begin{tabular}{rl}
     \multicolumn{2}{l}{\textbf{IMPERATIV} }\\
   & \textbf{lass(e) (du)!}\\
   & \textbf{lassen wir!}\\
   & \textbf{lasst (ihr)!}\\
   & \textbf{lassen Sie!}\\
   & \vspace{10pt}
  \end{tabular}

    \vspace{10pt}
 
    \begin{tabular}{rl}
     \multicolumn{2}{l}{\textbf{PERFEKT}}\\
   & haben + \textbf{gelassen}\\
  \end{tabular}
\hfill
   \begin{tabular}{rl}
   
    \multicolumn{2}{l}{\textbf{FUTURE I} }\\
    &werden + \textbf{lassen}\\
 
  \end{tabular}
\hfill
   \begin{tabular}{rl}
      \multicolumn{2}{l}{\textbf{PLUSQUAMPERFEKT}}\\
   & hatten + \textbf{gelassen}\\
  \end{tabular}
  
\vspace{10pt}

\begin{tabular}{lll}
\multicolumn{3}{l}{\textbf{MIT PRÄPOSITIONEN}}\\
& ---\\
\end{tabular}
\hfill
\begin{tabular}{lll}
\multicolumn{3}{l}{\textbf{TRENNBARE VERBEN}}\\
& ---\\
\end{tabular}

\vspace{-5pt}

\end{verbbox}

\section{Modal Verbs}

%dürfen
\begin{verbbox-modal}\addcontentsline{toc}{subsubsection}{dürfen (to be allowed to)}
    {\Large \bf dürfen} (to be allowed to) \hfill Type: \textbf{Modal} \hfill Auxiliary Verb: \textbf{haben} 
     \vspace{5pt}
    \hrule
    \vspace{5pt}
  \begin{tabular}{rl}
     \multicolumn{2}{l}{\textbf{PRÄSENS} }\\
    ich & \textbf{darf}\\
    du & \textbf{darfst}\\
    er/sie/es & \textbf{darf}\\
    wir & \textbf{dürfen}\\
    ihr & \textbf{dürft}\\
    sie/Sie & \textbf{dürfen}\\
  \end{tabular}
  \hfill
     \begin{tabular}{rl}
    \multicolumn{2}{l}{\textbf{PRÄTERITUM} }\\
    ich & \textbf{durfte}\\
    du & \textbf{durftest}\\
    er/sie/es & \textbf{durfte}\\
    wir & \textbf{durften}\\
    ihr & \textbf{durftet}\\
    sie/Sie & \textbf{durften}\\
  \end{tabular}
   \hfill
     \begin{tabular}{rl}
    \multicolumn{2}{l}{\textbf{KONJUNKTIV II} }\\
    ich & \textbf{dürfte}\\
    du & \textbf{dürftest}\\
    er/sie/es & \textbf{dürfte}\\
    wir & \textbf{dürften}\\
    ihr & \textbf{dürftet}\\
    sie/Sie & \textbf{dürften}\\
  \end{tabular}
  \hfill
  \begin{tabular}{rl}
     \multicolumn{2}{l}{\textbf{IMPERATIV} }\\
   & \textbf{-}\\
   & \textbf{-}\\
   & \textbf{-}\\
   & \textbf{-}\\
   & \vspace{10pt}
  \end{tabular}

    \vspace{10pt}
 
    \begin{tabular}{rl}
     \multicolumn{2}{l}{\textbf{PERFEKT}}\\
   & haben + \textbf{gedurft}\\
  \end{tabular}
\hfill
   \begin{tabular}{rl}
   
    \multicolumn{2}{l}{\textbf{FUTURE I} }\\
    &werden + \textbf{dürfen}\\
 
  \end{tabular}
\hfill
   \begin{tabular}{rl}
      \multicolumn{2}{l}{\textbf{PLUSQUAMPERFEKT}}\\
   & hatten + \textbf{gedurft}\\
  \end{tabular}
  
  \vspace{-5pt}


\end{verbbox-modal}

%können
\begin{verbbox-modal}\addcontentsline{toc}{subsubsection}{können (to be able to)}
    {\Large \bf können} (to be able to) \hfill Type: \textbf{Modal} \hfill Auxiliary Verb: \textbf{haben} 
     \vspace{5pt}
    \hrule
    \vspace{5pt}
  \begin{tabular}{rl}
     \multicolumn{2}{l}{\textbf{PRÄSENS} }\\
    ich & \textbf{kann}\\
    du & \textbf{kannst}\\
    er/sie/es & \textbf{kann}\\
    wir & \textbf{können}\\
    ihr & \textbf{könnt}\\
    sie/Sie & \textbf{können}\\
  \end{tabular}
  \hfill
     \begin{tabular}{rl}
    \multicolumn{2}{l}{\textbf{PRÄTERITUM} }\\
    ich & \textbf{konnte}\\
    du & \textbf{konntest}\\
    er/sie/es & \textbf{konnte}\\
    wir & \textbf{konnten}\\
    ihr & \textbf{konntet}\\
    sie/Sie & \textbf{konnten}\\
  \end{tabular}
   \hfill
     \begin{tabular}{rl}
    \multicolumn{2}{l}{\textbf{KONJUNKTIV II} }\\
    ich & \textbf{könnte}\\
    du & \textbf{könntest}\\
    er/sie/es & \textbf{könnte}\\
    wir & \textbf{könnten}\\
    ihr & \textbf{könntet}\\
    sie/Sie & \textbf{könnten}\\
  \end{tabular}
  \hfill
  \begin{tabular}{rl}
     \multicolumn{2}{l}{\textbf{IMPERATIV} }\\
   & \textbf{-}\\
   & \textbf{-}\\
   & \textbf{-}\\
   & \textbf{-}\\
   & \vspace{10pt}
  \end{tabular}

    \vspace{10pt}
 
    \begin{tabular}{rl}
     \multicolumn{2}{l}{\textbf{PERFEKT}}\\
   & haben + \textbf{gekonnt}\\
  \end{tabular}
\hfill
   \begin{tabular}{rl}
   
    \multicolumn{2}{l}{\textbf{FUTURE I} }\\
 &   werden + \textbf{können}\\
 
  \end{tabular}
\hfill
   \begin{tabular}{rl}
      \multicolumn{2}{l}{\textbf{PLUSQUAMPERFEKT}}\\
  &  hatten + \textbf{gekonnt}\\
  \end{tabular}
  
  \vspace{-5pt}

\end{verbbox-modal}

%mögen
\begin{verbbox-modal}\addcontentsline{toc}{subsubsection}{mögen (to like)}
    {\Large \bf mögen} (to like) \hfill Type: \textbf{Modal} \hfill Auxiliary Verb: \textbf{haben} 
     \vspace{5pt}
    \hrule
    \vspace{5pt}
  \begin{tabular}{rl}
     \multicolumn{2}{l}{\textbf{PRÄSENS} }\\
    ich & \textbf{mag}\\
    du & \textbf{magst}\\
    er/sie/es & \textbf{mag}\\
    wir & \textbf{mögen}\\
    ihr & \textbf{mögt}\\
    sie/Sie & \textbf{mögen}\\
  \end{tabular}
  \hfill
     \begin{tabular}{rl}
    \multicolumn{2}{l}{\textbf{PRÄTERITUM} }\\
    ich & \textbf{mochte}\\
    du & \textbf{mochtest}\\
    er/sie/es & \textbf{mochte}\\
    wir & \textbf{mochten}\\
    ihr & \textbf{mochtet}\\
    sie/Sie & \textbf{mochten}\\
  \end{tabular}
    \hfill
     \begin{tabular}{rl}
    \multicolumn{2}{l}{\textbf{KONJUNKTIV II} }\\
    ich & \textbf{möchte}\\
    du & \textbf{möchtest}\\
    er/sie/es & \textbf{möchte}\\
    wir & \textbf{möchten}\\
    ihr & \textbf{möchtet}\\
    sie/Sie & \textbf{möchten}\\
  \end{tabular}
  \hfill
  \begin{tabular}{rl}
     \multicolumn{2}{l}{\textbf{IMPERATIV} }\\
   & \textbf{-}\\
   & \textbf{-}\\
   & \textbf{-}\\
   & \textbf{-}\\
   & \vspace{10pt}
  \end{tabular}

    \vspace{10pt}
 
    \begin{tabular}{rl}
     \multicolumn{2}{l}{\textbf{PERFEKT}}\\
  &  haben + \textbf{gemocht}\\
  \end{tabular}
\hfill
   \begin{tabular}{rl}
   
    \multicolumn{2}{l}{\textbf{FUTURE I} }\\
  &  werden + \textbf{mögen}\\
 
  \end{tabular}
\hfill
   \begin{tabular}{rl}
      \multicolumn{2}{l}{\textbf{PLUSQUAMPERFEKT}}\\
   & hatten + \textbf{gemocht}\\
  \end{tabular}
  
  \vspace{-5pt}

\end{verbbox-modal}

%müssen
\begin{verbbox-modal}\addcontentsline{toc}{subsubsection}{müssen (to have to)}
    {\Large \bf müssen} (to have to) \hfill Type: \textbf{Modal} \hfill Auxiliary Verb: \textbf{haben} 
     \vspace{5pt}
    \hrule
    \vspace{5pt}
  \begin{tabular}{rl}
     \multicolumn{2}{l}{\textbf{PRÄSENS} }\\
    ich & \textbf{muss}\\
    du & \textbf{musst}\\
    er/sie/es & \textbf{muss}\\
    wir & \textbf{müssen}\\
    ihr & \textbf{müsst}\\
    sie/Sie & \textbf{müssen}\\
  \end{tabular}
  \hfill
     \begin{tabular}{rl}
    \multicolumn{2}{l}{\textbf{PRÄTERITUM} }\\
    ich & \textbf{musste}\\
    du & \textbf{musstest}\\
    er/sie/es & \textbf{musste}\\
    wir & \textbf{mussten}\\
    ihr & \textbf{musstet}\\
    sie/Sie & \textbf{mussten}\\
  \end{tabular}
   \hfill
     \begin{tabular}{rl}
    \multicolumn{2}{l}{\textbf{KONJUNKTIV II} }\\
    ich & \textbf{müsste}\\
    du & \textbf{müsstest}\\
    er/sie/es & \textbf{müsste}\\
    wir & \textbf{müssten}\\
    ihr & \textbf{müsstet}\\
    sie/Sie & \textbf{müssten}\\
  \end{tabular}
  \hfill
  \begin{tabular}{rl}
     \multicolumn{2}{l}{\textbf{IMPERATIV} }\\
   & \textbf{-}\\
   & \textbf{-}\\
   & \textbf{-}\\
   & \textbf{-}\\
   & \vspace{10pt}
  \end{tabular}

    \vspace{10pt}
 
    \begin{tabular}{rl}
     \multicolumn{2}{l}{\textbf{PERFEKT}}\\
  &  haben + \textbf{gemusst}\\
  \end{tabular}
\hfill
   \begin{tabular}{rl}
   
    \multicolumn{2}{l}{\textbf{FUTURE I} }\\
  &  werden + \textbf{müssen}\\
 
  \end{tabular}
\hfill
   \begin{tabular}{rl}
      \multicolumn{2}{l}{\textbf{PLUSQUAMPERFEKT}}\\
  &  hatten + \textbf{gemusst}\\
  \end{tabular}
  
  \vspace{-5pt}
\end{verbbox-modal}

%sollon
\begin{verbbox-modal}\addcontentsline{toc}{subsubsection}{sollen (should)}
    {\Large \bf sollen} (should) \hfill Type: \textbf{Modal} \hfill Auxiliary Verb: \textbf{haben} 
     \vspace{5pt}
    \hrule
    \vspace{5pt}
  \begin{tabular}{rl}
     \multicolumn{2}{l}{\textbf{PRÄSENS} }\\
    ich & \textbf{soll}\\
    du & \textbf{sollst}\\
    er/sie/es & \textbf{soll}\\
    wir & \textbf{sollen}\\
    ihr & \textbf{sollt}\\
    sie/Sie & \textbf{sollen}\\
  \end{tabular}
  \hfill
     \begin{tabular}{rl}
    \multicolumn{2}{l}{\textbf{PRÄTERITUM} }\\
    ich & \textbf{sollte}\\
    du & \textbf{solltest}\\
    er/sie/es & \textbf{sollte}\\
    wir & \textbf{sollten}\\
    ihr & \textbf{solltet}\\
    sie/Sie & \textbf{sollten}\\
  \end{tabular}
   \hfill
     \begin{tabular}{rl}
    \multicolumn{2}{l}{\textbf{KONJUNKTIV II} }\\
    ich & \textbf{sollte}\\
    du & \textbf{solltest}\\
    er/sie/es & \textbf{sollte}\\
    wir & \textbf{sollten}\\
    ihr & \textbf{solltet}\\
    sie/Sie & \textbf{sollten}\\
  \end{tabular}
  \hfill
  \begin{tabular}{rl}
     \multicolumn{2}{l}{\textbf{IMPERATIV} }\\
   & \textbf{soll(e) du!}\\
   & \textbf{sollen wir!}\\
   & \textbf{sollt (ihr)}\\
   & \textbf{sollen Sie!}\\
   & \vspace{10pt}
  \end{tabular}

    \vspace{10pt}
 
    \begin{tabular}{rl}
     \multicolumn{2}{l}{\textbf{PERFEKT}}\\
  &  haben + \textbf{gesollt}\\
  \end{tabular}
\hfill
   \begin{tabular}{rl}
   
    \multicolumn{2}{l}{\textbf{FUTURE I} }\\
  &  werden + \textbf{sollen}\\
 
  \end{tabular}
\hfill
   \begin{tabular}{rl}
      \multicolumn{2}{l}{\textbf{PLUSQUAMPERFEKT}}\\
  &  hatten + \textbf{gesollt}\\
  \end{tabular}
  
  \vspace{-5pt}

\end{verbbox-modal}

%wollen
\begin{verbbox-modal}\addcontentsline{toc}{subsubsection}{wollen (to want)}
    {\Large \bf wollen} (to want) \hfill Type: \textbf{Modal} \hfill Auxiliary Verb: \textbf{haben} 
     \vspace{5pt}
    \hrule
    \vspace{5pt}
  \begin{tabular}{rl}
     \multicolumn{2}{l}{\textbf{PRÄSENS} }\\
    ich & \textbf{will}\\
    du & \textbf{willst}\\
    er/sie/es & \textbf{will}\\
    wir & \textbf{wollen}\\
    ihr & \textbf{wollt}\\
    sie/Sie & \textbf{wollen}\\
  \end{tabular}
  \hfill
     \begin{tabular}{rl}
    \multicolumn{2}{l}{\textbf{PRÄTERITUM} }\\
    ich & \textbf{wollte}\\
    du & \textbf{wolltest}\\
    er/sie/es & \textbf{wollte}\\
    wir & \textbf{wollten}\\
    ihr & \textbf{wolltet}\\
    sie/Sie & \textbf{wollten}\\
  \end{tabular}
   \hfill
     \begin{tabular}{rl}
    \multicolumn{2}{l}{\textbf{KONJUNKTIV II} }\\
    ich & \textbf{wollte}\\
    du & \textbf{wolltest}\\
    er/sie/es & \textbf{wollte}\\
    wir & \textbf{wollten}\\
    ihr & \textbf{wolltet}\\
    sie/Sie & \textbf{wollten}\\
  \end{tabular}
  \hfill
  \begin{tabular}{rl}
     \multicolumn{2}{l}{\textbf{IMPERATIV} }\\
   & \textbf{wolle! (du)}\\
   & \textbf{wollen wir!}\\
   & \textbf{wollt! (ihr)}\\
   & \textbf{wollen Sie!}\\
   & \vspace{10pt}
  \end{tabular}

    \vspace{10pt}
 
    \begin{tabular}{rl}
     \multicolumn{2}{l}{\textbf{PERFEKT}}\\
  &  haben + \textbf{gewollt}\\
  \end{tabular}
\hfill
   \begin{tabular}{rl}
   
    \multicolumn{2}{l}{\textbf{FUTURE I} }\\
  &  werden + \textbf{wollen}\\
 
  \end{tabular}
\hfill
   \begin{tabular}{rl}
      \multicolumn{2}{l}{\textbf{PLUSQUAMPERFEKT}}\\
  &  hatten + \textbf{gewollt}\\
  \end{tabular}
  
  \vspace{-5pt}

\end{verbbox-modal}



\section{Reflexive Verbs}

%sich ändern [reflexive] - to change
\begin{verbbox-acc}\addcontentsline{toc}{subsubsection}{sich ändern (to change)}
        {\Large \bf sich ändern} (to change) \hfill Type: \textbf{regular, + Akkusativ} \hfill Auxiliary Verb: \textbf{haben} 
    
       \vspace{5pt}

    \hrule
    
    \vspace{5pt}

  \begin{tabular}{rl}
     \multicolumn{2}{l}{\textbf{PRÄSENS} }\\
    ich & \textbf{ändere mich}\\
    du & \textbf{änderst dich}\\
    er/sie/es & \textbf{ändert sich}\\
    wir & \textbf{änderen uns}\\
    ihr & \textbf{ändert euch}\\
    sie/Sie & \textbf{änderen sich}\\
  \end{tabular}
  \hfill
     \begin{tabular}{rl}
    \multicolumn{2}{l}{\textbf{PRÄTERITUM} }\\
    ich & \textbf{änderte mich}\\
    du & \textbf{ändertest dich}\\
    er/sie/es & \textbf{änderte sich}\\
    wir & \textbf{änderten uns}\\
    ihr & \textbf{ändertet euch}\\
    sie/Sie & \textbf{änderten sich}\\
  \end{tabular}
  \hfill
  \begin{tabular}{rl}
     \multicolumn{2}{l}{\textbf{IMPERATIV} }\\
   & \textbf{änder(e) (du) dich!}\\
   & \textbf{änderen wir uns!}\\
   & \textbf{ändert (ihr) euch!}\\
   & \textbf{änderen Sie sich!}\\
   & \vspace{10pt}
  \end{tabular}

    \vspace{10pt}

 \begin{tabular}{rl}
     \multicolumn{2}{l}{\textbf{PERFEKT}}\\
    &haben + sich + \textbf{geändert}\\
  \end{tabular}
\hfill
   \begin{tabular}{rl}
   
    \multicolumn{2}{l}{\textbf{FUTURE I} }\\
    &werden + sich + \textbf{änderen}\\
 
  \end{tabular}
\hfill
   \begin{tabular}{rl}
      \multicolumn{2}{l}{\textbf{PLUSQUAMPERFEKT}}\\
    &hatten + sich + \textbf{geändert}\\
  \end{tabular}
  
  
\vspace{10pt}

\begin{tabular}{lll}
\multicolumn{3}{l}{\textbf{MIT PRÄPOSITIONEN}}\\
& ---\\
\end{tabular}
\hfill
\begin{tabular}{lll}
\multicolumn{3}{l}{\textbf{TRENNBARE VERBEN}}\\
& ---\\
\end{tabular}

\vspace{-5pt}

\end{verbbox-acc}

%sich ärgern [reflexive] - to be annoyed
\begin{verbbox-acc}\addcontentsline{toc}{subsubsection}{sich ärgern (to be annoyed)}
        {\Large \bf sich ärgern} (to be annoyed) \hfill Type: \textbf{regular, + Akkusativ} \hfill Auxiliary Verb: \textbf{haben} 
    
       \vspace{5pt}

    \hrule
    
    \vspace{5pt}

  \begin{tabular}{rl}
     \multicolumn{2}{l}{\textbf{PRÄSENS} }\\
    ich & \textbf{ärgere mich}\\
    du & \textbf{ärgerst dich}\\
    er/sie/es & \textbf{ärgert sich}\\
    wir & \textbf{ärgeren uns}\\
    ihr & \textbf{ärgert euch}\\
    sie/Sie & \textbf{ärgeren sich}\\
  \end{tabular}
  \hfill
     \begin{tabular}{rl}
    \multicolumn{2}{l}{\textbf{PRÄTERITUM} }\\
    ich & \textbf{ärgerte mich}\\
    du & \textbf{ärgertest dich}\\
    er/sie/es & \textbf{ärgerte sich}\\
    wir & \textbf{ärgerten uns}\\
    ihr & \textbf{ärgertet euch}\\
    sie/Sie & \textbf{ärgerten sich}\\
  \end{tabular}
  \hfill
  \begin{tabular}{rl}
     \multicolumn{2}{l}{\textbf{IMPERATIV} }\\
   & \textbf{ärger(e) (du) dich!}\\
   & \textbf{ärgeren wir uns!}\\
   & \textbf{ärgert (ihr) euch!}\\
   & \textbf{ärgeren Sie sich!}\\
   & \vspace{10pt}
  \end{tabular}

    \vspace{10pt}

 \begin{tabular}{rl}
     \multicolumn{2}{l}{\textbf{PERFEKT}}\\
    &haben + sich + \textbf{geärgert}\\
  \end{tabular}
\hfill
   \begin{tabular}{rl}
   
    \multicolumn{2}{l}{\textbf{FUTURE I} }\\
    &werden + sich + \textbf{ärgeren}\\
 
  \end{tabular}
\hfill
   \begin{tabular}{rl}
      \multicolumn{2}{l}{\textbf{PLUSQUAMPERFEKT}}\\
    &hatten + sich + \textbf{geärgert}\\
  \end{tabular}
  
\vspace{10pt}

   \begin{tabular}{lll}
      \multicolumn{3}{l}{\textbf{MIT PRÄPOSITIONEN}}\\
& \textbf{sich ärgeren über} + Akk & (to be annoyed about)
  \end{tabular}
 % \hfill
 %  \begin{tabular}{lll}
  %    \multicolumn{3}{l}{\textbf{TRENNBARE VERBEN}}\\
%\vspace{10pt}
 % \end{tabular}
  
  \vspace{-5pt}

\end{verbbox-acc}

%sich bedanken [reflexive] - to thank
\begin{verbbox}\addcontentsline{toc}{subsubsection}{sich bedanken (to thank)}
        {\Large \bf sich bedanken} (to thank) \hfill Type: \textbf{regular} \hfill Auxiliary Verb: \textbf{haben} 
    
       \vspace{5pt}

    \hrule
    
    \vspace{5pt}

  \begin{tabular}{rl}
     \multicolumn{2}{l}{\textbf{PRÄSENS} }\\
    ich & \textbf{bedanke mich}\\
    du & \textbf{bedankst dich}\\
    er/sie/es & \textbf{bedankt sich}\\
    wir & \textbf{bedanken uns}\\
    ihr & \textbf{bedankt euch}\\
    sie/Sie & \textbf{bedanken sich}\\
  \end{tabular}
  \hfill
     \begin{tabular}{rl}
    \multicolumn{2}{l}{\textbf{PRÄTERITUM} }\\
    ich & \textbf{bedankte mich}\\
    du & \textbf{bedanktest dich}\\
    er/sie/es & \textbf{vbedankte sich}\\
    wir & \textbf{bedankten uns}\\
    ihr & \textbf{bedanktet euch}\\
    sie/Sie & \textbf{bedankten sich}\\
  \end{tabular}
  \hfill
  \begin{tabular}{rl}
     \multicolumn{2}{l}{\textbf{IMPERATIV} }\\
   & \textbf{bedank(e) (du) dich!}\\
   & \textbf{bedanken wir uns!}\\
   & \textbf{bedankt (ihr) euch!}\\
   & \textbf{bedanken Sie sich!}\\
   & \vspace{10pt}
  \end{tabular}

    \vspace{10pt}

 \begin{tabular}{rl}
     \multicolumn{2}{l}{\textbf{PERFEKT}}\\
    &haben + sich + \textbf{bedankt}\\
  \end{tabular}
\hfill
   \begin{tabular}{rl}
   
    \multicolumn{2}{l}{\textbf{FUTURE I} }\\
    &werden + sich + \textbf{bedanken}\\
 
  \end{tabular}
\hfill
   \begin{tabular}{rl}
      \multicolumn{2}{l}{\textbf{PLUSQUAMPERFEKT}}\\
    &hatten + sich + \textbf{bedankt}\\
  \end{tabular}
  
\vspace{10pt}

   \begin{tabular}{lll}
      \multicolumn{3}{l}{\textbf{MIT PRÄPOSITIONEN}}\\
& \textbf{sich bedanken für} + Akk & (to thank for)\\
& \textbf{sich bedanken bei} + Dat & (to be thankful to someone)
  \end{tabular}
 \hfill
  \begin{tabular}{lll}
     \multicolumn{3}{l}{\textbf{TRENNBARE VERBEN}}\\
     & --- & \\
\vspace{10pt}
  \end{tabular}
  
  \vspace{-5pt}

\end{verbbox}

%sich beeilen [reflexive] - to hurry
\begin{verbbox}\addcontentsline{toc}{subsubsection}{sich beeilen (to hurry)}
        {\Large \bf sich beeilen} (to hurry) \hfill Type: \textbf{regular} \hfill Auxiliary Verb: \textbf{haben} 
    
       \vspace{5pt}

    \hrule
    
    \vspace{5pt}

  \begin{tabular}{rl}
     \multicolumn{2}{l}{\textbf{PRÄSENS} }\\
    ich & \textbf{beeile mich}\\
    du & \textbf{beeilst dich}\\
    er/sie/es & \textbf{beeilt sich}\\
    wir & \textbf{beeilen uns}\\
    ihr & \textbf{beeilt euch}\\
    sie/Sie & \textbf{beeilen sich}\\
  \end{tabular}
  \hfill
     \begin{tabular}{rl}
    \multicolumn{2}{l}{\textbf{PRÄTERITUM} }\\
    ich & \textbf{beeilte mich}\\
    du & \textbf{beeiltest dich}\\
    er/sie/es & \textbf{beeilte sich}\\
    wir & \textbf{beeilten uns}\\
    ihr & \textbf{beeiltet euch}\\
    sie/Sie & \textbf{beeilten sich}\\
  \end{tabular}
  \hfill
  \begin{tabular}{rl}
     \multicolumn{2}{l}{\textbf{IMPERATIV} }\\
   & \textbf{beeil(e) (du) dich!}\\
   & \textbf{beeilen wir uns!}\\
   & \textbf{beeilt (ihr) euch!}\\
   & \textbf{beeilen Sie sich!}\\
   & \vspace{10pt}
  \end{tabular}

    \vspace{10pt}

 \begin{tabular}{rl}
     \multicolumn{2}{l}{\textbf{PERFEKT}}\\
    &haben + sich + \textbf{beeilt}\\
  \end{tabular}
\hfill
   \begin{tabular}{rl}
   
    \multicolumn{2}{l}{\textbf{FUTURE I} }\\
    &werden + sich + \textbf{beeilen}\\
 
  \end{tabular}
\hfill
   \begin{tabular}{rl}
      \multicolumn{2}{l}{\textbf{PLUSQUAMPERFEKT}}\\
    &hatten + sich + \textbf{beeilt}\\
  \end{tabular}
  
\vspace{10pt}

   \begin{tabular}{lll}
      \multicolumn{3}{l}{\textbf{MIT PRÄPOSITIONEN}}\\
& \textbf{sich beeilen mit} + Dat & (to hurry up with)
  \end{tabular}
 % \hfill
 %  \begin{tabular}{lll}
  %    \multicolumn{3}{l}{\textbf{TRENNBARE VERBEN}}\\
%\vspace{10pt}
 % \end{tabular}
  
  \vspace{-5pt}

\end{verbbox}

%sich befinden [reflexive] - to be located
\begin{verbbox}\addcontentsline{toc}{subsubsection}{sich befinden (to be located)}
        {\Large \bf sich befinden} (to be located) \hfill Type: \textbf{irregular} \hfill Auxiliary Verb: \textbf{haben} 
    
       \vspace{5pt}

    \hrule
    
    \vspace{5pt}

  \begin{tabular}{rl}
     \multicolumn{2}{l}{\textbf{PRÄSENS} }\\
    ich & \textbf{befinde mich}\\
    du & \textbf{befindest dich}\\
    er/sie/es & \textbf{befindet sich}\\
    wir & \textbf{befinden uns}\\
    ihr & \textbf{befindet euch}\\
    sie/Sie & \textbf{befinden sich}\\
  \end{tabular}
  \hfill
     \begin{tabular}{rl}
    \multicolumn{2}{l}{\textbf{PRÄTERITUM} }\\
    ich & \textbf{befand mich}\\
    du & \textbf{befandest dich}\\
    er/sie/es & \textbf{befand sich}\\
    wir & \textbf{befanden uns}\\
    ihr & \textbf{befandet euch}\\
    sie/Sie & \textbf{befanden sich}\\
  \end{tabular}
  \hfill
  \begin{tabular}{rl}
     \multicolumn{2}{l}{\textbf{IMPERATIV} }\\
   & \textbf{befind(e) (du) dich!}\\
   & \textbf{befinden wir uns!}\\
   & \textbf{befindet (ihr) euch!}\\
   & \textbf{befinden Sie sich!}\\
   & \vspace{10pt}
  \end{tabular}

    \vspace{10pt}

 \begin{tabular}{rl}
     \multicolumn{2}{l}{\textbf{PERFEKT}}\\
    &haben + sich + \textbf{befunden}\\
  \end{tabular}
\hfill
   \begin{tabular}{rl}
   
    \multicolumn{2}{l}{\textbf{FUTURE I} }\\
    &werden + sich + \textbf{befinden}\\
 
  \end{tabular}
\hfill
   \begin{tabular}{rl}
      \multicolumn{2}{l}{\textbf{PLUSQUAMPERFEKT}}\\
    &hatten + sich + \textbf{befunden}\\
  \end{tabular}
  
\vspace{10pt}

   \begin{tabular}{lll}
      \multicolumn{3}{l}{\textbf{MIT PRÄPOSITIONEN}}\\
& \textbf{sich befinden in} + Akk & (to be located in)
  \end{tabular}
  \hfill
   \begin{tabular}{lll}
      \multicolumn{3}{l}{\textbf{TRENNBARE VERBEN}}\\
& --- & \\
  \end{tabular}
  
  \vspace{-5pt}

\end{verbbox}

%sich bemühen - (to make an effort)
\begin{verbbox}\addcontentsline{toc}{subsubsection}{sich bemühen (to make an effort)}
    {\Large \bf sich bemühen} (to make an effort) \hfill Type: \textbf{regular, reflexive} \hfill Auxiliary Verb: \textbf{haben} 
    
     \vspace{5pt}

    \hrule
    
    \vspace{5pt}

  \begin{tabular}{rl}
     \multicolumn{2}{l}{\textbf{PRÄSENS} }\\
    ich & \textbf{bemühe mich}\\
    du & \textbf{bemühst dich}\\
    er/sie/es & \textbf{bemüht sich}\\
    wir & \textbf{bemühen uns}\\
    ihr & \textbf{bemüht euch}\\
    sie/Sie & \textbf{bemühen sich}\\
  \end{tabular}
  \hfill
     \begin{tabular}{rl}
    \multicolumn{2}{l}{\textbf{PRÄTERITUM} }\\
    ich & \textbf{bemühte mich}\\
    du & \textbf{bemühtest dich}\\
    er/sie/es & \textbf{bemühte sich}\\
    wir & \textbf{bemühten uns}\\
    ihr & \textbf{bemühtet euch}\\
    sie/Sie & \textbf{bemühten sich}\\
  \end{tabular}
  \hfill
  \begin{tabular}{rl}
     \multicolumn{2}{l}{\textbf{IMPERATIV} }\\
   & \textbf{bemüh(e) dich!}\\
   & \textbf{bemühen wir uns!}\\
   & \textbf{bemüht euch!}\\
   & \textbf{bemühen Sie sich!}\\
   & \vspace{10pt}
  \end{tabular}

    \vspace{10pt}

 \begin{tabular}{rl}
     \multicolumn{2}{l}{\textbf{PERFEKT}}\\
   & haben + sich + \textbf{bemüht}\\
  \end{tabular}
\hfill
   \begin{tabular}{rl}
   
    \multicolumn{2}{l}{\textbf{FUTURE I} }\\
   & werden + sich + \textbf{bemühen}\\
 
  \end{tabular}
\hfill
   \begin{tabular}{rl}
      \multicolumn{2}{l}{\textbf{PLUSQUAMPERFEKT}}\\
   & hatten + sich + \textbf{bemüht}\\
  \end{tabular}

   \vspace{10pt}

   \begin{tabular}{lll}
      \multicolumn{3}{l}{\textbf{MIT PRÄPOSITIONEN}}\\
& \textbf{sich bemühen um} + Akk & (to make an effort for / to seek)\\
& \textbf{sich bemühen zu} + Inf & (to try to)\\
\end{tabular}
  \hfill
   \begin{tabular}{lll}
      \multicolumn{3}{l}{\textbf{TRENNBARE VERBEN}}\\
 & \textbf{---} & (none common)\\
\vspace{10pt}
  \end{tabular}


  \vspace{-5pt}


\end{verbbox}

%sich beschäftigen (to occupy oneself / to deal with)
\begin{verbbox}\addcontentsline{toc}{subsubsection}{sich beschäftigen (to occupy oneself / to deal with)}
    {\Large \bf sich beschäftigen} (to occupy oneself / to deal with) \hfill Type: \textbf{regular} \hfill Auxiliary Verb: \textbf{haben} 
    
     \vspace{5pt}

    \hrule
    
    \vspace{5pt}

  \begin{tabular}{rl}
     \multicolumn{2}{l}{\textbf{PRÄSENS} }\\
    ich & \textbf{beschäftige mich}\\
    du & \textbf{beschäftigst dich}\\
    er/sie/es & \textbf{beschäftigt sich}\\
    wir & \textbf{beschäftigen uns}\\
    ihr & \textbf{beschäftigt euch}\\
    sie/Sie & \textbf{beschäftigen sich}\\
  \end{tabular}
  \hfill
     \begin{tabular}{rl}
    \multicolumn{2}{l}{\textbf{PRÄTERITUM} }\\
    ich & \textbf{beschäftigte mich}\\
    du & \textbf{beschäftigtest dich}\\
    er/sie/es & \textbf{beschäftigte sich}\\
    wir & \textbf{beschäftigten uns}\\
    ihr & \textbf{beschäftigtet euch}\\
    sie/Sie & \textbf{beschäftigten sich}\\
  \end{tabular}
  \hfill
  \begin{tabular}{rl}
     \multicolumn{2}{l}{\textbf{IMPERATIV} }\\
   & \textbf{beschäftige dich!}\\
   & \textbf{beschäftigen wir uns!}\\
   & \textbf{beschäftigt euch!}\\
   & \textbf{beschäftigen Sie sich!}\\
   & \vspace{10pt}
  \end{tabular}

    \vspace{10pt}

 \begin{tabular}{rl}
     \multicolumn{2}{l}{\textbf{PERFEKT}}\\
   & haben + sich + \textbf{beschäftigt}\\
  \end{tabular}
\hfill
   \begin{tabular}{rl}
    \multicolumn{2}{l}{\textbf{FUTURE I} }\\
   & werden + sich + \textbf{beschäftigen}\\
  \end{tabular}
\hfill
   \begin{tabular}{rl}
      \multicolumn{2}{l}{\textbf{PLUSQUAMPERFEKT}}\\
   & hatten + sich + \textbf{beschäftigt}\\
  \end{tabular}

   \vspace{10pt}

   \begin{tabular}{lll}
      \multicolumn{3}{l}{\textbf{MIT PRÄPOSITIONEN}}\\
& \textbf{sich beschäftigen mit} + Dat & (to deal with / be occupied with)\\
\end{tabular}
  \hfill
   \begin{tabular}{lll}
      \multicolumn{3}{l}{\textbf{TRENNBARE VERBEN}}\\
 & \textbf{---} & (none common)\\
\vspace{10pt}
  \end{tabular}

  \vspace{-5pt}

\end{verbbox}


%sich beschweren [reflexive] - to complain
\begin{verbbox}\addcontentsline{toc}{subsubsection}{sich beschweren (to complain)}
        {\Large \bf sich beschweren} (to complain) \hfill Type: \textbf{regular} \hfill Auxiliary Verb: \textbf{haben} 
    
       \vspace{5pt}

    \hrule
    
    \vspace{5pt}

  \begin{tabular}{rl}
     \multicolumn{2}{l}{\textbf{PRÄSENS} }\\
    ich & \textbf{beschwere mich}\\
    du & \textbf{beschwerst dich}\\
    er/sie/es & \textbf{beschwert sich}\\
    wir & \textbf{beschweren uns}\\
    ihr & \textbf{beschwert euch}\\
    sie/Sie & \textbf{beschweren sich}\\
  \end{tabular}
  \hfill
     \begin{tabular}{rl}
    \multicolumn{2}{l}{\textbf{PRÄTERITUM} }\\
    ich & \textbf{beschwerte mich}\\
    du & \textbf{beschwertest dich}\\
    er/sie/es & \textbf{beschwerte sich}\\
    wir & \textbf{beschwerten uns}\\
    ihr & \textbf{beschwertet euch}\\
    sie/Sie & \textbf{beschwerten sich}\\
  \end{tabular}
  \hfill
  \begin{tabular}{rl}
     \multicolumn{2}{l}{\textbf{IMPERATIV} }\\
   & \textbf{beschwer(e) (du) dich!}\\
   & \textbf{beschweren wir uns!}\\
   & \textbf{beschwert (ihr) euch!}\\
   & \textbf{beschweren Sie sich!}\\
   & \vspace{10pt}
  \end{tabular}

    \vspace{10pt}

 \begin{tabular}{rl}
     \multicolumn{2}{l}{\textbf{PERFEKT}}\\
    &haben + sich + \textbf{beschwert}\\
  \end{tabular}
\hfill
   \begin{tabular}{rl}
   
    \multicolumn{2}{l}{\textbf{FUTURE I} }\\
    &werden + sich + \textbf{beschweren}\\
 
  \end{tabular}
\hfill
   \begin{tabular}{rl}
      \multicolumn{2}{l}{\textbf{PLUSQUAMPERFEKT}}\\
    &hatten + sich + \textbf{beschwert}\\
  \end{tabular}
  
\vspace{10pt}

   \begin{tabular}{lll}
      \multicolumn{3}{l}{\textbf{MIT PRÄPOSITIONEN}}\\
& \textbf{sich beschweren über} + Akk & (to complain about something)
  \end{tabular}
  \hfill
   \begin{tabular}{lll}
      \multicolumn{3}{l}{\textbf{TRENNBARE VERBEN}}\\
& --- & \\
  \end{tabular}
  
  \vspace{-5pt}

\end{verbbox}

%sich bewegan - (to move oneself)
\begin{verbbox}\addcontentsline{toc}{subsubsection}{sich bewegen (to move oneself)}
    {\Large \bf sich bewegen} (to move oneself) \hfill Type: \textbf{irregular} \hfill Auxiliary Verb: \textbf{haben / sein} 
    
     \vspace{5pt}
    \hrule
    \vspace{5pt}

  \begin{tabular}{rl}
     \multicolumn{2}{l}{\textbf{PRÄSENS} }\\
    ich & \textbf{bewege mich}\\
    du & \textbf{bewegst dich}\\
    er/sie/es & \textbf{bewegt sich}\\
    wir & \textbf{bewegen uns}\\
    ihr & \textbf{bewegt euch}\\
    sie/Sie & \textbf{bewegen sich}\\
  \end{tabular}
  \hfill
     \begin{tabular}{rl}
    \multicolumn{2}{l}{\textbf{PRÄTERITUM} }\\
    ich & \textbf{bewegte mich}\\
    du & \textbf{bewegtest dich}\\
    er/sie/es & \textbf{bewegte sich}\\
    wir & \textbf{bewegten uns}\\
    ihr & \textbf{bewegtet euch}\\
    sie/Sie & \textbf{bewegten sich}\\
  \end{tabular}
  \hfill
  \begin{tabular}{rl}
     \multicolumn{2}{l}{\textbf{IMPERATIV} }\\
   & \textbf{beweg dich (du)!}\\
   & \textbf{bewegen wir uns!}\\
   & \textbf{bewegt euch (ihr)!}\\
   & \textbf{bewegen Sie sich!}\\
   & \vspace{10pt}
  \end{tabular}

 \vspace{10pt}

 \begin{tabular}{rl}
     \multicolumn{2}{l}{\textbf{PERFEKT}}\\
   & haben / sein + sich + \textbf{bewegt}
  \end{tabular}
\hfill
   \begin{tabular}{rl}
    \multicolumn{2}{l}{\textbf{FUTURE I}}\\
   & werden + sich + \textbf{bewegen}\\
  \end{tabular}
\hfill
   \begin{tabular}{rl}
      \multicolumn{2}{l}{\textbf{PLUSQUAMPERFEKT}}\\
   & haben / waren + sich + \textbf{bewegt}\\
  \end{tabular}

   \vspace{10pt}

   \begin{tabular}{lll}
      \multicolumn{3}{l}{\textbf{MIT PRÄPOSITIONEN}}\\
& \textbf{sich bewegen in} + Dat & (to move in)\\
& \textbf{sich bewegen auf} + Akk & (to move toward)\\
& \textbf{sich bewegen durch} + Akk & (to move through)\\
\end{tabular}
  \hfill
\begin{tabular}{lll}
\multicolumn{3}{l}{\textbf{TRENNBARE VERBEN}}\\
& ---\\
\vspace{20pt}
\end{tabular}

  \vspace{-5pt}
\end{verbbox}

%sich drehen (to turn / to rotate)
\begin{verbbox}\addcontentsline{toc}{subsubsection}{sich drehen (to turn / to rotate)}
    {\Large \bf sich drehen} (to turn / to rotate) \hfill Type: \textbf{regular} \hfill Auxiliary Verb: \textbf{haben} 
    
     \vspace{5pt}

    \hrule
    
    \vspace{5pt}

  \begin{tabular}{rl}
     \multicolumn{2}{l}{\textbf{PRÄSENS} }\\
    ich & \textbf{drehe mich}\\
    du & \textbf{drehst dich}\\
    er/sie/es & \textbf{dreht sich}\\
    wir & \textbf{drehen uns}\\
    ihr & \textbf{dreht euch}\\
    sie/Sie & \textbf{drehen sich}\\
  \end{tabular}
  \hfill
     \begin{tabular}{rl}
    \multicolumn{2}{l}{\textbf{PRÄTERITUM} }\\
    ich & \textbf{drehte mich}\\
    du & \textbf{drehtest dich}\\
    er/sie/es & \textbf{drehte sich}\\
    wir & \textbf{drehten uns}\\
    ihr & \textbf{drehtet euch}\\
    sie/Sie & \textbf{drehten sich}\\
  \end{tabular}
  \hfill
  \begin{tabular}{rl}
     \multicolumn{2}{l}{\textbf{IMPERATIV} }\\
   & \textbf{dreh(e) dich!}\\
   & \textbf{drehen wir uns!}\\
   & \textbf{dreht euch!}\\
   & \textbf{drehen Sie sich!}\\
   & \vspace{10pt}
  \end{tabular}

    \vspace{10pt}

 \begin{tabular}{rl}
     \multicolumn{2}{l}{\textbf{PERFEKT}}\\
   & haben + sich + \textbf{gedreht}\\
  \end{tabular}
\hfill
   \begin{tabular}{rl}
    \multicolumn{2}{l}{\textbf{FUTURE I} }\\
   & werden + sich + \textbf{drehen}\\
  \end{tabular}
\hfill
   \begin{tabular}{rl}
      \multicolumn{2}{l}{\textbf{PLUSQUAMPERFEKT}}\\
   & hatten + sich + \textbf{gedreht}\\
  \end{tabular}

   \vspace{10pt}

   \begin{tabular}{lll}
      \multicolumn{3}{l}{\textbf{MIT PRÄPOSITIONEN}}\\
& \textbf{sich drehen um} + Akk & (to revolve around / be about)\\
& \textbf{sich drehen nach} + Dat & (to turn toward)\\
\end{tabular}
  \hfill
   \begin{tabular}{lll}
      \multicolumn{3}{l}{\textbf{TRENNBARE VERBEN}}\\
 & \textbf{um$|$drehen} & (to turn around)\\
 & \textbf{ab$|$drehen} & (to turn off / shoot a film)\\
\vspace{10pt}
  \end{tabular}

  \vspace{-5pt}

\end{verbbox}

%sich eignen (to be suitable)
\begin{verbbox}\addcontentsline{toc}{subsubsection}{sich eignen (to be suitable)}
    {\Large \bf sich eignen} (to be suitable) \hfill Type: \textbf{regular} \hfill Auxiliary Verb: \textbf{haben} 
    
     \vspace{5pt}

    \hrule
    
    \vspace{5pt}

  \begin{tabular}{rl}
     \multicolumn{2}{l}{\textbf{PRÄSENS} }\\
    ich & \textbf{eigne mich}\\
    du & \textbf{eignest dich}\\
    er/sie/es & \textbf{eignet sich}\\
    wir & \textbf{eignen uns}\\
    ihr & \textbf{eignet euch}\\
    sie/Sie & \textbf{eignen sich}\\
  \end{tabular}
  \hfill
     \begin{tabular}{rl}
    \multicolumn{2}{l}{\textbf{PRÄTERITUM} }\\
    ich & \textbf{eignete mich}\\
    du & \textbf{eignetest dich}\\
    er/sie/es & \textbf{eignete sich}\\
    wir & \textbf{eigneten uns}\\
    ihr & \textbf{eignetet euch}\\
    sie/Sie & \textbf{eigneten sich}\\
  \end{tabular}
  \hfill
  \begin{tabular}{rl}
     \multicolumn{2}{l}{\textbf{IMPERATIV} }\\
   & \textbf{eigne dich!}\\
   & \textbf{eignen wir uns!}\\
   & \textbf{eignet euch!}\\
   & \textbf{eignen Sie sich!}\\
   & \vspace{10pt}
  \end{tabular}

    \vspace{10pt}

 \begin{tabular}{rl}
     \multicolumn{2}{l}{\textbf{PERFEKT}}\\
   & haben + sich + \textbf{geeignet}\\
  \end{tabular}
\hfill
   \begin{tabular}{rl}
    \multicolumn{2}{l}{\textbf{FUTURE I} }\\
   & werden + sich + \textbf{eignen}\\
  \end{tabular}
\hfill
   \begin{tabular}{rl}
      \multicolumn{2}{l}{\textbf{PLUSQUAMPERFEKT}}\\
   & hatten + sich + \textbf{geeignet}\\
  \end{tabular}

   \vspace{10pt}

   \begin{tabular}{lll}
      \multicolumn{3}{l}{\textbf{MIT PRÄPOSITIONEN}}\\
& \textbf{sich eignen für} + Akk & (to be suitable for)\\
& \textbf{sich eignen als} + Nom & (to be suitable as)\\
\end{tabular}
  \hfill
   \begin{tabular}{lll}
      \multicolumn{3}{l}{\textbf{TRENNBARE VERBEN}}\\
 & \textbf{---} & (none common)\\
\vspace{10pt}
  \end{tabular}

  \vspace{-5pt}

\end{verbbox}

%sich erinnern [reflexive] - to remenber
\begin{verbbox}\addcontentsline{toc}{subsubsection}{sich erinnern (to remember)}
        {\Large \bf sich erinnern} (to remember) \hfill Type: \textbf{regular} \hfill Auxiliary Verb: \textbf{haben} 
    
       \vspace{5pt}

    \hrule
    
    \vspace{5pt}

  \begin{tabular}{rl}
     \multicolumn{2}{l}{\textbf{PRÄSENS} }\\
    ich & \textbf{erinnere mich}\\
    du & \textbf{erinnerst dich}\\
    er/sie/es & \textbf{erinnert sich}\\
    wir & \textbf{erinnern uns}\\
    ihr & \textbf{erinnert euch}\\
    sie/Sie & \textbf{erinnern sich}\\
  \end{tabular}
  \hfill
     \begin{tabular}{rl}
    \multicolumn{2}{l}{\textbf{PRÄTERITUM} }\\
    ich & \textbf{erinnerte mich}\\
    du & \textbf{erinnertest dich}\\
    er/sie/es & \textbf{erinnerte sich}\\
    wir & \textbf{erinnerten uns}\\
    ihr & \textbf{erinnertet euch}\\
    sie/Sie & \textbf{erinnerten sich}\\
  \end{tabular}
  \hfill
  \begin{tabular}{rl}
     \multicolumn{2}{l}{\textbf{IMPERATIV} }\\
   & \textbf{erinner(e) (du) dich!}\\
   & \textbf{erinnern wir uns!}\\
   & \textbf{erinnert (ihr) euch!}\\
   & \textbf{erinnern Sie sich!}\\
   & \vspace{10pt}
  \end{tabular}

    \vspace{10pt}

 \begin{tabular}{rl}
     \multicolumn{2}{l}{\textbf{PERFEKT}}\\
    &haben + sich + \textbf{erinnert}\\
  \end{tabular}
\hfill
   \begin{tabular}{rl}
   
    \multicolumn{2}{l}{\textbf{FUTURE I} }\\
    &werden + sich + \textbf{erinnern}\\
 
  \end{tabular}
\hfill
   \begin{tabular}{rl}
      \multicolumn{2}{l}{\textbf{PLUSQUAMPERFEKT}}\\
    &hatten + sich + \textbf{erinnert}\\
  \end{tabular}
  
\vspace{10pt}

   \begin{tabular}{lll}
      \multicolumn{3}{l}{\textbf{MIT PRÄPOSITIONEN}}\\
& \textbf{sich erinnern an} + Akk & (to recall something)
  \end{tabular}
  \hfill
   \begin{tabular}{lll}
      \multicolumn{3}{l}{\textbf{TRENNBARE VERBEN}}\\
& --- & \\
 \end{tabular}
  
  \vspace{-5pt}

\end{verbbox}

%sich erkundigen (to inquire)
\begin{verbbox}\addcontentsline{toc}{subsubsection}{sich erkundigen (to inquire)}
    {\Large \bf sich erkundigen} (to inquire) \hfill Type: \textbf{regular} \hfill Auxiliary Verb: \textbf{haben} 
    
     \vspace{5pt}

    \hrule
    
    \vspace{5pt}

  \begin{tabular}{rl}
     \multicolumn{2}{l}{\textbf{PRÄSENS} }\\
    ich & \textbf{erkundige mich}\\
    du & \textbf{erkundigst dich}\\
    er/sie/es & \textbf{erkundigt sich}\\
    wir & \textbf{erkundigen uns}\\
    ihr & \textbf{erkundigt euch}\\
    sie/Sie & \textbf{erkundigen sich}\\
  \end{tabular}
  \hfill
     \begin{tabular}{rl}
    \multicolumn{2}{l}{\textbf{PRÄTERITUM} }\\
    ich & \textbf{erkundigte mich}\\
    du & \textbf{erkundigtest dich}\\
    er/sie/es & \textbf{erkundigte sich}\\
    wir & \textbf{erkundigten uns}\\
    ihr & \textbf{erkundigtet euch}\\
    sie/Sie & \textbf{erkundigten sich}\\
  \end{tabular}
  \hfill
  \begin{tabular}{rl}
     \multicolumn{2}{l}{\textbf{IMPERATIV} }\\
   & \textbf{erkundige dich!}\\
   & \textbf{erkundigen wir uns!}\\
   & \textbf{erkundigt euch!}\\
   & \textbf{erkundigen Sie sich!}\\
   & \vspace{10pt}
  \end{tabular}

    \vspace{10pt}

 \begin{tabular}{rl}
     \multicolumn{2}{l}{\textbf{PERFEKT}}\\
   & haben + sich + \textbf{erkundigt}\\
  \end{tabular}
\hfill
   \begin{tabular}{rl}
    \multicolumn{2}{l}{\textbf{FUTURE I} }\\
   & werden + sich + \textbf{erkundigen}\\
  \end{tabular}
\hfill
   \begin{tabular}{rl}
      \multicolumn{2}{l}{\textbf{PLUSQUAMPERFEKT}}\\
   & hatten + sich + \textbf{erkundigt}\\
  \end{tabular}

   \vspace{10pt}

   \begin{tabular}{lll}
      \multicolumn{3}{l}{\textbf{MIT PRÄPOSITIONEN}}\\
& \textbf{sich erkundigen nach} + Dat & (to inquire about)\\
& \textbf{sich erkundigen bei} + Dat & (to inquire with / ask at)\\
\end{tabular}
  \hfill
   \begin{tabular}{lll}
      \multicolumn{3}{l}{\textbf{TRENNBARE VERBEN}}\\
 & \textbf{---} & (none common)\\
\vspace{10pt}
  \end{tabular}

  \vspace{-5pt}

\end{verbbox}


%sich ernähren - to eat / to nourish oneself
\begin{verbbox}\addcontentsline{toc}{subsubsection}{sich ernähren (to eat / to nourish oneself)}
    {\Large \bf sich ernähren} (to eat / nourish oneself) \hfill Type: \textbf{regular} \hfill Auxiliary Verb: \textbf{haben} 
    
     \vspace{5pt}
    \hrule
    \vspace{5pt}

  \begin{tabular}{rl}
     \multicolumn{2}{l}{\textbf{PRÄSENS} }\\
    ich & \textbf{ernähre mich}\\
    du & \textbf{ernährst dich}\\
    er/sie/es & \textbf{ernährt sich}\\
    wir & \textbf{ernähren uns}\\
    ihr & \textbf{ernährt euch}\\
    sie/Sie & \textbf{ernähren sich}\\
  \end{tabular}
  \hfill
     \begin{tabular}{rl}
    \multicolumn{2}{l}{\textbf{PRÄTERITUM} }\\
    ich & \textbf{ernährte mich}\\
    du & \textbf{ernährtest dich}\\
    er/sie/es & \textbf{ernährte sich}\\
    wir & \textbf{ernährten uns}\\
    ihr & \textbf{ernährtet euch}\\
    sie/Sie & \textbf{ernährten sich}\\
  \end{tabular}
  \hfill
  \begin{tabular}{rl}
     \multicolumn{2}{l}{\textbf{IMPERATIV} }\\
   & \textbf{ernähr dich (du)!}\\
   & \textbf{ernähren wir uns!}\\
   & \textbf{ernährt euch (ihr)!}\\
   & \textbf{ernähren Sie sich!}\\
   & \vspace{10pt}
  \end{tabular}

 \vspace{10pt}

 \begin{tabular}{rl}
     \multicolumn{2}{l}{\textbf{PERFEKT}}\\
   & haben + sich + \textbf{ernährt}\\
  \end{tabular}
\hfill
   \begin{tabular}{rl}
    \multicolumn{2}{l}{\textbf{FUTURE I}}\\
   & werden + sich + \textbf{ernähren}\\
  \end{tabular}
\hfill
   \begin{tabular}{rl}
      \multicolumn{2}{l}{\textbf{PLUSQUAMPERFEKT}}\\
   & hatten + sich + \textbf{ernährt}\\
  \end{tabular}

   \vspace{10pt}

   \begin{tabular}{lll}
      \multicolumn{3}{l}{\textbf{MIT PRÄPOSITIONEN}}\\
& \textbf{sich ernähren von} + Dat & (to live on / eat mainly)\\
& \textbf{sich ernähren mit} + Dat & (to nourish oneself with)\\
\end{tabular}
  \hfill
\begin{tabular}{lll}
\multicolumn{3}{l}{\textbf{TRENNBARE VERBEN}}\\
& ---\\
\vspace{10pt}
\end{tabular}

\end{verbbox}

%sich engagieren (to get involved / commit oneself)
\begin{verbbox}\addcontentsline{toc}{subsubsection}{sich engagieren (to get involved / commit oneself)}
{\Large \bf sich engagieren} (to get involved / commit oneself) \hfill Type: \textbf{regular} \hfill Auxiliary Verb: \textbf{haben}

\vspace{5pt}
\hrule
\vspace{5pt}

\begin{tabular}{rl}
\multicolumn{2}{l}{\textbf{PRÄSENS}}\\
ich & \textbf{engagiere mich}\\
du & \textbf{engagierst dich}\\
er/sie/es & \textbf{engagiert sich}\\
wir & \textbf{engagieren uns}\\
ihr & \textbf{engagiert euch}\\
sie/Sie & \textbf{engagieren sich}\\
\end{tabular}
\hfill
\begin{tabular}{rl}
\multicolumn{2}{l}{\textbf{PRÄTERITUM}}\\
ich & \textbf{engagierte mich}\\
du & \textbf{engagiertest dich}\\
er/sie/es & \textbf{engagierte sich}\\
wir & \textbf{engagierten uns}\\
ihr & \textbf{engagiertet euch}\\
sie/Sie & \textbf{engagierten sich}\\
\end{tabular}
\hfill
\begin{tabular}{rl}
\multicolumn{2}{l}{\textbf{IMPERATIV}}\\
& \textbf{engagiere dich!}\\
& \textbf{engagieren wir uns!}\\
& \textbf{engagiert euch!}\\
& \textbf{engagieren Sie sich!}\\
\end{tabular}

\vspace{10pt}

\begin{tabular}{rl}
\multicolumn{2}{l}{\textbf{PERFEKT}}\\
& haben + sich + \textbf{engagiert}\\
\end{tabular}
\hfill
\begin{tabular}{rl}
\multicolumn{2}{l}{\textbf{FUTURE I}}\\
& werden + sich + \textbf{engagieren}\\
\end{tabular}
\hfill
\begin{tabular}{rl}
\multicolumn{2}{l}{\textbf{PLUSQUAMPERFEKT}}\\
& hatten + sich + \textbf{engagiert}\\
\end{tabular}

\vspace{10pt}

\begin{tabular}{lll}
\multicolumn{3}{l}{\textbf{MIT PRÄPOSITIONEN}}\\
& \textbf{sich engagieren für} + Akk & (to commit oneself to / support)\\
& \textbf{sich engagieren bei} + Dat & (to get involved at / with)\\
\end{tabular}
\hfill
\begin{tabular}{lll}
\multicolumn{3}{l}{\textbf{TRENNBARE VERBEN}}\\
& ---\\
\vspace{10pt}
\end{tabular}
\vspace{-5pt}

\end{verbbox}

%sich entscheiden [reflexive] - to decide
\begin{verbbox}\addcontentsline{toc}{subsubsection}{sich entscheiden (to decide)}
        {\Large \bf sich entscheiden} (to decide) \hfill Type: \textbf{regular} \hfill Auxiliary Verb: \textbf{haben} 
    
       \vspace{5pt}

    \hrule
    
    \vspace{5pt}

  \begin{tabular}{rl}
     \multicolumn{2}{l}{\textbf{PRÄSENS} }\\
    ich & \textbf{entscheide mich}\\
    du & \textbf{entscheidest dich}\\
    er/sie/es & \textbf{entscheidet sich}\\
    wir & \textbf{entscheiden uns}\\
    ihr & \textbf{entscheidet euch}\\
    sie/Sie & \textbf{entscheiden sich}\\
  \end{tabular}
  \hfill
     \begin{tabular}{rl}
    \multicolumn{2}{l}{\textbf{PRÄTERITUM} }\\
    ich & \textbf{entschied mich}\\
    du & \textbf{entschiedest dich}\\
    er/sie/es & \textbf{entschied sich}\\
    wir & \textbf{entschieden uns}\\
    ihr & \textbf{entschiedet euch}\\
    sie/Sie & \textbf{entschieden sich}\\
  \end{tabular}
  \hfill
  \begin{tabular}{rl}
     \multicolumn{2}{l}{\textbf{IMPERATIV} }\\
   & \textbf{entscheid(e) (du) dich!}\\
   & \textbf{entscheiden wir uns!}\\
   & \textbf{entscheidet (ihr) euch!}\\
   & \textbf{entscheiden Sie sich!}\\
   & \vspace{10pt}
  \end{tabular}

    \vspace{10pt}

 \begin{tabular}{rl}
     \multicolumn{2}{l}{\textbf{PERFEKT}}\\
    &haben + sich + \textbf{entschieden}\\
  \end{tabular}
\hfill
   \begin{tabular}{rl}
   
    \multicolumn{2}{l}{\textbf{FUTURE I} }\\
    &werden + sich + \textbf{entscheiden}\\
 
  \end{tabular}
\hfill
   \begin{tabular}{rl}
      \multicolumn{2}{l}{\textbf{PLUSQUAMPERFEKT}}\\
    &hatten + sich + \textbf{entschieden}\\
  \end{tabular}
  
\vspace{10pt}

   \begin{tabular}{lll}
      \multicolumn{3}{l}{\textbf{MIT PRÄPOSITIONEN}}\\
& \textbf{sich entscheiden für} + Akk & (to decide for something)\\
& \textbf{sich entscheiden gegen} + Akk & (to decide against something)
  \end{tabular}
  \hfill
   \begin{tabular}{lll}
      \multicolumn{3}{l}{\textbf{TRENNBARE VERBEN}}\\
      & ---& \\
\vspace{10pt}
  \end{tabular}
  
  \vspace{-5pt}

\end{verbbox}

%sich entschuldigen [reflexive] - to apologise
\begin{verbbox-acc}\addcontentsline{toc}{subsubsection}{sich entschuldigen (to apologise)}
        {\Large \bf sich entschuldigen} (to apologise) \hfill Type: \textbf{regular, + Akkusativ} \hfill Auxiliary Verb: \textbf{haben} 
    
       \vspace{5pt}

    \hrule
    
    \vspace{5pt}

  \begin{tabular}{rl}
     \multicolumn{2}{l}{\textbf{PRÄSENS} }\\
    ich & \textbf{entschuldige mich}\\
    du & \textbf{entschuldigst dich}\\
    er/sie/es & \textbf{entschuldigt sich}\\
    wir & \textbf{entschuldigen uns}\\
    ihr & \textbf{entschuldigt euch}\\
    sie/Sie & \textbf{entschuldigen sich}\\
  \end{tabular}
  \hfill
     \begin{tabular}{rl}
    \multicolumn{2}{l}{\textbf{PRÄTERITUM} }\\
    ich & \textbf{entschuldigte mich}\\
    du & \textbf{entschuldigtest dich}\\
    er/sie/es & \textbf{entschuldigte sich}\\
    wir & \textbf{entschuldigten uns}\\
    ihr & \textbf{entschuldigtet euch}\\
    sie/Sie & \textbf{entschuldigten sich}\\
  \end{tabular}
  \hfill
  \begin{tabular}{rl}
     \multicolumn{2}{l}{\textbf{IMPERATIV} }\\
   & \textbf{entschuldig(e) (du) dich!}\\
   & \textbf{entschuldigen wir uns!}\\
   & \textbf{entschuldigt (ihr) euch!}\\
   & \textbf{entschuldigen Sie sich!}\\
   & \vspace{10pt}
  \end{tabular}

    \vspace{10pt}

 \begin{tabular}{rl}
     \multicolumn{2}{l}{\textbf{PERFEKT}}\\
    &haben + sich + \textbf{entschuldigt}\\
  \end{tabular}
\hfill
   \begin{tabular}{rl}
   
    \multicolumn{2}{l}{\textbf{FUTURE I} }\\
    &werden + sich + \textbf{entschuldigen}\\
 
  \end{tabular}
\hfill
   \begin{tabular}{rl}
      \multicolumn{2}{l}{\textbf{PLUSQUAMPERFEKT}}\\
    &hatten + sich + \textbf{entschuldigt}\\
  \end{tabular}
  
\vspace{10pt}

   \begin{tabular}{lll}
      \multicolumn{3}{l}{\textbf{MIT PRÄPOSITIONEN}}\\
& \textbf{sich entschuldigen für} + Akk & (to apologise for)\\
& \textbf{sich entschuldigen bei} + Dat & (to apologise to)\\
& \textbf{sich entschuldigen mit} + Dat & (to plead something)
  \end{tabular}
  \hfill
   \begin{tabular}{lll}
      \multicolumn{3}{l}{\textbf{TRENNBARE VERBEN}}\\
\vspace{30pt}
 \end{tabular}
  
  \vspace{-5pt}

\end{verbbox-acc}

% sich entspannen  - to relax 
\begin{verbbox}\addcontentsline{toc}{subsubsection}{sich entspannen (to relax)}
    {\Large \bf sich entspannen} (to relax) \hfill Type: \textbf{regular} \hfill Auxiliary Verb: \textbf{haben} 
    
     \vspace{5pt}

    \hrule
    
    \vspace{5pt}

  \begin{tabular}{rl}
     \multicolumn{2}{l}{\textbf{PRÄSENS} }\\
    ich & \textbf{entspanne mich}\\
    du & \textbf{entspannst dich}\\
    er/sie/es & \textbf{entspannt sich}\\
    wir & \textbf{entspannen uns}\\
    ihr & \textbf{entspannt euch}\\
    sie/Sie & \textbf{entspannen sich}\\
  \end{tabular}
  \hfill
     \begin{tabular}{rl}
    \multicolumn{2}{l}{\textbf{PRÄTERITUM} }\\
    ich & \textbf{entspannte mich}\\
    du & \textbf{entspanntest dich}\\
    er/sie/es & \textbf{entspannte sich}\\
    wir & \textbf{entspannten uns}\\
    ihr & \textbf{entspanntet euch}\\
    sie/Sie & \textbf{entspannten sich}\\
  \end{tabular}
  \hfill
  \begin{tabular}{rl}
     \multicolumn{2}{l}{\textbf{IMPERATIV} }\\
   & \textbf{entspann dich (du)!}\\
   & \textbf{entspannen wir uns!}\\
   & \textbf{entspannt euch (ihr)!}\\
   & \textbf{entspannen Sie sich!}\\
   & \vspace{10pt}
  \end{tabular}

    \vspace{10pt}

 \begin{tabular}{rl}
     \multicolumn{2}{l}{\textbf{PERFEKT}}\\
   & haben + sich + \textbf{entspannt}\\
  \end{tabular}
\hfill
   \begin{tabular}{rl}
   
    \multicolumn{2}{l}{\textbf{FUTURE I} }\\
   & werden + sich + \textbf{entspannen}\\
 
  \end{tabular}
\hfill
   \begin{tabular}{rl}
      \multicolumn{2}{l}{\textbf{PLUSQUAMPERFEKT}}\\
   & hatten + sich + \textbf{entspannt}\\
  \end{tabular}

   \vspace{10pt}

   \begin{tabular}{lll}
      \multicolumn{3}{l}{\textbf{MIT PRÄPOSITIONEN}}\\
& \textbf{sich entspannen bei} + Dat & (to relax while / during)\\
& \textbf{sich entspannen mit} + Dat & (to relax with)\\
& \textbf{sich entspannen nach} + Dat & (to relax after)\\
\end{tabular}
  \hfill
\begin{tabular}{lll}
\multicolumn{3}{l}{\textbf{TRENNBARE VERBEN}}\\
& ---\\
\vspace{20pt}
\end{tabular}


  \vspace{-5pt}

\end{verbbox}


%sich entwickeln [reflexive] - to develop
\begin{verbbox}\addcontentsline{toc}{subsubsection}{sich entwickeln (to develop)}
        {\Large \bf sich entwickeln} (to develop) \hfill Type: \textbf{regular} \hfill Auxiliary Verb: \textbf{haben} 
    
       \vspace{5pt}

    \hrule
    
    \vspace{5pt}

  \begin{tabular}{rl}
     \multicolumn{2}{l}{\textbf{PRÄSENS} }\\
    ich & \textbf{entwickele mich}\\
    du & \textbf{entwickelst dich}\\
    er/sie/es & \textbf{entwickelt sich}\\
    wir & \textbf{entwickeln uns}\\
    ihr & \textbf{entwickelt euch}\\
    sie/Sie & \textbf{entwickeln sich}\\
  \end{tabular}
  \hfill
     \begin{tabular}{rl}
    \multicolumn{2}{l}{\textbf{PRÄTERITUM} }\\
    ich & \textbf{entwickelte mich}\\
    du & \textbf{ventwickeltest dich}\\
    er/sie/es & \textbf{entwickelte sich}\\
    wir & \textbf{entwickelten uns}\\
    ihr & \textbf{entwickeltet euch}\\
    sie/Sie & \textbf{entwickelten sich}\\
  \end{tabular}
  \hfill
  \begin{tabular}{rl}
     \multicolumn{2}{l}{\textbf{IMPERATIV} }\\
   & \textbf{entwickel(e) (du) dich!}\\
   & \textbf{entwickelen wir uns!}\\
   & \textbf{entwickelt (ihr) euch!}\\
   & \textbf{entwickelen Sie sich!}\\
   & \vspace{10pt}
  \end{tabular}

    \vspace{10pt}

 \begin{tabular}{rl}
     \multicolumn{2}{l}{\textbf{PERFEKT}}\\
    &haben + sich + \textbf{entwickelt}\\
  \end{tabular}
\hfill
   \begin{tabular}{rl}
   
    \multicolumn{2}{l}{\textbf{FUTURE I} }\\
    &werden + sich + \textbf{entwickeln}\\
 
  \end{tabular}
\hfill
   \begin{tabular}{rl}
      \multicolumn{2}{l}{\textbf{PLUSQUAMPERFEKT}}\\
    &hatten + sich + \textbf{entwickelt}\\
  \end{tabular}
  
\vspace{10pt}

   \begin{tabular}{lll}
      \multicolumn{3}{l}{\textbf{MIT PRÄPOSITIONEN}}\\
& \textbf{sich entwickeln aus} + Dat & (to evolve from something)\\
& \textbf{sich entwickeln zu} + Dat & (to evolve into something)
  \end{tabular}
  \hfill
   \begin{tabular}{lll}
      \multicolumn{3}{l}{\textbf{TRENNBARE VERBEN}}\\
      & --- &\\
\vspace{10pt}
  \end{tabular}
  
  \vspace{-5pt}

\end{verbbox}

%sich erholen (to recover / rest)
\begin{verbbox}\addcontentsline{toc}{subsubsection}{sich erholen (to recover / rest)}
{\Large \bf sich erholen} (to recover / rest) \hfill Type: \textbf{regular} \hfill Auxiliary Verb: \textbf{haben}

\vspace{5pt}
\hrule
\vspace{5pt}

\begin{tabular}{rl}
\multicolumn{2}{l}{\textbf{PRÄSENS}}\\
ich & \textbf{erhole mich}\\
du & \textbf{erholst dich}\\
er/sie/es & \textbf{erholt sich}\\
wir & \textbf{erholen uns}\\
ihr & \textbf{erholt euch}\\
sie/Sie & \textbf{erholen sich}\\
\end{tabular}
\hfill
\begin{tabular}{rl}
\multicolumn{2}{l}{\textbf{PRÄTERITUM}}\\
ich & \textbf{erholte mich}\\
du & \textbf{erholtest dich}\\
er/sie/es & \textbf{erholte sich}\\
wir & \textbf{erholten uns}\\
ihr & \textbf{erholtet euch}\\
sie/Sie & \textbf{erholten sich}\\
\end{tabular}
\hfill
\begin{tabular}{rl}
\multicolumn{2}{l}{\textbf{IMPERATIV}}\\
& \textbf{erhole dich!}\\
& \textbf{erholen wir uns!}\\
& \textbf{erholt euch!}\\
& \textbf{erholen Sie sich!}\\
\end{tabular}

\vspace{10pt}

\begin{tabular}{rl}
\multicolumn{2}{l}{\textbf{PERFEKT}}\\
& haben + sich + \textbf{erholt}\\
\end{tabular}
\hfill
\begin{tabular}{rl}
\multicolumn{2}{l}{\textbf{FUTURE I}}\\
& werden + sich + \textbf{erholen}\\
\end{tabular}
\hfill
\begin{tabular}{rl}
\multicolumn{2}{l}{\textbf{PLUSQUAMPERFEKT}}\\
& hatten + sich + \textbf{erholt}\\
\end{tabular}

\vspace{10pt}

\begin{tabular}{lll}
\multicolumn{3}{l}{\textbf{MIT PRÄPOSITIONEN}}\\
& \textbf{sich erholen von} + Dat & (to recover from)\\
& \textbf{sich erholen nach} + Dat & (to recover after)\\
\end{tabular}
\hfill
\begin{tabular}{lll}
\multicolumn{3}{l}{\textbf{TRENNBARE VERBEN}}\\
& ---\\
\vspace{10pt}
\end{tabular}

\vspace{-5pt}

\end{verbbox}

%sich freuen (A1) – to be happy
\begin{verbbox}\addcontentsline{toc}{subsubsection}{sich freuen (to be pleased about)}
    {\Large \bf sich freuen} (to be pleased about) \hfill Type: \textbf{regular} \hfill Auxiliary Verb: \textbf{haben} 
    
     \vspace{5pt}

    \hrule
    
    \vspace{5pt}

  \begin{tabular}{rl}
     \multicolumn{2}{l}{\textbf{PRÄSENS} }\\
    ich & \textbf{freue mich}\\
    du & \textbf{freust dich}\\
    er/sie/es & \textbf{freut sich}\\
    wir & \textbf{freuen uns}\\
    ihr & \textbf{freut euch}\\
    sie/Sie & \textbf{freuen sich}\\
  \end{tabular}
  \hfill
     \begin{tabular}{rl}
    \multicolumn{2}{l}{\textbf{PRÄTERITUM} }\\
    ich & \textbf{freute mich}\\
    du & \textbf{freutest dich}\\
    er/sie/es & \textbf{freute sich}\\
    wir & \textbf{freuten uns}\\
    ihr & \textbf{freutet euch}\\
    sie/Sie & \textbf{freiten sich}\\
  \end{tabular}
  \hfill
  \begin{tabular}{rl}
     \multicolumn{2}{l}{\textbf{IMPERATIV} }\\
   & \textbf{freu(e) (du) dich!}\\
   & \textbf{freuen wir uns!}\\
   & \textbf{freut (ihr) euch!}\\
   & \textbf{freuen Sie sich!}\\
   & \vspace{10pt}
  \end{tabular}

    \vspace{10pt}

 \begin{tabular}{rl}
     \multicolumn{2}{l}{\textbf{PERFEKT}}\\
   &  haben + sich + \textbf{gefreut}\\
  \end{tabular}
\hfill
   \begin{tabular}{rl}
    \multicolumn{2}{l}{\textbf{FUTURE I} }\\
   & werden + sich + \textbf{freuen}\\
  \end{tabular}
\hfill
   \begin{tabular}{rl}
      \multicolumn{2}{l}{\textbf{PLUSQUAMPERFEKT}}\\
    & hatten + sich + \textbf{gefreut}\\
  \end{tabular}
  
\vspace{10pt}

   \begin{tabular}{lll}
      \multicolumn{3}{l}{\textbf{MIT PRÄPOSITIONEN}}\\
& \textbf{sich freuen über} + Akk & (to be please about something)\\
& \textbf{sich freuen auf} + Akk & (to look forward to something)
  \end{tabular}
  \hfill
   \begin{tabular}{lll}
      \multicolumn{3}{l}{\textbf{TRENNBARE VERBEN}}\\
      & --- & \\
\vspace{10pt}
  \end{tabular}



  \vspace{-5pt}


\end{verbbox}

%sich fühlen [reflexive] - to feel
\begin{verbbox-acc}\addcontentsline{toc}{subsubsection}{sich fühlen (to feel)}
        {\Large \bf sich fühlen} (to feel) \hfill Type: \textbf{regular, + Akkusativ} \hfill Auxiliary Verb: \textbf{haben} 
    
       \vspace{5pt}

    \hrule
    
    \vspace{5pt}

  \begin{tabular}{rl}
     \multicolumn{2}{l}{\textbf{PRÄSENS} }\\
    ich & \textbf{fühle mich}\\
    du & \textbf{fühlst dich}\\
    er/sie/es & \textbf{fühlt sich}\\
    wir & \textbf{fühlen uns}\\
    ihr & \textbf{fühlt euch}\\
    sie/Sie & \textbf{fühlen sich}\\
  \end{tabular}
  \hfill
     \begin{tabular}{rl}
    \multicolumn{2}{l}{\textbf{PRÄTERITUM} }\\
    ich & \textbf{fühlte mich}\\
    du & \textbf{fühltest dich}\\
    er/sie/es & \textbf{fühlte sich}\\
    wir & \textbf{fühlten uns}\\
    ihr & \textbf{fühltet euch}\\
    sie/Sie & \textbf{fühlten sich}\\
  \end{tabular}
  \hfill
  \begin{tabular}{rl}
     \multicolumn{2}{l}{\textbf{IMPERATIV} }\\
   & \textbf{fühl(e) (du) dich!}\\
   & \textbf{fühlen wir uns!}\\
   & \textbf{fühlt (ihr) euch!}\\
   & \textbf{fühlen Sie sich!}\\
   & \vspace{10pt}
  \end{tabular}

    \vspace{10pt}

 \begin{tabular}{rl}
     \multicolumn{2}{l}{\textbf{PERFEKT}}\\
    &haben + sich + \textbf{gefühlt}\\
  \end{tabular}
\hfill
   \begin{tabular}{rl}
   
    \multicolumn{2}{l}{\textbf{FUTURE I} }\\
    &werden + sich + \textbf{fühlen}\\
 
  \end{tabular}
\hfill
   \begin{tabular}{rl}
      \multicolumn{2}{l}{\textbf{PLUSQUAMPERFEKT}}\\
    &hatten + sich + \textbf{gefühlt}\\
  \end{tabular}
  
\vspace{10pt}

\begin{tabular}{lll}
\multicolumn{3}{l}{\textbf{MIT PRÄPOSITIONEN}}\\
& ---\\
\end{tabular}
  \hfill
   \begin{tabular}{lll}
      \multicolumn{3}{l}{\textbf{TRENNBARE VERBEN}}\\
     & \textbf{sich wohl$|$fühlen} & (to be comfortable)
  \end{tabular}
  
  \vspace{-5pt}

\end{verbbox-acc}

%sich fürchten (to be afraid)
\begin{verbbox}\addcontentsline{toc}{subsubsection}{sich fürchten (to be afraid)}
    {\Large \bf sich fürchten} (to be afraid) \hfill Type: \textbf{regular} \hfill Auxiliary Verb: \textbf{haben} 
    
     \vspace{5pt}

    \hrule
    
    \vspace{5pt}

  \begin{tabular}{rl}
     \multicolumn{2}{l}{\textbf{PRÄSENS} }\\
    ich & \textbf{fürchte mich}\\
    du & \textbf{fürchtest dich}\\
    er/sie/es & \textbf{fürchtet sich}\\
    wir & \textbf{fürchten uns}\\
    ihr & \textbf{fürchtet euch}\\
    sie/Sie & \textbf{fürchten sich}\\
  \end{tabular}
  \hfill
     \begin{tabular}{rl}
    \multicolumn{2}{l}{\textbf{PRÄTERITUM} }\\
    ich & \textbf{fürchtete mich}\\
    du & \textbf{fürchtetest dich}\\
    er/sie/es & \textbf{fürchtete sich}\\
    wir & \textbf{fürchteten uns}\\
    ihr & \textbf{fürchtetet euch}\\
    sie/Sie & \textbf{fürchteten sich}\\
  \end{tabular}
  \hfill
  \begin{tabular}{rl}
     \multicolumn{2}{l}{\textbf{IMPERATIV} }\\
   & \textbf{fürchte dich!}\\
   & \textbf{fürchten wir uns!}\\
   & \textbf{fürchtet euch!}\\
   & \textbf{fürchten Sie sich!}\\
   & \vspace{10pt}
  \end{tabular}

    \vspace{10pt}

 \begin{tabular}{rl}
     \multicolumn{2}{l}{\textbf{PERFEKT}}\\
   & haben + sich + \textbf{gefürchtet}\\
  \end{tabular}
\hfill
   \begin{tabular}{rl}
    \multicolumn{2}{l}{\textbf{FUTURE I} }\\
   & werden + sich + \textbf{fürchten}\\
  \end{tabular}
\hfill
   \begin{tabular}{rl}
      \multicolumn{2}{l}{\textbf{PLUSQUAMPERFEKT}}\\
   & hatten + sich + \textbf{gefürchtet}\\
  \end{tabular}

   \vspace{10pt}

   \begin{tabular}{lll}
      \multicolumn{3}{l}{\textbf{MIT PRÄPOSITIONEN}}\\
& \textbf{sich fürchten vor} + Dat & (to be afraid of)\\
& \textbf{sich fürchten um} + Akk & (to fear for)\\
\end{tabular}
  \hfill
   \begin{tabular}{lll}
      \multicolumn{3}{l}{\textbf{TRENNBARE VERBEN}}\\
 & \textbf{---} & (none common)\\
\vspace{10pt}
  \end{tabular}

  \vspace{-5pt}

\end{verbbox}

%sich gewöhnen (to get used to)
\begin{verbbox}\addcontentsline{toc}{subsubsection}{sich gewöhnen (to get used to)}
    {\Large \bf sich gewöhnen} (to get used to) \hfill Type: \textbf{regular} \hfill Auxiliary Verb: \textbf{haben} 
    
     \vspace{5pt}

    \hrule
    
    \vspace{5pt}

  \begin{tabular}{rl}
     \multicolumn{2}{l}{\textbf{PRÄSENS} }\\
    ich & \textbf{gewöhne mich}\\
    du & \textbf{gewöhnst dich}\\
    er/sie/es & \textbf{gewöhnt sich}\\
    wir & \textbf{gewöhnen uns}\\
    ihr & \textbf{gewöhnt euch}\\
    sie/Sie & \textbf{gewöhnen sich}\\
  \end{tabular}
  \hfill
     \begin{tabular}{rl}
    \multicolumn{2}{l}{\textbf{PRÄTERITUM} }\\
    ich & \textbf{gewöhnte mich}\\
    du & \textbf{gewöhntest dich}\\
    er/sie/es & \textbf{gewöhnte sich}\\
    wir & \textbf{gewöhnten uns}\\
    ihr & \textbf{gewöhntet euch}\\
    sie/Sie & \textbf{gewöhnten sich}\\
  \end{tabular}
  \hfill
  \begin{tabular}{rl}
     \multicolumn{2}{l}{\textbf{IMPERATIV} }\\
   & \textbf{gewöhn(e) dich!}\\
   & \textbf{gewöhnen wir uns!}\\
   & \textbf{gewöhnt euch!}\\
   & \textbf{gewöhnen Sie sich!}\\
   & \vspace{10pt}
  \end{tabular}

    \vspace{10pt}

 \begin{tabular}{rl}
     \multicolumn{2}{l}{\textbf{PERFEKT}}\\
   & haben + sich + \textbf{gewöhnt}\\
  \end{tabular}
\hfill
   \begin{tabular}{rl}
    \multicolumn{2}{l}{\textbf{FUTURE I} }\\
   & werden + sich + \textbf{gewöhnen}\\
  \end{tabular}
\hfill
   \begin{tabular}{rl}
      \multicolumn{2}{l}{\textbf{PLUSQUAMPERFEKT}}\\
   & hatten + sich + \textbf{gewöhnt}\\
  \end{tabular}

   \vspace{10pt}

   \begin{tabular}{lll}
      \multicolumn{3}{l}{\textbf{MIT PRÄPOSITIONEN}}\\
& \textbf{sich gewöhnen an} + Akk & (to get used to)\\
\end{tabular}
  \hfill
   \begin{tabular}{lll}
      \multicolumn{3}{l}{\textbf{TRENNBARE VERBEN}}\\
 & \textbf{---} & (none common)\\
\vspace{10pt}
  \end{tabular}

  \vspace{-5pt}

\end{verbbox}

%sich hüten (to beware / to be careful)
\begin{verbbox}\addcontentsline{toc}{subsubsection}{sich hüten (to beware / to be careful)}
    {\Large \bf sich hüten} (to beware / to be careful) \hfill Type: \textbf{regular} \hfill Auxiliary Verb: \textbf{haben} 
    
     \vspace{5pt}

    \hrule
    
    \vspace{5pt}

  \begin{tabular}{rl}
     \multicolumn{2}{l}{\textbf{PRÄSENS} }\\
    ich & \textbf{hüte mich}\\
    du & \textbf{hütest dich}\\
    er/sie/es & \textbf{hütet sich}\\
    wir & \textbf{hüten uns}\\
    ihr & \textbf{hütet euch}\\
    sie/Sie & \textbf{hüten sich}\\
  \end{tabular}
  \hfill
     \begin{tabular}{rl}
    \multicolumn{2}{l}{\textbf{PRÄTERITUM} }\\
    ich & \textbf{hütete mich}\\
    du & \textbf{hütetest dich}\\
    er/sie/es & \textbf{hütete sich}\\
    wir & \textbf{hüteten uns}\\
    ihr & \textbf{hütetet euch}\\
    sie/Sie & \textbf{hüteten sich}\\
  \end{tabular}
  \hfill
  \begin{tabular}{rl}
     \multicolumn{2}{l}{\textbf{IMPERATIV} }\\
   & \textbf{hüte dich!}\\
   & \textbf{hüten wir uns!}\\
   & \textbf{hütet euch!}\\
   & \textbf{hüten Sie sich!}\\
   & \vspace{10pt}
  \end{tabular}

    \vspace{10pt}

 \begin{tabular}{rl}
     \multicolumn{2}{l}{\textbf{PERFEKT}}\\
   & haben + sich + \textbf{gehütet}\\
  \end{tabular}
\hfill
   \begin{tabular}{rl}
    \multicolumn{2}{l}{\textbf{FUTURE I} }\\
   & werden + sich + \textbf{hüten}\\
  \end{tabular}
\hfill
   \begin{tabular}{rl}
      \multicolumn{2}{l}{\textbf{PLUSQUAMPERFEKT}}\\
   & hatten + sich + \textbf{gehütet}\\
  \end{tabular}

   \vspace{10pt}

   \begin{tabular}{lll}
      \multicolumn{3}{l}{\textbf{MIT PRÄPOSITIONEN}}\\
& \textbf{sich hüten vor} + Dat & (to beware of)\\
& \textbf{sich hüten, zu} + Inf & (to be careful not to)\\
\end{tabular}
  \hfill
   \begin{tabular}{lll}
      \multicolumn{3}{l}{\textbf{TRENNBARE VERBEN}}\\
 & \textbf{---} & (none common)\\
\vspace{10pt}
  \end{tabular}

  \vspace{-5pt}

\end{verbbox}

%sich interessieren (A2) – to interest
\begin{verbbox}\addcontentsline{toc}{subsubsection}{sich interessieren (to be interested in)}
    {\Large \bf sich interessieren} (to be interest in) \hfill Type: \textbf{regular} \hfill Auxiliary Verb: \textbf{haben} 
    
     \vspace{5pt}

    \hrule
    
    \vspace{5pt}

  \begin{tabular}{rl}
     \multicolumn{2}{l}{\textbf{PRÄSENS} }\\
    ich & \textbf{interessiere mich}\\
    du & \textbf{interessierst dich}\\
    er/sie/es & \textbf{interessiert sich}\\
    wir & \textbf{interessieren uns}\\
    ihr & \textbf{interessiert euch}\\
    sie/Sie & \textbf{interessieren sich}\\
  \end{tabular}
  \hfill
     \begin{tabular}{rl}
    \multicolumn{2}{l}{\textbf{PRÄTERITUM} }\\
    ich & \textbf{interessierte mich}\\
    du & \textbf{interessiertest dich}\\
    er/sie/es & \textbf{interessierte sich}\\
    wir & \textbf{interessierten uns}\\
    ihr & \textbf{interessiertet euch}\\
    sie/Sie & \textbf{interessierten sich}\\
  \end{tabular}
  \hfill
  \begin{tabular}{rl}
     \multicolumn{2}{l}{\textbf{IMPERATIV} }\\
   & \textbf{interessier(e) (du) dich!}\\
   & \textbf{interessieren wir uns!}\\
   & \textbf{interessiert (ihr) euch!}\\
   & \textbf{interessieren Sie sich!}\\
   & \vspace{10pt}
  \end{tabular}

    \vspace{10pt}

 \begin{tabular}{rl}
     \multicolumn{2}{l}{\textbf{PERFEKT}}\\
   & haben + sich + \textbf{interessiert}\\
  \end{tabular}
\hfill
   \begin{tabular}{rl}
    \multicolumn{2}{l}{\textbf{FUTURE I} }\\
    &werden + sich + \textbf{interessieren}\\
  \end{tabular}
\hfill
   \begin{tabular}{rl}
      \multicolumn{2}{l}{\textbf{PLUSQUAMPERFEKT}}\\
   & hatten + sich + \textbf{interessiert}\\
  \end{tabular}
  
\vspace{10pt}

   \begin{tabular}{lll}
      \multicolumn{3}{l}{\textbf{MIT PRÄPOSITIONEN}}\\
& \textbf{sich interessieren für} + Akk & (to be interested in something)\\
  \end{tabular}
 % \hfill
 %  \begin{tabular}{lll}
  %    \multicolumn{3}{l}{\textbf{TRENNBARE VERBEN}}\\
%\vspace{10pt}
 % \end{tabular}

  \vspace{-5pt}

\end{verbbox}

%sich irren (to be mistaken)
\begin{verbbox}\addcontentsline{toc}{subsubsection}{sich irren (to be mistaken)}
    {\Large \bf sich irren} (to be mistaken) \hfill Type: \textbf{regular} \hfill Auxiliary Verb: \textbf{haben} 
    
     \vspace{5pt}

    \hrule
    
    \vspace{5pt}

  \begin{tabular}{rl}
     \multicolumn{2}{l}{\textbf{PRÄSENS} }\\
    ich & \textbf{irre mich}\\
    du & \textbf{irrst dich}\\
    er/sie/es & \textbf{irrt sich}\\
    wir & \textbf{irren uns}\\
    ihr & \textbf{irrt euch}\\
    sie/Sie & \textbf{irren sich}\\
  \end{tabular}
  \hfill
     \begin{tabular}{rl}
    \multicolumn{2}{l}{\textbf{PRÄTERITUM} }\\
    ich & \textbf{irrte mich}\\
    du & \textbf{irrtest dich}\\
    er/sie/es & \textbf{irrte sich}\\
    wir & \textbf{irrten uns}\\
    ihr & \textbf{irrtet euch}\\
    sie/Sie & \textbf{irrten sich}\\
  \end{tabular}
  \hfill
  \begin{tabular}{rl}
     \multicolumn{2}{l}{\textbf{IMPERATIV} }\\
   & \textbf{irr(e) dich!}\\
   & \textbf{irren wir uns!}\\
   & \textbf{irrt euch!}\\
   & \textbf{irren Sie sich!}\\
   & \vspace{10pt}
  \end{tabular}

    \vspace{10pt}

 \begin{tabular}{rl}
     \multicolumn{2}{l}{\textbf{PERFEKT}}\\
   & haben + sich + \textbf{geirrt}\\
  \end{tabular}
\hfill
   \begin{tabular}{rl}
    \multicolumn{2}{l}{\textbf{FUTURE I} }\\
   & werden + sich + \textbf{irren}\\
  \end{tabular}
\hfill
   \begin{tabular}{rl}
      \multicolumn{2}{l}{\textbf{PLUSQUAMPERFEKT}}\\
   & hatten + sich + \textbf{geirrt}\\
  \end{tabular}

   \vspace{10pt}

   \begin{tabular}{lll}
      \multicolumn{3}{l}{\textbf{MIT PRÄPOSITIONEN}}\\
& \textbf{sich irren in} + Dat & (to be mistaken about)\\
& \textbf{sich irren bei} + Dat & (to make a mistake in/with)\\
\end{tabular}
  \hfill
   \begin{tabular}{lll}
      \multicolumn{3}{l}{\textbf{TRENNBARE VERBEN}}\\
 & \textbf{---} & (none common)\\
\vspace{10pt}
  \end{tabular}

  \vspace{-5pt}

\end{verbbox}

%sich konzentrieren [reflexive] - to concentrate
\begin{verbbox}\addcontentsline{toc}{subsubsection}{sich konzentrieren (to concentrate)}
        {\Large \bf sich konzentrieren} (to concentrate) \hfill Type: \textbf{regular} \hfill Auxiliary Verb: \textbf{haben} 
    
       \vspace{5pt}

    \hrule
    
    \vspace{5pt}

  \begin{tabular}{rl}
     \multicolumn{2}{l}{\textbf{PRÄSENS} }\\
    ich & \textbf{konzentriere mich}\\
    du & \textbf{konzentrierst dich}\\
    er/sie/es & \textbf{konzentriert sich}\\
    wir & \textbf{konzentrieren uns}\\
    ihr & \textbf{konzentriert euch}\\
    sie/Sie & \textbf{konzentrieren sich}\\
  \end{tabular}
  \hfill
     \begin{tabular}{rl}
    \multicolumn{2}{l}{\textbf{PRÄTERITUM} }\\
    ich & \textbf{konzentrierte mich}\\
    du & \textbf{konzentriertest dich}\\
    er/sie/es & \textbf{konzentrierte sich}\\
    wir & \textbf{konzentrierten uns}\\
    ihr & \textbf{konzentriertet euch}\\
    sie/Sie & \textbf{konzentrierten sich}\\
  \end{tabular}
  \hfill
  \begin{tabular}{rl}
     \multicolumn{2}{l}{\textbf{IMPERATIV} }\\
   & \textbf{konzentrier(e) (du) dich!}\\
   & \textbf{konzentrieren wir uns!}\\
   & \textbf{konzentriert (ihr) euch!}\\
   & \textbf{konzentrieren Sie sich!}\\
   & \vspace{10pt}
  \end{tabular}

    \vspace{10pt}

 \begin{tabular}{rl}
     \multicolumn{2}{l}{\textbf{PERFEKT}}\\
    &haben + sich + \textbf{konzentriert}\\
  \end{tabular}
\hfill
   \begin{tabular}{rl}
   
    \multicolumn{2}{l}{\textbf{FUTURE I} }\\
    &werden + sich + \textbf{konzentrieren}\\
 
  \end{tabular}
\hfill
   \begin{tabular}{rl}
      \multicolumn{2}{l}{\textbf{PLUSQUAMPERFEKT}}\\
    &hatten + sich + \textbf{konzentriert}\\
  \end{tabular}
  
\vspace{10pt}

   \begin{tabular}{lll}
      \multicolumn{3}{l}{\textbf{MIT PRÄPOSITIONEN}}\\
& \textbf{sich konzentrieren auf} + Akk & (to concentrate on something)
  \end{tabular}
  \hfill
   \begin{tabular}{lll}
      \multicolumn{3}{l}{\textbf{TRENNBARE VERBEN}}\\
      & --- & \\
  \end{tabular}
  
  \vspace{-5pt}

\end{verbbox}

%sich kümmern [reflexive] - to take care of
\begin{verbbox}\addcontentsline{toc}{subsubsection}{sich kümmern (to take care of)}
        {\Large \bf sich kümmern} (to take care of) \hfill Type: \textbf{regular} \hfill Auxiliary Verb: \textbf{haben} 
    
       \vspace{5pt}

    \hrule
    
    \vspace{5pt}

  \begin{tabular}{rl}
     \multicolumn{2}{l}{\textbf{PRÄSENS} }\\
    ich & \textbf{kümmere mich}\\
    du & \textbf{kümmerst dich}\\
    er/sie/es & \textbf{kümmert sich}\\
    wir & \textbf{kümmern uns}\\
    ihr & \textbf{kümmert euch}\\
    sie/Sie & \textbf{kümmern sich}\\
  \end{tabular}
  \hfill
     \begin{tabular}{rl}
    \multicolumn{2}{l}{\textbf{PRÄTERITUM} }\\
    ich & \textbf{kümmerte mich}\\
    du & \textbf{kümmertest dich}\\
    er/sie/es & \textbf{kümmerte sich}\\
    wir & \textbf{kümmerten uns}\\
    ihr & \textbf{kümmertet euch}\\
    sie/Sie & \textbf{kümmerten sich}\\
  \end{tabular}
  \hfill
  \begin{tabular}{rl}
     \multicolumn{2}{l}{\textbf{IMPERATIV} }\\
   & \textbf{kümmer(e) (du) dich!}\\
   & \textbf{kümmern wir uns!}\\
   & \textbf{kümmert (ihr) euch!}\\
   & \textbf{kümmern Sie sich!}\\
   & \vspace{10pt}
  \end{tabular}

    \vspace{10pt}

 \begin{tabular}{rl}
     \multicolumn{2}{l}{\textbf{PERFEKT}}\\
    &haben + sich + \textbf{gekümmert}\\
  \end{tabular}
\hfill
   \begin{tabular}{rl}
   
    \multicolumn{2}{l}{\textbf{FUTURE I} }\\
    &werden + sich + \textbf{kümmern}\\
 
  \end{tabular}
\hfill
   \begin{tabular}{rl}
      \multicolumn{2}{l}{\textbf{PLUSQUAMPERFEKT}}\\
    &hatten + sich + \textbf{gekümmert}\\
  \end{tabular}
  
\vspace{10pt}

   \begin{tabular}{lll}
      \multicolumn{3}{l}{\textbf{MIT PRÄPOSITIONEN}}\\
& \textbf{sich kümmern um} + Akk & (to take care of)
  \end{tabular}
  \hfill
   \begin{tabular}{lll}
      \multicolumn{3}{l}{\textbf{TRENNBARE VERBEN}}\\
& --- & \\
  \end{tabular}
  
  \vspace{-5pt}

\end{verbbox}


%sich melden - to report something
\begin{verbbox}\addcontentsline{toc}{subsubsection}{sich melden (to get in touch / report)}
{\Large \bf sich melden} (to get in touch) \hfill Type: \textbf{regular} \hfill Auxiliary Verb: \textbf{haben}

\vspace{5pt}
\hrule
\vspace{5pt}

\begin{tabular}{rl}
\multicolumn{2}{l}{\textbf{PRÄSENS}}\\
ich & \textbf{melde mich}\\
du & \textbf{meldest dich}\\
er/sie/es & \textbf{meldet sich}\\
wir & \textbf{melden uns}\\
ihr & \textbf{meldet euch}\\
sie/Sie & \textbf{melden sich}\\
\end{tabular}
\hfill
\begin{tabular}{rl}
\multicolumn{2}{l}{\textbf{PRÄTERITUM}}\\
ich & \textbf{meldete mich}\\
du & \textbf{meldetest dich}\\
er/sie/es & \textbf{meldete sich}\\
wir & \textbf{meldeten uns}\\
ihr & \textbf{meldetet euch}\\
sie/Sie & \textbf{meldeten sich}\\
\end{tabular}
\hfill
\begin{tabular}{rl}
\multicolumn{2}{l}{\textbf{IMPERATIV}}\\
& \textbf{melde dich (du)!}\\
& \textbf{melden wir uns!}\\
& \textbf{meldet euch (ihr)!}\\
& \textbf{melden Sie sich!}\\
\end{tabular}

\vspace{10pt}

\begin{tabular}{rl}
\multicolumn{2}{l}{\textbf{PERFEKT}}\\
& haben + sich + \textbf{gemeldet}\\
\end{tabular}
\hfill
\begin{tabular}{rl}
\multicolumn{2}{l}{\textbf{FUTURE I}}\\
& werden + sich +\textbf{melden}\\
\end{tabular}
\hfill
\begin{tabular}{rl}
\multicolumn{2}{l}{\textbf{PLUSQUAMPERFEKT}}\\
& hatten + sich + \textbf{gemeldet}\\
\end{tabular}

\vspace{10pt}

\begin{tabular}{lll}
\multicolumn{3}{l}{\textbf{MIT PRÄPOSITIONEN}}\\
& \textbf{sich melden bei} + Dat & (to contact)\\
& \textbf{sich melden für} + Akk & (to register for)\\
\end{tabular}
\hfill
\begin{tabular}{lll}
\multicolumn{3}{l}{\textbf{TRENNBARE VERBEN}}\\
& \textbf{an$|$melden} & (to register)\\
& \textbf{ab$|$melden} & (to deregister)\\
\end{tabular}
\vspace{-5pt}
\end{verbbox}



%sich rächen (to take revenge)
\begin{verbbox}\addcontentsline{toc}{subsubsection}{sich rächen (to take revenge)}
    {\Large \bf sich rächen} (to take revenge) \hfill Type: \textbf{regular} \hfill Auxiliary Verb: \textbf{haben} 
    
     \vspace{5pt}

    \hrule
    
    \vspace{5pt}

  \begin{tabular}{rl}
     \multicolumn{2}{l}{\textbf{PRÄSENS} }\\
    ich & \textbf{räche mich}\\
    du & \textbf{rächst dich}\\
    er/sie/es & \textbf{rächt sich}\\
    wir & \textbf{rächen uns}\\
    ihr & \textbf{rächt euch}\\
    sie/Sie & \textbf{rächen sich}\\
  \end{tabular}
  \hfill
     \begin{tabular}{rl}
    \multicolumn{2}{l}{\textbf{PRÄTERITUM} }\\
    ich & \textbf{rächte mich}\\
    du & \textbf{rächtest dich}\\
    er/sie/es & \textbf{rächte sich}\\
    wir & \textbf{rächten uns}\\
    ihr & \textbf{rächtet euch}\\
    sie/Sie & \textbf{rächten sich}\\
  \end{tabular}
  \hfill
  \begin{tabular}{rl}
     \multicolumn{2}{l}{\textbf{IMPERATIV} }\\
   & \textbf{räch(e) dich!}\\
   & \textbf{rächen wir uns!}\\
   & \textbf{rächt euch!}\\
   & \textbf{rächen Sie sich!}\\
   & \vspace{10pt}
  \end{tabular}

    \vspace{10pt}

 \begin{tabular}{rl}
     \multicolumn{2}{l}{\textbf{PERFEKT}}\\
   & haben + sich + \textbf{gerächt}\\
  \end{tabular}
\hfill
   \begin{tabular}{rl}
    \multicolumn{2}{l}{\textbf{FUTURE I} }\\
   & werden + sich + \textbf{rächen}\\
  \end{tabular}
\hfill
   \begin{tabular}{rl}
      \multicolumn{2}{l}{\textbf{PLUSQUAMPERFEKT}}\\
   & hatten + sich + \textbf{gerächt}\\
  \end{tabular}

   \vspace{10pt}

   \begin{tabular}{lll}
      \multicolumn{3}{l}{\textbf{MIT PRÄPOSITIONEN}}\\
& \textbf{sich rächen an} + Dat & (to take revenge on)\\
& \textbf{sich rächen für} + Akk & (to take revenge for)\\
\end{tabular}
  \hfill
   \begin{tabular}{lll}
      \multicolumn{3}{l}{\textbf{TRENNBARE VERBEN}}\\
 & \textbf{---} & (none common)\\
\vspace{10pt}
  \end{tabular}

  \vspace{-5pt}

\end{verbbox}

%sich regen (to stir / to move)
\begin{verbbox}\addcontentsline{toc}{subsubsection}{sich regen (to stir / to move)}
    {\Large \bf sich regen} (to stir / to move) \hfill Type: \textbf{regular} \hfill Auxiliary Verb: \textbf{haben} 
    
     \vspace{5pt}

    \hrule
    
    \vspace{5pt}

  \begin{tabular}{rl}
     \multicolumn{2}{l}{\textbf{PRÄSENS} }\\
    ich & \textbf{rege mich}\\
    du & \textbf{regst dich}\\
    er/sie/es & \textbf{regt sich}\\
    wir & \textbf{regen uns}\\
    ihr & \textbf{regt euch}\\
    sie/Sie & \textbf{regen sich}\\
  \end{tabular}
  \hfill
     \begin{tabular}{rl}
    \multicolumn{2}{l}{\textbf{PRÄTERITUM} }\\
    ich & \textbf{regte mich}\\
    du & \textbf{regtest dich}\\
    er/sie/es & \textbf{regte sich}\\
    wir & \textbf{regten uns}\\
    ihr & \textbf{regtet euch}\\
    sie/Sie & \textbf{regten sich}\\
  \end{tabular}
  \hfill
  \begin{tabular}{rl}
     \multicolumn{2}{l}{\textbf{IMPERATIV} }\\
   & \textbf{reg(e) dich!}\\
   & \textbf{regen wir uns!}\\
   & \textbf{regt euch!}\\
   & \textbf{regen Sie sich!}\\
   & \vspace{10pt}
  \end{tabular}

    \vspace{10pt}

 \begin{tabular}{rl}
     \multicolumn{2}{l}{\textbf{PERFEKT}}\\
   & haben + sich + \textbf{geregt}\\
  \end{tabular}
\hfill
   \begin{tabular}{rl}
    \multicolumn{2}{l}{\textbf{FUTURE I} }\\
   & werden + sich + \textbf{regen}\\
  \end{tabular}
\hfill
   \begin{tabular}{rl}
      \multicolumn{2}{l}{\textbf{PLUSQUAMPERFEKT}}\\
   & hatten + sich + \textbf{geregt}\\
  \end{tabular}

   \vspace{10pt}

   \begin{tabular}{lll}
      \multicolumn{3}{l}{\textbf{MIT PRÄPOSITIONEN}}\\
& \textbf{sich regen in} + Dat & (to stir in / move in)\\
& \textbf{sich regen bei} + Dat & (to stir at / during)\\
\end{tabular}
  \hfill
   \begin{tabular}{lll}
      \multicolumn{3}{l}{\textbf{TRENNBARE VERBEN}}\\
 & \textbf{an$|$regen} & (to stimulate / suggest)\\
 & \textbf{auf$|$regen} & (to upset / excite)\\
\vspace{10pt}
  \end{tabular}

  \vspace{-5pt}

\end{verbbox}

%sich schämen [reflexive] - to be ashamed
\begin{verbbox-acc}\addcontentsline{toc}{subsubsection}{sich schämen (to be ashamed)}
        {\Large \bf sich schämen} (to be ashamed) \hfill Type: \textbf{regular, + Akkusativ} \hfill Auxiliary Verb: \textbf{haben} 
    
       \vspace{5pt}

    \hrule
    
    \vspace{5pt}

  \begin{tabular}{rl}
     \multicolumn{2}{l}{\textbf{PRÄSENS} }\\
    ich & \textbf{schäme mich}\\
    du & \textbf{schämst dich}\\
    er/sie/es & \textbf{schämt sich}\\
    wir & \textbf{schämen uns}\\
    ihr & \textbf{schämt euch}\\
    sie/Sie & \textbf{schämen sich}\\
  \end{tabular}
  \hfill
     \begin{tabular}{rl}
    \multicolumn{2}{l}{\textbf{PRÄTERITUM} }\\
    ich & \textbf{schämte mich}\\
    du & \textbf{schämtest dich}\\
    er/sie/es & \textbf{schämte sich}\\
    wir & \textbf{schämten uns}\\
    ihr & \textbf{schämtet euch}\\
    sie/Sie & \textbf{schämten sich}\\
  \end{tabular}
  \hfill
  \begin{tabular}{rl}
     \multicolumn{2}{l}{\textbf{IMPERATIV} }\\
   & \textbf{schäm(e) (du) dich!}\\
   & \textbf{schämen wir uns!}\\
   & \textbf{schämt (ihr) euch!}\\
   & \textbf{schämen Sie sich!}\\
   & \vspace{10pt}
  \end{tabular}

    \vspace{10pt}

 \begin{tabular}{rl}
     \multicolumn{2}{l}{\textbf{PERFEKT}}\\
    &haben + sich + \textbf{geschämt}\\
  \end{tabular}
\hfill
   \begin{tabular}{rl}
   
    \multicolumn{2}{l}{\textbf{FUTURE I} }\\
    &werden + sich + \textbf{schämen}\\
 
  \end{tabular}
\hfill
   \begin{tabular}{rl}
      \multicolumn{2}{l}{\textbf{PLUSQUAMPERFEKT}}\\
    &hatten + sich + \textbf{geschämt}\\
  \end{tabular}
  
\vspace{10pt}

   \begin{tabular}{lll}
      \multicolumn{3}{l}{\textbf{MIT PRÄPOSITIONEN}}\\
& \textbf{sich schämen für} + Akk & (to be ashamed of)
  \end{tabular}
 % \hfill
 %  \begin{tabular}{lll}
  %    \multicolumn{3}{l}{\textbf{TRENNBARE VERBEN}}\\
%\vspace{10pt}
 % \end{tabular}
  
  \vspace{-5pt}

\end{verbbox-acc}

%sich schminken (to put on makeup)
\begin{verbbox}\addcontentsline{toc}{subsubsection}{sich schminken (to put on makeup)}
{\Large \bf sich schminken} (to put on makeup) \hfill Type: \textbf{regular} \hfill Auxiliary Verb: \textbf{haben}

\vspace{5pt}
\hrule
\vspace{5pt}

\begin{tabular}{rl}
\multicolumn{2}{l}{\textbf{PRÄSENS}}\\
ich & \textbf{schminke mich}\\
du & \textbf{schminkst dich}\\
er/sie/es & \textbf{schminkt sich}\\
wir & \textbf{schminken uns}\\
ihr & \textbf{schminkt euch}\\
sie/Sie & \textbf{schminken sich}\\
\end{tabular}
\hfill
\begin{tabular}{rl}
\multicolumn{2}{l}{\textbf{PRÄTERITUM}}\\
ich & \textbf{schminkte mich}\\
du & \textbf{schminktest dich}\\
er/sie/es & \textbf{schminkte sich}\\
wir & \textbf{schminkten uns}\\
ihr & \textbf{schminktet euch}\\
sie/Sie & \textbf{schminkten sich}\\
\end{tabular}
\hfill
\begin{tabular}{rl}
\multicolumn{2}{l}{\textbf{IMPERATIV}}\\
& \textbf{schminke dich!}\\
& \textbf{schminken wir uns!}\\
& \textbf{schminkt euch!}\\
& \textbf{schminken Sie sich!}\\
\end{tabular}

\vspace{10pt}

\begin{tabular}{rl}
\multicolumn{2}{l}{\textbf{PERFEKT}}\\
& haben + sich + \textbf{geschminkt}\\
\end{tabular}
\hfill
\begin{tabular}{rl}
\multicolumn{2}{l}{\textbf{FUTURE I}}\\
& werden + sich + \textbf{schminken}\\
\end{tabular}
\hfill
\begin{tabular}{rl}
\multicolumn{2}{l}{\textbf{PLUSQUAMPERFEKT}}\\
& hatten + sich + \textbf{geschminkt}\\
\end{tabular}

\vspace{10pt}

\begin{tabular}{lll}
\multicolumn{3}{l}{\textbf{MIT PRÄPOSITIONEN}}\\
& \textbf{sich schminken für} + Akk & (to put on makeup for)\\
& \textbf{sich schminken mit} + Dat & (to put makeup on with)\\
\end{tabular}
\hfill
\begin{tabular}{lll}
\multicolumn{3}{l}{\textbf{TRENNBARE VERBEN}}\\
& \textbf{ab$|$schminken} & (to remove makeup)\\
& \textbf{an$|$schminken} & (to apply makeup)\\
\end{tabular}
\vspace{-5pt}
\end{verbbox}

%sich schützen (to protect oneself)
\begin{verbbox}\addcontentsline{toc}{subsubsection}{sich schützen (to protect oneself)}
    {\Large \bf sich schützen} (to protect oneself) \hfill Type: \textbf{regular} \hfill Auxiliary Verb: \textbf{haben} 
    
     \vspace{5pt}

    \hrule
    
    \vspace{5pt}

  \begin{tabular}{rl}
     \multicolumn{2}{l}{\textbf{PRÄSENS} }\\
    ich & \textbf{schütze mich}\\
    du & \textbf{schützt dich}\\
    er/sie/es & \textbf{schützt sich}\\
    wir & \textbf{schützen uns}\\
    ihr & \textbf{schützt euch}\\
    sie/Sie & \textbf{schützen sich}\\
  \end{tabular}
  \hfill
     \begin{tabular}{rl}
    \multicolumn{2}{l}{\textbf{PRÄTERITUM} }\\
    ich & \textbf{schützte mich}\\
    du & \textbf{schütztest dich}\\
    er/sie/es & \textbf{schützte sich}\\
    wir & \textbf{schützten uns}\\
    ihr & \textbf{schütztet euch}\\
    sie/Sie & \textbf{schützten sich}\\
  \end{tabular}
  \hfill
  \begin{tabular}{rl}
     \multicolumn{2}{l}{\textbf{IMPERATIV} }\\
   & \textbf{schütz(e) dich!}\\
   & \textbf{schützen wir uns!}\\
   & \textbf{schützt euch!}\\
   & \textbf{schützen Sie sich!}\\
   & \vspace{10pt}
  \end{tabular}

    \vspace{10pt}

 \begin{tabular}{rl}
     \multicolumn{2}{l}{\textbf{PERFEKT}}\\
   & haben + sich + \textbf{geschützt}\\
  \end{tabular}
\hfill
   \begin{tabular}{rl}
    \multicolumn{2}{l}{\textbf{FUTURE I} }\\
   & werden + sich + \textbf{schützen}\\
  \end{tabular}
\hfill
   \begin{tabular}{rl}
      \multicolumn{2}{l}{\textbf{PLUSQUAMPERFEKT}}\\
   & hatten + sich + \textbf{geschützt}\\
  \end{tabular}

   \vspace{10pt}

   \begin{tabular}{lll}
      \multicolumn{3}{l}{\textbf{MIT PRÄPOSITIONEN}}\\
& \textbf{sich schützen vor} + Dat & (to protect oneself from)\\
& \textbf{sich schützen gegen} + Akk & (to protect oneself against)\\
\end{tabular}
  \hfill
   \begin{tabular}{lll}
      \multicolumn{3}{l}{\textbf{TRENNBARE VERBEN}}\\
 & \textbf{---} & (none common)\\
\vspace{10pt}
  \end{tabular}

  \vspace{-5pt}

\end{verbbox}

%sich sehnen (to long for)
\begin{verbbox}\addcontentsline{toc}{subsubsection}{sich sehnen (to long for)}
    {\Large \bf sich sehnen} (to long for) \hfill Type: \textbf{regular, reflexive} \hfill Auxiliary Verb: \textbf{haben} 
    
     \vspace{5pt}

    \hrule
    
    \vspace{5pt}

  \begin{tabular}{rl}
     \multicolumn{2}{l}{\textbf{PRÄSENS} }\\
    ich & \textbf{sehne mich}\\
    du & \textbf{sehnst dich}\\
    er/sie/es & \textbf{sehnt sich}\\
    wir & \textbf{sehnen uns}\\
    ihr & \textbf{sehnt euch}\\
    sie/Sie & \textbf{sehnen sich}\\
  \end{tabular}
  \hfill
     \begin{tabular}{rl}
    \multicolumn{2}{l}{\textbf{PRÄTERITUM} }\\
    ich & \textbf{sehnte mich}\\
    du & \textbf{sehntest dich}\\
    er/sie/es & \textbf{sehnte sich}\\
    wir & \textbf{sehnten uns}\\
    ihr & \textbf{sehntet euch}\\
    sie/Sie & \textbf{sehnten sich}\\
  \end{tabular}
  \hfill
  \begin{tabular}{rl}
     \multicolumn{2}{l}{\textbf{IMPERATIV} }\\
   & \textbf{sehn(e) dich!}\\
   & \textbf{sehnen wir uns!}\\
   & \textbf{sehnt euch!}\\
   & \textbf{sehnen Sie sich!}\\
   & \vspace{10pt}
  \end{tabular}

    \vspace{10pt}

 \begin{tabular}{rl}
     \multicolumn{2}{l}{\textbf{PERFEKT}}\\
   & haben + sich + \textbf{gesehnt}\\
  \end{tabular}
\hfill
   \begin{tabular}{rl}
    \multicolumn{2}{l}{\textbf{FUTURE I} }\\
   & werden + sich + \textbf{sehnen}\\
  \end{tabular}
\hfill
   \begin{tabular}{rl}
      \multicolumn{2}{l}{\textbf{PLUSQUAMPERFEKT}}\\
   & hatten + sich + \textbf{gesehnt}\\
  \end{tabular}

   \vspace{10pt}

   \begin{tabular}{lll}
      \multicolumn{3}{l}{\textbf{MIT PRÄPOSITIONEN}}\\
& \textbf{sich sehnen nach} + Dat & (to long for)\\
\end{tabular}
  \hfill
   \begin{tabular}{lll}
      \multicolumn{3}{l}{\textbf{TRENNBARE VERBEN}}\\
 & \textbf{---} & (none common)\\
\vspace{10pt}
  \end{tabular}

  \vspace{-5pt}

\end{verbbox}

%sich setzen (to sit down / take a seat)
\begin{verbbox}\addcontentsline{toc}{subsubsection}{sich setzen (to sit down / take a seat)}
{\Large \bf sich setzen} (to sit down / take a seat) \hfill Type: \textbf{regular} \hfill Auxiliary Verb: \textbf{haben}

\vspace{5pt}
\hrule
\vspace{5pt}

\begin{tabular}{rl}
\multicolumn{2}{l}{\textbf{PRÄSENS}}\\
ich & \textbf{setze mich}\\
du & \textbf{setzt dich}\\
er/sie/es & \textbf{setzt sich}\\
wir & \textbf{setzen uns}\\
ihr & \textbf{setzt euch}\\
sie/Sie & \textbf{setzen sich}\\
\end{tabular}
\hfill
\begin{tabular}{rl}
\multicolumn{2}{l}{\textbf{PRÄTERITUM}}\\
ich & \textbf{setzte mich}\\
du & \textbf{setztest dich}\\
er/sie/es & \textbf{setzte sich}\\
wir & \textbf{setzten uns}\\
ihr & \textbf{setztet euch}\\
sie/Sie & \textbf{setzten sich}\\
\end{tabular}
\hfill
\begin{tabular}{rl}
\multicolumn{2}{l}{\textbf{IMPERATIV}}\\
& \textbf{setze dich!}\\
& \textbf{setzen wir uns!}\\
& \textbf{setzt euch!}\\
& \textbf{setzen Sie sich!}\\
\end{tabular}

\vspace{10pt}

\begin{tabular}{rl}
\multicolumn{2}{l}{\textbf{PERFEKT}}\\
& haben + sich + \textbf{gesetzt}\\
\end{tabular}
\hfill
\begin{tabular}{rl}
\multicolumn{2}{l}{\textbf{FUTURE I}}\\
& werden + sich + \textbf{setzen}\\
\end{tabular}
\hfill
\begin{tabular}{rl}
\multicolumn{2}{l}{\textbf{PLUSQUAMPERFEKT}}\\
& hatten + sich + \textbf{gesetzt}\\
\end{tabular}

\vspace{10pt}

\begin{tabular}{lll}
\multicolumn{3}{l}{\textbf{MIT PRÄPOSITIONEN}}\\
& \textbf{sich setzen auf} + Akk & (to sit down on)\\
& \textbf{sich setzen neben} + Akk & (to sit next to)\\
\end{tabular}
\hfill
\begin{tabular}{lll}
\multicolumn{3}{l}{\textbf{TRENNBARE VERBEN}}\\
& \textbf{hin$|$setzen} & (to sit down)\\
\end{tabular}
\vspace{-5pt}
\end{verbbox}

%sich sorgen [reflexive] - to worry
\begin{verbbox}\addcontentsline{toc}{subsubsection}{sich sorgen (to worry)}
        {\Large \bf sich sorgen} (to worry) \hfill Type: \textbf{regular} \hfill Auxiliary Verb: \textbf{haben} 
    
       \vspace{5pt}

    \hrule
    
    \vspace{5pt}

  \begin{tabular}{rl}
     \multicolumn{2}{l}{\textbf{PRÄSENS} }\\
    ich & \textbf{sorge mich}\\
    du & \textbf{sorgst dich}\\
    er/sie/es & \textbf{sorgt sich}\\
    wir & \textbf{sorgen uns}\\
    ihr & \textbf{sorgt euch}\\
    sie/Sie & \textbf{sorgen sich}\\
  \end{tabular}
  \hfill
     \begin{tabular}{rl}
    \multicolumn{2}{l}{\textbf{PRÄTERITUM} }\\
    ich & \textbf{sorgte mich}\\
    du & \textbf{sorgtest dich}\\
    er/sie/es & \textbf{sorgte sich}\\
    wir & \textbf{sorgten uns}\\
    ihr & \textbf{sorgtet euch}\\
    sie/Sie & \textbf{sorgten sich}\\
  \end{tabular}
  \hfill
  \begin{tabular}{rl}
     \multicolumn{2}{l}{\textbf{IMPERATIV} }\\
   & \textbf{sorg(e) (du) dich!}\\
   & \textbf{sorgen wir uns!}\\
   & \textbf{sorgt (ihr) euch!}\\
   & \textbf{sorgen Sie sich!}\\
   & \vspace{10pt}
  \end{tabular}

    \vspace{10pt}

 \begin{tabular}{rl}
     \multicolumn{2}{l}{\textbf{PERFEKT}}\\
    &haben + sich + \textbf{gesorgt}\\
  \end{tabular}
\hfill
   \begin{tabular}{rl}
   
    \multicolumn{2}{l}{\textbf{FUTURE I} }\\
    &werden + sich + \textbf{sorgen}\\
 
  \end{tabular}
\hfill
   \begin{tabular}{rl}
      \multicolumn{2}{l}{\textbf{PLUSQUAMPERFEKT}}\\
    &hatten + sich + \textbf{gesorgt}\\
  \end{tabular}
  
\vspace{10pt}

   \begin{tabular}{lll}
      \multicolumn{3}{l}{\textbf{MIT PRÄPOSITIONEN}}\\
& \textbf{sich sorgen um} + Akk & (to worry about something)
  \end{tabular}
  \hfill
   \begin{tabular}{lll}
      \multicolumn{3}{l}{\textbf{TRENNBARE VERBEN}}\\
& --- & \\
  \end{tabular}
  
  \vspace{-5pt}

\end{verbbox}

%sich stellen [reflexive] - to surrender
\begin{verbbox-acc}\addcontentsline{toc}{subsubsection}{sich stellen (to surrender)}
        {\Large \bf sich stellen} (to surrender) \hfill Type: \textbf{regular, + Akkusativ} \hfill Auxiliary Verb: \textbf{haben} 
    
       \vspace{5pt}

    \hrule
    
    \vspace{5pt}

  \begin{tabular}{rl}
     \multicolumn{2}{l}{\textbf{PRÄSENS} }\\
    ich & \textbf{stelle mich}\\
    du & \textbf{stellst dich}\\
    er/sie/es & \textbf{stellt sich}\\
    wir & \textbf{stellen uns}\\
    ihr & \textbf{stellt euch}\\
    sie/Sie & \textbf{stellen sich}\\
  \end{tabular}
  \hfill
     \begin{tabular}{rl}
    \multicolumn{2}{l}{\textbf{PRÄTERITUM} }\\
    ich & \textbf{stellte mich}\\
    du & \textbf{stelltest dich}\\
    er/sie/es & \textbf{stellte sich}\\
    wir & \textbf{stellten uns}\\
    ihr & \textbf{stelltet euch}\\
    sie/Sie & \textbf{stellten sich}\\
  \end{tabular}
  \hfill
  \begin{tabular}{rl}
     \multicolumn{2}{l}{\textbf{IMPERATIV} }\\
   & \textbf{stell(e) (du) dich!}\\
   & \textbf{stellen wir uns!}\\
   & \textbf{stellt (ihr) euch!}\\
   & \textbf{stellen Sie sich!}\\
   & \vspace{10pt}
  \end{tabular}

    \vspace{10pt}

 \begin{tabular}{rl}
     \multicolumn{2}{l}{\textbf{PERFEKT}}\\
    &haben + sich + \textbf{gestellt}\\
  \end{tabular}
\hfill
   \begin{tabular}{rl}
   
    \multicolumn{2}{l}{\textbf{FUTURE I} }\\
    &werden + sich + \textbf{stellen}\\
 
  \end{tabular}
\hfill
   \begin{tabular}{rl}
      \multicolumn{2}{l}{\textbf{PLUSQUAMPERFEKT}}\\
    &hatten + sich + \textbf{gestellt}\\
  \end{tabular}
  
\vspace{10pt}

   \begin{tabular}{lll}
      \multicolumn{3}{l}{\textbf{MIT PRÄPOSITIONEN}}\\
\vspace{0pt}
  \end{tabular}
  \hfill
   \begin{tabular}{lll}
      \multicolumn{3}{l}{\textbf{TRENNBARE VERBEN}}\\
& \textbf{sich vor$|$stellen} & (to introduce oneself)
  \end{tabular}
  
  \vspace{-5pt}

\end{verbbox-acc}

%sich teilen (A2) – to share
\begin{verbbox-dat}\addcontentsline{toc}{subsubsection}{sich teilen (to share)}
    {\Large \bf sich teilen} (to share) \hfill Type: \textbf{regular, + Dativ} \hfill Auxiliary Verb: \textbf{haben} 
    
     \vspace{5pt}

    \hrule
    
    \vspace{5pt}

  \begin{tabular}{rl}
     \multicolumn{2}{l}{\textbf{PRÄSENS} }\\
    ich & \textbf{teile mir}\\
    du & \textbf{teilst dir}\\
    er/sie/es & \textbf{teilt sich}\\
    wir & \textbf{teilen uns}\\
    ihr & \textbf{teilt euch}\\
    sie/Sie & \textbf{teilen sich}\\
  \end{tabular}
  \hfill
     \begin{tabular}{rl}
    \multicolumn{2}{l}{\textbf{PRÄTERITUM} }\\
    ich & \textbf{teilte mir}\\
    du & \textbf{teiltest dir}\\
    er/sie/es & \textbf{teilte sich}\\
    wir & \textbf{teilten uns}\\
    ihr & \textbf{teiltet euch}\\
    sie/Sie & \textbf{teilten sich}\\
  \end{tabular}
  \hfill
  \begin{tabular}{rl}
     \multicolumn{2}{l}{\textbf{IMPERATIV} }\\
   & \textbf{teil(e) (du) dir!}\\
   & \textbf{teilen wir uns!}\\
   & \textbf{teilt (ihr) euch!}\\
   & \textbf{teilen Sie sich!}\\
   & \vspace{10pt}
  \end{tabular}

    \vspace{10pt}

 \begin{tabular}{rl}
     \multicolumn{2}{l}{\textbf{PERFEKT}}\\
    &haben + sich + \textbf{geteilt}\\
  \end{tabular}
\hfill
   \begin{tabular}{rl}
   
    \multicolumn{2}{l}{\textbf{FUTURE I} }\\
    &werden + sich + \textbf{teilen}\\
 
  \end{tabular}
\hfill
   \begin{tabular}{rl}
      \multicolumn{2}{l}{\textbf{PLUSQUAMPERFEKT}}\\
    &hatten + sich + \textbf{geteilt}\\
  \end{tabular}
  
\vspace{10pt}

   \begin{tabular}{lll}
      \multicolumn{3}{l}{\textbf{MIT PRÄPOSITIONEN}}\\
& \textbf{sich teilen in} + Akk & (to split into parts)\\
& \textbf{sich teilen mit} + Dat & (to share with someone)
  \end{tabular}
 % \hfill
 %  \begin{tabular}{lll}
  %    \multicolumn{3}{l}{\textbf{TRENNBARE VERBEN}}\\
%\vspace{10pt}
 % \end{tabular}


  \vspace{-5pt}

\end{verbbox-dat}

%sich treffen [reflexive] - to meet up
\begin{verbbox-acc}\addcontentsline{toc}{subsubsection}{sich treffen (to meet up)}
        {\Large \bf sich treffen} (to come forward) \hfill Type: \textbf{irregular, + Akkusativ} \hfill Auxiliary Verb: \textbf{haben} 
    
       \vspace{5pt}

    \hrule
    
    \vspace{5pt}

  \begin{tabular}{rl}
     \multicolumn{2}{l}{\textbf{PRÄSENS} }\\
    ich & \textbf{treffe mich}\\
    du & \textbf{triffst dich}\\
    er/sie/es & \textbf{trifft sich}\\
    wir & \textbf{treffen uns}\\
    ihr & \textbf{trefft euch}\\
    sie/Sie & \textbf{treffen sich}\\
  \end{tabular}
  \hfill
     \begin{tabular}{rl}
    \multicolumn{2}{l}{\textbf{PRÄTERITUM} }\\
    ich & \textbf{traf mich}\\
    du & \textbf{trafst dich}\\
    er/sie/es & \textbf{traf sich}\\
    wir & \textbf{trafen uns}\\
    ihr & \textbf{traft euch}\\
    sie/Sie & \textbf{trafen sich}\\
  \end{tabular}
  \hfill
  \begin{tabular}{rl}
     \multicolumn{2}{l}{\textbf{IMPERATIV} }\\
   & \textbf{triff(e) (du) dich!}\\
   & \textbf{treffen wir uns!}\\
   & \textbf{trefft (ihr) euch!}\\
   & \textbf{treffen Sie sich!}\\
   & \vspace{10pt}
  \end{tabular}

    \vspace{10pt}

 \begin{tabular}{rl}
     \multicolumn{2}{l}{\textbf{PERFEKT}}\\
    &haben + sich + \textbf{getroffen}\\
  \end{tabular}
\hfill
   \begin{tabular}{rl}
   
    \multicolumn{2}{l}{\textbf{FUTURE I} }\\
    &werden + sich + \textbf{treffen}\\
 
  \end{tabular}
\hfill
   \begin{tabular}{rl}
      \multicolumn{2}{l}{\textbf{PLUSQUAMPERFEKT}}\\
    &hatten + sich + \textbf{getroffen}\\
  \end{tabular}
  
\vspace{10pt}

   \begin{tabular}{lll}
      \multicolumn{3}{l}{\textbf{MIT PRÄPOSITIONEN}}\\
& \textbf{sich treffen mit} + Dat & (to meet someone)\\
& \textbf{sich treffen zu} + Dat & (to meet for a purpose)\\
  \end{tabular}
  \hfill
   \begin{tabular}{lll}
      \multicolumn{3}{l}{\textbf{TRENNBARE VERBEN}}\\
      & --- & \\
\vspace{10pt}
 \end{tabular}


 
  
  \vspace{-5pt}


\end{verbbox-acc}


%sich trennen [reflexive] - to separate
\begin{verbbox-acc}\addcontentsline{toc}{subsubsection}{sich trennen (to separate)}
        {\Large \bf sich trennen} (to separate) \hfill Type: \textbf{regular, + Akkusativ} \hfill Auxiliary Verb: \textbf{haben} 
    
       \vspace{5pt}

    \hrule
    
    \vspace{5pt}

  \begin{tabular}{rl}
     \multicolumn{2}{l}{\textbf{PRÄSENS} }\\
    ich & \textbf{trenne mich}\\
    du & \textbf{trennst dich}\\
    er/sie/es & \textbf{trennt sich}\\
    wir & \textbf{trennen uns}\\
    ihr & \textbf{trennt euch}\\
    sie/Sie & \textbf{trennen sich}\\
  \end{tabular}
  \hfill
     \begin{tabular}{rl}
    \multicolumn{2}{l}{\textbf{PRÄTERITUM} }\\
    ich & \textbf{trennte mich}\\
    du & \textbf{trenntest dich}\\
    er/sie/es & \textbf{trennte sich}\\
    wir & \textbf{trennten uns}\\
    ihr & \textbf{trenntet euch}\\
    sie/Sie & \textbf{trennten sich}\\
  \end{tabular}
  \hfill
  \begin{tabular}{rl}
     \multicolumn{2}{l}{\textbf{IMPERATIV} }\\
   & \textbf{trenn(e) (du) dich!}\\
   & \textbf{trennen wir uns!}\\
   & \textbf{trennt (ihr) euch!}\\
   & \textbf{trennen Sie sich!}\\
   & \vspace{10pt}
  \end{tabular}

    \vspace{10pt}

 \begin{tabular}{rl}
     \multicolumn{2}{l}{\textbf{PERFEKT}}\\
    &haben + sich + \textbf{getrennt}\\
  \end{tabular}
\hfill
   \begin{tabular}{rl}
   
    \multicolumn{2}{l}{\textbf{FUTURE I} }\\
    &werden + sich + \textbf{trennen}\\
 
  \end{tabular}
\hfill
   \begin{tabular}{rl}
      \multicolumn{2}{l}{\textbf{PLUSQUAMPERFEKT}}\\
    &hatten + sich + \textbf{getrennt}\\
  \end{tabular}
  
\vspace{10pt}

   \begin{tabular}{lll}
      \multicolumn{3}{l}{\textbf{MIT PRÄPOSITIONEN}}\\
& \textbf{sich trennen von} + Dat & (to separate from someone / something)
  \end{tabular}
 % \hfill
 %  \begin{tabular}{lll}
  %    \multicolumn{3}{l}{\textbf{TRENNBARE VERBEN}}\\
%\vspace{10pt}
 % \end{tabular}
  
  \vspace{-5pt}

\end{verbbox-acc}

%sich unterhalten [reflexive] - to chat
\begin{verbbox}\addcontentsline{toc}{subsubsection}{sich unterhalten (to chat)}
        {\Large \bf sich unterhalten} (to chat) \hfill Type: \textbf{irregular} \hfill Auxiliary Verb: \textbf{haben} 
    
       \vspace{5pt}

    \hrule
    
    \vspace{5pt}

  \begin{tabular}{rl}
     \multicolumn{2}{l}{\textbf{PRÄSENS} }\\
    ich & \textbf{unterhalte mich}\\
    du & \textbf{unterhältst dich}\\
    er/sie/es & \textbf{unterhält sich}\\
    wir & \textbf{unterhalten uns}\\
    ihr & \textbf{unterhaltet euch}\\
    sie/Sie & \textbf{unterhalten sich}\\
  \end{tabular}
  \hfill
     \begin{tabular}{rl}
    \multicolumn{2}{l}{\textbf{PRÄTERITUM} }\\
    ich & \textbf{underhielt mich}\\
    du & \textbf{underhieltst dich}\\
    er/sie/es & \textbf{underhielt sich}\\
    wir & \textbf{underhielten uns}\\
    ihr & \textbf{underhieltet euch}\\
    sie/Sie & \textbf{underhielten sich}\\
  \end{tabular}
  \hfill
  \begin{tabular}{rl}
     \multicolumn{2}{l}{\textbf{IMPERATIV} }\\
   & \textbf{unterhalt(e) (du) dich!}\\
   & \textbf{unterhalten wir uns!}\\
   & \textbf{unterhaltt (ihr) euch!}\\
   & \textbf{unterhalten Sie sich!}\\
   & \vspace{10pt}
  \end{tabular}

    \vspace{10pt}

 \begin{tabular}{rl}
     \multicolumn{2}{l}{\textbf{PERFEKT}}\\
    &haben + sich + \textbf{unterhalten}\\
  \end{tabular}
\hfill
   \begin{tabular}{rl}
   
    \multicolumn{2}{l}{\textbf{FUTURE I} }\\
    &werden + sich + \textbf{unterhalten}\\
 
  \end{tabular}
\hfill
   \begin{tabular}{rl}
      \multicolumn{2}{l}{\textbf{PLUSQUAMPERFEKT}}\\
    &hatten + sich + \textbf{unterhalten}\\
  \end{tabular}
  
\vspace{10pt}

   \begin{tabular}{lll}
      \multicolumn{3}{l}{\textbf{MIT PRÄPOSITIONEN}}\\
& \textbf{sich unterhalten über} + Akk & (to chat about)\\
& \textbf{sich unterhalten mit} + Dat & (to chat mit)\\
  \end{tabular}
 \hfill
   \begin{tabular}{lll}
      \multicolumn{3}{l}{\textbf{TRENNBARE VERBEN}}\\
      & --- &\\
\vspace{20pt}
  \end{tabular}
  
  \vspace{-5pt}

\end{verbbox}

%sich verabreden (to arrange / make an appointment)
\begin{verbbox}\addcontentsline{toc}{subsubsection}{sich verabreden (to arrange / make an appointment)}
{\Large \bf sich verabreden} (to arrange / make an appointment) \hfill Type: \textbf{regular} \hfill Auxiliary Verb: \textbf{haben}

\vspace{5pt}
\hrule
\vspace{5pt}

\begin{tabular}{rl}
\multicolumn{2}{l}{\textbf{PRÄSENS}}\\
ich & \textbf{verabrede mich}\\
du & \textbf{verabredest dich}\\
er/sie/es & \textbf{verabredet sich}\\
wir & \textbf{verabreden uns}\\
ihr & \textbf{verabredet euch}\\
sie/Sie & \textbf{verabreden sich}\\
\end{tabular}
\hfill
\begin{tabular}{rl}
\multicolumn{2}{l}{\textbf{PRÄTERITUM}}\\
ich & \textbf{verabredete mich}\\
du & \textbf{verabredetest dich}\\
er/sie/es & \textbf{verabredete sich}\\
wir & \textbf{verabredeten uns}\\
ihr & \textbf{verabredetet euch}\\
sie/Sie & \textbf{verabredeten sich}\\
\end{tabular}
\hfill
\begin{tabular}{rl}
\multicolumn{2}{l}{\textbf{IMPERATIV}}\\
& \textbf{verabrede dich!}\\
& \textbf{verabreden wir uns!}\\
& \textbf{verabredet euch!}\\
& \textbf{verabreden Sie sich!}\\
\end{tabular}

\vspace{10pt}

\begin{tabular}{rl}
\multicolumn{2}{l}{\textbf{PERFEKT}}\\
& haben + sich + \textbf{verabredet}\\
\end{tabular}
\hfill
\begin{tabular}{rl}
\multicolumn{2}{l}{\textbf{FUTURE I}}\\
& werden + sich + \textbf{verabreden}\\
\end{tabular}
\hfill
\begin{tabular}{rl}
\multicolumn{2}{l}{\textbf{PLUSQUAMPERFEKT}}\\
& hatten + sich + \textbf{verabredet}\\
\end{tabular}

\vspace{10pt}

\begin{tabular}{lll}
\multicolumn{3}{l}{\textbf{MIT PRÄPOSITIONEN}}\\
& \textbf{sich verabreden mit} + Dat & (to arrange with)\\
\end{tabular}
\hfill
\begin{tabular}{lll}
\multicolumn{3}{l}{\textbf{TRENNBARE VERBEN}}\\
& ---\\
\end{tabular}
\vspace{-5pt}
\end{verbbox}

%sich verlassen (to rely)
\begin{verbbox}\addcontentsline{toc}{subsubsection}{sich verlassen (to rely)}
    {\Large \bf sich verlassen} (to rely) \hfill Type: \textbf{irregular} \hfill Auxiliary Verb: \textbf{haben} 
    
     \vspace{5pt}

    \hrule
    
    \vspace{5pt}

  \begin{tabular}{rl}
     \multicolumn{2}{l}{\textbf{PRÄSENS} }\\
    ich & \textbf{verlasse mich}\\
    du & \textbf{verlässt dich}\\
    er/sie/es & \textbf{verlässt sich}\\
    wir & \textbf{verlassen uns}\\
    ihr & \textbf{verlasst euch}\\
    sie/Sie & \textbf{verlassen sich}\\
  \end{tabular}
  \hfill
     \begin{tabular}{rl}
    \multicolumn{2}{l}{\textbf{PRÄTERITUM} }\\
    ich & \textbf{verließ mich}\\
    du & \textbf{verließest dich}\\
    er/sie/es & \textbf{verließ sich}\\
    wir & \textbf{verließen uns}\\
    ihr & \textbf{verließt euch}\\
    sie/Sie & \textbf{verließen sich}\\
  \end{tabular}
  \hfill
  \begin{tabular}{rl}
     \multicolumn{2}{l}{\textbf{IMPERATIV} }\\
   & \textbf{verlass(e) dich!}\\
   & \textbf{verlassen wir uns!}\\
   & \textbf{verlasst euch!}\\
   & \textbf{verlassen Sie sich!}\\
   & \vspace{10pt}
  \end{tabular}

    \vspace{10pt}

 \begin{tabular}{rl}
     \multicolumn{2}{l}{\textbf{PERFEKT}}\\
   & haben + sich + \textbf{verlassen}\\
  \end{tabular}
\hfill
   \begin{tabular}{rl}
    \multicolumn{2}{l}{\textbf{FUTURE I} }\\
   & werden + sich + \textbf{verlassen}\\
  \end{tabular}
\hfill
   \begin{tabular}{rl}
      \multicolumn{2}{l}{\textbf{PLUSQUAMPERFEKT}}\\
   & hatten + sich + \textbf{verlassen}\\
  \end{tabular}

   \vspace{10pt}

   \begin{tabular}{lll}
      \multicolumn{3}{l}{\textbf{MIT PRÄPOSITIONEN}}\\
& \textbf{sich verlassen auf} + Akk & (to rely on)\\
\end{tabular}
  \hfill
   \begin{tabular}{lll}
      \multicolumn{3}{l}{\textbf{TRENNBARE VERBEN}}\\
 & \textbf{---} & (none common)\\
\vspace{10pt}
  \end{tabular}

  \vspace{-5pt}

\end{verbbox}

%sich verletzen (to injure oneself)
\begin{verbbox}\addcontentsline{toc}{subsubsection}{sich verletzen (to injure oneself)}
{\Large \bf sich verletzen} (to injure oneself) \hfill Type: \textbf{regular} \hfill Auxiliary Verb: \textbf{haben}

\vspace{5pt}
\hrule
\vspace{5pt}

\begin{tabular}{rl}
\multicolumn{2}{l}{\textbf{PRÄSENS}}\\
ich & \textbf{verletze mich}\\
du & \textbf{verletzt dich}\\
er/sie/es & \textbf{verletzt sich}\\
wir & \textbf{verletzen uns}\\
ihr & \textbf{verletzt euch}\\
sie/Sie & \textbf{verletzen sich}\\
\end{tabular}
\hfill
\begin{tabular}{rl}
\multicolumn{2}{l}{\textbf{PRÄTERITUM}}\\
ich & \textbf{verletzte mich}\\
du & \textbf{verletztest dich}\\
er/sie/es & \textbf{verletzte sich}\\
wir & \textbf{verletzten uns}\\
ihr & \textbf{verletz tet euch}\\
sie/Sie & \textbf{verletzten sich}\\
\end{tabular}
\hfill
\begin{tabular}{rl}
\multicolumn{2}{l}{\textbf{IMPERATIV}}\\
& \textbf{verletze dich (du) nicht!}\\
& \textbf{verletzen wir uns nicht!}\\
& \textbf{verletzt euch (ihr) nicht!}\\
& \textbf{verletzen Sie sich nicht!}\\
\end{tabular}

\vspace{10pt}

\begin{tabular}{rl}
\multicolumn{2}{l}{\textbf{PERFEKT}}\\
& haben + sich + \textbf{verletzt}\\
\end{tabular}
\hfill
\begin{tabular}{rl}
\multicolumn{2}{l}{\textbf{FUTURE I}}\\
& werden + sich + \textbf{verletzen}\\
\end{tabular}
\hfill
\begin{tabular}{rl}
\multicolumn{2}{l}{\textbf{PLUSQUAMPERFEKT}}\\
& hatten + sich + \textbf{verletzt}\\
\end{tabular}

\vspace{10pt}

\begin{tabular}{lll}
\multicolumn{3}{l}{\textbf{MIT PRÄPOSITIONEN}}\\
& \textbf{sich verletzen an} + Dat & (to injure oneself on)\\
& \textbf{sich verletzen bei} + Dat & (to injure oneself while doing)\\
\end{tabular}
\hfill
\begin{tabular}{lll}
\multicolumn{3}{l}{\textbf{TRENNBARE VERBEN}}\\
& --- & (ver- prefix inseparable)\\
\end{tabular}
\vspace{-5pt}
\end{verbbox}

%sich verlieben [reflexive] - to fall in love
\begin{verbbox}\addcontentsline{toc}{subsubsection}{sich verlieben (to fall in love)}
        {\Large \bf sich verlieben} (to fall in love) \hfill Type: \textbf{regular} \hfill Auxiliary Verb: \textbf{haben} 
    
       \vspace{5pt}

    \hrule
    
    \vspace{5pt}

  \begin{tabular}{rl}
     \multicolumn{2}{l}{\textbf{PRÄSENS} }\\
    ich & \textbf{verliebe mir}\\
    du & \textbf{verliebst dir}\\
    er/sie/es & \textbf{verliebt sich}\\
    wir & \textbf{verlieben uns}\\
    ihr & \textbf{verliebt euch}\\
    sie/Sie & \textbf{verlieben sich}\\
  \end{tabular}
  \hfill
     \begin{tabular}{rl}
    \multicolumn{2}{l}{\textbf{PRÄTERITUM} }\\
    ich & \textbf{verliebte mir}\\
    du & \textbf{verliebtest dir}\\
    er/sie/es & \textbf{verliebte sich}\\
    wir & \textbf{verliebten uns}\\
    ihr & \textbf{verliebtet euch}\\
    sie/Sie & \textbf{verliebten sich}\\
  \end{tabular}
  \hfill
  \begin{tabular}{rl}
     \multicolumn{2}{l}{\textbf{IMPERATIV} }\\
   & \textbf{verlieb(e) (du) dir!}\\
   & \textbf{verlieben wir uns!}\\
   & \textbf{verliebt (ihr) euch!}\\
   & \textbf{verlieben Sie sich!}\\
   & \vspace{10pt}
  \end{tabular}

    \vspace{10pt}

 \begin{tabular}{rl}
     \multicolumn{2}{l}{\textbf{PERFEKT}}\\
    &haben + sich + \textbf{verliebt}\\
  \end{tabular}
\hfill
   \begin{tabular}{rl}
   
    \multicolumn{2}{l}{\textbf{FUTURE I} }\\
    &werden + sich + \textbf{verlieben}\\
 
  \end{tabular}
\hfill
   \begin{tabular}{rl}
      \multicolumn{2}{l}{\textbf{PLUSQUAMPERFEKT}}\\
    &hatten + sich + \textbf{verliebt}\\
  \end{tabular}
  
\vspace{10pt}

   \begin{tabular}{lll}
      \multicolumn{3}{l}{\textbf{MIT PRÄPOSITIONEN}}\\
& \textbf{sich verlieben in} + Akk & (to fall in love with someone)
  \end{tabular}
  \hfill
   \begin{tabular}{lll}
     \multicolumn{3}{l}{\textbf{TRENNBARE VERBEN}}\\
& --- & \\
 \end{tabular}
  
  \vspace{-5pt}

\end{verbbox}

%sich verschreiben  (write incorrectly)
\begin{verbbox}\addcontentsline{toc}{subsubsection}{sich verschreiben (to commit / write incorrectly)}
{\Large \bf sich verschreiben} (to commit / write incorrectly) \hfill Type: \textbf{irregular} \hfill Auxiliary Verb: \textbf{haben}

\vspace{5pt}
\hrule
\vspace{5pt}

\begin{tabular}{rl}
\multicolumn{2}{l}{\textbf{PRÄSENS}}\\
ich & \textbf{verschreibe mich}\\
du & \textbf{verschreibst dich}\\
er/sie/es & \textbf{verschreibt sich}\\
wir & \textbf{verschreiben uns}\\
ihr & \textbf{verschreibt euch}\\
sie/Sie & \textbf{verschreiben sich}\\
\end{tabular}
\hfill
\begin{tabular}{rl}
\multicolumn{2}{l}{\textbf{PRÄTERITUM}}\\
ich & \textbf{verschrieb mich}\\
du & \textbf{verschriebst dich}\\
er/sie/es & \textbf{verschrieb sich}\\
wir & \textbf{verschrieben uns}\\
ihr & \textbf{verschriebt euch}\\
sie/Sie & \textbf{verschrieben sich}\\
\end{tabular}
\hfill
\begin{tabular}{rl}
\multicolumn{2}{l}{\textbf{IMPERATIV}}\\
& \textbf{verschreibe (du) dich!}\\
& \textbf{verschreiben wir uns!}\\
& \textbf{verschreibt (ihr) euch!}\\
& \textbf{verschreiben Sie sich!}\\
\end{tabular}

\vspace{10pt}

\begin{tabular}{rl}
\multicolumn{2}{l}{\textbf{PERFEKT}}\\
&haben + sich + \textbf{verschrieben}\\
\end{tabular}
\hfill
\begin{tabular}{rl}
\multicolumn{2}{l}{\textbf{FUTURE I}}\\
& werden + sich + \textbf{verschreiben}\\
\end{tabular}
\hfill
\begin{tabular}{rl}
\multicolumn{2}{l}{\textbf{PLUSQUAMPERFEKT}}\\
& hatten + sich + \textbf{verschrieben}\\
\end{tabular}

\vspace{10pt}

\begin{tabular}{lll}
\multicolumn{3}{l}{\textbf{MIT PRÄPOSITIONEN}}\\
\end{tabular}
\hfill
\begin{tabular}{lll}
\multicolumn{3}{l}{\textbf{TRENNBARE VERBEN}}\\
\end{tabular}
\vspace{-5pt}
\end{verbbox}

%sich verstehen [reflexive] - to get along
\begin{verbbox-acc}\addcontentsline{toc}{subsubsection}{sich verstehen (to get along)}
        {\Large \bf sich verstehen} (to get along) \hfill Type: \textbf{irregular, + Akkusativ} \hfill Auxiliary Verb: \textbf{haben} 
    
       \vspace{5pt}

    \hrule
    
    \vspace{5pt}

  \begin{tabular}{rl}
     \multicolumn{2}{l}{\textbf{PRÄSENS} }\\
    ich & \textbf{verstehe mich}\\
    du & \textbf{verstehst dich}\\
    er/sie/es & \textbf{versteht sich}\\
    wir & \textbf{verstehen uns}\\
    ihr & \textbf{versteht euch}\\
    sie/Sie & \textbf{verstehen sich}\\
  \end{tabular}
  \hfill
     \begin{tabular}{rl}
    \multicolumn{2}{l}{\textbf{PRÄTERITUM} }\\
    ich & \textbf{verstand mich}\\
    du & \textbf{verstandest dich}\\
    er/sie/es & \textbf{verstand sich}\\
    wir & \textbf{verstanden uns}\\
    ihr & \textbf{verstandet euch}\\
    sie/Sie & \textbf{verstanden sich}\\
  \end{tabular}
  \hfill
  \begin{tabular}{rl}
     \multicolumn{2}{l}{\textbf{IMPERATIV} }\\
   & \textbf{versteh(e) (du) dich!}\\
   & \textbf{verstehen wir uns!}\\
   & \textbf{versteht (ihr) euch!}\\
   & \textbf{verstehen Sie sich!}\\
   & \vspace{10pt}
  \end{tabular}

    \vspace{10pt}

 \begin{tabular}{rl}
     \multicolumn{2}{l}{\textbf{PERFEKT}}\\
    &haben + sich + \textbf{verstanden}\\
  \end{tabular}
\hfill
   \begin{tabular}{rl}
   
    \multicolumn{2}{l}{\textbf{FUTURE I} }\\
    &werden + sich + \textbf{verstehen}\\
 
  \end{tabular}
\hfill
   \begin{tabular}{rl}
      \multicolumn{2}{l}{\textbf{PLUSQUAMPERFEKT}}\\
    &hatten + sich + \textbf{verstanden}\\
  \end{tabular}
  
\vspace{10pt}

   \begin{tabular}{lll}
      \multicolumn{3}{l}{\textbf{MIT PRÄPOSITIONEN}}\\
& \textbf{sich verlieben mit} + Dat & (to get along with)
  \end{tabular}
 % \hfill
 %  \begin{tabular}{lll}
  %    \multicolumn{3}{l}{\textbf{TRENNBARE VERBEN}}\\
%\vspace{10pt}
 % \end{tabular}
  
  \vspace{-5pt}

\end{verbbox-acc}

%sich verwählen (to misdial / miscall)
\begin{verbbox}\addcontentsline{toc}{subsubsection}{sich verwählen (to misdial / miscall)}
{\Large \bf sich verwählen} (to misdial / miscall) \hfill Type: \textbf{regular} \hfill Auxiliary Verb: \textbf{haben}

\vspace{5pt}
\hrule
\vspace{5pt}

\begin{tabular}{rl}
\multicolumn{2}{l}{\textbf{PRÄSENS}}\\
ich & \textbf{verwähle mich}\\
du & \textbf{verwählst dich}\\
er/sie/es & \textbf{verwählt sich}\\
wir & \textbf{verwählen uns}\\
ihr & \textbf{verwählt euch}\\
sie/Sie & \textbf{verwählen sich}\\
\end{tabular}
\hfill
\begin{tabular}{rl}
\multicolumn{2}{l}{\textbf{PRÄTERITUM}}\\
ich & \textbf{verwählte mich}\\
du & \textbf{verwähltest dich}\\
er/sie/es & \textbf{verwählte sich}\\
wir & \textbf{verwählten uns}\\
ihr & \textbf{verwähltet euch}\\
sie/Sie & \textbf{verwählten sich}\\
\end{tabular}
\hfill
\begin{tabular}{rl}
\multicolumn{2}{l}{\textbf{IMPERATIV}}\\
& \textbf{verwähle dich!}\\
& \textbf{verwählen wir uns!}\\
& \textbf{verwählt euch!}\\
& \textbf{verwählen Sie sich!}\\
\end{tabular}

\vspace{10pt}

\begin{tabular}{rl}
\multicolumn{2}{l}{\textbf{PERFEKT}}\\
& haben + sich + \textbf{verwählt}\\
\end{tabular}
\hfill
\begin{tabular}{rl}
\multicolumn{2}{l}{\textbf{FUTURE I}}\\
& werden + sich + \textbf{verwählen}\\
\end{tabular}
\hfill
\begin{tabular}{rl}
\multicolumn{2}{l}{\textbf{PLUSQUAMPERFEKT}}\\
& hatten + sich + \textbf{verwählt}\\
\end{tabular}

\vspace{10pt}

\begin{tabular}{lll}
\multicolumn{3}{l}{\textbf{MIT PRÄPOSITIONEN}}\\
& ---\\
\end{tabular}
\hfill
\begin{tabular}{lll}
\multicolumn{3}{l}{\textbf{TRENNBARE VERBEN}}\\
& ---\\
\end{tabular}
\vspace{-5pt}
\end{verbbox}

%sich waschen [reflexive] - to wash oneself
\begin{verbbox-acc}\addcontentsline{toc}{subsubsection}{sich waschen (to wash oneself)}
        {\Large \bf sich waschen} (to wash) \hfill Type: \textbf{irregular, + Akkusativ} \hfill Auxiliary Verb: \textbf{haben} 
    
       \vspace{5pt}

    \hrule
    
    \vspace{5pt}

  \begin{tabular}{rl}
     \multicolumn{2}{l}{\textbf{PRÄSENS} }\\
    ich & \textbf{wasche mich}\\
    du & \textbf{wäschst dich}\\
    er/sie/es & \textbf{wäscht sich}\\
    wir & \textbf{waschen uns}\\
    ihr & \textbf{wascht euch}\\
    sie/Sie & \textbf{waschen sich}\\
  \end{tabular}
  \hfill
     \begin{tabular}{rl}
    \multicolumn{2}{l}{\textbf{PRÄTERITUM} }\\
    ich & \textbf{wunch mich}\\
    du & \textbf{wunchst dich}\\
    er/sie/es & \textbf{wunch sich}\\
    wir & \textbf{wunchen uns}\\
    ihr & \textbf{wuncht euch}\\
    sie/Sie & \textbf{wunchen sich}\\
  \end{tabular}
  \hfill
  \begin{tabular}{rl}
     \multicolumn{2}{l}{\textbf{IMPERATIV} }\\
   & \textbf{wasch(e) (du) dich!}\\
   & \textbf{waschen wir uns!}\\
   & \textbf{wascht (ihr) euch!}\\
   & \textbf{waschen Sie sich!}\\
   & \vspace{10pt}
  \end{tabular}

    \vspace{10pt}

 \begin{tabular}{rl}
     \multicolumn{2}{l}{\textbf{PERFEKT}}\\
    &haben + sich + \textbf{gewaschen}\\
  \end{tabular}
\hfill
   \begin{tabular}{rl}
   
    \multicolumn{2}{l}{\textbf{FUTURE I} }\\
    &werden + sich + \textbf{waschen}\\
 
  \end{tabular}
\hfill
   \begin{tabular}{rl}
      \multicolumn{2}{l}{\textbf{PLUSQUAMPERFEKT}}\\
    &hatten + sich + \textbf{gewaschen}\\
  \end{tabular}
  
\vspace{10pt}

\begin{tabular}{lll}
\multicolumn{3}{l}{\textbf{MIT PRÄPOSITIONEN}}\\
& ---\\
\end{tabular}
\hfill
\begin{tabular}{lll}
\multicolumn{3}{l}{\textbf{TRENNBARE VERBEN}}\\
& ---\\
\end{tabular}

\vspace{-5pt}

\end{verbbox-acc}

%sich wundern [reflexive] - to be surprised
\begin{verbbox}\addcontentsline{toc}{subsubsection}{sich wundern (to wonder / be surprised)}
        {\Large \bf sich wundern} (to wonder / be suprised) \hfill Type: \textbf{regular} \hfill Auxiliary Verb: \textbf{haben} 
    
       \vspace{5pt}

    \hrule
    
    \vspace{5pt}

  \begin{tabular}{rl}
     \multicolumn{2}{l}{\textbf{PRÄSENS} }\\
    ich & \textbf{wundere mich}\\
    du & \textbf{wunderst dich}\\
    er/sie/es & \textbf{wundert sich}\\
    wir & \textbf{wundern uns}\\
    ihr & \textbf{wundert euch}\\
    sie/Sie & \textbf{wundern sich}\\
  \end{tabular}
  \hfill
     \begin{tabular}{rl}
    \multicolumn{2}{l}{\textbf{PRÄTERITUM} }\\
    ich & \textbf{wunderte mich}\\
    du & \textbf{wundertest dich}\\
    er/sie/es & \textbf{wunderte sich}\\
    wir & \textbf{wunderten uns}\\
    ihr & \textbf{wundertet euch}\\
    sie/Sie & \textbf{wunderten sich}\\
  \end{tabular}
  \hfill
  \begin{tabular}{rl}
     \multicolumn{2}{l}{\textbf{IMPERATIV} }\\
   & \textbf{wunder(e) (du) dich!}\\
   & \textbf{wundern wir uns!}\\
   & \textbf{wundert (ihr) euch!}\\
   & \textbf{wundern Sie sich!}\\
   & \vspace{10pt}
  \end{tabular}

    \vspace{10pt}

 \begin{tabular}{rl}
     \multicolumn{2}{l}{\textbf{PERFEKT}}\\
    &haben + sich + \textbf{gewundert}\\
  \end{tabular}
\hfill
   \begin{tabular}{rl}
   
    \multicolumn{2}{l}{\textbf{FUTURE I} }\\
    &werden + sich + \textbf{wundern}\\
 
  \end{tabular}
\hfill
   \begin{tabular}{rl}
      \multicolumn{2}{l}{\textbf{PLUSQUAMPERFEKT}}\\
    &hatten + sich + \textbf{gewundert}\\
  \end{tabular}
  
\vspace{10pt}

   \begin{tabular}{lll}
      \multicolumn{3}{l}{\textbf{MIT PRÄPOSITIONEN}}\\
& \textbf{sich wundern über} + Akk & (to wonder / be suprised about)
  \end{tabular}
  \hfill
  \begin{tabular}{lll}
      \multicolumn{3}{l}{\textbf{TRENNBARE VERBEN}}\\
& --- & \\
  \end{tabular}
  
  \vspace{-5pt}

\end{verbbox}

%sich ziehen [reflexive] - to drag
\begin{verbbox-acc}\addcontentsline{toc}{subsubsection}{sich ziehen (to drag)}
        {\Large \bf sich ziehen} (to drag) \hfill Type: \textbf{irregular, + Akkusativ} \hfill Auxiliary Verb: \textbf{haben} 
    
       \vspace{5pt}

    \hrule
    
    \vspace{5pt}

  \begin{tabular}{rl}
     \multicolumn{2}{l}{\textbf{PRÄSENS} }\\
    ich & \textbf{ziehe mich}\\
    du & \textbf{ziehst dich}\\
    er/sie/es & \textbf{zieht sich}\\
    wir & \textbf{ziehen uns}\\
    ihr & \textbf{zieht euch}\\
    sie/Sie & \textbf{ziehen sich}\\
  \end{tabular}
  \hfill
     \begin{tabular}{rl}
    \multicolumn{2}{l}{\textbf{PRÄTERITUM} }\\
    ich & \textbf{zog mich}\\
    du & \textbf{zogst dich}\\
    er/sie/es & \textbf{zog sich}\\
    wir & \textbf{zogen uns}\\
    ihr & \textbf{zogt euch}\\
    sie/Sie & \textbf{zogen sich}\\
  \end{tabular}
  \hfill
  \begin{tabular}{rl}
     \multicolumn{2}{l}{\textbf{IMPERATIV} }\\
   & \textbf{zieh(e) (du) dich!}\\
   & \textbf{ziehen wir uns!}\\
   & \textbf{zieht (ihr) euch!}\\
   & \textbf{ziehen Sie sich!}\\
   & \vspace{10pt}
  \end{tabular}

    \vspace{10pt}

 \begin{tabular}{rl}
     \multicolumn{2}{l}{\textbf{PERFEKT}}\\
    &haben + sich + \textbf{gezogen}\\
  \end{tabular}
\hfill
   \begin{tabular}{rl}
   
    \multicolumn{2}{l}{\textbf{FUTURE I} }\\
    &werden + sich + \textbf{ziehen}\\
 
  \end{tabular}
\hfill
   \begin{tabular}{rl}
      \multicolumn{2}{l}{\textbf{PLUSQUAMPERFEKT}}\\
    &hatten + sich + \textbf{gezogen}\\
  \end{tabular}
  
\vspace{10pt}

   \begin{tabular}{lll}
      \multicolumn{3}{l}{\textbf{MIT PRÄPOSITIONEN}}\\
\vspace{0pt}
  \end{tabular}
  \hfill
   \begin{tabular}{lll}
    \multicolumn{3}{l}{\textbf{TRENNBARE VERBEN}}\\
& \textbf{sich an$|$ziehan} + Akk & (to get dressed)
 \end{tabular}
  
  \vspace{-5pt}

\end{verbbox-acc}

\section{Verbs}

\subsection{A Verbs}

\begin{verbbox}\addcontentsline{toc}{subsubsection}{achten (to pay attention to / to respect)}
    {\Large \bf achten} (to pay attention to / to respect) \hfill Type: \textbf{regular} \hfill Auxiliary Verb: \textbf{haben} 
    
     \vspace{5pt}

    \hrule
    
    \vspace{5pt}

  \begin{tabular}{rl}
     \multicolumn{2}{l}{\textbf{PRÄSENS} }\\
    ich & \textbf{achte}\\
    du & \textbf{achtest}\\
    er/sie/es & \textbf{achtet}\\
    wir & \textbf{achten}\\
    ihr & \textbf{achtet}\\
    sie/Sie & \textbf{achten}\\
  \end{tabular}
  \hfill
     \begin{tabular}{rl}
    \multicolumn{2}{l}{\textbf{PRÄTERITUM} }\\
    ich & \textbf{achtete}\\
    du & \textbf{achtetest}\\
    er/sie/es & \textbf{achtete}\\
    wir & \textbf{achteten}\\
    ihr & \textbf{achtetet}\\
    sie/Sie & \textbf{achteten}\\
  \end{tabular}
  \hfill
  \begin{tabular}{rl}
     \multicolumn{2}{l}{\textbf{IMPERATIV} }\\
   & \textbf{acht(e) (du)!}\\
   & \textbf{achten wir!}\\
   & \textbf{achtet (ihr)!}\\
   & \textbf{achten Sie!}\\
   & \vspace{10pt}
  \end{tabular}

    \vspace{10pt}

 \begin{tabular}{rl}
     \multicolumn{2}{l}{\textbf{PERFEKT}}\\
   & haben + \textbf{geachtet}\\
  \end{tabular}
\hfill
   \begin{tabular}{rl}
   
    \multicolumn{2}{l}{\textbf{FUTURE I} }\\
   & werden + \textbf{achten}\\
 
  \end{tabular}
\hfill
   \begin{tabular}{rl}
      \multicolumn{2}{l}{\textbf{PLUSQUAMPERFEKT}}\\
   & hatten + \textbf{geachtet}\\
  \end{tabular}

   \vspace{10pt}

   \begin{tabular}{lll}
      \multicolumn{3}{l}{\textbf{MIT PRÄPOSITIONEN}}\\
& \textbf{achten auf} + Akk & (to pay attention to)\\
& \textbf{achten} + Akk & (to respect / value)\\
\end{tabular}
  \hfill
   \begin{tabular}{lll}
      \multicolumn{3}{l}{\textbf{TRENNBARE VERBEN}}\\
 & \textbf{---} & (none common)\\
\vspace{10pt}
  \end{tabular}


  \vspace{-5pt}


\end{verbbox}










%ahmen (to imitate)
\begin{verbbox}\addcontentsline{toc}{subsubsection}{ahmen (to imitate)}
{\Large \bf ahmen} (to imitate) \hfill Type: \textbf{regular} \hfill Auxiliary Verb: \textbf{haben}

\vspace{5pt}
\hrule
\vspace{5pt}

\begin{tabular}{rl}
\multicolumn{2}{l}{\textbf{PRÄSENS}}\\
ich & \textbf{ahme}\\
du & \textbf{ahmst}\\
er/sie/es & \textbf{ahmt}\\
wir & \textbf{ahmen}\\
ihr & \textbf{ahmt}\\
sie/Sie & \textbf{ahmen}\\
\end{tabular}
\hfill
\begin{tabular}{rl}
\multicolumn{2}{l}{\textbf{PRÄTERITUM}}\\
ich & \textbf{ahmte}\\
du & \textbf{ahmtest}\\
er/sie/es & \textbf{ahmte}\\
wir & \textbf{ahmten}\\
ihr & \textbf{ahmtet}\\
sie/Sie & \textbf{ahmten}\\
\end{tabular}
\hfill
\begin{tabular}{rl}
\multicolumn{2}{l}{\textbf{IMPERATIV}}\\
& \textbf{ahme (du)!}\\
& \textbf{ahmen wir!}\\
& \textbf{ahmt (ihr)!}\\
& \textbf{ahmen Sie!}\\
\end{tabular}

\vspace{10pt}

\begin{tabular}{rl}
\multicolumn{2}{l}{\textbf{PERFEKT}}\\
& haben + \textbf{geahmt}\\
\end{tabular}
\hfill
\begin{tabular}{rl}
\multicolumn{2}{l}{\textbf{FUTURE I}}\\
& werden + \textbf{ahmen}\\
\end{tabular}
\hfill
\begin{tabular}{rl}
\multicolumn{2}{l}{\textbf{PLUSQUAMPERFEKT}}\\
& hatten + \textbf{geahmt}\\
\end{tabular}

\vspace{10pt}

\begin{tabular}{lll}
\multicolumn{3}{l}{\textbf{MIT PRÄPOSITIONEN}}\\
\end{tabular}
\hfill
\begin{tabular}{lll}
\multicolumn{3}{l}{\textbf{TRENNBARE VERBEN}}\\
& \textbf{nach$|$ahmen} & (to imitate / emulate)\\
\end{tabular}
\vspace{-5pt}
\end{verbbox}

%akzeptieren (to accept)
\begin{verbbox}\addcontentsline{toc}{subsubsection}{akzeptieren (to accept)}
{\Large \bf akzeptieren} (to accept) \hfill Type: \textbf{regular} \hfill Auxiliary Verb: \textbf{haben} 

\vspace{5pt}
\hrule
\vspace{5pt}

\begin{tabular}{rl}
\multicolumn{2}{l}{\textbf{PRÄSENS}}\\
ich & \textbf{akzeptiere}\\
du & \textbf{akzeptierst}\\
er/sie/es & \textbf{akzeptiert}\\
wir & \textbf{akzeptieren}\\
ihr & \textbf{akzeptiert}\\
sie/Sie & \textbf{akzeptieren}\\
\end{tabular}
\hfill
\begin{tabular}{rl}
\multicolumn{2}{l}{\textbf{PRÄTERITUM}}\\
ich & \textbf{akzeptierte}\\
du & \textbf{akzeptiertest}\\
er/sie/es & \textbf{akzeptierte}\\
wir & \textbf{akzeptierten}\\
ihr & \textbf{akzeptiertet}\\
sie/Sie & \textbf{akzeptierten}\\
\end{tabular}
\hfill
\begin{tabular}{rl}
\multicolumn{2}{l}{\textbf{IMPERATIV}}\\
& \textbf{akzeptiere (du)!}\\
& \textbf{akzeptieren wir!}\\
& \textbf{akzeptiert (ihr)!}\\
& \textbf{akzeptieren Sie!}\\
\end{tabular}

\vspace{10pt}

\begin{tabular}{rl}
\multicolumn{2}{l}{\textbf{PERFEKT}}\\
& haben + \textbf{akzeptiert}\\
\end{tabular}
\hfill
\begin{tabular}{rl}
\multicolumn{2}{l}{\textbf{FUTURE I}}\\
& werden + \textbf{akzeptieren}\\
\end{tabular}
\hfill
\begin{tabular}{rl}
\multicolumn{2}{l}{\textbf{PLUSQUAMPERFEKT}}\\
& hatten + \textbf{akzeptiert}\\
\end{tabular}

\vspace{10pt}

\begin{tabular}{lll}
\multicolumn{3}{l}{\textbf{MIT PRÄPOSITIONEN}}\\
& \textbf{akzeptieren} + Akk & (to accept something)\\
\end{tabular}
\hfill
\begin{tabular}{lll}
\multicolumn{3}{l}{\textbf{TRENNBARE VERBEN}}\\
& ---\\
\end{tabular}

\vspace{-5pt}
\end{verbbox}

%ändern (A2) – to change
\begin{verbbox-acc}\addcontentsline{toc}{subsubsection}{ändern (to change)}
    {\Large \bf ändern} (to change) \hfill Type: \textbf{regular, + Akkusativ} \hfill Auxiliary Verb: \textbf{haben} 
    
     \vspace{5pt}

    \hrule
    
    \vspace{5pt}

  \begin{tabular}{rl}
     \multicolumn{2}{l}{\textbf{PRÄSENS} }\\
    ich & \textbf{ändere}\\
    du & \textbf{änderst}\\
    er/sie/es & \textbf{ändert}\\
    wir & \textbf{ändern}\\
    ihr & \textbf{ändert}\\
    sie/Sie & \textbf{ändern}\\
  \end{tabular}
  \hfill
     \begin{tabular}{rl}
    \multicolumn{2}{l}{\textbf{PRÄTERITUM} }\\
    ich & \textbf{änderte}\\
    du & \textbf{ändertest}\\
    er/sie/es & \textbf{änderte}\\
    wir & \textbf{änderten}\\
    ihr & \textbf{ändertet}\\
    sie/Sie & \textbf{änderten}\\
  \end{tabular}
  \hfill
  \begin{tabular}{rl}
     \multicolumn{2}{l}{\textbf{IMPERATIV} }\\
   & \textbf{änder(e) (du)!}\\
   & \textbf{ändern wir!}\\
   & \textbf{ändert (ihr)!}\\
   & \textbf{ändern Sie!}\\
   & \vspace{10pt}
  \end{tabular}

    \vspace{10pt}

 \begin{tabular}{rl}
     \multicolumn{2}{l}{\textbf{PERFEKT}}\\
  &  haben + \textbf{geändert}\\
  \end{tabular}
\hfill
   \begin{tabular}{rl}
   
    \multicolumn{2}{l}{\textbf{FUTURE I} }\\
  &  werden + \textbf{ändern}\\
 
  \end{tabular}
\hfill
   \begin{tabular}{rl}
      \multicolumn{2}{l}{\textbf{PLUSQUAMPERFEKT}}\\
  &  hatten + \textbf{geändert}\\
  \end{tabular}
  

  \vspace{10pt}

\begin{tabular}{lll}
\multicolumn{3}{l}{\textbf{MIT PRÄPOSITIONEN}}\\
& ---\\
\end{tabular}
\hfill
\begin{tabular}{lll}
\multicolumn{3}{l}{\textbf{TRENNBARE VERBEN}}\\
& ---\\
\end{tabular}

  \vspace{-5pt}


\end{verbbox-acc}
%antworten (A1) – to answer
\begin{verbbox-dat}\addcontentsline{toc}{subsubsection}{antworten (to answer)}
    {\Large \bf antworten} (to answer) \hfill Type: \textbf{regular, + Dativ} \hfill Auxiliary Verb: \textbf{haben} 
    
     \vspace{5pt}

    \hrule
    
    \vspace{5pt}

  \begin{tabular}{rl}
     \multicolumn{2}{l}{\textbf{PRÄSENS} }\\
    ich & \textbf{antworte}\\
    du & \textbf{antwortest}\\
    er/sie/es & \textbf{antwortet}\\
    wir & \textbf{antworten}\\
    ihr & \textbf{antwortet}\\
    sie/Sie & \textbf{antworten}\\
  \end{tabular}
  \hfill
     \begin{tabular}{rl}
    \multicolumn{2}{l}{\textbf{PRÄTERITUM} }\\
    ich & \textbf{antwortete}\\
    du & \textbf{antwortetest}\\
    er/sie/es & \textbf{antwortete}\\
    wir & \textbf{antworteten}\\
    ihr & \textbf{antwortetet}\\
    sie/Sie & \textbf{antworteten}\\
  \end{tabular}
  \hfill
  \begin{tabular}{rl}
     \multicolumn{2}{l}{\textbf{IMPERATIV} }\\
   & \textbf{antwort(e) (du)!}\\
   & \textbf{antworten wir!}\\
   & \textbf{antwortet (ihr)!}\\
   & \textbf{antworten Sie!}\\
   & \vspace{10pt}
  \end{tabular}

    \vspace{10pt}

 \begin{tabular}{rl}
     \multicolumn{2}{l}{\textbf{PERFEKT}}\\
  &  haben + \textbf{geantwortet}\\
  \end{tabular}
\hfill
   \begin{tabular}{rl}
   
    \multicolumn{2}{l}{\textbf{FUTURE I} }\\
   & werden + \textbf{antworten}\\
 
  \end{tabular}
\hfill
   \begin{tabular}{rl}
      \multicolumn{2}{l}{\textbf{PLUSQUAMPERFEKT}}\\
   & hatten + \textbf{geantwortet}\\
  \end{tabular}


    \vspace{10pt}

\begin{tabular}{lll}
\multicolumn{3}{l}{\textbf{MIT PRÄPOSITIONEN}}\\
& \textbf{antworten auf} + Akk & (to answer something)\\
\end{tabular}
  \hfill
\begin{tabular}{lll}
\multicolumn{3}{l}{\textbf{TRENNBARE VERBEN}}\\
& ---\\
\end{tabular}

  \vspace{-5pt}

\end{verbbox-dat}
%arbeiten (A1) – to work
\begin{verbbox}\addcontentsline{toc}{subsubsection}{arbeiten (to work)}
    {\Large \bf arbeiten} (to work) \hfill Type: \textbf{regular} \hfill Auxiliary Verb: \textbf{haben} 
    
     \vspace{5pt}

    \hrule
    
    \vspace{5pt}

  \begin{tabular}{rl}
     \multicolumn{2}{l}{\textbf{PRÄSENS} }\\
    ich & \textbf{arbeite}\\
    du & \textbf{arbeitest}\\
    er/sie/es & \textbf{arbeitet}\\
    wir & \textbf{arbeiten}\\
    ihr & \textbf{arbeitet}\\
    sie/Sie & \textbf{arbeiten}\\
  \end{tabular}
  \hfill
     \begin{tabular}{rl}
    \multicolumn{2}{l}{\textbf{PRÄTERITUM} }\\
    ich & \textbf{arbeitete}\\
    du & \textbf{arbeitetest}\\
    er/sie/es & \textbf{arbeitete}\\
    wir & \textbf{arbeiteten}\\
    ihr & \textbf{arbeitetet}\\
    sie/Sie & \textbf{arbeiteten}\\
  \end{tabular}
  \hfill
  \begin{tabular}{rl}
     \multicolumn{2}{l}{\textbf{IMPERATIV} }\\
   & \textbf{arbeit(e) (du)!}\\
   & \textbf{arbeiten wir!}\\
   & \textbf{arbeitet (ihr)!}\\
   & \textbf{arbeiten Sie!}\\
   & \vspace{10pt}
  \end{tabular}

    \vspace{10pt}

 \begin{tabular}{rl}
     \multicolumn{2}{l}{\textbf{PERFEKT}}\\
   & haben + \textbf{gearbeitet}\\
  \end{tabular}
\hfill
   \begin{tabular}{rl}
   
    \multicolumn{2}{l}{\textbf{FUTURE I} }\\
   & werden + \textbf{arbeiten}\\
 
  \end{tabular}
\hfill
   \begin{tabular}{rl}
      \multicolumn{2}{l}{\textbf{PLUSQUAMPERFEKT}}\\
   & hatten + \textbf{gearbeitet}\\
  \end{tabular}
  

    \vspace{10pt}
    
   \begin{tabular}{lll}
      \multicolumn{3}{l}{\textbf{MIT PRÄPOSITIONEN}}\\
& \textbf{arbeiten an} + Dat & (to work on)\\
& \textbf{arbeiten für} + Akk & (to work for)\\
   \end{tabular}
  \hfill
   \begin{tabular}{lll}
      \multicolumn{3}{l}{\textbf{TRENNBARE VERBEN}}\\
\vspace{20pt}
  \end{tabular}


  \vspace{-5pt}

\end{verbbox}

%ärgern (über, mit)
\begin{verbbox}\addcontentsline{toc}{subsubsection}{ärgern (to annoy)}
    {\Large \bf ärgern} (to annoy) \hfill Type: \textbf{regular} \hfill Auxiliary Verb: \textbf{haben} 
    
     \vspace{5pt}

    \hrule
    
    \vspace{5pt}

  \begin{tabular}{rl}
     \multicolumn{2}{l}{\textbf{PRÄSENS} }\\
    ich & \textbf{ärgere}\\
    du & \textbf{ärgerst}\\
    er/sie/es & \textbf{ärgert}\\
    wir & \textbf{ärgern}\\
    ihr & \textbf{ärgert}\\
    sie/Sie & \textbf{ärgern}\\
  \end{tabular}
  \hfill
     \begin{tabular}{rl}
    \multicolumn{2}{l}{\textbf{PRÄTERITUM} }\\
    ich & \textbf{ärgerte}\\
    du & \textbf{ärgertest}\\
    er/sie/es & \textbf{ärgerte}\\
    wir & \textbf{ärgerten}\\
    ihr & \textbf{ärgertet}\\
    sie/Sie & \textbf{ärgerten}\\
  \end{tabular}
  \hfill
  \begin{tabular}{rl}
     \multicolumn{2}{l}{\textbf{IMPERATIV} }\\
   & \textbf{ärger(e) (du)!}\\
   & \textbf{ärgern wir!}\\
   & \textbf{ärgert (ihr)!}\\
   & \textbf{ärgern Sie!}\\
   & \vspace{10pt}
  \end{tabular}

    \vspace{10pt}

 \begin{tabular}{rl}
     \multicolumn{2}{l}{\textbf{PERFEKT}}\\
   & haben + \textbf{geärgert}\\
  \end{tabular}
\hfill
   \begin{tabular}{rl}
   
    \multicolumn{2}{l}{\textbf{FUTURE I} }\\
   & werden + \textbf{ärgern}\\
 
  \end{tabular}
\hfill
   \begin{tabular}{rl}
      \multicolumn{2}{l}{\textbf{PLUSQUAMPERFEKT}}\\
   & hatten + \textbf{geärgert}\\
  \end{tabular}

   \vspace{10pt}

   \begin{tabular}{lll}
      \multicolumn{3}{l}{\textbf{MIT PRÄPOSITIONEN}}\\
& \textbf{ärgern über} + Akk & (to be annoyed about)\\
   \end{tabular}
  \hfill
   \begin{tabular}{lll}
      \multicolumn{3}{l}{\textbf{TRENNBARE VERBEN}}\\
\vspace{10pt}
  \end{tabular}

\vspace{-5pt}

\end{verbbox}

%atmen (to breathe)
\begin{verbbox}\addcontentsline{toc}{subsubsection}{atmen (to breathe)}
{\Large \bf atmen} (to breathe) \hfill Type: \textbf{regular} \hfill Auxiliary Verb: \textbf{haben}

\vspace{5pt}
\hrule
\vspace{5pt}

\begin{tabular}{rl}
\multicolumn{2}{l}{\textbf{PRÄSENS}}\\
ich & \textbf{atme}\\
du & \textbf{atmest}\\
er/sie/es & \textbf{atmet}\\
wir & \textbf{atmen}\\
ihr & \textbf{atmet}\\
sie/Sie & \textbf{atmen}\\
\end{tabular}
\hfill
\begin{tabular}{rl}
\multicolumn{2}{l}{\textbf{PRÄTERITUM}}\\
ich & \textbf{atmete}\\
du & \textbf{atmetest}\\
er/sie/es & \textbf{atmete}\\
wir & \textbf{atmeten}\\
ihr & \textbf{atmetet}\\
sie/Sie & \textbf{atmeten}\\
\end{tabular}
\hfill
\begin{tabular}{rl}
\multicolumn{2}{l}{\textbf{IMPERATIV}}\\
& \textbf{atme (du)!}\\
& \textbf{atmen wir!}\\
& \textbf{atmet (ihr)!}\\
& \textbf{atmen Sie!}\\
\end{tabular}

\vspace{10pt}

\begin{tabular}{rl}
\multicolumn{2}{l}{\textbf{PERFEKT}}\\
& haben + \textbf{geatmet}\\
\end{tabular}
\hfill
\begin{tabular}{rl}
\multicolumn{2}{l}{\textbf{FUTURE I}}\\
& werden + \textbf{atmen}\\
\end{tabular}
\hfill
\begin{tabular}{rl}
\multicolumn{2}{l}{\textbf{PLUSQUAMPERFEKT}}\\
& hatten + \textbf{geatmet}\\
\end{tabular}

\vspace{10pt}

\begin{tabular}{lll}
\multicolumn{3}{l}{\textbf{MIT PRÄPOSITIONEN}}\\
& \textbf{atmen in} + Akk & (to breathe into)\\
& \textbf{atmen durch} + Akk & (to breathe through)\\
& \textbf{atmen mit} + Dat & (to breathe with)\\
\end{tabular}
\hfill
\begin{tabular}{lll}
\multicolumn{3}{l}{\textbf{TRENNBARE VERBEN}}\\
& \textbf{ein$|$atmen} & (to inhale)\\
& \textbf{aus$|$atmen} & (to exhale)\\
& \textbf{auf$|$atmen} & (to breathe again / breathe easier)\\
\end{tabular}
\vspace{-5pt}
\end{verbbox}

\subsection{B Verb}
%baden (to bathe)
\begin{verbbox}\addcontentsline{toc}{subsubsection}{baden (to bathe)}
{\Large \bf baden} (to bathe) \hfill Type: \textbf{regular} \hfill Auxiliary Verb: \textbf{haben}

\vspace{5pt}
\hrule
\vspace{5pt}

\begin{tabular}{rl}
\multicolumn{2}{l}{\textbf{PRÄSENS}}\\
ich & \textbf{bade}\\
du & \textbf{badest}\\
er/sie/es & \textbf{badet}\\
wir & \textbf{baden}\\
ihr & \textbf{badet}\\
sie/Sie & \textbf{baden}\\
\end{tabular}
\hfill
\begin{tabular}{rl}
\multicolumn{2}{l}{\textbf{PRÄTERITUM}}\\
ich & \textbf{badete}\\
du & \textbf{badetest}\\
er/sie/es & \textbf{badete}\\
wir & \textbf{badeten}\\
ihr & \textbf{badetet}\\
sie/Sie & \textbf{badeten}\\
\end{tabular}
\hfill
\begin{tabular}{rl}
\multicolumn{2}{l}{\textbf{IMPERATIV}}\\
& \textbf{bade (du)!}\\
& \textbf{baden wir!}\\
& \textbf{badet (ihr)!}\\
& \textbf{baden Sie!}\\
\end{tabular}

\vspace{10pt}

\begin{tabular}{rl}
\multicolumn{2}{l}{\textbf{PERFEKT}}\\
& haben + \textbf{gebadet}\\
\end{tabular}
\hfill
\begin{tabular}{rl}
\multicolumn{2}{l}{\textbf{FUTURE I}}\\
& werden + \textbf{baden}\\
\end{tabular}
\hfill
\begin{tabular}{rl}
\multicolumn{2}{l}{\textbf{PLUSQUAMPERFEKT}}\\
& hatten + \textbf{gebadet}\\
\end{tabular}

\vspace{10pt}

\begin{tabular}{lll}
\multicolumn{3}{l}{\textbf{MIT PRÄPOSITIONEN}}\\
& \textbf{baden in} + Dat & (to bathe / swim in)\\
\end{tabular}
\hfill
\begin{tabular}{lll}
\multicolumn{3}{l}{\textbf{TRENNBARE VERBEN}}\\
& ---\\
\end{tabular}
\vspace{-5pt}
\end{verbbox}

%bauen - to build
\begin{verbbox}\addcontentsline{toc}{subsubsection}{bauen (to build)}
    {\Large \bf bauen} (to build) \hfill Type: \textbf{regular} \hfill Auxiliary Verb: \textbf{haben} 
    
     \vspace{5pt}

    \hrule
    
    \vspace{5pt}

  \begin{tabular}{rl}
     \multicolumn{2}{l}{\textbf{PRÄSENS} }\\
    ich & \textbf{baue}\\
    du & \textbf{baust}\\
    er/sie/es & \textbf{baut}\\
    wir & \textbf{bauen}\\
    ihr & \textbf{baut}\\
    sie/Sie & \textbf{bauen}\\
  \end{tabular}
  \hfill
     \begin{tabular}{rl}
    \multicolumn{2}{l}{\textbf{PRÄTERITUM} }\\
    ich & \textbf{baute}\\
    du & \textbf{bautest}\\
    er/sie/es & \textbf{baute}\\
    wir & \textbf{bauten}\\
    ihr & \textbf{bautet}\\
    sie/Sie & \textbf{bauten}\\
  \end{tabular}
  \hfill
  \begin{tabular}{rl}
     \multicolumn{2}{l}{\textbf{IMPERATIV} }\\
   & \textbf{bau(e) (du)!}\\
   & \textbf{bauen wir!}\\
   & \textbf{baut (ihr)!}\\
   & \textbf{bauen Sie!}\\
   & \vspace{10pt}
  \end{tabular}

    \vspace{10pt}

 \begin{tabular}{rl}
     \multicolumn{2}{l}{\textbf{PERFEKT}}\\
   & haben + \textbf{gebaut}\\
  \end{tabular}
\hfill
   \begin{tabular}{rl}
   
    \multicolumn{2}{l}{\textbf{FUTURE I} }\\
   & werden + \textbf{bauen}\\
 
  \end{tabular}
\hfill
   \begin{tabular}{rl}
      \multicolumn{2}{l}{\textbf{PLUSQUAMPERFEKT}}\\
   & hatten + \textbf{gebaut}\\
  \end{tabular}

   \vspace{10pt}

   \begin{tabular}{lll}
      \multicolumn{3}{l}{\textbf{MIT PRÄPOSITIONEN}}\\
& \textbf{bauen auf} + Akk & (to rely on / build on)\\
& \textbf{bauen an} + Dat & (to be working on building)\\
& \textbf{bauen in} + Akk & (to install into)\\
   \end{tabular}
  \hfill
   \begin{tabular}{lll}
      \multicolumn{3}{l}{\textbf{TRENNBARE VERBEN}}\\
 & \textbf{ab$|$bauen} & (to dismantle / reduce)\\
 & \textbf{auf$|$bauen} & (to build up / establish)\\
 & \textbf{ein$|$bauen} & (to install)\\
 & \textbf{um$|$bauen} & (to rebuild / remodel)
  \end{tabular}

\vspace{-5pt}

\end{verbbox}

%beantragen (to apply for)
\begin{verbbox-acc}\addcontentsline{toc}{subsubsection}{beantragen (to apply for)}
{\Large \bf beantragen} (to apply for) \hfill Type: \textbf{regular, + Akkusativ} \hfill Auxiliary Verb: \textbf{haben}

\vspace{5pt}
\hrule
\vspace{5pt}

\begin{tabular}{rl}
\multicolumn{2}{l}{\textbf{PRÄSENS}}\\
ich & \textbf{beantrage}\\
du & \textbf{beantragst}\\
er/sie/es & \textbf{beantragt}\\
wir & \textbf{beantragen}\\
ihr & \textbf{beantragt}\\
sie/Sie & \textbf{beantragen}\\
\end{tabular}
\hfill
\begin{tabular}{rl}
\multicolumn{2}{l}{\textbf{PRÄTERITUM}}\\
ich & \textbf{beantragte}\\
du & \textbf{beantragtest}\\
er/sie/es & \textbf{beantragte}\\
wir & \textbf{beantragten}\\
ihr & \textbf{beantragtet}\\
sie/Sie & \textbf{beantragten}\\
\end{tabular}
\hfill
\begin{tabular}{rl}
\multicolumn{2}{l}{\textbf{IMPERATIV}}\\
& \textbf{beantrage (du)!}\\
& \textbf{beantragen wir!}\\
& \textbf{beantragt (ihr)!}\\
& \textbf{beantragen Sie!}\\
\end{tabular}

\vspace{10pt}

\begin{tabular}{rl}
\multicolumn{2}{l}{\textbf{PERFEKT}}\\
& haben + \textbf{beantragt}\\
\end{tabular}
\hfill
\begin{tabular}{rl}
\multicolumn{2}{l}{\textbf{FUTURE I}}\\
& werden + \textbf{beantragen}\\
\end{tabular}
\hfill
\begin{tabular}{rl}
\multicolumn{2}{l}{\textbf{PLUSQUAMPERFEKT}}\\
& hatten + \textbf{beantragt}\\
\end{tabular}

\vspace{10pt}

\begin{tabular}{lll}
\multicolumn{3}{l}{\textbf{MIT PRÄPOSITIONEN}}\\
& \textbf{beantragen bei} + Dat & (to apply at)\\
\end{tabular}
\hfill
\begin{tabular}{lll}
\multicolumn{3}{l}{\textbf{TRENNBARE VERBEN}}\\
& ---\\
\end{tabular}
\vspace{-5pt}
\end{verbbox-acc}

%bedeuten - to mean something
\begin{verbbox-acc}\addcontentsline{toc}{subsubsection}{bedeuten (to mean something)}
    {\Large \bf bedeuten} (to mean something) \hfill Type: \textbf{regular, + Akkusativ} \hfill Auxiliary Verb: \textbf{haben} 
    
     \vspace{5pt}

    \hrule
    
    \vspace{5pt}

  \begin{tabular}{rl}
     \multicolumn{2}{l}{\textbf{PRÄSENS} }\\
    ich & \textbf{bedeute}\\
    du & \textbf{bedeutest}\\
    er/sie/es & \textbf{bedeutet}\\
    wir & \textbf{bedeuten}\\
    ihr & \textbf{bedeutet}\\
    sie/Sie & \textbf{bedeuten}\\
  \end{tabular}
  \hfill
     \begin{tabular}{rl}
    \multicolumn{2}{l}{\textbf{PRÄTERITUM} }\\
    ich & \textbf{bedeutete}\\
    du & \textbf{bedeutetest}\\
    er/sie/es & \textbf{bedeutete}\\
    wir & \textbf{bedeuteten}\\
    ihr & \textbf{bedeutetet}\\
    sie/Sie & \textbf{bedeuteten}\\
  \end{tabular}
  \hfill
  \begin{tabular}{rl}
     \multicolumn{2}{l}{\textbf{IMPERATIV} }\\
   & \textbf{bedeut(e) (du)!}\\
   & \textbf{bedeuten wir!}\\
   & \textbf{bedeutet (ihr)!}\\
   & \textbf{bedeuten Sie!}\\
   & \vspace{10pt}
  \end{tabular}

    \vspace{10pt}

 \begin{tabular}{rl}
     \multicolumn{2}{l}{\textbf{PERFEKT}}\\
   & haben + \textbf{bedeutet}\\
  \end{tabular}
\hfill
   \begin{tabular}{rl}
   
    \multicolumn{2}{l}{\textbf{FUTURE I} }\\
   & werden + \textbf{bedeuten}\\
 
  \end{tabular}
\hfill
   \begin{tabular}{rl}
      \multicolumn{2}{l}{\textbf{PLUSQUAMPERFEKT}}\\
   & hatten + \textbf{bedeutet}\\
  \end{tabular}
  
\vspace{10pt}

\begin{tabular}{lll}
\multicolumn{3}{l}{\textbf{MIT PRÄPOSITIONEN}}\\
& ---\\
\end{tabular}
\hfill
\begin{tabular}{lll}
\multicolumn{3}{l}{\textbf{TRENNBARE VERBEN}}\\
& ---\\
\end{tabular}

\vspace{-5pt}


\end{verbbox-acc}

%beginnen (A1) – to begin
\begin{verbbox-acc}\addcontentsline{toc}{subsubsection}{beginnen (to begin)}
    {\Large \bf beginnen} (to begin) \hfill Type: \textbf{irregular, + Akkusativ} \hfill Auxiliary Verb: \textbf{haben} 
    
     \vspace{5pt}

    \hrule
    
    \vspace{5pt}

  \begin{tabular}{rl}
     \multicolumn{2}{l}{\textbf{PRÄSENS} }\\
    ich & \textbf{beginne}\\
    du & \textbf{beginnst}\\
    er/sie/es & \textbf{beginnt}\\
    wir & \textbf{beginnen}\\
    ihr & \textbf{beginnt}\\
    sie/Sie & \textbf{beginnen}\\
  \end{tabular}
  \hfill
     \begin{tabular}{rl}
    \multicolumn{2}{l}{\textbf{PRÄTERITUM} }\\
    ich & \textbf{begann}\\
    du & \textbf{begannst}\\
    er/sie/es & \textbf{begann}\\
    wir & \textbf{begannen}\\
    ihr & \textbf{begannt}\\
    sie/Sie & \textbf{begannen}\\
  \end{tabular}
  \hfill
  \begin{tabular}{rl}
     \multicolumn{2}{l}{\textbf{IMPERATIV} }\\
   & \textbf{beginn(e) (du)!}\\
   & \textbf{beginnen wir!}\\
   & \textbf{beginnt (ihr)!}\\
   & \textbf{beginnen Sie!}\\
   & \vspace{10pt}
  \end{tabular}

    \vspace{10pt}

 \begin{tabular}{rl}
     \multicolumn{2}{l}{\textbf{PERFEKT}}\\
  &  haben + \textbf{begonnen}\\
  \end{tabular}
\hfill
   \begin{tabular}{rl}
   
    \multicolumn{2}{l}{\textbf{FUTURE I} }\\
  &  werden + \textbf{beginnen}\\
 
  \end{tabular}
\hfill
   \begin{tabular}{rl}
      \multicolumn{2}{l}{\textbf{PLUSQUAMPERFEKT}}\\
   & hatten + \textbf{begonnen}\\
  \end{tabular}
  
\vspace{10pt}

\begin{tabular}{lll}
\multicolumn{3}{l}{\textbf{MIT PRÄPOSITIONEN}}\\
& ---\\
\end{tabular}
\hfill
\begin{tabular}{lll}
\multicolumn{3}{l}{\textbf{TRENNBARE VERBEN}}\\
& ---\\
\end{tabular}

\vspace{-5pt}

\end{verbbox-acc}

%begrüßen (to greet)
\begin{verbbox-acc}\addcontentsline{toc}{subsubsection}{begrüßen (to greet)}
{\Large \bf begrüßen} (to greet) \hfill Type: \textbf{regular, + Akkusativ} \hfill Auxiliary Verb: \textbf{haben}

\vspace{5pt}
\hrule
\vspace{5pt}

\begin{tabular}{rl}
\multicolumn{2}{l}{\textbf{PRÄSENS}}\\
ich & \textbf{begrüße}\\
du & \textbf{begrüßt}\\
er/sie/es & \textbf{begrüßt}\\
wir & \textbf{begrüßen}\\
ihr & \textbf{begrüßt}\\
sie/Sie & \textbf{begrüßen}\\
\end{tabular}
\hfill
\begin{tabular}{rl}
\multicolumn{2}{l}{\textbf{PRÄTERITUM}}\\
ich & \textbf{begrüßte}\\
du & \textbf{begrüßtest}\\
er/sie/es & \textbf{begrüßte}\\
wir & \textbf{begrüßten}\\
ihr & \textbf{begrüßtet}\\
sie/Sie & \textbf{begrüßten}\\
\end{tabular}
\hfill
\begin{tabular}{rl}
\multicolumn{2}{l}{\textbf{IMPERATIV}}\\
& \textbf{begrüße (du)!}\\
& \textbf{begrüßen wir!}\\
& \textbf{begrüßt (ihr)!}\\
& \textbf{begrüßen Sie!}\\
\end{tabular}

\vspace{10pt}

\begin{tabular}{rl}
\multicolumn{2}{l}{\textbf{PERFEKT}}\\
& haben + \textbf{begrüßt}\\
\end{tabular}
\hfill
\begin{tabular}{rl}
\multicolumn{2}{l}{\textbf{FUTURE I}}\\
& werden + \textbf{begrüßen}\\
\end{tabular}
\hfill
\begin{tabular}{rl}
\multicolumn{2}{l}{\textbf{PLUSQUAMPERFEKT}}\\
& hatten + \textbf{begrüßt}\\
\end{tabular}

\vspace{10pt}

\begin{tabular}{lll}
\multicolumn{3}{l}{\textbf{MIT PRÄPOSITIONEN}}\\
& ---\\
\end{tabular}
\hfill
\begin{tabular}{lll}
\multicolumn{3}{l}{\textbf{TRENNBARE VERBEN}}\\
& ---\\
\end{tabular}

\vspace{-5pt}
\end{verbbox-acc}

%behandeln (to treat)
\begin{verbbox}\addcontentsline{toc}{subsubsection}{behandeln (to treat)}
    {\Large \bf behandeln} (to treat) \hfill Type: \textbf{regular} \hfill Auxiliary Verb: \textbf{haben} 
    
     \vspace{5pt}

    \hrule
    
    \vspace{5pt}

  \begin{tabular}{rl}
     \multicolumn{2}{l}{\textbf{PRÄSENS} }\\
    ich & \textbf{behandle}\\
    du & \textbf{behandelst}\\
    er/sie/es & \textbf{behandelt}\\
    wir & \textbf{behandeln}\\
    ihr & \textbf{behandelt}\\
    sie/Sie & \textbf{behandeln}\\
  \end{tabular}
  \hfill
     \begin{tabular}{rl}
    \multicolumn{2}{l}{\textbf{PRÄTERITUM} }\\
    ich & \textbf{behandelte}\\
    du & \textbf{behandeltest}\\
    er/sie/es & \textbf{behandelte}\\
    wir & \textbf{behandelten}\\
    ihr & \textbf{behandeltet}\\
    sie/Sie & \textbf{behandelten}\\
  \end{tabular}
  \hfill
  \begin{tabular}{rl}
     \multicolumn{2}{l}{\textbf{IMPERATIV} }\\
   & \textbf{behandle (du)!}\\
   & \textbf{behandeln wir!}\\
   & \textbf{behandelt (ihr)!}\\
   & \textbf{behandeln Sie!}\\
   & \vspace{10pt}
  \end{tabular}

    \vspace{10pt}

 \begin{tabular}{rl}
     \multicolumn{2}{l}{\textbf{PERFEKT}}\\
   & haben + \textbf{behandelt}\\
  \end{tabular}
\hfill
   \begin{tabular}{rl}
    \multicolumn{2}{l}{\textbf{FUTURE I} }\\
   & werde + \textbf{behandeln}\\
  \end{tabular}
\hfill
   \begin{tabular}{rl}
      \multicolumn{2}{l}{\textbf{PLUSQUAMPERFEKT}}\\
   & hatten + \textbf{behandelt}\\
  \end{tabular}

   \vspace{10pt}

   \begin{tabular}{lll}
      \multicolumn{3}{l}{\textbf{MIT PRÄPOSITIONEN}}\\
& \textbf{behandeln mit} + Dat & (to treat with)\\
& \textbf{behandeln wegen} + Gen & (to treat because of / for)\\
\end{tabular}
  \hfill
   \begin{tabular}{lll}
      \multicolumn{3}{l}{\textbf{TRENNBARE VERBEN}}\\
 & \textbf{---} & (none common)\\
\vspace{10pt}
  \end{tabular}

  \vspace{-5pt}

\end{verbbox}

%beißen (to bite)
\begin{verbbox}\addcontentsline{toc}{subsubsection}{beißen (to bite)}
    {\Large \bf beißen} (to bite) \hfill Type: \textbf{irregular, strong} \hfill Auxiliary Verb: \textbf{haben} 
    
     \vspace{5pt}

    \hrule
    
    \vspace{5pt}

  \begin{tabular}{rl}
     \multicolumn{2}{l}{\textbf{PRÄSENS} }\\
    ich & \textbf{beiße}\\
    du & \textbf{beißt}\\
    er/sie/es & \textbf{beißt}\\
    wir & \textbf{beißen}\\
    ihr & \textbf{beißt}\\
    sie/Sie & \textbf{beißen}\\
  \end{tabular}
  \hfill
     \begin{tabular}{rl}
    \multicolumn{2}{l}{\textbf{PRÄTERITUM} }\\
    ich & \textbf{biss}\\
    du & \textbf{bissest}\\
    er/sie/es & \textbf{biss}\\
    wir & \textbf{bissen}\\
    ihr & \textbf{bisst}\\
    sie/Sie & \textbf{bissen}\\
  \end{tabular}
  \hfill
  \begin{tabular}{rl}
     \multicolumn{2}{l}{\textbf{IMPERATIV} }\\
   & \textbf{beiß(e) (du)!}\\
   & \textbf{beißen wir!}\\
   & \textbf{beißt (ihr)!}\\
   & \textbf{beißen Sie!}\\
   & \vspace{10pt}
  \end{tabular}

    \vspace{10pt}

 \begin{tabular}{rl}
     \multicolumn{2}{l}{\textbf{PERFEKT}}\\
   & haben + \textbf{gebissen}\\
  \end{tabular}
\hfill
   \begin{tabular}{rl}
    \multicolumn{2}{l}{\textbf{FUTURE I} }\\
   & werde + \textbf{beißen}\\
  \end{tabular}
\hfill
   \begin{tabular}{rl}
      \multicolumn{2}{l}{\textbf{PLUSQUAMPERFEKT}}\\
   & hatten + \textbf{gebissen}\\
  \end{tabular}

   \vspace{10pt}

   \begin{tabular}{lll}
      \multicolumn{3}{l}{\textbf{MIT PRÄPOSITIONEN}}\\
& \textbf{beißen in} + Akk & (to bite into)\\
& \textbf{beißen nach} + Dat & (to snap at / bite at)\\
\end{tabular}
  \hfill
   \begin{tabular}{lll}
      \multicolumn{3}{l}{\textbf{TRENNBARE VERBEN}}\\
 & \textbf{ab$|$beißen} & (to bite off)\\
 & \textbf{zu$|$beißen} & (to bite down / snap)\\
\vspace{10pt}
  \end{tabular}

  \vspace{-5pt}

\end{verbbox}

%bekommen (A1) – to receive
\begin{verbbox-acc}\addcontentsline{toc}{subsubsection}{bekommen (to receive)}
    {\Large \bf bekommen} (to receive) \hfill Type: \textbf{irregular, + Akkusativ} \hfill Auxiliary Verb: \textbf{haben} 
    
     \vspace{5pt}

    \hrule
    
    \vspace{5pt}

  \begin{tabular}{rl}
     \multicolumn{2}{l}{\textbf{PRÄSENS} }\\
    ich & \textbf{bekomme}\\
    du & \textbf{bekommst}\\
    er/sie/es & \textbf{bekommt}\\
    wir & \textbf{bekommen}\\
    ihr & \textbf{bekommt}\\
    sie/Sie & \textbf{bekommen}\\
  \end{tabular}
  \hfill
     \begin{tabular}{rl}
    \multicolumn{2}{l}{\textbf{PRÄTERITUM} }\\
    ich & \textbf{bekam}\\
    du & \textbf{bekamst}\\
    er/sie/es & \textbf{bekam}\\
    wir & \textbf{bekamen}\\
    ihr & \textbf{bekamt}\\
    sie/Sie & \textbf{bekamen}\\
  \end{tabular}
  \hfill
  \begin{tabular}{rl}
     \multicolumn{2}{l}{\textbf{IMPERATIV} }\\
   & \textbf{bekomm(e) (du)!}\\
   & \textbf{bekommen wir!}\\
   & \textbf{bekommt (ihr)!}\\
   & \textbf{bekommen Sie!}\\
   & \vspace{10pt}
  \end{tabular}

    \vspace{10pt}

 \begin{tabular}{rl}
     \multicolumn{2}{l}{\textbf{PERFEKT}}\\
  &  haben + \textbf{bekommen}\\
  \end{tabular}
\hfill
   \begin{tabular}{rl}
   
    \multicolumn{2}{l}{\textbf{FUTURE I} }\\
  &  werden + \textbf{bekommen}\\
 
  \end{tabular}
\hfill
   \begin{tabular}{rl}
      \multicolumn{2}{l}{\textbf{PLUSQUAMPERFEKT}}\\
  &  hatten + \textbf{bekommen}\\
  \end{tabular}
  
\vspace{10pt}

\begin{tabular}{lll}
\multicolumn{3}{l}{\textbf{MIT PRÄPOSITIONEN}}\\
& ---\\
\end{tabular}
\hfill
\begin{tabular}{lll}
\multicolumn{3}{l}{\textbf{TRENNBARE VERBEN}}\\
& ---\\
\end{tabular}

\vspace{-5pt}


\end{verbbox-acc}

%belasten (to burden / to strain)
\begin{verbbox}\addcontentsline{toc}{subsubsection}{belasten (to burden / to strain)}
    {\Large \bf belasten} (to burden / to strain) \hfill Type: \textbf{regular} \hfill Auxiliary Verb: \textbf{haben} 
    
     \vspace{5pt}

    \hrule
    
    \vspace{5pt}

  \begin{tabular}{rl}
     \multicolumn{2}{l}{\textbf{PRÄSENS} }\\
    ich & \textbf{belaste}\\
    du & \textbf{belastest}\\
    er/sie/es & \textbf{belastet}\\
    wir & \textbf{belasten}\\
    ihr & \textbf{belastet}\\
    sie/Sie & \textbf{belasten}\\
  \end{tabular}
  \hfill
     \begin{tabular}{rl}
    \multicolumn{2}{l}{\textbf{PRÄTERITUM} }\\
    ich & \textbf{belastete}\\
    du & \textbf{belastetest}\\
    er/sie/es & \textbf{belastete}\\
    wir & \textbf{belasteten}\\
    ihr & \textbf{belastetet}\\
    sie/Sie & \textbf{belasteten}\\
  \end{tabular}
  \hfill
  \begin{tabular}{rl}
     \multicolumn{2}{l}{\textbf{IMPERATIV} }\\
   & \textbf{belaste (du)!}\\
   & \textbf{belasten wir!}\\
   & \textbf{belastet (ihr)!}\\
   & \textbf{belasten Sie!}\\
   & \vspace{10pt}
  \end{tabular}

    \vspace{10pt}

 \begin{tabular}{rl}
     \multicolumn{2}{l}{\textbf{PERFEKT}}\\
   & haben + \textbf{belastet}\\
  \end{tabular}
\hfill
   \begin{tabular}{rl}
    \multicolumn{2}{l}{\textbf{FUTURE I} }\\
   & werde + \textbf{belasten}\\
  \end{tabular}
\hfill
   \begin{tabular}{rl}
      \multicolumn{2}{l}{\textbf{PLUSQUAMPERFEKT}}\\
   & hatten + \textbf{belastet}\\
  \end{tabular}

   \vspace{10pt}

   \begin{tabular}{lll}
      \multicolumn{3}{l}{\textbf{MIT PRÄPOSITIONEN}}\\
& \textbf{belasten mit} + Dat & (to burden with)\\
& \textbf{belasten durch} + Akk & (to strain/burden through)\\
\end{tabular}
  \hfill
   \begin{tabular}{lll}
      \multicolumn{3}{l}{\textbf{TRENNBARE VERBEN}}\\
 & \textbf{---} & (none common)\\
\vspace{10pt}
  \end{tabular}

  \vspace{-5pt}

\end{verbbox}

%belegen (to occupy / register / cover)
\begin{verbbox-acc}\addcontentsline{toc}{subsubsection}{belegen (to occupy / register / cover)}
{\Large \bf belegen} (to occupy / register / cover) \hfill Type: \textbf{regular, + Akkusativ} \hfill Auxiliary Verb: \textbf{haben}

\vspace{5pt}
\hrule
\vspace{5pt}

\begin{tabular}{rl}
\multicolumn{2}{l}{\textbf{PRÄSENS}}\\
ich & \textbf{belege}\\
du & \textbf{belegst}\\
er/sie/es & \textbf{belegt}\\
wir & \textbf{belegen}\\
ihr & \textbf{belegt}\\
sie/Sie & \textbf{belegen}\\
\end{tabular}
\hfill
\begin{tabular}{rl}
\multicolumn{2}{l}{\textbf{PRÄTERITUM}}\\
ich & \textbf{belegte}\\
du & \textbf{belegtest}\\
er/sie/es & \textbf{belegte}\\
wir & \textbf{belegten}\\
ihr & \textbf{belegtet}\\
sie/Sie & \textbf{belegten}\\
\end{tabular}
\hfill
\begin{tabular}{rl}
\multicolumn{2}{l}{\textbf{IMPERATIV}}\\
& \textbf{belege (du)!}\\
& \textbf{belegen wir!}\\
& \textbf{belegt (ihr)!}\\
& \textbf{belegen Sie!}\\
\end{tabular}

\vspace{10pt}

\begin{tabular}{rl}
\multicolumn{2}{l}{\textbf{PERFEKT}}\\
& haben + \textbf{belegt}\\
\end{tabular}
\hfill
\begin{tabular}{rl}
\multicolumn{2}{l}{\textbf{FUTURE I}}\\
& werden + \textbf{belegen}\\
\end{tabular}
\hfill
\begin{tabular}{rl}
\multicolumn{2}{l}{\textbf{PLUSQUAMPERFEKT}}\\
& hatten + \textbf{belegt}\\
\end{tabular}

\vspace{10pt}

\begin{tabular}{lll}
\multicolumn{3}{l}{\textbf{MIT PRÄPOSITIONEN}}\\
& \textbf{belegen mit} + Dat & (to cover with)\\
\end{tabular}
\hfill
\begin{tabular}{lll}
\multicolumn{3}{l}{\textbf{TRENNBARE VERBEN}}\\
& ---\\
\end{tabular}
\vspace{-5pt}
\end{verbbox-acc}

%bemerken (to notice)
\begin{verbbox-acc}\addcontentsline{toc}{subsubsection}{bemerken (to notice)}
{\Large \bf bemerken} (to notice) \hfill Type: \textbf{regular, + Akkusativ} \hfill Auxiliary Verb: \textbf{haben}

\vspace{5pt}
\hrule
\vspace{5pt}

\begin{tabular}{rl}
\multicolumn{2}{l}{\textbf{PRÄSENS}}\\
ich & \textbf{bemerke}\\
du & \textbf{bemerkst}\\
er/sie/es & \textbf{bemerkt}\\
wir & \textbf{bemerken}\\
ihr & \textbf{bemerkt}\\
sie/Sie & \textbf{bemerken}\\
\end{tabular}
\hfill
\begin{tabular}{rl}
\multicolumn{2}{l}{\textbf{PRÄTERITUM}}\\
ich & \textbf{bemerkte}\\
du & \textbf{bemerktest}\\
er/sie/es & \textbf{bemerkte}\\
wir & \textbf{bemerkten}\\
ihr & \textbf{bemerktet}\\
sie/Sie & \textbf{bemerkten}\\
\end{tabular}
\hfill
\begin{tabular}{rl}
\multicolumn{2}{l}{\textbf{IMPERATIV}}\\
& \textbf{bemerke (du)!}\\
& \textbf{bemerken wir!}\\
& \textbf{bemerkt (ihr)!}\\
& \textbf{bemerken Sie!}\\
\end{tabular}

\vspace{10pt}

\begin{tabular}{rl}
\multicolumn{2}{l}{\textbf{PERFEKT}}\\
& haben + \textbf{bemerkt}\\
\end{tabular}
\hfill
\begin{tabular}{rl}
\multicolumn{2}{l}{\textbf{FUTURE I}}\\
& werden + \textbf{bemerken}\\
\end{tabular}
\hfill
\begin{tabular}{rl}
\multicolumn{2}{l}{\textbf{PLUSQUAMPERFEKT}}\\
& hatten + \textbf{bemerkt}\\
\end{tabular}

\vspace{10pt}

\begin{tabular}{lll}
\multicolumn{3}{l}{\textbf{MIT PRÄPOSITIONEN}}\\
& ---\\
\end{tabular}
\hfill
\begin{tabular}{lll}
\multicolumn{3}{l}{\textbf{TRENNBARE VERBEN}}\\
& ---\\
\end{tabular}

\vspace{-5pt}
\end{verbbox-acc}

%benutzen (A2) – to use
\begin{verbbox-acc}\addcontentsline{toc}{subsubsection}{benutzen (to use)}
    {\Large \bf benutzen} (to use) \hfill Type: \textbf{regular, + Akkusativ} \hfill Auxiliary Verb: \textbf{haben} 
    
     \vspace{5pt}

    \hrule
    
    \vspace{5pt}

  \begin{tabular}{rl}
     \multicolumn{2}{l}{\textbf{PRÄSENS} }\\
    ich & \textbf{benutze}\\
    du & \textbf{benutzt}\\
    er/sie/es & \textbf{benutz}\\
    wir & \textbf{benutzen}\\
    ihr & \textbf{benutzt}\\
    sie/Sie & \textbf{benutzen}\\
  \end{tabular}
  \hfill
     \begin{tabular}{rl}
    \multicolumn{2}{l}{\textbf{PRÄTERITUM} }\\
    ich & \textbf{benutzte}\\
    du & \textbf{benutztest}\\
    er/sie/es & \textbf{benutzte}\\
    wir & \textbf{benutzten}\\
    ihr & \textbf{benutztet}\\
    sie/Sie & \textbf{benutzten}\\
  \end{tabular}
  \hfill
  \begin{tabular}{rl}
     \multicolumn{2}{l}{\textbf{IMPERATIV} }\\
   & \textbf{benutz(e) (du)!}\\
   & \textbf{benutzen wir!}\\
   & \textbf{benutzt (ihr)!}\\
   & \textbf{benutzen Sie!}\\
   & \vspace{10pt}
  \end{tabular}

    \vspace{10pt}

 \begin{tabular}{rl}
     \multicolumn{2}{l}{\textbf{PERFEKT}}\\
   & haben + \textbf{benutzt}\\
  \end{tabular}
\hfill
   \begin{tabular}{rl}
   
    \multicolumn{2}{l}{\textbf{FUTURE I} }\\
   & werden + \textbf{benutzen}\\
 
  \end{tabular}
\hfill
   \begin{tabular}{rl}
      \multicolumn{2}{l}{\textbf{PLUSQUAMPERFEKT}}\\
   & hatten + \textbf{benutzt}\\
  \end{tabular}
  
  \vspace{-5pt}\vspace{10pt}

\begin{tabular}{lll}
\multicolumn{3}{l}{\textbf{MIT PRÄPOSITIONEN}}\\
& ---\\
\end{tabular}
\hfill
\begin{tabular}{lll}
\multicolumn{3}{l}{\textbf{TRENNBARE VERBEN}}\\
& ---\\
\end{tabular}

\vspace{-5pt}

\end{verbbox-acc}

%beobachten (to observe)
\begin{verbbox-acc}\addcontentsline{toc}{subsubsection}{beobachten (to observe)}
{\Large \bf beobachten} (to observe) \hfill Type: \textbf{regular, + Akkusativ} \hfill Auxiliary Verb: \textbf{haben} 

\vspace{5pt}
\hrule
\vspace{5pt}

\begin{tabular}{rl}
\multicolumn{2}{l}{\textbf{PRÄSENS}}\\
ich & \textbf{beobachte}\\
du & \textbf{beobachtest}\\
er/sie/es & \textbf{beobachtet}\\
wir & \textbf{beobachten}\\
ihr & \textbf{beobachtet}\\
sie/Sie & \textbf{beobachten}\\
\end{tabular}
\hfill
\begin{tabular}{rl}
\multicolumn{2}{l}{\textbf{PRÄTERITUM}}\\
ich & \textbf{beobachtete}\\
du & \textbf{beobachtetest}\\
er/sie/es & \textbf{beobachtete}\\
wir & \textbf{beobachteten}\\
ihr & \textbf{beobachtetet}\\
sie/Sie & \textbf{beobachteten}\\
\end{tabular}
\hfill
\begin{tabular}{rl}
\multicolumn{2}{l}{\textbf{IMPERATIV}}\\
& \textbf{beobachte (du)!}\\
& \textbf{beobachten wir!}\\
& \textbf{beobachtet (ihr)!}\\
& \textbf{beobachten Sie!}\\
\end{tabular}

\vspace{10pt}

\begin{tabular}{rl}
\multicolumn{2}{l}{\textbf{PERFEKT}}\\
& haben + \textbf{beobachtet}\\
\end{tabular}
\hfill
\begin{tabular}{rl}
\multicolumn{2}{l}{\textbf{FUTURE I}}\\
& werden + \textbf{beobachten}\\
\end{tabular}
\hfill
\begin{tabular}{rl}
\multicolumn{2}{l}{\textbf{PLUSQUAMPERFEKT}}\\
& hatten + \textbf{beobachtet}\\
\end{tabular}

\vspace{10pt}

\begin{tabular}{lll}
\multicolumn{3}{l}{\textbf{MIT PRÄPOSITIONEN}}\\
& \textbf{beobachten Akk} & (to observe something)\\
\end{tabular}
\hfill
\begin{tabular}{lll}
\multicolumn{3}{l}{\textbf{TRENNBARE VERBEN}}\\
& ---\\
\end{tabular}

\vspace{-5pt}
\end{verbbox-acc}

%beraten (to advise)
\begin{verbbox-acc}\addcontentsline{toc}{subsubsection}{beraten (to advise)}
{\Large \bf beraten} (to advise) \hfill Type: \textbf{irregular, + Akkusativ} \hfill Auxiliary Verb: \textbf{haben}

\vspace{5pt}
\hrule
\vspace{5pt}

\begin{tabular}{rl}
\multicolumn{2}{l}{\textbf{PRÄSENS}}\\
ich & \textbf{berate}\\
du & \textbf{berätst}\\
er/sie/es & \textbf{berät}\\
wir & \textbf{beraten}\\
ihr & \textbf{beratet}\\
sie/Sie & \textbf{beraten}\\
\end{tabular}
\hfill
\begin{tabular}{rl}
\multicolumn{2}{l}{\textbf{PRÄTERITUM}}\\
ich & \textbf{beriet}\\
du & \textbf{berietst}\\
er/sie/es & \textbf{beriet}\\
wir & \textbf{berieten}\\
ihr & \textbf{berietet}\\
sie/Sie & \textbf{berieten}\\
\end{tabular}
\hfill
\begin{tabular}{rl}
\multicolumn{2}{l}{\textbf{IMPERATIV}}\\
& \textbf{berate (du)!}\\
& \textbf{beraten wir!}\\
& \textbf{beratet (ihr)!}\\
& \textbf{beraten Sie!}\\
\end{tabular}

\vspace{10pt}

\begin{tabular}{rl}
\multicolumn{2}{l}{\textbf{PERFEKT}}\\
& haben + \textbf{beraten}\\
\end{tabular}
\hfill
\begin{tabular}{rl}
\multicolumn{2}{l}{\textbf{FUTURE I}}\\
& werden + \textbf{beraten}\\
\end{tabular}
\hfill
\begin{tabular}{rl}
\multicolumn{2}{l}{\textbf{PLUSQUAMPERFEKT}}\\
& hatten + \textbf{beraten}\\
\end{tabular}

\vspace{10pt}

\begin{tabular}{lll}
\multicolumn{3}{l}{\textbf{MIT PRÄPOSITIONEN}}\\
& \textbf{beraten über} + Akk & (to advise about)\\
& \textbf{beraten bei} + Dat & (to advise during / assist with)\\
\end{tabular}
\hfill
\begin{tabular}{lll}
\multicolumn{3}{l}{\textbf{TRENNBARE VERBEN}}\\
& --- & \\
\vspace{10pt}
\end{tabular}

\vspace{-5pt}

\end{verbbox-acc}

%bereiten (to prepare / cause)
\begin{verbbox-acc}\addcontentsline{toc}{subsubsection}{bereiten (to prepare / cause)}
{\Large \bf bereiten} (to prepare / cause) \hfill Type: \textbf{regular, + Akkusativ} \hfill Auxiliary Verb: \textbf{haben}

\vspace{5pt}
\hrule
\vspace{5pt}

\begin{tabular}{rl}
\multicolumn{2}{l}{\textbf{PRÄSENS}}\\
ich & \textbf{bereite}\\
du & \textbf{bereitest}\\
er/sie/es & \textbf{bereitet}\\
wir & \textbf{bereiten}\\
ihr & \textbf{bereitet}\\
sie/Sie & \textbf{bereiten}\\
\end{tabular}
\hfill
\begin{tabular}{rl}
\multicolumn{2}{l}{\textbf{PRÄTERITUM}}\\
ich & \textbf{bereitete}\\
du & \textbf{bereitetest}\\
er/sie/es & \textbf{bereitete}\\
wir & \textbf{bereiteten}\\
ihr & \textbf{bereitetet}\\
sie/Sie & \textbf{bereiteten}\\
\end{tabular}
\hfill
\begin{tabular}{rl}
\multicolumn{2}{l}{\textbf{IMPERATIV}}\\
& \textbf{bereite (du)!}\\
& \textbf{bereiten wir!}\\
& \textbf{bereitet (ihr)!}\\
& \textbf{bereiten Sie!}\\
\end{tabular}

\vspace{10pt}

\begin{tabular}{rl}
\multicolumn{2}{l}{\textbf{PERFEKT}}\\
& haben + \textbf{bereitet}\\
\end{tabular}
\hfill
\begin{tabular}{rl}
\multicolumn{2}{l}{\textbf{FUTURE I}}\\
& werden + \textbf{bereiten}\\
\end{tabular}
\hfill
\begin{tabular}{rl}
\multicolumn{2}{l}{\textbf{PLUSQUAMPERFEKT}}\\
& hatten + \textbf{bereitet}\\
\end{tabular}

\vspace{10pt}

\begin{tabular}{lll}
\multicolumn{3}{l}{\textbf{MIT PRÄPOSITIONEN}}\\
& \textbf{bereiten auf} + Akk & (to prepare for)\\
\end{tabular}
\hfill
\begin{tabular}{lll}
\multicolumn{3}{l}{\textbf{TRENNBARE VERBEN}}\\
& \textbf{vor$|$bereiten} & (to prepare)\\
\end{tabular}

\vspace{-5pt}
\end{verbbox-acc}

%besetzen (to occupy / fill a position)
\begin{verbbox-acc}\addcontentsline{toc}{subsubsection}{besetzen (to occupy / fill a position)}
{\Large \bf besetzen} (to occupy / fill a position) \hfill Type: \textbf{regular, + Akkusativ} \hfill Auxiliary Verb: \textbf{haben}

\vspace{5pt}
\hrule
\vspace{5pt}

\begin{tabular}{rl}
\multicolumn{2}{l}{\textbf{PRÄSENS}}\\
ich & \textbf{besetze}\\
du & \textbf{besetzt}\\
er/sie/es & \textbf{besetzt}\\
wir & \textbf{besetzen}\\
ihr & \textbf{besetzt}\\
sie/Sie & \textbf{besetzen}\\
\end{tabular}
\hfill
\begin{tabular}{rl}
\multicolumn{2}{l}{\textbf{PRÄTERITUM}}\\
ich & \textbf{besetzte}\\
du & \textbf{besetztest}\\
er/sie/es & \textbf{besetzte}\\
wir & \textbf{besetzten}\\
ihr & \textbf{besetztet}\\
sie/Sie & \textbf{besetzten}\\
\end{tabular}
\hfill
\begin{tabular}{rl}
\multicolumn{2}{l}{\textbf{IMPERATIV}}\\
& \textbf{besetze (du)!}\\
& \textbf{besetzen wir!}\\
& \textbf{besetzt (ihr)!}\\
& \textbf{besetzen Sie!}\\
\end{tabular}

\vspace{10pt}

\begin{tabular}{rl}
\multicolumn{2}{l}{\textbf{PERFEKT}}\\
& haben + \textbf{besetzt}\\
\end{tabular}
\hfill
\begin{tabular}{rl}
\multicolumn{2}{l}{\textbf{FUTURE I}}\\
& werden + \textbf{besetzen}\\
\end{tabular}
\hfill
\begin{tabular}{rl}
\multicolumn{2}{l}{\textbf{PLUSQUAMPERFEKT}}\\
& hatten + \textbf{besetzt}\\
\end{tabular}

\vspace{10pt}

\begin{tabular}{lll}
\multicolumn{3}{l}{\textbf{MIT PRÄPOSITIONEN}}\\
& \textbf{besetzen mit} + Dat & (to fill with / occupy with)\\
\end{tabular}
\hfill
\begin{tabular}{lll}
\multicolumn{3}{l}{\textbf{TRENNBARE VERBEN}}\\
& ---\\
\end{tabular}
\vspace{-5pt}
\end{verbbox-acc}

%beschließen (to decide)
\begin{verbbox-acc}\addcontentsline{toc}{subsubsection}{beschließen (to decide)}
{\Large \bf beschließen} (to decide) \hfill Type: \textbf{irregular, + Akkusativ} \hfill Auxiliary Verb: \textbf{haben}

\vspace{5pt}
\hrule
\vspace{5pt}

\begin{tabular}{rl}
\multicolumn{2}{l}{\textbf{PRÄSENS}}\\
ich & \textbf{beschließe}\\
du & \textbf{beschließt}\\
er/sie/es & \textbf{beschließt}\\
wir & \textbf{beschließen}\\
ihr & \textbf{beschließt}\\
sie/Sie & \textbf{beschließen}\\
\end{tabular}
\hfill
\begin{tabular}{rl}
\multicolumn{2}{l}{\textbf{PRÄTERITUM}}\\
ich & \textbf{beschloss}\\
du & \textbf{beschlossest}\\
er/sie/es & \textbf{beschloss}\\
wir & \textbf{beschlossen}\\
ihr & \textbf{beschlosst}\\
sie/Sie & \textbf{beschlossen}\\
\end{tabular}
\hfill
\begin{tabular}{rl}
\multicolumn{2}{l}{\textbf{IMPERATIV}}\\
& \textbf{beschließ(e) (du)!}\\
& \textbf{beschließen wir!}\\
& \textbf{beschließt (ihr)!}\\
& \textbf{beschließen Sie!}\\
\end{tabular}

\vspace{10pt}

\begin{tabular}{rl}
\multicolumn{2}{l}{\textbf{PERFEKT}}\\
& haben + \textbf{beschlossen}\\
\end{tabular}
\hfill
\begin{tabular}{rl}
\multicolumn{2}{l}{\textbf{FUTURE I}}\\
& werden + \textbf{beschließen}\\
\end{tabular}
\hfill
\begin{tabular}{rl}
\multicolumn{2}{l}{\textbf{PLUSQUAMPERFEKT}}\\
& hatten + \textbf{beschlossen}\\
\end{tabular}

\vspace{10pt}

\begin{tabular}{lll}
\multicolumn{3}{l}{\textbf{MIT PRÄPOSITIONEN}}\\
& \textbf{beschließen Akk} & (to decide something)\\
\end{tabular}
\hfill
\begin{tabular}{lll}
\multicolumn{3}{l}{\textbf{TRENNBARE VERBEN}}\\
& ---\\
\end{tabular}

\vspace{-5pt}
\end{verbbox-acc}

%beschreiben (B1) – to describe
\begin{verbbox-acc}\addcontentsline{toc}{subsubsection}{beschreiben (to describe)}
    {\Large \bf beschreiben} (to describe) \hfill Type: \textbf{irregular, + Akkusativ} \hfill Auxiliary Verb: \textbf{haben} 
    
     \vspace{5pt}

    \hrule
    
    \vspace{5pt}

  \begin{tabular}{rl}
     \multicolumn{2}{l}{\textbf{PRÄSENS} }\\
    ich & \textbf{beschreibe}\\
    du & \textbf{beschreibst}\\
    er/sie/es & \textbf{beschreibt}\\
    wir & \textbf{beschreiben}\\
    ihr & \textbf{beschreibt}\\
    sie/Sie & \textbf{beschreiben}\\
  \end{tabular}
  \hfill
     \begin{tabular}{rl}
    \multicolumn{2}{l}{\textbf{PRÄTERITUM} }\\
    ich & \textbf{beschrieb}\\
    du & \textbf{beschriebst}\\
    er/sie/es & \textbf{beschrieb}\\
    wir & \textbf{beschrieben}\\
    ihr & \textbf{beschriebt}\\
    sie/Sie & \textbf{beschrieben}\\
  \end{tabular}
  \hfill
  \begin{tabular}{rl}
     \multicolumn{2}{l}{\textbf{IMPERATIV} }\\
   & \textbf{beschreib(e) (du)!}\\
   & \textbf{beschreiben wir!}\\
   & \textbf{beschreibt (ihr)!}\\
   & \textbf{beschreiben Sie!}\\
   & \vspace{10pt}
  \end{tabular}

    \vspace{10pt}

 \begin{tabular}{rl}
     \multicolumn{2}{l}{\textbf{PERFEKT}}\\
   & haben + \textbf{beschrieben}\\
  \end{tabular}
\hfill
   \begin{tabular}{rl}
   
    \multicolumn{2}{l}{\textbf{FUTURE I} }\\
   & werden + \textbf{beschreiben}\\
 
  \end{tabular}
\hfill
   \begin{tabular}{rl}
      \multicolumn{2}{l}{\textbf{PLUSQUAMPERFEKT}}\\
   & hatten + \textbf{beschrieben}\\
  \end{tabular}
  
\vspace{10pt}

\begin{tabular}{lll}
\multicolumn{3}{l}{\textbf{MIT PRÄPOSITIONEN}}\\
& ---\\
\end{tabular}
\hfill
\begin{tabular}{lll}
\multicolumn{3}{l}{\textbf{TRENNBARE VERBEN}}\\
& ---\\
\end{tabular}

  \vspace{-5pt}

\end{verbbox-acc}

%bestätigen (to confirm)
\begin{verbbox-acc}\addcontentsline{toc}{subsubsection}{bestätigen (to confirm)}
{\Large \bf bestätigen} (to confirm) \hfill Type: \textbf{regular, + Akkusativ} \hfill Auxiliary Verb: \textbf{haben} 

\vspace{5pt}
\hrule
\vspace{5pt}

\begin{tabular}{rl}
\multicolumn{2}{l}{\textbf{PRÄSENS}}\\
ich & \textbf{bestätige}\\
du & \textbf{bestätigst}\\
er/sie/es & \textbf{bestätigt}\\
wir & \textbf{bestätigen}\\
ihr & \textbf{bestätigt}\\
sie/Sie & \textbf{bestätigen}\\
\end{tabular}
\hfill
\begin{tabular}{rl}
\multicolumn{2}{l}{\textbf{PRÄTERITUM}}\\
ich & \textbf{bestätigte}\\
du & \textbf{bestätigtest}\\
er/sie/es & \textbf{bestätigte}\\
wir & \textbf{bestätigten}\\
ihr & \textbf{bestätigtet}\\
sie/Sie & \textbf{bestätigten}\\
\end{tabular}
\hfill
\begin{tabular}{rl}
\multicolumn{2}{l}{\textbf{IMPERATIV}}\\
& \textbf{bestätige (du)!}\\
& \textbf{bestätigen wir!}\\
& \textbf{bestätigt (ihr)!}\\
& \textbf{bestätigen Sie!}\\
\end{tabular}

\vspace{10pt}

\begin{tabular}{rl}
\multicolumn{2}{l}{\textbf{PERFEKT}}\\
& haben + \textbf{bestätigt}\\
\end{tabular}
\hfill
\begin{tabular}{rl}
\multicolumn{2}{l}{\textbf{FUTURE I}}\\
& werden + \textbf{bestätigen}\\
\end{tabular}
\hfill
\begin{tabular}{rl}
\multicolumn{2}{l}{\textbf{PLUSQUAMPERFEKT}}\\
& hatten + \textbf{bestätigt}\\
\end{tabular}

\vspace{10pt}

\begin{tabular}{lll}
\multicolumn{3}{l}{\textbf{MIT PRÄPOSITIONEN}}\\
& \textbf{bestätigen gegenüber} + Dat & (to confirm to someone)\\
\end{tabular}
\hfill
\begin{tabular}{lll}
\multicolumn{3}{l}{\textbf{TRENNBARE VERBEN}}\\
& ---\\
\end{tabular}

\vspace{-5pt}
\end{verbbox-acc}

%bestehen (A1) – to pass
\begin{verbbox-acc}\addcontentsline{toc}{subsubsection}{bestehen (to pass)}
    {\Large \bf bestehen} (to pass) \hfill Type: \textbf{irregular, + Akkusativ} \hfill Auxiliary Verb: \textbf{haben} 
    
     \vspace{5pt}

    \hrule
    
    \vspace{5pt}

  \begin{tabular}{rl}
     \multicolumn{2}{l}{\textbf{PRÄSENS} }\\
    ich & \textbf{bestehe}\\
    du & \textbf{bestehst}\\
    er/sie/es & \textbf{besteht}\\
    wir & \textbf{bestehen}\\
    ihr & \textbf{besteht}\\
    sie/Sie & \textbf{bestehen}\\
  \end{tabular}
  \hfill
     \begin{tabular}{rl}
    \multicolumn{2}{l}{\textbf{PRÄTERITUM} }\\
    ich & \textbf{bestand}\\
    du & \textbf{bestandest}\\
    er/sie/es & \textbf{bestand}\\
    wir & \textbf{bestanden}\\
    ihr & \textbf{bestandet}\\
    sie/Sie & \textbf{bestanden}\\
  \end{tabular}
  \hfill
  \begin{tabular}{rl}
     \multicolumn{2}{l}{\textbf{IMPERATIV} }\\
   & \textbf{besteh(e) (du)!}\\
   & \textbf{bestehen wir!}\\
   & \textbf{besteht (ihr)!}\\
   & \textbf{bestehen Sie!}\\
   & \vspace{10pt}
  \end{tabular}

    \vspace{10pt}

 \begin{tabular}{rl}
     \multicolumn{2}{l}{\textbf{PERFEKT}}\\
  &  haben + \textbf{bestanden}\\
  \end{tabular}
\hfill
   \begin{tabular}{rl}
   
    \multicolumn{2}{l}{\textbf{FUTURE I} }\\
  &  werden + \textbf{bestehen}\\
 
  \end{tabular}
\hfill
   \begin{tabular}{rl}
      \multicolumn{2}{l}{\textbf{PLUSQUAMPERFEKT}}\\
   & hatten + \textbf{bestanden}\\
  \end{tabular}
  
 \vspace{10pt}

\begin{tabular}{lll}
\multicolumn{3}{l}{\textbf{MIT PRÄPOSITIONEN}}\\
& ---\\
\end{tabular}
\hfill
\begin{tabular}{lll}
\multicolumn{3}{l}{\textbf{TRENNBARE VERBEN}}\\
& ---\\
\end{tabular}

\vspace{-5pt}

\end{verbbox-acc}
%bestellen (A1) – to order
\begin{verbbox-acc}\addcontentsline{toc}{subsubsection}{bestellen (to order)}
    {\Large \bf bestellen} (to order) \hfill Type: \textbf{regular, + Akkusativ} \hfill Auxiliary Verb: \textbf{haben} 
    
     \vspace{5pt}

    \hrule
    
    \vspace{5pt}

  \begin{tabular}{rl}
     \multicolumn{2}{l}{\textbf{PRÄSENS} }\\
    ich & \textbf{bestelle}\\
    du & \textbf{bestellst}\\
    er/sie/es & \textbf{bestellt}\\
    wir & \textbf{bestellen}\\
    ihr & \textbf{bestellt}\\
    sie/Sie & \textbf{bestellen}\\
  \end{tabular}
  \hfill
     \begin{tabular}{rl}
    \multicolumn{2}{l}{\textbf{PRÄTERITUM} }\\
    ich & \textbf{bestellte}\\
    du & \textbf{bestelltest}\\
    er/sie/es & \textbf{bestellte}\\
    wir & \textbf{bestellten}\\
    ihr & \textbf{bestelltet}\\
    sie/Sie & \textbf{bestellten}\\
  \end{tabular}
  \hfill
  \begin{tabular}{rl}
     \multicolumn{2}{l}{\textbf{IMPERATIV} }\\
   & \textbf{bestell(e) (du)!}\\
   & \textbf{bestellen wir!}\\
   & \textbf{bestellt (ihr)!}\\
   & \textbf{bestellen Sie!}\\
   & \vspace{10pt}
  \end{tabular}

    \vspace{10pt}

 \begin{tabular}{rl}
     \multicolumn{2}{l}{\textbf{PERFEKT}}\\
   & haben + \textbf{bestellt}\\
  \end{tabular}
\hfill
   \begin{tabular}{rl}
   
    \multicolumn{2}{l}{\textbf{FUTURE I} }\\
   & werden + \textbf{bestellen}\\
 
  \end{tabular}
\hfill
   \begin{tabular}{rl}
      \multicolumn{2}{l}{\textbf{PLUSQUAMPERFEKT}}\\
   & hatten + \textbf{bestellt}\\
  \end{tabular}
  
 \vspace{10pt}

\begin{tabular}{lll}
\multicolumn{3}{l}{\textbf{MIT PRÄPOSITIONEN}}\\
& ---\\
\end{tabular}
\hfill
\begin{tabular}{lll}
\multicolumn{3}{l}{\textbf{TRENNBARE VERBEN}}\\
& ---\\
\end{tabular}

\vspace{-5pt}


\end{verbbox-acc}

%betragen - to compromise
\begin{verbbox-acc}\addcontentsline{toc}{subsubsection}{betragen (to amount to / to be)}
{\Large \bf betragen} (to amount to / to be) \hfill Type: \textbf{irregular, + Akkusativ} \hfill Auxiliary Verb: \textbf{haben}

\vspace{5pt}
\hrule
\vspace{5pt}

\begin{tabular}{rl}
\multicolumn{2}{l}{\textbf{PRÄSENS}}\\
ich & \textbf{betrage}\\
du & \textbf{beträgst}\\
er/sie/es & \textbf{beträgt}\\
wir & \textbf{betragen}\\
ihr & \textbf{betragt}\\
sie/Sie & \textbf{betragen}\\
\end{tabular}
\hfill
\begin{tabular}{rl}
\multicolumn{2}{l}{\textbf{PRÄTERITUM}}\\
ich & \textbf{betrug}\\
du & \textbf{betrugst}\\
er/sie/es & \textbf{betrug}\\
wir & \textbf{betrugen}\\
ihr & \textbf{betrugt}\\
sie/Sie & \textbf{betrugen}\\
\end{tabular}
\hfill
\begin{tabular}{rl}
\multicolumn{2}{l}{\textbf{IMPERATIV}}\\
\end{tabular}

\vspace{10pt}

\begin{tabular}{rl}
\multicolumn{2}{l}{\textbf{PERFEKT}}\\
& haben + \textbf{betragen}\\
\end{tabular}
\hfill
\begin{tabular}{rl}
\multicolumn{2}{l}{\textbf{FUTURE I}}\\
& werden + \textbf{betragen}\\
\end{tabular}
\hfill
\begin{tabular}{rl}
\multicolumn{2}{l}{\textbf{PLUSQUAMPERFEKT}}\\
& hatten + \textbf{betragen}\\
\end{tabular}

\vspace{10pt}

\begin{tabular}{lll}
\multicolumn{3}{l}{\textbf{MIT PRÄPOSITIONEN}}\\
\end{tabular}
\hfill
\begin{tabular}{lll}
\multicolumn{3}{l}{\textbf{TRENNBARE VERBEN}}\\
& --- & (no separable forms)\\
\end{tabular}
\vspace{-5pt}
\end{verbbox-acc}

%betreuen - to supervise
\begin{verbbox-acc}\addcontentsline{toc}{subsubsection}{betreuen (to supervise / take care of)}
{\Large \bf betreuen} (to supervise / take care of) \hfill Type: \textbf{regular, + Akkusativ} \hfill Auxiliary Verb: \textbf{haben}

\vspace{5pt}
\hrule
\vspace{5pt}

\begin{tabular}{rl}
\multicolumn{2}{l}{\textbf{PRÄSENS}}\\
ich & \textbf{betreue}\\
du & \textbf{betreust}\\
er/sie/es & \textbf{betreut}\\
wir & \textbf{betreuen}\\
ihr & \textbf{betreut}\\
sie/Sie & \textbf{betreuen}\\
\end{tabular}
\hfill
\begin{tabular}{rl}
\multicolumn{2}{l}{\textbf{PRÄTERITUM}}\\
ich & \textbf{betreute}\\
du & \textbf{betreutest}\\
er/sie/es & \textbf{betreute}\\
wir & \textbf{betreuten}\\
ihr & \textbf{betreutet}\\
sie/Sie & \textbf{betreuten}\\
\end{tabular}
\hfill
\begin{tabular}{rl}
\multicolumn{2}{l}{\textbf{IMPERATIV}}\\
& \textbf{betreue (du)!}\\
& \textbf{betreuen wir!}\\
& \textbf{betreut (ihr)!}\\
& \textbf{betreuen Sie!}\\
\end{tabular}

\vspace{10pt}

\begin{tabular}{rl}
\multicolumn{2}{l}{\textbf{PERFEKT}}\\
& haben + \textbf{betreut}\\
\end{tabular}
\hfill
\begin{tabular}{rl}
\multicolumn{2}{l}{\textbf{FUTURE I}}\\
& werden + \textbf{betreuen}\\
\end{tabular}
\hfill
\begin{tabular}{rl}
\multicolumn{2}{l}{\textbf{PLUSQUAMPERFEKT}}\\
& hatten + \textbf{betreut}\\
\end{tabular}

\vspace{10pt}

\begin{tabular}{lll}
\multicolumn{3}{l}{\textbf{MIT PRÄPOSITIONEN}}\\
& \textbf{betreuen bei} + Dat & (to assist during)\\
& \textbf{betreuen in} + Dat & (to supervise in)\\
\end{tabular}
\hfill
\begin{tabular}{lll}
\multicolumn{3}{l}{\textbf{TRENNBARE VERBEN}}\\
& --- & (no separable forms)\\
\end{tabular}
\vspace{-5pt}
\end{verbbox-acc}


%bewerben (to advertise / promote)
\begin{verbbox-acc}\addcontentsline{toc}{subsubsection}{bewerben (to advertise / promote)}
{\Large \bf bewerben} (to advertise / promote) \hfill Type: \textbf{irregular, + Akkusativ} \hfill Auxiliary Verb: \textbf{haben}

\vspace{5pt}
\hrule
\vspace{5pt}

\begin{tabular}{rl}
\multicolumn{2}{l}{\textbf{PRÄSENS}}\\
ich & \textbf{bewerbe}\\
du & \textbf{bewirbst}\\
er/sie/es & \textbf{bewirbt}\\
wir & \textbf{bewerben}\\
ihr & \textbf{bewerbt}\\
sie/Sie & \textbf{bewerben}\\
\end{tabular}
\hfill
\begin{tabular}{rl}
\multicolumn{2}{l}{\textbf{PRÄTERITUM}}\\
ich & \textbf{bewarb}\\
du & \textbf{bewarbst}\\
er/sie/es & \textbf{bewarb}\\
wir & \textbf{bewarben}\\
ihr & \textbf{bewarbt}\\
sie/Sie & \textbf{bewarben}\\
\end{tabular}
\hfill
\begin{tabular}{rl}
\multicolumn{2}{l}{\textbf{IMPERATIV}}\\
& \textbf{bewirb (du)!}\\
& \textbf{bewerben wir!}\\
& \textbf{bewerbt (ihr)!}\\
& \textbf{bewerben Sie!}\\
\end{tabular}

\vspace{10pt}

\begin{tabular}{rl}
\multicolumn{2}{l}{\textbf{PERFEKT}}\\
& haben + \textbf{beworben}\\
\end{tabular}
\hfill
\begin{tabular}{rl}
\multicolumn{2}{l}{\textbf{FUTURE I}}\\
& werden + \textbf{bewerben}\\
\end{tabular}
\hfill
\begin{tabular}{rl}
\multicolumn{2}{l}{\textbf{PLUSQUAMPERFEKT}}\\
& hatten + \textbf{beworben}\\
\end{tabular}

\vspace{10pt}

\begin{tabular}{lll}
\multicolumn{3}{l}{\textbf{MIT PRÄPOSITIONEN}}\\
& ---\\
\end{tabular}
\hfill
\begin{tabular}{lll}
\multicolumn{3}{l}{\textbf{TRENNBARE VERBEN}}\\
& ---\\
\end{tabular}

\vspace{-5pt}
\end{verbbox-acc}


%bezahlen (A1) – to pay
\begin{verbbox-acc}\addcontentsline{toc}{subsubsection}{bezahlen (to pay)}
    {\Large \bf bezahlen} (to pay) \hfill Type: \textbf{regular, + Akkusativ} \hfill Auxiliary Verb: \textbf{haben} 
    
     \vspace{5pt}

    \hrule
    
    \vspace{5pt}

  \begin{tabular}{rl}
     \multicolumn{2}{l}{\textbf{PRÄSENS} }\\
    ich & \textbf{bezahle}\\
    du & \textbf{bezahlst}\\
    er/sie/es & \textbf{bezahlt}\\
    wir & \textbf{bezahlen}\\
    ihr & \textbf{bezahlt}\\
    sie/Sie & \textbf{bezahlen}\\
  \end{tabular}
  \hfill
     \begin{tabular}{rl}
    \multicolumn{2}{l}{\textbf{PRÄTERITUM} }\\
    ich & \textbf{bezahlte}\\
    du & \textbf{bezahltest}\\
    er/sie/es & \textbf{bezahlte}\\
    wir & \textbf{bezahlten}\\
    ihr & \textbf{bezahltet}\\
    sie/Sie & \textbf{bezahlten}\\
  \end{tabular}
  \hfill
  \begin{tabular}{rl}
     \multicolumn{2}{l}{\textbf{IMPERATIV} }\\
   & \textbf{bezahl(e) (du)!}\\
   & \textbf{bezahlen wir!}\\
   & \textbf{bezahlt (ihr)!}\\
   & \textbf{bezahlen Sie!}\\
   & \vspace{10pt}
  \end{tabular}

    \vspace{10pt}

 \begin{tabular}{rl}
     \multicolumn{2}{l}{\textbf{PERFEKT}}\\
  &  haben + \textbf{bezahlt}\\
  \end{tabular}
\hfill
   \begin{tabular}{rl}
   
    \multicolumn{2}{l}{\textbf{FUTURE I} }\\
  &  werden + \textbf{bezahlen}\\
 
  \end{tabular}
\hfill
   \begin{tabular}{rl}
      \multicolumn{2}{l}{\textbf{PLUSQUAMPERFEKT}}\\
  &  hatten + \textbf{bezahlt}\\
  \end{tabular}
  
  \vspace{10pt}

\begin{tabular}{lll}
\multicolumn{3}{l}{\textbf{MIT PRÄPOSITIONEN}}\\
& ---\\
\end{tabular}
\hfill
\begin{tabular}{lll}
\multicolumn{3}{l}{\textbf{TRENNBARE VERBEN}}\\
& ---\\
\end{tabular}

\vspace{-5pt}


\end{verbbox-acc}

%beziehen (to obtain / refer to)
\begin{verbbox-acc}\addcontentsline{toc}{subsubsection}{beziehen (to obtain / refer to)}
{\Large \bf beziehen} (to obtain / refer to) \hfill Type: \textbf{irregular, + Akkusativ} \hfill Auxiliary Verb: \textbf{haben}

\vspace{5pt}
\hrule
\vspace{5pt}

\begin{tabular}{rl}
\multicolumn{2}{l}{\textbf{PRÄSENS}}\\
ich & \textbf{beziehe}\\
du & \textbf{beziehst}\\
er/sie/es & \textbf{bezieht}\\
wir & \textbf{beziehen}\\
ihr & \textbf{bezieht}\\
sie/Sie & \textbf{beziehen}\\
\end{tabular}
\hfill
\begin{tabular}{rl}
\multicolumn{2}{l}{\textbf{PRÄTERITUM}}\\
ich & \textbf{bezog}\\
du & \textbf{bezogst}\\
er/sie/es & \textbf{bezog}\\
wir & \textbf{bezogen}\\
ihr & \textbf{bezogt}\\
sie/Sie & \textbf{bezogen}\\
\end{tabular}
\hfill
\begin{tabular}{rl}
\multicolumn{2}{l}{\textbf{IMPERATIV}}\\
& \textbf{bezieh(e) (du)!}\\
& \textbf{beziehen wir!}\\
& \textbf{bezieht (ihr)!}\\
& \textbf{beziehen Sie!}\\
\end{tabular}

\vspace{10pt}

\begin{tabular}{rl}
\multicolumn{2}{l}{\textbf{PERFEKT}}\\
& haben + \textbf{bezogen}\\
\end{tabular}
\hfill
\begin{tabular}{rl}
\multicolumn{2}{l}{\textbf{FUTURE I}}\\
& werden + \textbf{beziehen}\\
\end{tabular}
\hfill
\begin{tabular}{rl}
\multicolumn{2}{l}{\textbf{PLUSQUAMPERFEKT}}\\
& hatten + \textbf{bezogen}\\
\end{tabular}

\vspace{10pt}

\begin{tabular}{lll}
\multicolumn{3}{l}{\textbf{MIT PRÄPOSITIONEN}}\\
& \textbf{beziehen Akk} & (to obtain something)\\
& \textbf{beziehen auf} + Akk & (to refer to)\\
\end{tabular}
\hfill
\begin{tabular}{lll}
\multicolumn{3}{l}{\textbf{TRENNBARE VERBEN}}\\
& \textbf{ein$|$beziehen} & (to involve / include)\\
& \textbf{bevor$|$ziehen} & (to prefer)\\
\end{tabular}

\vspace{-5pt}
\end{verbbox-acc}

%bieten (to offer)
\begin{verbbox}\addcontentsline{toc}{subsubsection}{bieten (to offer)}
{\Large \bf bieten} (to offer) \hfill Type: \textbf{irregular} \hfill Auxiliary Verb: \textbf{haben}

\vspace{5pt}
\hrule
\vspace{5pt}

\begin{tabular}{rl}
\multicolumn{2}{l}{\textbf{PRÄSENS}}\\
ich & \textbf{biete}\\
du & \textbf{bietest}\\
er/sie/es & \textbf{bietet}\\
wir & \textbf{bieten}\\
ihr & \textbf{bietet}\\
sie/Sie & \textbf{bieten}\\
\end{tabular}
\hfill
\begin{tabular}{rl}
\multicolumn{2}{l}{\textbf{PRÄTERITUM}}\\
ich & \textbf{bot}\\
du & \textbf{botest}\\
er/sie/es & \textbf{bot}\\
wir & \textbf{boten}\\
ihr & \textbf{botet}\\
sie/Sie & \textbf{boten}\\
\end{tabular}
\hfill
\begin{tabular}{rl}
\multicolumn{2}{l}{\textbf{IMPERATIV}}\\
& \textbf{biete (du)!}\\
& \textbf{bieten wir!}\\
& \textbf{bietet (ihr)!}\\
& \textbf{bieten Sie!}\\
\end{tabular}

\vspace{10pt}

\begin{tabular}{rl}
\multicolumn{2}{l}{\textbf{PERFEKT}}\\
& haben + \textbf{geboten}\\
\end{tabular}
\hfill
\begin{tabular}{rl}
\multicolumn{2}{l}{\textbf{FUTURE I}}\\
& werden + \textbf{bieten}\\
\end{tabular}
\hfill
\begin{tabular}{rl}
\multicolumn{2}{l}{\textbf{PLUSQUAMPERFEKT}}\\
& hatten + \textbf{geboten}\\
\end{tabular}

\vspace{10pt}

\begin{tabular}{lll}
\multicolumn{3}{l}{\textbf{MIT PRÄPOSITIONEN}}\\
& \textbf{bieten} + Dat + Akk & (to offer someone something)\\
& \textbf{bieten für} + Akk & (to bid for)\\
\end{tabular}
\hfill
\begin{tabular}{lll}
\multicolumn{3}{l}{\textbf{TRENNBARE VERBEN}}\\
& \textbf{an$|$bieten} & (to offer)\\
\vspace{10pt}
\end{tabular}

\vspace{-5pt}
\end{verbbox}

%bitten (to ask / request)
\begin{verbbox}\addcontentsline{toc}{subsubsection}{bitten (to ask / request)}
{\Large \bf bitten} (to ask / request) \hfill Type: \textbf{irregular} \hfill Auxiliary Verb: \textbf{haben}

\vspace{5pt}
\hrule
\vspace{5pt}

\begin{tabular}{rl}
\multicolumn{2}{l}{\textbf{PRÄSENS}}\\
ich & \textbf{bitte}\\
du & \textbf{bittest}\\
er/sie/es & \textbf{bittet}\\
wir & \textbf{bitten}\\
ihr & \textbf{bittet}\\
sie/Sie & \textbf{bitten}\\
\end{tabular}
\hfill
\begin{tabular}{rl}
\multicolumn{2}{l}{\textbf{PRÄTERITUM}}\\
ich & \textbf{bat}\\
du & \textbf{batst}\\
er/sie/es & \textbf{bat}\\
wir & \textbf{baten}\\
ihr & \textbf{batet}\\
sie/Sie & \textbf{baten}\\
\end{tabular}
\hfill
\begin{tabular}{rl}
\multicolumn{2}{l}{\textbf{IMPERATIV}}\\
& \textbf{bitte (du)!}\\
& \textbf{bitten wir!}\\
& \textbf{bittet (ihr)!}\\
& \textbf{bitten Sie!}\\
\end{tabular}

\vspace{10pt}

\begin{tabular}{rl}
\multicolumn{2}{l}{\textbf{PERFEKT}}\\
& haben + \textbf{gebeten}\\
\end{tabular}
\hfill
\begin{tabular}{rl}
\multicolumn{2}{l}{\textbf{FUTURE I}}\\
& werden + \textbf{bitten}\\
\end{tabular}
\hfill
\begin{tabular}{rl}
\multicolumn{2}{l}{\textbf{PLUSQUAMPERFEKT}}\\
& hatten + \textbf{gebeten}\\
\end{tabular}

\vspace{10pt}

\begin{tabular}{lll}
\multicolumn{3}{l}{\textbf{MIT PRÄPOSITIONEN}}\\
& \textbf{bitten um} + Akk & (to ask for / request)\\
\end{tabular}
\hfill
\begin{tabular}{lll}
\multicolumn{3}{l}{\textbf{TRENNBARE VERBEN}}\\
& --- & (no separable forms)\\
\end{tabular}
\vspace{-5pt}
\end{verbbox}


%bleiben (A1) – to stay
\begin{verbbox}\addcontentsline{toc}{subsubsection}{bleiben (to stay)}
    {\Large \bf bleiben} (to stay) \hfill Type: \textbf{irregular} \hfill Auxiliary Verb: \textbf{sein} 
    
     \vspace{5pt}

    \hrule
    
    \vspace{5pt}

  \begin{tabular}{rl}
     \multicolumn{2}{l}{\textbf{PRÄSENS} }\\
    ich & \textbf{bleibe}\\
    du & \textbf{bleibst}\\
    er/sie/es & \textbf{bleibt}\\
    wir & \textbf{bleiben}\\
    ihr & \textbf{bleibt}\\
    sie/Sie & \textbf{bleiben}\\
  \end{tabular}
  \hfill
     \begin{tabular}{rl}
    \multicolumn{2}{l}{\textbf{PRÄTERITUM} }\\
    ich & \textbf{blieb}\\
    du & \textbf{bliebst}\\
    er/sie/es & \textbf{blieb}\\
    wir & \textbf{blieben}\\
    ihr & \textbf{bliebt}\\
    sie/Sie & \textbf{blieben}\\
  \end{tabular}
  \hfill
  \begin{tabular}{rl}
     \multicolumn{2}{l}{\textbf{IMPERATIV} }\\
   & \textbf{bleib(e) (du)!}\\
   & \textbf{bleiben wir!}\\
   & \textbf{bleibt (ihr)!}\\
   & \textbf{bleiben Sie!}\\
   & \vspace{10pt}
  \end{tabular}

    \vspace{10pt}

 \begin{tabular}{rl}
     \multicolumn{2}{l}{\textbf{PERFEKT}}\\
  &  sein + \textbf{geblieben}\\
  \end{tabular}
\hfill
   \begin{tabular}{rl}
   
    \multicolumn{2}{l}{\textbf{FUTURE I} }\\
  &  werden + \textbf{bleiben}\\
 
  \end{tabular}
\hfill
   \begin{tabular}{rl}
      \multicolumn{2}{l}{\textbf{PLUSQUAMPERFEKT}}\\
  &  waren + \textbf{geblieben}\\
  \end{tabular}
  
  \vspace{10pt}

\begin{tabular}{lll}
\multicolumn{3}{l}{\textbf{MIT PRÄPOSITIONEN}}\\
& ---\\
\end{tabular}
\hfill
\begin{tabular}{lll}
\multicolumn{3}{l}{\textbf{TRENNBARE VERBEN}}\\
& ---\\
\end{tabular}

  \vspace{-5pt}


\end{verbbox}

%blühen (to bloom / flourish)
\begin{verbbox}\addcontentsline{toc}{subsubsection}{blühen (to bloom / flourish)}
{\Large \bf blühen} (to bloom / flourish) \hfill Type: \textbf{regular} \hfill Auxiliary Verb: \textbf{haben}

\vspace{5pt}
\hrule
\vspace{5pt}

\begin{tabular}{rl}
\multicolumn{2}{l}{\textbf{PRÄSENS}}\\
ich & \textbf{blühe}\\
du & \textbf{blühst}\\
er/sie/es & \textbf{blüht}\\
wir & \textbf{blühen}\\
ihr & \textbf{blüht}\\
sie/Sie & \textbf{blühen}\\
\end{tabular}
\hfill
\begin{tabular}{rl}
\multicolumn{2}{l}{\textbf{PRÄTERITUM}}\\
ich & \textbf{blühte}\\
du & \textbf{blühtest}\\
er/sie/es & \textbf{blühte}\\
wir & \textbf{blühten}\\
ihr & \textbf{blühtet}\\
sie/Sie & \textbf{blühten}\\
\end{tabular}
\hfill
\begin{tabular}{rl}
\multicolumn{2}{l}{\textbf{IMPERATIV}}\\
& \textbf{blühe (du)!}\\
& \textbf{blühen wir!}\\
& \textbf{blüht (ihr)!}\\
& \textbf{blühen Sie!}\\
\end{tabular}

\vspace{10pt}

\begin{tabular}{rl}
\multicolumn{2}{l}{\textbf{PERFEKT}}\\
& haben + \textbf{geblüht}\\
\end{tabular}
\hfill
\begin{tabular}{rl}
\multicolumn{2}{l}{\textbf{FUTURE I}}\\
& werden + \textbf{blühen}\\
\end{tabular}
\hfill
\begin{tabular}{rl}
\multicolumn{2}{l}{\textbf{PLUSQUAMPERFEKT}}\\
& hatten + \textbf{geblüht}\\
\end{tabular}

\vspace{10pt}

\begin{tabular}{lll}
\multicolumn{3}{l}{\textbf{MIT PRÄPOSITIONEN}}\\
& ---\\
\end{tabular}
\hfill
\begin{tabular}{lll}
\multicolumn{3}{l}{\textbf{TRENNBARE VERBEN}}\\
& \textbf{auf$|$blühen} & (to blossom / flourish)\\
\end{tabular}
\vspace{-5pt}
\end{verbbox}

%bohren (to drill)
\begin{verbbox}\addcontentsline{toc}{subsubsection}{bohren (to drill)}
{\Large \bf bohren} (to drill) \hfill Type: \textbf{regular} \hfill Auxiliary Verb: \textbf{haben}

\vspace{5pt}
\hrule
\vspace{5pt}

\begin{tabular}{rl}
\multicolumn{2}{l}{\textbf{PRÄSENS}}\\
ich & \textbf{bohre}\\
du & \textbf{bohrst}\\
er/sie/es & \textbf{bohrt}\\
wir & \textbf{bohren}\\
ihr & \textbf{bohrt}\\
sie/Sie & \textbf{bohren}\\
\end{tabular}
\hfill
\begin{tabular}{rl}
\multicolumn{2}{l}{\textbf{PRÄTERITUM}}\\
ich & \textbf{bohrte}\\
du & \textbf{bohrtest}\\
er/sie/es & \textbf{bohrte}\\
wir & \textbf{bohrten}\\
ihr & \textbf{bohrtet}\\
sie/Sie & \textbf{bohrten}\\
\end{tabular}
\hfill
\begin{tabular}{rl}
\multicolumn{2}{l}{\textbf{IMPERATIV}}\\
& \textbf{bohre (du)!}\\
& \textbf{bohren wir!}\\
& \textbf{bohrt (ihr)!}\\
& \textbf{bohren Sie!}\\
\end{tabular}

\vspace{10pt}

\begin{tabular}{rl}
\multicolumn{2}{l}{\textbf{PERFEKT}}\\
& haben + \textbf{gebohrt}\\
\end{tabular}
\hfill
\begin{tabular}{rl}
\multicolumn{2}{l}{\textbf{FUTURE I}}\\
& werden + \textbf{bohren}\\
\end{tabular}
\hfill
\begin{tabular}{rl}
\multicolumn{2}{l}{\textbf{PLUSQUAMPERFEKT}}\\
& hatten + \textbf{gebohrt}\\
\end{tabular}

\vspace{10pt}

\begin{tabular}{lll}
\multicolumn{3}{l}{\textbf{MIT PRÄPOSITIONEN}}\\
& \textbf{bohren in} + Dat & (to drill in / into)\\
& \textbf{bohren durch} + Akk & (to drill through)\\
\end{tabular}
\hfill
\begin{tabular}{lll}
\multicolumn{3}{l}{\textbf{TRENNBARE VERBEN}}\\
& \textbf{auf$|$bohren} & (to drill on / onto surface)\\
\end{tabular}
\vspace{-5pt}
\end{verbbox}

%brauchen (A1) – to need
\begin{verbbox-acc}\addcontentsline{toc}{subsubsection}{brauchen (to need)}
    {\Large \bf brauchen} (to need) \hfill Type: \textbf{regular, + Akkusativ} \hfill Auxiliary Verb: \textbf{haben} 
    
     \vspace{5pt}

    \hrule
    
    \vspace{5pt}

  \begin{tabular}{rl}
     \multicolumn{2}{l}{\textbf{PRÄSENS} }\\
    ich & \textbf{brauche}\\
    du & \textbf{brauchst}\\
    er/sie/es & \textbf{braucht}\\
    wir & \textbf{brauchen}\\
    ihr & \textbf{braucht}\\
    sie/Sie & \textbf{brauchen}\\
  \end{tabular}
  \hfill
     \begin{tabular}{rl}
    \multicolumn{2}{l}{\textbf{PRÄTERITUM} }\\
    ich & \textbf{brauchte}\\
    du & \textbf{brauchtest}\\
    er/sie/es & \textbf{brauchte}\\
    wir & \textbf{brauchten}\\
    ihr & \textbf{brauchtet}\\
    sie/Sie & \textbf{brauchten}\\
  \end{tabular}
  \hfill
  \begin{tabular}{rl}
     \multicolumn{2}{l}{\textbf{IMPERATIV} }\\
   & \textbf{brauch(e) (du)!}\\
   & \textbf{brauchen wir!}\\
   & \textbf{braucht (ihr)!}\\
   & \textbf{brauchen Sie!}\\
   & \vspace{10pt}
  \end{tabular}

    \vspace{10pt}

 \begin{tabular}{rl}
     \multicolumn{2}{l}{\textbf{PERFEKT}}\\
   & haben + \textbf{gebraucht}\\
  \end{tabular}
\hfill
   \begin{tabular}{rl}
   
    \multicolumn{2}{l}{\textbf{FUTURE I} }\\
  &  werden + \textbf{brauchen}\\
 
  \end{tabular}
\hfill
   \begin{tabular}{rl}
      \multicolumn{2}{l}{\textbf{PLUSQUAMPERFEKT}}\\
  &  hatten + \textbf{gebraucht}\\
  \end{tabular}
  
  \vspace{10pt}

\begin{tabular}{lll}
\multicolumn{3}{l}{\textbf{MIT PRÄPOSITIONEN}}\\
& ---\\
\end{tabular}
\hfill
\begin{tabular}{lll}
\multicolumn{3}{l}{\textbf{TRENNBARE VERBEN}}\\
& ---\\
\end{tabular}

  \vspace{-5pt}


\end{verbbox-acc}
%bringen (A1) – to bring
\begin{verbbox-acc}\addcontentsline{toc}{subsubsection}{bringen (to bring)}
    {\Large \bf bringen} (to bring) \hfill Type: \textbf{irregular, + Akkusativ} \hfill Auxiliary Verb: \textbf{haben} 
    
     \vspace{5pt}

    \hrule
    
    \vspace{5pt}

  \begin{tabular}{rl}
     \multicolumn{2}{l}{\textbf{PRÄSENS} }\\
    ich & \textbf{bringe}\\
    du & \textbf{bringst}\\
    er/sie/es & \textbf{bringt}\\
    wir & \textbf{bringen}\\
    ihr & \textbf{bringt}\\
    sie/Sie & \textbf{bringen}\\
  \end{tabular}
  \hfill
     \begin{tabular}{rl}
    \multicolumn{2}{l}{\textbf{PRÄTERITUM} }\\
    ich & \textbf{brachte}\\
    du & \textbf{brachtest}\\
    er/sie/es & \textbf{brachte}\\
    wir & \textbf{brachten}\\
    ihr & \textbf{brachtet}\\
    sie/Sie & \textbf{brachten}\\
  \end{tabular}
  \hfill
  \begin{tabular}{rl}
     \multicolumn{2}{l}{\textbf{IMPERATIV} }\\
   & \textbf{bring(e) (du)!}\\
   & \textbf{bringen wir!}\\
   & \textbf{bringt (ihr)!}\\
   & \textbf{bringen Sie!}\\
   & \vspace{10pt}
  \end{tabular}

    \vspace{10pt}

 \begin{tabular}{rl}
     \multicolumn{2}{l}{\textbf{PERFEKT}}\\
 &   haben + \textbf{gebracht}\\
  \end{tabular}
\hfill
   \begin{tabular}{rl}
   
    \multicolumn{2}{l}{\textbf{FUTURE I} }\\
  &  werden + \textbf{bringen}\\
 
  \end{tabular}
\hfill
   \begin{tabular}{rl}
      \multicolumn{2}{l}{\textbf{PLUSQUAMPERFEKT}}\\
   & hatten + \textbf{gebracht}\\
  \end{tabular}

  
  \vspace{10pt}

   \begin{tabular}{lll}
      \multicolumn{3}{l}{\textbf{MIT PRÄPOSITIONEN}}\\
& \textbf{bringen zu} + Dat & (to bring to someone / place)\\
&\textbf{bringen auf} + Akk & (to raise / summon)\\
&\textbf{bringen für} + Akk & (to bring for)\\
&\textbf{bringen von} + Dat & (to bring from)\\
&\textbf{bringen in} + Akk & (to move into)\\
\vspace{70pt}
  \end{tabular}
  \hfill
   \begin{tabular}{lll}
      \multicolumn{3}{l}{\textbf{TRENNBARE VERBEN}}\\
 & \textbf{um$|$bringen} & (to kill)\\
 & \textbf{mit$|$bringen} & (to bring along)\\
  & \textbf{nach$|$bringen} & (to bring later)\\
  & \textbf{hin$|$bringen} & (to bring there)\\
  & \textbf{her$|$bringen} & (to bring here)\\
     & \textbf{auf$|$bringen} & (to raise money / summon)\\
     & \textbf{ein$|$bringen} & (to yield / introduce)\\
     & \textbf{bei$|$bringen} & (to teach / instill)\\
     & \textbf{unter$|$bringen} & (to accomodate)\\
     & \textbf{weg$|$bringen} & (to take away)\\
      & \textbf{zurück$|$bringen} & (to bring back)\\

      & \textbf{vorbei$|$bringen} & (to bring by / drop off)\\
  \end{tabular}


  \vspace{-5pt}

\end{verbbox-acc}

%buchen (to book / reserve)
\begin{verbbox-acc}\addcontentsline{toc}{subsubsection}{buchen (to book / reserve)}
{\Large \bf buchen} (to book / reserve) \hfill Type: \textbf{regular, + Akkusativ} \hfill Auxiliary Verb: \textbf{haben}

\vspace{5pt}
\hrule
\vspace{5pt}

\begin{tabular}{rl}
\multicolumn{2}{l}{\textbf{PRÄSENS}}\\
ich & \textbf{buche}\\
du & \textbf{buchst}\\
er/sie/es & \textbf{bucht}\\
wir & \textbf{buchen}\\
ihr & \textbf{bucht}\\
sie/Sie & \textbf{buchen}\\
\end{tabular}
\hfill
\begin{tabular}{rl}
\multicolumn{2}{l}{\textbf{PRÄTERITUM}}\\
ich & \textbf{buchte}\\
du & \textbf{buchtest}\\
er/sie/es & \textbf{buchte}\\
wir & \textbf{buchten}\\
ihr & \textbf{buchtet}\\
sie/Sie & \textbf{buchten}\\
\end{tabular}
\hfill
\begin{tabular}{rl}
\multicolumn{2}{l}{\textbf{IMPERATIV}}\\
& \textbf{buche (du)!}\\
& \textbf{buchen wir!}\\
& \textbf{bucht (ihr)!}\\
& \textbf{buchen Sie!}\\
\end{tabular}

\vspace{10pt}

\begin{tabular}{rl}
\multicolumn{2}{l}{\textbf{PERFEKT}}\\
& haben + \textbf{gebucht}\\
\end{tabular}
\hfill
\begin{tabular}{rl}
\multicolumn{2}{l}{\textbf{FUTURE I}}\\
& werden + \textbf{buchen}\\
\end{tabular}
\hfill
\begin{tabular}{rl}
\multicolumn{2}{l}{\textbf{PLUSQUAMPERFEKT}}\\
& hatten + \textbf{gebucht}\\
\end{tabular}

\vspace{10pt}

\begin{tabular}{lll}
\multicolumn{3}{l}{\textbf{MIT PRÄPOSITIONEN}}\\
& \textbf{buchen für} + Akk & (to book for)\\
& \textbf{buchen bei} + Dat & (to book at / with)\\
\end{tabular}
\hfill
\begin{tabular}{lll}
\multicolumn{3}{l}{\textbf{TRENNBARE VERBEN}}\\
& ---\\
\end{tabular}
\vspace{-5pt}
\end{verbbox-acc}

\subsection{C Verbs}

%checken (to check / verify, informal)
\begin{verbbox}\addcontentsline{toc}{subsubsection}{checken (to check / verify, informal)}
{\Large \bf checken} (to check / verify) \hfill Type: \textbf{regular, colloquial} \hfill Auxiliary Verb: \textbf{haben}

\vspace{5pt}
\hrule
\vspace{5pt}

\begin{tabular}{rl}
\multicolumn{2}{l}{\textbf{PRÄSENS}}\\
ich & \textbf{checke}\\
du & \textbf{checkst}\\
er/sie/es & \textbf{checkt}\\
wir & \textbf{checken}\\
ihr & \textbf{checkt}\\
sie/Sie & \textbf{checken}\\
\end{tabular}
\hfill
\begin{tabular}{rl}
\multicolumn{2}{l}{\textbf{PRÄTERITUM}}\\
ich & \textbf{checkte}\\
du & \textbf{checktest}\\
er/sie/es & \textbf{checkte}\\
wir & \textbf{checkten}\\
ihr & \textbf{checktet}\\
sie/Sie & \textbf{checkten}\\
\end{tabular}
\hfill
\begin{tabular}{rl}
\multicolumn{2}{l}{\textbf{IMPERATIV}}\\
& \textbf{checke (du)!}\\
& \textbf{checken wir!}\\
& \textbf{checkt (ihr)!}\\
& \textbf{checken Sie!}\\
\end{tabular}

\vspace{10pt}

\begin{tabular}{rl}
\multicolumn{2}{l}{\textbf{PERFEKT}}\\
& haben + \textbf{gecheckt}\\
\end{tabular}
\hfill
\begin{tabular}{rl}
\multicolumn{2}{l}{\textbf{FUTURE I}}\\
& werden + \textbf{checken}\\
\end{tabular}
\hfill
\begin{tabular}{rl}
\multicolumn{2}{l}{\textbf{PLUSQUAMPERFEKT}}\\
& hatten + \textbf{gecheckt}\\
\end{tabular}

\vspace{10pt}

\begin{tabular}{lll}
\multicolumn{3}{l}{\textbf{MIT PRÄPOSITIONEN}}\\
& ---\\
\end{tabular}
\hfill
\begin{tabular}{lll}
\multicolumn{3}{l}{\textbf{TRENNBARE VERBEN}}\\
& ---\\
\end{tabular}
\vspace{-5pt}
\end{verbbox}

\subsection{D Verbs}

%dauern - (to last / take time)
\begin{verbbox}\addcontentsline{toc}{subsubsection}{dauern (to last / take time)}
    {\Large \bf dauern} (to last / take time) \hfill Type: \textbf{regular} \hfill Auxiliary Verb: \textbf{haben} 
    
     \vspace{5pt}

    \hrule
    
    \vspace{5pt}

  \begin{tabular}{rl}
     \multicolumn{2}{l}{\textbf{PRÄSENS} }\\
    ich & \textbf{dauere}\\
    du & \textbf{dauerst}\\
    er/sie/es & \textbf{dauert}\\
    wir & \textbf{dauern}\\
    ihr & \textbf{dauert}\\
    sie/Sie & \textbf{dauern}\\
  \end{tabular}
  \hfill
     \begin{tabular}{rl}
    \multicolumn{2}{l}{\textbf{PRÄTERITUM} }\\
    ich & \textbf{dauerte}\\
    du & \textbf{dauertest}\\
    er/sie/es & \textbf{dauerte}\\
    wir & \textbf{dauerten}\\
    ihr & \textbf{dauertet}\\
    sie/Sie & \textbf{dauerten}\\
  \end{tabular}
  \hfill
  \begin{tabular}{rl}
     \multicolumn{2}{l}{\textbf{IMPERATIV} }\\
   & \textbf{dauere (du)!}\\
   & \textbf{dauern wir!}\\
   & \textbf{dauert (ihr)!}\\
   & \textbf{dauern Sie!}\\
   & \vspace{10pt}
  \end{tabular}

    \vspace{10pt}

 \begin{tabular}{rl}
     \multicolumn{2}{l}{\textbf{PERFEKT}}\\
   & haben + \textbf{gedauert}\\
  \end{tabular}
\hfill
   \begin{tabular}{rl}
    \multicolumn{2}{l}{\textbf{FUTURE I} }\\
   & werden + \textbf{dauern}\\
  \end{tabular}
\hfill
   \begin{tabular}{rl}
      \multicolumn{2}{l}{\textbf{PLUSQUAMPERFEKT}}\\
   & hatten + \textbf{gedauert}\\
  \end{tabular}

   \vspace{10pt}

   \begin{tabular}{lll}
      \multicolumn{3}{l}{\textbf{MIT PRÄPOSITIONEN}}\\
& \textbf{dauern bis} + Akk & (to last until)\\
& \textbf{dauern über} + Akk & (to last over / longer than)\\
   \end{tabular}
  \hfill
   \begin{tabular}{lll}
      \multicolumn{3}{l}{\textbf{TRENNBARE VERBEN}}\\
 & \textbf{an$|$dauern} & (to continue / persist)\\
\vspace{10pt}
  \end{tabular}

\vspace{-5pt}
\end{verbbox}

%demonstrieren (to demonstrate / protest)
\begin{verbbox}\addcontentsline{toc}{subsubsection}{demonstrieren (to demonstrate / protest)}
{\Large \bf demonstrieren} (to demonstrate / protest) \hfill Type: \textbf{regular} \hfill Auxiliary Verb: \textbf{haben}

\vspace{5pt}
\hrule
\vspace{5pt}

\begin{tabular}{rl}
\multicolumn{2}{l}{\textbf{PRÄSENS}}\\
ich & \textbf{demonstriere}\\
du & \textbf{demonstrierst}\\
er/sie/es & \textbf{demonstriert}\\
wir & \textbf{demonstrieren}\\
ihr & \textbf{demonstriert}\\
sie/Sie & \textbf{demonstrieren}\\
\end{tabular}
\hfill
\begin{tabular}{rl}
\multicolumn{2}{l}{\textbf{PRÄTERITUM}}\\
ich & \textbf{demonstrierte}\\
du & \textbf{demonstriertest}\\
er/sie/es & \textbf{demonstrierte}\\
wir & \textbf{demonstrierten}\\
ihr & \textbf{demonstriertet}\\
sie/Sie & \textbf{demonstrierten}\\
\end{tabular}
\hfill
\begin{tabular}{rl}
\multicolumn{2}{l}{\textbf{IMPERATIV}}\\
& \textbf{demonstriere (du)!}\\
& \textbf{demonstrieren wir!}\\
& \textbf{demonstriert (ihr)!}\\
& \textbf{demonstrieren Sie!}\\
\end{tabular}

\vspace{10pt}

\begin{tabular}{rl}
\multicolumn{2}{l}{\textbf{PERFEKT}}\\
& haben + \textbf{demonstriert}\\
\end{tabular}
\hfill
\begin{tabular}{rl}
\multicolumn{2}{l}{\textbf{FUTURE I}}\\
& werden + \textbf{demonstrieren}\\
\end{tabular}
\hfill
\begin{tabular}{rl}
\multicolumn{2}{l}{\textbf{PLUSQUAMPERFEKT}}\\
& hatten + \textbf{demonstriert}\\
\end{tabular}

\vspace{10pt}

\begin{tabular}{lll}
\multicolumn{3}{l}{\textbf{MIT PRÄPOSITIONEN}}\\
& \textbf{demonstrieren gegen} + Akk & (to demonstrate against)\\
& \textbf{demonstrieren für} + Akk & (to demonstrate for)\\
\end{tabular}
\hfill
\begin{tabular}{lll}
\multicolumn{3}{l}{\textbf{TRENNBARE VERBEN}}\\
& ---\\
\vspace{10pt}
\end{tabular}

\vspace{-5pt}
\end{verbbox}

%denken (A1) – to think
\begin{verbbox-acc}\addcontentsline{toc}{subsubsection}{denken (to think)}
    {\Large \bf denken} (to think) \hfill Type: \textbf{mixed, + Akkusativ} \hfill Auxiliary Verb: \textbf{haben} 
    
     \vspace{5pt}

    \hrule
    
    \vspace{5pt}

  \begin{tabular}{rl}
     \multicolumn{2}{l}{\textbf{PRÄSENS} }\\
    ich & \textbf{denke}\\
    du & \textbf{denkst}\\
    er/sie/es & \textbf{denkt}\\
    wir & \textbf{denken}\\
    ihr & \textbf{denkt}\\
    sie/Sie & \textbf{denken}\\
  \end{tabular}
  \hfill
     \begin{tabular}{rl}
    \multicolumn{2}{l}{\textbf{PRÄTERITUM} }\\
    ich & \textbf{dachte}\\
    du & \textbf{dachtest}\\
    er/sie/es & \textbf{dachte}\\
    wir & \textbf{dachten}\\
    ihr & \textbf{dachtet}\\
    sie/Sie & \textbf{dachten}\\
  \end{tabular}
  \hfill
  \begin{tabular}{rl}
     \multicolumn{2}{l}{\textbf{IMPERATIV} }\\
   & \textbf{denk(e) (du)!}\\
   & \textbf{denken wir!}\\
   & \textbf{denkt (ihr)!}\\
   & \textbf{denken Sie!}\\
   & \vspace{10pt}
  \end{tabular}

    \vspace{10pt}

 \begin{tabular}{rl}
     \multicolumn{2}{l}{\textbf{PERFEKT}}\\
  &  haben + \textbf{gedacht}\\
  \end{tabular}
\hfill
   \begin{tabular}{rl}
   
    \multicolumn{2}{l}{\textbf{FUTURE I} }\\
  &  werden + \textbf{denken}\\
 
  \end{tabular}
\hfill
   \begin{tabular}{rl}
      \multicolumn{2}{l}{\textbf{PLUSQUAMPERFEKT}}\\
  &  hatten + \textbf{gedacht}\\
  \end{tabular}
  
  \vspace{10pt}

\begin{tabular}{lll}
\multicolumn{3}{l}{\textbf{MIT PRÄPOSITIONEN}}\\
& ---\\
\end{tabular}
\hfill
\begin{tabular}{lll}
\multicolumn{3}{l}{\textbf{TRENNBARE VERBEN}}\\
& ---\\
\end{tabular}

  \vspace{-5pt}


\end{verbbox-acc}

%diskutieren (B1) – to discuss
\begin{verbbox-acc}\addcontentsline{toc}{subsubsection}{diskutieren (to discuss)}
    {\Large \bf diskutieren} (to discuss) \hfill Type: \textbf{irregular, + Akkusativ} \hfill Auxiliary Verb: \textbf{haben} 
    
     \vspace{5pt}

    \hrule
    
    \vspace{5pt}

  \begin{tabular}{rl}
     \multicolumn{2}{l}{\textbf{PRÄSENS} }\\
    ich & \textbf{diskutiere}\\
    du & \textbf{diskutierst}\\
    er/sie/es & \textbf{diskutiert}\\
    wir & \textbf{diskutieren}\\
    ihr & \textbf{diskutiert}\\
    sie/Sie & \textbf{diskutieren}\\
  \end{tabular}
  \hfill
     \begin{tabular}{rl}
    \multicolumn{2}{l}{\textbf{PRÄTERITUM} }\\
    ich & \textbf{diskutierte}\\
    du & \textbf{diskutiertest}\\
    er/sie/es & \textbf{diskutierte}\\
    wir & \textbf{diskutierten}\\
    ihr & \textbf{diskutiertet}\\
    sie/Sie & \textbf{diskutierten}\\
  \end{tabular}
  \hfill
  \begin{tabular}{rl}
     \multicolumn{2}{l}{\textbf{IMPERATIV} }\\
   & \textbf{diskutier(e) (du)!}\\
   & \textbf{diskutieren wir!}\\
   & \textbf{diskutiert (ihr)!}\\
   & \textbf{diskutieren Sie!}\\
   & \vspace{10pt}
  \end{tabular}

    \vspace{10pt}

 \begin{tabular}{rl}
     \multicolumn{2}{l}{\textbf{PERFEKT}}\\
  &  haben + \textbf{diskutiert}\\
  \end{tabular}
\hfill
   \begin{tabular}{rl}
   
    \multicolumn{2}{l}{\textbf{FUTURE I} }\\
  &  werden + \textbf{diskutieren}\\
 
  \end{tabular}
\hfill
   \begin{tabular}{rl}
      \multicolumn{2}{l}{\textbf{PLUSQUAMPERFEKT}}\\
  &  hatten + \textbf{diskutiert}\\
  \end{tabular}
  
  \vspace{10pt}

\begin{tabular}{lll}
\multicolumn{3}{l}{\textbf{MIT PRÄPOSITIONEN}}\\
& ---\\
\end{tabular}
\hfill
\begin{tabular}{lll}
\multicolumn{3}{l}{\textbf{TRENNBARE VERBEN}}\\
& ---\\
\end{tabular}

  \vspace{-5pt}


\end{verbbox-acc}

%drehen (to turn / rotate)
\begin{verbbox}\addcontentsline{toc}{subsubsection}{drehen (to turn / rotate)}
{\Large \bf drehen} (to turn / rotate) \hfill Type: \textbf{regular} \hfill Auxiliary Verb: \textbf{haben}

\vspace{5pt}
\hrule
\vspace{5pt}

\begin{tabular}{rl}
\multicolumn{2}{l}{\textbf{PRÄSENS}}\\
ich & \textbf{drehe}\\
du & \textbf{drehst}\\
er/sie/es & \textbf{dreht}\\
wir & \textbf{drehen}\\
ihr & \textbf{dreht}\\
sie/Sie & \textbf{drehen}\\
\end{tabular}
\hfill
\begin{tabular}{rl}
\multicolumn{2}{l}{\textbf{PRÄTERITUM}}\\
ich & \textbf{drehte}\\
du & \textbf{drehtest}\\
er/sie/es & \textbf{drehte}\\
wir & \textbf{drehten}\\
ihr & \textbf{drehtet}\\
sie/Sie & \textbf{drehten}\\
\end{tabular}
\hfill
\begin{tabular}{rl}
\multicolumn{2}{l}{\textbf{IMPERATIV}}\\
& \textbf{dreh(e) (du)!}\\
& \textbf{drehen wir!}\\
& \textbf{dreht (ihr)!}\\
& \textbf{drehen Sie!}\\
\end{tabular}

\vspace{10pt}

\begin{tabular}{rl}
\multicolumn{2}{l}{\textbf{PERFEKT}}\\
& haben + \textbf{gedreht}\\
\end{tabular}
\hfill
\begin{tabular}{rl}
\multicolumn{2}{l}{\textbf{FUTURE I}}\\
& werden + \textbf{drehen}\\
\end{tabular}
\hfill
\begin{tabular}{rl}
\multicolumn{2}{l}{\textbf{PLUSQUAMPERFEKT}}\\
& hatten + \textbf{gedreht}\\
\end{tabular}

\vspace{10pt}

\begin{tabular}{lll}
\multicolumn{3}{l}{\textbf{MIT PRÄPOSITIONEN}}\\
& \textbf{drehen an} + Dat & (to turn / adjust)\\
& \textbf{drehen um} + Akk & (to turn around)\\
& \textbf{drehen nach} + Dat & (to turn toward)\\
\end{tabular}
\hfill
\begin{tabular}{lll}
\multicolumn{3}{l}{\textbf{TRENNBARE VERBEN}}\\
& \textbf{ab$|$drehen} & (to turn off)\\
& \textbf{auf$|$drehen} & (to turn up)\\
& \textbf{um$|$drehen} & (to turn around)\\
& \textbf{zu$|$drehen} & (to close / screw shut)\\
\end{tabular}
\vspace{-5pt}
\end{verbbox}

%duzen (to address someone with "du")
\begin{verbbox}\addcontentsline{toc}{subsubsection}{duzen (to address someone with ``du")}
{\Large \bf duzen} (to address someone with ``du") \hfill Type: \textbf{regular} \hfill Auxiliary Verb: \textbf{haben}

\vspace{5pt}
\hrule
\vspace{5pt}

\begin{tabular}{rl}
\multicolumn{2}{l}{\textbf{PRÄSENS}}\\
ich & \textbf{duze}\\
du & \textbf{duzt}\\
er/sie/es & \textbf{duzt}\\
wir & \textbf{duzen}\\
ihr & \textbf{duzt}\\
sie/Sie & \textbf{duzen}\\
\end{tabular}
\hfill
\begin{tabular}{rl}
\multicolumn{2}{l}{\textbf{PRÄTERITUM}}\\
ich & \textbf{duzte}\\
du & \textbf{duztest}\\
er/sie/es & \textbf{duzte}\\
wir & \textbf{duzten}\\
ihr & \textbf{duztet}\\
sie/Sie & \textbf{duzten}\\
\end{tabular}
\hfill
\begin{tabular}{rl}
\multicolumn{2}{l}{\textbf{IMPERATIV}}\\
& \textbf{duze (du)!}\\
& \textbf{duzen wir!}\\
& \textbf{duzt (ihr)!}\\
& \textbf{duzen Sie!}\\
\end{tabular}

\vspace{10pt}

\begin{tabular}{rl}
\multicolumn{2}{l}{\textbf{PERFEKT}}\\
& haben + \textbf{geduzt}\\
\end{tabular}
\hfill
\begin{tabular}{rl}
\multicolumn{2}{l}{\textbf{FUTURE I}}\\
& werden + \textbf{duzen}\\
\end{tabular}
\hfill
\begin{tabular}{rl}
\multicolumn{2}{l}{\textbf{PLUSQUAMPERFEKT}}\\
& hatten + \textbf{geduzt}\\
\end{tabular}

\vspace{10pt}

\begin{tabular}{lll}
\multicolumn{3}{l}{\textbf{MIT PRÄPOSITIONEN}}\\
& \textbf{duzen mit} + Dat & (to use "du" with someone)\\
\end{tabular}
\hfill
\begin{tabular}{lll}
\multicolumn{3}{l}{\textbf{TRENNBARE VERBEN}}\\
& ---\\
\end{tabular}
\vspace{-5pt}
\end{verbbox}

\subsection{E Verbs}

%empfehlen (to recommend)
\begin{verbbox}\addcontentsline{toc}{subsubsection}{empfehlen (to recommend)}
{\Large \bf empfehlen} (to recommend) \hfill Type: \textbf{irregular} \hfill Auxiliary Verb: \textbf{haben}

\vspace{5pt}
\hrule
\vspace{5pt}

\begin{tabular}{rl}
\multicolumn{2}{l}{\textbf{PRÄSENS}}\\
ich & \textbf{empfehle}\\
du & \textbf{empfiehlst}\\
er/sie/es & \textbf{empfiehlt}\\
wir & \textbf{empfehlen}\\
ihr & \textbf{empfehlt}\\
sie/Sie & \textbf{empfehlen}\\
\end{tabular}
\hfill
\begin{tabular}{rl}
\multicolumn{2}{l}{\textbf{PRÄTERITUM}}\\
ich & \textbf{empfahl}\\
du & \textbf{empfahlst}\\
er/sie/es & \textbf{empfahl}\\
wir & \textbf{empfahlen}\\
ihr & \textbf{empfahlt}\\
sie/Sie & \textbf{empfahlen}\\
\end{tabular}
\hfill
\begin{tabular}{rl}
\multicolumn{2}{l}{\textbf{IMPERATIV}}\\
& \textbf{empfiehl (du)!}\\
& \textbf{empfehlen wir!}\\
& \textbf{empfehlt (ihr)!}\\
& \textbf{empfehlen Sie!}\\
\end{tabular}

\vspace{10pt}

\begin{tabular}{rl}
\multicolumn{2}{l}{\textbf{PERFEKT}}\\
& haben + \textbf{empfohlen}\\
\end{tabular}
\hfill
\begin{tabular}{rl}
\multicolumn{2}{l}{\textbf{FUTURE I}}\\
& werden + \textbf{empfehlen}\\
\end{tabular}
\hfill
\begin{tabular}{rl}
\multicolumn{2}{l}{\textbf{PLUSQUAMPERFEKT}}\\
& hatten + \textbf{empfohlen}\\
\end{tabular}

\vspace{10pt}

\begin{tabular}{lll}
\multicolumn{3}{l}{\textbf{MIT PRÄPOSITIONEN}}\\
& \textbf{empfehlen für} + Akk & (to recommend for)\\
& \textbf{empfehlen an} + Akk & (to recommend to)\\
\end{tabular}
\hfill
\begin{tabular}{lll}
\multicolumn{3}{l}{\textbf{TRENNBARE VERBEN}}\\
& ---\\
\end{tabular}
\vspace{-5pt}
\end{verbbox}

%engagieren (to hire / engage)
\begin{verbbox-acc}\addcontentsline{toc}{subsubsection}{engagieren (to hire / engage)}
{\Large \bf engagieren} (to hire / engage) \hfill Type: \textbf{regular, + Akkusativ} \hfill Auxiliary Verb: \textbf{haben}

\vspace{5pt}
\hrule
\vspace{5pt}

\begin{tabular}{rl}
\multicolumn{2}{l}{\textbf{PRÄSENS}}\\
ich & \textbf{engagiere}\\
du & \textbf{engagierst}\\
er/sie/es & \textbf{engagiert}\\
wir & \textbf{engagieren}\\
ihr & \textbf{engagiert}\\
sie/Sie & \textbf{engagieren}\\
\end{tabular}
\hfill
\begin{tabular}{rl}
\multicolumn{2}{l}{\textbf{PRÄTERITUM}}\\
ich & \textbf{engagierte}\\
du & \textbf{engagiertest}\\
er/sie/es & \textbf{engagierte}\\
wir & \textbf{engagierten}\\
ihr & \textbf{engagiertet}\\
sie/Sie & \textbf{engagierten}\\
\end{tabular}
\hfill
\begin{tabular}{rl}
\multicolumn{2}{l}{\textbf{IMPERATIV}}\\
& \textbf{engagiere (du)!}\\
& \textbf{engagieren wir!}\\
& \textbf{engagiert (ihr)!}\\
& \textbf{engagieren Sie!}\\
\end{tabular}

\vspace{10pt}

\begin{tabular}{rl}
\multicolumn{2}{l}{\textbf{PERFEKT}}\\
& haben + \textbf{engagiert}\\
\end{tabular}
\hfill
\begin{tabular}{rl}
\multicolumn{2}{l}{\textbf{FUTURE I}}\\
& werden + \textbf{engagieren}\\
\end{tabular}
\hfill
\begin{tabular}{rl}
\multicolumn{2}{l}{\textbf{PLUSQUAMPERFEKT}}\\
& hatten + \textbf{engagiert}\\
\end{tabular}

\vspace{10pt}

\begin{tabular}{lll}
\multicolumn{3}{l}{\textbf{MIT PRÄPOSITIONEN}}\\
& ---\\
\end{tabular}
\hfill
\begin{tabular}{lll}
\multicolumn{3}{l}{\textbf{TRENNBARE VERBEN}}\\
& ---\\
\end{tabular}

\vspace{-5pt}
\end{verbbox-acc}

%entdecken (to discover)
\begin{verbbox-acc}\addcontentsline{toc}{subsubsection}{entdecken (to discover)}
{\Large \bf entdecken} (to discover) \hfill Type: \textbf{regular, + Akkusativ} \hfill Auxiliary Verb: \textbf{haben}

\vspace{5pt}
\hrule
\vspace{5pt}

\begin{tabular}{rl}
\multicolumn{2}{l}{\textbf{PRÄSENS}}\\
ich & \textbf{entdecke}\\
du & \textbf{entdeckst}\\
er/sie/es & \textbf{entdeckt}\\
wir & \textbf{entdecken}\\
ihr & \textbf{entdeckt}\\
sie/Sie & \textbf{entdecken}\\
\end{tabular}
\hfill
\begin{tabular}{rl}
\multicolumn{2}{l}{\textbf{PRÄTERITUM}}\\
ich & \textbf{entdeckte}\\
du & \textbf{entdecktest}\\
er/sie/es & \textbf{entdeckte}\\
wir & \textbf{entdeckten}\\
ihr & \textbf{entdecktet}\\
sie/Sie & \textbf{entdeckten}\\
\end{tabular}
\hfill
\begin{tabular}{rl}
\multicolumn{2}{l}{\textbf{IMPERATIV}}\\
& \textbf{entdecke (du)!}\\
& \textbf{entdecken wir!}\\
& \textbf{entdeckt (ihr)!}\\
& \textbf{entdecken Sie!}\\
\end{tabular}

\vspace{10pt}

\begin{tabular}{rl}
\multicolumn{2}{l}{\textbf{PERFEKT}}\\
& haben + \textbf{entdeckt}\\
\end{tabular}
\hfill
\begin{tabular}{rl}
\multicolumn{2}{l}{\textbf{FUTURE I}}\\
& werden + \textbf{entdecken}\\
\end{tabular}
\hfill
\begin{tabular}{rl}
\multicolumn{2}{l}{\textbf{PLUSQUAMPERFEKT}}\\
& hatten + \textbf{entdeckt}\\
\end{tabular}

\vspace{10pt}

\begin{tabular}{lll}
\multicolumn{3}{l}{\textbf{MIT PRÄPOSITIONEN}}\\
& ---\\
\end{tabular}
\hfill
\begin{tabular}{lll}
\multicolumn{3}{l}{\textbf{TRENNBARE VERBEN}}\\
& ---\\
\end{tabular}

\vspace{-5pt}
\end{verbbox-acc}


%entfernen (to remove)
\begin{verbbox-acc}\addcontentsline{toc}{subsubsection}{entfernen (to remove)}
{\Large \bf entfernen} (to remove) \hfill Type: \textbf{regular, + Akkusativ} \hfill Auxiliary Verb: \textbf{haben}

\vspace{5pt}
\hrule
\vspace{5pt}

\begin{tabular}{rl}
\multicolumn{2}{l}{\textbf{PRÄSENS}}\\
ich & \textbf{entferne}\\
du & \textbf{entfernst}\\
er/sie/es & \textbf{entfernt}\\
wir & \textbf{entfernen}\\
ihr & \textbf{entfernt}\\
sie/Sie & \textbf{entfernen}\\
\end{tabular}
\hfill
\begin{tabular}{rl}
\multicolumn{2}{l}{\textbf{PRÄTERITUM}}\\
ich & \textbf{entfernte}\\
du & \textbf{entferntest}\\
er/sie/es & \textbf{entfernte}\\
wir & \textbf{entfernten}\\
ihr & \textbf{entferntet}\\
sie/Sie & \textbf{entfernten}\\
\end{tabular}
\hfill
\begin{tabular}{rl}
\multicolumn{2}{l}{\textbf{IMPERATIV}}\\
& \textbf{entferne (du)!}\\
& \textbf{entfernen wir!}\\
& \textbf{entfernt (ihr)!}\\
& \textbf{entfernen Sie!}\\
\end{tabular}

\vspace{10pt}

\begin{tabular}{rl}
\multicolumn{2}{l}{\textbf{PERFEKT}}\\
& haben + \textbf{entfernt}\\
\end{tabular}
\hfill
\begin{tabular}{rl}
\multicolumn{2}{l}{\textbf{FUTURE I}}\\
& werden + \textbf{entfernen}\\
\end{tabular}
\hfill
\begin{tabular}{rl}
\multicolumn{2}{l}{\textbf{PLUSQUAMPERFEKT}}\\
& hatten + \textbf{entfernt}\\
\end{tabular}

\vspace{10pt}

\begin{tabular}{lll}
\multicolumn{3}{l}{\textbf{MIT PRÄPOSITIONEN}}\\
& \textbf{entfernen aus} + Dat & (to remove from)\\
\end{tabular}
\hfill
\begin{tabular}{lll}
\multicolumn{3}{l}{\textbf{TRENNBARE VERBEN}}\\
& ---\\
\end{tabular}

\vspace{-5pt}
\end{verbbox-acc}

%entlassen (to dismiss / to release)
\begin{verbbox}\addcontentsline{toc}{subsubsection}{entlassen (to dismiss / to release)}
    {\Large \bf entlassen} (to dismiss / to release) \hfill Type: \textbf{irregular} \hfill Auxiliary Verb: \textbf{haben} 
    
     \vspace{5pt}

    \hrule
    
    \vspace{5pt}

  \begin{tabular}{rl}
     \multicolumn{2}{l}{\textbf{PRÄSENS} }\\
    ich & \textbf{entlasse}\\
    du & \textbf{entlässt}\\
    er/sie/es & \textbf{entlässt}\\
    wir & \textbf{entlassen}\\
    ihr & \textbf{entlasst}\\
    sie/Sie & \textbf{entlassen}\\
  \end{tabular}
  \hfill
     \begin{tabular}{rl}
    \multicolumn{2}{l}{\textbf{PRÄTERITUM} }\\
    ich & \textbf{entließ}\\
    du & \textbf{entließt}\\
    er/sie/es & \textbf{entließ}\\
    wir & \textbf{entließen}\\
    ihr & \textbf{entließt}\\
    sie/Sie & \textbf{entließen}\\
  \end{tabular}
  \hfill
  \begin{tabular}{rl}
     \multicolumn{2}{l}{\textbf{IMPERATIV} }\\
   & \textbf{entlass(e) (du)!}\\
   & \textbf{entlassen wir!}\\
   & \textbf{entlasst (ihr)!}\\
   & \textbf{entlassen Sie!}\\
   & \vspace{10pt}
  \end{tabular}

    \vspace{10pt}

 \begin{tabular}{rl}
     \multicolumn{2}{l}{\textbf{PERFEKT}}\\
   & haben + \textbf{entlassen}\\
  \end{tabular}
\hfill
   \begin{tabular}{rl}
    \multicolumn{2}{l}{\textbf{FUTURE I} }\\
   & werden + \textbf{entlassen}\\
  \end{tabular}
\hfill
   \begin{tabular}{rl}
      \multicolumn{2}{l}{\textbf{PLUSQUAMPERFEKT}}\\
   & hatten + \textbf{entlassen}\\
  \end{tabular}

   \vspace{10pt}

   \begin{tabular}{lll}
      \multicolumn{3}{l}{\textbf{MIT PRÄPOSITIONEN}}\\
& \textbf{entlassen aus} + Dat & (to release from / dismiss from)\\
& \textbf{entlassen wegen} + Gen & (to dismiss because of)\\
\end{tabular}
  \hfill
   \begin{tabular}{lll}
      \multicolumn{3}{l}{\textbf{TRENNBARE VERBEN}}\\
 & \textbf{---} & (none common)\\
\vspace{10pt}
  \end{tabular}

  \vspace{-5pt}

\end{verbbox}
%entscheiden (A2) – to decide
\begin{verbbox-acc}\addcontentsline{toc}{subsubsection}{entscheiden (to decide)}
    {\Large \bf entscheiden} (to decide) \hfill Type: \textbf{irregular, + Akkusativ} \hfill Auxiliary Verb: \textbf{haben} 
    
     \vspace{5pt}

    \hrule
    
    \vspace{5pt}

  \begin{tabular}{rl}
     \multicolumn{2}{l}{\textbf{PRÄSENS} }\\
    ich & \textbf{entscheide}\\
    du & \textbf{entscheidest}\\
    er/sie/es & \textbf{entscheidet}\\
    wir & \textbf{entscheiden}\\
    ihr & \textbf{entscheidet}\\
    sie/Sie & \textbf{entscheiden}\\
  \end{tabular}
  \hfill
     \begin{tabular}{rl}
    \multicolumn{2}{l}{\textbf{PRÄTERITUM} }\\
    ich & \textbf{entschied}\\
    du & \textbf{entschiedst}\\
    er/sie/es & \textbf{entschied}\\
    wir & \textbf{entschieden}\\
    ihr & \textbf{entschiedet}\\
    sie/Sie & \textbf{entschieden}\\
  \end{tabular}
  \hfill
  \begin{tabular}{rl}
     \multicolumn{2}{l}{\textbf{IMPERATIV} }\\
   & \textbf{entscheid(e) (du)!}\\
   & \textbf{entscheiden wir!}\\
   & \textbf{entscheidet (ihr)!}\\
   & \textbf{entscheiden Sie!}\\
   & \vspace{10pt}
  \end{tabular}

    \vspace{10pt}

 \begin{tabular}{rl}
     \multicolumn{2}{l}{\textbf{PERFEKT}}\\
  &  haben + \textbf{entschieden}\\
  \end{tabular}
\hfill
   \begin{tabular}{rl}
   
    \multicolumn{2}{l}{\textbf{FUTURE I} }\\
  &  werden + \textbf{entscheiden}\\
 
  \end{tabular}
\hfill
   \begin{tabular}{rl}
      \multicolumn{2}{l}{\textbf{PLUSQUAMPERFEKT}}\\
  &  hatten + \textbf{entschieden}\\
  \end{tabular}
  
   \vspace{10pt}

\begin{tabular}{lll}
\multicolumn{3}{l}{\textbf{MIT PRÄPOSITIONEN}}\\
& ---\\
\end{tabular}
\hfill
\begin{tabular}{lll}
\multicolumn{3}{l}{\textbf{TRENNBARE VERBEN}}\\
& ---\\
\end{tabular}

  \vspace{-5pt}


\end{verbbox-acc}

%entsorgen (to dispose of)
\begin{verbbox}\addcontentsline{toc}{subsubsection}{entsorgen (to dispose of)}
    {\Large \bf entsorgen} (to dispose of) \hfill Type: \textbf{regular} \hfill Auxiliary Verb: \textbf{haben} 
    
     \vspace{5pt}

    \hrule
    
    \vspace{5pt}

  \begin{tabular}{rl}
     \multicolumn{2}{l}{\textbf{PRÄSENS} }\\
    ich & \textbf{entsorge}\\
    du & \textbf{entsorgst}\\
    er/sie/es & \textbf{entsorgt}\\
    wir & \textbf{entsorgen}\\
    ihr & \textbf{entsorgt}\\
    sie/Sie & \textbf{entsorgen}\\
  \end{tabular}
  \hfill
     \begin{tabular}{rl}
    \multicolumn{2}{l}{\textbf{PRÄTERITUM} }\\
    ich & \textbf{entsorgte}\\
    du & \textbf{entsorgtest}\\
    er/sie/es & \textbf{entsorgte}\\
    wir & \textbf{entsorgten}\\
    ihr & \textbf{entsorgtet}\\
    sie/Sie & \textbf{entsorgten}\\
  \end{tabular}
  \hfill
  \begin{tabular}{rl}
     \multicolumn{2}{l}{\textbf{IMPERATIV} }\\
   & \textbf{entsorg(e) (du)!}\\
   & \textbf{entsorgen wir!}\\
   & \textbf{entsorgt (ihr)!}\\
   & \textbf{entsorgen Sie!}\\
   & \vspace{10pt}
  \end{tabular}

    \vspace{10pt}

 \begin{tabular}{rl}
     \multicolumn{2}{l}{\textbf{PERFEKT}}\\
   & haben + \textbf{entsorgt}\\
  \end{tabular}
\hfill
   \begin{tabular}{rl}
    \multicolumn{2}{l}{\textbf{FUTURE I} }\\
   & werde + \textbf{entsorgen}\\
  \end{tabular}
\hfill
   \begin{tabular}{rl}
      \multicolumn{2}{l}{\textbf{PLUSQUAMPERFEKT}}\\
   & hatten + \textbf{entsorgt}\\
  \end{tabular}

   \vspace{10pt}

   \begin{tabular}{lll}
      \multicolumn{3}{l}{\textbf{MIT PRÄPOSITIONEN}}\\
& \textbf{entsorgen über} + Akk & (to dispose of via)\\
& \textbf{entsorgen in} + Dat & (to dispose of in)\\
\end{tabular}
  \hfill
   \begin{tabular}{lll}
      \multicolumn{3}{l}{\textbf{TRENNBARE VERBEN}}\\
 & \textbf{---} & (none common)\\
\vspace{10pt}
  \end{tabular}

  \vspace{-5pt}

\end{verbbox}

%entstehen (B1) – to arise
\begin{verbbox}\addcontentsline{toc}{subsubsection}{entstehen (to arise)}
    {\Large \bf entstehen} (to arise) \hfill Type: \textbf{irregular} \hfill Auxiliary Verb: \textbf{sein} 
    
     \vspace{5pt}

    \hrule
    
    \vspace{5pt}

  \begin{tabular}{rl}
     \multicolumn{2}{l}{\textbf{PRÄSENS} }\\
    ich & \textbf{entstehe}\\
    du & \textbf{entstehst}\\
    er/sie/es & \textbf{entsteht}\\
    wir & \textbf{entstehen}\\
    ihr & \textbf{entsteht}\\
    sie/Sie & \textbf{entstehen}\\
  \end{tabular}
  \hfill
     \begin{tabular}{rl}
    \multicolumn{2}{l}{\textbf{PRÄTERITUM} }\\
    ich & \textbf{entstand}\\
    du & \textbf{entstandest}\\
    er/sie/es & \textbf{entstand}\\
    wir & \textbf{entstanden}\\
    ihr & \textbf{entstandet}\\
    sie/Sie & \textbf{entstanden}\\
  \end{tabular}
  \hfill
  \begin{tabular}{rl}
     \multicolumn{2}{l}{\textbf{IMPERATIV} }\\
   & \textbf{entsteh(e) (du)!}\\
   & \textbf{entstehen wir!}\\
   & \textbf{entsteht (ihr)!}\\
   & \textbf{entstehen Sie!}\\
   & \vspace{10pt}
  \end{tabular}

    \vspace{10pt}

 \begin{tabular}{ll}
     \multicolumn{2}{l}{\textbf{PERFEKT}}\\
   & sein + \textbf{entstanden}\\
  \end{tabular}
\hfill
   \begin{tabular}{rl}
    \multicolumn{2}{l}{\textbf{FUTURE I} }\\
   & werden + \textbf{entstehen}\\
 
  \end{tabular}
\hfill
   \begin{tabular}{rl}
      \multicolumn{2}{l}{\textbf{PLUSQUAMPERFEKT}}\\
   & waren + \textbf{entstanden}\\
  \end{tabular}
  
\vspace{10pt}

\begin{tabular}{lll}
\multicolumn{3}{l}{\textbf{MIT PRÄPOSITIONEN}}\\
& --- & \\
\end{tabular}
\hfill
\begin{tabular}{lll}
\multicolumn{3}{l}{\textbf{TRENNBARE VERBEN}}\\
& --- & \\
\end{tabular}

  \vspace{10pt}

\begin{tabular}{lll}
\multicolumn{3}{l}{\textbf{MIT PRÄPOSITIONEN}}\\
& ---\\
\end{tabular}
\hfill
\begin{tabular}{lll}
\multicolumn{3}{l}{\textbf{TRENNBARE VERBEN}}\\
& ---\\
\end{tabular}

  \vspace{-5pt}


\end{verbbox}

%entzünden (to ignite / inflame)
\begin{verbbox}\addcontentsline{toc}{subsubsection}{entzünden (to ignite / inflame)}
{\Large \bf entzünden} (to ignite / inflame) \hfill Type: \textbf{regular} \hfill Auxiliary Verb: \textbf{haben}

\vspace{5pt}
\hrule
\vspace{5pt}

\begin{tabular}{rl}
\multicolumn{2}{l}{\textbf{PRÄSENS}}\\
ich & \textbf{entzünde}\\
du & \textbf{entzündest}\\
er/sie/es & \textbf{entzündet}\\
wir & \textbf{entzünden}\\
ihr & \textbf{entzündet}\\
sie/Sie & \textbf{entzünden}\\
\end{tabular}
\hfill
\begin{tabular}{rl}
\multicolumn{2}{l}{\textbf{PRÄTERITUM}}\\
ich & \textbf{entzündete}\\
du & \textbf{entzündetest}\\
er/sie/es & \textbf{entzündete}\\
wir & \textbf{entzündeten}\\
ihr & \textbf{entzündetet}\\
sie/Sie & \textbf{entzündeten}\\
\end{tabular}
\hfill
\begin{tabular}{rl}
\multicolumn{2}{l}{\textbf{IMPERATIV}}\\
& \textbf{entzünde (du)!}\\
& \textbf{entzünden wir!}\\
& \textbf{entzündet (ihr)!}\\
& \textbf{entzünden Sie!}\\
\end{tabular}

\vspace{10pt}

\begin{tabular}{rl}
\multicolumn{2}{l}{\textbf{PERFEKT}}\\
& haben + \textbf{entzündet}\\
\end{tabular}
\hfill
\begin{tabular}{rl}
\multicolumn{2}{l}{\textbf{FUTURE I}}\\
& werden + \textbf{entzünden}\\
\end{tabular}
\hfill
\begin{tabular}{rl}
\multicolumn{2}{l}{\textbf{PLUSQUAMPERFEKT}}\\
& hatten + \textbf{entzündet}\\
\end{tabular}

\vspace{10pt}

\begin{tabular}{lll}
\multicolumn{3}{l}{\textbf{MIT PRÄPOSITIONEN}}\\
& \textbf{entzünden an} + Dat & (to ignite on / by)\\
\end{tabular}
\hfill
\begin{tabular}{lll}
\multicolumn{3}{l}{\textbf{TRENNBARE VERBEN}}\\
& --- & \\
\end{tabular}
\vspace{-5pt}
\end{verbbox}


%erfahren (B1) – to find out, experience
\begin{verbbox-acc}\addcontentsline{toc}{subsubsection}{erfahren (to experience)}
    {\Large \bf erfahren} (to experience) \hfill Type: \textbf{irregular, + Akkusativ} \hfill Auxiliary Verb: \textbf{haben} 
    
     \vspace{5pt}

    \hrule
    
    \vspace{5pt}

  \begin{tabular}{rl}
     \multicolumn{2}{l}{\textbf{PRÄSENS} }\\
    ich & \textbf{erfahre}\\
    du & \textbf{erf\"ahrst}\\
    er/sie/es & \textbf{erf\"ahrt}\\
    wir & \textbf{erfahren}\\
    ihr & \textbf{erfahrt}\\
    sie/Sie & \textbf{erfahren}\\
  \end{tabular}
  \hfill
     \begin{tabular}{rl}
    \multicolumn{2}{l}{\textbf{PRÄTERITUM} }\\
    ich & \textbf{erfuhr}\\
    du & \textbf{erfuhrst}\\
    er/sie/es & \textbf{erfuhr}\\
    wir & \textbf{erfuhren}\\
    ihr & \textbf{erfuhrt}\\
    sie/Sie & \textbf{erfuhren}\\
  \end{tabular}
  \hfill
  \begin{tabular}{rl}
     \multicolumn{2}{l}{\textbf{IMPERATIV} }\\
   & \textbf{erfahr(e) (du)!}\\
   & \textbf{erfahren wir!}\\
   & \textbf{erfahrt (ihr)!}\\
   & \textbf{erfahren Sie!}\\
   & \vspace{10pt}
  \end{tabular}

    \vspace{10pt}

 \begin{tabular}{rl}
     \multicolumn{2}{l}{\textbf{PERFEKT}}\\
  &  haben + \textbf{erfahren}\\
  \end{tabular}
\hfill
   \begin{tabular}{rl}
   
    \multicolumn{2}{l}{\textbf{FUTURE I} }\\
  &  werden + \textbf{erfahren}\\
 
  \end{tabular}
\hfill
   \begin{tabular}{rl}
      \multicolumn{2}{l}{\textbf{PLUSQUAMPERFEKT}}\\
  &  hatten + \textbf{erfahren}\\
  \end{tabular}
  
  \vspace{-5pt}

\end{verbbox-acc}

%ergänzen (to complete / supplement)
\begin{verbbox}\addcontentsline{toc}{subsubsection}{ergänzen (to complete / supplement)}
{\Large \bf ergänzen} (to complete / supplement) \hfill Type: \textbf{regular} \hfill Auxiliary Verb: \textbf{haben} 

\vspace{5pt}
\hrule
\vspace{5pt}

\begin{tabular}{rl}
\multicolumn{2}{l}{\textbf{PRÄSENS}}\\
ich & \textbf{ergänze}\\
du & \textbf{ergänzt}\\
er/sie/es & \textbf{ergänzt}\\
wir & \textbf{ergänzen}\\
ihr & \textbf{ergänzt}\\
sie/Sie & \textbf{ergänzen}\\
\end{tabular}
\hfill
\begin{tabular}{rl}
\multicolumn{2}{l}{\textbf{PRÄTERITUM}}\\
ich & \textbf{ergänzte}\\
du & \textbf{ergänztest}\\
er/sie/es & \textbf{ergänzte}\\
wir & \textbf{ergänzten}\\
ihr & \textbf{ergänztet}\\
sie/Sie & \textbf{ergänzten}\\
\end{tabular}
\hfill
\begin{tabular}{rl}
\multicolumn{2}{l}{\textbf{IMPERATIV}}\\
& \textbf{ergänze (du)!}\\
& \textbf{ergänzen wir!}\\
& \textbf{ergänzt (ihr)!}\\
& \textbf{ergänzen Sie!}\\
\end{tabular}

\vspace{10pt}

\begin{tabular}{rl}
\multicolumn{2}{l}{\textbf{PERFEKT}}\\
& haben + \textbf{ergänzt}\\
\end{tabular}
\hfill
\begin{tabular}{rl}
\multicolumn{2}{l}{\textbf{FUTURE I}}\\
& werden + \textbf{ergänzen}\\
\end{tabular}
\hfill
\begin{tabular}{rl}
\multicolumn{2}{l}{\textbf{PLUSQUAMPERFEKT}}\\
& hatten + \textbf{ergänzt}\\
\end{tabular}

\vspace{10pt}

\begin{tabular}{lll}
\multicolumn{3}{l}{\textbf{MIT PRÄPOSITIONEN}}\\
& \textbf{ergänzen} + Akk & (to complete something)\\
\end{tabular}
\hfill
\begin{tabular}{lll}
\multicolumn{3}{l}{\textbf{TRENNBARE VERBEN}}\\
& ---\\
\end{tabular}

\vspace{-5pt}
\end{verbbox}

%erhalten (to receive / to preserve)
\begin{verbbox}\addcontentsline{toc}{subsubsection}{erhalten (to receive / to preserve)}
    {\Large \bf erhalten} (to receive / to preserve) \hfill Type: \textbf{irregular} \hfill Auxiliary Verb: \textbf{haben} 
    
     \vspace{5pt}

    \hrule
    
    \vspace{5pt}

  \begin{tabular}{rl}
     \multicolumn{2}{l}{\textbf{PRÄSENS} }\\
    ich & \textbf{erhalte}\\
    du & \textbf{erhältst}\\
    er/sie/es & \textbf{erhält}\\
    wir & \textbf{erhalten}\\
    ihr & \textbf{erhaltet}\\
    sie/Sie & \textbf{erhalten}\\
  \end{tabular}
  \hfill
     \begin{tabular}{rl}
    \multicolumn{2}{l}{\textbf{PRÄTERITUM} }\\
    ich & \textbf{erhielt}\\
    du & \textbf{erhieltest}\\
    er/sie/es & \textbf{erhielt}\\
    wir & \textbf{erhielten}\\
    ihr & \textbf{erhieltet}\\
    sie/Sie & \textbf{erhielten}\\
  \end{tabular}
  \hfill
  \begin{tabular}{rl}
     \multicolumn{2}{l}{\textbf{IMPERATIV} }\\
   & \textbf{erhalte (du)!}\\
   & \textbf{erhalten wir!}\\
   & \textbf{erhaltet (ihr)!}\\
   & \textbf{erhalten Sie!}\\
   & \vspace{10pt}
  \end{tabular}

    \vspace{10pt}

 \begin{tabular}{rl}
     \multicolumn{2}{l}{\textbf{PERFEKT}}\\
   & haben + \textbf{erhalten}\\
  \end{tabular}
\hfill
   \begin{tabular}{rl}
    \multicolumn{2}{l}{\textbf{FUTURE I} }\\
   & werde + \textbf{erhalten}\\
  \end{tabular}
\hfill
   \begin{tabular}{rl}
      \multicolumn{2}{l}{\textbf{PLUSQUAMPERFEKT}}\\
   & hatten + \textbf{erhalten}\\
  \end{tabular}

   \vspace{10pt}

   \begin{tabular}{lll}
      \multicolumn{3}{l}{\textbf{MIT PRÄPOSITIONEN}}\\
& \textbf{erhalten von} + Dat & (to receive from)\\
& \textbf{erhalten für} + Akk & (to receive for)\\
\end{tabular}
  \hfill
   \begin{tabular}{lll}
      \multicolumn{3}{l}{\textbf{TRENNBARE VERBEN}}\\
 & \textbf{---} & (none common)\\
\vspace{10pt}
  \end{tabular}

  \vspace{-5pt}

\end{verbbox}

%erinnern (A2) – to remember
\begin{verbbox-acc}\addcontentsline{toc}{subsubsection}{erinnern (to remember)}
    {\Large \bf erinnern} (to remember) \hfill Type: \textbf{irregular, + Akkusativ} \hfill Auxiliary Verb: \textbf{haben} 
    
     \vspace{5pt}

    \hrule
    
    \vspace{5pt}

  \begin{tabular}{rl}
     \multicolumn{2}{l}{\textbf{PRÄSENS} }\\
    ich & \textbf{erinnere}\\
    du & \textbf{erinnerst}\\
    er/sie/es & \textbf{erinnert}\\
    wir & \textbf{erinnern}\\
    ihr & \textbf{erinnert}\\
    sie/Sie & \textbf{erinnern}\\
  \end{tabular}
  \hfill
     \begin{tabular}{rl}
    \multicolumn{2}{l}{\textbf{PRÄTERITUM} }\\
    ich & \textbf{erinnerte}\\
    du & \textbf{erinnertest}\\
    er/sie/es & \textbf{erinnerte}\\
    wir & \textbf{erinnerten}\\
    ihr & \textbf{erinnertet}\\
    sie/Sie & \textbf{erinnerten}\\
  \end{tabular}
  \hfill
  \begin{tabular}{rl}
     \multicolumn{2}{l}{\textbf{IMPERATIV} }\\
   & \textbf{erinner(e) (du)!}\\
   & \textbf{erinnern wir!}\\
   & \textbf{erinnert (ihr)!}\\
   & \textbf{erinnern Sie!}\\
   & \vspace{10pt}
  \end{tabular}

    \vspace{10pt}

 \begin{tabular}{rl}
     \multicolumn{2}{l}{\textbf{PERFEKT}}\\
   & haben + \textbf{erinnert}\\
  \end{tabular}
\hfill
   \begin{tabular}{rl}
   
    \multicolumn{2}{l}{\textbf{FUTURE I} }\\
   & werden + \textbf{erinnern}\\
 
  \end{tabular}
\hfill
   \begin{tabular}{rl}
      \multicolumn{2}{l}{\textbf{PLUSQUAMPERFEKT}}\\
   & hatten + \textbf{erinnert}\\
  \end{tabular}
  
  \vspace{10pt}

\begin{tabular}{lll}
\multicolumn{3}{l}{\textbf{MIT PRÄPOSITIONEN}}\\
& ---\\
\end{tabular}
\hfill
\begin{tabular}{lll}
\multicolumn{3}{l}{\textbf{TRENNBARE VERBEN}}\\
& ---\\
\end{tabular}

  \vspace{-5pt}


\end{verbbox-acc}

%erklären (A2) – to explain
\begin{verbbox-dat}\addcontentsline{toc}{subsubsection}{erklären (to explain)}
    {\Large \bf erklären} (to explain) \hfill Type: \textbf{irregular, + Dativ} \hfill Auxiliary Verb: \textbf{haben} 
    
     \vspace{5pt}

    \hrule
    
    \vspace{5pt}

  \begin{tabular}{rl}
     \multicolumn{2}{l}{\textbf{PRÄSENS} }\\
    ich & \textbf{erkläre}\\
    du & \textbf{erklärst}\\
    er/sie/es & \textbf{erklärt}\\
    wir & \textbf{erklären}\\
    ihr & \textbf{erklärt}\\
    sie/Sie & \textbf{erklären}\\
  \end{tabular}
  \hfill
     \begin{tabular}{rl}
    \multicolumn{2}{l}{\textbf{PRÄTERITUM} }\\
    ich & \textbf{erklärte}\\
    du & \textbf{erklärtest}\\
    er/sie/es & \textbf{erklärte}\\
    wir & \textbf{erklärten}\\
    ihr & \textbf{erklärtet}\\
    sie/Sie & \textbf{erklärten}\\
  \end{tabular}
  \hfill
  \begin{tabular}{rl}
     \multicolumn{2}{l}{\textbf{IMPERATIV} }\\
   & \textbf{erklär(e) (du)!}\\
   & \textbf{erklären wir!}\\
   & \textbf{erklärt (ihr)!}\\
   & \textbf{erklären Sie!}\\
   & \vspace{10pt}
  \end{tabular}

    \vspace{10pt}

 \begin{tabular}{rl}
     \multicolumn{2}{l}{\textbf{PERFEKT}}\\
  &  haben + \textbf{erklärt}\\
  \end{tabular}
\hfill
   \begin{tabular}{rl}
   
    \multicolumn{2}{l}{\textbf{FUTURE I} }\\
  &  werden + \textbf{erklären}\\
 
  \end{tabular}
\hfill
   \begin{tabular}{rl}
      \multicolumn{2}{l}{\textbf{PLUSQUAMPERFEKT}}\\
  &  hatten + \textbf{erklärt}\\
  \end{tabular}
  
  \vspace{-5pt}


\end{verbbox-dat}
%erkunden (to explore)
\begin{verbbox-acc}\addcontentsline{toc}{subsubsection}{erkunden (to explore)}
{\Large \bf erkunden} (to explore) \hfill Type: \textbf{regular, + Akkusativ} \hfill Auxiliary Verb: \textbf{haben}

\vspace{5pt}
\hrule
\vspace{5pt}

\begin{tabular}{rl}
\multicolumn{2}{l}{\textbf{PRÄSENS}}\\
ich & \textbf{erkunde}\\
du & \textbf{erkundest}\\
er/sie/es & \textbf{erkundet}\\
wir & \textbf{erkunden}\\
ihr & \textbf{erkundet}\\
sie/Sie & \textbf{erkunden}\\
\end{tabular}
\hfill
\begin{tabular}{rl}
\multicolumn{2}{l}{\textbf{PRÄTERITUM}}\\
ich & \textbf{erkundete}\\
du & \textbf{erkundetest}\\
er/sie/es & \textbf{erkundete}\\
wir & \textbf{erkundeten}\\
ihr & \textbf{erkundetet}\\
sie/Sie & \textbf{erkundeten}\\
\end{tabular}
\hfill
\begin{tabular}{rl}
\multicolumn{2}{l}{\textbf{IMPERATIV}}\\
& \textbf{erkunde (du)!}\\
& \textbf{erkunden wir!}\\
& \textbf{erkundet (ihr)!}\\
& \textbf{erkunden Sie!}\\
\end{tabular}

\vspace{10pt}

\begin{tabular}{rl}
\multicolumn{2}{l}{\textbf{PERFEKT}}\\
& haben + \textbf{erkundet}\\
\end{tabular}
\hfill
\begin{tabular}{rl}
\multicolumn{2}{l}{\textbf{FUTURE I}}\\
& werden + \textbf{erkunden}\\
\end{tabular}
\hfill
\begin{tabular}{rl}
\multicolumn{2}{l}{\textbf{PLUSQUAMPERFEKT}}\\
& hatten + \textbf{erkundet}\\
\end{tabular}

\vspace{10pt}

\begin{tabular}{lll}
\multicolumn{3}{l}{\textbf{MIT PRÄPOSITIONEN}}\\
& ---\\
\end{tabular}
\hfill
\begin{tabular}{lll}
\multicolumn{3}{l}{\textbf{TRENNBARE VERBEN}}\\
& ---\\
\end{tabular}

\vspace{-5pt}
\end{verbbox-acc}


%erlauben (B1) – to allow
\begin{verbbox-acc}\addcontentsline{toc}{subsubsection}{erlauben (to allow)}
    {\Large \bf erlauben} (to allow) \hfill Type: \textbf{regular, + Akkusativ} \hfill Auxiliary Verb: \textbf{haben} 
    
     \vspace{5pt}

    \hrule
    
    \vspace{5pt}

  \begin{tabular}{rl}
     \multicolumn{2}{l}{\textbf{PRÄSENS} }\\
    ich & \textbf{erlaube}\\
    du & \textbf{erlaubst}\\
    er/sie/es & \textbf{erlaubt}\\
    wir & \textbf{erlauben}\\
    ihr & \textbf{erlaubt}\\
    sie/Sie & \textbf{erlauben}\\
  \end{tabular}
  \hfill
     \begin{tabular}{rl}
    \multicolumn{2}{l}{\textbf{PRÄTERITUM} }\\
    ich & \textbf{erlaubte}\\
    du & \textbf{erlaubtest}\\
    er/sie/es & \textbf{erlaubte}\\
    wir & \textbf{erlaubten}\\
    ihr & \textbf{erlaubtet}\\
    sie/Sie & \textbf{erlaubten}\\
  \end{tabular}
  \hfill
  \begin{tabular}{rl}
     \multicolumn{2}{l}{\textbf{IMPERATIV} }\\
   & \textbf{erlaub(e) (du)!}\\
   & \textbf{erlauben wir!}\\
   & \textbf{erlaubt (ihr)!}\\
   & \textbf{erlauben Sie!}\\
   & \vspace{10pt}
  \end{tabular}

    \vspace{10pt}

 \begin{tabular}{rl}
     \multicolumn{2}{l}{\textbf{PERFEKT}}\\
   & haben + \textbf{erlaubt}\\
  \end{tabular}
\hfill
   \begin{tabular}{rl}
   
    \multicolumn{2}{l}{\textbf{FUTURE I} }\\
  &  werden + \textbf{erlauben}\\
 
  \end{tabular}
\hfill
   \begin{tabular}{rl}
      \multicolumn{2}{l}{\textbf{PLUSQUAMPERFEKT}}\\
  &  hatten + \textbf{erlaubt}\\
  \end{tabular}
  
  \vspace{10pt}

\begin{tabular}{lll}
\multicolumn{3}{l}{\textbf{MIT PRÄPOSITIONEN}}\\
& ---\\
\end{tabular}
\hfill
\begin{tabular}{lll}
\multicolumn{3}{l}{\textbf{TRENNBARE VERBEN}}\\
& ---\\
\end{tabular}

  \vspace{-5pt}


\end{verbbox-acc}

%erleben (to experience)
\begin{verbbox-acc}\addcontentsline{toc}{subsubsection}{erleben (to experience)}
{\Large \bf erleben} (to experience) \hfill Type: \textbf{regular, + Akkusativ} \hfill Auxiliary Verb: \textbf{haben}

\vspace{5pt}
\hrule
\vspace{5pt}

\begin{tabular}{rl}
\multicolumn{2}{l}{\textbf{PRÄSENS}}\\
ich & \textbf{erlebe}\\
du & \textbf{erlebst}\\
er/sie/es & \textbf{erlebt}\\
wir & \textbf{erleben}\\
ihr & \textbf{erlebt}\\
sie/Sie & \textbf{erleben}\\
\end{tabular}
\hfill
\begin{tabular}{rl}
\multicolumn{2}{l}{\textbf{PRÄTERITUM}}\\
ich & \textbf{erlebte}\\
du & \textbf{erlebtest}\\
er/sie/es & \textbf{erlebte}\\
wir & \textbf{erlebten}\\
ihr & \textbf{erlebtet}\\
sie/Sie & \textbf{erlebten}\\
\end{tabular}
\hfill
\begin{tabular}{rl}
\multicolumn{2}{l}{\textbf{IMPERATIV}}\\
& \textbf{erleb(e) (du)!}\\
& \textbf{erleben wir!}\\
& \textbf{erlebt (ihr)!}\\
& \textbf{erleben Sie!}\\
\end{tabular}

\vspace{10pt}

\begin{tabular}{rl}
\multicolumn{2}{l}{\textbf{PERFEKT}}\\
& haben + \textbf{erlebt}\\
\end{tabular}
\hfill
\begin{tabular}{rl}
\multicolumn{2}{l}{\textbf{FUTURE I}}\\
& werden + \textbf{erleben}\\
\end{tabular}
\hfill
\begin{tabular}{rl}
\multicolumn{2}{l}{\textbf{PLUSQUAMPERFEKT}}\\
& hatten + \textbf{erlebt}\\
\end{tabular}

\vspace{10pt}

\begin{tabular}{lll}
\multicolumn{3}{l}{\textbf{MIT PRÄPOSITIONEN}}\\
\end{tabular}
\hfill
\begin{tabular}{lll}
\multicolumn{3}{l}{\textbf{TRENNBARE VERBEN}}\\
& ---\\
\end{tabular}

\vspace{-5pt}
\end{verbbox-acc}


%erledigen (to complete / take care of)
\begin{verbbox}\addcontentsline{toc}{subsubsection}{erledigen (to complete / take care of)}
{\Large \bf erledigen} (to complete / take care of) \hfill Type: \textbf{regular} \hfill Auxiliary Verb: \textbf{haben} 

\vspace{5pt}
\hrule
\vspace{5pt}

\begin{tabular}{rl}
\multicolumn{2}{l}{\textbf{PRÄSENS}}\\
ich & \textbf{erledige}\\
du & \textbf{erledigst}\\
er/sie/es & \textbf{erledigt}\\
wir & \textbf{erledigen}\\
ihr & \textbf{erledigt}\\
sie/Sie & \textbf{erledigen}\\
\end{tabular}
\hfill
\begin{tabular}{rl}
\multicolumn{2}{l}{\textbf{PRÄTERITUM}}\\
ich & \textbf{erledigte}\\
du & \textbf{erledigtest}\\
er/sie/es & \textbf{erledigte}\\
wir & \textbf{erledigten}\\
ihr & \textbf{erledigtet}\\
sie/Sie & \textbf{erledigten}\\
\end{tabular}
\hfill
\begin{tabular}{rl}
\multicolumn{2}{l}{\textbf{IMPERATIV}}\\
& \textbf{erledige (du)!}\\
& \textbf{erledigen wir!}\\
& \textbf{erledigt (ihr)!}\\
& \textbf{erledigen Sie!}\\
\end{tabular}

\vspace{10pt}

\begin{tabular}{rl}
\multicolumn{2}{l}{\textbf{PERFEKT}}\\
& haben + \textbf{erledigt}\\
\end{tabular}
\hfill
\begin{tabular}{rl}
\multicolumn{2}{l}{\textbf{FUTURE I}}\\
& werden + \textbf{erledigen}\\
\end{tabular}
\hfill
\begin{tabular}{rl}
\multicolumn{2}{l}{\textbf{PLUSQUAMPERFEKT}}\\
& hatten + \textbf{erledigt}\\
\end{tabular}

\vspace{10pt}

\begin{tabular}{lll}
\multicolumn{3}{l}{\textbf{MIT PRÄPOSITIONEN}}\\
& \textbf{erledigen} + Akk & (to complete something)\\
\end{tabular}
\hfill
\begin{tabular}{lll}
\multicolumn{3}{l}{\textbf{TRENNBARE VERBEN}}\\
& ---\\
\end{tabular}

\vspace{-5pt}
\end{verbbox}

%erlernen (to learn)
\begin{verbbox-acc}\addcontentsline{toc}{subsubsection}{erlernen (to learn)}
{\Large \bf erlernen} (to learn / acquire a skill) \hfill Type: \textbf{regular, + Akkusativ} \hfill Auxiliary Verb: \textbf{haben}

\vspace{5pt}
\hrule
\vspace{5pt}

\begin{tabular}{rl}
\multicolumn{2}{l}{\textbf{PRÄSENS}}\\
ich & \textbf{erlerne}\\
du & \textbf{erlernst}\\
er/sie/es & \textbf{erlernt}\\
wir & \textbf{erlernen}\\
ihr & \textbf{erlernt}\\
sie/Sie & \textbf{erlernen}\\
\end{tabular}
\hfill
\begin{tabular}{rl}
\multicolumn{2}{l}{\textbf{PRÄTERITUM}}\\
ich & \textbf{erlernte}\\
du & \textbf{erlerntest}\\
er/sie/es & \textbf{erlernte}\\
wir & \textbf{erlernten}\\
ihr & \textbf{erlerntet}\\
sie/Sie & \textbf{erlernten}\\
\end{tabular}
\hfill
\begin{tabular}{rl}
\multicolumn{2}{l}{\textbf{IMPERATIV}}\\
& \textbf{erlern(e) (du)!}\\
& \textbf{erlernen wir!}\\
& \textbf{erlernt (ihr)!}\\
& \textbf{erlernen Sie!}\\
\end{tabular}

\vspace{10pt}

\begin{tabular}{rl}
\multicolumn{2}{l}{\textbf{PERFEKT}}\\
& haben + \textbf{erlernt}\\
\end{tabular}
\hfill
\begin{tabular}{rl}
\multicolumn{2}{l}{\textbf{FUTURE I}}\\
& werden + \textbf{erlernen}\\
\end{tabular}
\hfill
\begin{tabular}{rl}
\multicolumn{2}{l}{\textbf{PLUSQUAMPERFEKT}}\\
& hatten + \textbf{erlernt}\\
\end{tabular}

\vspace{10pt}

\begin{tabular}{lll}
\multicolumn{3}{l}{\textbf{MIT PRÄPOSITIONEN}}\\
& ---\\
\end{tabular}
\hfill
\begin{tabular}{lll}
\multicolumn{3}{l}{\textbf{TRENNBARE VERBEN}}\\
& ---\\
\end{tabular}

\vspace{-5pt}
\end{verbbox-acc}

%eröffnen (to open / inaugurate)
\begin{verbbox}\addcontentsline{toc}{subsubsection}{eröffnen (to open / inaugurate)}
{\Large \bf eröffnen} (to open / inaugurate) \hfill Type: \textbf{regular} \hfill Auxiliary Verb: \textbf{haben}

\vspace{5pt}
\hrule
\vspace{5pt}

\begin{tabular}{rl}
\multicolumn{2}{l}{\textbf{PRÄSENS}}\\
ich & \textbf{eröffne}\\
du & \textbf{eröffnest}\\
er/sie/es & \textbf{eröffnet}\\
wir & \textbf{eröffnen}\\
ihr & \textbf{eröffnet}\\
sie/Sie & \textbf{eröffnen}\\
\end{tabular}
\hfill
\begin{tabular}{rl}
\multicolumn{2}{l}{\textbf{PRÄTERITUM}}\\
ich & \textbf{eröffnete}\\
du & \textbf{eröffnetest}\\
er/sie/es & \textbf{eröffnete}\\
wir & \textbf{eröffneten}\\
ihr & \textbf{eröffnetet}\\
sie/Sie & \textbf{eröffneten}\\
\end{tabular}
\hfill
\begin{tabular}{rl}
\multicolumn{2}{l}{\textbf{IMPERATIV}}\\
& \textbf{eröffne (du)!}\\
& \textbf{eröffnen wir!}\\
& \textbf{eröffnet (ihr)!}\\
& \textbf{eröffnen Sie!}\\
\end{tabular}

\vspace{10pt}

\begin{tabular}{rl}
\multicolumn{2}{l}{\textbf{PERFEKT}}\\
& haben + \textbf{eröffnet}\\
\end{tabular}
\hfill
\begin{tabular}{rl}
\multicolumn{2}{l}{\textbf{FUTURE I}}\\
& werden + \textbf{eröffnen}\\
\end{tabular}
\hfill
\begin{tabular}{rl}
\multicolumn{2}{l}{\textbf{PLUSQUAMPERFEKT}}\\
& hatten + \textbf{eröffnet}\\
\end{tabular}

\vspace{10pt}

\begin{tabular}{lll}
\multicolumn{3}{l}{\textbf{MIT PRÄPOSITIONEN}}\\
& \textbf{eröffnen für} + Akk & (to open for)\\
& \textbf{eröffnen zu} + Dat & (to inaugurate at / on)\\
\end{tabular}
\hfill
\begin{tabular}{lll}
\multicolumn{3}{l}{\textbf{TRENNBARE VERBEN}}\\
& --- & \\
\end{tabular}
\vspace{-5pt}
\end{verbbox}

%erreichen (B1) – to reach, achieve
\begin{verbbox-acc}\addcontentsline{toc}{subsubsection}{erreichen (to reach, achieve)}
    {\Large \bf erreichen} (to reach, achieve) \hfill Type: \textbf{regular, + Akkusativ} \hfill Auxiliary Verb: \textbf{haben} 
    
     \vspace{5pt}

    \hrule
    
    \vspace{5pt}

  \begin{tabular}{rl}
     \multicolumn{2}{l}{\textbf{PRÄSENS} }\\
    ich & \textbf{erreiche}\\
    du & \textbf{erreichst}\\
    er/sie/es & \textbf{erreicht}\\
    wir & \textbf{erreichen}\\
    ihr & \textbf{erreicht}\\
    sie/Sie & \textbf{erreichen}\\
  \end{tabular}
  \hfill
     \begin{tabular}{rl}
    \multicolumn{2}{l}{\textbf{PRÄTERITUM} }\\
    ich & \textbf{erreichte}\\
    du & \textbf{erreichtest}\\
    er/sie/es & \textbf{erreichte}\\
    wir & \textbf{erreichten}\\
    ihr & \textbf{erreichtet}\\
    sie/Sie & \textbf{erreichten}\\
  \end{tabular}
  \hfill
  \begin{tabular}{rl}
     \multicolumn{2}{l}{\textbf{IMPERATIV} }\\
   & \textbf{erreich(e) (du)!}\\
   & \textbf{erreichen wir!}\\
   & \textbf{erreicht (ihr)!}\\
   & \textbf{erreichen Sie!}\\
   & \vspace{10pt}
  \end{tabular}

    \vspace{10pt}

 \begin{tabular}{rl}
     \multicolumn{2}{l}{\textbf{PERFEKT}}\\
  &  haben + \textbf{erreicht}\\
  \end{tabular}
\hfill
   \begin{tabular}{rl}
   
    \multicolumn{2}{l}{\textbf{FUTURE I} }\\
  &  werden + \textbf{erreichen}\\
 
  \end{tabular}
\hfill
   \begin{tabular}{rl}
      \multicolumn{2}{l}{\textbf{PLUSQUAMPERFEKT}}\\
  &  hatten + \textbf{erreicht}\\
  \end{tabular}
  
  \vspace{-5pt}


\end{verbbox-acc}

%erreuern - to renew
\begin{verbbox}\addcontentsline{toc}{subsubsection}{erneuern (to renew / renovate)}
    {\Large \bf erneuern} (to renew / renovate) \hfill Type: \textbf{regular} \hfill Auxiliary Verb: \textbf{haben} 
    
     \vspace{5pt}

    \hrule
    
    \vspace{5pt}

  \begin{tabular}{rl}
     \multicolumn{2}{l}{\textbf{PRÄSENS} }\\
    ich & \textbf{erneuere}\\
    du & \textbf{erneuerst}\\
    er/sie/es & \textbf{erneuert}\\
    wir & \textbf{erneuern}\\
    ihr & \textbf{erneuert}\\
    sie/Sie & \textbf{erneuern}\\
  \end{tabular}
  \hfill
  \begin{tabular}{rl}
    \multicolumn{2}{l}{\textbf{PRÄTERITUM} }\\
    ich & \textbf{erneuerte}\\
    du & \textbf{erneuertest}\\
    er/sie/es & \textbf{erneuerte}\\
    wir & \textbf{erneuerten}\\
    ihr & \textbf{erneuertet}\\
    sie/Sie & \textbf{erneuerten}\\
  \end{tabular}
  \hfill
  \begin{tabular}{rl}
     \multicolumn{2}{l}{\textbf{IMPERATIV} }\\
   & \textbf{erneuere (du)!}\\
   & \textbf{erneuern wir!}\\
   & \textbf{erneuert (ihr)!}\\
   & \textbf{erneuern Sie!}\\
   & \vspace{10pt}
  \end{tabular}

 \vspace{10pt}

 \begin{tabular}{rl}
     \multicolumn{2}{l}{\textbf{PERFEKT}}\\
   & haben + \textbf{erneuert}\\
  \end{tabular}
\hfill
\begin{tabular}{rl}
    \multicolumn{2}{l}{\textbf{FUTURE I}}\\
   & werden + \textbf{erneuern}\\
  \end{tabular}
\hfill
\begin{tabular}{rl}
    \multicolumn{2}{l}{\textbf{PLUSQUAMPERFEKT}}\\
   & hatten + \textbf{erneuert}\\
  \end{tabular}

 \vspace{10pt}

\begin{tabular}{lll}
      \multicolumn{3}{l}{\textbf{MIT PRÄPOSITIONEN}}\\
& \textbf{erneuern von} + Dat & (to renew from / of)\\
& \textbf{erneuern durch} + Akk & (to renew through / by)\\
& \textbf{erneuern zu} + Dat & (to renew into)\\
\end{tabular}
  \hfill
\begin{tabular}{lll}
      \multicolumn{3}{l}{\textbf{TRENNBARE VERBEN}}\\
 & \textbf{ein$|$erneuern} & (to renew in / reinstall)\\
\vspace{20pt}
  \end{tabular}


\vspace{-5pt}
\end{verbbox}

%erscheinen (to appear)
\begin{verbbox-dat}\addcontentsline{toc}{subsubsection}{erscheinen (to appear)}
    {\Large \bf erscheinen} (to appear) \hfill Type: \textbf{irregular, inseparable} \hfill Auxiliary Verb: \textbf{sein} 
    
     \vspace{5pt}

    \hrule
    
    \vspace{5pt}

  \begin{tabular}{rl}
     \multicolumn{2}{l}{\textbf{PRÄSENS} }\\
    ich & \textbf{erscheine}\\
    du & \textbf{erscheinst}\\
    er/sie/es & \textbf{erscheint}\\
    wir & \textbf{erscheinen}\\
    ihr & \textbf{erscheint}\\
    sie/Sie & \textbf{erscheinen}\\
  \end{tabular}
  \hfill
     \begin{tabular}{rl}
    \multicolumn{2}{l}{\textbf{PRÄTERITUM} }\\
    ich & \textbf{erschien}\\
    du & \textbf{erschienst}\\
    er/sie/es & \textbf{erschien}\\
    wir & \textbf{erschienen}\\
    ihr & \textbf{erschient}\\
    sie/Sie & \textbf{erschienen}\\
  \end{tabular}
  \hfill
  \begin{tabular}{rl}
     \multicolumn{2}{l}{\textbf{IMPERATIV} }\\
   & \textbf{erschein(e) (du)!}\\
   & \textbf{erscheinen wir!}\\
   & \textbf{erscheint (ihr)!}\\
   & \textbf{erscheinen Sie!}\\
   & \vspace{10pt}
  \end{tabular}

    \vspace{10pt}

 \begin{tabular}{rl}
     \multicolumn{2}{l}{\textbf{PERFEKT}}\\
   & sein + \textbf{erschienen}\\
  \end{tabular}
\hfill
   \begin{tabular}{rl}
   
    \multicolumn{2}{l}{\textbf{FUTURE I} }\\
   & werden + \textbf{erscheinen}\\
 
  \end{tabular}
\hfill
   \begin{tabular}{rl}
      \multicolumn{2}{l}{\textbf{PLUSQUAMPERFEKT}}\\
   & waren + \textbf{erschienen}\\
  \end{tabular}

   \vspace{10pt}

   \begin{tabular}{lll}
      \multicolumn{3}{l}{\textbf{MIT PRÄPOSITIONEN}}\\
& \textbf{erscheinen in} + Dat & (to appear in)\\
& \textbf{erscheinen vor} + Dat & (to appear before)\\
\end{tabular}
  \hfill
   \begin{tabular}{lll}
      \multicolumn{3}{l}{\textbf{TRENNBARE VERBEN}}\\
 & -- & --\\
\vspace{10pt}
  \end{tabular}


  \vspace{-5pt}


\end{verbbox-dat}

%ersetzen (to replace)
\begin{verbbox}\addcontentsline{toc}{subsubsection}{ersetzen (to replace)}
    {\Large \bf ersetzen} (to replace) \hfill Type: \textbf{regular, inseparable} \hfill Auxiliary Verb: \textbf{haben} 
    
     \vspace{5pt}

    \hrule
    
    \vspace{5pt}

  \begin{tabular}{rl}
     \multicolumn{2}{l}{\textbf{PRÄSENS} }\\
    ich & \textbf{ersetze}\\
    du & \textbf{ersetzt}\\
    er/sie/es & \textbf{ersetzt}\\
    wir & \textbf{ersetzen}\\
    ihr & \textbf{ersetzt}\\
    sie/Sie & \textbf{ersetzen}\\
  \end{tabular}
  \hfill
     \begin{tabular}{rl}
    \multicolumn{2}{l}{\textbf{PRÄTERITUM} }\\
    ich & \textbf{ersetzte}\\
    du & \textbf{ersetztest}\\
    er/sie/es & \textbf{ersetzte}\\
    wir & \textbf{ersetzten}\\
    ihr & \textbf{ersetztet}\\
    sie/Sie & \textbf{ersetzten}\\
  \end{tabular}
  \hfill
  \begin{tabular}{rl}
     \multicolumn{2}{l}{\textbf{IMPERATIV} }\\
   & \textbf{ersetz(e) (du)!}\\
   & \textbf{ersetzen wir!}\\
   & \textbf{ersetzt (ihr)!}\\
   & \textbf{ersetzen Sie!}\\
   & \vspace{10pt}
  \end{tabular}

    \vspace{10pt}

 \begin{tabular}{rl}
     \multicolumn{2}{l}{\textbf{PERFEKT}}\\
   & haben + \textbf{ersetzt}\\
  \end{tabular}
\hfill
   \begin{tabular}{rl}
    \multicolumn{2}{l}{\textbf{FUTURE I} }\\
   & werden + \textbf{ersetzen}\\
  \end{tabular}
\hfill
   \begin{tabular}{rl}
      \multicolumn{2}{l}{\textbf{PLUSQUAMPERFEKT}}\\
   & hatten + \textbf{ersetzt}\\
  \end{tabular}

   \vspace{10pt}

   \begin{tabular}{lll}
      \multicolumn{3}{l}{\textbf{MIT PRÄPOSITIONEN}}\\
& \textbf{ersetzen durch} + Akk & (to replace with / by)\\
& \textbf{ersetzen mit} + Dat & (to replace with)\\
\end{tabular}
  \hfill
   \begin{tabular}{lll}
      \multicolumn{3}{l}{\textbf{TRENNBARE VERBEN}}\\
 & \textbf{---} & (none common)\\
\vspace{10pt}
  \end{tabular}

  \vspace{-5pt}

\end{verbbox}

%erstellen (to create / to prepare)
\begin{verbbox}\addcontentsline{toc}{subsubsection}{erstellen (to create / to prepare)}
    {\Large \bf erstellen} (to create / to prepare) \hfill Type: \textbf{regular, inseparable} \hfill Auxiliary Verb: \textbf{haben} 
    
     \vspace{5pt}

    \hrule
    
    \vspace{5pt}

  \begin{tabular}{rl}
     \multicolumn{2}{l}{\textbf{PRÄSENS} }\\
    ich & \textbf{erstelle}\\
    du & \textbf{erstellst}\\
    er/sie/es & \textbf{erstellt}\\
    wir & \textbf{erstellen}\\
    ihr & \textbf{erstellt}\\
    sie/Sie & \textbf{erstellen}\\
  \end{tabular}
  \hfill
     \begin{tabular}{rl}
    \multicolumn{2}{l}{\textbf{PRÄTERITUM} }\\
    ich & \textbf{erstellte}\\
    du & \textbf{erstelltest}\\
    er/sie/es & \textbf{erstellte}\\
    wir & \textbf{erstellten}\\
    ihr & \textbf{erstelltet}\\
    sie/Sie & \textbf{erstellten}\\
  \end{tabular}
  \hfill
  \begin{tabular}{rl}
     \multicolumn{2}{l}{\textbf{IMPERATIV} }\\
   & \textbf{erstelle (du)!}\\
   & \textbf{erstellen wir!}\\
   & \textbf{erstellt (ihr)!}\\
   & \textbf{erstellen Sie!}\\
   & \vspace{10pt}
  \end{tabular}

    \vspace{10pt}

 \begin{tabular}{rl}
     \multicolumn{2}{l}{\textbf{PERFEKT}}\\
   & haben + \textbf{erstellt}\\
  \end{tabular}
\hfill
   \begin{tabular}{rl}
    \multicolumn{2}{l}{\textbf{FUTURE I} }\\
   & werden + \textbf{erstellen}\\
  \end{tabular}
\hfill
   \begin{tabular}{rl}
      \multicolumn{2}{l}{\textbf{PLUSQUAMPERFEKT}}\\
   & hatten + \textbf{erstellt}\\
  \end{tabular}

   \vspace{10pt}

   \begin{tabular}{lll}
      \multicolumn{3}{l}{\textbf{MIT PRÄPOSITIONEN}}\\
& \textbf{erstellen für} + Akk & (to create for)\\
& \textbf{erstellen mit} + Dat & (to create with)\\
\end{tabular}
  \hfill
   \begin{tabular}{lll}
      \multicolumn{3}{l}{\textbf{TRENNBARE VERBEN}}\\
 & \textbf{---} & (none common)\\
\vspace{10pt}
  \end{tabular}

  \vspace{-5pt}

\end{verbbox}



%erwartet - to expect something
\begin{verbbox}\addcontentsline{toc}{subsubsection}{erwarten (to expect)}
{\Large \bf erwarten} (to expect) \hfill Type: \textbf{regular} \hfill Auxiliary Verb: \textbf{haben}

\vspace{5pt}
\hrule
\vspace{5pt}

\begin{tabular}{rl}
\multicolumn{2}{l}{\textbf{PRÄSENS}}\\
ich & \textbf{erwarte}\\
du & \textbf{erwartest}\\
er/sie/es & \textbf{erwartet}\\
wir & \textbf{erwarten}\\
ihr & \textbf{erwartet}\\
sie/Sie & \textbf{erwarten}\\
\end{tabular}
\hfill
\begin{tabular}{rl}
\multicolumn{2}{l}{\textbf{PRÄTERITUM}}\\
ich & \textbf{erwartete}\\
du & \textbf{erwartetest}\\
er/sie/es & \textbf{erwartete}\\
wir & \textbf{erwarteten}\\
ihr & \textbf{erwartetet}\\
sie/Sie & \textbf{erwarteten}\\
\end{tabular}
\hfill
\begin{tabular}{rl}
\multicolumn{2}{l}{\textbf{IMPERATIV}}\\
& \textbf{erwarte (du)!}\\
& \textbf{erwarten wir!}\\
& \textbf{erwartet (ihr)!}\\
& \textbf{erwarten Sie!}\\
\end{tabular}

\vspace{10pt}

\begin{tabular}{rl}
\multicolumn{2}{l}{\textbf{PERFEKT}}\\
& haben + \textbf{erwartet}\\
\end{tabular}
\hfill
\begin{tabular}{rl}
\multicolumn{2}{l}{\textbf{FUTURE I}}\\
& werden + \textbf{erwarten}\\
\end{tabular}
\hfill
\begin{tabular}{rl}
\multicolumn{2}{l}{\textbf{PLUSQUAMPERFEKT}}\\
& hatten + \textbf{erwartet}\\
\end{tabular}

\vspace{10pt}

\begin{tabular}{lll}
\multicolumn{3}{l}{\textbf{MIT PRÄPOSITIONEN}}\\
& \textbf{erwarten von} + Dat & (to expect from)\\
\end{tabular}
\hfill
\begin{tabular}{lll}
\multicolumn{3}{l}{\textbf{TRENNBARE VERBEN}}\\
& --- & (no separable forms)\\
\end{tabular}
\vspace{-5pt}
\end{verbbox}

%erweitern (to expand)
\begin{verbbox-acc}\addcontentsline{toc}{subsubsection}{erweitern (to expand)}
{\Large \bf erweitern} (to expand) \hfill Type: \textbf{regular, + Akkusativ} \hfill Auxiliary Verb: \textbf{haben}

\vspace{5pt}
\hrule
\vspace{5pt}

\begin{tabular}{rl}
\multicolumn{2}{l}{\textbf{PRÄSENS}}\\
ich & \textbf{erweitere}\\
du & \textbf{erweiterst}\\
er/sie/es & \textbf{erweitert}\\
wir & \textbf{erweitern}\\
ihr & \textbf{erweitert}\\
sie/Sie & \textbf{erweitern}\\
\end{tabular}
\hfill
\begin{tabular}{rl}
\multicolumn{2}{l}{\textbf{PRÄTERITUM}}\\
ich & \textbf{erweiterte}\\
du & \textbf{erweitertest}\\
er/sie/es & \textbf{erweiterte}\\
wir & \textbf{erweiterten}\\
ihr & \textbf{erweitertet}\\
sie/Sie & \textbf{erweiterten}\\
\end{tabular}
\hfill
\begin{tabular}{rl}
\multicolumn{2}{l}{\textbf{IMPERATIV}}\\
& \textbf{erweiter(e) (du)!}\\
& \textbf{erweitern wir!}\\
& \textbf{erweitert (ihr)!}\\
& \textbf{erweitern Sie!}\\
\end{tabular}

\vspace{10pt}

\begin{tabular}{rl}
\multicolumn{2}{l}{\textbf{PERFEKT}}\\
& haben + \textbf{erweitert}\\
\end{tabular}
\hfill
\begin{tabular}{rl}
\multicolumn{2}{l}{\textbf{FUTURE I}}\\
& werden + \textbf{erweitern}\\
\end{tabular}
\hfill
\begin{tabular}{rl}
\multicolumn{2}{l}{\textbf{PLUSQUAMPERFEKT}}\\
& hatten + \textbf{erweitert}\\
\end{tabular}

\vspace{10pt}

\begin{tabular}{lll}
\multicolumn{3}{l}{\textbf{MIT PRÄPOSITIONEN}}\\
& \textbf{erweitern um} + Akk & (to expand by)\\
\end{tabular}
\hfill
\begin{tabular}{lll}
\multicolumn{3}{l}{\textbf{TRENNBARE VERBEN}}\\
& ---\\
\end{tabular}

\vspace{-5pt}
\end{verbbox-acc}

%erzählen (A2) – to tell
\begin{verbbox-dat}\addcontentsline{toc}{subsubsection}{erzählen (to tell)}
    {\Large \bf erzählen} (to tell) \hfill Type: \textbf{regular, + Dativ} \hfill Auxiliary Verb: \textbf{haben} 
    
     \vspace{5pt}

    \hrule
    
    \vspace{5pt}

  \begin{tabular}{rl}
     \multicolumn{2}{l}{\textbf{PRÄSENS} }\\
    ich & \textbf{erzähle}\\
    du & \textbf{erzählst}\\
    er/sie/es & \textbf{erzählt}\\
    wir & \textbf{erzählen}\\
    ihr & \textbf{erzählt}\\
    sie/Sie & \textbf{erzählen}\\
  \end{tabular}
  \hfill
     \begin{tabular}{rl}
    \multicolumn{2}{l}{\textbf{PRÄTERITUM} }\\
    ich & \textbf{erzählte}\\
    du & \textbf{erzähltest}\\
    er/sie/es & \textbf{erzählte}\\
    wir & \textbf{erzählten}\\
    ihr & \textbf{erzähltet}\\
    sie/Sie & \textbf{erzählten}\\
  \end{tabular}
  \hfill
  \begin{tabular}{rl}
     \multicolumn{2}{l}{\textbf{IMPERATIV} }\\
   & \textbf{erzähl(e) (du)!}\\
   & \textbf{erzählen wir!}\\
   & \textbf{erzählt (ihr)!}\\
   & \textbf{erzählen Sie!}\\
   & \vspace{10pt}
  \end{tabular}

    \vspace{10pt}

 \begin{tabular}{rl}
     \multicolumn{2}{l}{\textbf{PERFEKT}}\\
  &  haben + \textbf{erzählt}\\
  \end{tabular}
\hfill
   \begin{tabular}{rl}
   
    \multicolumn{2}{l}{\textbf{FUTURE I} }\\
  &  werden + \textbf{erzählen}\\
 
  \end{tabular}
\hfill
   \begin{tabular}{rl}
      \multicolumn{2}{l}{\textbf{PLUSQUAMPERFEKT}}\\
   & hatten + \textbf{erzählen}\\
  \end{tabular}
  
  \vspace{-5pt}


\end{verbbox-dat}

%erziehen (to raise / to educate)
\begin{verbbox}\addcontentsline{toc}{subsubsection}{erziehen (to raise / to educate)}
    {\Large \bf erziehen} (to raise / to educate) \hfill Type: \textbf{irregular} \hfill Auxiliary Verb: \textbf{haben} 
    
     \vspace{5pt}

    \hrule
    
    \vspace{5pt}

  \begin{tabular}{rl}
     \multicolumn{2}{l}{\textbf{PRÄSENS} }\\
    ich & \textbf{erziehe}\\
    du & \textbf{erziehst}\\
    er/sie/es & \textbf{erzieht}\\
    wir & \textbf{erziehen}\\
    ihr & \textbf{erzieht}\\
    sie/Sie & \textbf{erziehen}\\
  \end{tabular}
  \hfill
     \begin{tabular}{rl}
    \multicolumn{2}{l}{\textbf{PRÄTERITUM} }\\
    ich & \textbf{erzog}\\
    du & \textbf{erzogst}\\
    er/sie/es & \textbf{erzog}\\
    wir & \textbf{erzogen}\\
    ihr & \textbf{erzogt}\\
    sie/Sie & \textbf{erzogen}\\
  \end{tabular}
  \hfill
  \begin{tabular}{rl}
     \multicolumn{2}{l}{\textbf{IMPERATIV} }\\
   & \textbf{erzieh(e) (du)!}\\
   & \textbf{erziehen wir!}\\
   & \textbf{erzieht (ihr)!}\\
   & \textbf{erziehen Sie!}\\
   & \vspace{10pt}
  \end{tabular}

    \vspace{10pt}

 \begin{tabular}{rl}
     \multicolumn{2}{l}{\textbf{PERFEKT}}\\
   & haben + \textbf{erzogen}\\
  \end{tabular}
\hfill
   \begin{tabular}{rl}
    \multicolumn{2}{l}{\textbf{FUTURE I} }\\
   & werde + \textbf{erziehen}\\
  \end{tabular}
\hfill
   \begin{tabular}{rl}
      \multicolumn{2}{l}{\textbf{PLUSQUAMPERFEKT}}\\
   & hatten + \textbf{erzogen}\\
  \end{tabular}

   \vspace{10pt}

   \begin{tabular}{lll}
      \multicolumn{3}{l}{\textbf{MIT PRÄPOSITIONEN}}\\
& \textbf{erziehen zu} + Dat & (to raise/educate to be)\\
& \textbf{erziehen mit} + Dat & (to raise/educate with)\\
\end{tabular}
  \hfill
   \begin{tabular}{lll}
      \multicolumn{3}{l}{\textbf{TRENNBARE VERBEN}}\\
 & \textbf{---} & (none common)\\
\vspace{10pt}
  \end{tabular}

  \vspace{-5pt}

\end{verbbox}


%essen (A1) – to eat
\begin{verbbox-acc}\addcontentsline{toc}{subsubsection}{essen (to eat)}
    {\Large \bf essen} (to eat) \hfill Type: \textbf{irregular, + Akkusativ} \hfill Auxiliary Verb: \textbf{haben} 
    
     \vspace{5pt}

    \hrule
    
    \vspace{5pt}

  \begin{tabular}{rl}
     \multicolumn{2}{l}{\textbf{PRÄSENS} }\\
    ich & \textbf{esse}\\
    du & \textbf{isst}\\
    er/sie/es & \textbf{isst}\\
    wir & \textbf{essen}\\
    ihr & \textbf{esst}\\
    sie/Sie & \textbf{essen}\\
  \end{tabular}
  \hfill
     \begin{tabular}{rl}
    \multicolumn{2}{l}{\textbf{PRÄTERITUM} }\\
    ich & \textbf{aß}\\
    du & \textbf{aßest}\\
    er/sie/es & \textbf{aß}\\
    wir & \textbf{aßen}\\
    ihr & \textbf{aßt}\\
    sie/Sie & \textbf{aßen}\\
  \end{tabular}
  \hfill
  \begin{tabular}{rl}
     \multicolumn{2}{l}{\textbf{IMPERATIV} }\\
   & \textbf{iss (du)!}\\
   & \textbf{essen wir!}\\
   & \textbf{esst (ihr)!}\\
   & \textbf{essen Sie!}\\
   & \vspace{10pt}
  \end{tabular}

    \vspace{10pt}

 \begin{tabular}{rl}
     \multicolumn{2}{l}{\textbf{PERFEKT}}\\
  &  haben + \textbf{gegessen}\\
  \end{tabular}
\hfill
   \begin{tabular}{rl}
   
    \multicolumn{2}{l}{\textbf{FUTURE I} }\\
  &  werden + \textbf{essen}\\
 
  \end{tabular}
\hfill
   \begin{tabular}{rl}
      \multicolumn{2}{l}{\textbf{PLUSQUAMPERFEKT}}\\
   & hatten + \textbf{gegessen}\\
  \end{tabular}
  
  \vspace{-5pt}


\end{verbbox-acc}

\subsection{F Verbs}

%fahren (A1) – to drive, go
\begin{verbbox-acc}\addcontentsline{toc}{subsubsection}{fahren (to drive)}
    {\Large \bf fahren} (to drive) \hfill Type: \textbf{irregular, + Akkusativ} \hfill Auxiliary Verb: \textbf{haben/sein} 
    
     \vspace{5pt}

    \hrule
    
    \vspace{5pt}

  \begin{tabular}{rl}
     \multicolumn{2}{l}{\textbf{PRÄSENS} }\\
    ich & \textbf{fahre}\\
    du & \textbf{f\"ahrst}\\
    er/sie/es & \textbf{f\"ahrt}\\
    wir & \textbf{fahren}\\
    ihr & \textbf{fahrt}\\
    sie/Sie & \textbf{fahren}\\
  \end{tabular}
  \hfill
     \begin{tabular}{rl}
    \multicolumn{2}{l}{\textbf{PRÄTERITUM} }\\
    ich & \textbf{fuhr}\\
    du & \textbf{fuhrst}\\
    er/sie/es & \textbf{fuhr}\\
    wir & \textbf{fuhren}\\
    ihr & \textbf{fuhrt}\\
    sie/Sie & \textbf{fuhren}\\
  \end{tabular}
  \hfill
  \begin{tabular}{rl}
     \multicolumn{2}{l}{\textbf{IMPERATIV} }\\
   & \textbf{fahr(e) (du)!}\\
   & \textbf{fahren wir!}\\
   & \textbf{fahrt (ihr)!}\\
   & \textbf{fahren Sie!}\\
   & \vspace{10pt}
  \end{tabular}

    \vspace{10pt}

 \begin{tabular}{rl}
     \multicolumn{2}{l}{\textbf{PERFEKT}}\\
  &  sein/haben + \textbf{gefahren}\\
  \end{tabular}
\hfill
   \begin{tabular}{rl}
   
    \multicolumn{2}{l}{\textbf{FUTURE I} }\\
  &  werden + \textbf{fahren}\\
 
  \end{tabular}
\hfill
   \begin{tabular}{rl}
      \multicolumn{2}{l}{\textbf{PLUSQUAMPERFEKT}}\\
   & sein/hatten + \textbf{gefahren}\\
  \end{tabular}
  
\vspace{10pt}

   \begin{tabular}{lll}
      \multicolumn{3}{l}{\textbf{MIT PRÄPOSITIONEN}}\\
& \textbf{fahren nach} + Ort & (to go to city / country)\\
&\textbf{fahren in} + Akk & (to go into)\\
&\textbf{fahren zu} + Dat & (to go to place / person)\\
&\textbf{fahren mit} + Dat & (to go with)\\
&\textbf{fahren über} + Akk & (to go via)\\
&\textbf{fahren von} + Akk  \textbf{nach} + Akk & (to go from ... to)\\
&\textbf{fahren durch} + Akk & (to go through)\\
\vspace{20pt}
  \end{tabular}
  \hfill
   \begin{tabular}{lll}
      \multicolumn{3}{l}{\textbf{TRENNBARE VERBEN}}\\
 & \textbf{ab$|$fahren} & (to depart)\\
 & \textbf{an$|$fahren} & (to start driving)\\
  & \textbf{aus$|$fahren} & (to drive out)\\
  & \textbf{ein$|$fahren} & (to enter / pull in)\\
  & \textbf{los$|$fahren} & (to set off)\\
     & \textbf{mit$|$fahren} & (to ride along)\\
     & \textbf{vor$|$fahren} & (to pull up)\\
     & \textbf{weg$|$fahren} & (to drive away)\\
      & \textbf{zurück$|$fahren} & (to bring back)\\
  \end{tabular}


  \vspace{-5pt}

\end{verbbox-acc}

%fallen - to fall
\begin{verbbox}\addcontentsline{toc}{subsubsection}{fallen (to fall)}
    {\Large \bf fallen} (to fall) \hfill Type: \textbf{irregular} \hfill Auxiliary Verb: \textbf{sein} 
    
     \vspace{5pt}

    \hrule
    
    \vspace{5pt}

  \begin{tabular}{rl}
     \multicolumn{2}{l}{\textbf{PRÄSENS} }\\
    ich & \textbf{falle}\\
    du & \textbf{fällst}\\
    er/sie/es & \textbf{fällt}\\
    wir & \textbf{fallen}\\
    ihr & \textbf{fallt}\\
    sie/Sie & \textbf{fallen}\\
  \end{tabular}
  \hfill
     \begin{tabular}{rl}
    \multicolumn{2}{l}{\textbf{PRÄTERITUM} }\\
    ich & \textbf{fiel}\\
    du & \textbf{fielst}\\
    er/sie/es & \textbf{fiel}\\
    wir & \textbf{fielen}\\
    ihr & \textbf{fielt}\\
    sie/Sie & \textbf{fielen}\\
  \end{tabular}
  \hfill
  \begin{tabular}{rl}
     \multicolumn{2}{l}{\textbf{IMPERATIV} }\\
   & \textbf{fall(e) (du)!}\\
   & \textbf{fallen wir!}\\
   & \textbf{fallt (ihr)!}\\
   & \textbf{fallen Sie!}\\
   & \vspace{10pt}
  \end{tabular}

    \vspace{10pt}

 \begin{tabular}{rl}
     \multicolumn{2}{l}{\textbf{PERFEKT}}\\
    & sein + \textbf{gefallen}\\
  \end{tabular}
\hfill
   \begin{tabular}{rl}
   
    \multicolumn{2}{l}{\textbf{FUTURE I} }\\
    &werden + \textbf{fallen}\\
 
  \end{tabular}
\hfill
   \begin{tabular}{rl}
      \multicolumn{2}{l}{\textbf{PLUSQUAMPERFEKT}}\\
    &waren + \textbf{gefallen}\\
  \end{tabular}

   \vspace{10pt}

   \begin{tabular}{lll}
      \multicolumn{3}{l}{\textbf{MIT PRÄPOSITIONEN}}\\
& \textbf{fallen in} + Akk & (to fall into)\\
& \textbf{fallen auf} + Akk & (to fall onto)\\
&\textbf{fallen von} + Dat & (to fall from)\\
&\textbf{fallen über} + Akk & (to fall over)\\
&\textbf{fallen an} + Akk & (to accrue / arise costs or interest)\\
\vspace{10pt} 
\end{tabular}
  \hfill
   \begin{tabular}{lll}
      \multicolumn{3}{l}{\textbf{TRENNBARE VERBEN}}\\
 & \textbf{ab$|$fallen} & (to fall / drop off)\\
  & \textbf{auf$|$fallen} & (to stand out)\\
   & \textbf{aus$|$fallen} & (to be cancelled / to fail)\\
   & \textbf{ein$|$fallen} & (to occur to someone)\\
    & \textbf{um$|$fallen} & (to fall over)\\
       & \textbf{zurück$|$fallen} & (to fall behind)\\
  \end{tabular}


  \vspace{-5pt}


\end{verbbox}

%fangen (A2) – to catch / capture
\begin{verbbox-acc}\addcontentsline{toc}{subsubsection}{fangen (to catch / capture)}
    {\Large \bf fangen} (to catch / capture) \hfill Type: \textbf{irregular, + Akkusativ} \hfill Auxiliary Verb: \textbf{haben} 
    
     \vspace{5pt}

    \hrule
    
    \vspace{5pt}

  \begin{tabular}{rl}
     \multicolumn{2}{l}{\textbf{PRÄSENS} }\\
    ich & \textbf{fange}\\
    du & \textbf{fängst}\\
    er/sie/es & \textbf{fängt}\\
    wir & \textbf{fangen}\\
    ihr & \textbf{fangt}\\
    sie/Sie & \textbf{fangen}\\
  \end{tabular}
  \hfill
     \begin{tabular}{rl}
    \multicolumn{2}{l}{\textbf{PRÄTERITUM} }\\
    ich & \textbf{fing}\\
    du & \textbf{fingst}\\
    er/sie/es & \textbf{fing}\\
    wir & \textbf{fingen}\\
    ihr & \textbf{fingt}\\
    sie/Sie & \textbf{fingen}\\
  \end{tabular}
  \hfill
  \begin{tabular}{rl}
     \multicolumn{2}{l}{\textbf{IMPERATIV} }\\
   & \textbf{fang(e) (du)!}\\
   & \textbf{fangen wir!}\\
   & \textbf{fangt (ihr)!}\\
   & \textbf{fangen Sie!}\\
   & \vspace{10pt}
  \end{tabular}

    \vspace{10pt}

 \begin{tabular}{rl}
     \multicolumn{2}{l}{\textbf{PERFEKT}}\\
  &  haben + \textbf{gefangen}\\
  \end{tabular}
\hfill
   \begin{tabular}{rl}
   
    \multicolumn{2}{l}{\textbf{FUTURE I} }\\
  &  werden + \textbf{fangen}\\
 
  \end{tabular}
\hfill
   \begin{tabular}{rl}
      \multicolumn{2}{l}{\textbf{PLUSQUAMPERFEKT}}\\
  &  hatten + \textbf{gefangen}\\
  \end{tabular}
  
  \vspace{10pt}

   \begin{tabular}{lll}
      \multicolumn{3}{l}{\textbf{MIT PRÄPOSITIONEN}}\\
      & \textbf{fangen nach} + Dat & (to reach for / try to catch)\\
&\textbf{anfangen mit} + Dat & (to begin with)\\
&\textbf{anfangen bei} + Dat & (to begin at (point / person))\\
& \textbf{einfangen in} + Akk & (to capture in)\\
& \textbf{abfangen an} + Dat & (to intercept at)
  \end{tabular}
  \hfill
   \begin{tabular}{lll}
      \multicolumn{3}{l}{\textbf{TRENNBARE VERBEN}}\\
 & \textbf{an$|$fangen} & (to begin)\\
 & \textbf{auf$|$fangen} & (to catch / absorb / cushion)\\
 & \textbf{ein$|$fangen} & (to catch / capture)\\
  & \textbf{ab$|$fangen} & (to intercept)\\
  & \textbf{weg$|$fangen} & (to snatch away)\\
  & \textbf{über$|$fangen} & (to cover / overtake)
  \end{tabular}


  \vspace{-5pt}

\end{verbbox-acc}


%faszinieren (to fascinate)
\begin{verbbox-acc}\addcontentsline{toc}{subsubsection}{faszinieren (to fascinate)}
{\Large \bf faszinieren} (to fascinate) \hfill Type: \textbf{regular, + Akkusativ} \hfill Auxiliary Verb: \textbf{haben}

\vspace{5pt}
\hrule
\vspace{5pt}

\begin{tabular}{rl}
\multicolumn{2}{l}{\textbf{PRÄSENS}}\\
ich & \textbf{fasziniere}\\
du & \textbf{faszinierst}\\
er/sie/es & \textbf{fasziniert}\\
wir & \textbf{faszinieren}\\
ihr & \textbf{fasziniert}\\
sie/Sie & \textbf{faszinieren}\\
\end{tabular}
\hfill
\begin{tabular}{rl}
\multicolumn{2}{l}{\textbf{PRÄTERITUM}}\\
ich & \textbf{faszinierte}\\
du & \textbf{fasziniertest}\\
er/sie/es & \textbf{faszinierte}\\
wir & \textbf{faszinierten}\\
ihr & \textbf{fasziniertet}\\
sie/Sie & \textbf{faszinierten}\\
\end{tabular}
\hfill
\begin{tabular}{rl}
\multicolumn{2}{l}{\textbf{IMPERATIV}}\\
& \textbf{fasziniere (du)!}\\
& \textbf{faszinieren wir!}\\
& \textbf{fasziniert (ihr)!}\\
& \textbf{faszinieren Sie!}\\
\end{tabular}

\vspace{10pt}

\begin{tabular}{rl}
\multicolumn{2}{l}{\textbf{PERFEKT}}\\
& haben + \textbf{fasziniert}\\
\end{tabular}
\hfill
\begin{tabular}{rl}
\multicolumn{2}{l}{\textbf{FUTURE I}}\\
& werden + \textbf{faszinieren}\\
\end{tabular}
\hfill
\begin{tabular}{rl}
\multicolumn{2}{l}{\textbf{PLUSQUAMPERFEKT}}\\
& hatten + \textbf{fasziniert}\\
\end{tabular}

\vspace{10pt}

\begin{tabular}{lll}
\multicolumn{3}{l}{\textbf{MIT PRÄPOSITIONEN}}\\
& ---\\
\end{tabular}
\hfill
\begin{tabular}{lll}
\multicolumn{3}{l}{\textbf{TRENNBARE VERBEN}}\\
& ---\\
\end{tabular}

\vspace{-5pt}
\end{verbbox-acc}

%fehlen (A2) – to be missing
\begin{verbbox}\addcontentsline{toc}{subsubsection}{fehlen (to be missing)}
    {\Large \bf fehlen} (to be missing) \hfill Type: \textbf{regular} \hfill Auxiliary Verb: \textbf{haben} 
    
     \vspace{5pt}

    \hrule
    
    \vspace{5pt}

  \begin{tabular}{rl}
     \multicolumn{2}{l}{\textbf{PRÄSENS} }\\
    ich & \textbf{fehle}\\
    du & \textbf{fehlst}\\
    er/sie/es & \textbf{fehlt}\\
    wir & \textbf{fehlen}\\
    ihr & \textbf{fehlt}\\
    sie/Sie & \textbf{fehlen}\\
  \end{tabular}
  \hfill
     \begin{tabular}{rl}
    \multicolumn{2}{l}{\textbf{PRÄTERITUM} }\\
    ich & \textbf{fehlte}\\
    du & \textbf{fehltest}\\
    er/sie/es & \textbf{fehlte}\\
    wir & \textbf{fehlten}\\
    ihr & \textbf{fehltet}\\
    sie/Sie & \textbf{fehlten}\\
  \end{tabular}
  \hfill
  \begin{tabular}{rl}
     \multicolumn{2}{l}{\textbf{IMPERATIV} }\\
   & \textbf{fehl(e) (du)!}\\
   & \textbf{fehlen wir!}\\
   & \textbf{fehlt (ihr)!}\\
   & \textbf{fehlen Sie!}\\
   & \vspace{10pt}
  \end{tabular}

    \vspace{10pt}

 \begin{tabular}{rl}
     \multicolumn{2}{l}{\textbf{PERFEKT}}\\
    &haben + \textbf{gefehlt}\\
  \end{tabular}
\hfill
   \begin{tabular}{rl}
   
    \multicolumn{2}{l}{\textbf{FUTURE I} }\\
   & werden + \textbf{fehlen}\\
 
  \end{tabular}
\hfill
   \begin{tabular}{rl}
      \multicolumn{2}{l}{\textbf{PLUSQUAMPERFEKT}}\\
    &hatten + \textbf{gefehlt}\\
  \end{tabular}
  
  \vspace{10pt}

\begin{tabular}{lll}
\multicolumn{3}{l}{\textbf{MIT PRÄPOSITIONEN}}\\
& ---\\
\end{tabular}
\hfill
\begin{tabular}{lll}
\multicolumn{3}{l}{\textbf{TRENNBARE VERBEN}}\\
& ---\\
\end{tabular}

  \vspace{-5pt}

\end{verbbox}
%feiern (A2) – to celebrate
\begin{verbbox-acc}\addcontentsline{toc}{subsubsection}{feiern (to celebrate)}
    {\Large \bf feiern} (to celebrate) \hfill Type: \textbf{regular, + Akkusativ} \hfill Auxiliary Verb: \textbf{haben} 
    
     \vspace{5pt}

    \hrule
    
    \vspace{5pt}

  \begin{tabular}{rl}
     \multicolumn{2}{l}{\textbf{PRÄSENS} }\\
    ich & \textbf{feiere}\\
    du & \textbf{feierst}\\
    er/sie/es & \textbf{feiert}\\
    wir & \textbf{feiern}\\
    ihr & \textbf{feiert}\\
    sie/Sie & \textbf{feiern}\\
  \end{tabular}
  \hfill
     \begin{tabular}{rl}
    \multicolumn{2}{l}{\textbf{PRÄTERITUM} }\\
    ich & \textbf{feierte}\\
    du & \textbf{feiertest}\\
    er/sie/es & \textbf{feierte}\\
    wir & \textbf{feierten}\\
    ihr & \textbf{feiertet}\\
    sie/Sie & \textbf{feierten}\\
  \end{tabular}
  \hfill
  \begin{tabular}{rl}
     \multicolumn{2}{l}{\textbf{IMPERATIV} }\\
   & \textbf{feier(e) (du)!}\\
   & \textbf{feiern wir!}\\
   & \textbf{feiert (ihr)!}\\
   & \textbf{feiern Sie!}\\
   & \vspace{10pt}
  \end{tabular}

    \vspace{10pt}

 \begin{tabular}{rl}
     \multicolumn{2}{l}{\textbf{PERFEKT}}\\
 &   haben + \textbf{gefeiert}\\
  \end{tabular}
\hfill
   \begin{tabular}{rl}
   
    \multicolumn{2}{l}{\textbf{FUTURE I} }\\
  &  werden + \textbf{feiern}\\
 
  \end{tabular}
\hfill
   \begin{tabular}{rl}
      \multicolumn{2}{l}{\textbf{PLUSQUAMPERFEKT}}\\
  &  hatten + \textbf{gefeiert}\\
  \end{tabular}
  
  \vspace{-5pt}


\end{verbbox-acc}

%finanzieren (to finance / fund)
\begin{verbbox}\addcontentsline{toc}{subsubsection}{finanzieren (to finance / fund)}
{\Large \bf finanzieren} (to finance / fund) \hfill Type: \textbf{regular} \hfill Auxiliary Verb: \textbf{haben}

\vspace{5pt}
\hrule
\vspace{5pt}

\begin{tabular}{rl}
\multicolumn{2}{l}{\textbf{PRÄSENS}}\\
ich & \textbf{finanziere}\\
du & \textbf{finanzierst}\\
er/sie/es & \textbf{finanziert}\\
wir & \textbf{finanzieren}\\
ihr & \textbf{finanziert}\\
sie/Sie & \textbf{finanzieren}\\
\end{tabular}
\hfill
\begin{tabular}{rl}
\multicolumn{2}{l}{\textbf{PRÄTERITUM}}\\
ich & \textbf{finanzierte}\\
du & \textbf{finanziertest}\\
er/sie/es & \textbf{finanzierte}\\
wir & \textbf{finanzierten}\\
ihr & \textbf{finanziertet}\\
sie/Sie & \textbf{finanzierten}\\
\end{tabular}
\hfill
\begin{tabular}{rl}
\multicolumn{2}{l}{\textbf{IMPERATIV}}\\
& \textbf{finanziere (du)!}\\
& \textbf{finanzieren wir!}\\
& \textbf{finanziert (ihr)!}\\
& \textbf{finanzieren Sie!}\\
\end{tabular}

\vspace{10pt}

\begin{tabular}{rl}
\multicolumn{2}{l}{\textbf{PERFEKT}}\\
& haben + \textbf{finanziert}\\
\end{tabular}
\hfill
\begin{tabular}{rl}
\multicolumn{2}{l}{\textbf{FUTURE I}}\\
& werden + \textbf{finanzieren}\\
\end{tabular}
\hfill
\begin{tabular}{rl}
\multicolumn{2}{l}{\textbf{PLUSQUAMPERFEKT}}\\
& hatten + \textbf{finanziert}\\
\end{tabular}

\vspace{10pt}

\begin{tabular}{lll}
\multicolumn{3}{l}{\textbf{MIT PRÄPOSITIONEN}}\\
& \textbf{finanzieren mit} + Dat & (to finance with)\\
& \textbf{finanzieren durch} + Akk & (to finance through / by)\\
\end{tabular}
\hfill
\begin{tabular}{lll}
\multicolumn{3}{l}{\textbf{TRENNBARE VERBEN}}\\
& --- & (no separable forms)\\
\end{tabular}
\vspace{-5pt}
\end{verbbox}

%finden (A1) – to find
\begin{verbbox-acc}\addcontentsline{toc}{subsubsection}{finden (to find)}
    {\Large \bf finden} (to find) \hfill Type: \textbf{irregular, + Akkusativ} \hfill Auxiliary Verb: \textbf{haben} 
    
     \vspace{5pt}

    \hrule
    
    \vspace{5pt}

  \begin{tabular}{rl}
     \multicolumn{2}{l}{\textbf{PRÄSENS} }\\
    ich & \textbf{finde}\\
    du & \textbf{findest}\\
    er/sie/es & \textbf{findet}\\
    wir & \textbf{finden}\\
    ihr & \textbf{findet}\\
    sie/Sie & \textbf{finden}\\
  \end{tabular}
  \hfill
     \begin{tabular}{rl}
    \multicolumn{2}{l}{\textbf{PRÄTERITUM} }\\
    ich & \textbf{fand}\\
    du & \textbf{fandest}\\
    er/sie/es & \textbf{fand}\\
    wir & \textbf{fanden}\\
    ihr & \textbf{fandet}\\
    sie/Sie & \textbf{fanden}\\
  \end{tabular}
  \hfill
  \begin{tabular}{rl}
     \multicolumn{2}{l}{\textbf{IMPERATIV} }\\
   & \textbf{find(e) (du)!}\\
   & \textbf{finden wir!}\\
   & \textbf{findet (ihr)!}\\
   & \textbf{finden Sie!}\\
   & \vspace{10pt}
  \end{tabular}

    \vspace{10pt}

 \begin{tabular}{rl}
     \multicolumn{2}{l}{\textbf{PERFEKT}}\\
  &  haben + \textbf{gefunden}\\
  \end{tabular}
\hfill
   \begin{tabular}{rl}
   
    \multicolumn{2}{l}{\textbf{FUTURE I} }\\
  &  werden + \textbf{finden}\\
 
  \end{tabular}
\hfill
   \begin{tabular}{rl}
      \multicolumn{2}{l}{\textbf{PLUSQUAMPERFEKT}}\\
   & hatten + \textbf{gefunden}\\
  \end{tabular}
  
   \vspace{10pt}

\begin{tabular}{lll}
      \multicolumn{3}{l}{\textbf{MIT PRÄPOSITIONEN}}\\
& \textbf{finden für} + Akk & (to find suitable for)\\
& \textbf{finden zu} + Dat & (to find to / tend to)\\
& \textbf{finden in} + Dat & (to find in / inside)\\
\end{tabular}
  \hfill
\begin{tabular}{lll}
      \multicolumn{3}{l}{\textbf{TRENNBARE VERBEN}}\\
 & \textbf{wieder$|$finden} & (to recover / find again)\\
 & \textbf{statt$|$finden} & (to take place / happen)\\
\vspace{10pt}
  \end{tabular}

  \vspace{-5pt}

\end{verbbox-acc}

%fliegen - to fly
\begin{verbbox}\addcontentsline{toc}{subsubsection}{fliegen (to fly)}
    {\Large \bf fliegen} (to fly) \hfill Type: \textbf{irregular} \hfill Auxiliary Verb: \textbf{sein} 
    
     \vspace{5pt}

    \hrule
    
    \vspace{5pt}

  \begin{tabular}{rl}
     \multicolumn{2}{l}{\textbf{PRÄSENS} }\\
    ich & \textbf{fliege}\\
    du & \textbf{fliegst}\\
    er/sie/es & \textbf{fliegt}\\
    wir & \textbf{fliegen}\\
    ihr & \textbf{fliegt}\\
    sie/Sie & \textbf{fliegen}\\
  \end{tabular}
  \hfill
     \begin{tabular}{rl}
    \multicolumn{2}{l}{\textbf{PRÄTERITUM} }\\
    ich & \textbf{flog}\\
    du & \textbf{flogst}\\
    er/sie/es & \textbf{flog}\\
    wir & \textbf{flogen}\\
    ihr & \textbf{flogt}\\
    sie/Sie & \textbf{flogen}\\
  \end{tabular}
  \hfill
  \begin{tabular}{rl}
     \multicolumn{2}{l}{\textbf{IMPERATIV} }\\
   & \textbf{flieg(e) (du)!}\\
   & \textbf{fliegen wir!}\\
   & \textbf{fliegt (ihr)!}\\
   & \textbf{fliegen Sie!}\\
   & \vspace{10pt}
  \end{tabular}

    \vspace{10pt}

 \begin{tabular}{rl}
     \multicolumn{2}{l}{\textbf{PERFEKT}}\\
   & sein + \textbf{geflogen}\\
  \end{tabular}
\hfill
   \begin{tabular}{rl}
    \multicolumn{2}{l}{\textbf{FUTURE I} }\\
   & werden + \textbf{fliegen}\\
  \end{tabular}
\hfill
   \begin{tabular}{rl}
      \multicolumn{2}{l}{\textbf{PLUSQUAMPERFEKT}}\\
   & waren + \textbf{geflogen}\\
  \end{tabular}

   \vspace{10pt}

   \begin{tabular}{lll}
      \multicolumn{3}{l}{\textbf{MIT PRÄPOSITIONEN}}\\
& \textbf{fliegen nach} + Dat & (to fly to – cities/countries)\\
& \textbf{fliegen in} + Akk & (to fly into)\\
& \textbf{fliegen über} + Akk & (to fly over)\\
& \textbf{fliegen von} + Dat & (to fly from)\\
   \end{tabular}
  \hfill
   \begin{tabular}{lll}
      \multicolumn{3}{l}{\textbf{TRENNBARE VERBEN}}\\
 & \textbf{ab$|$fliegen} & (to depart by plane)\\
 & \textbf{an$|$fliegen} & (to approach by flying)\\
 & \textbf{ein$|$fliegen} & (to fly in)\\
 & \textbf{weg$|$fliegen} & (to fly away)\\
 & \textbf{über$|$fliegen} & (to fly over)\\
  \end{tabular}

\vspace{-5pt}
\end{verbbox}

%fliehen (to flee)
\begin{verbbox}\addcontentsline{toc}{subsubsection}{fliehen (to flee)}
    {\Large \bf fliehen} (to flee) \hfill Type: \textbf{irregular} \hfill Auxiliary Verb: \textbf{sein} 
    
     \vspace{5pt}

    \hrule
    
    \vspace{5pt}

  \begin{tabular}{rl}
     \multicolumn{2}{l}{\textbf{PRÄSENS} }\\
    ich & \textbf{fliehe}\\
    du & \textbf{fliehst}\\
    er/sie/es & \textbf{flieht}\\
    wir & \textbf{fliehen}\\
    ihr & \textbf{flieht}\\
    sie/Sie & \textbf{fliehen}\\
  \end{tabular}
  \hfill
     \begin{tabular}{rl}
    \multicolumn{2}{l}{\textbf{PRÄTERITUM} }\\
    ich & \textbf{floh}\\
    du & \textbf{flohst}\\
    er/sie/es & \textbf{floh}\\
    wir & \textbf{flohen}\\
    ihr & \textbf{floht}\\
    sie/Sie & \textbf{flohen}\\
  \end{tabular}
  \hfill
  \begin{tabular}{rl}
     \multicolumn{2}{l}{\textbf{IMPERATIV} }\\
   & \textbf{flieh(e) (du)!}\\
   & \textbf{fliehen wir!}\\
   & \textbf{flieht (ihr)!}\\
   & \textbf{fliehen Sie!}\\
   & \vspace{10pt}
  \end{tabular}

    \vspace{10pt}

 \begin{tabular}{rl}
     \multicolumn{2}{l}{\textbf{PERFEKT}}\\
   & sein + \textbf{geflohen}\\
  \end{tabular}
\hfill
   \begin{tabular}{rl}
    \multicolumn{2}{l}{\textbf{FUTURE I} }\\
   & werden + \textbf{fliehen}\\
  \end{tabular}
\hfill
   \begin{tabular}{rl}
      \multicolumn{2}{l}{\textbf{PLUSQUAMPERFEKT}}\\
   & waren + \textbf{geflohen}\\
  \end{tabular}

   \vspace{10pt}

   \begin{tabular}{lll}
      \multicolumn{3}{l}{\textbf{MIT PRÄPOSITIONEN}}\\
& \textbf{fliehen vor} + Dat & (to flee from)\\
& \textbf{fliehen aus} + Dat & (to flee out of / from)\\
\end{tabular}
  \hfill
   \begin{tabular}{lll}
      \multicolumn{3}{l}{\textbf{TRENNBARE VERBEN}}\\
 & \textbf{ent$|$fliehen} & (to escape from)\\
\vspace{10pt}
  \end{tabular}

  \vspace{-5pt}

\end{verbbox}

%folgen (to follow)
\begin{verbbox-dat}\addcontentsline{toc}{subsubsection}{folgen (to follow)}
    {\Large \bf folgen} (to follow) \hfill Type: \textbf{regular} \hfill Auxiliary Verb: \textbf{sein} 
    
     \vspace{5pt}

    \hrule
    
    \vspace{5pt}

  \begin{tabular}{rl}
     \multicolumn{2}{l}{\textbf{PRÄSENS} }\\
    ich & \textbf{folge}\\
    du & \textbf{folgst}\\
    er/sie/es & \textbf{folgt}\\
    wir & \textbf{folgen}\\
    ihr & \textbf{folgt}\\
    sie/Sie & \textbf{folgen}\\
  \end{tabular}
  \hfill
     \begin{tabular}{rl}
    \multicolumn{2}{l}{\textbf{PRÄTERITUM} }\\
    ich & \textbf{folgte}\\
    du & \textbf{folgtest}\\
    er/sie/es & \textbf{folgte}\\
    wir & \textbf{folgten}\\
    ihr & \textbf{folgtet}\\
    sie/Sie & \textbf{folgten}\\
  \end{tabular}
  \hfill
  \begin{tabular}{rl}
     \multicolumn{2}{l}{\textbf{IMPERATIV} }\\
   & \textbf{folg(e) (du)!}\\
   & \textbf{folgen wir!}\\
   & \textbf{folgt (ihr)!}\\
   & \textbf{folgen Sie!}\\
   & \vspace{10pt}
  \end{tabular}

    \vspace{10pt}

 \begin{tabular}{rl}
     \multicolumn{2}{l}{\textbf{PERFEKT}}\\
   & sein + \textbf{gefolgt}\\
  \end{tabular}
\hfill
   \begin{tabular}{rl}
   
    \multicolumn{2}{l}{\textbf{FUTURE I} }\\
   & werden + \textbf{folgen}\\
 
  \end{tabular}
\hfill
   \begin{tabular}{rl}
      \multicolumn{2}{l}{\textbf{PLUSQUAMPERFEKT}}\\
   & waren + \textbf{gefolgt}\\
  \end{tabular}

   \vspace{10pt}

   \begin{tabular}{lll}
      \multicolumn{3}{l}{\textbf{MIT PRÄPOSITIONEN}}\\
& \textbf{folgen auf} + Akk & (to follow after)\\
\vspace{10pt}
\end{tabular}
  \hfill
   \begin{tabular}{lll}
      \multicolumn{3}{l}{\textbf{TRENNBARE VERBEN}}\\
 & \textbf{nach$|$folgen} & (to succeed / follow after)\\
 & \textbf{ver$|$folgen} & (not separable: to pursue)\\
  \end{tabular}


  \vspace{-5pt}


\end{verbbox-dat}

%fragen (A1) – to ask
\begin{verbbox-acc}\addcontentsline{toc}{subsubsection}{fragen (to ask)}
    {\Large \bf fragen} (to ask) \hfill Type: \textbf{regular, + Akkusativ} \hfill Auxiliary Verb: \textbf{haben} 
    
     \vspace{5pt}

    \hrule
    
    \vspace{5pt}

  \begin{tabular}{rl}
     \multicolumn{2}{l}{\textbf{PRÄSENS} }\\
    ich & \textbf{frage}\\
    du & \textbf{fragst}\\
    er/sie/es & \textbf{fragt}\\
    wir & \textbf{fragen}\\
    ihr & \textbf{fragt}\\
    sie/Sie & \textbf{fragen}\\
  \end{tabular}
  \hfill
     \begin{tabular}{rl}
    \multicolumn{2}{l}{\textbf{PRÄTERITUM} }\\
    ich & \textbf{fragte}\\
    du & \textbf{fragtest}\\
    er/sie/es & \textbf{fragte}\\
    wir & \textbf{fragten}\\
    ihr & \textbf{fragtet}\\
    sie/Sie & \textbf{fragten}\\
  \end{tabular}
  \hfill
  \begin{tabular}{rl}
     \multicolumn{2}{l}{\textbf{IMPERATIV} }\\
   & \textbf{frag(e) (du)!}\\
   & \textbf{fragen wir!}\\
   & \textbf{fragt (ihr)!}\\
   & \textbf{fragen Sie!}\\
   & \vspace{10pt}
  \end{tabular}

    \vspace{10pt}

 \begin{tabular}{rl}
     \multicolumn{2}{l}{\textbf{PERFEKT}}\\
   & haben + \textbf{gefragt}\\
  \end{tabular}
\hfill
   \begin{tabular}{rl}
   
    \multicolumn{2}{l}{\textbf{FUTURE I} }\\
   & werden + \textbf{fragen}\\
 
  \end{tabular}
\hfill
   \begin{tabular}{rl}
      \multicolumn{2}{l}{\textbf{PLUSQUAMPERFEKT}}\\
   & hatten + \textbf{gefragt}\\
  \end{tabular}
  
  \vspace{-5pt}


\end{verbbox-acc}

%fressen (to eat — animals)
\begin{verbbox}\addcontentsline{toc}{subsubsection}{fressen (to eat — animals)}
    {\Large \bf fressen} (to eat — animals) \hfill Type: \textbf{irregular} \hfill Auxiliary Verb: \textbf{haben} 
    
     \vspace{5pt}

    \hrule
    
    \vspace{5pt}

  \begin{tabular}{rl}
     \multicolumn{2}{l}{\textbf{PRÄSENS} }\\
    ich & \textbf{fresse}\\
    du & \textbf{frisst}\\
    er/sie/es & \textbf{frisst}\\
    wir & \textbf{fressen}\\
    ihr & \textbf{fresst}\\
    sie/Sie & \textbf{fressen}\\
  \end{tabular}
  \hfill
     \begin{tabular}{rl}
    \multicolumn{2}{l}{\textbf{PRÄTERITUM} }\\
    ich & \textbf{fraß}\\
    du & \textbf{fraßest}\\
    er/sie/es & \textbf{fraß}\\
    wir & \textbf{fraßen}\\
    ihr & \textbf{fraßt}\\
    sie/Sie & \textbf{fraßen}\\
  \end{tabular}
  \hfill
  \begin{tabular}{rl}
     \multicolumn{2}{l}{\textbf{IMPERATIV} }\\
   & \textbf{friss (du)!}\\
   & \textbf{fressen wir!}\\
   & \textbf{fresst (ihr)!}\\
   & \textbf{fressen Sie!}\\
   & \vspace{10pt}
  \end{tabular}

    \vspace{10pt}

 \begin{tabular}{rl}
     \multicolumn{2}{l}{\textbf{PERFEKT}}\\
   & haben + \textbf{gefressen}\\
  \end{tabular}
\hfill
   \begin{tabular}{rl}
    \multicolumn{2}{l}{\textbf{FUTURE I} }\\
   & werde + \textbf{fressen}\\
  \end{tabular}
\hfill
   \begin{tabular}{rl}
      \multicolumn{2}{l}{\textbf{PLUSQUAMPERFEKT}}\\
   & hatten + \textbf{gefressen}\\
  \end{tabular}

   \vspace{10pt}

   \begin{tabular}{lll}
      \multicolumn{3}{l}{\textbf{MIT PRÄPOSITIONEN}}\\
& \textbf{fressen von} + Dat & (to eat from / of)\\
& \textbf{fressen aus} + Dat & (to eat from)\\
\end{tabular}
  \hfill
   \begin{tabular}{lll}
      \multicolumn{3}{l}{\textbf{TRENNBARE VERBEN}}\\
 & \textbf{auf$|$fressen} & (to eat up)\\
 & \textbf{an$|$fressen} & (to nibble at / eat into)\\
\vspace{10pt}
  \end{tabular}

  \vspace{-5pt}

\end{verbbox}

%fügen (to add / insert)
\begin{verbbox}\addcontentsline{toc}{subsubsection}{fügen (to add / insert)}
{\Large \bf fügen} (to add / insert) \hfill Type: \textbf{regular} \hfill Auxiliary Verb: \textbf{haben}

\vspace{5pt}
\hrule
\vspace{5pt}

\begin{tabular}{rl}
\multicolumn{2}{l}{\textbf{PRÄSENS}}\\
ich & \textbf{füge}\\
du & \textbf{fügst}\\
er/sie/es & \textbf{fügt}\\
wir & \textbf{fügen}\\
ihr & \textbf{fügt}\\
sie/Sie & \textbf{fügen}\\
\end{tabular}
\hfill
\begin{tabular}{rl}
\multicolumn{2}{l}{\textbf{PRÄTERITUM}}\\
ich & \textbf{fügte}\\
du & \textbf{fügtest}\\
er/sie/es & \textbf{fügte}\\
wir & \textbf{fügten}\\
ihr & \textbf{fügtet}\\
sie/Sie & \textbf{fügten}\\
\end{tabular}
\hfill
\begin{tabular}{rl}
\multicolumn{2}{l}{\textbf{IMPERATIV}}\\
& \textbf{füg(e) (du)!}\\
& \textbf{fügen wir!}\\
& \textbf{fügt (ihr)!}\\
& \textbf{fügen Sie!}\\
\end{tabular}

\vspace{10pt}

\begin{tabular}{rl}
\multicolumn{2}{l}{\textbf{PERFEKT}}\\
& haben + \textbf{gefügt}\\
\end{tabular}
\hfill
\begin{tabular}{rl}
\multicolumn{2}{l}{\textbf{FUTURE I}}\\
& werden + \textbf{fügen}\\
\end{tabular}
\hfill
\begin{tabular}{rl}
\multicolumn{2}{l}{\textbf{PLUSQUAMPERFEKT}}\\
& hatten + \textbf{gefügt}\\
\end{tabular}

\vspace{10pt}

\begin{tabular}{lll}
\multicolumn{3}{l}{\textbf{MIT PRÄPOSITIONEN}}\\
& \textbf{sich fügen in} + Akk & (to submit to)\\
\end{tabular}
\hfill
\begin{tabular}{lll}
\multicolumn{3}{l}{\textbf{TRENNBARE VERBEN}}\\
& \textbf{ein$|$fügen} & (to insert)\\
& \textbf{hinzu$|$fügen} & (to add)\\
\end{tabular}

\vspace{-5pt}
\end{verbbox}

%fuhren - to guide
\begin{verbbox}\addcontentsline{toc}{subsubsection}{führen (to lead / guide / drive)}
    {\Large \bf führen} (to lead / guide / drive) \hfill Type: \textbf{regular} \hfill Auxiliary Verb: \textbf{haben} 
    
     \vspace{5pt}

    \hrule
    
    \vspace{5pt}

  \begin{tabular}{rl}
     \multicolumn{2}{l}{\textbf{PRÄSENS} }\\
    ich & \textbf{führe}\\
    du & \textbf{führst}\\
    er/sie/es & \textbf{führt}\\
    wir & \textbf{führen}\\
    ihr & \textbf{führt}\\
    sie/Sie & \textbf{führen}\\
  \end{tabular}
  \hfill
  \begin{tabular}{rl}
    \multicolumn{2}{l}{\textbf{PRÄTERITUM} }\\
    ich & \textbf{führte}\\
    du & \textbf{führtest}\\
    er/sie/es & \textbf{führte}\\
    wir & \textbf{führten}\\
    ihr & \textbf{führtet}\\
    sie/Sie & \textbf{führten}\\
  \end{tabular}
  \hfill
  \begin{tabular}{rl}
     \multicolumn{2}{l}{\textbf{IMPERATIV} }\\
   & \textbf{führe (du)!}\\
   & \textbf{führen wir!}\\
   & \textbf{führt (ihr)!}\\
   & \textbf{führen Sie!}\\
   & \vspace{10pt}
  \end{tabular}

 \vspace{10pt}

 \begin{tabular}{rl}
     \multicolumn{2}{l}{\textbf{PERFEKT}}\\
   & haben + \textbf{geführt}\\
  \end{tabular}
\hfill
\begin{tabular}{rl}
    \multicolumn{2}{l}{\textbf{FUTURE I}}\\
   & werden + \textbf{führen}\\
  \end{tabular}
\hfill
\begin{tabular}{rl}
    \multicolumn{2}{l}{\textbf{PLUSQUAMPERFEKT}}\\
   & hatten + \textbf{geführt}\\
  \end{tabular}

 \vspace{10pt}

\begin{tabular}{lll}
      \multicolumn{3}{l}{\textbf{MIT PRÄPOSITIONEN}}\\
& \textbf{führen zu} + Dat & (to lead to / result in)\\
& \textbf{führen durch} + Akk & (to guide through)\\
& \textbf{führen über} + Akk & (to lead across / over)\\
\end{tabular}
  \hfill
\begin{tabular}{lll}
      \multicolumn{3}{l}{\textbf{TRENNBARE VERBEN}}\\
 & \textbf{an$|$führen} & (to lead to / bring to)\\
 & \textbf{durch$|$führen} & (to guide through)\\
\vspace{10pt}
  \end{tabular}


\vspace{-5pt}
\end{verbbox}

%füllen (to fill)
\begin{verbbox}\addcontentsline{toc}{subsubsection}{füllen (to fill)}
    {\Large \bf füllen} (to fill) \hfill Type: \textbf{regular} \hfill Auxiliary Verb: \textbf{haben} 
    
     \vspace{5pt}

    \hrule
    
    \vspace{5pt}

  \begin{tabular}{rl}
     \multicolumn{2}{l}{\textbf{PRÄSENS} }\\
    ich & \textbf{fülle}\\
    du & \textbf{füllst}\\
    er/sie/es & \textbf{füllt}\\
    wir & \textbf{füllen}\\
    ihr & \textbf{füllt}\\
    sie/Sie & \textbf{füllen}\\
  \end{tabular}
  \hfill
     \begin{tabular}{rl}
    \multicolumn{2}{l}{\textbf{PRÄTERITUM} }\\
    ich & \textbf{füllte}\\
    du & \textbf{fülltest}\\
    er/sie/es & \textbf{füllte}\\
    wir & \textbf{füllten}\\
    ihr & \textbf{fülltet}\\
    sie/Sie & \textbf{füllten}\\
  \end{tabular}
  \hfill
  \begin{tabular}{rl}
     \multicolumn{2}{l}{\textbf{IMPERATIV} }\\
   & \textbf{füll(e) (du)!}\\
   & \textbf{füllen wir!}\\
   & \textbf{füllt (ihr)!}\\
   & \textbf{füllen Sie!}\\
   & \vspace{10pt}
  \end{tabular}

    \vspace{10pt}

 \begin{tabular}{rl}
     \multicolumn{2}{l}{\textbf{PERFEKT}}\\
   & haben + \textbf{gefüllt}\\
  \end{tabular}
\hfill
   \begin{tabular}{rl}
    \multicolumn{2}{l}{\textbf{FUTURE I} }\\
   & werde + \textbf{füllen}\\
  \end{tabular}
\hfill
   \begin{tabular}{rl}
      \multicolumn{2}{l}{\textbf{PLUSQUAMPERFEKT}}\\
   & hatten + \textbf{gefüllt}\\
  \end{tabular}

   \vspace{10pt}

   \begin{tabular}{lll}
      \multicolumn{3}{l}{\textbf{MIT PRÄPOSITIONEN}}\\
& \textbf{füllen mit} + Dat & (to fill with)\\
& \textbf{füllen in} + Akk & (to fill into)\\
\end{tabular}
  \hfill
   \begin{tabular}{lll}
      \multicolumn{3}{l}{\textbf{TRENNBARE VERBEN}}\\
 & \textbf{aus$|$füllen} & (to fill out)\\
 & \textbf{auf$|$füllen} & (to refill / top up)\\
\vspace{10pt}
  \end{tabular}

  \vspace{-5pt}

\end{verbbox}

%funktionieren (to function / work)
\begin{verbbox}\addcontentsline{toc}{subsubsection}{funktionieren (to function / work)}
{\Large \bf funktionieren} (to function / work) \hfill Type: \textbf{regular} \hfill Auxiliary Verb: \textbf{haben}

\vspace{5pt}
\hrule
\vspace{5pt}

\begin{tabular}{rl}
\multicolumn{2}{l}{\textbf{PRÄSENS}}\\
ich & \textbf{funktioniere}\\
du & \textbf{funktionierst}\\
er/sie/es & \textbf{funktioniert}\\
wir & \textbf{funktionieren}\\
ihr & \textbf{funktioniert}\\
sie/Sie & \textbf{funktionieren}\\
\end{tabular}
\hfill
\begin{tabular}{rl}
\multicolumn{2}{l}{\textbf{PRÄTERITUM}}\\
ich & \textbf{funktionierte}\\
du & \textbf{funktioniertest}\\
er/sie/es & \textbf{funktionierte}\\
wir & \textbf{funktionierten}\\
ihr & \textbf{funktioniertet}\\
sie/Sie & \textbf{funktionierten}\\
\end{tabular}
\hfill
\begin{tabular}{rl}
\multicolumn{2}{l}{\textbf{IMPERATIV}}\\
\end{tabular}

\vspace{10pt}

\begin{tabular}{rl}
\multicolumn{2}{l}{\textbf{PERFEKT}}\\
& haben + \textbf{funktioniert}\\
\end{tabular}
\hfill
\begin{tabular}{rl}
\multicolumn{2}{l}{\textbf{FUTURE I}}\\
& werden + \textbf{funktionieren}\\
\end{tabular}
\hfill
\begin{tabular}{rl}
\multicolumn{2}{l}{\textbf{PLUSQUAMPERFEKT}}\\
& hatten + \textbf{funktioniert}\\
\end{tabular}

\vspace{10pt}

\begin{tabular}{lll}
\multicolumn{3}{l}{\textbf{MIT PRÄPOSITIONEN}}\\
& \textbf{funktionieren mit} + Dat & (to function with)\\
\end{tabular}
\hfill
\begin{tabular}{lll}
\multicolumn{3}{l}{\textbf{TRENNBARE VERBEN}}\\
& --- & (no separable forms)\\
\end{tabular}
\vspace{-5pt}
\end{verbbox}

%füttern (to feed)
\begin{verbbox}\addcontentsline{toc}{subsubsection}{füttern (to feed)}
    {\Large \bf füttern} (to feed) \hfill Type: \textbf{regular} \hfill Auxiliary Verb: \textbf{haben} 
    
     \vspace{5pt}

    \hrule
    
    \vspace{5pt}

  \begin{tabular}{rl}
     \multicolumn{2}{l}{\textbf{PRÄSENS} }\\
    ich & \textbf{füttere}\\
    du & \textbf{fütterst}\\
    er/sie/es & \textbf{füttert}\\
    wir & \textbf{füttern}\\
    ihr & \textbf{füttert}\\
    sie/Sie & \textbf{füttern}\\
  \end{tabular}
  \hfill
     \begin{tabular}{rl}
    \multicolumn{2}{l}{\textbf{PRÄTERITUM} }\\
    ich & \textbf{fütterte}\\
    du & \textbf{füttertest}\\
    er/sie/es & \textbf{fütterte}\\
    wir & \textbf{fütterten}\\
    ihr & \textbf{füttertet}\\
    sie/Sie & \textbf{fütterten}\\
  \end{tabular}
  \hfill
  \begin{tabular}{rl}
     \multicolumn{2}{l}{\textbf{IMPERATIV} }\\
   & \textbf{füttere (du)!}\\
   & \textbf{füttern wir!}\\
   & \textbf{füttert (ihr)!}\\
   & \textbf{füttern Sie!}\\
   & \vspace{10pt}
  \end{tabular}

    \vspace{10pt}

 \begin{tabular}{rl}
     \multicolumn{2}{l}{\textbf{PERFEKT}}\\
   & haben + \textbf{gefüttert}\\
  \end{tabular}
\hfill
   \begin{tabular}{rl}
    \multicolumn{2}{l}{\textbf{FUTURE I} }\\
   & werde + \textbf{füttern}\\
  \end{tabular}
\hfill
   \begin{tabular}{rl}
      \multicolumn{2}{l}{\textbf{PLUSQUAMPERFEKT}}\\
   & hatten + \textbf{gefüttert}\\
  \end{tabular}

   \vspace{10pt}

   \begin{tabular}{lll}
      \multicolumn{3}{l}{\textbf{MIT PRÄPOSITIONEN}}\\
& \textbf{füttern mit} + Dat & (to feed with)\\
& \textbf{füttern an} + Akk & (to feed to)\\
\end{tabular}
  \hfill
   \begin{tabular}{lll}
      \multicolumn{3}{l}{\textbf{TRENNBARE VERBEN}}\\
 & \textbf{an$|$füttern} & (to start feeding / bait)\\
 & \textbf{zu$|$füttern} & (to give supplementary food)\\
\vspace{10pt}
  \end{tabular}

  \vspace{-5pt}

\end{verbbox}

\subsection{G Verbs}

%geben (A1) – to give
\begin{verbbox-dat}\addcontentsline{toc}{subsubsection}{geben (to give)}
    {\Large \bf geben} (to give) \hfill Type: \textbf{irregular, + Dativ} \hfill Auxiliary Verb: \textbf{haben} 
    
     \vspace{5pt}

    \hrule
    
    \vspace{5pt}

  \begin{tabular}{rl}
     \multicolumn{2}{l}{\textbf{PRÄSENS} }\\
    ich & \textbf{gebe}\\
    du & \textbf{gibst}\\
    er/sie/es & \textbf{gibt}\\
    wir & \textbf{geben}\\
    ihr & \textbf{gebt}\\
    sie/Sie & \textbf{geben}\\
  \end{tabular}
  \hfill
     \begin{tabular}{rl}
    \multicolumn{2}{l}{\textbf{PRÄTERITUM} }\\
    ich & \textbf{gab}\\
    du & \textbf{gabst}\\
    er/sie/es & \textbf{gab}\\
    wir & \textbf{gaben}\\
    ihr & \textbf{gabt}\\
    sie/Sie & \textbf{gaben}\\
  \end{tabular}
  \hfill
  \begin{tabular}{rl}
     \multicolumn{2}{l}{\textbf{IMPERATIV} }\\
   & \textbf{gib (du)!}\\
   & \textbf{geben wir!}\\
   & \textbf{gebt (ihr)!}\\
   & \textbf{geben Sie!}\\
   & \vspace{10pt}
  \end{tabular}

    \vspace{10pt}

 \begin{tabular}{rl}
     \multicolumn{2}{l}{\textbf{PERFEKT}}\\
   & haben + \textbf{gegeben}\\
  \end{tabular}
\hfill
   \begin{tabular}{rl}
   
    \multicolumn{2}{l}{\textbf{FUTURE I} }\\
  &  werden + \textbf{geben}\\
 
  \end{tabular}
\hfill
   \begin{tabular}{rl}
      \multicolumn{2}{l}{\textbf{PLUSQUAMPERFEKT}}\\
  &  hatten + \textbf{gegeben}\\
  \end{tabular}
  
  \vspace{-5pt}


\end{verbbox-dat}
%gefallen (A2) – to like
\begin{verbbox-dat}\addcontentsline{toc}{subsubsection}{gefallen (to like)}
    {\Large \bf gefallen} (to like) \hfill Type: \textbf{irregular, + Dativ} \hfill Auxiliary Verb: \textbf{haben} 
    
     \vspace{5pt}

    \hrule
    
    \vspace{5pt}

  \begin{tabular}{rl}
     \multicolumn{2}{l}{\textbf{PRÄSENS} }\\
    ich & \textbf{gefalle}\\
    du & \textbf{gefällst}\\
    er/sie/es & \textbf{gefällt}\\
    wir & \textbf{gefallen}\\
    ihr & \textbf{gefallt}\\
    sie/Sie & \textbf{gefallen}\\
  \end{tabular}
  \hfill
     \begin{tabular}{rl}
    \multicolumn{2}{l}{\textbf{PRÄTERITUM} }\\
    ich & \textbf{gefiel}\\
    du & \textbf{gefielst}\\
    er/sie/es & \textbf{gefiel}\\
    wir & \textbf{gefielen}\\
    ihr & \textbf{gefielt}\\
    sie/Sie & \textbf{gefielen}\\
  \end{tabular}
  \hfill
  \begin{tabular}{rl}
     \multicolumn{2}{l}{\textbf{IMPERATIV} }\\
   & \textbf{gefall(e) (du)!}\\
   & \textbf{gefallen wir!}\\
   & \textbf{gefallt (ihr)!}\\
   & \textbf{gefallen Sie!}\\
   & \vspace{10pt}
  \end{tabular}

    \vspace{10pt}

 \begin{tabular}{rl}
     \multicolumn{2}{l}{\textbf{PERFEKT}}\\
  &  haben + \textbf{gefallen}\\
  \end{tabular}
\hfill
   \begin{tabular}{rl}
   
    \multicolumn{2}{l}{\textbf{FUTURE I} }\\
  &  werden + \textbf{gefallen}\\
 
  \end{tabular}
\hfill
   \begin{tabular}{rl}
      \multicolumn{2}{l}{\textbf{PLUSQUAMPERFEKT}}\\
   & hatten + \textbf{gefallen}\\
  \end{tabular}
  
  \vspace{-5pt}


\end{verbbox-dat}
%gehören (A2) – to belong
\begin{verbbox-dat}\addcontentsline{toc}{subsubsection}{gehören (to belong)}
    {\Large \bf gehören} (to belong) \hfill Type: \textbf{regular, + Dativ} \hfill Auxiliary Verb: \textbf{haben} 
    
     \vspace{5pt}

    \hrule
    
    \vspace{5pt}

  \begin{tabular}{rl}
     \multicolumn{2}{l}{\textbf{PRÄSENS} }\\
    ich & \textbf{gehöre}\\
    du & \textbf{gehörst}\\
    er/sie/es & \textbf{gehört}\\
    wir & \textbf{gehören}\\
    ihr & \textbf{gehört}\\
    sie/Sie & \textbf{gehören}\\
  \end{tabular}
  \hfill
     \begin{tabular}{rl}
    \multicolumn{2}{l}{\textbf{PRÄTERITUM} }\\
    ich & \textbf{gehörte}\\
    du & \textbf{gehörtest}\\
    er/sie/es & \textbf{gehörte}\\
    wir & \textbf{gehörten}\\
    ihr & \textbf{gehörtet}\\
    sie/Sie & \textbf{gehörten}\\
  \end{tabular}
  \hfill
  \begin{tabular}{rl}
     \multicolumn{2}{l}{\textbf{IMPERATIV} }\\
   & \textbf{gehör(e) (du)!}\\
   & \textbf{gehören wir!}\\
   & \textbf{gehört (ihr)!}\\
   & \textbf{gehören Sie!}\\
   & \vspace{10pt}
  \end{tabular}

    \vspace{10pt}

 \begin{tabular}{rl}
     \multicolumn{2}{l}{\textbf{PERFEKT}}\\
  &  haben + \textbf{gehört}\\
  \end{tabular}
\hfill
   \begin{tabular}{rl}
   
    \multicolumn{2}{l}{\textbf{FUTURE I} }\\
   & werden + \textbf{gehören}\\
 
  \end{tabular}
\hfill
   \begin{tabular}{rl}
      \multicolumn{2}{l}{\textbf{PLUSQUAMPERFEKT}}\\
   & hatten + \textbf{gehört}\\
  \end{tabular}
  
  \vspace{-5pt}


\end{verbbox-dat}
%gehen (A1) – to go
\begin{verbbox-acc}\addcontentsline{toc}{subsubsection}{gehen (to go)}
    {\Large \bf gehen} (to go) \hfill Type: \textbf{irregular} \hfill Auxiliary Verb: \textbf{sein} 
    
     \vspace{5pt}

    \hrule
    
    \vspace{5pt}

  \begin{tabular}{rl}
     \multicolumn{2}{l}{\textbf{PRÄSENS} }\\
    ich & \textbf{gehe}\\
    du & \textbf{gehst}\\
    er/sie/es & \textbf{geht}\\
    wir & \textbf{gehen}\\
    ihr & \textbf{geht}\\
    sie/Sie & \textbf{gehen}\\
  \end{tabular}
  \hfill
     \begin{tabular}{rl}
    \multicolumn{2}{l}{\textbf{PRÄTERITUM} }\\
    ich & \textbf{ging}\\
    du & \textbf{gingst}\\
    er/sie/es & \textbf{ging}\\
    wir & \textbf{gingen}\\
    ihr & \textbf{gingt}\\
    sie/Sie & \textbf{gingen}\\
  \end{tabular}
  \hfill
  \begin{tabular}{rl}
     \multicolumn{2}{l}{\textbf{IMPERATIV} }\\
   & \textbf{geh(e) (du)!}\\
   & \textbf{gehen wir!}\\
   & \textbf{geht (ihr)!}\\
   & \textbf{gehen Sie!}\\
   & \vspace{10pt}
  \end{tabular}

    \vspace{10pt}

 \begin{tabular}{rl}
     \multicolumn{2}{l}{\textbf{PERFEKT}}\\
   & sein + \textbf{gegangen}\\
  \end{tabular}
\hfill
   \begin{tabular}{rl}
   
    \multicolumn{2}{l}{\textbf{FUTURE I} }\\
   & werden + \textbf{gehen}\\
 
  \end{tabular}
\hfill
   \begin{tabular}{rl}
      \multicolumn{2}{l}{\textbf{PLUSQUAMPERFEKT}}\\
   & waren + \textbf{gegangen}\\
  \end{tabular}
  
  \vspace{-5pt}


\end{verbbox-acc}

%gelten (to be valid / to apply)
\begin{verbbox}\addcontentsline{toc}{subsubsection}{gelten (to be valid / to apply)}
    {\Large \bf gelten} (to be valid / to apply) \hfill Type: \textbf{irregular} \hfill Auxiliary Verb: \textbf{haben} 
    
     \vspace{5pt}

    \hrule
    
    \vspace{5pt}

  \begin{tabular}{rl}
     \multicolumn{2}{l}{\textbf{PRÄSENS} }\\
    ich & \textbf{gelte}\\
    du & \textbf{giltst}\\
    er/sie/es & \textbf{gilt}\\
    wir & \textbf{gelten}\\
    ihr & \textbf{geltet}\\
    sie/Sie & \textbf{gelten}\\
  \end{tabular}
  \hfill
     \begin{tabular}{rl}
    \multicolumn{2}{l}{\textbf{PRÄTERITUM} }\\
    ich & \textbf{galt}\\
    du & \textbf{galtest}\\
    er/sie/es & \textbf{galt}\\
    wir & \textbf{galten}\\
    ihr & \textbf{galtet}\\
    sie/Sie & \textbf{galten}\\
  \end{tabular}
  \hfill
  \begin{tabular}{rl}
     \multicolumn{2}{l}{\textbf{IMPERATIV} }\\
   & \textbf{gilt (du)!}\\
   & \textbf{gelten wir!}\\
   & \textbf{geltet (ihr)!}\\
   & \textbf{gelten Sie!}\\
   & \vspace{10pt}
  \end{tabular}

    \vspace{10pt}

 \begin{tabular}{rl}
     \multicolumn{2}{l}{\textbf{PERFEKT}}\\
   & haben + \textbf{gegolten}\\
  \end{tabular}
\hfill
   \begin{tabular}{rl}
    \multicolumn{2}{l}{\textbf{FUTURE I} }\\
   & werden + \textbf{gelten}\\
  \end{tabular}
\hfill
   \begin{tabular}{rl}
      \multicolumn{2}{l}{\textbf{PLUSQUAMPERFEKT}}\\
   & hatten + \textbf{gegolten}\\
  \end{tabular}

   \vspace{10pt}

   \begin{tabular}{lll}
      \multicolumn{3}{l}{\textbf{MIT PRÄPOSITIONEN}}\\
& \textbf{gelten für} + Akk & (to apply to / be valid for)\\
& \textbf{gelten als} + Nom & (to be considered as)\\
\end{tabular}
  \hfill
   \begin{tabular}{lll}
      \multicolumn{3}{l}{\textbf{TRENNBARE VERBEN}}\\
 & \textbf{---} & (none common)\\
\vspace{10pt}
  \end{tabular}

  \vspace{-5pt}

\end{verbbox}

%gießen (to pour / to water)
\begin{verbbox}\addcontentsline{toc}{subsubsection}{gießen (to pour / to water)}
    {\Large \bf gießen} (to pour / to water) \hfill Type: \textbf{irregular} \hfill Auxiliary Verb: \textbf{haben} 
    
     \vspace{5pt}

    \hrule
    
    \vspace{5pt}

  \begin{tabular}{rl}
     \multicolumn{2}{l}{\textbf{PRÄSENS} }\\
    ich & \textbf{gieße}\\
    du & \textbf{gießt}\\
    er/sie/es & \textbf{gießt}\\
    wir & \textbf{gießen}\\
    ihr & \textbf{gießt}\\
    sie/Sie & \textbf{gießen}\\
  \end{tabular}
  \hfill
     \begin{tabular}{rl}
    \multicolumn{2}{l}{\textbf{PRÄTERITUM} }\\
    ich & \textbf{goss}\\
    du & \textbf{gossest}\\
    er/sie/es & \textbf{goss}\\
    wir & \textbf{gossen}\\
    ihr & \textbf{gosst}\\
    sie/Sie & \textbf{gossen}\\
  \end{tabular}
  \hfill
  \begin{tabular}{rl}
     \multicolumn{2}{l}{\textbf{IMPERATIV} }\\
   & \textbf{gieß(e) (du)!}\\
   & \textbf{gießen wir!}\\
   & \textbf{gießt (ihr)!}\\
   & \textbf{gießen Sie!}\\
   & \vspace{10pt}
  \end{tabular}

    \vspace{10pt}

 \begin{tabular}{rl}
     \multicolumn{2}{l}{\textbf{PERFEKT}}\\
   & haben + \textbf{gegossen}\\
  \end{tabular}
\hfill
   \begin{tabular}{rl}
   
    \multicolumn{2}{l}{\textbf{FUTURE I} }\\
   & werden + \textbf{gießen}\\
 
  \end{tabular}
\hfill
   \begin{tabular}{rl}
      \multicolumn{2}{l}{\textbf{PLUSQUAMPERFEKT}}\\
   & hatten + \textbf{gegossen}\\
  \end{tabular}

   \vspace{10pt}

   \begin{tabular}{lll}
      \multicolumn{3}{l}{\textbf{MIT PRÄPOSITIONEN}}\\
& \textbf{gießen in} + Akk & (to pour into)\\
& \textbf{gießen über} + Akk & (to pour over)\\
\end{tabular}
  \hfill
   \begin{tabular}{lll}
      \multicolumn{3}{l}{\textbf{TRENNBARE VERBEN}}\\
 & \textbf{ein$|$gießen} & (to pour in)\\
 & \textbf{aus$|$gießen} & (to pour out)\\
\vspace{10pt}
  \end{tabular}


  \vspace{-5pt}


\end{verbbox}

%genießen (B1) – to enjoy
\begin{verbbox-acc}\addcontentsline{toc}{subsubsection}{genießen (to enjoy)}
    {\Large \bf genießen} (to enjoy) \hfill Type: \textbf{irregular, + Akkusativ} \hfill Auxiliary Verb: \textbf{haben} 
    
     \vspace{5pt}

    \hrule
    
    \vspace{5pt}

  \begin{tabular}{rl}
     \multicolumn{2}{l}{\textbf{PRÄSENS} }\\
    ich & \textbf{genieße}\\
    du & \textbf{genießt}\\
    er/sie/es & \textbf{genießt}\\
    wir & \textbf{genießen}\\
    ihr & \textbf{genießt}\\
    sie/Sie & \textbf{genießen}\\
  \end{tabular}
  \hfill
     \begin{tabular}{rl}
    \multicolumn{2}{l}{\textbf{PRÄTERITUM} }\\
    ich & \textbf{genoss}\\
    du & \textbf{genossest}\\
    er/sie/es & \textbf{genoss}\\
    wir & \textbf{kgenossen}\\
    ihr & \textbf{genosst}\\
    sie/Sie & \textbf{genossen}\\
  \end{tabular}
  \hfill
  \begin{tabular}{rl}
     \multicolumn{2}{l}{\textbf{IMPERATIV} }\\
   & \textbf{genieß(e) (du)!}\\
   & \textbf{genießen wir!}\\
   & \textbf{genießt (ihr)!}\\
   & \textbf{genießen Sie!}\\
   & \vspace{10pt}
  \end{tabular}

    \vspace{10pt}

 \begin{tabular}{rl}
     \multicolumn{2}{l}{\textbf{PERFEKT}}\\
  &  haben + \textbf{genossen}\\
  \end{tabular}
\hfill
   \begin{tabular}{rl}
   
    \multicolumn{2}{l}{\textbf{FUTURE I} }\\
  &  werden + \textbf{genießen}\\
 
  \end{tabular}
\hfill
   \begin{tabular}{rl}
      \multicolumn{2}{l}{\textbf{PLUSQUAMPERFEKT}}\\
  &  hatten + \textbf{genossen}\\
  \end{tabular}
  
  \vspace{-5pt}


\end{verbbox-acc}

%gestikulieren (to gesticulate)
\begin{verbbox}\addcontentsline{toc}{subsubsection}{gestikulieren (to gesticulate)}
{\Large \bf gestikulieren} (to gesticulate) \hfill Type: \textbf{regular} \hfill Auxiliary Verb: \textbf{haben}

\vspace{5pt}
\hrule
\vspace{5pt}

\begin{tabular}{rl}
\multicolumn{2}{l}{\textbf{PRÄSENS}}\\
ich & \textbf{gestikuliere}\\
du & \textbf{gestikulierst}\\
er/sie/es & \textbf{gestikuliert}\\
wir & \textbf{gestikulieren}\\
ihr & \textbf{gestikuliert}\\
sie/Sie & \textbf{gestikulieren}\\
\end{tabular}
\hfill
\begin{tabular}{rl}
\multicolumn{2}{l}{\textbf{PRÄTERITUM}}\\
ich & \textbf{gestikulierte}\\
du & \textbf{gestikuliertest}\\
er/sie/es & \textbf{gestikulierte}\\
wir & \textbf{gestikulierten}\\
ihr & \textbf{gestikuliertet}\\
sie/Sie & \textbf{gestikulierten}\\
\end{tabular}
\hfill
\begin{tabular}{rl}
\multicolumn{2}{l}{\textbf{IMPERATIV}}\\
& \textbf{gestikuliere (du)!}\\
& \textbf{gestikulieren wir!}\\
& \textbf{gestikuliert (ihr)!}\\
& \textbf{gestikulieren Sie!}\\
\end{tabular}

\vspace{10pt}

\begin{tabular}{rl}
\multicolumn{2}{l}{\textbf{PERFEKT}}\\
& haben + \textbf{gestikuliert}\\
\end{tabular}
\hfill
\begin{tabular}{rl}
\multicolumn{2}{l}{\textbf{FUTURE I}}\\
& werden + \textbf{gestikulieren}\\
\end{tabular}
\hfill
\begin{tabular}{rl}
\multicolumn{2}{l}{\textbf{PLUSQUAMPERFEKT}}\\
& hatten + \textbf{gestikuliert}\\
\end{tabular}

\vspace{10pt}

\begin{tabular}{lll}
\multicolumn{3}{l}{\textbf{MIT PRÄPOSITIONEN}}\\
& \textbf{gestikulieren mit} + Dat & (to gesticulate with)\\
\end{tabular}
\hfill
\begin{tabular}{lll}
\multicolumn{3}{l}{\textbf{TRENNBARE VERBEN}}\\
& --- & (no separable forms)\\
\end{tabular}
\vspace{-5pt}
\end{verbbox}


%gewinnen (B1) – to win
\begin{verbbox-acc}\addcontentsline{toc}{subsubsection}{gewinnen (to win)}
    {\Large \bf gewinnen} (to win) \hfill Type: \textbf{irregular, + Akkusativ} \hfill Auxiliary Verb: \textbf{haben} 
    
     \vspace{5pt}

    \hrule
    
    \vspace{5pt}

  \begin{tabular}{rl}
     \multicolumn{2}{l}{\textbf{PRÄSENS} }\\
    ich & \textbf{gewinne}\\
    du & \textbf{gewinnst}\\
    er/sie/es & \textbf{gewinnt}\\
    wir & \textbf{gewinnen}\\
    ihr & \textbf{gewinnt}\\
    sie/Sie & \textbf{gewinnen}\\
  \end{tabular}
  \hfill
     \begin{tabular}{rl}
    \multicolumn{2}{l}{\textbf{PRÄTERITUM} }\\
    ich & \textbf{gewann}\\
    du & \textbf{gewannst}\\
    er/sie/es & \textbf{gewann}\\
    wir & \textbf{gewannen}\\
    ihr & \textbf{gewannt}\\
    sie/Sie & \textbf{gewannen}\\
  \end{tabular}
  \hfill
  \begin{tabular}{rl}
     \multicolumn{2}{l}{\textbf{IMPERATIV} }\\
   & \textbf{gewinn(e) (du)!}\\
   & \textbf{gewinnen wir!}\\
   & \textbf{gewinnt (ihr)!}\\
   & \textbf{gewinnen Sie!}\\
   & \vspace{10pt}
  \end{tabular}

    \vspace{10pt}

 \begin{tabular}{rl}
     \multicolumn{2}{l}{\textbf{PERFEKT}}\\
  &  haben + \textbf{gewonnen}\\
  \end{tabular}
\hfill
   \begin{tabular}{rl}
   
    \multicolumn{2}{l}{\textbf{FUTURE I} }\\
  &  werden + \textbf{gewinnen}\\
 
  \end{tabular}
\hfill
   \begin{tabular}{rl}
      \multicolumn{2}{l}{\textbf{PLUSQUAMPERFEKT}}\\
  &  hatten + \textbf{gewonnen}\\
  \end{tabular}
  
  \vspace{-5pt}


\end{verbbox-acc}
%glauben (A1) – to believe
\begin{verbbox-dat}\addcontentsline{toc}{subsubsection}{glauben (to believe)}
    {\Large \bf glauben} (to believe) \hfill Type: \textbf{regular, + Dativ} \hfill Auxiliary Verb: \textbf{haben} 
    
     \vspace{5pt}

    \hrule
    
    \vspace{5pt}

  \begin{tabular}{rl}
     \multicolumn{2}{l}{\textbf{PRÄSENS} }\\
    ich & \textbf{glaube}\\
    du & \textbf{glaubst}\\
    er/sie/es & \textbf{glaubt}\\
    wir & \textbf{glauben}\\
    ihr & \textbf{glaubt}\\
    sie/Sie & \textbf{glauben}\\
  \end{tabular}
  \hfill
     \begin{tabular}{rl}
    \multicolumn{2}{l}{\textbf{PRÄTERITUM} }\\
    ich & \textbf{glaubte}\\
    du & \textbf{glaubtest}\\
    er/sie/es & \textbf{glaubte}\\
    wir & \textbf{glaubten}\\
    ihr & \textbf{glaubtet}\\
    sie/Sie & \textbf{glaubten}\\
  \end{tabular}
  \hfill
  \begin{tabular}{rl}
     \multicolumn{2}{l}{\textbf{IMPERATIV} }\\
   & \textbf{glaub(e) (du)!}\\
   & \textbf{glauben wir!}\\
   & \textbf{glaubt (ihr)!}\\
   & \textbf{glauben Sie!}\\
   & \vspace{10pt}
  \end{tabular}

    \vspace{10pt}

 \begin{tabular}{rl}
     \multicolumn{2}{l}{\textbf{PERFEKT}}\\
   & haben + \textbf{geglaubt}\\
  \end{tabular}
\hfill
   \begin{tabular}{rl}
   
    \multicolumn{2}{l}{\textbf{FUTURE I} }\\
   & werden + \textbf{glauben}\\
 
  \end{tabular}
\hfill
   \begin{tabular}{rl}
      \multicolumn{2}{l}{\textbf{PLUSQUAMPERFEKT}}\\
   & hatten + \textbf{geglaubt}\\
  \end{tabular}
  
\vspace{10pt}

\begin{tabular}{lll}
\multicolumn{3}{l}{\textbf{MIT PRÄPOSITIONEN}}\\
& \textbf{glauben an} + Akk & (to believe in)\\
\end{tabular}
\hfill
\begin{tabular}{lll}
\multicolumn{3}{l}{\textbf{TRENNBARE VERBEN}}\\
& --- & \\
\end{tabular}

  \vspace{-5pt}


\end{verbbox-dat}


%grillen (to grill / barbecue)
\begin{verbbox-acc}\addcontentsline{toc}{subsubsection}{grillen (to grill / barbecue)}
{\Large \bf grillen} (to grill / barbecue) \hfill Type: \textbf{regular, + Akkusativ} \hfill Auxiliary Verb: \textbf{haben}

\vspace{5pt}
\hrule
\vspace{5pt}

\begin{tabular}{rl}
\multicolumn{2}{l}{\textbf{PRÄSENS}}\\
ich & \textbf{grille}\\
du & \textbf{grillst}\\
er/sie/es & \textbf{grillt}\\
wir & \textbf{grillen}\\
ihr & \textbf{grillt}\\
sie/Sie & \textbf{grillen}\\
\end{tabular}
\hfill
\begin{tabular}{rl}
\multicolumn{2}{l}{\textbf{PRÄTERITUM}}\\
ich & \textbf{grillte}\\
du & \textbf{grilltest}\\
er/sie/es & \textbf{grillte}\\
wir & \textbf{grillten}\\
ihr & \textbf{grilltet}\\
sie/Sie & \textbf{grillten}\\
\end{tabular}
\hfill
\begin{tabular}{rl}
\multicolumn{2}{l}{\textbf{IMPERATIV}}\\
& \textbf{grille (du)!}\\
& \textbf{grillen wir!}\\
& \textbf{grillt (ihr)!}\\
& \textbf{grillen Sie!}\\
\end{tabular}

\vspace{10pt}

\begin{tabular}{rl}
\multicolumn{2}{l}{\textbf{PERFEKT}}\\
& haben + \textbf{gegrillt}\\
\end{tabular}
\hfill
\begin{tabular}{rl}
\multicolumn{2}{l}{\textbf{FUTURE I}}\\
& werden + \textbf{grillen}\\
\end{tabular}
\hfill
\begin{tabular}{rl}
\multicolumn{2}{l}{\textbf{PLUSQUAMPERFEKT}}\\
& hatten + \textbf{gegrillt}\\
\end{tabular}

\vspace{10pt}

\begin{tabular}{lll}
\multicolumn{3}{l}{\textbf{MIT PRÄPOSITIONEN}}\\
& \textbf{grillen auf} + Dat & (to grill on)\\
\end{tabular}
\hfill
\begin{tabular}{lll}
\multicolumn{3}{l}{\textbf{TRENNBARE VERBEN}}\\
& an$|$grillen & (to lightly grill / start grilling)\\
\end{tabular}

\vspace{-5pt}
\end{verbbox-acc}

%gurgeln (to gargle)
\begin{verbbox}\addcontentsline{toc}{subsubsection}{gurgeln (to gargle)}
{\Large \bf gurgeln} (to gargle) \hfill Type: \textbf{regular} \hfill Auxiliary Verb: \textbf{haben}

\vspace{5pt}
\hrule
\vspace{5pt}

\begin{tabular}{rl}
\multicolumn{2}{l}{\textbf{PRÄSENS}}\\
ich & \textbf{gurgle}\\
du & \textbf{gurgelst}\\
er/sie/es & \textbf{gurgelt}\\
wir & \textbf{gurgeln}\\
ihr & \textbf{gurgelt}\\
sie/Sie & \textbf{gurgeln}\\
\end{tabular}
\hfill
\begin{tabular}{rl}
\multicolumn{2}{l}{\textbf{PRÄTERITUM}}\\
ich & \textbf{gurgelte}\\
du & \textbf{gurgeltest}\\
er/sie/es & \textbf{gurgelte}\\
wir & \textbf{gurgelten}\\
ihr & \textbf{gurgeltet}\\
sie/Sie & \textbf{gurgelten}\\
\end{tabular}
\hfill
\begin{tabular}{rl}
\multicolumn{2}{l}{\textbf{IMPERATIV}}\\
& \textbf{gurgle (du)!}\\
& \textbf{gurgeln wir!}\\
& \textbf{gurgelt (ihr)!}\\
& \textbf{gurgeln Sie!}\\
\end{tabular}

\vspace{10pt}

\begin{tabular}{rl}
\multicolumn{2}{l}{\textbf{PERFEKT}}\\
& haben + \textbf{gegurgelt}\\
\end{tabular}
\hfill
\begin{tabular}{rl}
\multicolumn{2}{l}{\textbf{FUTURE I}}\\
& werden + \textbf{gurgeln}\\
\end{tabular}
\hfill
\begin{tabular}{rl}
\multicolumn{2}{l}{\textbf{PLUSQUAMPERFEKT}}\\
& hatten + \textbf{gegurgelt}\\
\end{tabular}

\vspace{10pt}

\begin{tabular}{lll}
\multicolumn{3}{l}{\textbf{MIT PRÄPOSITIONEN}}\\
& \textbf{gurgeln mit} + Dat & (to gargle with)\\
\end{tabular}
\hfill
\begin{tabular}{lll}
\multicolumn{3}{l}{\textbf{TRENNBARE VERBEN}}\\
& ---\\
\end{tabular}
\vspace{-5pt}
\end{verbbox}

\subsection{H Verbs}

%halten (B1) – to hold, stop
\begin{verbbox}\addcontentsline{toc}{subsubsection}{halten (to hold, stop)}
    {\Large \bf halten} (to hold, stop) \hfill Type: \textbf{irregular} \hfill Auxiliary Verb: \textbf{haben} 
    
     \vspace{5pt}

    \hrule
    
    \vspace{5pt}

  \begin{tabular}{rl}
     \multicolumn{2}{l}{\textbf{PRÄSENS} }\\
    ich & \textbf{halte}\\
    du & \textbf{hältst}\\
    er/sie/es & \textbf{hält}\\
    wir & \textbf{halten}\\
    ihr & \textbf{haltet}\\
    sie/Sie & \textbf{halten}\\
  \end{tabular}
  \hfill
     \begin{tabular}{rl}
    \multicolumn{2}{l}{\textbf{PRÄTERITUM} }\\
    ich & \textbf{hielt}\\
    du & \textbf{hieltst}\\
    er/sie/es & \textbf{hielt}\\
    wir & \textbf{hielten}\\
    ihr & \textbf{hieltet}\\
    sie/Sie & \textbf{hielten}\\
  \end{tabular}
  \hfill
  \begin{tabular}{rl}
     \multicolumn{2}{l}{\textbf{IMPERATIV} }\\
   & \textbf{halt(e) (du)!}\\
   & \textbf{halten wir!}\\
   & \textbf{haltet (ihr)!}\\
   & \textbf{halten Sie!}\\
   & \vspace{10pt}
  \end{tabular}

    \vspace{10pt}

 \begin{tabular}{rl}
     \multicolumn{2}{l}{\textbf{PERFEKT}}\\
  &  haben + \textbf{gehalten}\\
  \end{tabular}
\hfill
   \begin{tabular}{rl}
   
    \multicolumn{2}{l}{\textbf{FUTURE I} }\\
   & werden + \textbf{halten}\\
 
  \end{tabular}
\hfill
   \begin{tabular}{rl}
      \multicolumn{2}{l}{\textbf{PLUSQUAMPERFEKT}}\\
   & hatten + \textbf{gehalten}\\
  \end{tabular}
  
  \vspace{10pt}



\begin{tabular}{lll}
\multicolumn{3}{l}{\textbf{MIT PRÄPOSITIONEN}}\\
& ---\\
\end{tabular}
\hfill
\begin{tabular}{lll}
\multicolumn{3}{l}{\textbf{TRENNBARE VERBEN}}\\
& ---\\
\end{tabular}

  \vspace{-5pt}

\end{verbbox}

%heben - to lift/raise
\begin{verbbox}\addcontentsline{toc}{subsubsection}{heben (to lift / raise)}
    {\Large \bf heben} (to lift / raise) \hfill Type: \textbf{irregular} \hfill Auxiliary Verb: \textbf{haben} 
    
     \vspace{5pt}

    \hrule
    
    \vspace{5pt}

  \begin{tabular}{rl}
     \multicolumn{2}{l}{\textbf{PRÄSENS} }\\
    ich & \textbf{hebe}\\
    du & \textbf{hebst}\\
    er/sie/es & \textbf{hebt}\\
    wir & \textbf{heben}\\
    ihr & \textbf{hebt}\\
    sie/Sie & \textbf{heben}\\
  \end{tabular}
  \hfill
     \begin{tabular}{rl}
    \multicolumn{2}{l}{\textbf{PRÄTERITUM} }\\
    ich & \textbf{hob}\\
    du & \textbf{hobst}\\
    er/sie/es & \textbf{hob}\\
    wir & \textbf{hoben}\\
    ihr & \textbf{hobt}\\
    sie/Sie & \textbf{hoben}\\
  \end{tabular}
  \hfill
  \begin{tabular}{rl}
     \multicolumn{2}{l}{\textbf{IMPERATIV} }\\
   & \textbf{heb(e) (du)!}\\
   & \textbf{heben wir!}\\
   & \textbf{hebt (ihr)!}\\
   & \textbf{heben Sie!}\\
   & \vspace{10pt}
  \end{tabular}

    \vspace{10pt}

 \begin{tabular}{rl}
     \multicolumn{2}{l}{\textbf{PERFEKT}}\\
   & haben + \textbf{gehoben}\\
  \end{tabular}
\hfill
   \begin{tabular}{rl}
   
    \multicolumn{2}{l}{\textbf{FUTURE I} }\\
   & werden + \textbf{heben}\\
 
  \end{tabular}
\hfill
   \begin{tabular}{rl}
      \multicolumn{2}{l}{\textbf{PLUSQUAMPERFEKT}}\\
   & hatten + \textbf{gehoben}\\
  \end{tabular}

   \vspace{10pt}

   \begin{tabular}{lll}
      \multicolumn{3}{l}{\textbf{MIT PRÄPOSITIONEN}}\\
& \textbf{heben auf} + Akk & (to pick up / keep / store)\\
& \textbf{heben von} + Dat & (to lift from)\\
& \textbf{heben aus} + Dat & (to lift out of)\\
\vspace{10pt}
   \end{tabular}
  \hfill
   \begin{tabular}{lll}
      \multicolumn{3}{l}{\textbf{TRENNBARE VERBEN}}\\
 & \textbf{ab$|$heben} & (to lift off / withdraw money)\\
 & \textbf{auf$|$heben} & (to pick up / cancel / keep)\\
 & \textbf{an$|$heben} & (to begin / raise slightly)\\
 & \textbf{heraus$|$heben} & (to lift out)\\
  \end{tabular}

\vspace{-5pt}


\end{verbbox}

%heiraten (to marry)
\begin{verbbox}\addcontentsline{toc}{subsubsection}{heiraten (to marry)}
{\Large \bf heiraten} (to marry) \hfill Type: \textbf{regular} \hfill Auxiliary Verb: \textbf{haben} 

\vspace{5pt}
\hrule
\vspace{5pt}

\begin{tabular}{rl}
\multicolumn{2}{l}{\textbf{PRÄSENS}}\\
ich & \textbf{heirate}\\
du & \textbf{heiratest}\\
er/sie/es & \textbf{heiratet}\\
wir & \textbf{heiraten}\\
ihr & \textbf{heiratet}\\
sie/Sie & \textbf{heiraten}\\
\end{tabular}
\hfill
\begin{tabular}{rl}
\multicolumn{2}{l}{\textbf{PRÄTERITUM}}\\
ich & \textbf{heiratete}\\
du & \textbf{heiratetest}\\
er/sie/es & \textbf{heiratete}\\
wir & \textbf{heirateten}\\
ihr & \textbf{heiratetet}\\
sie/Sie & \textbf{heirateten}\\
\end{tabular}
\hfill
\begin{tabular}{rl}
\multicolumn{2}{l}{\textbf{IMPERATIV}}\\
& \textbf{heirate (du)!}\\
& \textbf{heiraten wir!}\\
& \textbf{heiratet (ihr)!}\\
& \textbf{heiraten Sie!}\\
\end{tabular}

\vspace{10pt}

\begin{tabular}{rl}
\multicolumn{2}{l}{\textbf{PERFEKT}}\\
& haben + \textbf{geheiratet}\\
\end{tabular}
\hfill
\begin{tabular}{rl}
\multicolumn{2}{l}{\textbf{FUTURE I}}\\
& werden + \textbf{heiraten}\\
\end{tabular}
\hfill
\begin{tabular}{rl}
\multicolumn{2}{l}{\textbf{PLUSQUAMPERFEKT}}\\
& hatten + \textbf{geheiratet}\\
\end{tabular}

\vspace{10pt}

\begin{tabular}{lll}
\multicolumn{3}{l}{\textbf{MIT PRÄPOSITIONEN}}\\
& \textbf{heiraten} + Akk & (to marry someone)\\
& \textbf{heiraten mit} + Dat & (to marry, formal)\\
\end{tabular}
\hfill
\begin{tabular}{lll}
\multicolumn{3}{l}{\textbf{TRENNBARE VERBEN}}\\
& ---\\
\vspace{10pt}
\end{tabular}

\vspace{-5pt}
\end{verbbox}



%heißen (A1) – to be called
\begin{verbbox}\addcontentsline{toc}{subsubsection}{heißen (to be called)}
    {\Large \bf heißen} (to be called) \hfill Type: \textbf{irregular} \hfill Auxiliary Verb: \textbf{haben} 
    
     \vspace{5pt}

    \hrule
    
    \vspace{5pt}

  \begin{tabular}{rl}
     \multicolumn{2}{l}{\textbf{PRÄSENS} }\\
    ich & \textbf{heiße}\\
    du & \textbf{heißt}\\
    er/sie/es & \textbf{heißt}\\
    wir & \textbf{heißen}\\
    ihr & \textbf{heißt}\\
    sie/Sie & \textbf{heißen}\\
  \end{tabular}
  \hfill
     \begin{tabular}{rl}
    \multicolumn{2}{l}{\textbf{PRÄTERITUM} }\\
    ich & \textbf{hieß}\\
    du & \textbf{hießest}\\
    er/sie/es & \textbf{hieß}\\
    wir & \textbf{hießen}\\
    ihr & \textbf{hießt}\\
    sie/Sie & \textbf{hießen}\\
  \end{tabular}
  \hfill
  \begin{tabular}{rl}
     \multicolumn{2}{l}{\textbf{IMPERATIV} }\\
   & \textbf{heiß(e) (du)!}\\
   & \textbf{heißen wir!}\\
   & \textbf{heißt (ihr)!}\\
   & \textbf{heißen Sie!}\\
   & \vspace{10pt}
  \end{tabular}

    \vspace{10pt}

 \begin{tabular}{rl}
     \multicolumn{2}{l}{\textbf{PERFEKT}}\\
   & haben + \textbf{geheißen}\\
  \end{tabular}
\hfill
   \begin{tabular}{rl}
   
    \multicolumn{2}{l}{\textbf{FUTURE I} }\\
   & werden + \textbf{heißen}\\
 
  \end{tabular}
\hfill
   \begin{tabular}{rl}
      \multicolumn{2}{l}{\textbf{PLUSQUAMPERFEKT}}\\
   & hatten + \textbf{geheißen}\\
  \end{tabular}
  
  \vspace{10pt}



\begin{tabular}{lll}
\multicolumn{3}{l}{\textbf{MIT PRÄPOSITIONEN}}\\
& ---\\
\end{tabular}
\hfill
\begin{tabular}{lll}
\multicolumn{3}{l}{\textbf{TRENNBARE VERBEN}}\\
& ---\\
\end{tabular}

  \vspace{-5pt}

\end{verbbox}
%helfen (A1) – to help
\begin{verbbox-dat}\addcontentsline{toc}{subsubsection}{helfen (to help)}
    {\Large \bf helfen} (to help) \hfill Type: \textbf{irregular, + Dativ} \hfill Auxiliary Verb: \textbf{haben} 
    
     \vspace{5pt}

    \hrule
    
    \vspace{5pt}

  \begin{tabular}{rl}
     \multicolumn{2}{l}{\textbf{PRÄSENS} }\\
    ich & \textbf{helfe}\\
    du & \textbf{hilfst}\\
    er/sie/es & \textbf{hilft}\\
    wir & \textbf{helfen}\\
    ihr & \textbf{helft}\\
    sie/Sie & \textbf{helfen}\\
  \end{tabular}
  \hfill
     \begin{tabular}{rl}
    \multicolumn{2}{l}{\textbf{PRÄTERITUM} }\\
    ich & \textbf{half}\\
    du & \textbf{halfst}\\
    er/sie/es & \textbf{half}\\
    wir & \textbf{halfen}\\
    ihr & \textbf{halft}\\
    sie/Sie & \textbf{halfen}\\
  \end{tabular}
  \hfill
  \begin{tabular}{rl}
     \multicolumn{2}{l}{\textbf{IMPERATIV} }\\
   & \textbf{hilf (du)!}\\
   & \textbf{helfen wir!}\\
   & \textbf{helft (ihr)!}\\
   & \textbf{helfen Sie!}\\
   & \vspace{10pt}
  \end{tabular}

    \vspace{10pt}

 \begin{tabular}{rl}
     \multicolumn{2}{l}{\textbf{PERFEKT}}\\
  &  haben + \textbf{geholfen}\\
  \end{tabular}
\hfill
   \begin{tabular}{rl}
   
    \multicolumn{2}{l}{\textbf{FUTURE I} }\\
   & werden + \textbf{helfen}\\
 
  \end{tabular}
\hfill
   \begin{tabular}{rl}
      \multicolumn{2}{l}{\textbf{PLUSQUAMPERFEKT}}\\
   & hatten + \textbf{geholfen}\\
  \end{tabular}
  
  \vspace{-5pt}


\end{verbbox-dat}
%holen (A1) – to fetch
\begin{verbbox-acc}\addcontentsline{toc}{subsubsection}{holen (to fetch)}
    {\Large \bf holen} (to fetch) \hfill Type: \textbf{regular, + Akkusativ} \hfill Auxiliary Verb: \textbf{haben} 
    
     \vspace{5pt}

    \hrule
    
    \vspace{5pt}

  \begin{tabular}{rl}
     \multicolumn{2}{l}{\textbf{PRÄSENS} }\\
    ich & \textbf{hole}\\
    du & \textbf{holst}\\
    er/sie/es & \textbf{holt}\\
    wir & \textbf{holen}\\
    ihr & \textbf{holt}\\
    sie/Sie & \textbf{holen}\\
  \end{tabular}
  \hfill
     \begin{tabular}{rl}
    \multicolumn{2}{l}{\textbf{PRÄTERITUM} }\\
    ich & \textbf{gholte}\\
    du & \textbf{holtest}\\
    er/sie/es & \textbf{holte}\\
    wir & \textbf{holten}\\
    ihr & \textbf{holtet}\\
    sie/Sie & \textbf{holten}\\
  \end{tabular}
  \hfill
  \begin{tabular}{rl}
     \multicolumn{2}{l}{\textbf{IMPERATIV} }\\
   & \textbf{hol(e) (du)!}\\
   & \textbf{holen wir!}\\
   & \textbf{holt (ihr)!}\\
   & \textbf{holen Sie!}\\
   & \vspace{10pt}
  \end{tabular}

    \vspace{10pt}

 \begin{tabular}{rl}
     \multicolumn{2}{l}{\textbf{PERFEKT}}\\
   & haben + \textbf{geholt}\\
  \end{tabular}
\hfill
   \begin{tabular}{rl}
   
    \multicolumn{2}{l}{\textbf{FUTURE I} }\\
   & werden + \textbf{holen}\\
 
  \end{tabular}
\hfill
   \begin{tabular}{rl}
      \multicolumn{2}{l}{\textbf{PLUSQUAMPERFEKT}}\\
   & hatten + \textbf{geholt}\\
  \end{tabular}
  
  \vspace{10pt}

   \begin{tabular}{lll}
      \multicolumn{3}{l}{\textbf{MIT PRÄPOSITIONEN}}\\
& \textbf{holen für} + Akk & (to fetch something for someone)\\
& \textbf{holen aus} + Dat & (to take out of / fetch from inside)\\
&\textbf{holen nach} + Dat & (to fetch to a location)\\
&\textbf{abholen von} + Akk & (to pick up from)
  \end{tabular}
  \hfill
   \begin{tabular}{lll}
      \multicolumn{3}{l}{\textbf{TRENNBARE VERBEN}}\\
 & \textbf{ab$|$holen} & (to pick up / collect)\\
 & \textbf{auf$|$holen} & (to catch up)\\
     & \textbf{ein$|$holen} & (to catch up / retrieve)\\
     & \textbf{nach$|$holen} & (to make up / catch up later)
  \end{tabular}


  \vspace{-5pt}

\end{verbbox-acc}
%hoffen (A2) – to hope
\begin{verbbox-acc}\addcontentsline{toc}{subsubsection}{hoffen (to hope)}
    {\Large \bf hoffen} (to hope) \hfill Type: \textbf{regular, + Akkusativ} \hfill Auxiliary Verb: \textbf{haben} 
    
     \vspace{5pt}

    \hrule
    
    \vspace{5pt}

  \begin{tabular}{rl}
     \multicolumn{2}{l}{\textbf{PRÄSENS} }\\
    ich & \textbf{hoffe}\\
    du & \textbf{hoffst}\\
    er/sie/es & \textbf{hofft}\\
    wir & \textbf{hoffen}\\
    ihr & \textbf{hofft}\\
    sie/Sie & \textbf{hoffen}\\
  \end{tabular}
  \hfill
     \begin{tabular}{rl}
    \multicolumn{2}{l}{\textbf{PRÄTERITUM} }\\
    ich & \textbf{hoffte}\\
    du & \textbf{hofftest}\\
    er/sie/es & \textbf{hoffte}\\
    wir & \textbf{hofften}\\
    ihr & \textbf{hofftet}\\
    sie/Sie & \textbf{hofften}\\
  \end{tabular}
  \hfill
  \begin{tabular}{rl}
     \multicolumn{2}{l}{\textbf{IMPERATIV} }\\
   & \textbf{hoff(e) (du)!}\\
   & \textbf{hoffen wir!}\\
   & \textbf{hofft (ihr)!}\\
   & \textbf{hoffen Sie!}\\
   & \vspace{10pt}
  \end{tabular}

    \vspace{10pt}

 \begin{tabular}{rl}
     \multicolumn{2}{l}{\textbf{PERFEKT}}\\
   & haben + \textbf{gehofft}\\
  \end{tabular}
\hfill
   \begin{tabular}{rl}
   
    \multicolumn{2}{l}{\textbf{FUTURE I} }\\
   & werden + \textbf{hoffen}\\
 
  \end{tabular}
\hfill
   \begin{tabular}{rl}
      \multicolumn{2}{l}{\textbf{PLUSQUAMPERFEKT}}\\
   & hatten + \textbf{gehofft}\\
  \end{tabular}
  
  \vspace{-5pt}


\end{verbbox-acc}
%hören (A1) – to hear
\begin{verbbox-acc}\addcontentsline{toc}{subsubsection}{hören (to hear)}
    {\Large \bf hören} (to hear) \hfill Type: \textbf{regular, + Akkusativ} \hfill Auxiliary Verb: \textbf{haben} 
    
     \vspace{5pt}

    \hrule
    
    \vspace{5pt}

  \begin{tabular}{rl}
     \multicolumn{2}{l}{\textbf{PRÄSENS} }\\
    ich & \textbf{höre}\\
    du & \textbf{hörst}\\
    er/sie/es & \textbf{hört}\\
    wir & \textbf{hören}\\
    ihr & \textbf{hört}\\
    sie/Sie & \textbf{hören}\\
  \end{tabular}
  \hfill
     \begin{tabular}{rl}
    \multicolumn{2}{l}{\textbf{PRÄTERITUM} }\\
    ich & \textbf{hörte}\\
    du & \textbf{hörtest}\\
    er/sie/es & \textbf{hörte}\\
    wir & \textbf{hörten}\\
    ihr & \textbf{hörtet}\\
    sie/Sie & \textbf{hörten}\\
  \end{tabular}
  \hfill
  \begin{tabular}{rl}
     \multicolumn{2}{l}{\textbf{IMPERATIV} }\\
   & \textbf{hör(e) (du)!}\\
   & \textbf{hören wir!}\\
   & \textbf{hört (ihr)!}\\
   & \textbf{hören Sie!}\\
   & \vspace{10pt}
  \end{tabular}

    \vspace{10pt}

 \begin{tabular}{rl}
     \multicolumn{2}{l}{\textbf{PERFEKT}}\\
    &haben + \textbf{gehört}\\
  \end{tabular}
\hfill
   \begin{tabular}{rl}
   
    \multicolumn{2}{l}{\textbf{FUTURE I} }\\
    &werden + \textbf{hören}\\
 
  \end{tabular}
\hfill
   \begin{tabular}{rl}
      \multicolumn{2}{l}{\textbf{PLUSQUAMPERFEKT}}\\
    &hatten + \textbf{gehört}\\
  \end{tabular}
  

  \vspace{10pt}

   \begin{tabular}{lll}
      \multicolumn{3}{l}{\textbf{MIT PRÄPOSITIONEN}}\\
& \textbf{hören auf} + Akk & (to listen to / obey / pay attention to)\\
& \textbf{hören von} + Dat & (to hear about / receive news)\\
\vspace{30pt}
  \end{tabular}
  \hfill
   \begin{tabular}{lll}
      \multicolumn{3}{l}{\textbf{TRENNBARE VERBEN}}\\
 & \textbf{ab$|$hören} & (to listen to / examine)\\
 & \textbf{an$|$hören} & (to hear completely)\\
 & \textbf{auf$|$hören} & (to stop / quit)\\
  & \textbf{mit$|$hören} & (to listen in / eavesdrop)\\
  & \textbf{nach$|$hören} & (to check / listen again)\\
   & \textbf{zu$|$hören} & (to listen carefully)\\
  \end{tabular}

  \vspace{-5pt}

\end{verbbox-acc}

%husten (to cough)
\begin{verbbox}\addcontentsline{toc}{subsubsection}{husten (to cough)}
{\Large \bf husten} (to cough) \hfill Type: \textbf{regular} \hfill Auxiliary Verb: \textbf{haben}

\vspace{5pt}
\hrule
\vspace{5pt}

\begin{tabular}{rl}
\multicolumn{2}{l}{\textbf{PRÄSENS}}\\
ich & \textbf{huste}\\
du & \textbf{hustest}\\
er/sie/es & \textbf{hustet}\\
wir & \textbf{husten}\\
ihr & \textbf{hustet}\\
sie/Sie & \textbf{husten}\\
\end{tabular}
\hfill
\begin{tabular}{rl}
\multicolumn{2}{l}{\textbf{PRÄTERITUM}}\\
ich & \textbf{hustete}\\
du & \textbf{hustetest}\\
er/sie/es & \textbf{hustete}\\
wir & \textbf{husteten}\\
ihr & \textbf{hustetet}\\
sie/Sie & \textbf{husteten}\\
\end{tabular}
\hfill
\begin{tabular}{rl}
\multicolumn{2}{l}{\textbf{IMPERATIV}}\\
& \textbf{huste (du)!}\\
& \textbf{husten wir!}\\
& \textbf{hustet (ihr)!}\\
& \textbf{husten Sie!}\\
\end{tabular}

\vspace{10pt}

\begin{tabular}{rl}
\multicolumn{2}{l}{\textbf{PERFEKT}}\\
& haben + \textbf{gehustet}\\
\end{tabular}
\hfill
\begin{tabular}{rl}
\multicolumn{2}{l}{\textbf{FUTURE I}}\\
& werden + \textbf{husten}\\
\end{tabular}
\hfill
\begin{tabular}{rl}
\multicolumn{2}{l}{\textbf{PLUSQUAMPERFEKT}}\\
& hatten + \textbf{gehustet}\\
\end{tabular}

\vspace{10pt}

\begin{tabular}{lll}
\multicolumn{3}{l}{\textbf{MIT PRÄPOSITIONEN}}\\
& \textbf{husten in} + Akk & (to cough into)\\
& \textbf{husten vor} + Dat & (to cough from / because of)\\
\end{tabular}
\hfill
\begin{tabular}{lll}
\multicolumn{3}{l}{\textbf{TRENNBARE VERBEN}}\\
& \textbf{aus$|$husten} & (to cough out)\\
\end{tabular}
\vspace{-5pt}
\end{verbbox}

\subsection{I Verbs}

%impfen (to vaccinate / inoculate)
\begin{verbbox}\addcontentsline{toc}{subsubsection}{impfen (to vaccinate / inoculate)}
{\Large \bf impfen} (to vaccinate / inoculate) \hfill Type: \textbf{regular} \hfill Auxiliary Verb: \textbf{haben}

\vspace{5pt}
\hrule
\vspace{5pt}

\begin{tabular}{rl}
\multicolumn{2}{l}{\textbf{PRÄSENS}}\\
ich & \textbf{impfe}\\
du & \textbf{impfst}\\
er/sie/es & \textbf{impft}\\
wir & \textbf{impfen}\\
ihr & \textbf{impft}\\
sie/Sie & \textbf{impfen}\\
\end{tabular}
\hfill
\begin{tabular}{rl}
\multicolumn{2}{l}{\textbf{PRÄTERITUM}}\\
ich & \textbf{impfte}\\
du & \textbf{impftest}\\
er/sie/es & \textbf{impfte}\\
wir & \textbf{impften}\\
ihr & \textbf{impftet}\\
sie/Sie & \textbf{impften}\\
\end{tabular}
\hfill
\begin{tabular}{rl}
\multicolumn{2}{l}{\textbf{IMPERATIV}}\\
& \textbf{impfe (du)!}\\
& \textbf{impfen wir!}\\
& \textbf{impft (ihr)!}\\
& \textbf{impfen Sie!}\\
\end{tabular}

\vspace{10pt}

\begin{tabular}{rl}
\multicolumn{2}{l}{\textbf{PERFEKT}}\\
& haben + \textbf{geimpft}\\
\end{tabular}
\hfill
\begin{tabular}{rl}
\multicolumn{2}{l}{\textbf{FUTURE I}}\\
& werden + \textbf{impfen}\\
\end{tabular}
\hfill
\begin{tabular}{rl}
\multicolumn{2}{l}{\textbf{PLUSQUAMPERFEKT}}\\
& hatten + \textbf{geimpft}\\
\end{tabular}

\vspace{10pt}

\begin{tabular}{lll}
\multicolumn{3}{l}{\textbf{MIT PRÄPOSITIONEN}}\\
& \textbf{impfen gegen} + Akk & (to vaccinate against)\\
& \textbf{impfen für} + Akk & (to vaccinate for / on behalf of)\\
\end{tabular}
\hfill
\begin{tabular}{lll}
\multicolumn{3}{l}{\textbf{TRENNBARE VERBEN}}\\
& --- & (no separable forms)\\
\end{tabular}
\vspace{-5pt}
\end{verbbox}

%informaterien
\begin{verbbox}\addcontentsline{toc}{subsubsection}{informieren (to inform)}
    {\Large \bf informieren} (to inform) \hfill Type: \textbf{regular} \hfill Auxiliary Verb: \textbf{haben} 
    
     \vspace{5pt}

    \hrule
    
    \vspace{5pt}

  \begin{tabular}{rl}
     \multicolumn{2}{l}{\textbf{PRÄSENS} }\\
    ich & \textbf{informiere}\\
    du & \textbf{informierst}\\
    er/sie/es & \textbf{informiert}\\
    wir & \textbf{informieren}\\
    ihr & \textbf{informiert}\\
    sie/Sie & \textbf{informieren}\\
  \end{tabular}
  \hfill
  \begin{tabular}{rl}
    \multicolumn{2}{l}{\textbf{PRÄTERITUM} }\\
    ich & \textbf{informierte}\\
    du & \textbf{informiertest}\\
    er/sie/es & \textbf{informierte}\\
    wir & \textbf{informierten}\\
    ihr & \textbf{informiertet}\\
    sie/Sie & \textbf{informierten}\\
  \end{tabular}
  \hfill
  \begin{tabular}{rl}
     \multicolumn{2}{l}{\textbf{IMPERATIV} }\\
   & \textbf{informiere (du)!}\\
   & \textbf{informieren wir!}\\
   & \textbf{informiert (ihr)!}\\
   & \textbf{informieren Sie!}\\
   & \vspace{10pt}
  \end{tabular}

 \vspace{10pt}

 \begin{tabular}{rl}
     \multicolumn{2}{l}{\textbf{PERFEKT}}\\
   & haben + \textbf{informiert}\\
  \end{tabular}
\hfill
\begin{tabular}{rl}
    \multicolumn{2}{l}{\textbf{FUTURE I}}\\
   & werden + \textbf{informieren}\\
  \end{tabular}
\hfill
\begin{tabular}{rl}
    \multicolumn{2}{l}{\textbf{PLUSQUAMPERFEKT}}\\
   & hatten + \textbf{informiert}\\
  \end{tabular}

 \vspace{10pt}

\begin{tabular}{lll}
      \multicolumn{3}{l}{\textbf{MIT PRÄPOSITIONEN}}\\
& \textbf{informieren über} + Akk & (to inform about)\\
& \textbf{informieren von} + Dat & (to inform of / from)\\
& \textbf{informieren bei} + Dat & (to inform at / with)\\
\end{tabular}
  \hfill
\begin{tabular}{lll}
      \multicolumn{3}{l}{\textbf{TRENNBARE VERBEN}}\\
 & \textbf{ein$|$informieren} & (to inform in / register information)\\
\vspace{20pt}
  \end{tabular}


\vspace{-5pt}
\end{verbbox}

%installerient
\begin{verbbox}\addcontentsline{toc}{subsubsection}{installieren (to install)}
    {\Large \bf installieren} (to install) \hfill Type: \textbf{regular} \hfill Auxiliary Verb: \textbf{haben} 
    
     \vspace{5pt}

    \hrule
    
    \vspace{5pt}

  \begin{tabular}{rl}
     \multicolumn{2}{l}{\textbf{PRÄSENS} }\\
    ich & \textbf{installiere}\\
    du & \textbf{installierst}\\
    er/sie/es & \textbf{installiert}\\
    wir & \textbf{installieren}\\
    ihr & \textbf{installiert}\\
    sie/Sie & \textbf{installieren}\\
  \end{tabular}
  \hfill
  \begin{tabular}{rl}
    \multicolumn{2}{l}{\textbf{PRÄTERITUM} }\\
    ich & \textbf{installierte}\\
    du & \textbf{installiertest}\\
    er/sie/es & \textbf{installierte}\\
    wir & \textbf{installierten}\\
    ihr & \textbf{installiertet}\\
    sie/Sie & \textbf{installierten}\\
  \end{tabular}
  \hfill
  \begin{tabular}{rl}
     \multicolumn{2}{l}{\textbf{IMPERATIV} }\\
   & \textbf{installiere (du)!}\\
   & \textbf{installieren wir!}\\
   & \textbf{installiert (ihr)!}\\
   & \textbf{installieren Sie!}\\
   & \vspace{10pt}
  \end{tabular}

 \vspace{10pt}

 \begin{tabular}{rl}
     \multicolumn{2}{l}{\textbf{PERFEKT}}\\
   & haben + \textbf{installiert}\\
  \end{tabular}
\hfill
\begin{tabular}{rl}
    \multicolumn{2}{l}{\textbf{FUTURE I}}\\
   & werden + \textbf{installieren}\\
  \end{tabular}
\hfill
\begin{tabular}{rl}
    \multicolumn{2}{l}{\textbf{PLUSQUAMPERFEKT}}\\
   & hatten + \textbf{installiert}\\
  \end{tabular}

 \vspace{10pt}

\begin{tabular}{lll}
      \multicolumn{3}{l}{\textbf{MIT PRÄPOSITIONEN}}\\
& \textbf{installieren auf} + Dat & (to install on)\\
& \textbf{installieren in} + Akk & (to install in)\\
\end{tabular}
  \hfill
\begin{tabular}{lll}
      \multicolumn{3}{l}{\textbf{TRENNBARE VERBEN}}\\
 & \textbf{ein$|$installieren} & (to install into)\\
\vspace{10pt}
  \end{tabular}

\vspace{-5pt}
\end{verbbox}

\subsection{J Verbs}

\subsection{K Verbs}

%kaufen (A1) – to buy
\begin{verbbox-acc}\addcontentsline{toc}{subsubsection}{kaufen (to buy)}
    {\Large \bf kaufen} (to buy) \hfill Type: \textbf{regular, + Akkusativ} \hfill Auxiliary Verb: \textbf{haben} 
    
     \vspace{5pt}

    \hrule
    
    \vspace{5pt}

  \begin{tabular}{rl}
     \multicolumn{2}{l}{\textbf{PRÄSENS} }\\
    ich & \textbf{kaufe}\\
    du & \textbf{kaufst}\\
    er/sie/es & \textbf{kauft}\\
    wir & \textbf{kaufen}\\
    ihr & \textbf{kauft}\\
    sie/Sie & \textbf{kaufen}\\
  \end{tabular}
  \hfill
     \begin{tabular}{rl}
    \multicolumn{2}{l}{\textbf{PRÄTERITUM} }\\
    ich & \textbf{kaufte}\\
    du & \textbf{kauftest}\\
    er/sie/es & \textbf{kaufte}\\
    wir & \textbf{kauften}\\
    ihr & \textbf{kauftet}\\
    sie/Sie & \textbf{kauften}\\
  \end{tabular}
  \hfill
  \begin{tabular}{rl}
     \multicolumn{2}{l}{\textbf{IMPERATIV} }\\
   & \textbf{kauf(e) (du)!}\\
   & \textbf{gkaufen wir!}\\
   & \textbf{kauft (ihr)!}\\
   & \textbf{kaufen Sie!}\\
   & \vspace{10pt}
  \end{tabular}

    \vspace{10pt}

 \begin{tabular}{rl}
     \multicolumn{2}{l}{\textbf{PERFEKT}}\\
   & haben + \textbf{gekauft}\\
  \end{tabular}
\hfill
   \begin{tabular}{rl}
   
    \multicolumn{2}{l}{\textbf{FUTURE I} }\\
   & werden + \textbf{kaufen}\\
 
  \end{tabular}
\hfill
   \begin{tabular}{rl}
      \multicolumn{2}{l}{\textbf{PLUSQUAMPERFEKT}}\\
   & hatten + \textbf{gekauft}\\
  \end{tabular}
  
  \vspace{-5pt}


\end{verbbox-acc}

%kennen (A1) – to know (person/place)
\begin{verbbox-acc}\addcontentsline{toc}{subsubsection}{kennen (to know (person/place))}
    {\Large \bf kennen} (to know (person/place)) \hfill Type: \textbf{regular, + Akkusativ} \hfill Auxiliary Verb: \textbf{haben} 
    
     \vspace{5pt}

    \hrule
    
    \vspace{5pt}

  \begin{tabular}{rl}
     \multicolumn{2}{l}{\textbf{PRÄSENS} }\\
    ich & \textbf{kenne}\\
    du & \textbf{kennst}\\
    er/sie/es & \textbf{kennt}\\
    wir & \textbf{kennen}\\
    ihr & \textbf{kennt}\\
    sie/Sie & \textbf{kennen}\\
  \end{tabular}
  \hfill
     \begin{tabular}{rl}
    \multicolumn{2}{l}{\textbf{PRÄTERITUM} }\\
    ich & \textbf{kannte}\\
    du & \textbf{kanntest}\\
    er/sie/es & \textbf{kannte}\\
    wir & \textbf{kannten}\\
    ihr & \textbf{kanntet}\\
    sie/Sie & \textbf{kannten}\\
  \end{tabular}
  \hfill
  \begin{tabular}{rl}
     \multicolumn{2}{l}{\textbf{IMPERATIV} }\\
   & \textbf{kenn(e) (du)!}\\
   & \textbf{kennen wir!}\\
   & \textbf{kennt (ihr)!}\\
   & \textbf{kennen Sie!}\\
   & \vspace{10pt}
  \end{tabular}

    \vspace{10pt}

 \begin{tabular}{rl}
     \multicolumn{2}{l}{\textbf{PERFEKT}}\\
  &  haben + \textbf{gekannt}\\
  \end{tabular}
\hfill
   \begin{tabular}{rl}
   
    \multicolumn{2}{l}{\textbf{FUTURE I} }\\
  &  werden + \textbf{kennen}\\
 
  \end{tabular}
\hfill
   \begin{tabular}{rl}
      \multicolumn{2}{l}{\textbf{PLUSQUAMPERFEKT}}\\
   & hatten + \textbf{gekannt}\\
  \end{tabular}
  
  \vspace{-5pt}


\end{verbbox-acc}

%klären (B1) – to clarify
\begin{verbbox-acc}\addcontentsline{toc}{subsubsection}{klären (to clarify)}
    {\Large \bf klären} (to clarify) \hfill Type: \textbf{regular, + Akkusativ} \hfill Auxiliary Verb: \textbf{haben} 
    
     \vspace{5pt}

    \hrule
    
    \vspace{5pt}

  \begin{tabular}{rl}
     \multicolumn{2}{l}{\textbf{PRÄSENS} }\\
    ich & \textbf{kläre}\\
    du & \textbf{klärst}\\
    er/sie/es & \textbf{klärt}\\
    wir & \textbf{klären}\\
    ihr & \textbf{klärt}\\
    sie/Sie & \textbf{klären}\\
  \end{tabular}
  \hfill
     \begin{tabular}{rl}
    \multicolumn{2}{l}{\textbf{PRÄTERITUM} }\\
    ich & \textbf{klärte}\\
    du & \textbf{klärtest}\\
    er/sie/es & \textbf{klärte}\\
    wir & \textbf{klärten}\\
    ihr & \textbf{klärtet}\\
    sie/Sie & \textbf{klärten}\\
  \end{tabular}
  \hfill
  \begin{tabular}{rl}
     \multicolumn{2}{l}{\textbf{IMPERATIV} }\\
   & \textbf{klär(e) (du)!}\\
   & \textbf{klären wir!}\\
   & \textbf{klärt (ihr)!}\\
   & \textbf{klären Sie!}\\
   & \vspace{10pt}
  \end{tabular}

    \vspace{10pt}

 \begin{tabular}{rl}
     \multicolumn{2}{l}{\textbf{PERFEKT}}\\
   & haben + \textbf{geklärt}\\
  \end{tabular}
\hfill
   \begin{tabular}{rl}
   
    \multicolumn{2}{l}{\textbf{FUTURE I} }\\
   & werden + \textbf{klären}\\
 
  \end{tabular}
\hfill
   \begin{tabular}{rl}
      \multicolumn{2}{l}{\textbf{PLUSQUAMPERFEKT}}\\
   & hatten + \textbf{geklärt}\\
  \end{tabular}
  
  \vspace{-5pt}


\end{verbbox-acc}

%kleiden (to dress / to clothe)
\begin{verbbox}\addcontentsline{toc}{subsubsection}{kleiden (to dress / to clothe)}
    {\Large \bf kleiden} (to dress / to clothe) \hfill Type: \textbf{regular} \hfill Auxiliary Verb: \textbf{haben} 
    
     \vspace{5pt}

    \hrule
    
    \vspace{5pt}

  \begin{tabular}{rl}
     \multicolumn{2}{l}{\textbf{PRÄSENS} }\\
    ich & \textbf{kleide}\\
    du & \textbf{kleidest}\\
    er/sie/es & \textbf{kleidet}\\
    wir & \textbf{kleiden}\\
    ihr & \textbf{kleidet}\\
    sie/Sie & \textbf{kleiden}\\
  \end{tabular}
  \hfill
     \begin{tabular}{rl}
    \multicolumn{2}{l}{\textbf{PRÄTERITUM} }\\
    ich & \textbf{kleidete}\\
    du & \textbf{kleidetest}\\
    er/sie/es & \textbf{kleidete}\\
    wir & \textbf{kleideten}\\
    ihr & \textbf{kleidetet}\\
    sie/Sie & \textbf{kleideten}\\
  \end{tabular}
  \hfill
  \begin{tabular}{rl}
     \multicolumn{2}{l}{\textbf{IMPERATIV} }\\
   & \textbf{kleide (du)!}\\
   & \textbf{kleiden wir!}\\
   & \textbf{kleidet (ihr)!}\\
   & \textbf{kleiden Sie!}\\
   & \vspace{10pt}
  \end{tabular}

    \vspace{10pt}

 \begin{tabular}{rl}
     \multicolumn{2}{l}{\textbf{PERFEKT}}\\
   & haben + \textbf{gekleidet}\\
  \end{tabular}
\hfill
   \begin{tabular}{rl}
    \multicolumn{2}{l}{\textbf{FUTURE I} }\\
   & werde + \textbf{kleiden}\\
  \end{tabular}
\hfill
   \begin{tabular}{rl}
      \multicolumn{2}{l}{\textbf{PLUSQUAMPERFEKT}}\\
   & hatten + \textbf{gekleidet}\\
  \end{tabular}

   \vspace{10pt}

   \begin{tabular}{lll}
      \multicolumn{3}{l}{\textbf{MIT PRÄPOSITIONEN}}\\
& \textbf{sich kleiden in} + Akk & (to dress in)\\
& \textbf{kleiden mit} + Dat & (to clothe with)\\
\end{tabular}
  \hfill
   \begin{tabular}{lll}
      \multicolumn{3}{l}{\textbf{TRENNBARE VERBEN}}\\
 & \textbf{an$|$kleiden} & (to dress someone)\\
 & \textbf{ver$|$kleiden} & (to disguise / dress up)\\
\vspace{10pt}
  \end{tabular}

  \vspace{-5pt}

\end{verbbox}

%klettern (to climb)
\begin{verbbox}\addcontentsline{toc}{subsubsection}{klettern (to climb)}
    {\Large \bf klettern} (to climb) \hfill Type: \textbf{regular} \hfill Auxiliary Verb: \textbf{sein / haben} 
    
     \vspace{5pt}

    \hrule
    
    \vspace{5pt}

  \begin{tabular}{rl}
     \multicolumn{2}{l}{\textbf{PRÄSENS} }\\
    ich & \textbf{klettere}\\
    du & \textbf{kletterst}\\
    er/sie/es & \textbf{klettert}\\
    wir & \textbf{klettern}\\
    ihr & \textbf{klettert}\\
    sie/Sie & \textbf{klettern}\\
  \end{tabular}
  \hfill
     \begin{tabular}{rl}
    \multicolumn{2}{l}{\textbf{PRÄTERITUM} }\\
    ich & \textbf{kletterte}\\
    du & \textbf{klettertest}\\
    er/sie/es & \textbf{kletterte}\\
    wir & \textbf{kletterten}\\
    ihr & \textbf{klettertet}\\
    sie/Sie & \textbf{kletterten}\\
  \end{tabular}
  \hfill
  \begin{tabular}{rl}
     \multicolumn{2}{l}{\textbf{IMPERATIV} }\\
   & \textbf{klettere (du)!}\\
   & \textbf{klettern wir!}\\
   & \textbf{klettert (ihr)!}\\
   & \textbf{klettern Sie!}\\
   & \vspace{10pt}
  \end{tabular}

    \vspace{10pt}

 \begin{tabular}{rl}
     \multicolumn{2}{l}{\textbf{PERFEKT}}\\
   & sein + \textbf{geklettert}\\
  \end{tabular}
\hfill
   \begin{tabular}{rl}
    \multicolumn{2}{l}{\textbf{FUTURE I} }\\
   & werde + \textbf{klettern}\\
  \end{tabular}
\hfill
   \begin{tabular}{rl}
      \multicolumn{2}{l}{\textbf{PLUSQUAMPERFEKT}}\\
   & waren + \textbf{geklettert}\\
  \end{tabular}

   \vspace{10pt}

   \begin{tabular}{lll}
      \multicolumn{3}{l}{\textbf{MIT PRÄPOSITIONEN}}\\
& \textbf{klettern auf} + Akk & (to climb onto)\\
& \textbf{klettern über} + Akk & (to climb over)\\
& \textbf{klettern in} + Akk & (to climb into)\\
\end{tabular}
  \hfill
   \begin{tabular}{lll}
      \multicolumn{3}{l}{\textbf{TRENNBARE VERBEN}}\\
 & \textbf{hinauf$|$klettern} & (to climb up)\\
 & \textbf{herunter$|$klettern} & (to climb down)\\
\vspace{10pt}
  \end{tabular}

  \vspace{-5pt}

\end{verbbox}

%klicken (to click)
\begin{verbbox}\addcontentsline{toc}{subsubsection}{klicken (to click)}
{\Large \bf klicken} (to click) \hfill Type: \textbf{regular} \hfill Auxiliary Verb: \textbf{haben} 

\vspace{5pt}
\hrule
\vspace{5pt}

\begin{tabular}{rl}
\multicolumn{2}{l}{\textbf{PRÄSENS}}\\
ich & \textbf{klicke}\\
du & \textbf{klickst}\\
er/sie/es & \textbf{klickt}\\
wir & \textbf{klicken}\\
ihr & \textbf{klickt}\\
sie/Sie & \textbf{klicken}\\
\end{tabular}
\hfill
\begin{tabular}{rl}
\multicolumn{2}{l}{\textbf{PRÄTERITUM}}\\
ich & \textbf{klickte}\\
du & \textbf{klicktest}\\
er/sie/es & \textbf{klickte}\\
wir & \textbf{klickten}\\
ihr & \textbf{klicktet}\\
sie/Sie & \textbf{klickten}\\
\end{tabular}
\hfill
\begin{tabular}{rl}
\multicolumn{2}{l}{\textbf{IMPERATIV}}\\
& \textbf{klicke (du)!}\\
& \textbf{klicken wir!}\\
& \textbf{klickt (ihr)!}\\
& \textbf{klicken Sie!}\\
\end{tabular}

\vspace{10pt}

\begin{tabular}{rl}
\multicolumn{2}{l}{\textbf{PERFEKT}}\\
& haben + \textbf{geklickt}\\
\end{tabular}
\hfill
\begin{tabular}{rl}
\multicolumn{2}{l}{\textbf{FUTURE I}}\\
& werden + \textbf{klicken}\\
\end{tabular}
\hfill
\begin{tabular}{rl}
\multicolumn{2}{l}{\textbf{PLUSQUAMPERFEKT}}\\
& hatten + \textbf{geklickt}\\
\end{tabular}

\vspace{10pt}

\begin{tabular}{lll}
\multicolumn{3}{l}{\textbf{MIT PRÄPOSITIONEN}}\\
& \textbf{klicken auf} + Akk & (to click on)\\
& \textbf{klicken in} + Akk & (to click into)\\
\end{tabular}
\hfill
\begin{tabular}{lll}
\multicolumn{3}{l}{\textbf{TRENNBARE VERBEN}}\\
& an$|$klicken & (to click on / activate)\\
\vspace{10pt}
\end{tabular}

\vspace{-5pt}
\end{verbbox}

%knien - to kneel
\begin{verbbox}\addcontentsline{toc}{subsubsection}{knien (to kneel)}
    {\Large \bf  knien} (to kneel) \hfill Type: \textbf{regular} \hfill Auxiliary Verb: \textbf{haben / sein} 
    
     \vspace{5pt}

    \hrule
    
    \vspace{5pt}

  \begin{tabular}{rl}
     \multicolumn{2}{l}{\textbf{PRÄSENS} }\\
    ich & \textbf{knie}\\
    du & \textbf{kniest}\\
    er/sie/es & \textbf{kniet}\\
    wir & \textbf{knien}\\
    ihr & \textbf{kniet}\\
    sie/Sie & \textbf{knien}\\
  \end{tabular}
  \hfill
     \begin{tabular}{rl}
    \multicolumn{2}{l}{\textbf{PRÄTERITUM} }\\
    ich & \textbf{kniete}\\
    du & \textbf{knietest}\\
    er/sie/es & \textbf{kniete}\\
    wir & \textbf{knieten}\\
    ihr & \textbf{knietet}\\
    sie/Sie & \textbf{knieten}\\
  \end{tabular}
  \hfill
  \begin{tabular}{rl}
     \multicolumn{2}{l}{\textbf{IMPERATIV} }\\
   & \textbf{knie (du)!}\\
   & \textbf{knien wir!}\\
   & \textbf{kniet (ihr)!}\\
   & \textbf{knien Sie!}\\
   & \vspace{10pt}
  \end{tabular}

    \vspace{10pt}

 \begin{tabular}{rl}
     \multicolumn{2}{l}{\textbf{PERFEKT}}\\
  &  haben / sein + \textbf{gekniet}\\
  \end{tabular}
\hfill
   \begin{tabular}{rl}
   
    \multicolumn{2}{l}{\textbf{FUTURE I} }\\
   & werden + \textbf{knien}\\
 
  \end{tabular}
\hfill
   \begin{tabular}{rl}
      \multicolumn{2}{l}{\textbf{PLUSQUAMPERFEKT}}\\
   & hatten / waren + \textbf{gekniet}\\
  \end{tabular}
  
  \vspace{10pt}



\begin{tabular}{lll}
\multicolumn{3}{l}{\textbf{MIT PRÄPOSITIONEN}}\\
& ---\\
\end{tabular}
\hfill
\begin{tabular}{lll}
\multicolumn{3}{l}{\textbf{TRENNBARE VERBEN}}\\
& ---\\
\end{tabular}

  \vspace{-5pt}

\end{verbbox}


%knüpfen (to tie / establish)
\begin{verbbox}\addcontentsline{toc}{subsubsection}{knüpfen (to tie / establish)}
{\Large \bf knüpfen} (to tie / establish) \hfill Type: \textbf{regular} \hfill Auxiliary Verb: \textbf{haben} 

\vspace{5pt}
\hrule
\vspace{5pt}

\begin{tabular}{rl}
\multicolumn{2}{l}{\textbf{PRÄSENS}}\\
ich & \textbf{knüpfe}\\
du & \textbf{knüpfst}\\
er/sie/es & \textbf{knüpft}\\
wir & \textbf{knüpfen}\\
ihr & \textbf{knüpft}\\
sie/Sie & \textbf{knüpfen}\\
\end{tabular}
\hfill
\begin{tabular}{rl}
\multicolumn{2}{l}{\textbf{PRÄTERITUM}}\\
ich & \textbf{knüpfte}\\
du & \textbf{knüpftest}\\
er/sie/es & \textbf{knüpfte}\\
wir & \textbf{knüpften}\\
ihr & \textbf{knüpftet}\\
sie/Sie & \textbf{knüpften}\\
\end{tabular}
\hfill
\begin{tabular}{rl}
\multicolumn{2}{l}{\textbf{IMPERATIV}}\\
& \textbf{knüpfe (du)!}\\
& \textbf{knüpfen wir!}\\
& \textbf{knüpft (ihr)!}\\
& \textbf{knüpfen Sie!}\\
\end{tabular}

\vspace{10pt}

\begin{tabular}{rl}
\multicolumn{2}{l}{\textbf{PERFEKT}}\\
& haben + \textbf{geknüpft}\\
\end{tabular}
\hfill
\begin{tabular}{rl}
\multicolumn{2}{l}{\textbf{FUTURE I}}\\
& werden + \textbf{knüpfen}\\
\end{tabular}
\hfill
\begin{tabular}{rl}
\multicolumn{2}{l}{\textbf{PLUSQUAMPERFEKT}}\\
& hatten + \textbf{geknüpft}\\
\end{tabular}

\vspace{10pt}

\begin{tabular}{lll}
\multicolumn{3}{l}{\textbf{MIT PRÄPOSITIONEN}}\\
& \textbf{knüpfen an} + Akk & (to link to)\\
\end{tabular}
\hfill
\begin{tabular}{lll}
\multicolumn{3}{l}{\textbf{TRENNBARE VERBEN}}\\
& \textbf{an$|$knüpfen} & (to connect to)\\
\end{tabular}

\vspace{-5pt}
\end{verbbox}

%kochen (A1) – to cook
\begin{verbbox-acc}\addcontentsline{toc}{subsubsection}{kochen (to cook)}
    {\Large \bf kochen} (to cook) \hfill Type: \textbf{regular, + Akkusativ} \hfill Auxiliary Verb: \textbf{haben} 
    
     \vspace{5pt}

    \hrule
    
    \vspace{5pt}

  \begin{tabular}{rl}
     \multicolumn{2}{l}{\textbf{PRÄSENS} }\\
    ich & \textbf{koche}\\
    du & \textbf{kochst}\\
    er/sie/es & \textbf{kocht}\\
    wir & \textbf{kochen}\\
    ihr & \textbf{kocht}\\
    sie/Sie & \textbf{kochen}\\
  \end{tabular}
  \hfill
     \begin{tabular}{rl}
    \multicolumn{2}{l}{\textbf{PRÄTERITUM} }\\
    ich & \textbf{kochte}\\
    du & \textbf{kochtest}\\
    er/sie/es & \textbf{kochte}\\
    wir & \textbf{kochten}\\
    ihr & \textbf{kochtet}\\
    sie/Sie & \textbf{kochten}\\
  \end{tabular}
  \hfill
  \begin{tabular}{rl}
     \multicolumn{2}{l}{\textbf{IMPERATIV} }\\
   & \textbf{koch(e) (du)!}\\
   & \textbf{kochen wir!}\\
   & \textbf{kocht (ihr)!}\\
   & \textbf{kochen Sie!}\\
   & \vspace{10pt}
  \end{tabular}

    \vspace{10pt}

 \begin{tabular}{rl}
     \multicolumn{2}{l}{\textbf{PERFEKT}}\\
   & haben + \textbf{gekocht}\\
  \end{tabular}
\hfill
   \begin{tabular}{rl}
   
    \multicolumn{2}{l}{\textbf{FUTURE I} }\\
   & werden + \textbf{kochen}\\
 
  \end{tabular}
\hfill
   \begin{tabular}{rl}
      \multicolumn{2}{l}{\textbf{PLUSQUAMPERFEKT}}\\
    & hatten + \textbf{gekocht}\\
  \end{tabular}
  
  \vspace{-5pt}


\end{verbbox-acc}
%kommen (A1) – to come
\begin{verbbox}\addcontentsline{toc}{subsubsection}{kommen (to come)}
    {\Large \bf kommen} (to come) \hfill Type: \textbf{irregular} \hfill Auxiliary Verb: \textbf{sein} 
    
     \vspace{5pt}

    \hrule
    
    \vspace{5pt}

  \begin{tabular}{rl}
     \multicolumn{2}{l}{\textbf{PRÄSENS} }\\
    ich & \textbf{komme}\\
    du & \textbf{kommst}\\
    er/sie/es & \textbf{kommt}\\
    wir & \textbf{kommen}\\
    ihr & \textbf{kommt}\\
    sie/Sie & \textbf{kommen}\\
  \end{tabular}
  \hfill
     \begin{tabular}{rl}
    \multicolumn{2}{l}{\textbf{PRÄTERITUM} }\\
    ich & \textbf{kam}\\
    du & \textbf{kamst}\\
    er/sie/es & \textbf{kam}\\
    wir & \textbf{kamen}\\
    ihr & \textbf{kamt}\\
    sie/Sie & \textbf{kamen}\\
  \end{tabular}
  \hfill
  \begin{tabular}{rl}
     \multicolumn{2}{l}{\textbf{IMPERATIV} }\\
   & \textbf{komm(e) (du)!}\\
   & \textbf{kommen wir!}\\
   & \textbf{kommt (ihr)!}\\
   & \textbf{kommen Sie!}\\
   & \vspace{10pt}
  \end{tabular}

    \vspace{10pt}

 \begin{tabular}{rl}
     \multicolumn{2}{l}{\textbf{PERFEKT}}\\
   & sein + \textbf{gekommen}\\
  \end{tabular}
\hfill
   \begin{tabular}{rl}
   
    \multicolumn{2}{l}{\textbf{FUTURE I} }\\
   & werden + \textbf{kommen}\\
 
  \end{tabular}
\hfill
   \begin{tabular}{rl}
      \multicolumn{2}{l}{\textbf{PLUSQUAMPERFEKT}}\\
   & waren + \textbf{gekommen}\\
  \end{tabular}

  \vspace{10pt}

   \begin{tabular}{lll}
      \multicolumn{3}{l}{\textbf{MIT PRÄPOSITIONEN}}\\
& \textbf{kommen aus} + Dat & (to come from)\\
& \textbf{kommen von} + Dat & (origin / source of)\\
&\textbf{kommen zu} + Dat & (to approach to)\\
&\textbf{kommen nach} + Dat & (come to a city / country)\\
&\textbf{kommen in} + Akk & (to come into)\\
\vspace{50pt}
  \end{tabular}
  \hfill
   \begin{tabular}{lll}
      \multicolumn{3}{l}{\textbf{TRENNBARE VERBEN}}\\
 & \textbf{an$|$kommen} & (to arrive)\\
 & \textbf{mit$|$kommen} & (to come along)\\
     & \textbf{auf$|$kommen} & (to arise / come up)\\
     & \textbf{vor$|$kommen} & (to occur)\\
     & \textbf{nach$|$kommen} & (to follow / comply)\\
     & \textbf{zu$|$kommen} & (to approach)\\
      & \textbf{herein$|$kommen} & (to come in)\\
     & \textbf{hinein$|$kommen} & (to get inside / enter)\\
     & \textbf{dazu$|$kommen} & (to join / be added)\\
      & \textbf{fort$|$kommen} & (to make progress)
  \end{tabular}

  \vspace{-5pt}

  \end{verbbox}

%kommunizieren (to communicate)
\begin{verbbox}\addcontentsline{toc}{subsubsection}{kommunizieren (to communicate)}
{\Large \bf kommunizieren} (to communicate) \hfill Type: \textbf{regular} \hfill Auxiliary Verb: \textbf{haben} 

\vspace{5pt}
\hrule
\vspace{5pt}

\begin{tabular}{rl}
\multicolumn{2}{l}{\textbf{PRÄSENS}}\\
ich & \textbf{kommuniziere}\\
du & \textbf{kommunizierst}\\
er/sie/es & \textbf{kommuniziert}\\
wir & \textbf{kommunizieren}\\
ihr & \textbf{kommuniziert}\\
sie/Sie & \textbf{kommunizieren}\\
\end{tabular}
\hfill
\begin{tabular}{rl}
\multicolumn{2}{l}{\textbf{PRÄTERITUM}}\\
ich & \textbf{kommunizierte}\\
du & \textbf{kommuniziertest}\\
er/sie/es & \textbf{kommunizierte}\\
wir & \textbf{kommunizierten}\\
ihr & \textbf{kommuniziertet}\\
sie/Sie & \textbf{kommunizierten}\\
\end{tabular}
\hfill
\begin{tabular}{rl}
\multicolumn{2}{l}{\textbf{IMPERATIV}}\\
& \textbf{kommuniziere (du)!}\\
& \textbf{kommunizieren wir!}\\
& \textbf{kommuniziert (ihr)!}\\
& \textbf{kommunizieren Sie!}\\
\end{tabular}

\vspace{10pt}

\begin{tabular}{rl}
\multicolumn{2}{l}{\textbf{PERFEKT}}\\
& haben + \textbf{kommuniziert}\\
\end{tabular}
\hfill
\begin{tabular}{rl}
\multicolumn{2}{l}{\textbf{FUTURE I}}\\
& werden + \textbf{kommunizieren}\\
\end{tabular}
\hfill
\begin{tabular}{rl}
\multicolumn{2}{l}{\textbf{PLUSQUAMPERFEKT}}\\
& hatten + \textbf{kommuniziert}\\
\end{tabular}

\vspace{10pt}

\begin{tabular}{lll}
\multicolumn{3}{l}{\textbf{MIT PRÄPOSITIONEN}}\\
& \textbf{kommunizieren mit} + Dat & (to communicate with)\\
& \textbf{kommunizieren über} + Akk & (to communicate about)\\
\end{tabular}
\hfill
\begin{tabular}{lll}
\multicolumn{3}{l}{\textbf{TRENNBARE VERBEN}}\\
& ---\\
\vspace{10pt}
\end{tabular}

\vspace{-5pt}
\end{verbbox}



%kosten (A1) – to cost
\begin{verbbox-acc}\addcontentsline{toc}{subsubsection}{kosten (to cost)}
    {\Large \bf kosten} (to cost) \hfill Type: \textbf{regular, + Akkusativ} \hfill Auxiliary Verb: \textbf{haben} 
    
     \vspace{5pt}

    \hrule
    
    \vspace{5pt}

  \begin{tabular}{rl}
     \multicolumn{2}{l}{\textbf{PRÄSENS} }\\
    ich & \textbf{-}\\
    du & \textbf{-}\\
    er/sie/es & \textbf{kostet}\\
    wir & \textbf{-}\\
    ihr & \textbf{-}\\
    sie/Sie & \textbf{kosten}\\
  \end{tabular}
  \hfill
     \begin{tabular}{rl}
    \multicolumn{2}{l}{\textbf{PRÄTERITUM} }\\
    ich & \textbf{-}\\
    du & \textbf{-}\\
    er/sie/es & \textbf{kostete}\\
    wir & \textbf{-}\\
    ihr & \textbf{-}\\
    sie/Sie & \textbf{kosteten}\\
  \end{tabular}
  \hfill
  \begin{tabular}{rl}
     \multicolumn{2}{l}{\textbf{IMPERATIV} }\\
   & \textbf{-}\\
   & \textbf{-}\\
   & \textbf{-}\\
   & \textbf{-}\\
   & \vspace{10pt}
  \end{tabular}

    \vspace{10pt}

 \begin{tabular}{rl}
     \multicolumn{2}{l}{\textbf{PERFEKT}}\\
   & haben + \textbf{gekostet}\\
  \end{tabular}
\hfill
   \begin{tabular}{rl}
   
    \multicolumn{2}{l}{\textbf{FUTURE I} }\\
   & werden + \textbf{kosten}\\
 
  \end{tabular}
\hfill
   \begin{tabular}{rl}
      \multicolumn{2}{l}{\textbf{PLUSQUAMPERFEKT}}\\
   & hatten + \textbf{gekostet}\\
  \end{tabular}
  
  \vspace{-5pt}


\end{verbbox-acc}

%kucken (B1) – to look
\begin{verbbox}\addcontentsline{toc}{subsubsection}{kucken (to look)}
    {\Large \bf kucken} (to look) \hfill Type: \textbf{regular} \hfill Auxiliary Verb: \textbf{haben} 
    
     \vspace{5pt}

    \hrule
    
    \vspace{5pt}

  \begin{tabular}{rl}
     \multicolumn{2}{l}{\textbf{PRÄSENS} }\\
    ich & \textbf{kucke}\\
    du & \textbf{kuckst}\\
    er/sie/es & \textbf{kuckt}\\
    wir & \textbf{kucken}\\
    ihr & \textbf{kuckt}\\
    sie/Sie & \textbf{kucken}\\
  \end{tabular}
  \hfill
     \begin{tabular}{rl}
    \multicolumn{2}{l}{\textbf{PRÄTERITUM} }\\
    ich & \textbf{kuckte}\\
    du & \textbf{kucktest}\\
    er/sie/es & \textbf{kuckte}\\
    wir & \textbf{kuckten}\\
    ihr & \textbf{kucktet}\\
    sie/Sie & \textbf{kuckten}\\
  \end{tabular}
  \hfill
  \begin{tabular}{rl}
     \multicolumn{2}{l}{\textbf{IMPERATIV} }\\
   & \textbf{kuck(e) du!}\\
   & \textbf{kucken wir!}\\
   & \textbf{kuckt (ihr)!}\\
   & \textbf{kucken Sie!}\\
   & \vspace{10pt}
  \end{tabular}

    \vspace{10pt}

 \begin{tabular}{rl}
     \multicolumn{2}{l}{\textbf{PERFEKT}}\\
   & haben + \textbf{gekuckt}\\
  \end{tabular}
\hfill
   \begin{tabular}{rl}
   
    \multicolumn{2}{l}{\textbf{FUTURE I} }\\
   & werden + \textbf{kucken}\\
 
  \end{tabular}
\hfill
   \begin{tabular}{rl}
      \multicolumn{2}{l}{\textbf{PLUSQUAMPERFEKT}}\\
   & hatten + \textbf{gekuckt}\\
  \end{tabular}
  
  \vspace{10pt}



\begin{tabular}{lll}
\multicolumn{3}{l}{\textbf{MIT PRÄPOSITIONEN}}\\
& ---\\
\end{tabular}
\hfill
\begin{tabular}{lll}
\multicolumn{3}{l}{\textbf{TRENNBARE VERBEN}}\\
& ---\\
\end{tabular}

  \vspace{-5pt}
\end{verbbox}

%kündigen (to quit / to terminate)
\begin{verbbox}\addcontentsline{toc}{subsubsection}{kündigen (to quit / to terminate)}
    {\Large \bf kündigen} (to quit / to terminate) \hfill Type: \textbf{regular} \hfill Auxiliary Verb: \textbf{haben} 
    
     \vspace{5pt}

    \hrule
    
    \vspace{5pt}

  \begin{tabular}{rl}
     \multicolumn{2}{l}{\textbf{PRÄSENS} }\\
    ich & \textbf{kündige}\\
    du & \textbf{kündigst}\\
    er/sie/es & \textbf{kündigt}\\
    wir & \textbf{kündigen}\\
    ihr & \textbf{kündigt}\\
    sie/Sie & \textbf{kündigen}\\
  \end{tabular}
  \hfill
     \begin{tabular}{rl}
    \multicolumn{2}{l}{\textbf{PRÄTERITUM} }\\
    ich & \textbf{kündigte}\\
    du & \textbf{kündigtest}\\
    er/sie/es & \textbf{kündigte}\\
    wir & \textbf{kündigten}\\
    ihr & \textbf{kündigtet}\\
    sie/Sie & \textbf{kündigten}\\
  \end{tabular}
  \hfill
  \begin{tabular}{rl}
     \multicolumn{2}{l}{\textbf{IMPERATIV} }\\
   & \textbf{kündige (du)!}\\
   & \textbf{kündigen wir!}\\
   & \textbf{kündigt (ihr)!}\\
   & \textbf{kündigen Sie!}\\
   & \vspace{10pt}
  \end{tabular}

    \vspace{10pt}

 \begin{tabular}{rl}
     \multicolumn{2}{l}{\textbf{PERFEKT}}\\
   & haben + \textbf{gekündigt}\\
  \end{tabular}
\hfill
   \begin{tabular}{rl}
    \multicolumn{2}{l}{\textbf{FUTURE I} }\\
   & werde + \textbf{kündigen}\\
  \end{tabular}
\hfill
   \begin{tabular}{rl}
      \multicolumn{2}{l}{\textbf{PLUSQUAMPERFEKT}}\\
   & hatten + \textbf{gekündigt}\\
  \end{tabular}

   \vspace{10pt}

   \begin{tabular}{lll}
      \multicolumn{3}{l}{\textbf{MIT PRÄPOSITIONEN}}\\
& \textbf{kündigen bei} + Dat & (to quit at / resign from)\\
& \textbf{kündigen wegen} + Gen & (to terminate because of)\\
\end{tabular}
  \hfill
   \begin{tabular}{lll}
      \multicolumn{3}{l}{\textbf{TRENNBARE VERBEN}}\\
 & \textbf{an$|$kündigen} & (to announce)\\
 & \textbf{auf$|$kündigen} & (to terminate / cancel)\\
\vspace{10pt}
  \end{tabular}

  \vspace{-5pt}

\end{verbbox}

\subsection{L Verbs}

%lachen (to laugh)
\begin{verbbox}\addcontentsline{toc}{subsubsection}{lachen (to laugh)}
{\Large \bf lachen} (to laugh) \hfill Type: \textbf{regular} \hfill Auxiliary Verb: \textbf{haben} 

\vspace{5pt}
\hrule
\vspace{5pt}

\begin{tabular}{rl}
\multicolumn{2}{l}{\textbf{PRÄSENS}}\\
ich & \textbf{lache}\\
du & \textbf{lachst}\\
er/sie/es & \textbf{lacht}\\
wir & \textbf{lachen}\\
ihr & \textbf{lacht}\\
sie/Sie & \textbf{lachen}\\
\end{tabular}
\hfill
\begin{tabular}{rl}
\multicolumn{2}{l}{\textbf{PRÄTERITUM}}\\
ich & \textbf{lachte}\\
du & \textbf{lachtest}\\
er/sie/es & \textbf{lachte}\\
wir & \textbf{lachten}\\
ihr & \textbf{lachtet}\\
sie/Sie & \textbf{lachten}\\
\end{tabular}
\hfill
\begin{tabular}{rl}
\multicolumn{2}{l}{\textbf{IMPERATIV}}\\
& \textbf{lach(e) (du)!}\\
& \textbf{lachen wir!}\\
& \textbf{lacht (ihr)!}\\
& \textbf{lachen Sie!}\\
\end{tabular}

\vspace{10pt}

\begin{tabular}{rl}
\multicolumn{2}{l}{\textbf{PERFEKT}}\\
& haben + \textbf{gelacht}\\
\end{tabular}
\hfill
\begin{tabular}{rl}
\multicolumn{2}{l}{\textbf{FUTURE I}}\\
& werden + \textbf{lachen}\\
\end{tabular}
\hfill
\begin{tabular}{rl}
\multicolumn{2}{l}{\textbf{PLUSQUAMPERFEKT}}\\
& hatten + \textbf{gelacht}\\
\end{tabular}

\vspace{10pt}

\begin{tabular}{lll}
\multicolumn{3}{l}{\textbf{MIT PRÄPOSITIONEN}}\\
& \textbf{lachen über} + Akk & (to laugh about)\\
\end{tabular}
\hfill
\begin{tabular}{lll}
\multicolumn{3}{l}{\textbf{TRENNBARE VERBEN}}\\
& aus$|$lachen & (to laugh at / mock)\\
\end{tabular}
\vspace{-5pt}
\end{verbbox}

%lächeln (to smile)
\begin{verbbox}\addcontentsline{toc}{subsubsection}{lächeln (to smile)}
{\Large \bf lächeln} (to smile) \hfill Type: \textbf{regular} \hfill Auxiliary Verb: \textbf{haben}

\vspace{5pt}
\hrule
\vspace{5pt}

\begin{tabular}{rl}
\multicolumn{2}{l}{\textbf{PRÄSENS}}\\
ich & \textbf{lächele}\\
du & \textbf{lächelst}\\
er/sie/es & \textbf{lächelt}\\
wir & \textbf{lächeln}\\
ihr & \textbf{lächelt}\\
sie/Sie & \textbf{lächeln}\\
\end{tabular}
\hfill
\begin{tabular}{rl}
\multicolumn{2}{l}{\textbf{PRÄTERITUM}}\\
ich & \textbf{lächelte}\\
du & \textbf{lächeltest}\\
er/sie/es & \textbf{lächelte}\\
wir & \textbf{lächelten}\\
ihr & \textbf{lächeltet}\\
sie/Sie & \textbf{lächelten}\\
\end{tabular}
\hfill
\begin{tabular}{rl}
\multicolumn{2}{l}{\textbf{IMPERATIV}}\\
& \textbf{lächle (du)!}\\
& \textbf{lächeln wir!}\\
& \textbf{lächelt (ihr)!}\\
& \textbf{lächeln Sie!}\\
\end{tabular}

\vspace{10pt}

\begin{tabular}{rl}
\multicolumn{2}{l}{\textbf{PERFEKT}}\\
& haben + \textbf{gelächelt}\\
\end{tabular}
\hfill
\begin{tabular}{rl}
\multicolumn{2}{l}{\textbf{FUTURE I}}\\
& werden + \textbf{lächeln}\\
\end{tabular}
\hfill
\begin{tabular}{rl}
\multicolumn{2}{l}{\textbf{PLUSQUAMPERFEKT}}\\
& hatten + \textbf{gelächelt}\\
\end{tabular}

\vspace{10pt}

\begin{tabular}{lll}
\multicolumn{3}{l}{\textbf{MIT PRÄPOSITIONEN}}\\
& \textbf{lächeln über} + Akk & (to smile about)\\
& \textbf{lächeln zu} + Dat & (to smile at / toward)\\
\end{tabular}
\hfill
\begin{tabular}{lll}
\multicolumn{3}{l}{\textbf{TRENNBARE VERBEN}}\\
& an$|$lächeln & (to smile at)\\
\vspace{10pt}
\end{tabular}

\vspace{-5pt}
\end{verbbox}

%laden (A2) – to load
\begin{verbbox-acc}\addcontentsline{toc}{subsubsection}{laden (to load / invite)}
    {\Large \bf laden} (to load / invite) \hfill Type: \textbf{irregular} \hfill Auxiliary Verb: \textbf{haben} 
    
     \vspace{5pt}

    \hrule
    
    \vspace{5pt}

  \begin{tabular}{rl}
     \multicolumn{2}{l}{\textbf{PRÄSENS} }\\
    ich & \textbf{lade}\\
    du & \textbf{lädst}\\
    er/sie/es & \textbf{lädt}\\
    wir & \textbf{laden}\\
    ihr & \textbf{ladet}\\
    sie/Sie & \textbf{laden}\\
  \end{tabular}
  \hfill
  \begin{tabular}{rl}
    \multicolumn{2}{l}{\textbf{PRÄTERITUM} }\\
    ich & \textbf{lud}\\
    du & \textbf{ludst}\\
    er/sie/es & \textbf{lud}\\
    wir & \textbf{luden}\\
    ihr & \textbf{ludet}\\
    sie/Sie & \textbf{luden}\\
  \end{tabular}
  \hfill
  \begin{tabular}{rl}
     \multicolumn{2}{l}{\textbf{IMPERATIV} }\\
   & \textbf{lad(e) (du)!}\\
   & \textbf{laden wir!}\\
   & \textbf{ladet (ihr)!}\\
   & \textbf{laden Sie!}\\
   & \vspace{10pt}
  \end{tabular}

 \vspace{10pt}

 \begin{tabular}{rl}
     \multicolumn{2}{l}{\textbf{PERFEKT}}\\
   & haben + \textbf{geladen}\\
  \end{tabular}
\hfill
\begin{tabular}{rl}
    \multicolumn{2}{l}{\textbf{FUTURE I}}\\
   & werden + \textbf{laden}\\
  \end{tabular}
\hfill
\begin{tabular}{rl}
    \multicolumn{2}{l}{\textbf{PLUSQUAMPERFEKT}}\\
   & hatten + \textbf{geladen}\\
  \end{tabular}

 \vspace{10pt}

\begin{tabular}{lll}
      \multicolumn{3}{l}{\textbf{MIT PRÄPOSITIONEN}}\\
& \textbf{laden auf} + Akk & (to load onto)\\
& \textbf{laden in} + Akk & (to load into)\\
& \textbf{laden zu} + Dat & (to invite to)\\
\end{tabular}
  \hfill
\begin{tabular}{lll}
      \multicolumn{3}{l}{\textbf{TRENNBARE VERBEN}}\\
 & \textbf{ein$|$laden} & (to invite)\\
 & \textbf{auf$|$laden} & (to load up)\\
\vspace{10pt}
  \end{tabular}
  
  \vspace{-5pt}

\end{verbbox-acc}

%landen - to land
\begin{verbbox}\addcontentsline{toc}{subsubsection}{landen (to land)}
    {\Large \bf landen} (to land) \hfill Type: \textbf{regular} \hfill Auxiliary Verb: \textbf{sein} 
    
     \vspace{5pt}

    \hrule
    
    \vspace{5pt}

  \begin{tabular}{rl}
     \multicolumn{2}{l}{\textbf{PRÄSENS} }\\
    ich & \textbf{lande}\\
    du & \textbf{landest}\\
    er/sie/es & \textbf{landet}\\
    wir & \textbf{landen}\\
    ihr & \textbf{landet}\\
    sie/Sie & \textbf{landen}\\
  \end{tabular}
  \hfill
     \begin{tabular}{rl}
    \multicolumn{2}{l}{\textbf{PRÄTERITUM} }\\
    ich & \textbf{landete}\\
    du & \textbf{landetest}\\
    er/sie/es & \textbf{landete}\\
    wir & \textbf{landeten}\\
    ihr & \textbf{landetet}\\
    sie/Sie & \textbf{landeten}\\
  \end{tabular}
  \hfill
  \begin{tabular}{rl}
     \multicolumn{2}{l}{\textbf{IMPERATIV} }\\
   & \textbf{land(e) (du)!}\\
   & \textbf{landen wir!}\\
   & \textbf{landet (ihr)!}\\
   & \textbf{landen Sie!}\\
   & \vspace{10pt}
  \end{tabular}

    \vspace{10pt}

 \begin{tabular}{rl}
     \multicolumn{2}{l}{\textbf{PERFEKT}}\\
   & sein + \textbf{gelandet}\\
  \end{tabular}
\hfill
   \begin{tabular}{rl}
   
    \multicolumn{2}{l}{\textbf{FUTURE I} }\\
   & werden + \textbf{landen}\\
 
  \end{tabular}
\hfill
   \begin{tabular}{rl}
      \multicolumn{2}{l}{\textbf{PLUSQUAMPERFEKT}}\\
   & waren + \textbf{gelandet}\\
  \end{tabular}

   \vspace{10pt}

   \begin{tabular}{lll}
      \multicolumn{3}{l}{\textbf{MIT PRÄPOSITIONEN}}\\
& \textbf{landen auf} + Dat & (to land on – surface / area)\\
& \textbf{landen in} + Dat & (to land in – country / city)\\
& \textbf{landen bei} + Dat & (to land near / at)\\
\vspace{10pt}
   \end{tabular}
  \hfill
   \begin{tabular}{lll}
      \multicolumn{3}{l}{\textbf{TRENNBARE VERBEN}}\\
 & \textbf{an$|$landen} & (to arrive / land at a place)\\
 & \textbf{auf$|$landen} & (to land on – surface)\\
 & \textbf{ein$|$landen} & (to land in / enter)\\
 & \textbf{zurück$|$landen} & (to land back / return)\\
  \end{tabular}

\vspace{-5pt}

\end{verbbox}

%laufen (A1) – to run
\begin{verbbox}\addcontentsline{toc}{subsubsection}{laufen (to run)}
    {\Large \bf laufen} (to run) \hfill Type: \textbf{irregular} \hfill Auxiliary Verb: \textbf{sein} 
    
     \vspace{5pt}

    \hrule
    
    \vspace{5pt}

  \begin{tabular}{rl}
     \multicolumn{2}{l}{\textbf{PRÄSENS} }\\
    ich & \textbf{laufe}\\
    du & \textbf{läufst}\\
    er/sie/es & \textbf{läuft}\\
    wir & \textbf{laufen}\\
    ihr & \textbf{lauft}\\
    sie/Sie & \textbf{laufen}\\
  \end{tabular}
  \hfill
     \begin{tabular}{rl}
    \multicolumn{2}{l}{\textbf{PRÄTERITUM} }\\
    ich & \textbf{lief}\\
    du & \textbf{liefst}\\
    er/sie/es & \textbf{lief}\\
    wir & \textbf{liefen}\\
    ihr & \textbf{lieft}\\
    sie/Sie & \textbf{liefen}\\
  \end{tabular}
  \hfill
  \begin{tabular}{rl}
     \multicolumn{2}{l}{\textbf{IMPERATIV} }\\
   & \textbf{lauf(e) (du)!}\\
   & \textbf{laufen wir!}\\
   & \textbf{lauft (ihr)!}\\
   & \textbf{laufen Sie!}\\
   & \vspace{10pt}
  \end{tabular}

    \vspace{10pt}

 \begin{tabular}{rl}
     \multicolumn{2}{l}{\textbf{PERFEKT}}\\
  &  sein + \textbf{gelaufen}\\
  \end{tabular}
\hfill
   \begin{tabular}{rl}
   
    \multicolumn{2}{l}{\textbf{FUTURE I} }\\
  &  werden + \textbf{laufen}\\
 
  \end{tabular}
\hfill
   \begin{tabular}{rl}
      \multicolumn{2}{l}{\textbf{PLUSQUAMPERFEKT}}\\
  &  waren + \textbf{gelaufen}\\
  \end{tabular}
  
  \vspace{10pt}



\begin{tabular}{lll}
\multicolumn{3}{l}{\textbf{MIT PRÄPOSITIONEN}}\\
& ---\\
\end{tabular}
\hfill
\begin{tabular}{lll}
\multicolumn{3}{l}{\textbf{TRENNBARE VERBEN}}\\
& ---\\
\end{tabular}
  
\vspace{-5pt}


\end{verbbox}
%leben (A1) – to live
\begin{verbbox}\addcontentsline{toc}{subsubsection}{leben (to live)}
    {\Large \bf leben} (to live) \hfill Type: \textbf{regular} \hfill Auxiliary Verb: \textbf{haben} 
    
     \vspace{5pt}

    \hrule
    
    \vspace{5pt}

  \begin{tabular}{rl}
     \multicolumn{2}{l}{\textbf{PRÄSENS} }\\
    ich & \textbf{lebe}\\
    du & \textbf{lebst}\\
    er/sie/es & \textbf{lebt}\\
    wir & \textbf{leben}\\
    ihr & \textbf{lebt}\\
    sie/Sie & \textbf{leben}\\
  \end{tabular}
  \hfill
     \begin{tabular}{rl}
    \multicolumn{2}{l}{\textbf{PRÄTERITUM} }\\
    ich & \textbf{lebte}\\
    du & \textbf{lebtest}\\
    er/sie/es & \textbf{lebte}\\
    wir & \textbf{lebten}\\
    ihr & \textbf{lebtet}\\
    sie/Sie & \textbf{lebten}\\
  \end{tabular}
  \hfill
  \begin{tabular}{rl}
     \multicolumn{2}{l}{\textbf{IMPERATIV} }\\
   & \textbf{leb(e) (du)!}\\
   & \textbf{leben wir!}\\
   & \textbf{lebt (ihr)!}\\
   & \textbf{leben Sie!}\\
   & \vspace{10pt}
  \end{tabular}

    \vspace{10pt}

 \begin{tabular}{rl}
     \multicolumn{2}{l}{\textbf{PERFEKT}}\\
   & haben + \textbf{gelebt}\\
  \end{tabular}
\hfill
   \begin{tabular}{rl}
   
    \multicolumn{2}{l}{\textbf{FUTURE I} }\\
   & werden + \textbf{leben}\\
 
  \end{tabular}
\hfill
   \begin{tabular}{rl}
      \multicolumn{2}{l}{\textbf{PLUSQUAMPERFEKT}}\\
  &  hatten + \textbf{gelebt}\\
  \end{tabular}
  
  \vspace{10pt}

\begin{tabular}{lll}
\multicolumn{3}{l}{\textbf{MIT PRÄPOSITIONEN}}\\
& ---\\
\end{tabular}
\hfill
\begin{tabular}{lll}
\multicolumn{3}{l}{\textbf{TRENNBARE VERBEN}}\\
& ---\\
\end{tabular}

  \vspace{-5pt}

\end{verbbox}

%legen (to lay / put)
\begin{verbbox}\addcontentsline{toc}{subsubsection}{legen (to lay / put)}
{\Large \bf legen} (to lay / put) \hfill Type: \textbf{regular} \hfill Auxiliary Verb: \textbf{haben}

\vspace{5pt}
\hrule
\vspace{5pt}

\begin{tabular}{rl}
\multicolumn{2}{l}{\textbf{PRÄSENS}}\\
ich & \textbf{lege}\\
du & \textbf{legst}\\
er/sie/es & \textbf{legt}\\
wir & \textbf{legen}\\
ihr & \textbf{legt}\\
sie/Sie & \textbf{legen}\\
\end{tabular}
\hfill
\begin{tabular}{rl}
\multicolumn{2}{l}{\textbf{PRÄTERITUM}}\\
ich & \textbf{legte}\\
du & \textbf{legtest}\\
er/sie/es & \textbf{legte}\\
wir & \textbf{legten}\\
ihr & \textbf{legtet}\\
sie/Sie & \textbf{legten}\\
\end{tabular}
\hfill
\begin{tabular}{rl}
\multicolumn{2}{l}{\textbf{IMPERATIV}}\\
& \textbf{lege (du)!}\\
& \textbf{legen wir!}\\
& \textbf{legt (ihr)!}\\
& \textbf{legen Sie!}\\
\end{tabular}

\vspace{10pt}

\begin{tabular}{rl}
\multicolumn{2}{l}{\textbf{PERFEKT}}\\
& haben + \textbf{gelegt}\\
\end{tabular}
\hfill
\begin{tabular}{rl}
\multicolumn{2}{l}{\textbf{FUTURE I}}\\
& werden + \textbf{legen}\\
\end{tabular}
\hfill
\begin{tabular}{rl}
\multicolumn{2}{l}{\textbf{PLUSQUAMPERFEKT}}\\
& hatten + \textbf{gelegt}\\
\end{tabular}

\vspace{10pt}

\begin{tabular}{lll}
\multicolumn{3}{l}{\textbf{MIT PRÄPOSITIONEN}}\\
& \textbf{legen auf} + Akk & (to lay on)\\
& \textbf{legen in} + Akk & (to put into)\\
& \textbf{legen unter} + Akk & (to lay under)\\
\end{tabular}
\hfill
\begin{tabular}{lll}
\multicolumn{3}{l}{\textbf{TRENNBARE VERBEN}}\\
& \textbf{hin$|$legen} & (to lay down)\\
& \textbf{auf$|$legen} & (to put on / overlay)\\
\end{tabular}
\vspace{-5pt}
\end{verbbox}

%lehnen (to lean)
\begin{verbbox}\addcontentsline{toc}{subsubsection}{lehnen (to lean)}
{\Large \bf lehnen} (to lean) \hfill Type: \textbf{regular} \hfill Auxiliary Verb: \textbf{haben / sein reflexive}

\vspace{5pt}
\hrule
\vspace{5pt}

\begin{tabular}{rl}
\multicolumn{2}{l}{\textbf{PRÄSENS}}\\
ich & \textbf{lehne}\\
du & \textbf{lehnst}\\
er/sie/es & \textbf{lehnt}\\
wir & \textbf{lehnen}\\
ihr & \textbf{lehnt}\\
sie/Sie & \textbf{lehnen}\\
\end{tabular}
\hfill
\begin{tabular}{rl}
\multicolumn{2}{l}{\textbf{PRÄTERITUM}}\\
ich & \textbf{lehnte}\\
du & \textbf{lehntest}\\
er/sie/es & \textbf{lehnte}\\
wir & \textbf{lehnten}\\
ihr & \textbf{lehntet}\\
sie/Sie & \textbf{lehnten}\\
\end{tabular}
\hfill
\begin{tabular}{rl}
\multicolumn{2}{l}{\textbf{IMPERATIV}}\\
& \textbf{lehne (du)!}\\
& \textbf{lehnen wir!}\\
& \textbf{lehnt (ihr)!}\\
& \textbf{lehnen Sie!}\\
\end{tabular}

\vspace{10pt}

\begin{tabular}{rl}
\multicolumn{2}{l}{\textbf{PERFEKT}}\\
& haben + \textbf{gelehnt}\\
\end{tabular}
\hfill
\begin{tabular}{rl}
\multicolumn{2}{l}{\textbf{FUTURE I}}\\
& werden + \textbf{lehnen}\\
\end{tabular}
\hfill
\begin{tabular}{rl}
\multicolumn{2}{l}{\textbf{PLUSQUAMPERFEKT}}\\
& hatten + \textbf{gelehnt}\\
\end{tabular}

\vspace{10pt}

\begin{tabular}{lll}
\multicolumn{3}{l}{\textbf{MIT PRÄPOSITIONEN}}\\
& \textbf{lehnen an} + Dat & (to lean against)\\
& \textbf{sich lehnen gegen} + Akk & (to lean against)\\
\end{tabular}
\hfill
\begin{tabular}{lll}
\multicolumn{3}{l}{\textbf{TRENNBARE VERBEN}}\\
& ---\\
\end{tabular}
\vspace{-5pt}
\end{verbbox}

%leihen (A2) – to borrow something
\begin{verbbox-acc}\addcontentsline{toc}{subsubsection}{leihen (to borrow)}
    {\Large \bf leihen} (to borrow) \hfill Type: \textbf{irregular, + Akkusativ} \hfill Auxiliary Verb: \textbf{haben} 
    
     \vspace{5pt}

    \hrule
    
    \vspace{5pt}

  \begin{tabular}{rl}
     \multicolumn{2}{l}{\textbf{PRÄSENS} }\\
    ich & \textbf{leihe}\\
    du & \textbf{leihst}\\
    er/sie/es & \textbf{leiht}\\
    wir & \textbf{leihen}\\
    ihr & \textbf{leiht}\\
    sie/Sie & \textbf{leihen}\\
  \end{tabular}
  \hfill
     \begin{tabular}{rl}
    \multicolumn{2}{l}{\textbf{PRÄTERITUM} }\\
    ich & \textbf{lieh}\\
    du & \textbf{liehst}\\
    er/sie/es & \textbf{lieh}\\
    wir & \textbf{liehen}\\
    ihr & \textbf{lieht}\\
    sie/Sie & \textbf{liehen}\\
  \end{tabular}
  \hfill
  \begin{tabular}{rl}
     \multicolumn{2}{l}{\textbf{IMPERATIV} }\\
   & \textbf{leih(e) (du)!}\\
   & \textbf{leihen wir!}\\
   & \textbf{leiht (ihr)!}\\
   & \textbf{leihen Sie!}\\
   & \vspace{10pt}
  \end{tabular}

    \vspace{10pt}

 \begin{tabular}{rl}
     \multicolumn{2}{l}{\textbf{PERFEKT}}\\
   & sein + \textbf{geliehen}\\
  \end{tabular}
\hfill
   \begin{tabular}{rl}
   
    \multicolumn{2}{l}{\textbf{FUTURE I} }\\
   & werden + \textbf{leihen}\\
 
  \end{tabular}
\hfill
   \begin{tabular}{rl}
      \multicolumn{2}{l}{\textbf{PLUSQUAMPERFEKT}}\\
   & waren + \textbf{geliehen}\\
  \end{tabular}
  


 \vspace{10pt}

   \begin{tabular}{lll}
      \multicolumn{3}{l}{\textbf{MIT PRÄPOSITIONEN}}\\
& \textbf{leihen von} + Dat & (to borrow from)\\
& \textbf{leihen an} + Dat & (to lend to)\\
  \end{tabular}
  \hfill
   \begin{tabular}{lll}
      \multicolumn{3}{l}{\textbf{TRENNBARE VERBEN}}\\
 & \textbf{aus$|$leihen} & (to lend out (library / shop))\\
  \end{tabular}


  \vspace{-5pt}

\end{verbbox-acc}
%lernen (A1) – to learn
\begin{verbbox-acc}\addcontentsline{toc}{subsubsection}{lernen (to learn)}
    {\Large \bf lernen} (to learn) \hfill Type: \textbf{regular, + Akkusativ} \hfill Auxiliary Verb: \textbf{haben} 
    
     \vspace{5pt}

    \hrule
    
    \vspace{5pt}

  \begin{tabular}{rl}
     \multicolumn{2}{l}{\textbf{PRÄSENS} }\\
    ich & \textbf{lerne}\\
    du & \textbf{lernst}\\
    er/sie/es & \textbf{lernt}\\
    wir & \textbf{lernen}\\
    ihr & \textbf{lernt}\\
    sie/Sie & \textbf{lernen}\\
  \end{tabular}
  \hfill
     \begin{tabular}{rl}
    \multicolumn{2}{l}{\textbf{PRÄTERITUM} }\\
    ich & \textbf{lernte}\\
    du & \textbf{lerntest}\\
    er/sie/es & \textbf{lernte}\\
    wir & \textbf{lernten}\\
    ihr & \textbf{lerntet}\\
    sie/Sie & \textbf{lernten}\\
  \end{tabular}
  \hfill
  \begin{tabular}{rl}
     \multicolumn{2}{l}{\textbf{IMPERATIV} }\\
   & \textbf{glern(e) (du)!}\\
   & \textbf{lernen wir!}\\
   & \textbf{lernt (ihr)!}\\
   & \textbf{lernen Sie!}\\
   & \vspace{10pt}
  \end{tabular}

    \vspace{10pt}

 \begin{tabular}{rl}
     \multicolumn{2}{l}{\textbf{PERFEKT}}\\
  &  haben + \textbf{gelernt}\\
  \end{tabular}
\hfill
   \begin{tabular}{rl}
   
    \multicolumn{2}{l}{\textbf{FUTURE I} }\\
  &  werden + \textbf{lernen}\\
 
  \end{tabular}
\hfill
   \begin{tabular}{rl}
      \multicolumn{2}{l}{\textbf{PLUSQUAMPERFEKT}}\\
   & hatten + \textbf{gelernt}\\
  \end{tabular}
  
  \vspace{-5pt}


\end{verbbox-acc}
%lesen (A1) – to read
\begin{verbbox-acc}\addcontentsline{toc}{subsubsection}{lesen (to read)}
    {\Large \bf lesen} (to read) \hfill Type: \textbf{irregular, + acusative} \hfill Auxiliary Verb: \textbf{haben} 
    
     \vspace{5pt}

    \hrule
    
    \vspace{5pt}

  \begin{tabular}{rl}
     \multicolumn{2}{l}{\textbf{PRÄSENS} }\\
    ich & \textbf{lese}\\
    du & \textbf{liest}\\
    er/sie/es & \textbf{liest}\\
    wir & \textbf{lesen}\\
    ihr & \textbf{lest}\\
    sie/Sie & \textbf{lesen}\\
  \end{tabular}
  \hfill
     \begin{tabular}{rl}
    \multicolumn{2}{l}{\textbf{PRÄTERITUM} }\\
    ich & \textbf{las}\\
    du & \textbf{lastest}\\
    er/sie/es & \textbf{las}\\
    wir & \textbf{lasen}\\
    ihr & \textbf{last}\\
    sie/Sie & \textbf{lasen}\\
  \end{tabular}
  \hfill
  \begin{tabular}{rl}
     \multicolumn{2}{l}{\textbf{IMPERATIV} }\\
   & \textbf{lies (du)!}\\
   & \textbf{lesen wir!}\\
   & \textbf{lest (ihr)!}\\
   & \textbf{lesen Sie!}\\
   & \vspace{10pt}
  \end{tabular}

    \vspace{10pt}

 \begin{tabular}{rl}
     \multicolumn{2}{l}{\textbf{PERFEKT}}\\
  &  haben + \textbf{gelesen}\\
  \end{tabular}
\hfill
   \begin{tabular}{rl}
   
    \multicolumn{2}{l}{\textbf{FUTURE I} }\\
  &  werden + \textbf{lesen}\\
 
  \end{tabular}
\hfill
   \begin{tabular}{rl}
      \multicolumn{2}{l}{\textbf{PLUSQUAMPERFEKT}}\\
   & hatten + \textbf{gelesen}\\
  \end{tabular}
  
  \vspace{-5pt}


\end{verbbox-acc}
%lieben (A1) – to love
\begin{verbbox-acc}\addcontentsline{toc}{subsubsection}{lieben (to love)}
    {\Large \bf lieben} (to love) \hfill Type: \textbf{regular, + Akkusativ} \hfill Auxiliary Verb: \textbf{haben} 
    
     \vspace{5pt}

    \hrule
    
    \vspace{5pt}

  \begin{tabular}{rl}
     \multicolumn{2}{l}{\textbf{PRÄSENS} }\\
    ich & \textbf{liebe}\\
    du & \textbf{liebst}\\
    er/sie/es & \textbf{liebt}\\
    wir & \textbf{lieben}\\
    ihr & \textbf{liebt}\\
    sie/Sie & \textbf{lieben}\\
  \end{tabular}
  \hfill
     \begin{tabular}{rl}
    \multicolumn{2}{l}{\textbf{PRÄTERITUM} }\\
    ich & \textbf{liebte}\\
    du & \textbf{liebtest}\\
    er/sie/es & \textbf{liebte}\\
    wir & \textbf{liebten}\\
    ihr & \textbf{liebtet}\\
    sie/Sie & \textbf{liebten}\\
  \end{tabular}
  \hfill
  \begin{tabular}{rl}
     \multicolumn{2}{l}{\textbf{IMPERATIV} }\\
   & \textbf{lieb(e) (du)!}\\
   & \textbf{lieben wir!}\\
   & \textbf{liebt (ihr)!}\\
   & \textbf{lieben Sie!}\\
   & \vspace{10pt}
  \end{tabular}

    \vspace{10pt}

 \begin{tabular}{rl}
     \multicolumn{2}{l}{\textbf{PERFEKT}}\\
   & haben + \textbf{geliebt}\\
  \end{tabular}
\hfill
   \begin{tabular}{rl}
   
    \multicolumn{2}{l}{\textbf{FUTURE I} }\\
   & werden + \textbf{lieben}\\
 
  \end{tabular}
\hfill
   \begin{tabular}{rl}
      \multicolumn{2}{l}{\textbf{PLUSQUAMPERFEKT}}\\
   & hatten + \textbf{geliebt}\\
  \end{tabular}
  
  \vspace{-5pt}


\end{verbbox-acc}

%loben (to praise)
\begin{verbbox}\addcontentsline{toc}{subsubsection}{loben (to praise)}
    {\Large \bf loben} (to praise) \hfill Type: \textbf{regular} \hfill Auxiliary Verb: \textbf{haben} 
    
     \vspace{5pt}

    \hrule
    
    \vspace{5pt}

  \begin{tabular}{rl}
     \multicolumn{2}{l}{\textbf{PRÄSENS} }\\
    ich & \textbf{lobe}\\
    du & \textbf{lobst}\\
    er/sie/es & \textbf{lobt}\\
    wir & \textbf{loben}\\
    ihr & \textbf{lobt}\\
    sie/Sie & \textbf{loben}\\
  \end{tabular}
  \hfill
     \begin{tabular}{rl}
    \multicolumn{2}{l}{\textbf{PRÄTERITUM} }\\
    ich & \textbf{lobte}\\
    du & \textbf{lobtest}\\
    er/sie/es & \textbf{lobte}\\
    wir & \textbf{lobten}\\
    ihr & \textbf{lobtet}\\
    sie/Sie & \textbf{lobten}\\
  \end{tabular}
  \hfill
  \begin{tabular}{rl}
     \multicolumn{2}{l}{\textbf{IMPERATIV} }\\
   & \textbf{lob(e) (du)!}\\
   & \textbf{loben wir!}\\
   & \textbf{lobt (ihr)!}\\
   & \textbf{loben Sie!}\\
   & \vspace{10pt}
  \end{tabular}

    \vspace{10pt}

 \begin{tabular}{rl}
     \multicolumn{2}{l}{\textbf{PERFEKT}}\\
   & haben + \textbf{gelobt}\\
  \end{tabular}
\hfill
   \begin{tabular}{rl}
    \multicolumn{2}{l}{\textbf{FUTURE I} }\\
   & werden + \textbf{loben}\\
  \end{tabular}
\hfill
   \begin{tabular}{rl}
      \multicolumn{2}{l}{\textbf{PLUSQUAMPERFEKT}}\\
   & hatten + \textbf{gelobt}\\
  \end{tabular}

   \vspace{10pt}

   \begin{tabular}{lll}
      \multicolumn{3}{l}{\textbf{MIT PRÄPOSITIONEN}}\\
& \textbf{loben für} + Akk & (to praise for)\\
& \textbf{loben wegen} + Gen & (to praise because of)\\
\end{tabular}
  \hfill
   \begin{tabular}{lll}
      \multicolumn{3}{l}{\textbf{TRENNBARE VERBEN}}\\
 & \textbf{---} & (none common)\\
\vspace{10pt}
  \end{tabular}

  \vspace{-5pt}

\end{verbbox}


\subsection{M Verbs}


%machen (A1) – to make / do
\begin{verbbox-acc}\addcontentsline{toc}{subsubsection}{machen (to make)}
    {\Large \bf machen} (to make) \hfill Type: \textbf{regular, + Akkusative} \hfill Auxiliary Verb: \textbf{haben} 
    
     \vspace{5pt}

    \hrule
    
    \vspace{5pt}

  \begin{tabular}{rl}
     \multicolumn{2}{l}{\textbf{PRÄSENS} }\\
    ich & \textbf{mache}\\
    du & \textbf{machst}\\
    er/sie/es & \textbf{macht}\\
    wir & \textbf{machen}\\
    ihr & \textbf{macht}\\
    sie/Sie & \textbf{machen}\\
  \end{tabular}
  \hfill
     \begin{tabular}{rl}
    \multicolumn{2}{l}{\textbf{PRÄTERITUM} }\\
    ich & \textbf{machte}\\
    du & \textbf{machtest}\\
    er/sie/es & \textbf{machte}\\
    wir & \textbf{machten}\\
    ihr & \textbf{machtet}\\
    sie/Sie & \textbf{machten}\\
  \end{tabular}
  \hfill
  \begin{tabular}{rl}
     \multicolumn{2}{l}{\textbf{IMPERATIV} }\\
   & \textbf{mach(e) (du)!}\\
   & \textbf{machen wir!}\\
   & \textbf{macht (ihr)!}\\
   & \textbf{machen Sie!}\\
   & \vspace{10pt}
  \end{tabular}

    \vspace{10pt}

 \begin{tabular}{rl}
     \multicolumn{2}{l}{\textbf{PERFEKT}}\\
   & haben + \textbf{gemacht}\\
  \end{tabular}
\hfill
   \begin{tabular}{rl}
   
    \multicolumn{2}{l}{\textbf{FUTURE I} }\\
   & werden + \textbf{machen}\\
 
  \end{tabular}
\hfill
   \begin{tabular}{rl}
      \multicolumn{2}{l}{\textbf{PLUSQUAMPERFEKT}}\\
   & hatten + \textbf{gemacht}\\
  \end{tabular}
  

  \vspace{10pt}

   \begin{tabular}{lll}
      \multicolumn{3}{l}{\textbf{MIT PRÄPOSITIONEN}}\\
& \textbf{machen mit} + Dat & (to participate in)\\
& \textbf{machen für} + Akk & (to do something for someone / purpose)\\
&\textbf{machen aus} + Dat & (to turn something into something)\\
&\textbf{machen über} + Akk & (to complain / joke about)\\
&\textbf{machen von} + Dat & (to make of / think of)\\
&\textbf{machen bei} + Dat & (to take part in)\\
\vspace{40pt}
  \end{tabular}
  \hfill
   \begin{tabular}{lll}
      \multicolumn{3}{l}{\textbf{TRENNBARE VERBEN}}\\
 & \textbf{an$|$machen} & (to turn on)\\
 & \textbf{aus$|$machen} & (to turn off)\\
 & \textbf{auf$|$machen} & (to open)\\
  & \textbf{zu$|$machen} & (to close)\\
  & \textbf{mit$|$machen} & (to participate)\\
  & \textbf{nach$|$machen} & (to imitate)\\
  & \textbf{weiter$|$machen} & (to continue)\\
     & \textbf{fertig$|$machen} & (to finish /exhaust)\\
     & \textbf{kaputt$|$machen} & (to break)\\
     & \textbf{los$|$machen} & (to loosen / detach)\\
  \end{tabular}


  \vspace{-5pt}

\end{verbbox-acc}

%mauern - to build walls
\begin{verbbox}\addcontentsline{toc}{subsubsection}{mauern (to build walls / brick)}
    {\Large \bf mauern} (to build walls / brick) \hfill Type: \textbf{regular} \hfill Auxiliary Verb: \textbf{haben} 
    
     \vspace{5pt}

    \hrule
    
    \vspace{5pt}

  \begin{tabular}{rl}
     \multicolumn{2}{l}{\textbf{PRÄSENS} }\\
    ich & \textbf{mauere}\\
    du & \textbf{mauerst}\\
    er/sie/es & \textbf{mauert}\\
    wir & \textbf{mauern}\\
    ihr & \textbf{mauert}\\
    sie/Sie & \textbf{mauern}\\
  \end{tabular}
  \hfill
  \begin{tabular}{rl}
    \multicolumn{2}{l}{\textbf{PRÄTERITUM} }\\
    ich & \textbf{mauerte}\\
    du & \textbf{mauertest}\\
    er/sie/es & \textbf{mauerte}\\
    wir & \textbf{mauerten}\\
    ihr & \textbf{mauertet}\\
    sie/Sie & \textbf{mauerten}\\
  \end{tabular}
  \hfill
  \begin{tabular}{rl}
     \multicolumn{2}{l}{\textbf{IMPERATIV} }\\
   & \textbf{mauere (du)!}\\
   & \textbf{mauern wir!}\\
   & \textbf{mauert (ihr)!}\\
   & \textbf{mauern Sie!}\\
   & \vspace{10pt}
  \end{tabular}

 \vspace{10pt}

 \begin{tabular}{rl}
     \multicolumn{2}{l}{\textbf{PERFEKT}}\\
   & haben + \textbf{gemauert}\\
  \end{tabular}
\hfill
\begin{tabular}{rl}
    \multicolumn{2}{l}{\textbf{FUTURE I}}\\
   & werden + \textbf{mauern}\\
  \end{tabular}
\hfill
\begin{tabular}{rl}
    \multicolumn{2}{l}{\textbf{PLUSQUAMPERFEKT}}\\
   & hatten + \textbf{gemauert}\\
  \end{tabular}

 \vspace{10pt}

\begin{tabular}{lll}
      \multicolumn{3}{l}{\textbf{MIT PRÄPOSITIONEN}}\\
& \textbf{mauern mit} + Dat & (to build with / using)\\
& \textbf{mauern an} + Dat & (to build onto / on)\\
\end{tabular}
  \hfill
\begin{tabular}{lll}
      \multicolumn{3}{l}{\textbf{TRENNBARE VERBEN}}\\
 & \textbf{ein$|$mauern} & (to wall in / enclose)\\
\vspace{10pt}
  \end{tabular}

  \vspace{-5pt}

\end{verbbox}

%meinen (B1) – to mean something
\begin{verbbox-acc}\addcontentsline{toc}{subsubsection}{meinen (to mean something)}
    {\Large \bf meinen} (to mean something) \hfill Type: \textbf{regular, + Akkusativ} \hfill Auxiliary Verb: \textbf{haben} 
    
     \vspace{5pt}

    \hrule
    
    \vspace{5pt}

  \begin{tabular}{rl}
     \multicolumn{2}{l}{\textbf{PRÄSENS} }\\
    ich & \textbf{meine}\\
    du & \textbf{meinst}\\
    er/sie/es & \textbf{meint}\\
    wir & \textbf{meinen}\\
    ihr & \textbf{meint}\\
    sie/Sie & \textbf{meinen}\\
  \end{tabular}
  \hfill
     \begin{tabular}{rl}
    \multicolumn{2}{l}{\textbf{PRÄTERITUM} }\\
    ich & \textbf{meinte}\\
    du & \textbf{meintest}\\
    er/sie/es & \textbf{meinte}\\
    wir & \textbf{meinten}\\
    ihr & \textbf{meintet}\\
    sie/Sie & \textbf{meinten}\\
  \end{tabular}
  \hfill
  \begin{tabular}{rl}
     \multicolumn{2}{l}{\textbf{IMPERATIV} }\\
   & \textbf{mein(e) (du)!}\\
   & \textbf{meinen wir!}\\
   & \textbf{meint (ihr)!}\\
   & \textbf{meinen Sie!}\\
   & \vspace{10pt}
  \end{tabular}

    \vspace{10pt}

 \begin{tabular}{rl}
     \multicolumn{2}{l}{\textbf{PERFEKT}}\\
  &  haben + \textbf{gemeint}\\
  \end{tabular}
\hfill
   \begin{tabular}{rl}
   
    \multicolumn{2}{l}{\textbf{FUTURE I} }\\
   & werden + \textbf{meinen}\\
 
  \end{tabular}
\hfill
   \begin{tabular}{rl}
      \multicolumn{2}{l}{\textbf{PLUSQUAMPERFEKT}}\\
   & hatten + \textbf{gemeint}\\
  \end{tabular}
  
  \vspace{-5pt}


\end{verbbox-acc}

%melken - to milk
\begin{verbbox}\addcontentsline{toc}{subsubsection}{melken (to milk)}
    {\Large \bf melken} (to milk) \hfill Type: \textbf{irregular} \hfill Auxiliary Verb: \textbf{haben} 
    
     \vspace{5pt}

    \hrule
    
    \vspace{5pt}

  \begin{tabular}{rl}
     \multicolumn{2}{l}{\textbf{PRÄSENS} }\\
    ich & \textbf{melke}\\
    du & \textbf{melkst}\\
    er/sie/es & \textbf{melkt}\\
    wir & \textbf{melken}\\
    ihr & \textbf{melkt}\\
    sie/Sie & \textbf{melken}\\
  \end{tabular}
  \hfill
     \begin{tabular}{rl}
    \multicolumn{2}{l}{\textbf{PRÄTERITUM} }\\
    ich & \textbf{molk}\\
    du & \textbf{molkst}\\
    er/sie/es & \textbf{molk}\\
    wir & \textbf{molken}\\
    ihr & \textbf{molkt}\\
    sie/Sie & \textbf{molken}\\
  \end{tabular}
  \hfill
  \begin{tabular}{rl}
     \multicolumn{2}{l}{\textbf{IMPERATIV} }\\
   & \textbf{melk(e) (du)!}\\
   & \textbf{melken wir!}\\
   & \textbf{melkt (ihr)!}\\
   & \textbf{melken Sie!}\\
   & \vspace{10pt}
  \end{tabular}

    \vspace{10pt}

 \begin{tabular}{rl}
     \multicolumn{2}{l}{\textbf{PERFEKT}}\\
   & haben + \textbf{gemolken}\\
  \end{tabular}
\hfill
   \begin{tabular}{rl}
   
    \multicolumn{2}{l}{\textbf{FUTURE I} }\\
   & werden + \textbf{melken}\\
 
  \end{tabular}
\hfill
   \begin{tabular}{rl}
      \multicolumn{2}{l}{\textbf{PLUSQUAMPERFEKT}}\\
   & hatten + \textbf{gemolken}\\
  \end{tabular}

   \vspace{10pt}

   \begin{tabular}{lll}
      \multicolumn{3}{l}{\textbf{MIT PRÄPOSITIONEN}}\\
& \textbf{melken von} + Dat & (to extract from)\\
& \textbf{melken aus} + Dat & (to extract out of)\\
   \end{tabular}
  \hfill
   \begin{tabular}{lll}
      \multicolumn{3}{l}{\textbf{TRENNBARE VERBEN}}\\
 & \textbf{ab$|$melken} & (to siphon off / exploit)\\
 & \textbf{aus$|$melken} & (to milk out / exploit fully)\\
  \end{tabular}

\vspace{-5pt}

\end{verbbox}

%merken (to notice / remember)
\begin{verbbox}\addcontentsline{toc}{subsubsection}{merken (to notice / remember)}
{\Large \bf merken} (to notice / remember) \hfill Type: \textbf{regular} \hfill Auxiliary Verb: \textbf{haben}

\vspace{5pt}
\hrule
\vspace{5pt}

\begin{tabular}{rl}
\multicolumn{2}{l}{\textbf{PRÄSENS}}\\
ich & \textbf{merke}\\
du & \textbf{merkst}\\
er/sie/es & \textbf{merkt}\\
wir & \textbf{merken}\\
ihr & \textbf{merkt}\\
sie/Sie & \textbf{merken}\\
\end{tabular}
\hfill
\begin{tabular}{rl}
\multicolumn{2}{l}{\textbf{PRÄTERITUM}}\\
ich & \textbf{merkte}\\
du & \textbf{merktest}\\
er/sie/es & \textbf{merkte}\\
wir & \textbf{merkten}\\
ihr & \textbf{merktet}\\
sie/Sie & \textbf{merkten}\\
\end{tabular}
\hfill
\begin{tabular}{rl}
\multicolumn{2}{l}{\textbf{IMPERATIV}}\\
& \textbf{merk(e) (du)!}\\
& \textbf{merken wir!}\\
& \textbf{merkt (ihr)!}\\
& \textbf{merken Sie!}\\
\end{tabular}

\vspace{10pt}

\begin{tabular}{rl}
\multicolumn{2}{l}{\textbf{PERFEKT}}\\
& haben + \textbf{gemerkt}\\
\end{tabular}
\hfill
\begin{tabular}{rl}
\multicolumn{2}{l}{\textbf{FUTURE I}}\\
& werden + \textbf{merken}\\
\end{tabular}
\hfill
\begin{tabular}{rl}
\multicolumn{2}{l}{\textbf{PLUSQUAMPERFEKT}}\\
& hatten + \textbf{gemerkt}\\
\end{tabular}

\vspace{10pt}

\begin{tabular}{lll}
\multicolumn{3}{l}{\textbf{MIT PRÄPOSITIONEN}}\\
& \textbf{merken an} + Dat & (to notice in)\\
& \textbf{merken sich} + Akk & (to remember)\\
\end{tabular}
\hfill
\begin{tabular}{lll}
\multicolumn{3}{l}{\textbf{TRENNBARE VERBEN}}\\
& \textbf{fest$|$merken} & (to take note of)\\
& \textbf{zurück$|$merken} & (to remember for later)\\
\end{tabular}
\vspace{-5pt}
\end{verbbox}


%missverstehen (to misunderstand)
\begin{verbbox}\addcontentsline{toc}{subsubsection}{missverstehen (to misunderstand)}
    {\Large \bf missverstehen} (to misunderstand) \hfill Type: \textbf{irregular} \hfill Auxiliary Verb: \textbf{haben} 
    
     \vspace{5pt}

    \hrule
    
    \vspace{5pt}

  \begin{tabular}{rl}
     \multicolumn{2}{l}{\textbf{PRÄSENS} }\\
    ich & \textbf{missverstehe}\\
    du & \textbf{missverstehst}\\
    er/sie/es & \textbf{missversteht}\\
    wir & \textbf{missverstehen}\\
    ihr & \textbf{missversteht}\\
    sie/Sie & \textbf{missverstehen}\\
  \end{tabular}
  \hfill
     \begin{tabular}{rl}
    \multicolumn{2}{l}{\textbf{PRÄTERITUM} }\\
    ich & \textbf{missverstand}\\
    du & \textbf{missverstandest}\\
    er/sie/es & \textbf{missverstand}\\
    wir & \textbf{missverstanden}\\
    ihr & \textbf{missverstandet}\\
    sie/Sie & \textbf{missverstanden}\\
  \end{tabular}
  \hfill
  \begin{tabular}{rl}
     \multicolumn{2}{l}{\textbf{IMPERATIV} }\\
   & \textbf{missversteh(e) (du)!}\\
   & \textbf{missverstehen wir!}\\
   & \textbf{missversteht (ihr)!}\\
   & \textbf{missverstehen Sie!}\\
   & \vspace{10pt}
  \end{tabular}

    \vspace{10pt}

 \begin{tabular}{rl}
     \multicolumn{2}{l}{\textbf{PERFEKT}}\\
   & haben + \textbf{missverstanden}\\
  \end{tabular}
\hfill
   \begin{tabular}{rl}
    \multicolumn{2}{l}{\textbf{FUTURE I} }\\
   & werde + \textbf{missverstehen}\\
  \end{tabular}
\hfill
   \begin{tabular}{rl}
      \multicolumn{2}{l}{\textbf{PLUSQUAMPERFEKT}}\\
   & hatten + \textbf{missverstanden}\\
  \end{tabular}

   \vspace{10pt}

   \begin{tabular}{lll}
      \multicolumn{3}{l}{\textbf{MIT PRÄPOSITIONEN}}\\
& \textbf{missverstehen als} + Akk & (to misunderstand as)\\
& \textbf{missverstehen wegen} + Gen & (to misunderstand because of)\\
\end{tabular}
  \hfill
   \begin{tabular}{lll}
      \multicolumn{3}{l}{\textbf{TRENNBARE VERBEN}}\\
 & \textbf{---} & (none common)\\
\vspace{10pt}
  \end{tabular}

  \vspace{-5pt}

\end{verbbox}

\subsection{N Verbs}

%nähen (to sew)
\begin{verbbox}\addcontentsline{toc}{subsubsection}{nähen (to sew)}
{\Large \bf nähen} (to sew) \hfill Type: \textbf{regular} \hfill Auxiliary Verb: \textbf{haben}

\vspace{5pt}
\hrule
\vspace{5pt}

\begin{tabular}{rl}
\multicolumn{2}{l}{\textbf{PRÄSENS}}\\
ich & \textbf{nähe}\\
du & \textbf{nähst}\\
er/sie/es & \textbf{näht}\\
wir & \textbf{nähen}\\
ihr & \textbf{näht}\\
sie/Sie & \textbf{nähen}\\
\end{tabular}
\hfill
\begin{tabular}{rl}
\multicolumn{2}{l}{\textbf{PRÄTERITUM}}\\
ich & \textbf{nähte}\\
du & \textbf{nähntest}\\
er/sie/es & \textbf{nähte}\\
wir & \textbf{nähten}\\
ihr & \textbf{nähtet}\\
sie/Sie & \textbf{nähten}\\
\end{tabular}
\hfill
\begin{tabular}{rl}
\multicolumn{2}{l}{\textbf{IMPERATIV}}\\
& \textbf{nähe (du)!}\\
& \textbf{nähen wir!}\\
& \textbf{näht (ihr)!}\\
& \textbf{nähen Sie!}\\
\end{tabular}

\vspace{10pt}

\begin{tabular}{rl}
\multicolumn{2}{l}{\textbf{PERFEKT}}\\
& haben + \textbf{genäht}\\
\end{tabular}
\hfill
\begin{tabular}{rl}
\multicolumn{2}{l}{\textbf{FUTURE I}}\\
& werden + \textbf{nähen}\\
\end{tabular}
\hfill
\begin{tabular}{rl}
\multicolumn{2}{l}{\textbf{PLUSQUAMPERFEKT}}\\
& hatten + \textbf{genäht}\\
\end{tabular}

\vspace{10pt}

\begin{tabular}{lll}
\multicolumn{3}{l}{\textbf{MIT PRÄPOSITIONEN}}\\
& \textbf{nähen an} + Dat & (to sew on / at)\\
& \textbf{nähen für} + Akk & (to sew for / on behalf of)\\
\end{tabular}
\hfill
\begin{tabular}{lll}
\multicolumn{3}{l}{\textbf{TRENNBARE VERBEN}}\\
& \textbf{auf$|$nähen} & (to sew onto / attach by sewing)\\
\end{tabular}
\vspace{-5pt}
\end{verbbox}


%nehmen (A1) – to take
\begin{verbbox-acc}\addcontentsline{toc}{subsubsection}{nehmen (to take)}
    {\Large \bf nehmen} (to take) \hfill Type: \textbf{irregular, + Akkusative} \hfill Auxiliary Verb: \textbf{haben} 
    
     \vspace{5pt}

    \hrule
    
    \vspace{5pt}

  \begin{tabular}{rl}
     \multicolumn{2}{l}{\textbf{PRÄSENS} }\\
    ich & \textbf{nehme}\\
    du & \textbf{nimmst}\\
    er/sie/es & \textbf{nimmt}\\
    wir & \textbf{nehmen}\\
    ihr & \textbf{nehmt}\\
    sie/Sie & \textbf{nehmen}\\
  \end{tabular}
  \hfill
     \begin{tabular}{rl}
    \multicolumn{2}{l}{\textbf{PRÄTERITUM} }\\
    ich & \textbf{nahm}\\
    du & \textbf{nahmst}\\
    er/sie/es & \textbf{nahm}\\
    wir & \textbf{nahmen}\\
    ihr & \textbf{nahmt}\\
    sie/Sie & \textbf{nahmen}\\
  \end{tabular}
  \hfill
  \begin{tabular}{rl}
     \multicolumn{2}{l}{\textbf{IMPERATIV} }\\
   & \textbf{nimm (du)!}\\
   & \textbf{nehmen wir!}\\
   & \textbf{nehmt (ihr)!}\\
   & \textbf{nehmen Sie!}\\
   & \vspace{10pt}
  \end{tabular}

    \vspace{10pt}

 \begin{tabular}{rl}
     \multicolumn{2}{l}{\textbf{PERFEKT}}\\
  &  haben + \textbf{genommen}\\
  \end{tabular}
\hfill
   \begin{tabular}{rl}
   
    \multicolumn{2}{l}{\textbf{FUTURE I} }\\
  &  werden + \textbf{nehmen}\\
 
  \end{tabular}
\hfill
   \begin{tabular}{rl}
      \multicolumn{2}{l}{\textbf{PLUSQUAMPERFEKT}}\\
  &  hatten + \textbf{genommen}\\
  \end{tabular}

    \vspace{10pt}

   \begin{tabular}{lll}
      \multicolumn{3}{l}{\textbf{MIT PRÄPOSITIONEN}}\\
& \textbf{nehmen aus} + Dat & (to take from)\\
& \textbf{nehmen von} + Dat & (to take from)\\
&\textbf{nehmen für} + Akk & (to take for)\\
&\textbf{nehmen über} + Akk & (to take over / control of)\\
\vspace{50pt}
  \end{tabular}
  \hfill
   \begin{tabular}{lll}
      \multicolumn{3}{l}{\textbf{TRENNBARE VERBEN}}\\
      & \textbf{ab$|$nehmen} & (to turn off / lose weight)\\
 & \textbf{an$|$nehmen} & (to accept / assume)\\
 & \textbf{auf$|$nehmen} & (to pick up / record)\\
 & \textbf{ein$|$nehmen} & (to take something / to occupy)\\
  & \textbf{mit$|$nehmen} & (to take along)\\
  & \textbf{über$|$nehmen} & (to take over)\\
  & \textbf{weg$|$nehmen} & (to take away)\\
     & \textbf{zurück$|$nehmen} & (to take back)\\
     & \textbf{teil$|$nehmen} & (to participate)\\
       & \textbf{zu$|$nehmen} & (to gain weight)\\
  \end{tabular}



  \vspace{-5pt}

\end{verbbox-acc}
%nennen (B1) – to name
\begin{verbbox-acc}\addcontentsline{toc}{subsubsection}{nennen (to name)}
    {\Large \bf nennen} (to name) \hfill Type: \textbf{irregular, + Akkusativ} \hfill Auxiliary Verb: \textbf{haben} 
    
     \vspace{5pt}

    \hrule
    
    \vspace{5pt}

  \begin{tabular}{rl}
     \multicolumn{2}{l}{\textbf{PRÄSENS} }\\
    ich & \textbf{nenne}\\
    du & \textbf{nennst}\\
    er/sie/es & \textbf{nennt}\\
    wir & \textbf{nennen}\\
    ihr & \textbf{nennt}\\
    sie/Sie & \textbf{nennen}\\
  \end{tabular}
  \hfill
     \begin{tabular}{rl}
    \multicolumn{2}{l}{\textbf{PRÄTERITUM} }\\
    ich & \textbf{nannte}\\
    du & \textbf{nanntest}\\
    er/sie/es & \textbf{nannte}\\
    wir & \textbf{nannten}\\
    ihr & \textbf{nanntet}\\
    sie/Sie & \textbf{nannten}\\
  \end{tabular}
  \hfill
  \begin{tabular}{rl}
     \multicolumn{2}{l}{\textbf{IMPERATIV} }\\
   & \textbf{gnenn(e) (du)!}\\
   & \textbf{nennen wir!}\\
   & \textbf{nennt (ihr)!}\\
   & \textbf{gnennen Sie!}\\
   & \vspace{10pt}
  \end{tabular}

    \vspace{10pt}

 \begin{tabular}{rl}
     \multicolumn{2}{l}{\textbf{PERFEKT}}\\
  &  haben + \textbf{genannt}\\
  \end{tabular}
\hfill
   \begin{tabular}{rl}
   
    \multicolumn{2}{l}{\textbf{FUTURE I} }\\
   & werden + \textbf{nennen}\\
 
  \end{tabular}
\hfill
   \begin{tabular}{rl}
      \multicolumn{2}{l}{\textbf{PLUSQUAMPERFEKT}}\\
   & hatten + \textbf{genannt}\\
  \end{tabular}
  
  \vspace{-5pt}


\end{verbbox-acc}

%niesen (to sneeze)
\begin{verbbox}\addcontentsline{toc}{subsubsection}{niesen (to sneeze)}
{\Large \bf niesen} (to sneeze) \hfill Type: \textbf{regular} \hfill Auxiliary Verb: \textbf{haben}

\vspace{5pt}
\hrule
\vspace{5pt}

\begin{tabular}{rl}
\multicolumn{2}{l}{\textbf{PRÄSENS}}\\
ich & \textbf{niese}\\
du & \textbf{niest}\\
er/sie/es & \textbf{niest}\\
wir & \textbf{niesen}\\
ihr & \textbf{niest}\\
sie/Sie & \textbf{niesen}\\
\end{tabular}
\hfill
\begin{tabular}{rl}
\multicolumn{2}{l}{\textbf{PRÄTERITUM}}\\
ich & \textbf{nieste}\\
du & \textbf{niestest}\\
er/sie/es & \textbf{nieste}\\
wir & \textbf{niesten}\\
ihr & \textbf{niestet}\\
sie/Sie & \textbf{niesten}\\
\end{tabular}
\hfill
\begin{tabular}{rl}
\multicolumn{2}{l}{\textbf{IMPERATIV}}\\
& \textbf{niese (du)!}\\
& \textbf{niesen wir!}\\
& \textbf{niest (ihr)!}\\
& \textbf{niesen Sie!}\\
\end{tabular}

\vspace{10pt}

\begin{tabular}{rl}
\multicolumn{2}{l}{\textbf{PERFEKT}}\\
& haben + \textbf{geniest}\\
\end{tabular}
\hfill
\begin{tabular}{rl}
\multicolumn{2}{l}{\textbf{FUTURE I}}\\
& werden + \textbf{niesen}\\
\end{tabular}
\hfill
\begin{tabular}{rl}
\multicolumn{2}{l}{\textbf{PLUSQUAMPERFEKT}}\\
& hatten + \textbf{geniest}\\
\end{tabular}

\vspace{10pt}

\begin{tabular}{lll}
\multicolumn{3}{l}{\textbf{MIT PRÄPOSITIONEN}}\\
& \textbf{niesen in} + Akk & (to sneeze into)\\
\end{tabular}
\hfill
\begin{tabular}{lll}
\multicolumn{3}{l}{\textbf{TRENNBARE VERBEN}}\\
& ---\\
\end{tabular}
\vspace{-5pt}
\end{verbbox}

%nutzen (A2) – to use
\begin{verbbox-acc}\addcontentsline{toc}{subsubsection}{nutzen (to use)}
    {\Large \bf nutzen} (to use) \hfill Type: \textbf{regular, + Akkusativ} \hfill Auxiliary Verb: \textbf{haben} 
    
     \vspace{5pt}

    \hrule
    
    \vspace{5pt}

  \begin{tabular}{rl}
     \multicolumn{2}{l}{\textbf{PRÄSENS} }\\
    ich & \textbf{nutze}\\
    du & \textbf{nutzt}\\
    er/sie/es & \textbf{nutzt}\\
    wir & \textbf{nutzen}\\
    ihr & \textbf{nutzt}\\
    sie/Sie & \textbf{nutzen}\\
  \end{tabular}
  \hfill
     \begin{tabular}{rl}
    \multicolumn{2}{l}{\textbf{PRÄTERITUM} }\\
    ich & \textbf{nutzte}\\
    du & \textbf{nutztest}\\
    er/sie/es & \textbf{nutzte}\\
    wir & \textbf{nutzten}\\
    ihr & \textbf{nutztet}\\
    sie/Sie & \textbf{nutzten}\\
  \end{tabular}
  \hfill
  \begin{tabular}{rl}
     \multicolumn{2}{l}{\textbf{IMPERATIV} }\\
   & \textbf{nutz(e) (du)!}\\
   & \textbf{ven wir!}\\
   & \textbf{nutzt (ihr)!}\\
   & \textbf{nutzen Sie!}\\
   & \vspace{10pt}
  \end{tabular}

    \vspace{10pt}

 \begin{tabular}{rl}
     \multicolumn{2}{l}{\textbf{PERFEKT}}\\
  &  haben + \textbf{genutzt}\\
  \end{tabular}
\hfill
   \begin{tabular}{rl}
   
    \multicolumn{2}{l}{\textbf{FUTURE I} }\\
  &  werden + \textbf{nutzen}\\
 
  \end{tabular}
\hfill
   \begin{tabular}{rl}
      \multicolumn{2}{l}{\textbf{PLUSQUAMPERFEKT}}\\
  &  hatten + \textbf{genutzt}\\
  \end{tabular}
  
  \vspace{-5pt}


\end{verbbox-acc}

\subsection{O Verbs}

%öffnen (A1) – to open
\begin{verbbox-acc}\addcontentsline{toc}{subsubsection}{öffnen (to open)}
    {\Large \bf öffnen} (to open) \hfill Type: \textbf{regular, + Akkusativ} \hfill Auxiliary Verb: \textbf{haben} 
    
     \vspace{5pt}

    \hrule
    
    \vspace{5pt}

  \begin{tabular}{rl}
     \multicolumn{2}{l}{\textbf{PRÄSENS} }\\
    ich & \textbf{öffne}\\
    du & \textbf{öffnest}\\
    er/sie/es & \textbf{öffnet}\\
    wir & \textbf{öffnen}\\
    ihr & \textbf{öffnet}\\
    sie/Sie & \textbf{öffnen}\\
  \end{tabular}
  \hfill
     \begin{tabular}{rl}
    \multicolumn{2}{l}{\textbf{PRÄTERITUM} }\\
    ich & \textbf{öffnete}\\
    du & \textbf{öffnetest}\\
    er/sie/es & \textbf{öffnete}\\
    wir & \textbf{öffneten}\\
    ihr & \textbf{öffnetet}\\
    sie/Sie & \textbf{öffneten}\\
  \end{tabular}
  \hfill
  \begin{tabular}{rl}
     \multicolumn{2}{l}{\textbf{IMPERATIV} }\\
   & \textbf{göffne (du)!}\\
   & \textbf{öffnen wir!}\\
   & \textbf{öffnet (ihr)!}\\
   & \textbf{öffnen Sie!}\\
   & \vspace{10pt}
  \end{tabular}

    \vspace{10pt}

 \begin{tabular}{rl}
     \multicolumn{2}{l}{\textbf{PERFEKT}}\\
   & haben + \textbf{geöffnet}\\
  \end{tabular}
\hfill
   \begin{tabular}{rl}
   
    \multicolumn{2}{l}{\textbf{FUTURE I} }\\
  &  werden + \textbf{öffnen}\\
 
  \end{tabular}
\hfill
   \begin{tabular}{rl}
      \multicolumn{2}{l}{\textbf{PLUSQUAMPERFEKT}}\\
  &  hatten + \textbf{geöffnet}\\
  \end{tabular}
  
  \vspace{-5pt}


\end{verbbox-acc}

%organisieren (to organize)
\begin{verbbox-acc}\addcontentsline{toc}{subsubsection}{organisieren (to organize)}
{\Large \bf organisieren} (to organize) \hfill Type: \textbf{regular, + Akkusativ} \hfill Auxiliary Verb: \textbf{haben}

\vspace{5pt}
\hrule
\vspace{5pt}

\begin{tabular}{rl}
\multicolumn{2}{l}{\textbf{PRÄSENS}}\\
ich & \textbf{organisiere}\\
du & \textbf{organisierst}\\
er/sie/es & \textbf{organisiert}\\
wir & \textbf{organisieren}\\
ihr & \textbf{organisiert}\\
sie/Sie & \textbf{organisieren}\\
\end{tabular}
\hfill
\begin{tabular}{rl}
\multicolumn{2}{l}{\textbf{PRÄTERITUM}}\\
ich & \textbf{organisierte}\\
du & \textbf{organisiertest}\\
er/sie/es & \textbf{organisierte}\\
wir & \textbf{organisierten}\\
ihr & \textbf{organisiertet}\\
sie/Sie & \textbf{organisierten}\\
\end{tabular}
\hfill
\begin{tabular}{rl}
\multicolumn{2}{l}{\textbf{IMPERATIV}}\\
& \textbf{organisiere (du)!}\\
& \textbf{organisieren wir!}\\
& \textbf{organisiert (ihr)!}\\
& \textbf{organisieren Sie!}\\
\end{tabular}

\vspace{10pt}

\begin{tabular}{rl}
\multicolumn{2}{l}{\textbf{PERFEKT}}\\
& haben + \textbf{organisiert}\\
\end{tabular}
\hfill
\begin{tabular}{rl}
\multicolumn{2}{l}{\textbf{FUTURE I}}\\
& werden + \textbf{organisieren}\\
\end{tabular}
\hfill
\begin{tabular}{rl}
\multicolumn{2}{l}{\textbf{PLUSQUAMPERFEKT}}\\
& hatten + \textbf{organisiert}\\
\end{tabular}

\vspace{10pt}

\begin{tabular}{lll}
\multicolumn{3}{l}{\textbf{MIT PRÄPOSITIONEN}}\\
& \textbf{organisieren für} + Akk & (to organize for)\\
\end{tabular}
\hfill
\begin{tabular}{lll}
\multicolumn{3}{l}{\textbf{TRENNBARE VERBEN}}\\
& ---\\
\end{tabular}

\vspace{-5pt}
\end{verbbox-acc}

\subsection{P Verbs}

%packen (to pack)
\begin{verbbox}\addcontentsline{toc}{subsubsection}{packen (to pack)}
{\Large \bf packen} (to pack) \hfill Type: \textbf{regular} \hfill Auxiliary Verb: \textbf{haben}

\vspace{5pt}
\hrule
\vspace{5pt}

\begin{tabular}{rl}
\multicolumn{2}{l}{\textbf{PRÄSENS}}\\
ich & \textbf{packe}\\
du & \textbf{packst}\\
er/sie/es & \textbf{packt}\\
wir & \textbf{packen}\\
ihr & \textbf{packt}\\
sie/Sie & \textbf{packen}\\
\end{tabular}
\hfill
\begin{tabular}{rl}
\multicolumn{2}{l}{\textbf{PRÄTERITUM}}\\
ich & \textbf{packte}\\
du & \textbf{packtest}\\
er/sie/es & \textbf{packte}\\
wir & \textbf{packten}\\
ihr & \textbf{packtet}\\
sie/Sie & \textbf{packten}\\
\end{tabular}
\hfill
\begin{tabular}{rl}
\multicolumn{2}{l}{\textbf{IMPERATIV}}\\
& \textbf{pack(e) (du)!}\\
& \textbf{packen wir!}\\
& \textbf{packt (ihr)!}\\
& \textbf{packen Sie!}\\
\end{tabular}

\vspace{10pt}

\begin{tabular}{rl}
\multicolumn{2}{l}{\textbf{PERFEKT}}\\
& haben + \textbf{gepackt}\\
\end{tabular}
\hfill
\begin{tabular}{rl}
\multicolumn{2}{l}{\textbf{FUTURE I}}\\
& werden + \textbf{packen}\\
\end{tabular}
\hfill
\begin{tabular}{rl}
\multicolumn{2}{l}{\textbf{PLUSQUAMPERFEKT}}\\
& hatten + \textbf{gepackt}\\
\end{tabular}

\vspace{10pt}

\begin{tabular}{lll}
\multicolumn{3}{l}{\textbf{MIT PRÄPOSITIONEN}}\\
& \textbf{packen in} + Akk & (to pack into)\\
& \textbf{packen auf} + Akk & (to load onto)\\
\end{tabular}
\hfill
\begin{tabular}{lll}
\multicolumn{3}{l}{\textbf{TRENNBARE VERBEN}}\\
& \textbf{ein$|$packen} & (to pack up)\\
& \textbf{aus$|$packen} & (to unpack)\\
& \textbf{an$|$packen} & (to tackle / get started)\\
& \textbf{mit$|$packen} & (to help pack)\\
\end{tabular}
\vspace{-5pt}
\end{verbbox}

%parken - park
\begin{verbbox}\addcontentsline{toc}{subsubsection}{parken (to park)}
    {\Large \bf parken} (to park) \hfill Type: \textbf{regular} \hfill Auxiliary Verb: \textbf{haben} 
    
     \vspace{5pt}

    \hrule
    
    \vspace{5pt}

  \begin{tabular}{rl}
     \multicolumn{2}{l}{\textbf{PRÄSENS} }\\
    ich & \textbf{parke}\\
    du & \textbf{parkst}\\
    er/sie/es & \textbf{parkt}\\
    wir & \textbf{parken}\\
    ihr & \textbf{parkt}\\
    sie/Sie & \textbf{parken}\\
  \end{tabular}
  \hfill
     \begin{tabular}{rl}
    \multicolumn{2}{l}{\textbf{PRÄTERITUM} }\\
    ich & \textbf{parkte}\\
    du & \textbf{parktest}\\
    er/sie/es & \textbf{parkte}\\
    wir & \textbf{parkten}\\
    ihr & \textbf{parktet}\\
    sie/Sie & \textbf{parkten}\\
  \end{tabular}
  \hfill
  \begin{tabular}{rl}
     \multicolumn{2}{l}{\textbf{IMPERATIV} }\\
   & \textbf{park(e) (du)!}\\
   & \textbf{parken wir!}\\
   & \textbf{parket (ihr)!}\\
   & \textbf{parken Sie!}\\
   & \vspace{10pt}
  \end{tabular}

    \vspace{10pt}

 \begin{tabular}{rl}
     \multicolumn{2}{l}{\textbf{PERFEKT}}\\
  &  haben + \textbf{geparkt}\\
  \end{tabular}
\hfill
   \begin{tabular}{rl}
   
    \multicolumn{2}{l}{\textbf{FUTURE I} }\\
  &  werden + \textbf{parken}\\
 
  \end{tabular}
\hfill
   \begin{tabular}{rl}
      \multicolumn{2}{l}{\textbf{PLUSQUAMPERFEKT}}\\
   & hatten + \textbf{geparkt}\\
  \end{tabular}
  
  \vspace{10pt}

\begin{tabular}{lll}
\multicolumn{3}{l}{\textbf{MIT PRÄPOSITIONEN}}\\
& ---\\
\end{tabular}
\hfill
\begin{tabular}{lll}
\multicolumn{3}{l}{\textbf{TRENNBARE VERBEN}}\\
& ---\\
\end{tabular}

\end{verbbox}

%passen (to fit / to suit)
\begin{verbbox-dat}\addcontentsline{toc}{subsubsection}{passen (to fit / to suit)}
    {\Large \bf passen} (to fit / to suit) \hfill Type: \textbf{regular} \hfill Auxiliary Verb: \textbf{haben} 
    
     \vspace{5pt}

    \hrule
    
    \vspace{5pt}

  \begin{tabular}{rl}
     \multicolumn{2}{l}{\textbf{PRÄSENS} }\\
    ich & \textbf{passe}\\
    du & \textbf{passt}\\
    er/sie/es & \textbf{passt}\\
    wir & \textbf{passen}\\
    ihr & \textbf{passt}\\
    sie/Sie & \textbf{passen}\\
  \end{tabular}
  \hfill
     \begin{tabular}{rl}
    \multicolumn{2}{l}{\textbf{PRÄTERITUM} }\\
    ich & \textbf{passte}\\
    du & \textbf{passtest}\\
    er/sie/es & \textbf{passte}\\
    wir & \textbf{passten}\\
    ihr & \textbf{passtet}\\
    sie/Sie & \textbf{passten}\\
  \end{tabular}
  \hfill
  \begin{tabular}{rl}
     \multicolumn{2}{l}{\textbf{IMPERATIV} }\\
   & \textbf{pass(e) (du)!}\\
   & \textbf{passen wir!}\\
   & \textbf{passt (ihr)!}\\
   & \textbf{passen Sie!}\\
   & \vspace{10pt}
  \end{tabular}

    \vspace{10pt}

 \begin{tabular}{rl}
     \multicolumn{2}{l}{\textbf{PERFEKT}}\\
   & haben + \textbf{gepasst}\\
  \end{tabular}
\hfill
   \begin{tabular}{rl}
   
    \multicolumn{2}{l}{\textbf{FUTURE I} }\\
   & werden + \textbf{passen}\\
 
  \end{tabular}
\hfill
   \begin{tabular}{rl}
      \multicolumn{2}{l}{\textbf{PLUSQUAMPERFEKT}}\\
   & hatten + \textbf{gepasst}\\
  \end{tabular}

   \vspace{10pt}

   \begin{tabular}{lll}
      \multicolumn{3}{l}{\textbf{MIT PRÄPOSITIONEN}}\\
& \textbf{passen zu} + Dat & (to suit / go with)\\
& \textbf{passen in} + Akk & (to fit into)\\
& \textbf{passen auf} + Akk & (to watch / take care of)
\end{tabular}
  \hfill
   \begin{tabular}{lll}
      \multicolumn{3}{l}{\textbf{TRENNBARE VERBEN}}\\
 & \textbf{auf$|$passen} & (to pay attention / watch)\\
 & \textbf{an$|$passen} & (to adapt / adjust)\\
  \end{tabular}


  \vspace{-5pt}


\end{verbbox-dat}


%passieren (A2) – to happen
\begin{verbbox}\addcontentsline{toc}{subsubsection}{passieren (to happen)}
    {\Large \bf passieren} (to happen) \hfill Type: \textbf{irregular} \hfill Auxiliary Verb: \textbf{sein} 
    
     \vspace{5pt}

    \hrule
    
    \vspace{5pt}

  \begin{tabular}{rl}
     \multicolumn{2}{l}{\textbf{PRÄSENS} }\\
    ich & \textbf{-}\\
    du & \textbf{-}\\
    er/sie/es & \textbf{passiert}\\
    wir & \textbf{-}\\
    ihr & \textbf{-}\\
    sie/Sie & \textbf{passieren}\\
  \end{tabular}
  \hfill
     \begin{tabular}{rl}
    \multicolumn{2}{l}{\textbf{PRÄTERITUM} }\\
    ich & \textbf{-}\\
    du & \textbf{-}\\
    er/sie/es & \textbf{passierte}\\
    wir & \textbf{-}\\
    ihr & \textbf{-}\\
    sie/Sie & \textbf{passierten}\\
  \end{tabular}
  \hfill
  \begin{tabular}{rl}
     \multicolumn{2}{l}{\textbf{IMPERATIV} }\\
   & \textbf{-}\\
   & \textbf{-}\\
   & \textbf{-}\\
   & \textbf{-}\\
   & \vspace{10pt}
  \end{tabular}

    \vspace{10pt}

 \begin{tabular}{rl}
     \multicolumn{2}{l}{\textbf{PERFEKT}}\\
   & sein + \textbf{passiert}\\
  \end{tabular}
\hfill
   \begin{tabular}{rl}
   
    \multicolumn{2}{l}{\textbf{FUTURE I} }\\
   & werden + \textbf{passieren}\\
 
  \end{tabular}
\hfill
   \begin{tabular}{rl}
      \multicolumn{2}{l}{\textbf{PLUSQUAMPERFEKT}}\\
   & waren + \textbf{passiert}\\
  \end{tabular}
  
  \vspace{10pt}



\begin{tabular}{lll}
\multicolumn{3}{l}{\textbf{MIT PRÄPOSITIONEN}}\\
& ---\\
\end{tabular}
\hfill
\begin{tabular}{lll}
\multicolumn{3}{l}{\textbf{TRENNBARE VERBEN}}\\
& ---\\
\end{tabular}

\vspace{-5pt}

\end{verbbox}

%picknicken (to have a picnic)
\begin{verbbox}\addcontentsline{toc}{subsubsection}{picknicken (to have a picnic)}
{\Large \bf picknicken} (to have a picnic) \hfill Type: \textbf{regular} \hfill Auxiliary Verb: \textbf{haben}

\vspace{5pt}
\hrule
\vspace{5pt}

\begin{tabular}{rl}
\multicolumn{2}{l}{\textbf{PRÄSENS}}\\
ich & \textbf{picknicke}\\
du & \textbf{picknickst}\\
er/sie/es & \textbf{picknickt}\\
wir & \textbf{picknicken}\\
ihr & \textbf{picknickt}\\
sie/Sie & \textbf{picknicken}\\
\end{tabular}
\hfill
\begin{tabular}{rl}
\multicolumn{2}{l}{\textbf{PRÄTERITUM}}\\
ich & \textbf{picknickte}\\
du & \textbf{picknicktest}\\
er/sie/es & \textbf{picknickte}\\
wir & \textbf{picknickten}\\
ihr & \textbf{picknicktet}\\
sie/Sie & \textbf{picknickten}\\
\end{tabular}
\hfill
\begin{tabular}{rl}
\multicolumn{2}{l}{\textbf{IMPERATIV}}\\
& \textbf{picknicke (du)!}\\
& \textbf{picknicken wir!}\\
& \textbf{picknickt (ihr)!}\\
& \textbf{picknicken Sie!}\\
\end{tabular}

\vspace{10pt}

\begin{tabular}{rl}
\multicolumn{2}{l}{\textbf{PERFEKT}}\\
& haben + \textbf{gepicknickt}\\
\end{tabular}
\hfill
\begin{tabular}{rl}
\multicolumn{2}{l}{\textbf{FUTURE I}}\\
& werden + \textbf{picknicken}\\
\end{tabular}
\hfill
\begin{tabular}{rl}
\multicolumn{2}{l}{\textbf{PLUSQUAMPERFEKT}}\\
& hatten + \textbf{gepicknickt}\\
\end{tabular}

\vspace{10pt}

\begin{tabular}{lll}
\multicolumn{3}{l}{\textbf{MIT PRÄPOSITIONEN}}\\
& \textbf{picknicken im} + Dat & (to have a picnic in)\\
& \textbf{picknicken auf} + Dat & (to have a picnic on)\\
\end{tabular}
\hfill
\begin{tabular}{lll}
\multicolumn{3}{l}{\textbf{TRENNBARE VERBEN}}\\
& ---\\
\vspace{10pt}
\end{tabular}

\vspace{-5pt}
\end{verbbox}


%planen (A2) – to plan
\begin{verbbox-acc}\addcontentsline{toc}{subsubsection}{planen (to plan)}
    {\Large \bf planen} (to plan) \hfill Type: \textbf{regular, + Akkusativ} \hfill Auxiliary Verb: \textbf{haben} 
    
     \vspace{5pt}

    \hrule
    
    \vspace{5pt}

  \begin{tabular}{rl}
     \multicolumn{2}{l}{\textbf{PRÄSENS} }\\
    ich & \textbf{plane}\\
    du & \textbf{planst}\\
    er/sie/es & \textbf{plant}\\
    wir & \textbf{planen}\\
    ihr & \textbf{plant}\\
    sie/Sie & \textbf{planen}\\
  \end{tabular}
  \hfill
     \begin{tabular}{rl}
    \multicolumn{2}{l}{\textbf{PRÄTERITUM} }\\
    ich & \textbf{plante}\\
    du & \textbf{plantest}\\
    er/sie/es & \textbf{plante}\\
    wir & \textbf{planten}\\
    ihr & \textbf{plantet}\\
    sie/Sie & \textbf{planten}\\
  \end{tabular}
  \hfill
  \begin{tabular}{rl}
     \multicolumn{2}{l}{\textbf{IMPERATIV} }\\
   & \textbf{plan(e) (du)!}\\
   & \textbf{planen wir!}\\
   & \textbf{plant (ihr)!}\\
   & \textbf{planen Sie!}\\
   & \vspace{10pt}
  \end{tabular}

    \vspace{10pt}

 \begin{tabular}{rl}
     \multicolumn{2}{l}{\textbf{PERFEKT}}\\
   & haben + \textbf{geplant}\\
  \end{tabular}
\hfill
   \begin{tabular}{rl}
   
    \multicolumn{2}{l}{\textbf{FUTURE I} }\\
   & werden + \textbf{planen}\\
 
  \end{tabular}
\hfill
   \begin{tabular}{rl}
      \multicolumn{2}{l}{\textbf{PLUSQUAMPERFEKT}}\\
   & hatten + \textbf{geplant}\\
  \end{tabular}
  
  \vspace{-5pt}


\end{verbbox-acc}

%plaudern (to chat)
\begin{verbbox}\addcontentsline{toc}{subsubsection}{plaudern (to chat)}
{\Large \bf plaudern} (to chat) \hfill Type: \textbf{regular} \hfill Auxiliary Verb: \textbf{haben}

\vspace{5pt}
\hrule
\vspace{5pt}

\begin{tabular}{rl}
\multicolumn{2}{l}{\textbf{PRÄSENS}}\\
ich & \textbf{plaudere}\\
du & \textbf{plauderst}\\
er/sie/es & \textbf{plaudert}\\
wir & \textbf{plaudern}\\
ihr & \textbf{plaudert}\\
sie/Sie & \textbf{plaudern}\\
\end{tabular}
\hfill
\begin{tabular}{rl}
\multicolumn{2}{l}{\textbf{PRÄTERITUM}}\\
ich & \textbf{plauderte}\\
du & \textbf{plaudertest}\\
er/sie/es & \textbf{plauderte}\\
wir & \textbf{plauderten}\\
ihr & \textbf{plaudertet}\\
sie/Sie & \textbf{plauderten}\\
\end{tabular}
\hfill
\begin{tabular}{rl}
\multicolumn{2}{l}{\textbf{IMPERATIV}}\\
& \textbf{plaudere (du)!}\\
& \textbf{plaudern wir!}\\
& \textbf{plaudert (ihr)!}\\
& \textbf{plaudern Sie!}\\
\end{tabular}

\vspace{10pt}

\begin{tabular}{rl}
\multicolumn{2}{l}{\textbf{PERFEKT}}\\
& haben + \textbf{geplaudert}\\
\end{tabular}
\hfill
\begin{tabular}{rl}
\multicolumn{2}{l}{\textbf{FUTURE I}}\\
& werden + \textbf{plaudern}\\
\end{tabular}
\hfill
\begin{tabular}{rl}
\multicolumn{2}{l}{\textbf{PLUSQUAMPERFEKT}}\\
& hatten + \textbf{geplaudert}\\
\end{tabular}

\vspace{10pt}

\begin{tabular}{lll}
\multicolumn{3}{l}{\textbf{MIT PRÄPOSITIONEN}}\\
& \textbf{plaudern mit} + Dat & (to chat with)\\
& \textbf{plaudern über} + Akk & (to chat about)\\
\end{tabular}
\hfill
\begin{tabular}{lll}
\multicolumn{3}{l}{\textbf{TRENNBARE VERBEN}}\\
& aus$|$plaudern & (to let slip / blurt out)\\
\vspace{10pt}
\end{tabular}

\vspace{-5pt}
\end{verbbox}


%posten (to post)
\begin{verbbox}\addcontentsline{toc}{subsubsection}{posten (to post)}
{\Large \bf posten} (to post) \hfill Type: \textbf{regular} \hfill Auxiliary Verb: \textbf{haben} 

\vspace{5pt}
\hrule
\vspace{5pt}

\begin{tabular}{rl}
\multicolumn{2}{l}{\textbf{PRÄSENS}}\\
ich & \textbf{poste}\\
du & \textbf{postest}\\
er/sie/es & \textbf{postet}\\
wir & \textbf{posten}\\
ihr & \textbf{postet}\\
sie/Sie & \textbf{posten}\\
\end{tabular}
\hfill
\begin{tabular}{rl}
\multicolumn{2}{l}{\textbf{PRÄTERITUM}}\\
ich & \textbf{postete}\\
du & \textbf{postetest}\\
er/sie/es & \textbf{postete}\\
wir & \textbf{posteten}\\
ihr & \textbf{postetet}\\
sie/Sie & \textbf{posteten}\\
\end{tabular}
\hfill
\begin{tabular}{rl}
\multicolumn{2}{l}{\textbf{IMPERATIV}}\\
& \textbf{poste (du)!}\\
& \textbf{posten wir!}\\
& \textbf{postet (ihr)!}\\
& \textbf{posten Sie!}\\
\end{tabular}

\vspace{10pt}

\begin{tabular}{rl}
\multicolumn{2}{l}{\textbf{PERFEKT}}\\
& haben + \textbf{gepostet}\\
\end{tabular}
\hfill
\begin{tabular}{rl}
\multicolumn{2}{l}{\textbf{FUTURE I}}\\
& werden + \textbf{posten}\\
\end{tabular}
\hfill
\begin{tabular}{rl}
\multicolumn{2}{l}{\textbf{PLUSQUAMPERFEKT}}\\
& hatten + \textbf{gepostet}\\
\end{tabular}

\vspace{10pt}

\begin{tabular}{lll}
\multicolumn{3}{l}{\textbf{MIT PRÄPOSITIONEN}}\\
& \textbf{posten auf} + Dat & (to post on)\\
\end{tabular}
\hfill
\begin{tabular}{lll}
\multicolumn{3}{l}{\textbf{TRENNBARE VERBEN}}\\
& ---\\
\end{tabular}

\vspace{-5pt}
\end{verbbox}




%probieren (A2) – to try
\begin{verbbox-acc}\addcontentsline{toc}{subsubsection}{probieren (to try)}
    {\Large \bf probieren} (to try) \hfill Type: \textbf{irregular, + Akkusativ} \hfill Auxiliary Verb: \textbf{haben} 
    
     \vspace{5pt}

    \hrule
    
    \vspace{5pt}

  \begin{tabular}{rl}
     \multicolumn{2}{l}{\textbf{PRÄSENS} }\\
    ich & \textbf{probiere}\\
    du & \textbf{probierst}\\
    er/sie/es & \textbf{probiert}\\
    wir & \textbf{probieren}\\
    ihr & \textbf{probiert}\\
    sie/Sie & \textbf{probieren}\\
  \end{tabular}
  \hfill
     \begin{tabular}{rl}
    \multicolumn{2}{l}{\textbf{PRÄTERITUM} }\\
    ich & \textbf{probierte}\\
    du & \textbf{probiertest}\\
    er/sie/es & \textbf{probierte}\\
    wir & \textbf{probierten}\\
    ihr & \textbf{probiertet}\\
    sie/Sie & \textbf{probierten}\\
  \end{tabular}
  \hfill
  \begin{tabular}{rl}
     \multicolumn{2}{l}{\textbf{IMPERATIV} }\\
   & \textbf{probier(e) (du)!}\\
   & \textbf{gprobieren wir!}\\
   & \textbf{probiert (ihr)!}\\
   & \textbf{probieren Sie!}\\
   & \vspace{10pt}
  \end{tabular}

    \vspace{10pt}

 \begin{tabular}{rl}
     \multicolumn{2}{l}{\textbf{PERFEKT}}\\
   & haben + \textbf{probiert}\\
  \end{tabular}
\hfill
   \begin{tabular}{rl}
   
    \multicolumn{2}{l}{\textbf{FUTURE I} }\\
   & werden + \textbf{probieren}\\
 
  \end{tabular}
\hfill
   \begin{tabular}{rl}
      \multicolumn{2}{l}{\textbf{PLUSQUAMPERFEKT}}\\
   & hatten + \textbf{probiert}\\
  \end{tabular}
  
  \vspace{-5pt}


\end{verbbox-acc}

%prüfen (to check / examine)
\begin{verbbox}\addcontentsline{toc}{subsubsection}{prüfen (to check / examine)}
{\Large \bf prüfen} (to check / examine) \hfill Type: \textbf{regular} \hfill Auxiliary Verb: \textbf{haben}

\vspace{5pt}
\hrule
\vspace{5pt}

\begin{tabular}{rl}
\multicolumn{2}{l}{\textbf{PRÄSENS}}\\
ich & \textbf{prüfe}\\
du & \textbf{prüfst}\\
er/sie/es & \textbf{prüft}\\
wir & \textbf{prüfen}\\
ihr & \textbf{prüft}\\
sie/Sie & \textbf{prüfen}\\
\end{tabular}
\hfill
\begin{tabular}{rl}
\multicolumn{2}{l}{\textbf{PRÄTERITUM}}\\
ich & \textbf{prüfte}\\
du & \textbf{prüftest}\\
er/sie/es & \textbf{prüfte}\\
wir & \textbf{prüften}\\
ihr & \textbf{prüftet}\\
sie/Sie & \textbf{prüften}\\
\end{tabular}
\hfill
\begin{tabular}{rl}
\multicolumn{2}{l}{\textbf{IMPERATIV}}\\
& \textbf{prüfe (du)!}\\
& \textbf{prüfen wir!}\\
& \textbf{prüft (ihr)!}\\
& \textbf{prüfen Sie!}\\
\end{tabular}

\vspace{10pt}

\begin{tabular}{rl}
\multicolumn{2}{l}{\textbf{PERFEKT}}\\
& haben + \textbf{geprüft}\\
\end{tabular}
\hfill
\begin{tabular}{rl}
\multicolumn{2}{l}{\textbf{FUTURE I}}\\
& werden + \textbf{prüfen}\\
\end{tabular}
\hfill
\begin{tabular}{rl}
\multicolumn{2}{l}{\textbf{PLUSQUAMPERFEKT}}\\
& hatten + \textbf{geprüft}\\
\end{tabular}

\vspace{10pt}

\begin{tabular}{lll}
\multicolumn{3}{l}{\textbf{MIT PRÄPOSITIONEN}}\\
& \textbf{prüfen auf} + Akk & (to check for)\\
\end{tabular}
\hfill
\begin{tabular}{lll}
\multicolumn{3}{l}{\textbf{TRENNBARE VERBEN}}\\
& ---\\
\end{tabular}
\vspace{-5pt}
\end{verbbox}

%putzen (A1) – to clean
\begin{verbbox-acc}\addcontentsline{toc}{subsubsection}{putzen (to clean)}
    {\Large \bf putzen} (to clean) \hfill Type: \textbf{regular, + Akkusativ} \hfill Auxiliary Verb: \textbf{haben} 
    
     \vspace{5pt}

    \hrule
    
    \vspace{5pt}

  \begin{tabular}{rl}
     \multicolumn{2}{l}{\textbf{PRÄSENS} }\\
    ich & \textbf{putze}\\
    du & \textbf{putzt}\\
    er/sie/es & \textbf{putzt}\\
    wir & \textbf{putzen}\\
    ihr & \textbf{putzt}\\
    sie/Sie & \textbf{putzen}\\
  \end{tabular}
  \hfill
     \begin{tabular}{rl}
    \multicolumn{2}{l}{\textbf{PRÄTERITUM} }\\
    ich & \textbf{putzte}\\
    du & \textbf{putztest}\\
    er/sie/es & \textbf{putzte}\\
    wir & \textbf{putzten}\\
    ihr & \textbf{putztet}\\
    sie/Sie & \textbf{putzten}\\
  \end{tabular}
  \hfill
  \begin{tabular}{rl}
     \multicolumn{2}{l}{\textbf{IMPERATIV} }\\
   & \textbf{putz(e) (du)!}\\
   & \textbf{putzen wir!}\\
   & \textbf{putzt (ihr)!}\\
   & \textbf{putzen Sie!}\\
   & \vspace{10pt}
  \end{tabular}

    \vspace{10pt}

 \begin{tabular}{rl}
     \multicolumn{2}{l}{\textbf{PERFEKT}}\\
 &  haben + \textbf{geputzt}\\
  \end{tabular}
\hfill
   \begin{tabular}{rl}
   
    \multicolumn{2}{l}{\textbf{FUTURE I} }\\
  &  werden + \textbf{putzen}\\
 
  \end{tabular}
\hfill
   \begin{tabular}{rl}
      \multicolumn{2}{l}{\textbf{PLUSQUAMPERFEKT}}\\
   & hatten + \textbf{geputzt}\\
  \end{tabular}
  
  \vspace{-5pt}


\end{verbbox-acc}

\subsection{Q Verbs}

%qualifizieren  - to qualify
\begin{verbbox-acc}\addcontentsline{toc}{subsubsection}{qualifizieren (to qualify)}
{\Large \bf qualifizieren} (to qualify) \hfill Type: \textbf{regular, + Akkusativ} \hfill Auxiliary Verb: \textbf{haben}

\vspace{5pt}
\hrule
\vspace{5pt}

\begin{tabular}{rl}
\multicolumn{2}{l}{\textbf{PRÄSENS}}\\
ich & \textbf{qualifiziere}\\
du & \textbf{qualifizierst}\\
er/sie/es & \textbf{qualifiziert}\\
wir & \textbf{qualifizieren}\\
ihr & \textbf{qualifiziert}\\
sie/Sie & \textbf{qualifizieren}\\
\end{tabular}
\hfill
\begin{tabular}{rl}
\multicolumn{2}{l}{\textbf{PRÄTERITUM}}\\
ich & \textbf{qualifizierte}\\
du & \textbf{qualifiziertest}\\
er/sie/es & \textbf{qualifizierte}\\
wir & \textbf{qualifizierten}\\
ihr & \textbf{qualifiziertet}\\
sie/Sie & \textbf{qualifizierten}\\
\end{tabular}
\hfill
\begin{tabular}{rl}
\multicolumn{2}{l}{\textbf{IMPERATIV}}\\
& \textbf{qualifizier(e) (du)!}\\
& \textbf{qualifizieren wir!}\\
& \textbf{qualifiziert (ihr)!}\\
& \textbf{qualifizieren Sie!}\\
\end{tabular}

\vspace{10pt}

\begin{tabular}{rl}
\multicolumn{2}{l}{\textbf{PERFEKT}}\\
& haben + \textbf{qualifiziert}\\
\end{tabular}
\hfill
\begin{tabular}{rl}
\multicolumn{2}{l}{\textbf{FUTURE I}}\\
& werden + \textbf{qualifizieren}\\
\end{tabular}
\hfill
\begin{tabular}{rl}
\multicolumn{2}{l}{\textbf{PLUSQUAMPERFEKT}}\\
& hatten + \textbf{qualifiziert}\\
\end{tabular}

\vspace{10pt}

\begin{tabular}{lll}
\multicolumn{3}{l}{\textbf{MIT PRÄPOSITIONEN}}\\
& \textbf{sich qualifizieren für} + Akk & (to qualify for)\\
\end{tabular}
\hfill
\begin{tabular}{lll}
\multicolumn{3}{l}{\textbf{TRENNBARE VERBEN}}\\
& ---\\
\end{tabular}

\vspace{-5pt}
\end{verbbox-acc}

\subsection{R Verbs}

%raten (to advise / to guess)
\begin{verbbox-dat}\addcontentsline{toc}{subsubsection}{raten (to advise / to guess)}
    {\Large \bf raten} (to advise / to guess) \hfill Type: \textbf{irregular, + Dativ} \hfill Auxiliary Verb: \textbf{haben} 
    
     \vspace{5pt}

    \hrule
    
    \vspace{5pt}

  \begin{tabular}{rl}
     \multicolumn{2}{l}{\textbf{PRÄSENS} }\\
    ich & \textbf{rate}\\
    du & \textbf{rätst}\\
    er/sie/es & \textbf{rät}\\
    wir & \textbf{raten}\\
    ihr & \textbf{ratet}\\
    sie/Sie & \textbf{raten}\\
  \end{tabular}
  \hfill
     \begin{tabular}{rl}
    \multicolumn{2}{l}{\textbf{PRÄTERITUM} }\\
    ich & \textbf{riet}\\
    du & \textbf{rietest}\\
    er/sie/es & \textbf{riet}\\
    wir & \textbf{rieten}\\
    ihr & \textbf{rietet}\\
    sie/Sie & \textbf{rieten}\\
  \end{tabular}
  \hfill
  \begin{tabular}{rl}
     \multicolumn{2}{l}{\textbf{IMPERATIV} }\\
   & \textbf{rat(e) (du)!}\\
   & \textbf{raten wir!}\\
   & \textbf{ratet (ihr)!}\\
   & \textbf{raten Sie!}\\
   & \vspace{10pt}
  \end{tabular}

    \vspace{10pt}

 \begin{tabular}{rl}
     \multicolumn{2}{l}{\textbf{PERFEKT}}\\
   & haben + \textbf{geraten}\\
  \end{tabular}
\hfill
   \begin{tabular}{rl}
   
    \multicolumn{2}{l}{\textbf{FUTURE I} }\\
   & werden + \textbf{raten}\\
 
  \end{tabular}
\hfill
   \begin{tabular}{rl}
      \multicolumn{2}{l}{\textbf{PLUSQUAMPERFEKT}}\\
   & hatten + \textbf{geraten}\\
  \end{tabular}

   \vspace{10pt}

   \begin{tabular}{lll}
      \multicolumn{3}{l}{\textbf{MIT PRÄPOSITIONEN}}\\
& \textbf{raten zu} + Dat & (to advise to do / recommend)\\
& \textbf{raten von} + Dat & (to advise against)\\
\end{tabular}
  \hfill
   \begin{tabular}{lll}
      \multicolumn{3}{l}{\textbf{TRENNBARE VERBEN}}\\
 & \textbf{ab$|$raten} & (to advise against)\\
 & \textbf{zu$|$raten} & (to advise to do)\\
  \end{tabular}


  \vspace{-5pt}


\end{verbbox-dat}

%reagieren (B1) – to react
\begin{verbbox-acc}\addcontentsline{toc}{subsubsection}{reagieren (to react)}
    {\Large \bf reagieren} (to react) \hfill Type: \textbf{irregular, + Akkusativ} \hfill Auxiliary Verb: \textbf{haben} 
    
     \vspace{5pt}

    \hrule
    
    \vspace{5pt}

  \begin{tabular}{rl}
     \multicolumn{2}{l}{\textbf{PRÄSENS} }\\
    ich & \textbf{reagiere}\\
    du & \textbf{greagierst}\\
    er/sie/es & \textbf{reagiert}\\
    wir & \textbf{reagieren}\\
    ihr & \textbf{reagiert}\\
    sie/Sie & \textbf{greagieren}\\
  \end{tabular}
  \hfill
     \begin{tabular}{rl}
    \multicolumn{2}{l}{\textbf{PRÄTERITUM} }\\
    ich & \textbf{reagierte}\\
    du & \textbf{reagiertest}\\
    er/sie/es & \textbf{reagierte}\\
    wir & \textbf{reagierten}\\
    ihr & \textbf{reagiertet}\\
    sie/Sie & \textbf{reagierten}\\
  \end{tabular}
  \hfill
  \begin{tabular}{rl}
     \multicolumn{2}{l}{\textbf{IMPERATIV} }\\
   & \textbf{reagier(e) (du)!}\\
   & \textbf{reagieren wir!}\\
   & \textbf{reagiert (ihr)!}\\
   & \textbf{reagieren Sie!}\\
   & \vspace{10pt}
  \end{tabular}

    \vspace{10pt}

 \begin{tabular}{rl}
     \multicolumn{2}{l}{\textbf{PERFEKT}}\\
  &  haben + \textbf{reagiert}\\
  \end{tabular}
\hfill
   \begin{tabular}{rl}
   
    \multicolumn{2}{l}{\textbf{FUTURE I} }\\
  &  werden + \textbf{reagieren}\\
 
  \end{tabular}
\hfill
   \begin{tabular}{rl}
      \multicolumn{2}{l}{\textbf{PLUSQUAMPERFEKT}}\\
  &  hatten + \textbf{reagiert}\\
  \end{tabular}
  
  \vspace{-5pt}


\end{verbbox-acc}
%reden (A1) – to talk
\begin{verbbox}\addcontentsline{toc}{subsubsection}{reden (to talk)}
    {\Large \bf reden} (to talk) \hfill Type: \textbf{regular} \hfill Auxiliary Verb: \textbf{haben} 
    
     \vspace{5pt}

    \hrule
    
    \vspace{5pt}

  \begin{tabular}{rl}
     \multicolumn{2}{l}{\textbf{PRÄSENS} }\\
    ich & \textbf{rede}\\
    du & \textbf{redest}\\
    er/sie/es & \textbf{redet}\\
    wir & \textbf{reden}\\
    ihr & \textbf{redet}\\
    sie/Sie & \textbf{reden}\\
  \end{tabular}
  \hfill
     \begin{tabular}{rl}
    \multicolumn{2}{l}{\textbf{PRÄTERITUM} }\\
    ich & \textbf{redete}\\
    du & \textbf{redetest}\\
    er/sie/es & \textbf{redete}\\
    wir & \textbf{redeten}\\
    ihr & \textbf{redetet}\\
    sie/Sie & \textbf{redeten}\\
  \end{tabular}
  \hfill
  \begin{tabular}{rl}
     \multicolumn{2}{l}{\textbf{IMPERATIV} }\\
   & \textbf{red(e) (du)!}\\
   & \textbf{reden wir!}\\
   & \textbf{redet (ihr)!}\\
   & \textbf{reden Sie!}\\
   & \vspace{10pt}
  \end{tabular}

    \vspace{10pt}

 \begin{tabular}{rl}
     \multicolumn{2}{l}{\textbf{PERFEKT}}\\
  &  haben + \textbf{geredet}\\
  \end{tabular}
\hfill
   \begin{tabular}{rl}
   
    \multicolumn{2}{l}{\textbf{FUTURE I} }\\
  &  werden + \textbf{reden}\\
 
  \end{tabular}
\hfill
   \begin{tabular}{rl}
      \multicolumn{2}{l}{\textbf{PLUSQUAMPERFEKT}}\\
  &  hatten + \textbf{geredet}\\
  \end{tabular}
  
  \vspace{10pt}



\begin{tabular}{lll}
\multicolumn{3}{l}{\textbf{MIT PRÄPOSITIONEN}}\\
& ---\\
\end{tabular}
\hfill
\begin{tabular}{lll}
\multicolumn{3}{l}{\textbf{TRENNBARE VERBEN}}\\
& ---\\
\end{tabular}

  \vspace{-5pt}

\end{verbbox}

%reduzieren (to reduce)
\begin{verbbox-acc}\addcontentsline{toc}{subsubsection}{reduzieren (to reduce)}
{\Large \bf reduzieren} (to reduce) \hfill Type: \textbf{regular, + Akkusativ} \hfill Auxiliary Verb: \textbf{haben}

\vspace{5pt}
\hrule
\vspace{5pt}

\begin{tabular}{rl}
\multicolumn{2}{l}{\textbf{PRÄSENS}}\\
ich & \textbf{reduziere}\\
du & \textbf{reduzierst}\\
er/sie/es & \textbf{reduziert}\\
wir & \textbf{reduzieren}\\
ihr & \textbf{reduziert}\\
sie/Sie & \textbf{reduzieren}\\
\end{tabular}
\hfill
\begin{tabular}{rl}
\multicolumn{2}{l}{\textbf{PRÄTERITUM}}\\
ich & \textbf{reduzierte}\\
du & \textbf{reduziertest}\\
er/sie/es & \textbf{reduzierte}\\
wir & \textbf{reduzierten}\\
ihr & \textbf{reduziertet}\\
sie/Sie & \textbf{reduzierten}\\
\end{tabular}
\hfill
\begin{tabular}{rl}
\multicolumn{2}{l}{\textbf{IMPERATIV}}\\
& \textbf{reduziere (du)!}\\
& \textbf{reduzieren wir!}\\
& \textbf{reduziert (ihr)!}\\
& \textbf{reduzieren Sie!}\\
\end{tabular}

\vspace{10pt}

\begin{tabular}{rl}
\multicolumn{2}{l}{\textbf{PERFEKT}}\\
& haben + \textbf{reduziert}\\
\end{tabular}
\hfill
\begin{tabular}{rl}
\multicolumn{2}{l}{\textbf{FUTURE I}}\\
& werden + \textbf{reduzieren}\\
\end{tabular}
\hfill
\begin{tabular}{rl}
\multicolumn{2}{l}{\textbf{PLUSQUAMPERFEKT}}\\
& hatten + \textbf{reduziert}\\
\end{tabular}

\vspace{10pt}

\begin{tabular}{lll}
\multicolumn{3}{l}{\textbf{MIT PRÄPOSITIONEN}}\\
& \textbf{reduzieren} + Akk & (to reduce something)\\
& \textbf{reduzieren auf} + Akk & (to reduce to)\\
& \textbf{reduzieren um} + Akk & (to reduce by)\\
\end{tabular}
\hfill
\begin{tabular}{lll}
\multicolumn{3}{l}{\textbf{TRENNBARE VERBEN}}\\
& ---\\
\vspace{20pt}
\end{tabular}

\vspace{-5pt}
\end{verbbox-acc}

%registrieren (to register)
\begin{verbbox}\addcontentsline{toc}{subsubsection}{registrieren (to register)}
    {\Large \bf registrieren} (to register) \hfill Type: \textbf{regular, -ieren verb} \hfill Auxiliary Verb: \textbf{haben} 
    
     \vspace{5pt}

    \hrule
    
    \vspace{5pt}

  \begin{tabular}{rl}
     \multicolumn{2}{l}{\textbf{PRÄSENS} }\\
    ich & \textbf{registriere}\\
    du & \textbf{registrierst}\\
    er/sie/es & \textbf{registriert}\\
    wir & \textbf{registrieren}\\
    ihr & \textbf{registriert}\\
    sie/Sie & \textbf{registrieren}\\
  \end{tabular}
  \hfill
     \begin{tabular}{rl}
    \multicolumn{2}{l}{\textbf{PRÄTERITUM} }\\
    ich & \textbf{registrierte}\\
    du & \textbf{registriertest}\\
    er/sie/es & \textbf{registrierte}\\
    wir & \textbf{registrierten}\\
    ihr & \textbf{registriertet}\\
    sie/Sie & \textbf{registrierten}\\
  \end{tabular}
  \hfill
  \begin{tabular}{rl}
     \multicolumn{2}{l}{\textbf{IMPERATIV} }\\
   & \textbf{registriere (du)!}\\
   & \textbf{registrieren wir!}\\
   & \textbf{registriert (ihr)!}\\
   & \textbf{registrieren Sie!}\\
   & \vspace{10pt}
  \end{tabular}

    \vspace{10pt}

 \begin{tabular}{rl}
     \multicolumn{2}{l}{\textbf{PERFEKT}}\\
   & haben + \textbf{registriert}\\
  \end{tabular}
\hfill
   \begin{tabular}{rl}
    \multicolumn{2}{l}{\textbf{FUTURE I} }\\
   & werden + \textbf{registrieren}\\
  \end{tabular}
\hfill
   \begin{tabular}{rl}
      \multicolumn{2}{l}{\textbf{PLUSQUAMPERFEKT}}\\
   & hatten + \textbf{registriert}\\
  \end{tabular}

   \vspace{10pt}

   \begin{tabular}{lll}
      \multicolumn{3}{l}{\textbf{MIT PRÄPOSITIONEN}}\\
& \textbf{sich registrieren bei} + Dat & (to register with)\\
& \textbf{sich registrieren für} + Akk & (to register for)\\
\end{tabular}
  \hfill
   \begin{tabular}{lll}
      \multicolumn{3}{l}{\textbf{TRENNBARE VERBEN}}\\
 & \textbf{---} & (none common)\\
\vspace{10pt}
  \end{tabular}

  \vspace{-5pt}

\end{verbbox}

%regeln (to regulate / manage)
\begin{verbbox}\addcontentsline{toc}{subsubsection}{regeln (to regulate / manage)}
{\Large \bf regeln} (to regulate / manage) \hfill Type: \textbf{regular} \hfill Auxiliary Verb: \textbf{haben} 

\vspace{5pt}
\hrule
\vspace{5pt}

\begin{tabular}{rl}
\multicolumn{2}{l}{\textbf{PRÄSENS}}\\
ich & \textbf{regle}\\
du & \textbf{regelst}\\
er/sie/es & \textbf{regelt}\\
wir & \textbf{regeln}\\
ihr & \textbf{regelt}\\
sie/Sie & \textbf{regeln}\\
\end{tabular}
\hfill
\begin{tabular}{rl}
\multicolumn{2}{l}{\textbf{PRÄTERITUM}}\\
ich & \textbf{regelte}\\
du & \textbf{regeltest}\\
er/sie/es & \textbf{regelte}\\
wir & \textbf{regelten}\\
ihr & \textbf{regeltet}\\
sie/Sie & \textbf{regelten}\\
\end{tabular}
\hfill
\begin{tabular}{rl}
\multicolumn{2}{l}{\textbf{IMPERATIV}}\\
& \textbf{regle (du)!}\\
& \textbf{regeln wir!}\\
& \textbf{regelt (ihr)!}\\
& \textbf{regeln Sie!}\\
\end{tabular}

\vspace{10pt}

\begin{tabular}{rl}
\multicolumn{2}{l}{\textbf{PERFEKT}}\\
& haben + \textbf{geregelt}\\
\end{tabular}
\hfill
\begin{tabular}{rl}
\multicolumn{2}{l}{\textbf{FUTURE I}}\\
& werden + \textbf{regeln}\\
\end{tabular}
\hfill
\begin{tabular}{rl}
\multicolumn{2}{l}{\textbf{PLUSQUAMPERFEKT}}\\
& hatten + \textbf{geregelt}\\
\end{tabular}

\vspace{10pt}

\begin{tabular}{lll}
\multicolumn{3}{l}{\textbf{MIT PRÄPOSITIONEN}}\\
& \textbf{regeln} + Akk & (to manage something)\\
\end{tabular}
\hfill
\begin{tabular}{lll}
\multicolumn{3}{l}{\textbf{TRENNBARE VERBEN}}\\
& ---\\
\end{tabular}

\vspace{-5pt}
\end{verbbox}



%regnen (A1) – to rain
\begin{verbbox}\addcontentsline{toc}{subsubsection}{regnen (to rain)}
    {\Large \bf regnen} (to rain) \hfill Type: \textbf{regular} \hfill Auxiliary Verb: \textbf{haben} 
    
     \vspace{5pt}

    \hrule
    
    \vspace{5pt}

  \begin{tabular}{rl}
     \multicolumn{2}{l}{\textbf{PRÄSENS} }\\
    ich & \textbf{-}\\
    du & \textbf{-}\\
    es & \textbf{regnet}\\
    wir & \textbf{-}\\
    ihr & \textbf{-}\\
    sie/Sie & \textbf{-}\\
  \end{tabular}
  \hfill
     \begin{tabular}{rl}
    \multicolumn{2}{l}{\textbf{PRÄTERITUM} }\\
    ich & \textbf{-}\\
    du & \textbf{-}\\
    es & \textbf{regnete}\\
    wir & \textbf{-}\\
    ihr & \textbf{-}\\
    sie/Sie & \textbf{-}\\
  \end{tabular}
  \hfill
  \begin{tabular}{rl}
     \multicolumn{2}{l}{\textbf{IMPERATIV} }\\
   & \textbf{-}\\
   & \textbf{-}\\
   & \textbf{-}\\
   & \textbf{-}\\
   & \vspace{10pt}
  \end{tabular}

    \vspace{10pt}

 \begin{tabular}{rl}
     \multicolumn{2}{l}{\textbf{PERFEKT}}\\
  &  haben + \textbf{geredet}\\
  \end{tabular}
\hfill
   \begin{tabular}{rl}
   
    \multicolumn{2}{l}{\textbf{FUTURE I} }\\
  &  werden + \textbf{regnen}\\
 
  \end{tabular}
\hfill
   \begin{tabular}{rl}
      \multicolumn{2}{l}{\textbf{PLUSQUAMPERFEKT}}\\
  &  hatten + \textbf{geregnet}\\
  \end{tabular}
  
  \vspace{10pt}



\begin{tabular}{lll}
\multicolumn{3}{l}{\textbf{MIT PRÄPOSITIONEN}}\\
& ---\\
\end{tabular}
\hfill
\begin{tabular}{lll}
\multicolumn{3}{l}{\textbf{TRENNBARE VERBEN}}\\
& ---\\
\end{tabular}

  \vspace{-5pt}

\end{verbbox}

%reichen (to be enough / to hand)
\begin{verbbox}\addcontentsline{toc}{subsubsection}{reichen (to be enough / to hand)}
{\Large \bf reichen} (to be enough / to hand) \hfill Type: \textbf{regular} \hfill Auxiliary Verb: \textbf{haben} 

\vspace{5pt}
\hrule
\vspace{5pt}

\begin{tabular}{rl}
\multicolumn{2}{l}{\textbf{PRÄSENS}}\\
ich & \textbf{reiche}\\
du & \textbf{reichst}\\
er/sie/es & \textbf{reicht}\\
wir & \textbf{reichen}\\
ihr & \textbf{reicht}\\
sie/Sie & \textbf{reichen}\\
\end{tabular}
\hfill
\begin{tabular}{rl}
\multicolumn{2}{l}{\textbf{PRÄTERITUM}}\\
ich & \textbf{reichte}\\
du & \textbf{reichtest}\\
er/sie/es & \textbf{reichte}\\
wir & \textbf{reichten}\\
ihr & \textbf{reichtet}\\
sie/Sie & \textbf{reichten}\\
\end{tabular}
\hfill
\begin{tabular}{rl}
\multicolumn{2}{l}{\textbf{IMPERATIV}}\\
& \textbf{reiche (du)!}\\
& \textbf{reichen wir!}\\
& \textbf{reicht (ihr)!}\\
& \textbf{reichen Sie!}\\
\end{tabular}

\vspace{10pt}

\begin{tabular}{rl}
\multicolumn{2}{l}{\textbf{PERFEKT}}\\
& haben + \textbf{gereicht}\\
\end{tabular}
\hfill
\begin{tabular}{rl}
\multicolumn{2}{l}{\textbf{FUTURE I}}\\
& werden + \textbf{reichen}\\
\end{tabular}
\hfill
\begin{tabular}{rl}
\multicolumn{2}{l}{\textbf{PLUSQUAMPERFEKT}}\\
& hatten + \textbf{gereicht}\\
\end{tabular}

\vspace{10pt}

\begin{tabular}{lll}
\multicolumn{3}{l}{\textbf{MIT PRÄPOSITIONEN}}\\
& \textbf{reichen für} + Akk & (to be enough for)\\
& \textbf{reichen an} + Akk & (to hand to)\\
\end{tabular}
\hfill
\begin{tabular}{lll}
\multicolumn{3}{l}{\textbf{TRENNBARE VERBEN}}\\
& \textbf{ein$|$reichen} & (to submit)\\
& \textbf{nach$|$reichen} & (to hand later)\\
\end{tabular}

\vspace{-5pt}
\end{verbbox}


%reinigen (to clean)
\begin{verbbox}\addcontentsline{toc}{subsubsection}{reinigen (to clean)}
{\Large \bf reinigen} (to clean) \hfill Type: \textbf{regular} \hfill Auxiliary Verb: \textbf{haben}

\vspace{5pt}
\hrule
\vspace{5pt}

\begin{tabular}{rl}
\multicolumn{2}{l}{\textbf{PRÄSENS}}\\
ich & \textbf{reinige}\\
du & \textbf{reinigst}\\
er/sie/es & \textbf{reinigt}\\
wir & \textbf{reinigen}\\
ihr & \textbf{reinigt}\\
sie/Sie & \textbf{reinigen}\\
\end{tabular}
\hfill
\begin{tabular}{rl}
\multicolumn{2}{l}{\textbf{PRÄTERITUM}}\\
ich & \textbf{reinigte}\\
du & \textbf{reinigtest}\\
er/sie/es & \textbf{reinigte}\\
wir & \textbf{reinigten}\\
ihr & \textbf{reinigtet}\\
sie/Sie & \textbf{reinigten}\\
\end{tabular}
\hfill
\begin{tabular}{rl}
\multicolumn{2}{l}{\textbf{IMPERATIV}}\\
& \textbf{reinige (du)!}\\
& \textbf{reinigen wir!}\\
& \textbf{reinigt (ihr)!}\\
& \textbf{reinigen Sie!}\\
\end{tabular}

\vspace{10pt}

\begin{tabular}{rl}
\multicolumn{2}{l}{\textbf{PERFEKT}}\\
& haben + \textbf{gereinigt}\\
\end{tabular}
\hfill
\begin{tabular}{rl}
\multicolumn{2}{l}{\textbf{FUTURE I}}\\
& werden + \textbf{reinigen}\\
\end{tabular}
\hfill
\begin{tabular}{rl}
\multicolumn{2}{l}{\textbf{PLUSQUAMPERFEKT}}\\
& hatten + \textbf{gereinigt}\\
\end{tabular}

\vspace{10pt}

\begin{tabular}{lll}
\multicolumn{3}{l}{\textbf{MIT PRÄPOSITIONEN}}\\
& \textbf{reinigen von} + Dat & (to clean from)\\
& \textbf{reinigen mit} + Dat & (to clean with)\\
\end{tabular}
\hfill
\begin{tabular}{lll}
\multicolumn{3}{l}{\textbf{TRENNBARE VERBEN}}\\
& \textbf{ab$|$reinigen} & (to clean off)\\
& \textbf{aus$|$reinigen} & (to clean out)\\
\end{tabular}
\vspace{-5pt}
\end{verbbox}

%reisen (A1) – to travel
\begin{verbbox}\addcontentsline{toc}{subsubsection}{reisen (to travel)}
    {\Large \bf reisen} (to travel) \hfill Type: \textbf{regular} \hfill Auxiliary Verb: \textbf{sein} 
    
     \vspace{5pt}

    \hrule
    
    \vspace{5pt}

  \begin{tabular}{rl}
     \multicolumn{2}{l}{\textbf{PRÄSENS} }\\
    ich & \textbf{reise}\\
    du & \textbf{reist}\\
    er/sie/es & \textbf{reist}\\
    wir & \textbf{reisen}\\
    ihr & \textbf{reist}\\
    sie/Sie & \textbf{reisen}\\
  \end{tabular}
  \hfill
     \begin{tabular}{rl}
    \multicolumn{2}{l}{\textbf{PRÄTERITUM} }\\
    ich & \textbf{reiste}\\
    du & \textbf{reistest}\\
    er/sie/es & \textbf{reiste}\\
    wir & \textbf{reisten}\\
    ihr & \textbf{reistet}\\
    sie/Sie & \textbf{reisten}\\
  \end{tabular}
  \hfill
  \begin{tabular}{rl}
     \multicolumn{2}{l}{\textbf{IMPERATIV} }\\
   & \textbf{reis(e) (du)!}\\
   & \textbf{reisen wir!}\\
   & \textbf{reist (ihr)!}\\
   & \textbf{reisen Sie!}\\
   & \vspace{10pt}
  \end{tabular}

    \vspace{10pt}

 \begin{tabular}{rl}
     \multicolumn{2}{l}{\textbf{PERFEKT}}\\
   & sein + \textbf{gereist}\\
  \end{tabular}
\hfill
   \begin{tabular}{rl}
   
    \multicolumn{2}{l}{\textbf{FUTURE I} }\\
  &  werden + \textbf{reisen}\\
 
  \end{tabular}
\hfill
   \begin{tabular}{rl}
      \multicolumn{2}{l}{\textbf{PLUSQUAMPERFEKT}}\\
  &  waren + \textbf{gereist}\\
  \end{tabular}
  
  \vspace{10pt}



\begin{tabular}{lll}
\multicolumn{3}{l}{\textbf{MIT PRÄPOSITIONEN}}\\
& ---\\
\end{tabular}
\hfill
\begin{tabular}{lll}
\multicolumn{3}{l}{\textbf{TRENNBARE VERBEN}}\\
& ---\\
\end{tabular}

  \vspace{-5pt}

\end{verbbox}

%reklamieren (to complain / claim)
\begin{verbbox}\addcontentsline{toc}{subsubsection}{reklamieren (to complain / claim)}
{\Large \bf reklamieren} (to complain / claim) \hfill Type: \textbf{regular} \hfill Auxiliary Verb: \textbf{haben}

\vspace{5pt}
\hrule
\vspace{5pt}

\begin{tabular}{rl}
\multicolumn{2}{l}{\textbf{PRÄSENS}}\\
ich & \textbf{reklamiere}\\
du & \textbf{reklamierst}\\
er/sie/es & \textbf{reklamiert}\\
wir & \textbf{reklamieren}\\
ihr & \textbf{reklamiert}\\
sie/Sie & \textbf{reklamieren}\\
\end{tabular}
\hfill
\begin{tabular}{rl}
\multicolumn{2}{l}{\textbf{PRÄTERITUM}}\\
ich & \textbf{reklamierte}\\
du & \textbf{reklamiertest}\\
er/sie/es & \textbf{reklamierte}\\
wir & \textbf{reklamierten}\\
ihr & \textbf{reklamiertet}\\
sie/Sie & \textbf{reklamierten}\\
\end{tabular}
\hfill
\begin{tabular}{rl}
\multicolumn{2}{l}{\textbf{IMPERATIV}}\\
& \textbf{reklamiere (du)!}\\
& \textbf{reklamieren wir!}\\
& \textbf{reklamiert (ihr)!}\\
& \textbf{reklamieren Sie!}\\
\end{tabular}

\vspace{10pt}

\begin{tabular}{rl}
\multicolumn{2}{l}{\textbf{PERFEKT}}\\
& haben + \textbf{reklamiert}\\
\end{tabular}
\hfill
\begin{tabular}{rl}
\multicolumn{2}{l}{\textbf{FUTURE I}}\\
& werden + \textbf{reklamieren}\\
\end{tabular}
\hfill
\begin{tabular}{rl}
\multicolumn{2}{l}{\textbf{PLUSQUAMPERFEKT}}\\
& hatten + \textbf{reklamiert}\\
\end{tabular}

\vspace{10pt}

\begin{tabular}{lll}
\multicolumn{3}{l}{\textbf{MIT PRÄPOSITIONEN}}\\
& \textbf{reklamieren bei} + Dat & (to complain to)\\
& \textbf{reklamieren für} + Akk & (to claim for)\\
\end{tabular}
\hfill
\begin{tabular}{lll}
\multicolumn{3}{l}{\textbf{TRENNBARE VERBEN}}\\
& ---\\
\end{tabular}
\vspace{-5pt}
\end{verbbox}

%rufen (A2)  – to call / shout
\begin{verbbox-acc}\addcontentsline{toc}{subsubsection}{rufen (to call / shout)}
    {\Large \bf rufen} (to call / shout) \hfill Type: \textbf{irregular, + Akkusativ} \hfill Auxiliary Verb: \textbf{haben} 
    
     \vspace{5pt}

    \hrule
    
    \vspace{5pt}

  \begin{tabular}{rl}
     \multicolumn{2}{l}{\textbf{PRÄSENS} }\\
    ich & \textbf{rufe}\\
    du & \textbf{rufst}\\
    er/sie/es & \textbf{ruft}\\
    wir & \textbf{rufen}\\
    ihr & \textbf{ruft}\\
    sie/Sie & \textbf{rufen}\\
  \end{tabular}
  \hfill
     \begin{tabular}{rl}
    \multicolumn{2}{l}{\textbf{PRÄTERITUM} }\\
    ich & \textbf{rief}\\
    du & \textbf{riefst}\\
    er/sie/es & \textbf{rief}\\
    wir & \textbf{riefen}\\
    ihr & \textbf{rieft}\\
    sie/Sie & \textbf{riefen}\\
  \end{tabular}
  \hfill
  \begin{tabular}{rl}
     \multicolumn{2}{l}{\textbf{IMPERATIV} }\\
   & \textbf{ruf(e) (du)!}\\
   & \textbf{rufen wir!}\\
   & \textbf{ruft (ihr)!}\\
   & \textbf{rufen Sie!}\\
   & \vspace{10pt}
  \end{tabular}

    \vspace{10pt}

 \begin{tabular}{rl}
     \multicolumn{2}{l}{\textbf{PERFEKT}}\\
   & haben + \textbf{gerufen}\\
  \end{tabular}
\hfill
   \begin{tabular}{rl}
   
    \multicolumn{2}{l}{\textbf{FUTURE I} }\\
   & werden + \textbf{rufen}\\
 
  \end{tabular}
\hfill
   \begin{tabular}{rl}
      \multicolumn{2}{l}{\textbf{PLUSQUAMPERFEKT}}\\
   & hatten + \textbf{gerufen}\\
  \end{tabular}
  

 \vspace{10pt}

   \begin{tabular}{lll}
      \multicolumn{3}{l}{\textbf{MIT PRÄPOSITIONEN}}\\
& \textbf{rufen nach} + Dat & (to cry out for)\\
& \textbf{rufen zu} + Dat & (to call to)\\
&\textbf{rufen um} + Akk & (to plead for)\\
&\textbf{anrufen bei} + Dat & (to call at (office))\\
&\textbf{anrufen wegen} + Gen/Dat & (to call about)\\
&\textbf{zurufen zu} + Dat & (to shout to)\\
&\textbf{aufrufen zu} + Dat & (to call on to do)\\
&\textbf{abrufen von} + Dat & (to retrieve from)\\
  \end{tabular}
  \hfill
   \begin{tabular}{lll}
      \multicolumn{3}{l}{\textbf{TRENNBARE VERBEN}}\\
 & \textbf{an$|$rufen} & (to call (phone))\\
 & \textbf{ab$|$rufen} & (to retrieve / access)\\
 & \textbf{zu$|$rufen} & (to call out)\\
  & \textbf{heraus$|$rufen} & (to call out (from inside))\\
  & \textbf{hinein$|$rufen} & (to call into)\\
  & \textbf{zurück$|$rufen} & (to call back)\\
  & \textbf{auf$|$rufen} & (to call up / apppeal)\\
     & \textbf{herbei$|$rufen} & (to summon)\\  
  \end{tabular}


  \vspace{-5pt}

\end{verbbox-acc}


%ruhen (to rest)
\begin{verbbox}\addcontentsline{toc}{subsubsection}{ruhen (to rest)}
{\Large \bf ruhen} (to rest) \hfill Type: \textbf{regular} \hfill Auxiliary Verb: \textbf{haben} 

\vspace{5pt}
\hrule
\vspace{5pt}

\begin{tabular}{rl}
\multicolumn{2}{l}{\textbf{PRÄSENS}}\\
ich & \textbf{ruhe}\\
du & \textbf{ruhst}\\
er/sie/es & \textbf{ruht}\\
wir & \textbf{ruhen}\\
ihr & \textbf{ruht}\\
sie/Sie & \textbf{ruhen}\\
\end{tabular}
\hfill
\begin{tabular}{rl}
\multicolumn{2}{l}{\textbf{PRÄTERITUM}}\\
ich & \textbf{ruhte}\\
du & \textbf{ruhtest}\\
er/sie/es & \textbf{ruhte}\\
wir & \textbf{ruhten}\\
ihr & \textbf{ruhtet}\\
sie/Sie & \textbf{ruhten}\\
\end{tabular}
\hfill
\begin{tabular}{rl}
\multicolumn{2}{l}{\textbf{IMPERATIV}}\\
& \textbf{ruhe (du)!}\\
& \textbf{ruhen wir!}\\
& \textbf{ruht (ihr)!}\\
& \textbf{ruhen Sie!}\\
\end{tabular}

\vspace{10pt}

\begin{tabular}{rl}
\multicolumn{2}{l}{\textbf{PERFEKT}}\\
& haben + \textbf{geruht}\\
\end{tabular}
\hfill
\begin{tabular}{rl}
\multicolumn{2}{l}{\textbf{FUTURE I}}\\
& werden + \textbf{ruhen}\\
\end{tabular}
\hfill
\begin{tabular}{rl}
\multicolumn{2}{l}{\textbf{PLUSQUAMPERFEKT}}\\
& hatten + \textbf{geruht}\\
\end{tabular}

\vspace{10pt}

\begin{tabular}{lll}
\multicolumn{3}{l}{\textbf{MIT PRÄPOSITIONEN}}\\
& \textbf{ruhen auf} + Dat & (to rest on)\\
\end{tabular}
\hfill
\begin{tabular}{lll}
\multicolumn{3}{l}{\textbf{TRENNBARE VERBEN}}\\
& ---\\
\end{tabular}

\vspace{-5pt}
\end{verbbox}

\subsection{S Verbs}

%sagen (A1) – to say
\begin{verbbox-acc}\addcontentsline{toc}{subsubsection}{sagen (to say)}
    {\Large \bf sagen} (to say) \hfill Type: \textbf{regular, + Akkusativ} \hfill Auxiliary Verb: \textbf{haben} 
    
     \vspace{5pt}

    \hrule
    
    \vspace{5pt}

  \begin{tabular}{rl}
     \multicolumn{2}{l}{\textbf{PRÄSENS} }\\
    ich & \textbf{sage}\\
    du & \textbf{sagst}\\
    er/sie/es & \textbf{sagt}\\
    wir & \textbf{sagen}\\
    ihr & \textbf{sagt}\\
    sie/Sie & \textbf{sagen}\\
  \end{tabular}
  \hfill
     \begin{tabular}{rl}
    \multicolumn{2}{l}{\textbf{PRÄTERITUM} }\\
    ich & \textbf{sagte}\\
    du & \textbf{sagtest}\\
    er/sie/es & \textbf{sagte}\\
    wir & \textbf{sagten}\\
    ihr & \textbf{sagtet}\\
    sie/Sie & \textbf{sagten}\\
  \end{tabular}
  \hfill
  \begin{tabular}{rl}
     \multicolumn{2}{l}{\textbf{IMPERATIV} }\\
   & \textbf{sag(e) (du)!}\\
   & \textbf{sagen wir!}\\
   & \textbf{sagt (ihr)!}\\
   & \textbf{sagen Sie!}\\
   & \vspace{10pt}
  \end{tabular}

    \vspace{10pt}

 \begin{tabular}{rl}
     \multicolumn{2}{l}{\textbf{PERFEKT}}\\
   & haben + \textbf{gesagt}\\
  \end{tabular}
\hfill
   \begin{tabular}{rl}
   
    \multicolumn{2}{l}{\textbf{FUTURE I} }\\
  &  werden + \textbf{sagen}\\
 
  \end{tabular}
\hfill
   \begin{tabular}{rl}
      \multicolumn{2}{l}{\textbf{PLUSQUAMPERFEKT}}\\
  &  hatten + \textbf{gesagt}\\
  \end{tabular}
  
  \vspace{-5pt}


\end{verbbox-acc}

%sammeln (to collect / gather)
\begin{verbbox-acc}\addcontentsline{toc}{subsubsection}{sammeln (to collect / gather)}
    {\Large \bf sammeln} (to collect / gather) \hfill Type: \textbf{regular} \hfill Auxiliary Verb: \textbf{haben} 
    
     \vspace{5pt}

    \hrule
    
    \vspace{5pt}

  \begin{tabular}{rl}
     \multicolumn{2}{l}{\textbf{PRÄSENS} }\\
    ich & \textbf{sammle}\\
    du & \textbf{sammelst}\\
    er/sie/es & \textbf{sammelt}\\
    wir & \textbf{sammeln}\\
    ihr & \textbf{sammelt}\\
    sie/Sie & \textbf{sammeln}\\
  \end{tabular}
  \hfill
  \begin{tabular}{rl}
    \multicolumn{2}{l}{\textbf{PRÄTERITUM} }\\
    ich & \textbf{sammelte}\\
    du & \textbf{sammeltest}\\
    er/sie/es & \textbf{sammelte}\\
    wir & \textbf{sammelten}\\
    ihr & \textbf{sammeltet}\\
    sie/Sie & \textbf{sammelten}\\
  \end{tabular}
  \hfill
  \begin{tabular}{rl}
     \multicolumn{2}{l}{\textbf{IMPERATIV} }\\
   & \textbf{sammle (du)!}\\
   & \textbf{sammeln wir!}\\
   & \textbf{sammelt (ihr)!}\\
   & \textbf{sammeln Sie!}\\
   & \vspace{10pt}
  \end{tabular}

 \vspace{10pt}

 \begin{tabular}{rl}
     \multicolumn{2}{l}{\textbf{PERFEKT}}\\
   & haben + \textbf{gesammelt}\\
  \end{tabular}
\hfill
\begin{tabular}{rl}
    \multicolumn{2}{l}{\textbf{FUTURE I}}\\
   & werden + \textbf{sammeln}\\
  \end{tabular}
\hfill
\begin{tabular}{rl}
    \multicolumn{2}{l}{\textbf{PLUSQUAMPERFEKT}}\\
   & hatten + \textbf{gesammelt}\\
  \end{tabular}

 \vspace{10pt}

\begin{tabular}{lll}
      \multicolumn{3}{l}{\textbf{MIT PRÄPOSITIONEN}}\\
& \textbf{sammeln für} + Akk & (to collect for / gather for)\\
& \textbf{sammeln von} + Dat & (to collect from / gather from)\\
& \textbf{sammeln in} + Dat & (to collect in / within)\\
\end{tabular}
  \hfill
\begin{tabular}{lll}
      \multicolumn{3}{l}{\textbf{TRENNBARE VERBEN}}\\
 & \textbf{zusammen$|$sammeln} & (to gather together / collect together)\\
 & \textbf{ein$|$sammeln} & (to collect in / pick up)\\
\vspace{10pt}
  \end{tabular}

  \vspace{-5pt}


\end{verbbox-acc}

%saugen (to vacuum / suck)
\begin{verbbox-acc}\addcontentsline{toc}{subsubsection}{saugen (to vacuum / suck)}
{\Large \bf saugen} (to vacuum / suck) \hfill Type: \textbf{regular, + Akkusativ} \hfill Auxiliary Verb: \textbf{haben}

\vspace{5pt}
\hrule
\vspace{5pt}

\begin{tabular}{rl}
\multicolumn{2}{l}{\textbf{PRÄSENS}}\\
ich & \textbf{sauge}\\
du & \textbf{saugst}\\
er/sie/es & \textbf{saugt}\\
wir & \textbf{saugen}\\
ihr & \textbf{saugt}\\
sie/Sie & \textbf{saugen}\\
\end{tabular}
\hfill
\begin{tabular}{rl}
\multicolumn{2}{l}{\textbf{PRÄTERITUM}}\\
ich & \textbf{saugte}\\
du & \textbf{saugtest}\\
er/sie/es & \textbf{saugte}\\
wir & \textbf{saugten}\\
ihr & \textbf{saugtet}\\
sie/Sie & \textbf{saugten}\\
\end{tabular}
\hfill
\begin{tabular}{rl}
\multicolumn{2}{l}{\textbf{IMPERATIV}}\\
& \textbf{saug(e) (du)!}\\
& \textbf{saugen wir!}\\
& \textbf{saugt (ihr)!}\\
& \textbf{saugen Sie!}\\
\end{tabular}

\vspace{10pt}

\begin{tabular}{rl}
\multicolumn{2}{l}{\textbf{PERFEKT}}\\
& haben + \textbf{gesaugt}\\
\end{tabular}
\hfill
\begin{tabular}{rl}
\multicolumn{2}{l}{\textbf{FUTURE I}}\\
& werden + \textbf{saugen}\\
\end{tabular}
\hfill
\begin{tabular}{rl}
\multicolumn{2}{l}{\textbf{PLUSQUAMPERFEKT}}\\
& hatten + \textbf{gesaugt}\\
\end{tabular}

\vspace{10pt}

\begin{tabular}{lll}
\multicolumn{3}{l}{\textbf{MIT PRÄPOSITIONEN}}\\
\end{tabular}
\hfill
\begin{tabular}{lll}
\multicolumn{3}{l}{\textbf{TRENNBARE VERBEN}}\\
& \textbf{auf$|$saugen} & (to absorb)\\
& \textbf{ein$|$saugen} & (to inhale / absorb)\\
\end{tabular}

\vspace{-5pt}
\end{verbbox-acc}

%schaden (to harm / damage)
\begin{verbbox}\addcontentsline{toc}{subsubsection}{schaden (to harm / damage)}
{\Large \bf schaden} (to harm / damage) \hfill Type: \textbf{regular} \hfill Auxiliary Verb: \textbf{haben}

\vspace{5pt}
\hrule
\vspace{5pt}

\begin{tabular}{rl}
\multicolumn{2}{l}{\textbf{PRÄSENS}}\\
ich & \textbf{schade}\\
du & \textbf{schadest}\\
er/sie/es & \textbf{schadet}\\
wir & \textbf{schaden}\\
ihr & \textbf{schadet}\\
sie/Sie & \textbf{schaden}\\
\end{tabular}
\hfill
\begin{tabular}{rl}
\multicolumn{2}{l}{\textbf{PRÄTERITUM}}\\
ich & \textbf{schadete}\\
du & \textbf{schadetest}\\
er/sie/es & \textbf{schadete}\\
wir & \textbf{schadeten}\\
ihr & \textbf{schadetet}\\
sie/Sie & \textbf{schadeten}\\
\end{tabular}
\hfill
\begin{tabular}{rl}
\multicolumn{2}{l}{\textbf{IMPERATIV}}\\
& \textbf{schade (du)!}\\
& \textbf{schaden wir!}\\
& \textbf{schadet (ihr)!}\\
& \textbf{schaden Sie!}\\
\end{tabular}

\vspace{10pt}

\begin{tabular}{rl}
\multicolumn{2}{l}{\textbf{PERFEKT}}\\
& haben + \textbf{geschadet}\\
\end{tabular}
\hfill
\begin{tabular}{rl}
\multicolumn{2}{l}{\textbf{FUTURE I}}\\
& werden + \textbf{schaden}\\
\end{tabular}
\hfill
\begin{tabular}{rl}
\multicolumn{2}{l}{\textbf{PLUSQUAMPERFEKT}}\\
& hatten + \textbf{geschadet}\\
\end{tabular}

\vspace{10pt}

\begin{tabular}{lll}
\multicolumn{3}{l}{\textbf{MIT PRÄPOSITIONEN}}\\
& \textbf{schaden durch} + Akk & (to be harmed by)\\
\end{tabular}
\hfill
\begin{tabular}{lll}
\multicolumn{3}{l}{\textbf{TRENNBARE VERBEN}}\\
& ---\\
\end{tabular}
\vspace{-5pt}
\end{verbbox}

%schaffen (B1) – to manage, create
\begin{verbbox-acc}\addcontentsline{toc}{subsubsection}{schaffen (to manage, create)}
    {\Large \bf schaffen} (to manage, create) \hfill Type: \textbf{irregular, + Akkusativ} \hfill Auxiliary Verb: \textbf{haben} 
    
     \vspace{5pt}

    \hrule
    
    \vspace{5pt}

  \begin{tabular}{rl}
     \multicolumn{2}{l}{\textbf{PRÄSENS} }\\
    ich & \textbf{schaffe}\\
    du & \textbf{schaffst}\\
    er/sie/es & \textbf{schafft}\\
    wir & \textbf{schaffen}\\
    ihr & \textbf{schafft}\\
    sie/Sie & \textbf{schaffen}\\
  \end{tabular}
  \hfill
     \begin{tabular}{rl}
    \multicolumn{2}{l}{\textbf{PRÄTERITUM} }\\
    ich & \textbf{schuf}\\
    du & \textbf{schufst}\\
    er/sie/es & \textbf{schuf}\\
    wir & \textbf{schufen}\\
    ihr & \textbf{schuft}\\
    sie/Sie & \textbf{schufen}\\
  \end{tabular}
  \hfill
  \begin{tabular}{rl}
     \multicolumn{2}{l}{\textbf{IMPERATIV} }\\
   & \textbf{schaff(e) (du)!}\\
   & \textbf{schaffen wir!}\\
   & \textbf{schafft (ihr)!}\\
   & \textbf{schaffen Sie!}\\
   & \vspace{10pt}
  \end{tabular}

    \vspace{10pt}

 \begin{tabular}{rl}
     \multicolumn{2}{l}{\textbf{PERFEKT}}\\
   & haben + \textbf{geschaffen}\\
  \end{tabular}
\hfill
   \begin{tabular}{rl}
   
    \multicolumn{2}{l}{\textbf{FUTURE I} }\\
  &  werden + \textbf{schaffen}\\
 
  \end{tabular}
\hfill
   \begin{tabular}{rl}
      \multicolumn{2}{l}{\textbf{PLUSQUAMPERFEKT}}\\
  &  hatten + \textbf{geschaffen}\\
  \end{tabular}
  
  \vspace{-5pt}


\end{verbbox-acc}

%schalten (to switch)
\begin{verbbox}\addcontentsline{toc}{subsubsection}{schalten (to switch)}
{\Large \bf schalten} (to switch) \hfill Type: \textbf{regular} \hfill Auxiliary Verb: \textbf{haben} 

\vspace{5pt}
\hrule
\vspace{5pt}

\begin{tabular}{rl}
\multicolumn{2}{l}{\textbf{PRÄSENS}}\\
ich & \textbf{schalte}\\
du & \textbf{schaltest}\\
er/sie/es & \textbf{schaltet}\\
wir & \textbf{schalten}\\
ihr & \textbf{schaltet}\\
sie/Sie & \textbf{schalten}\\
\end{tabular}
\hfill
\begin{tabular}{rl}
\multicolumn{2}{l}{\textbf{PRÄTERITUM}}\\
ich & \textbf{schaltete}\\
du & \textbf{schaltetest}\\
er/sie/es & \textbf{schaltete}\\
wir & \textbf{schalteten}\\
ihr & \textbf{schaltetet}\\
sie/Sie & \textbf{schalteten}\\
\end{tabular}
\hfill
\begin{tabular}{rl}
\multicolumn{2}{l}{\textbf{IMPERATIV}}\\
& \textbf{schalte (du)!}\\
& \textbf{schalten wir!}\\
& \textbf{schaltet (ihr)!}\\
& \textbf{schalten Sie!}\\
\end{tabular}

\vspace{10pt}

\begin{tabular}{rl}
\multicolumn{2}{l}{\textbf{PERFEKT}}\\
& haben + \textbf{geschaltet}\\
\end{tabular}
\hfill
\begin{tabular}{rl}
\multicolumn{2}{l}{\textbf{FUTURE I}}\\
& werden + \textbf{schalten}\\
\end{tabular}
\hfill
\begin{tabular}{rl}
\multicolumn{2}{l}{\textbf{PLUSQUAMPERFEKT}}\\
& hatten + \textbf{geschaltet}\\
\end{tabular}

\vspace{10pt}

\begin{tabular}{lll}
\multicolumn{3}{l}{\textbf{MIT PRÄPOSITIONEN}}\\
& \textbf{schalten auf} + Akk & (to switch to)\\
\vspace{10pt}
\end{tabular}
\hfill
\begin{tabular}{lll}
\multicolumn{3}{l}{\textbf{TRENNBARE VERBEN}}\\
& \textbf{ein$|$schalten} & (to switch on)\\
& \textbf{aus$|$schalten} & (to switch off)\\
\end{tabular}

\vspace{-5pt}
\end{verbbox}

%schauen (to look / watch)
\begin{verbbox}\addcontentsline{toc}{subsubsection}{schauen (to look / watch)}
{\Large \bf schauen} (to look / watch) \hfill Type: \textbf{regular} \hfill Auxiliary Verb: \textbf{haben} 

\vspace{5pt}
\hrule
\vspace{5pt}

\begin{tabular}{rl}
\multicolumn{2}{l}{\textbf{PRÄSENS}}\\
ich & \textbf{schaue}\\
du & \textbf{schaust}\\
er/sie/es & \textbf{schaut}\\
wir & \textbf{schauen}\\
ihr & \textbf{schaut}\\
sie/Sie & \textbf{schauen}\\
\end{tabular}
\hfill
\begin{tabular}{rl}
\multicolumn{2}{l}{\textbf{PRÄTERITUM}}\\
ich & \textbf{schaute}\\
du & \textbf{schautest}\\
er/sie/es & \textbf{schaute}\\
wir & \textbf{schauten}\\
ihr & \textbf{schautet}\\
sie/Sie & \textbf{schauten}\\
\end{tabular}
\hfill
\begin{tabular}{rl}
\multicolumn{2}{l}{\textbf{IMPERATIV}}\\
& \textbf{schaue (du)!}\\
& \textbf{schauen wir!}\\
& \textbf{schaut (ihr)!}\\
& \textbf{schauen Sie!}\\
\end{tabular}

\vspace{10pt}

\begin{tabular}{rl}
\multicolumn{2}{l}{\textbf{PERFEKT}}\\
& haben + \textbf{geschaut}\\
\end{tabular}
\hfill
\begin{tabular}{rl}
\multicolumn{2}{l}{\textbf{FUTURE I}}\\
& werden + \textbf{schauen}\\
\end{tabular}
\hfill
\begin{tabular}{rl}
\multicolumn{2}{l}{\textbf{PLUSQUAMPERFEKT}}\\
& hatten + \textbf{geschaut}\\
\end{tabular}

\vspace{10pt}

\begin{tabular}{lll}
\multicolumn{3}{l}{\textbf{MIT PRÄPOSITIONEN}}\\
& \textbf{schauen auf} + Akk & (to look at)\\
\end{tabular}
\hfill
\begin{tabular}{lll}
\multicolumn{3}{l}{\textbf{TRENNBARE VERBEN}}\\
& \textbf{an$|$schauen} & (to look at)\\
\end{tabular}

\vspace{-5pt}
\end{verbbox}


%scheinen - to shine
\begin{verbbox}\addcontentsline{toc}{subsubsection}{scheinen (to shine)}
    {\Large \bf scheinen} (to shine) \hfill Type: \textbf{irregular} \hfill Auxiliary Verb: \textbf{haben} 
    
     \vspace{5pt}

    \hrule
    
    \vspace{5pt}

  \begin{tabular}{rl}
     \multicolumn{2}{l}{\textbf{PRÄSENS} }\\
    ich & \textbf{scheine}\\
    du & \textbf{scheinst}\\
    er/sie/es & \textbf{scheint}\\
    wir & \textbf{scheinen}\\
    ihr & \textbf{scheint}\\
    sie/Sie & \textbf{scheinen}\\
  \end{tabular}
  \hfill
     \begin{tabular}{rl}
    \multicolumn{2}{l}{\textbf{PRÄTERITUM} }\\
    ich & \textbf{schien}\\
    du & \textbf{schienst}\\
    er/sie/es & \textbf{schien}\\
    wir & \textbf{schienen}\\
    ihr & \textbf{schient}\\
    sie/Sie & \textbf{schienen}\\
  \end{tabular}
  \hfill
  \begin{tabular}{rl}
     \multicolumn{2}{l}{\textbf{IMPERATIV} }\\
   & \textbf{schein(e) (du)!}\\
   & \textbf{scheinen wir!}\\
   & \textbf{scheint (ihr)!}\\
   & \textbf{scheinen Sie!}\\
   & \vspace{10pt}
  \end{tabular}

    \vspace{10pt}

 \begin{tabular}{rl}
     \multicolumn{2}{l}{\textbf{PERFEKT}}\\
   & haben + \textbf{geschienen}\\
  \end{tabular}
\hfill
   \begin{tabular}{rl}
   
    \multicolumn{2}{l}{\textbf{FUTURE I} }\\
  &  werden + \textbf{scheinen}\\
 
  \end{tabular}
\hfill
   \begin{tabular}{rl}
      \multicolumn{2}{l}{\textbf{PLUSQUAMPERFEKT}}\\
  &  hatten + \textbf{geschienen}\\
  \end{tabular}
  
  \vspace{10pt}



\begin{tabular}{lll}
\multicolumn{3}{l}{\textbf{MIT PRÄPOSITIONEN}}\\
& ---\\
\end{tabular}
\hfill
\begin{tabular}{lll}
\multicolumn{3}{l}{\textbf{TRENNBARE VERBEN}}\\
& ---\\
\end{tabular}

\vspace{-5pt}

\end{verbbox}

%schicken (to send)
\begin{verbbox}\addcontentsline{toc}{subsubsection}{schicken (to send)}
{\Large \bf schicken} (to send) \hfill Type: \textbf{regular} \hfill Auxiliary Verb: \textbf{haben}

\vspace{5pt}
\hrule
\vspace{5pt}

\begin{tabular}{rl}
\multicolumn{2}{l}{\textbf{PRÄSENS}}\\
ich & \textbf{schicke}\\
du & \textbf{schickst}\\
er/sie/es & \textbf{schickt}\\
wir & \textbf{schicken}\\
ihr & \textbf{schickt}\\
sie/Sie & \textbf{schicken}\\
\end{tabular}
\hfill
\begin{tabular}{rl}
\multicolumn{2}{l}{\textbf{PRÄTERITUM}}\\
ich & \textbf{schickte}\\
du & \textbf{schicktest}\\
er/sie/es & \textbf{schickte}\\
wir & \textbf{schickten}\\
ihr & \textbf{schicktet}\\
sie/Sie & \textbf{schickten}\\
\end{tabular}
\hfill
\begin{tabular}{rl}
\multicolumn{2}{l}{\textbf{IMPERATIV}}\\
& \textbf{schick(e) (du)!}\\
& \textbf{schicken wir!}\\
& \textbf{schickt (ihr)!}\\
& \textbf{schicken Sie!}\\
\end{tabular}

\vspace{10pt}

\begin{tabular}{rl}
\multicolumn{2}{l}{\textbf{PERFEKT}}\\
& haben + \textbf{geschickt}\\
\end{tabular}
\hfill
\begin{tabular}{rl}
\multicolumn{2}{l}{\textbf{FUTURE I}}\\
& werden + \textbf{schicken}\\
\end{tabular}
\hfill
\begin{tabular}{rl}
\multicolumn{2}{l}{\textbf{PLUSQUAMPERFEKT}}\\
& hatten + \textbf{geschickt}\\
\end{tabular}

\vspace{10pt}

\begin{tabular}{lll}
\multicolumn{3}{l}{\textbf{MIT PRÄPOSITIONEN}}\\
& \textbf{schicken an} + Akk & (to send to)\\
& \textbf{schicken mit} + Dat & (to send with)\\
\end{tabular}
\hfill
\begin{tabular}{lll}
\multicolumn{3}{l}{\textbf{TRENNBARE VERBEN}}\\
& \textbf{ab$|$schicken} & (to send off)\\
& \textbf{ein$|$schicken} & (to submit / send in)\\
& \textbf{zurück$|$schicken} & (to send back)\\
\end{tabular}
\vspace{-5pt}
\end{verbbox}

%schlafen (A1) – to sleep
\begin{verbbox}\addcontentsline{toc}{subsubsection}{schlafen (to sleep)}
    {\Large \bf schlafen} (to sleep) \hfill Type: \textbf{irregular} \hfill Auxiliary Verb: \textbf{haben} 
    
     \vspace{5pt}

    \hrule
    
    \vspace{5pt}

  \begin{tabular}{rl}
     \multicolumn{2}{l}{\textbf{PRÄSENS} }\\
    ich & \textbf{schlafe}\\
    du & \textbf{schläft}\\
    er/sie/es & \textbf{schläft}\\
    wir & \textbf{schlafen}\\
    ihr & \textbf{schlaft}\\
    sie/Sie & \textbf{schlafen}\\
  \end{tabular}
  \hfill
     \begin{tabular}{rl}
    \multicolumn{2}{l}{\textbf{PRÄTERITUM} }\\
    ich & \textbf{schlief}\\
    du & \textbf{schliefest}\\
    er/sie/es & \textbf{schlief}\\
    wir & \textbf{schliefen}\\
    ihr & \textbf{schlieft}\\
    sie/Sie & \textbf{schlieften}\\
  \end{tabular}
  \hfill
  \begin{tabular}{rl}
     \multicolumn{2}{l}{\textbf{IMPERATIV} }\\
   & \textbf{schlaf(e) (du)!}\\
   & \textbf{schlafen wir!}\\
   & \textbf{schlaft (ihr)!}\\
   & \textbf{schlafen Sie!}\\
   & \vspace{10pt}
  \end{tabular}

    \vspace{10pt}

 \begin{tabular}{rl}
     \multicolumn{2}{l}{\textbf{PERFEKT}}\\
  &  haben + \textbf{geschlafen}\\
  \end{tabular}
\hfill
   \begin{tabular}{rl}
   
    \multicolumn{2}{l}{\textbf{FUTURE I} }\\
  &  werden + \textbf{schlafen}\\
 
  \end{tabular}
\hfill
   \begin{tabular}{rl}
      \multicolumn{2}{l}{\textbf{PLUSQUAMPERFEKT}}\\
  &  hatten + \textbf{geschlafen}\\
  \end{tabular}
  
  \vspace{10pt}



\begin{tabular}{lll}
\multicolumn{3}{l}{\textbf{MIT PRÄPOSITIONEN}}\\
& ---\\
\end{tabular}
\hfill
\begin{tabular}{lll}
\multicolumn{3}{l}{\textbf{TRENNBARE VERBEN}}\\
& ---\\
\end{tabular}

  \vspace{-5pt}

\end{verbbox}

%schlagen (to hit / beat / strike)
\begin{verbbox-acc} \addcontentsline{toc}{subsubsection}{schlagen (to hit / beat / strike)}
    {\Large \bf schlagen} (to hit / beat / strike) \hfill Type: \textbf{irregular, + Akkusativ} \hfill Auxiliary Verb: \textbf{haben} 
    
     \vspace{5pt}

    \hrule
    
    \vspace{5pt}

  \begin{tabular}{rl}
     \multicolumn{2}{l}{\textbf{PRÄSENS} }\\
    ich & \textbf{schlage}\\
    du & \textbf{schlägst}\\
    er/sie/es & \textbf{schlägt}\\
    wir & \textbf{schlagen}\\
    ihr & \textbf{schlagt}\\
    sie/Sie & \textbf{schlagen}\\
  \end{tabular}
  \hfill
  \begin{tabular}{rl}
    \multicolumn{2}{l}{\textbf{PRÄTERITUM} }\\
    ich & \textbf{schlug}\\
    du & \textbf{schlugst}\\
    er/sie/es & \textbf{schlug}\\
    wir & \textbf{schlugen}\\
    ihr & \textbf{schlugt}\\
    sie/Sie & \textbf{schlugen}\\
  \end{tabular}
  \hfill
  \begin{tabular}{rl}
     \multicolumn{2}{l}{\textbf{IMPERATIV} }\\
   & \textbf{schlag (du)!}\\
   & \textbf{schlagen wir!}\\
   & \textbf{schlagt (ihr)!}\\
   & \textbf{schlagen Sie!}\\
   & \vspace{10pt}
  \end{tabular}

 \vspace{10pt}

 \begin{tabular}{rl}
     \multicolumn{2}{l}{\textbf{PERFEKT}}\\
   & haben + \textbf{geschlagen}\\
  \end{tabular}
\hfill
\begin{tabular}{rl}
    \multicolumn{2}{l}{\textbf{FUTURE I}}\\
   & werden + \textbf{schlagen}\\
  \end{tabular}
\hfill
\begin{tabular}{rl}
    \multicolumn{2}{l}{\textbf{PLUSQUAMPERFEKT}}\\
   & hatten + \textbf{geschlagen}\\
  \end{tabular}

 \vspace{10pt}

\begin{tabular}{lll}
      \multicolumn{3}{l}{\textbf{MIT PRÄPOSITIONEN}}\\
& \textbf{schlagen auf} + Akk & (to hit on / strike)\\
& \textbf{schlagen gegen} + Akk & (to hit against / collide)\\
& \textbf{schlagen für} + Akk & (to vote / advocate for)\\
\vspace{20pt}
\end{tabular}
  \hfill
\begin{tabular}{lll}
      \multicolumn{3}{l}{\textbf{TRENNBARE VERBEN}}\\
 & \textbf{auf$|$schlagen} & (to hit open / pitch / unfold)\\
 & \textbf{ein$|$schlagen} & (to break / strike / hammer in)\\
 & \textbf{vor$|$schlagen} & (to propose / suggest)\\
 & \textbf{zurück$|$schlagen} & (to hit back / retaliate)\\
 & \textbf{zusammen$|$schlagen} & (to beat / knock together)\\
  \end{tabular}

  \vspace{-5pt}

\end{verbbox-acc}

%schleifen
\begin{verbbox}\addcontentsline{toc}{subsubsection}{schleifen (to sand / grind / sharpen)}
    {\Large \bf schleifen} (to sand / grind / sharpen) \hfill Type: \textbf{regular} \hfill Auxiliary Verb: \textbf{haben} 
    
     \vspace{5pt}

    \hrule
    
    \vspace{5pt}

  \begin{tabular}{rl}
     \multicolumn{2}{l}{\textbf{PRÄSENS} }\\
    ich & \textbf{schleife}\\
    du & \textbf{schleifst}\\
    er/sie/es & \textbf{schleift}\\
    wir & \textbf{schleifen}\\
    ihr & \textbf{schleift}\\
    sie/Sie & \textbf{schleifen}\\
  \end{tabular}
  \hfill
  \begin{tabular}{rl}
    \multicolumn{2}{l}{\textbf{PRÄTERITUM} }\\
    ich & \textbf{schliff}\\
    du & \textbf{schliffst}\\
    er/sie/es & \textbf{schliff}\\
    wir & \textbf{schliffen}\\
    ihr & \textbf{schlifft}\\
    sie/Sie & \textbf{schliffen}\\
  \end{tabular}
  \hfill
  \begin{tabular}{rl}
     \multicolumn{2}{l}{\textbf{IMPERATIV} }\\
   & \textbf{schleife (du)!}\\
   & \textbf{schleifen wir!}\\
   & \textbf{schleift (ihr)!}\\
   & \textbf{schleifen Sie!}\\
   & \vspace{10pt}
  \end{tabular}

 \vspace{10pt}

 \begin{tabular}{rl}
     \multicolumn{2}{l}{\textbf{PERFEKT}}\\
   & haben + \textbf{geschliffen}\\
  \end{tabular}
\hfill
\begin{tabular}{rl}
    \multicolumn{2}{l}{\textbf{FUTURE I}}\\
   & werden + \textbf{schleifen}\\
  \end{tabular}
\hfill
\begin{tabular}{rl}
    \multicolumn{2}{l}{\textbf{PLUSQUAMPERFEKT}}\\
   & hatten + \textbf{geschliffen}\\
  \end{tabular}

 \vspace{10pt}

\begin{tabular}{lll}
      \multicolumn{3}{l}{\textbf{MIT PRÄPOSITIONEN}}\\
& \textbf{schleifen an} + Dat & (to sharpen on / work on)\\
& \textbf{schleifen mit} + Dat & (to sand / grind with)\\
\end{tabular}
  \hfill
\begin{tabular}{lll}
      \multicolumn{3}{l}{\textbf{TRENNBARE VERBEN}}\\
 & \textbf{ab$|$schleifen} & (to grind off / remove)\\
 & \textbf{ein$|$schleifen} & (to polish / fit in)\\
  \end{tabular}


\vspace{-5pt}
\end{verbbox}

%schlendern (to stroll)
\begin{verbbox}\addcontentsline{toc}{subsubsection}{schlendern (to stroll)}
{\Large \bf schlendern} (to stroll) \hfill Type: \textbf{regular} \hfill Auxiliary Verb: \textbf{sein}

\vspace{5pt}
\hrule
\vspace{5pt}

\begin{tabular}{rl}
\multicolumn{2}{l}{\textbf{PRÄSENS}}\\
ich & \textbf{schlendere}\\
du & \textbf{schlenderst}\\
er/sie/es & \textbf{schlendert}\\
wir & \textbf{schlendern}\\
ihr & \textbf{schlendert}\\
sie/Sie & \textbf{schlendern}\\
\end{tabular}
\hfill
\begin{tabular}{rl}
\multicolumn{2}{l}{\textbf{PRÄTERITUM}}\\
ich & \textbf{schlenderte}\\
du & \textbf{schlendertest}\\
er/sie/es & \textbf{schlenderte}\\
wir & \textbf{schlenderten}\\
ihr & \textbf{schlendertet}\\
sie/Sie & \textbf{schlenderten}\\
\end{tabular}
\hfill
\begin{tabular}{rl}
\multicolumn{2}{l}{\textbf{IMPERATIV}}\\
& \textbf{schlendere (du)!}\\
& \textbf{schlendern wir!}\\
& \textbf{schlendert (ihr)!}\\
& \textbf{schlendern Sie!}\\
\end{tabular}

\vspace{10pt}

\begin{tabular}{rl}
\multicolumn{2}{l}{\textbf{PERFEKT}}\\
& sein + \textbf{geschlendert}\\
\end{tabular}
\hfill
\begin{tabular}{rl}
\multicolumn{2}{l}{\textbf{FUTURE I}}\\
& werden + \textbf{schlendern}\\
\end{tabular}
\hfill
\begin{tabular}{rl}
\multicolumn{2}{l}{\textbf{PLUSQUAMPERFEKT}}\\
& waren + \textbf{geschlendert}\\
\end{tabular}

\vspace{10pt}

\begin{tabular}{lll}
\multicolumn{3}{l}{\textbf{MIT PRÄPOSITIONEN}}\\
& \textbf{schlendern durch} + Akk & (to stroll through)\\
& \textbf{schlendern über} + Akk & (to stroll across)\\
\end{tabular}
\hfill
\begin{tabular}{lll}
\multicolumn{3}{l}{\textbf{TRENNBARE VERBEN}}\\
& ---\\
\vspace{10pt}
\end{tabular}

\vspace{-5pt}
\end{verbbox}

%schließen (A1) – to close
\begin{verbbox-acc}\addcontentsline{toc}{subsubsection}{schließen (to close)}
    {\Large \bf schließen} (to close) \hfill Type: \textbf{irregular, + Akkusativ} \hfill Auxiliary Verb: \textbf{haben} 
    
     \vspace{5pt}

    \hrule
    
    \vspace{5pt}

  \begin{tabular}{rl}
     \multicolumn{2}{l}{\textbf{PRÄSENS} }\\
    ich & \textbf{schließe}\\
    du & \textbf{schließt}\\
    er/sie/es & \textbf{schließt}\\
    wir & \textbf{schließen}\\
    ihr & \textbf{schließt}\\
    sie/Sie & \textbf{schließen}\\
  \end{tabular}
  \hfill
     \begin{tabular}{rl}
    \multicolumn{2}{l}{\textbf{PRÄTERITUM} }\\
    ich & \textbf{schloss}\\
    du & \textbf{schlossest}\\
    er/sie/es & \textbf{schloss}\\
    wir & \textbf{schlossen}\\
    ihr & \textbf{schlosst}\\
    sie/Sie & \textbf{schlossen}\\
  \end{tabular}
  \hfill
  \begin{tabular}{rl}
     \multicolumn{2}{l}{\textbf{IMPERATIV} }\\
   & \textbf{schließ(e) (du)!}\\
   & \textbf{schließen wir!}\\
   & \textbf{schließt (ihr)!}\\
   & \textbf{schließen Sie!}\\
   & \vspace{10pt}
  \end{tabular}

    \vspace{10pt}

 \begin{tabular}{rl}
     \multicolumn{2}{l}{\textbf{PERFEKT}}\\
   & haben + \textbf{geschlossen}\\
  \end{tabular}
\hfill
   \begin{tabular}{rl}
   
    \multicolumn{2}{l}{\textbf{FUTURE I} }\\
   & werden + \textbf{schließen}\\
 
  \end{tabular}
\hfill
   \begin{tabular}{rl}
      \multicolumn{2}{l}{\textbf{PLUSQUAMPERFEKT}}\\
   & hatten + \textbf{geschlossen}\\
  \end{tabular}
  
  \vspace{10pt}

   \begin{tabular}{lll}
      \multicolumn{3}{l}{\textbf{MIT PRÄPOSITIONEN}}\\
      & \textbf{schließen auf} + Akk & (to infer / conclude)\\
&\textbf{schließen aus} + Dat & (to infer from)\\
&\textbf{schließen mit} + Dat & (to close with / end with)\\
& \textbf{anschließen an} + Akk & (to conect to)\\
& \textbf{abschließen mit} + Dat & (to finish with)\\
&\textbf{einschließen in} + Akk & (to include in)\\
&\textbf{ausschließen von} + Dat & (to exlude from)
  \end{tabular}
  \hfill
   \begin{tabular}{lll}
      \multicolumn{3}{l}{\textbf{TRENNBARE VERBEN}}\\
 & \textbf{ab$|$schließen} & (to lock)\\
 & \textbf{auf$|$schließen} & (to unlock)\\
 & \textbf{ein$|$schließen} & (to include / enclose)\\
  & \textbf{aus$|$schließen} & (to exclude)\\
  & \textbf{an$|$schließen} & (to connect / join)\\
  & \textbf{zu$|$schließen} & (to lock (shut))
  \end{tabular}




  \vspace{-5pt}

\end{verbbox-acc}

%schmecken (A1) – to taste
\begin{verbbox-dat}\addcontentsline{toc}{subsubsection}{schmecken (to taste)}
    {\Large \bf schmecken} (to taste) \hfill Type: \textbf{regular, + Akkusativ/Dativ} \hfill Auxiliary Verb: \textbf{haben} 
    
     \vspace{5pt}

    \hrule
    
    \vspace{5pt}

  \begin{tabular}{rl}
     \multicolumn{2}{l}{\textbf{PRÄSENS} }\\
    ich & \textbf{schmecke}\\
    du & \textbf{schmeckst}\\
    er/sie/es & \textbf{schmeckt}\\
    wir & \textbf{schmecken}\\
    ihr & \textbf{schmeckt}\\
    sie/Sie & \textbf{schmecken}\\
  \end{tabular}
  \hfill
     \begin{tabular}{rl}
    \multicolumn{2}{l}{\textbf{PRÄTERITUM} }\\
    ich & \textbf{schmeckte}\\
    du & \textbf{schmecktest}\\
    er/sie/es & \textbf{schmeckte}\\
    wir & \textbf{schmeckten}\\
    ihr & \textbf{schmecktet}\\
    sie/Sie & \textbf{schmeckten}\\
  \end{tabular}
  \hfill
  \begin{tabular}{rl}
     \multicolumn{2}{l}{\textbf{IMPERATIV} }\\
   & \textbf{schmeck(e) (du)!}\\
   & \textbf{schmecken wir!}\\
   & \textbf{schmeckt (ihr)!}\\
   & \textbf{schmecken Sie!}\\
   & \vspace{10pt}
  \end{tabular}

    \vspace{10pt}

 \begin{tabular}{rl}
     \multicolumn{2}{l}{\textbf{PERFEKT}}\\
  &  haben + \textbf{geschmeckt}\\
  \end{tabular}
\hfill
   \begin{tabular}{rl}
   
    \multicolumn{2}{l}{\textbf{FUTURE I} }\\
  &  werden + \textbf{schmecken}\\
 
  \end{tabular}
\hfill
   \begin{tabular}{rl}
      \multicolumn{2}{l}{\textbf{PLUSQUAMPERFEKT}}\\
  &  hatten + \textbf{geschmeckt}\\
  \end{tabular}
  
  \vspace{-5pt}


\end{verbbox-dat}

%schmücken (to decorate)
\begin{verbbox}\addcontentsline{toc}{subsubsection}{schmücken (to decorate)}
{\Large \bf schmücken} (to decorate) \hfill Type: \textbf{regular} \hfill Auxiliary Verb: \textbf{haben}

\vspace{5pt}
\hrule
\vspace{5pt}

\begin{tabular}{rl}
\multicolumn{2}{l}{\textbf{PRÄSENS}}\\
ich & \textbf{schmücke}\\
du & \textbf{schmückst}\\
er/sie/es & \textbf{schmückt}\\
wir & \textbf{schmücken}\\
ihr & \textbf{schmückt}\\
sie/Sie & \textbf{schmücken}\\
\end{tabular}
\hfill
\begin{tabular}{rl}
\multicolumn{2}{l}{\textbf{PRÄTERITUM}}\\
ich & \textbf{schmückte}\\
du & \textbf{schmücktest}\\
er/sie/es & \textbf{schmückte}\\
wir & \textbf{schmückten}\\
ihr & \textbf{schmücktet}\\
sie/Sie & \textbf{schmückten}\\
\end{tabular}
\hfill
\begin{tabular}{rl}
\multicolumn{2}{l}{\textbf{IMPERATIV}}\\
& \textbf{schmücke (du)!}\\
& \textbf{schmücken wir!}\\
& \textbf{schmückt (ihr)!}\\
& \textbf{schmücken Sie!}\\
\end{tabular}

\vspace{10pt}

\begin{tabular}{rl}
\multicolumn{2}{l}{\textbf{PERFEKT}}\\
& haben + \textbf{geschmückt}\\
\end{tabular}
\hfill
\begin{tabular}{rl}
\multicolumn{2}{l}{\textbf{FUTURE I}}\\
& werden + \textbf{schmücken}\\
\end{tabular}
\hfill
\begin{tabular}{rl}
\multicolumn{2}{l}{\textbf{PLUSQUAMPERFEKT}}\\
& hatten + \textbf{geschmückt}\\
\end{tabular}

\vspace{10pt}

\begin{tabular}{lll}
\multicolumn{3}{l}{\textbf{MIT PRÄPOSITIONEN}}\\
& \textbf{schmücken mit} + Dat & (to decorate with)\\
& \textbf{schmücken für} + Akk & (to decorate for)\\
\end{tabular}
\hfill
\begin{tabular}{lll}
\multicolumn{3}{l}{\textbf{TRENNBARE VERBEN}}\\
& \textbf{aus$|$schmücken} & (to decorate elaborately)\\
\end{tabular}
\vspace{-5pt}
\end{verbbox}

%schneien (A1) – to snow
\begin{verbbox}\addcontentsline{toc}{subsubsection}{schneien (to snow)}
    {\Large \bf schneien} (to snow) \hfill Type: \textbf{irregular} \hfill Auxiliary Verb: \textbf{haben} 
    
     \vspace{5pt}

    \hrule
    
    \vspace{5pt}

  \begin{tabular}{rl}
     \multicolumn{2}{l}{\textbf{PRÄSENS} }\\
    ich & \textbf{-}\\
    du & \textbf{-}\\
    es & \textbf{schneit}\\
    wir & \textbf{-}\\
    ihr & \textbf{-}\\
    sie/Sie & \textbf{-}\\
  \end{tabular}
  \hfill
     \begin{tabular}{rl}
    \multicolumn{2}{l}{\textbf{PRÄTERITUM} }\\
    ich & \textbf{-}\\
    du & \textbf{-}\\
    es & \textbf{schneite}\\
    wir & \textbf{-}\\
    ihr & \textbf{-}\\
    sie/Sie & \textbf{-}\\
  \end{tabular}
  \hfill
  \begin{tabular}{rl}
     \multicolumn{2}{l}{\textbf{IMPERATIV} }\\
   & \textbf{-}\\
   & \textbf{-}\\
   & \textbf{-}\\
   & \textbf{-}\\
   & \vspace{10pt}
  \end{tabular}

    \vspace{10pt}

 \begin{tabular}{rl}
     \multicolumn{2}{l}{\textbf{PERFEKT}}\\
  &  haben + \textbf{geschneit}\\
  \end{tabular}
\hfill
   \begin{tabular}{rl}
   
    \multicolumn{2}{l}{\textbf{FUTURE I} }\\
  &  werden + \textbf{schneien}\\
 
  \end{tabular}
\hfill
   \begin{tabular}{rl}
      \multicolumn{2}{l}{\textbf{PLUSQUAMPERFEKT}}\\
  &  hatten + \textbf{geschneit}\\
  \end{tabular}
  
  \vspace{10pt}



\begin{tabular}{lll}
\multicolumn{3}{l}{\textbf{MIT PRÄPOSITIONEN}}\\
& ---\\
\end{tabular}
\hfill
\begin{tabular}{lll}
\multicolumn{3}{l}{\textbf{TRENNBARE VERBEN}}\\
& ---\\
\end{tabular}

  \vspace{-5pt}

\end{verbbox}


%schreiben (A1) – to write
\begin{verbbox-acc}\addcontentsline{toc}{subsubsection}{schreiben (to write)}
    {\Large \bf schreiben} (to write) \hfill Type: \textbf{irregular, + Akkusativ} \hfill Auxiliary Verb: \textbf{haben} 
    
     \vspace{5pt}

    \hrule
    
    \vspace{5pt}

  \begin{tabular}{rl}
     \multicolumn{2}{l}{\textbf{PRÄSENS} }\\
    ich & \textbf{schreibe}\\
    du & \textbf{schreibst}\\
    er/sie/es & \textbf{schreibt}\\
    wir & \textbf{schreiben}\\
    ihr & \textbf{schreibt}\\
    sie/Sie & \textbf{schreiben}\\
  \end{tabular}
  \hfill
     \begin{tabular}{rl}
    \multicolumn{2}{l}{\textbf{PRÄTERITUM} }\\
    ich & \textbf{schrieb}\\
    du & \textbf{schriebst}\\
    er/sie/es & \textbf{schrieb}\\
    wir & \textbf{schrieben}\\
    ihr & \textbf{schriebt}\\
    sie/Sie & \textbf{schrieben}\\
  \end{tabular}
  \hfill
  \begin{tabular}{rl}
     \multicolumn{2}{l}{\textbf{IMPERATIV} }\\
   & \textbf{schreib(e) (du)!}\\
   & \textbf{schreiben wir!}\\
   & \textbf{schreibt (ihr)!}\\
   & \textbf{schreiben Sie!}\\
   & \vspace{10pt}
  \end{tabular}

    \vspace{10pt}

 \begin{tabular}{rl}
     \multicolumn{2}{l}{\textbf{PERFEKT}}\\
  &  haben + \textbf{geschrieben}\\
  \end{tabular}
\hfill
   \begin{tabular}{rl}
   
    \multicolumn{2}{l}{\textbf{FUTURE I} }\\
  &  werden + \textbf{schreiben}\\
 
  \end{tabular}
\hfill
   \begin{tabular}{rl}
      \multicolumn{2}{l}{\textbf{PLUSQUAMPERFEKT}}\\
  &  hatten + \textbf{geschrieben}\\
  \end{tabular}
  
  \vspace{-5pt}


\end{verbbox-acc}


%schwitzen (to sweat)
\begin{verbbox}\addcontentsline{toc}{subsubsection}{schwitzen (to sweat)}
{\Large \bf schwitzen} (to sweat) \hfill Type: \textbf{regular} \hfill Auxiliary Verb: \textbf{haben}

\vspace{5pt}
\hrule
\vspace{5pt}

\begin{tabular}{rl}
\multicolumn{2}{l}{\textbf{PRÄSENS}}\\
ich & \textbf{schwitze}\\
du & \textbf{schwitzt}\\
er/sie/es & \textbf{schwitzt}\\
wir & \textbf{schwitzen}\\
ihr & \textbf{schwitzt}\\
sie/Sie & \textbf{schwitzen}\\
\end{tabular}
\hfill
\begin{tabular}{rl}
\multicolumn{2}{l}{\textbf{PRÄTERITUM}}\\
ich & \textbf{schwitzte}\\
du & \textbf{schwitztest}\\
er/sie/es & \textbf{schwitzte}\\
wir & \textbf{schwitzten}\\
ihr & \textbf{schwitztet}\\
sie/Sie & \textbf{schwitzten}\\
\end{tabular}
\hfill
\begin{tabular}{rl}
\multicolumn{2}{l}{\textbf{IMPERATIV}}\\
& \textbf{schwitze (du)!}\\
& \textbf{schwitzen wir!}\\
& \textbf{schwitzt (ihr)!}\\
& \textbf{schwitzen Sie!}\\
\end{tabular}

\vspace{10pt}

\begin{tabular}{rl}
\multicolumn{2}{l}{\textbf{PERFEKT}}\\
& haben + \textbf{geschwitzt}\\
\end{tabular}
\hfill
\begin{tabular}{rl}
\multicolumn{2}{l}{\textbf{FUTURE I}}\\
& werden + \textbf{schwitzen}\\
\end{tabular}
\hfill
\begin{tabular}{rl}
\multicolumn{2}{l}{\textbf{PLUSQUAMPERFEKT}}\\
& hatten + \textbf{geschwitzt}\\
\end{tabular}

\vspace{10pt}

\begin{tabular}{lll}
\multicolumn{3}{l}{\textbf{MIT PRÄPOSITIONEN}}\\
& \textbf{schwitzen bei} + Dat & (to sweat during)\\
& \textbf{schwitzen vor} + Dat & (to sweat from)\\
\end{tabular}
\hfill
\begin{tabular}{lll}
\multicolumn{3}{l}{\textbf{TRENNBARE VERBEN}}\\
& ---\\
\end{tabular}

\vspace{-5pt}

\end{verbbox}

%sehen (A1) – to see
\begin{verbbox-acc}\addcontentsline{toc}{subsubsection}{sehen (to see)}
    {\Large \bf sehen} (to see) \hfill Type: \textbf{irregular, + Akkusativ} \hfill Auxiliary Verb: \textbf{haben} 
    
     \vspace{5pt}

    \hrule
    
    \vspace{5pt}

  \begin{tabular}{rl}
     \multicolumn{2}{l}{\textbf{PRÄSENS} }\\
    ich & \textbf{sehe}\\
    du & \textbf{siehst}\\
    er/sie/es & \textbf{sieht}\\
    wir & \textbf{sehen}\\
    ihr & \textbf{seht}\\
    sie/Sie & \textbf{sehen}\\
  \end{tabular}
  \hfill
     \begin{tabular}{rl}
    \multicolumn{2}{l}{\textbf{PRÄTERITUM} }\\
    ich & \textbf{sah}\\
    du & \textbf{sahst}\\
    er/sie/es & \textbf{sah}\\
    wir & \textbf{sahen}\\
    ihr & \textbf{saht}\\
    sie/Sie & \textbf{sahen}\\
  \end{tabular}
  \hfill
  \begin{tabular}{rl}
     \multicolumn{2}{l}{\textbf{IMPERATIV} }\\
   & \textbf{sieh(e) (du)!}\\
   & \textbf{sehen wir!}\\
   & \textbf{seht (ihr)!}\\
   & \textbf{sehen Sie!}\\
   & \vspace{10pt}
  \end{tabular}

    \vspace{10pt}

 \begin{tabular}{rl}
     \multicolumn{2}{l}{\textbf{PERFEKT}}\\
  &  haben + \textbf{gesehen}\\
  \end{tabular}
\hfill
   \begin{tabular}{rl}
   
    \multicolumn{2}{l}{\textbf{FUTURE I} }\\
  &  werden + \textbf{sehen}\\
 
  \end{tabular}
\hfill
   \begin{tabular}{rl}
      \multicolumn{2}{l}{\textbf{PLUSQUAMPERFEKT}}\\
  &  hatten + \textbf{gesehen}\\
  \end{tabular}
  


 \vspace{10pt}

   \begin{tabular}{lll}
      \multicolumn{3}{l}{\textbf{MIT PRÄPOSITIONEN}}\\
& \textbf{sehen auf} + Akk & (to pay attention to)\\
& \textbf{sehen in} + Akk & (to look into)\\
& \textbf{sehen nach} + Dat & (to check / look after)\\
& \textbf{absehen von} + Dat & (to refrain from)\\
& \textbf{ansehen als} + Akk & (to regard as)\\
& \textbf{ansehen von} + Dat & (to look from)\\
& \textbf{hinsehen auf} + Akk & (to look at)\\
& \textbf{wegsehen bei} + Dat & (to ignore during)\\
& \textbf{zusehen bei} + Dat & (to watch while)\\
  \end{tabular}
  \hfill
   \begin{tabular}{lll}
      \multicolumn{3}{l}{\textbf{TRENNBARE VERBEN}}\\
 & \textbf{an$|$sehen} & (to regard as)\\
  & \textbf{ab$|$sehen} & (to foresee / refrain)\\
   & \textbf{auf$|$sehen} & (to look up)\\
     & \textbf{aus$|$sehen} & (to appear)\\
    & \textbf{durch$|$sehen} & (to look through / examine)\\
      & \textbf{ein$|$sehen} & (to review)\\
    & \textbf{hin$|$sehen} & (to look (in a direction))\\
      & \textbf{um$|$sehen} & (to look around)\\
     & \textbf{weg$|$sehen} & (to look away / ignore)\\
      & \textbf{zu$|$sehen} & (to watch)\\
  \end{tabular}


  \vspace{-5pt}

\end{verbbox-acc}

%setzen - 
\begin{verbbox}\addcontentsline{toc}{subsubsection}{setzen (to set / place / put)}
    {\Large \bf setzen} (to set / place / put) \hfill Type: \textbf{regular} \hfill Auxiliary Verb: \textbf{haben} 
    
     \vspace{5pt}

    \hrule
    
    \vspace{5pt}

  \begin{tabular}{rl}
     \multicolumn{2}{l}{\textbf{PRÄSENS} }\\
    ich & \textbf{setze}\\
    du & \textbf{setzt}\\
    er/sie/es & \textbf{setzt}\\
    wir & \textbf{setzen}\\
    ihr & \textbf{setzt}\\
    sie/Sie & \textbf{setzen}\\
  \end{tabular}
  \hfill
  \begin{tabular}{rl}
    \multicolumn{2}{l}{\textbf{PRÄTERITUM} }\\
    ich & \textbf{setzte}\\
    du & \textbf{setztest}\\
    er/sie/es & \textbf{setzte}\\
    wir & \textbf{setzten}\\
    ihr & \textbf{setztet}\\
    sie/Sie & \textbf{setzten}\\
  \end{tabular}
  \hfill
  \begin{tabular}{rl}
     \multicolumn{2}{l}{\textbf{IMPERATIV} }\\
   & \textbf{setze (du)!}\\
   & \textbf{setzen wir!}\\
   & \textbf{setzt (ihr)!}\\
   & \textbf{setzen Sie!}\\
   & \vspace{10pt}
  \end{tabular}

 \vspace{10pt}

 \begin{tabular}{rl}
     \multicolumn{2}{l}{\textbf{PERFEKT}}\\
   & haben + \textbf{gesetzt}\\
  \end{tabular}
\hfill
\begin{tabular}{rl}
    \multicolumn{2}{l}{\textbf{FUTURE I}}\\
   & werden + \textbf{setzen}\\
  \end{tabular}
\hfill
\begin{tabular}{rl}
    \multicolumn{2}{l}{\textbf{PLUSQUAMPERFEKT}}\\
   & hatten + \textbf{gesetzt}\\
  \end{tabular}

 \vspace{10pt}

\begin{tabular}{lll}
      \multicolumn{3}{l}{\textbf{MIT PRÄPOSITIONEN}}\\
& \textbf{setzen auf} + Akk & (to place on / set on)\\
& \textbf{setzen in} + Akk & (to place into / put in)\\
& \textbf{sich setzen auf} + Akk & (to sit down on)\\
& \textbf{setzen für} + Akk & (to set for / dedicate)\\
\end{tabular}
  \hfill
\begin{tabular}{lll}
      \multicolumn{3}{l}{\textbf{TRENNBARE VERBEN}}\\
 & \textbf{ein$|$setzen} & (to insert / deploy / put in)\\
 & \textbf{auf$|$setzen} & (to place on / put onto)\\
 & \textbf{durch$|$setzen} & (to enforce / push through)\\
\vspace{10pt}
  \end{tabular}


\vspace{-5pt}
\end{verbbox}

%siezen (to address someone with "Sie")
\begin{verbbox}\addcontentsline{toc}{subsubsection}{siezen (to address someone with ``Sie")}
{\Large \bf siezen} (to address someone with ``Sie") \hfill Type: \textbf{regular} \hfill Auxiliary Verb: \textbf{haben}

\vspace{5pt}
\hrule
\vspace{5pt}

\begin{tabular}{rl}
\multicolumn{2}{l}{\textbf{PRÄSENS}}\\
ich & \textbf{sieze}\\
du & \textbf{siezt}\\
er/sie/es & \textbf{siezt}\\
wir & \textbf{siezen}\\
ihr & \textbf{siezt}\\
sie/Sie & \textbf{siezen}\\
\end{tabular}
\hfill
\begin{tabular}{rl}
\multicolumn{2}{l}{\textbf{PRÄTERITUM}}\\
ich & \textbf{siezte}\\
du & \textbf{sieztest}\\
er/sie/es & \textbf{siezte}\\
wir & \textbf{siezten}\\
ihr & \textbf{sieztet}\\
sie/Sie & \textbf{siezten}\\
\end{tabular}
\hfill
\begin{tabular}{rl}
\multicolumn{2}{l}{\textbf{IMPERATIV}}\\
& \textbf{sieze (du)!}\\
& \textbf{siezen wir!}\\
& \textbf{siezt (ihr)!}\\
& \textbf{siezen Sie!}\\
\end{tabular}

\vspace{10pt}

\begin{tabular}{rl}
\multicolumn{2}{l}{\textbf{PERFEKT}}\\
& haben + \textbf{gesiezt}\\
\end{tabular}
\hfill
\begin{tabular}{rl}
\multicolumn{2}{l}{\textbf{FUTURE I}}\\
& werden + \textbf{siezen}\\
\end{tabular}
\hfill
\begin{tabular}{rl}
\multicolumn{2}{l}{\textbf{PLUSQUAMPERFEKT}}\\
& hatten + \textbf{gesiezt}\\
\end{tabular}

\vspace{10pt}

\begin{tabular}{lll}
\multicolumn{3}{l}{\textbf{MIT PRÄPOSITIONEN}}\\
& \textbf{siezen mit} + Dat & (to use "Sie" with someone)\\
\end{tabular}
\hfill
\begin{tabular}{lll}
\multicolumn{3}{l}{\textbf{TRENNBARE VERBEN}}\\
& ---\\
\end{tabular}
\vspace{-5pt}
\end{verbbox}

%sitzen (A1) – to sit
\begin{verbbox}\addcontentsline{toc}{subsubsection}{sitzen (to sit)}
    {\Large \bf sitzen} (to sit) \hfill Type: \textbf{irregular} \hfill Auxiliary Verb: \textbf{haben / sein context-dependent} 
    
     \vspace{5pt}

    \hrule
    
    \vspace{5pt}

  \begin{tabular}{rl}
     \multicolumn{2}{l}{\textbf{PRÄSENS} }\\
    ich & \textbf{sitze}\\
    du & \textbf{sitzt}\\
    er/sie/es & \textbf{sitzt}\\
    wir & \textbf{sitzen}\\
    ihr & \textbf{sitzt}\\
    sie/Sie & \textbf{sitzen}\\
  \end{tabular}
  \hfill
  \begin{tabular}{rl}
    \multicolumn{2}{l}{\textbf{PRÄTERITUM} }\\
    ich & \textbf{saß}\\
    du & \textbf{saßt}\\
    er/sie/es & \textbf{saß}\\
    wir & \textbf{saßen}\\
    ihr & \textbf{saßt}\\
    sie/Sie & \textbf{saßen}\\
  \end{tabular}
  \hfill
  \begin{tabular}{rl}
     \multicolumn{2}{l}{\textbf{IMPERATIV} }\\
   & \textbf{sitze (du)!}\\
   & \textbf{sitzen wir!}\\
   & \textbf{sitzt (ihr)!}\\
   & \textbf{sitzen Sie!}\\
   & \vspace{10pt}
  \end{tabular}

 \vspace{10pt}

 \begin{tabular}{rl}
     \multicolumn{2}{l}{\textbf{PERFEKT}}\\
   & haben + \textbf{gesessen}\\
  \end{tabular}
\hfill
\begin{tabular}{rl}
    \multicolumn{2}{l}{\textbf{FUTURE I}}\\
   & werden + \textbf{sitzen}\\
  \end{tabular}
\hfill
\begin{tabular}{rl}
    \multicolumn{2}{l}{\textbf{PLUSQUAMPERFEKT}}\\
   & hatten + \textbf{gesessen}\\
  \end{tabular}

 \vspace{10pt}

\begin{tabular}{lll}
      \multicolumn{3}{l}{\textbf{MIT PRÄPOSITIONEN}}\\
& \textbf{sitzen auf} + Dat & (to sit on)\\
& \textbf{sitzen in} + Dat & (to sit in / inside)\\
& \textbf{sitzen an} + Dat & (to sit at / by)\\
\end{tabular}
  \hfill
\begin{tabular}{lll}
      \multicolumn{3}{l}{\textbf{TRENNBARE VERBEN}}\\
 & \textbf{abs$|$sitzen} & (to be seated / sit down)\\
 & \textbf{zusammen$|$sitzen} & (to sit together / be grouped)\\
 & \textbf{nach$|$sitzen} & (to stay after school / serve detention)\\
  \end{tabular}
  \vspace{-5pt}

\end{verbbox}

%sparen - to save
\begin{verbbox-acc}\addcontentsline{toc}{subsubsection}{sparen (to save)}
    {\Large \bf sparen} (to save) \hfill Type: \textbf{regular, + Akkusativ} \hfill Auxiliary Verb: \textbf{haben} 
    
     \vspace{5pt}

    \hrule
    
    \vspace{5pt}

  \begin{tabular}{rl}
     \multicolumn{2}{l}{\textbf{PRÄSENS} }\\
    ich & \textbf{spare}\\
    du & \textbf{sparst}\\
    er/sie/es & \textbf{spart}\\
    wir & \textbf{sparen}\\
    ihr & \textbf{spart}\\
    sie/Sie & \textbf{sparen}\\
  \end{tabular}
  \hfill
     \begin{tabular}{rl}
    \multicolumn{2}{l}{\textbf{PRÄTERITUM} }\\
    ich & \textbf{sparte}\\
    du & \textbf{spartest}\\
    er/sie/es & \textbf{sparte}\\
    wir & \textbf{sparten}\\
    ihr & \textbf{spartet}\\
    sie/Sie & \textbf{sparten}\\
  \end{tabular}
  \hfill
  \begin{tabular}{rl}
     \multicolumn{2}{l}{\textbf{IMPERATIV} }\\
   & \textbf{spar(e) (du)!}\\
   & \textbf{sparen wir!}\\
   & \textbf{sparet (ihr)!}\\
   & \textbf{sparen Sie!}\\
   & \vspace{10pt}
  \end{tabular}

    \vspace{10pt}

 \begin{tabular}{rl}
     \multicolumn{2}{l}{\textbf{PERFEKT}}\\
  &  haben + \textbf{gespart}\\
  \end{tabular}
\hfill
   \begin{tabular}{rl}
   
    \multicolumn{2}{l}{\textbf{FUTURE I} }\\
  &  werden + \textbf{sparen}\\
 
  \end{tabular}
\hfill
   \begin{tabular}{rl}
      \multicolumn{2}{l}{\textbf{PLUSQUAMPERFEKT}}\\
  &  hatten + \textbf{gesparet}\\
  \end{tabular}
  
  \vspace{-5pt}


\end{verbbox-acc}


%spielen (A1) – to play
\begin{verbbox-acc}\addcontentsline{toc}{subsubsection}{spielen (to play)}
    {\Large \bf spielen} (to play) \hfill Type: \textbf{regular, + Akkusativ} \hfill Auxiliary Verb: \textbf{haben} 
    
     \vspace{5pt}

    \hrule
    
    \vspace{5pt}

  \begin{tabular}{rl}
     \multicolumn{2}{l}{\textbf{PRÄSENS} }\\
    ich & \textbf{spiele}\\
    du & \textbf{spielst}\\
    er/sie/es & \textbf{spielt}\\
    wir & \textbf{spielen}\\
    ihr & \textbf{spielt}\\
    sie/Sie & \textbf{spielen}\\
  \end{tabular}
  \hfill
     \begin{tabular}{rl}
    \multicolumn{2}{l}{\textbf{PRÄTERITUM} }\\
    ich & \textbf{spielte}\\
    du & \textbf{spieltest}\\
    er/sie/es & \textbf{spielte}\\
    wir & \textbf{spielten}\\
    ihr & \textbf{spieltet}\\
    sie/Sie & \textbf{spielten}\\
  \end{tabular}
  \hfill
  \begin{tabular}{rl}
     \multicolumn{2}{l}{\textbf{IMPERATIV} }\\
   & \textbf{spiel(e) (du)!}\\
   & \textbf{spielen wir!}\\
   & \textbf{spielt (ihr)!}\\
   & \textbf{spielen Sie!}\\
   & \vspace{10pt}
  \end{tabular}

    \vspace{10pt}

 \begin{tabular}{rl}
     \multicolumn{2}{l}{\textbf{PERFEKT}}\\
   & haben + \textbf{gespielt}\\
  \end{tabular}
\hfill
   \begin{tabular}{rl}
   
    \multicolumn{2}{l}{\textbf{FUTURE I} }\\
  &  werden + \textbf{spielen}\\
 
  \end{tabular}
\hfill
   \begin{tabular}{rl}
      \multicolumn{2}{l}{\textbf{PLUSQUAMPERFEKT}}\\
   & hatten + \textbf{gespielt}\\
  \end{tabular}
  
  \vspace{-5pt}


\end{verbbox-acc}

%sprechen (A1) – to speak
\begin{verbbox}\addcontentsline{toc}{subsubsection}{sprechen (to speak)}
    {\Large \bf sprechen} (to speak) \hfill Type: \textbf{irregular} \hfill Auxiliary Verb: \textbf{haben} 
    
     \vspace{5pt}

    \hrule
    
    \vspace{5pt}

  \begin{tabular}{rl}
     \multicolumn{2}{l}{\textbf{PRÄSENS} }\\
    ich & \textbf{spreche}\\
    du & \textbf{sprichst}\\
    er/sie/es & \textbf{spricht}\\
    wir & \textbf{sprechen}\\
    ihr & \textbf{sprecht}\\
    sie/Sie & \textbf{sprechen}\\
  \end{tabular}
  \hfill
     \begin{tabular}{rl}
    \multicolumn{2}{l}{\textbf{PRÄTERITUM} }\\
    ich & \textbf{sprach}\\
    du & \textbf{sprachst}\\
    er/sie/es & \textbf{sprach}\\
    wir & \textbf{sprachen}\\
    ihr & \textbf{spracht}\\
    sie/Sie & \textbf{sprachen}\\
  \end{tabular}
  \hfill
  \begin{tabular}{rl}
     \multicolumn{2}{l}{\textbf{IMPERATIV} }\\
   & \textbf{sprich (du)!}\\
   & \textbf{sprechen wir!}\\
   & \textbf{sprecht (ihr)!}\\
   & \textbf{sprechen Sie!}\\
   & \vspace{10pt}
  \end{tabular}

    \vspace{10pt}

 \begin{tabular}{rl}
     \multicolumn{2}{l}{\textbf{PERFEKT}}\\
  &  haben + \textbf{gesprochen}\\
  \end{tabular}
\hfill
   \begin{tabular}{rl}
   
    \multicolumn{2}{l}{\textbf{FUTURE I} }\\
   & werden + \textbf{sprechen}\\
 
  \end{tabular}
\hfill
   \begin{tabular}{rl}
      \multicolumn{2}{l}{\textbf{PLUSQUAMPERFEKT}}\\
   & hatten + \textbf{gesprochen}\\
  \end{tabular}
  
   \vspace{10pt}

   \begin{tabular}{lll}
      \multicolumn{3}{l}{\textbf{MIT PRÄPOSITIONEN}}\\
&\textbf{sprechen für} + Akk & (to speak for)\\
& \textbf{sprechen über} + Akk & (to speak about)\\
& \textbf{sprechen mit} + Dat & (to speak with)\\
&\textbf{sprechen von} + Dat & (to speak of)
  \end{tabular}
  \hfill
   \begin{tabular}{lll}
      \multicolumn{3}{l}{\textbf{TRENNBARE VERBEN}}\\
 & \textbf{aus$|$sprechen} & (to pronounce)\\
&\textbf{ab$|$sprechen} & (to agree / cancel)\\
&\textbf{nach$|$sprechen} & (to repeat)
\vspace{10pt}
  \end{tabular}
  

  \vspace{-5pt}


\end{verbbox}

%stapeln (to stack / pile up)
\begin{verbbox}\addcontentsline{toc}{subsubsection}{stapeln (to stack / pile up)}
{\Large \bf stapeln} (to stack / pile up) \hfill Type: \textbf{regular} \hfill Auxiliary Verb: \textbf{haben}

\vspace{5pt}
\hrule
\vspace{5pt}

\begin{tabular}{rl}
\multicolumn{2}{l}{\textbf{PRÄSENS}}\\
ich & \textbf{staple}\\
du & \textbf{stapelst}\\
er/sie/es & \textbf{stapelt}\\
wir & \textbf{stapeln}\\
ihr & \textbf{stapelt}\\
sie/Sie & \textbf{stapeln}\\
\end{tabular}
\hfill
\begin{tabular}{rl}
\multicolumn{2}{l}{\textbf{PRÄTERITUM}}\\
ich & \textbf{stapelte}\\
du & \textbf{stapeltest}\\
er/sie/es & \textbf{stapelte}\\
wir & \textbf{stapelten}\\
ihr & \textbf{stapeltet}\\
sie/Sie & \textbf{stapelten}\\
\end{tabular}
\hfill
\begin{tabular}{rl}
\multicolumn{2}{l}{\textbf{IMPERATIV}}\\
& \textbf{staple (du)!}\\
& \textbf{stapeln wir!}\\
& \textbf{stapelt (ihr)!}\\
& \textbf{stapeln Sie!}\\
\end{tabular}

\vspace{10pt}

\begin{tabular}{rl}
\multicolumn{2}{l}{\textbf{PERFEKT}}\\
& haben + \textbf{gestapelt}\\
\end{tabular}
\hfill
\begin{tabular}{rl}
\multicolumn{2}{l}{\textbf{FUTURE I}}\\
& werden + \textbf{stapeln}\\
\end{tabular}
\hfill
\begin{tabular}{rl}
\multicolumn{2}{l}{\textbf{PLUSQUAMPERFEKT}}\\
& hatten + \textbf{gestapelt}\\
\end{tabular}

\vspace{10pt}

\begin{tabular}{lll}
\multicolumn{3}{l}{\textbf{MIT PRÄPOSITIONEN}}\\
& \textbf{stapeln auf} + Akk & (to stack onto)\\
& \textbf{stapeln in} + Akk & (to stack into)\\
\end{tabular}
\hfill
\begin{tabular}{lll}
\multicolumn{3}{l}{\textbf{TRENNBARE VERBEN}}\\
& \textbf{auf$|$stapeln} & (to stack up)\\
\end{tabular}
\vspace{-5pt}
\end{verbbox}

%starten
\begin{verbbox}\addcontentsline{toc}{subsubsection}{starten (to start / launch)}
    {\Large \bf starten} (to start / launch) \hfill Type: \textbf{regular} \hfill Auxiliary Verb: \textbf{haben} 
    
     \vspace{5pt}

    \hrule
    
    \vspace{5pt}

  \begin{tabular}{rl}
     \multicolumn{2}{l}{\textbf{PRÄSENS} }\\
    ich & \textbf{starte}\\
    du & \textbf{startest}\\
    er/sie/es & \textbf{startet}\\
    wir & \textbf{starten}\\
    ihr & \textbf{startet}\\
    sie/Sie & \textbf{starten}\\
  \end{tabular}
  \hfill
  \begin{tabular}{rl}
    \multicolumn{2}{l}{\textbf{PRÄTERITUM} }\\
    ich & \textbf{startete}\\
    du & \textbf{startetest}\\
    er/sie/es & \textbf{startete}\\
    wir & \textbf{starteten}\\
    ihr & \textbf{startetet}\\
    sie/Sie & \textbf{starteten}\\
  \end{tabular}
  \hfill
  \begin{tabular}{rl}
     \multicolumn{2}{l}{\textbf{IMPERATIV} }\\
   & \textbf{starte (du)!}\\
   & \textbf{starten wir!}\\
   & \textbf{startet (ihr)!}\\
   & \textbf{starten Sie!}\\
   & \vspace{10pt}
  \end{tabular}

 \vspace{10pt}

 \begin{tabular}{rl}
     \multicolumn{2}{l}{\textbf{PERFEKT}}\\
   & haben + \textbf{gestartet}\\
  \end{tabular}
\hfill
\begin{tabular}{rl}
    \multicolumn{2}{l}{\textbf{FUTURE I}}\\
   & werden + \textbf{starten}\\
  \end{tabular}
\hfill
\begin{tabular}{rl}
    \multicolumn{2}{l}{\textbf{PLUSQUAMPERFEKT}}\\
   & hatten + \textbf{gestartet}\\
  \end{tabular}

 \vspace{10pt}

\begin{tabular}{lll}
      \multicolumn{3}{l}{\textbf{MIT PRÄPOSITIONEN}}\\
& \textbf{starten zu} + Dat & (to start to / head towards)\\
& \textbf{starten mit} + Dat & (to start with)\\
& \textbf{starten von} + Dat & (to start from)\\
\end{tabular}
  \hfill
\begin{tabular}{lll}
      \multicolumn{3}{l}{\textbf{TRENNBARE VERBEN}}\\
 & \textbf{ab$|$starten} & (to take off / depart)\\
 & \textbf{auf$|$starten} & (to launch / initiate)\\
\vspace{10pt}
  \end{tabular}


\vspace{-5pt}
\end{verbbox}


%staubsaugen (to vacuum)
\begin{verbbox}\addcontentsline{toc}{subsubsection}{staubsaugen (to vacuum)}
{\Large \bf staubsaugen} (to vacuum) \hfill Type: \textbf{mixed / separable} \hfill Auxiliary Verb: \textbf{haben} 

\vspace{5pt}
\hrule
\vspace{5pt}

\begin{tabular}{rl}
\multicolumn{2}{l}{\textbf{PRÄSENS}}\\
ich & \textbf{sauge Staub}\\
du & \textbf{saugst Staub}\\
er/sie/es & \textbf{saugt Staub}\\
wir & \textbf{saugen Staub}\\
ihr & \textbf{saugt Staub}\\
sie/Sie & \textbf{saugen Staub}\\
\end{tabular}
\hfill
\begin{tabular}{rl}
\multicolumn{2}{l}{\textbf{PRÄTERITUM}}\\
ich & \textbf{saugte Staub}\\
du & \textbf{saugtest Staub}\\
er/sie/es & \textbf{saugte Staub}\\
wir & \textbf{saugten Staub}\\
ihr & \textbf{saugtet Staub}\\
sie/Sie & \textbf{saugten Staub}\\
\end{tabular}
\hfill
\begin{tabular}{rl}
\multicolumn{2}{l}{\textbf{IMPERATIV}}\\
& \textbf{saug Staub (du)!}\\
& \textbf{saugen wir Staub!}\\
& \textbf{saugt Staub (ihr)!}\\
& \textbf{saugen Sie Staub!}\\
\end{tabular}

\vspace{10pt}

\begin{tabular}{rl}
\multicolumn{2}{l}{\textbf{PERFEKT}}\\
& haben + \textbf{staubgesaugt}\\
\end{tabular}
\hfill
\begin{tabular}{rl}
\multicolumn{2}{l}{\textbf{FUTURE I}}\\
& werden + \textbf{staubsaugen}\\
\end{tabular}
\hfill
\begin{tabular}{rl}
\multicolumn{2}{l}{\textbf{PLUSQUAMPERFEKT}}\\
& hatten + \textbf{staubgesaugt}\\
\end{tabular}

\vspace{10pt}

\begin{tabular}{lll}
\multicolumn{3}{l}{\textbf{MIT PRÄPOSITIONEN}}\\
& \textbf{staubsaugen} + Akk & (to vacuum something)\\
\end{tabular}
\hfill
\begin{tabular}{lll}
\multicolumn{3}{l}{\textbf{TRENNBARE VERBEN}}\\
& ---\\
\end{tabular}

\vspace{-5pt}
\end{verbbox}

%stechen (to sting / to stab)
\begin{verbbox}\addcontentsline{toc}{subsubsection}{stechen (to sting / to stab)}
    {\Large \bf stechen} (to sting / to stab) \hfill Type: \textbf{irregular} \hfill Auxiliary Verb: \textbf{haben} 
    
     \vspace{5pt}

    \hrule
    
    \vspace{5pt}

  \begin{tabular}{rl}
     \multicolumn{2}{l}{\textbf{PRÄSENS} }\\
    ich & \textbf{steche}\\
    du & \textbf{stichst}\\
    er/sie/es & \textbf{sticht}\\
    wir & \textbf{stechen}\\
    ihr & \textbf{stecht}\\
    sie/Sie & \textbf{stechen}\\
  \end{tabular}
  \hfill
     \begin{tabular}{rl}
    \multicolumn{2}{l}{\textbf{PRÄTERITUM} }\\
    ich & \textbf{stach}\\
    du & \textbf{stachst}\\
    er/sie/es & \textbf{stach}\\
    wir & \textbf{stachen}\\
    ihr & \textbf{stacht}\\
    sie/Sie & \textbf{stachen}\\
  \end{tabular}
  \hfill
  \begin{tabular}{rl}
     \multicolumn{2}{l}{\textbf{IMPERATIV} }\\
   & \textbf{stich (du)!}\\
   & \textbf{stechen wir!}\\
   & \textbf{stecht (ihr)!}\\
   & \textbf{stechen Sie!}\\
   & \vspace{10pt}
  \end{tabular}

    \vspace{10pt}

 \begin{tabular}{rl}
     \multicolumn{2}{l}{\textbf{PERFEKT}}\\
   & haben + \textbf{gestochen}\\
  \end{tabular}
\hfill
   \begin{tabular}{rl}
    \multicolumn{2}{l}{\textbf{FUTURE I} }\\
   & werde + \textbf{stechen}\\
  \end{tabular}
\hfill
   \begin{tabular}{rl}
      \multicolumn{2}{l}{\textbf{PLUSQUAMPERFEKT}}\\
   & hatten + \textbf{gestochen}\\
  \end{tabular}

   \vspace{10pt}

   \begin{tabular}{lll}
      \multicolumn{3}{l}{\textbf{MIT PRÄPOSITIONEN}}\\
& \textbf{stechen in} + Akk & (to sting/stab into)\\
& \textbf{stechen mit} + Dat & (to sting/stab with)\\
\end{tabular}
  \hfill
   \begin{tabular}{lll}
      \multicolumn{3}{l}{\textbf{TRENNBARE VERBEN}}\\
 & \textbf{zu$|$stechen} & (to stab)\\
 & \textbf{ein$|$stechen} & (to stab into / pierce)\\
\vspace{10pt}
  \end{tabular}

  \vspace{-5pt}

\end{verbbox}

%stehen (A1) – to stand
\begin{verbbox}\addcontentsline{toc}{subsubsection}{stehen (to stand)}
    {\Large \bf stehen} (to stand) \hfill Type: \textbf{irregular} \hfill Auxiliary Verb: \textbf{haben/sein} 
    
     \vspace{5pt}

    \hrule
    
    \vspace{5pt}

  \begin{tabular}{rl}
     \multicolumn{2}{l}{\textbf{PRÄSENS} }\\
    ich & \textbf{stehe}\\
    du & \textbf{stehst}\\
    er/sie/es & \textbf{steht}\\
    wir & \textbf{stehen}\\
    ihr & \textbf{steht}\\
    sie/Sie & \textbf{stehen}\\
  \end{tabular}
  \hfill
     \begin{tabular}{rl}
    \multicolumn{2}{l}{\textbf{PRÄTERITUM} }\\
    ich & \textbf{stand}\\
    du & \textbf{standst}\\
    er/sie/es & \textbf{stand}\\
    wir & \textbf{standen}\\
    ihr & \textbf{standet}\\
    sie/Sie & \textbf{standen}\\
  \end{tabular}
  \hfill
  \begin{tabular}{rl}
     \multicolumn{2}{l}{\textbf{IMPERATIV} }\\
   & \textbf{steh(e) (du)!}\\
   & \textbf{stehen wir!}\\
   & \textbf{steht (ihr)!}\\
   & \textbf{stehen Sie!}\\
   & \vspace{10pt}
  \end{tabular}

    \vspace{10pt}

 \begin{tabular}{rl}
     \multicolumn{2}{l}{\textbf{PERFEKT}}\\
  &  haben / sein + \textbf{gestanden}\\
  \end{tabular}
\hfill
   \begin{tabular}{rl}
   
    \multicolumn{2}{l}{\textbf{FUTURE I} }\\
  &  werden + \textbf{stehen}\\
 
  \end{tabular}
\hfill
   \begin{tabular}{rl}
      \multicolumn{2}{l}{\textbf{PLUSQUAMPERFEKT}}\\
  &  hatten / waren + \textbf{gestanden}\\
  \end{tabular}

   \vspace{10pt}

   \begin{tabular}{lll}
      \multicolumn{3}{l}{\textbf{MIT PRÄPOSITIONEN}}\\
& \textbf{stehen für} + Akk & (to represent)\\
& \textbf{stehen zu} + Dat & (to stand by)\\
&\textbf{stehen vor} + Dat & (to stand facing / in front of)\\
&\textbf{stehen in} + Dat & (to stand in)\\
&\textbf{stehen auf} + Dat & (to stand on)\\
&\textbf{stehen an} + Dat & (to stand at)\\
&\textbf{stehen neben} + Dat & (to stand next to)\\
&\textbf{stehen mit} + Dat & (to be connected with)\\
  \end{tabular}
  \hfill
   \begin{tabular}{lll}
      \multicolumn{3}{l}{\textbf{TRENNBARE VERBEN}}\\
 & \textbf{auf$|$stehen} & (to stand up, get up)\\
 & \textbf{an$|$stehen} & (to queue)\\
     & \textbf{fest$|$stehen} & (to be certain)\\
     & \textbf{zu$|$stehen} & (to be entitled to)\\
\vspace{40pt}
  \end{tabular}


  \vspace{-5pt}


\end{verbbox}

%steigen
\begin{verbbox}
\addcontentsline{toc}{subsubsection}{steigen (to climb / rise)}
    {\Large \bf steigen} (to climb / rise) \hfill Type: \textbf{irregular} \hfill Auxiliary Verb: \textbf{sein} 
    
     \vspace{5pt}

    \hrule
    
    \vspace{5pt}

  \begin{tabular}{rl}
     \multicolumn{2}{l}{\textbf{PRÄSENS} }\\
    ich & \textbf{steige}\\
    du & \textbf{steigst}\\
    er/sie/es & \textbf{steigt}\\
    wir & \textbf{steigen}\\
    ihr & \textbf{steigt}\\
    sie/Sie & \textbf{steigen}\\
  \end{tabular}
  \hfill
  \begin{tabular}{rl}
    \multicolumn{2}{l}{\textbf{PRÄTERITUM} }\\
    ich & \textbf{stieg}\\
    du & \textbf{stiegst}\\
    er/sie/es & \textbf{stieg}\\
    wir & \textbf{stiegen}\\
    ihr & \textbf{stiegt}\\
    sie/Sie & \textbf{stiegen}\\
  \end{tabular}
  \hfill
  \begin{tabular}{rl}
     \multicolumn{2}{l}{\textbf{IMPERATIV} }\\
   & \textbf{steige (du)!}\\
   & \textbf{steigen wir!}\\
   & \textbf{steigt (ihr)!}\\
   & \textbf{steigen Sie!}\\
   & \vspace{10pt}
  \end{tabular}

 \vspace{10pt}

 \begin{tabular}{rl}
     \multicolumn{2}{l}{\textbf{PERFEKT}}\\
   & sein + \textbf{gestiegen}\\
  \end{tabular}
\hfill
\begin{tabular}{rl}
    \multicolumn{2}{l}{\textbf{FUTURE I}}\\
   & werden + \textbf{steigen}\\
  \end{tabular}
\hfill
\begin{tabular}{rl}
    \multicolumn{2}{l}{\textbf{PLUSQUAMPERFEKT}}\\
   & waren + \textbf{gestiegen}\\
  \end{tabular}

 \vspace{10pt}

\begin{tabular}{lll}
      \multicolumn{3}{l}{\textbf{MIT PRÄPOSITIONEN}}\\
& \textbf{steigen auf} + Akk & (to climb onto / rise to)\\
& \textbf{steigen in} + Akk & (to climb into / rise into)\\
& \textbf{steigen aus} + Dat & (to climb out of)\\
\vspace{20pt}
\end{tabular}
  \hfill
\begin{tabular}{lll}
      \multicolumn{3}{l}{\textbf{TRENNBARE VERBEN}}\\
 & \textbf{auf$|$steigen} & (to take off / rise up)\\
 & \textbf{ein$|$steigen} & (to get on / board)\\
 & \textbf{aus$|$steigen} & (to get off / disembark)\\
 & \textbf{ab$|$steigen} & (to descend / go down)\\
  & \textbf{um$|$steigen} & (to transfer)\\
  \end{tabular}


  \vspace{-5pt}
  \end{verbbox}

%stellen (B1) – to place
\begin{verbbox-acc}\addcontentsline{toc}{subsubsection}{stellen (to place)}
    {\Large \bf stellen} (to place) \hfill Type: \textbf{regular, + Akkusativ} \hfill Auxiliary Verb: \textbf{haben} 
    
     \vspace{5pt}

    \hrule
    
    \vspace{5pt}

  \begin{tabular}{rl}
     \multicolumn{2}{l}{\textbf{PRÄSENS} }\\
    ich & \textbf{stelle}\\
    du & \textbf{stellst}\\
    er/sie/es & \textbf{stellt}\\
    wir & \textbf{stellen}\\
    ihr & \textbf{stellt}\\
    sie/Sie & \textbf{stellen}\\
  \end{tabular}
  \hfill
     \begin{tabular}{rl}
    \multicolumn{2}{l}{\textbf{PRÄTERITUM} }\\
    ich & \textbf{stellte}\\
    du & \textbf{stelltest}\\
    er/sie/es & \textbf{stellte}\\
    wir & \textbf{stellten}\\
    ihr & \textbf{stelltet}\\
    sie/Sie & \textbf{stellten}\\
  \end{tabular}
  \hfill
  \begin{tabular}{rl}
     \multicolumn{2}{l}{\textbf{IMPERATIV} }\\
   & \textbf{stell(e) (du)!}\\
   & \textbf{stellen wir!}\\
   & \textbf{stellt (ihr)!}\\
   & \textbf{stellen Sie!}\\
   & \vspace{10pt}
  \end{tabular}

    \vspace{10pt}

 \begin{tabular}{rl}
     \multicolumn{2}{l}{\textbf{PERFEKT}}\\
   & haben + \textbf{gestellt}\\
  \end{tabular}
\hfill
   \begin{tabular}{rl}
   
    \multicolumn{2}{l}{\textbf{FUTURE I} }\\
   & werden + \textbf{stellen}\\
 
  \end{tabular}
\hfill
   \begin{tabular}{rl}
      \multicolumn{2}{l}{\textbf{PLUSQUAMPERFEKT}}\\
   & hatten + \textbf{gestellt}\\
  \end{tabular}
  
\vspace{10pt}

   \begin{tabular}{lll}
      \multicolumn{3}{l}{\textbf{MIT PRÄPOSITIONEN}}\\
&\textbf{sich stellen auf} + Akk & (to prepare for / brace for)\\
& \textbf{sich stellen gegan} + Akk & (to oppose)\\
& \textbf{sich stellen vor} + Akk & (to side with)\\
&\textbf{sich stellen zu} + Dat & (to stand in front of / defend)\\
\vspace{60pt}
  \end{tabular}
  \hfill
   \begin{tabular}{lll}
      \multicolumn{3}{l}{\textbf{TRENNBARE VERBEN}}\\
 & \textbf{auf$|$stellen} & (to set up / put up)\\
&\textbf{hin$|$stellen} & (to put (something) there)\\
&\textbf{um$|$stellen} & (to rearrange)\\
&\textbf{vor$|$stellen} & (to introduce)\\
&\textbf{ver$|$stellen} & (to block)\\
&\textbf{ein$|$stellen} & (to adjust, hire)\\
&\textbf{her$|$stellen} & (to produce)\\
&\textbf{dar$|$stellen} & (to portray)\\
&\textbf{fest$|$stellen} & (to determine, establish)\\
&\textbf{zurück$|$stellen} & (to postpone)
  \end{tabular}


  \vspace{-5pt}


\end{verbbox-acc}

%stimmen (to be correct / to vote)
\begin{verbbox-dat}\addcontentsline{toc}{subsubsection}{stimmen (to be correct / to vote)}
    {\Large \bf stimmen} (to be correct / to vote) \hfill Type: \textbf{regular} \hfill Auxiliary Verb: \textbf{haben} 
    
     \vspace{5pt}

    \hrule
    
    \vspace{5pt}

  \begin{tabular}{rl}
     \multicolumn{2}{l}{\textbf{PRÄSENS} }\\
    ich & \textbf{stimme}\\
    du & \textbf{stimmst}\\
    er/sie/es & \textbf{stimmt}\\
    wir & \textbf{stimmen}\\
    ihr & \textbf{stimmt}\\
    sie/Sie & \textbf{stimmen}\\
  \end{tabular}
  \hfill
     \begin{tabular}{rl}
    \multicolumn{2}{l}{\textbf{PRÄTERITUM} }\\
    ich & \textbf{stimmte}\\
    du & \textbf{stimmtest}\\
    er/sie/es & \textbf{stimmte}\\
    wir & \textbf{stimmten}\\
    ihr & \textbf{stimmtet}\\
    sie/Sie & \textbf{stimmten}\\
  \end{tabular}
  \hfill
  \begin{tabular}{rl}
     \multicolumn{2}{l}{\textbf{IMPERATIV} }\\
   & \textbf{stimm(e) (du)!}\\
   & \textbf{stimmen wir!}\\
   & \textbf{stimmt (ihr)!}\\
   & \textbf{stimmen Sie!}\\
   & \vspace{10pt}
  \end{tabular}

    \vspace{10pt}

 \begin{tabular}{rl}
     \multicolumn{2}{l}{\textbf{PERFEKT}}\\
   & haben + \textbf{gestimmt}\\
  \end{tabular}
\hfill
   \begin{tabular}{rl}
   
    \multicolumn{2}{l}{\textbf{FUTURE I} }\\
   & werden + \textbf{stimmen}\\
 
  \end{tabular}
\hfill
   \begin{tabular}{rl}
      \multicolumn{2}{l}{\textbf{PLUSQUAMPERFEKT}}\\
   & hatten + \textbf{gestimmt}\\
  \end{tabular}

   \vspace{10pt}

   \begin{tabular}{lll}
      \multicolumn{3}{l}{\textbf{MIT PRÄPOSITIONEN}}\\
& \textbf{stimmen für} + Akk & (to vote for)\\
& \textbf{stimmen gegen} + Akk & (to vote against)\\
\end{tabular}
  \hfill
   \begin{tabular}{lll}
      \multicolumn{3}{l}{\textbf{TRENNBARE VERBEN}}\\
 & \textbf{ab$|$stimmen} & (to vote / coordinate)\\
 & \textbf{überein$|$stimmen} & (to agree / correspond)\\
  \end{tabular}


  \vspace{-5pt}


\end{verbbox-dat}

%stoppen (to stop)
\begin{verbbox}\addcontentsline{toc}{subsubsection}{stoppen (to stop)}
    {\Large \bf stoppen} (to stop) \hfill Type: \textbf{regular} \hfill Auxiliary Verb: \textbf{haben} 
    
     \vspace{5pt}

    \hrule
    
    \vspace{5pt}

  \begin{tabular}{rl}
     \multicolumn{2}{l}{\textbf{PRÄSENS} }\\
    ich & \textbf{stoppe}\\
    du & \textbf{stoppst}\\
    er/sie/es & \textbf{stoppt}\\
    wir & \textbf{stoppen}\\
    ihr & \textbf{stoppt}\\
    sie/Sie & \textbf{stoppen}\\
  \end{tabular}
  \hfill
     \begin{tabular}{rl}
    \multicolumn{2}{l}{\textbf{PRÄTERITUM} }\\
    ich & \textbf{stoppte}\\
    du & \textbf{stopptest}\\
    er/sie/es & \textbf{stoppte}\\
    wir & \textbf{stoppten}\\
    ihr & \textbf{stopptet}\\
    sie/Sie & \textbf{stoppten}\\
  \end{tabular}
  \hfill
  \begin{tabular}{rl}
     \multicolumn{2}{l}{\textbf{IMPERATIV} }\\
   & \textbf{stopp(e) (du)!}\\
   & \textbf{stoppen wir!}\\
   & \textbf{stoppt (ihr)!}\\
   & \textbf{stoppen Sie!}\\
   & \vspace{10pt}
  \end{tabular}

    \vspace{10pt}

 \begin{tabular}{rl}
     \multicolumn{2}{l}{\textbf{PERFEKT}}\\
   & haben + \textbf{gestoppt}\\
  \end{tabular}
\hfill
   \begin{tabular}{rl}
    \multicolumn{2}{l}{\textbf{FUTURE I} }\\
   & werden + \textbf{stoppen}\\
  \end{tabular}
\hfill
   \begin{tabular}{rl}
      \multicolumn{2}{l}{\textbf{PLUSQUAMPERFEKT}}\\
   & hatten + \textbf{gestoppt}\\
  \end{tabular}

   \vspace{10pt}

   \begin{tabular}{lll}
      \multicolumn{3}{l}{\textbf{MIT PRÄPOSITIONEN}}\\
& \textbf{stoppen vor} + Dat & (to stop in front of)\\
& \textbf{stoppen bei} + Dat & (to stop at)\\
\end{tabular}
  \hfill
   \begin{tabular}{lll}
      \multicolumn{3}{l}{\textbf{TRENNBARE VERBEN}}\\
 & \textbf{an$|$stoppen} & (to stop / halt briefly)\\
\vspace{10pt}
  \end{tabular}

  \vspace{-5pt}

\end{verbbox}


%streichen
\begin{verbbox}\addcontentsline{toc}{subsubsection}{streichen (to paint / stroke)}
    {\Large \bf streichen} (to paint / stroke) \hfill Type: \textbf{irregular} \hfill Auxiliary Verb: \textbf{haben} 
    
     \vspace{5pt}

    \hrule
    
    \vspace{5pt}

  \begin{tabular}{rl}
     \multicolumn{2}{l}{\textbf{PRÄSENS} }\\
    ich & \textbf{streiche}\\
    du & \textbf{streichst}\\
    er/sie/es & \textbf{streicht}\\
    wir & \textbf{streichen}\\
    ihr & \textbf{streicht}\\
    sie/Sie & \textbf{streichen}\\
  \end{tabular}
  \hfill
  \begin{tabular}{rl}
    \multicolumn{2}{l}{\textbf{PRÄTERITUM} }\\
    ich & \textbf{strich}\\
    du & \textbf{strichst}\\
    er/sie/es & \textbf{strich}\\
    wir & \textbf{strichen}\\
    ihr & \textbf{stricht}\\
    sie/Sie & \textbf{strichen}\\
  \end{tabular}
  \hfill
  \begin{tabular}{rl}
     \multicolumn{2}{l}{\textbf{IMPERATIV} }\\
   & \textbf{streiche (du)!}\\
   & \textbf{streichen wir!}\\
   & \textbf{streicht (ihr)!}\\
   & \textbf{streichen Sie!}\\
   & \vspace{10pt}
  \end{tabular}

 \vspace{10pt}

 \begin{tabular}{rl}
     \multicolumn{2}{l}{\textbf{PERFEKT}}\\
   & haben + \textbf{gestrichen}\\
  \end{tabular}
\hfill
\begin{tabular}{rl}
    \multicolumn{2}{l}{\textbf{FUTURE I}}\\
   & werden + \textbf{streichen}\\
  \end{tabular}
\hfill
\begin{tabular}{rl}
    \multicolumn{2}{l}{\textbf{PLUSQUAMPERFEKT}}\\
   & hatten + \textbf{gestrichen}\\
  \end{tabular}

 \vspace{10pt}

\begin{tabular}{lll}
      \multicolumn{3}{l}{\textbf{MIT PRÄPOSITIONEN}}\\
& \textbf{streichen über} + Akk & (to stroke over / to paint over)\\
& \textbf{streichen mit} + Dat & (to paint / smear with)\\
\vspace{10pt}
\end{tabular}
  \hfill
\begin{tabular}{lll}
      \multicolumn{3}{l}{\textbf{TRENNBARE VERBEN}}\\
 & \textbf{ab$|$streichen} & (to cancel / remove / coat off)\\
 & \textbf{ein$|$streichen} & (to apply / cover in)\\
 & \textbf{auf$|$streichen} & (to spread / coat onto)\\
  \end{tabular}


\vspace{-5pt}
\end{verbbox}

%studieren (A1) – to study
\begin{verbbox}\addcontentsline{toc}{subsubsection}{studieren (to study)}
    {\Large \bf studieren} (to study) \hfill Type: \textbf{regular} \hfill Auxiliary Verb: \textbf{haben} 
    
     \vspace{5pt}

    \hrule
    
    \vspace{5pt}

  \begin{tabular}{rl}
     \multicolumn{2}{l}{\textbf{PRÄSENS} }\\
    ich & \textbf{studiere}\\
    du & \textbf{studierst}\\
    er/sie/es & \textbf{studiert}\\
    wir & \textbf{studieren}\\
    ihr & \textbf{studiert}\\
    sie/Sie & \textbf{studieren}\\
  \end{tabular}
  \hfill
     \begin{tabular}{rl}
    \multicolumn{2}{l}{\textbf{PRÄTERITUM} }\\
    ich & \textbf{studierte}\\
    du & \textbf{studiertest}\\
    er/sie/es & \textbf{studierte}\\
    wir & \textbf{studierten}\\
    ihr & \textbf{studiertet}\\
    sie/Sie & \textbf{studierten}\\
  \end{tabular}
  \hfill
  \begin{tabular}{rl}
     \multicolumn{2}{l}{\textbf{IMPERATIV} }\\
   & \textbf{studier(e) (du)!}\\
   & \textbf{studieren wir!}\\
   & \textbf{studiert (ihr)!}\\
   & \textbf{studieren Sie!}\\
   & \vspace{10pt}
  \end{tabular}

    \vspace{10pt}

 \begin{tabular}{rl}
     \multicolumn{2}{l}{\textbf{PERFEKT}}\\
   & haben + \textbf{studiert}\\
  \end{tabular}
\hfill
   \begin{tabular}{rl}
   
    \multicolumn{2}{l}{\textbf{FUTURE I} }\\
   & werden + \textbf{studieren}\\
 
  \end{tabular}
\hfill
   \begin{tabular}{rl}
      \multicolumn{2}{l}{\textbf{PLUSQUAMPERFEKT}}\\
   & hatten + \textbf{studiert}\\
  \end{tabular}
  
  \vspace{10pt}

\begin{tabular}{lll}
\multicolumn{3}{l}{\textbf{MIT PRÄPOSITIONEN}}\\
& ---\\
\end{tabular}
\hfill
\begin{tabular}{lll}
\multicolumn{3}{l}{\textbf{TRENNBARE VERBEN}}\\
& ---\\
\end{tabular}

  \vspace{-5pt}

\end{verbbox}
%suchen (A1) – to search
\begin{verbbox-acc}\addcontentsline{toc}{subsubsection}{suchen (to search)}
    {\Large \bf suchen} (to search) \hfill Type: \textbf{regular, + Akkusativ} \hfill Auxiliary Verb: \textbf{haben} 
    
     \vspace{5pt}

    \hrule
    
    \vspace{5pt}

  \begin{tabular}{rl}
     \multicolumn{2}{l}{\textbf{PRÄSENS} }\\
    ich & \textbf{suche}\\
    du & \textbf{suchst}\\
    er/sie/es & \textbf{sucht}\\
    wir & \textbf{suchen}\\
    ihr & \textbf{sucht}\\
    sie/Sie & \textbf{suchen}\\
  \end{tabular}
  \hfill
     \begin{tabular}{rl}
    \multicolumn{2}{l}{\textbf{PRÄTERITUM} }\\
    ich & \textbf{suchte}\\
    du & \textbf{suchtest}\\
    er/sie/es & \textbf{suchte}\\
    wir & \textbf{suchten}\\
    ihr & \textbf{suchtet}\\
    sie/Sie & \textbf{suchten}\\
  \end{tabular}
  \hfill
  \begin{tabular}{rl}
     \multicolumn{2}{l}{\textbf{IMPERATIV} }\\
   & \textbf{such(e) (du)!}\\
   & \textbf{suchen wir!}\\
   & \textbf{sucht (ihr)!}\\
   & \textbf{suchen Sie!}\\
   & \vspace{10pt}
  \end{tabular}

    \vspace{10pt}

 \begin{tabular}{rl}
     \multicolumn{2}{l}{\textbf{PERFEKT}}\\
   & haben + \textbf{gesucht}\\
  \end{tabular}
\hfill
   \begin{tabular}{rl}
   
    \multicolumn{2}{l}{\textbf{FUTURE I} }\\
   & werden + \textbf{suchen}\\
 
  \end{tabular}
\hfill
   \begin{tabular}{rl}
      \multicolumn{2}{l}{\textbf{PLUSQUAMPERFEKT}}\\
   & hatten + \textbf{gesucht}\\
  \end{tabular}
  
  \vspace{-5pt}


\end{verbbox-acc}

\subsection{T Verbs}

%tauschen (to exchange / swap)
\begin{verbbox}\addcontentsline{toc}{subsubsection}{tauschen (to exchange / swap)}
{\Large \bf tauschen} (to exchange / swap) \hfill Type: \textbf{regular} \hfill Auxiliary Verb: \textbf{haben}

\vspace{5pt}
\hrule
\vspace{5pt}

\begin{tabular}{rl}
\multicolumn{2}{l}{\textbf{PRÄSENS}}\\
ich & \textbf{tausche}\\
du & \textbf{tauschst}\\
er/sie/es & \textbf{tauscht}\\
wir & \textbf{tauschen}\\
ihr & \textbf{tauscht}\\
sie/Sie & \textbf{tauschen}\\
\end{tabular}
\hfill
\begin{tabular}{rl}
\multicolumn{2}{l}{\textbf{PRÄTERITUM}}\\
ich & \textbf{tauschte}\\
du & \textbf{tauschtest}\\
er/sie/es & \textbf{tauschte}\\
wir & \textbf{tauschten}\\
ihr & \textbf{tauschtet}\\
sie/Sie & \textbf{tauschten}\\
\end{tabular}
\hfill
\begin{tabular}{rl}
\multicolumn{2}{l}{\textbf{IMPERATIV}}\\
& \textbf{tausche (du)!}\\
& \textbf{tauschen wir!}\\
& \textbf{tauscht (ihr)!}\\
& \textbf{tauschen Sie!}\\
\end{tabular}

\vspace{10pt}

\begin{tabular}{rl}
\multicolumn{2}{l}{\textbf{PERFEKT}}\\
& haben + \textbf{getauscht}\\
\end{tabular}
\hfill
\begin{tabular}{rl}
\multicolumn{2}{l}{\textbf{FUTURE I}}\\
& werden + \textbf{tauschen}\\
\end{tabular}
\hfill
\begin{tabular}{rl}
\multicolumn{2}{l}{\textbf{PLUSQUAMPERFEKT}}\\
& hatten + \textbf{getauscht}\\
\end{tabular}

\vspace{10pt}

\begin{tabular}{lll}
\multicolumn{3}{l}{\textbf{MIT PRÄPOSITIONEN}}\\
& \textbf{tauschen gegen} + Akk & (to exchange for)\\
& \textbf{tauschen mit} + Dat & (to exchange with)\\
\end{tabular}
\hfill
\begin{tabular}{lll}
\multicolumn{3}{l}{\textbf{TRENNBARE VERBEN}}\\
& \textbf{aus$|$tauschen} & (to exchange / replace)\\
& \textbf{um$|$tauschen} & (to return (a purchased item))\\
\end{tabular}
\vspace{-5pt}
\end{verbbox}

%teilen (to share / divide)
\begin{verbbox}\addcontentsline{toc}{subsubsection}{teilen (to share / divide)}
{\Large \bf teilen} (to share / divide) \hfill Type: \textbf{regular} \hfill Auxiliary Verb: \textbf{haben}

\vspace{5pt}
\hrule
\vspace{5pt}

\begin{tabular}{rl}
\multicolumn{2}{l}{\textbf{PRÄSENS}}\\
ich & \textbf{teile}\\
du & \textbf{teilst}\\
er/sie/es & \textbf{teilt}\\
wir & \textbf{teilen}\\
ihr & \textbf{teilt}\\
sie/Sie & \textbf{teilen}\\
\end{tabular}
\hfill
\begin{tabular}{rl}
\multicolumn{2}{l}{\textbf{PRÄTERITUM}}\\
ich & \textbf{teilte}\\
du & \textbf{teiltest}\\
er/sie/es & \textbf{teilte}\\
wir & \textbf{teilten}\\
ihr & \textbf{teiltet}\\
sie/Sie & \textbf{teilten}\\
\end{tabular}
\hfill
\begin{tabular}{rl}
\multicolumn{2}{l}{\textbf{IMPERATIV}}\\
& \textbf{teil(e) (du)!}\\
& \textbf{teilen wir!}\\
& \textbf{teilt (ihr)!}\\
& \textbf{teilen Sie!}\\
\end{tabular}

\vspace{10pt}

\begin{tabular}{rl}
\multicolumn{2}{l}{\textbf{PERFEKT}}\\
& haben + \textbf{geteilt}\\
\end{tabular}
\hfill
\begin{tabular}{rl}
\multicolumn{2}{l}{\textbf{FUTURE I}}\\
& werden + \textbf{teilen}\\
\end{tabular}
\hfill
\begin{tabular}{rl}
\multicolumn{2}{l}{\textbf{PLUSQUAMPERFEKT}}\\
& hatten + \textbf{geteilt}\\
\end{tabular}

\vspace{10pt}

\begin{tabular}{lll}
\multicolumn{3}{l}{\textbf{MIT PRÄPOSITIONEN}}\\
& \textbf{teilen mit} + Dat & (to share with)\\
& \textbf{teilen in} + Akk & (to divide into)\\
\end{tabular}
\hfill
\begin{tabular}{lll}
\multicolumn{3}{l}{\textbf{TRENNBARE VERBEN}}\\
& \textbf{auf$|$teilen} & (to split up / divide)\\
& \textbf{ab$|$teilen} & (to allocate / section off)\\
\end{tabular}
\vspace{-5pt}
\end{verbbox}

%tragen (A2) – to wear, carry
\begin{verbbox-acc}\addcontentsline{toc}{subsubsection}{tragen (to wear, carry)}
    {\Large \bf tragen} (to wear, carry) \hfill Type: \textbf{irregular, + Akkusativ} \hfill Auxiliary Verb: \textbf{haben} 
    
     \vspace{5pt}

    \hrule
    
    \vspace{5pt}

  \begin{tabular}{rl}
     \multicolumn{2}{l}{\textbf{PRÄSENS} }\\
    ich & \textbf{trage}\\
    du & \textbf{trägst}\\
    er/sie/es & \textbf{trägt}\\
    wir & \textbf{tragen}\\
    ihr & \textbf{tragt}\\
    sie/Sie & \textbf{tragen}\\
  \end{tabular}
  \hfill
     \begin{tabular}{rl}
    \multicolumn{2}{l}{\textbf{PRÄTERITUM} }\\
    ich & \textbf{trug}\\
    du & \textbf{trugst}\\
    er/sie/es & \textbf{trug}\\
    wir & \textbf{trugen}\\
    ihr & \textbf{trugt}\\
    sie/Sie & \textbf{trugen}\\
  \end{tabular}
  \hfill
  \begin{tabular}{rl}
     \multicolumn{2}{l}{\textbf{IMPERATIV} }\\
   & \textbf{trag(e) (du)!}\\
   & \textbf{tragen wir!}\\
   & \textbf{tragt (ihr)!}\\
   & \textbf{tragen Sie!}\\
   & \vspace{10pt}
  \end{tabular}

    \vspace{10pt}

 \begin{tabular}{rl}
     \multicolumn{2}{l}{\textbf{PERFEKT}}\\
   & haben + \textbf{getragen}\\
  \end{tabular}
\hfill
   \begin{tabular}{rl}
   
    \multicolumn{2}{l}{\textbf{FUTURE I} }\\
   & werden + \textbf{tragen}\\
 
  \end{tabular}
\hfill
   \begin{tabular}{rl}
      \multicolumn{2}{l}{\textbf{PLUSQUAMPERFEKT}}\\
   & hatten + \textbf{getragen}\\
  \end{tabular}
  
  \vspace{-5pt}


\end{verbbox-acc}


%trainieren - to train
\begin{verbbox}\addcontentsline{toc}{subsubsection}{trainieren (to train / practice)}
    {\Large \bf trainieren} (to train / practice) \hfill Type: \textbf{regular} \hfill Auxiliary Verb: \textbf{haben} 
    
     \vspace{5pt}
    \hrule
    \vspace{5pt}

  \begin{tabular}{rl}
     \multicolumn{2}{l}{\textbf{PRÄSENS} }\\
    ich & \textbf{trainiere}\\
    du & \textbf{trainierst}\\
    er/sie/es & \textbf{trainiert}\\
    wir & \textbf{trainieren}\\
    ihr & \textbf{trainiert}\\
    sie/Sie & \textbf{trainieren}\\
  \end{tabular}
  \hfill
     \begin{tabular}{rl}
    \multicolumn{2}{l}{\textbf{PRÄTERITUM} }\\
    ich & \textbf{trainierte}\\
    du & \textbf{trainiertest}\\
    er/sie/es & \textbf{trainierte}\\
    wir & \textbf{trainierten}\\
    ihr & \textbf{trainiertet}\\
    sie/Sie & \textbf{trainierten}\\
  \end{tabular}
  \hfill
  \begin{tabular}{rl}
     \multicolumn{2}{l}{\textbf{IMPERATIV} }\\
   & \textbf{trainier(e) (du)!}\\
   & \textbf{trainieren wir!}\\
   & \textbf{trainiert (ihr)!}\\
   & \textbf{trainieren Sie!}\\
   & \vspace{10pt}
  \end{tabular}

 \vspace{10pt}

 \begin{tabular}{rl}
     \multicolumn{2}{l}{\textbf{PERFEKT}}\\
   & haben + \textbf{trainiert}\\
  \end{tabular}
\hfill
   \begin{tabular}{rl}
    \multicolumn{2}{l}{\textbf{FUTURE I}}\\
   & werden + \textbf{trainieren}\\
  \end{tabular}
\hfill
   \begin{tabular}{rl}
      \multicolumn{2}{l}{\textbf{PLUSQUAMPERFEKT}}\\
   & hatten + \textbf{trainiert}\\
  \end{tabular}

   \vspace{10pt}

   \begin{tabular}{lll}
      \multicolumn{3}{l}{\textbf{MIT PRÄPOSITIONEN}}\\
& \textbf{trainieren für} + Akk & (to train for)\\
& \textbf{trainieren mit} + Dat & (to train with)\\
& \textbf{trainieren gegen} + Akk & (to train against / play against)\\
\end{tabular}
  \hfill
   \begin{tabular}{lll}
      \multicolumn{3}{l}{\textbf{TRENNBARE VERBEN}}\\
 & \textbf{ab$|$trainieren} & (to train off / reduce through training)\\
\vspace{20pt}
  \end{tabular}
  \vspace{-5pt}
\end{verbbox}

%träumen - to dream
\begin{verbbox}\addcontentsline{toc}{subsubsection}{träumen (to dream)}
    {\Large \bf träumen} (to dream) \hfill Type: \textbf{regular} \hfill Auxiliary Verb: \textbf{haben} 
    
     \vspace{5pt}

    \hrule
    
    \vspace{5pt}

  \begin{tabular}{rl}
     \multicolumn{2}{l}{\textbf{PRÄSENS} }\\
    ich & \textbf{träume}\\
    du & \textbf{träumst}\\
    er/sie/es & \textbf{träumt}\\
    wir & \textbf{träumen}\\
    ihr & \textbf{träumt}\\
    sie/Sie & \textbf{träumen}\\
  \end{tabular}
  \hfill
     \begin{tabular}{rl}
    \multicolumn{2}{l}{\textbf{PRÄTERITUM} }\\
    ich & \textbf{träumte}\\
    du & \textbf{träumtest}\\
    er/sie/es & \textbf{träumte}\\
    wir & \textbf{träumten}\\
    ihr & \textbf{träumtet}\\
    sie/Sie & \textbf{träumten}\\
  \end{tabular}
  \hfill
  \begin{tabular}{rl}
     \multicolumn{2}{l}{\textbf{IMPERATIV} }\\
   & \textbf{träum(e) (du)!}\\
   & \textbf{träumen wir!}\\
   & \textbf{träumt (ihr)!}\\
   & \textbf{träumen Sie!}\\
   & \vspace{10pt}
  \end{tabular}

    \vspace{10pt}

 \begin{tabular}{rl}
     \multicolumn{2}{l}{\textbf{PERFEKT}}\\
   & haben + \textbf{geträumt}\\
  \end{tabular}
\hfill
   \begin{tabular}{rl}
   
    \multicolumn{2}{l}{\textbf{FUTURE I} }\\
   & werden + \textbf{träumen}\\
 
  \end{tabular}
\hfill
   \begin{tabular}{rl}
      \multicolumn{2}{l}{\textbf{PLUSQUAMPERFEKT}}\\
   & hatten + \textbf{geträumt}\\
  \end{tabular}
  
  \vspace{10pt}


\begin{tabular}{lll}
\multicolumn{3}{l}{\textbf{MIT PRÄPOSITIONEN}}\\
& ---\\
\end{tabular}
\hfill
\begin{tabular}{lll}
\multicolumn{3}{l}{\textbf{TRENNBARE VERBEN}}\\
& ---\\
\end{tabular}

  \vspace{-5pt}

\end{verbbox}

%treten (A2) - to step
\begin{verbbox}\addcontentsline{toc}{subsubsection}{treten (to step)}
    {\Large \bf treten} (to step) \hfill Type: \textbf{irregular} \hfill Auxiliary Verb: \textbf{sein} 
    
     \vspace{5pt}

    \hrule
    
    \vspace{5pt}

  \begin{tabular}{rl}
     \multicolumn{2}{l}{\textbf{PRÄSENS} }\\
    ich & \textbf{trete}\\
    du & \textbf{trittst}\\
    er/sie/es & \textbf{tritt}\\
    wir & \textbf{treten}\\
    ihr & \textbf{tretet}\\
    sie/Sie & \textbf{treten}\\
  \end{tabular}
  \hfill
     \begin{tabular}{rl}
    \multicolumn{2}{l}{\textbf{PRÄTERITUM} }\\
    ich & \textbf{trat}\\
    du & \textbf{tratst}\\
    er/sie/es & \textbf{trat}\\
    wir & \textbf{traten}\\
    ihr & \textbf{tratet}\\
    sie/Sie & \textbf{traten}\\
  \end{tabular}
  \hfill
  \begin{tabular}{rl}
     \multicolumn{2}{l}{\textbf{IMPERATIV} }\\
   & \textbf{tritt (du)!}\\
   & \textbf{treten wir!}\\
   & \textbf{tretet (ihr)!}\\
   & \textbf{treten Sie!}\\
   & \vspace{10pt}
  \end{tabular}

    \vspace{10pt}

 \begin{tabular}{rl}
     \multicolumn{2}{l}{\textbf{PERFEKT}}\\
   & sein + \textbf{getreten}\\
  \end{tabular}
\hfill
   \begin{tabular}{rl}
   
    \multicolumn{2}{l}{\textbf{FUTURE I} }\\
   & werden + \textbf{treten}\\
 
  \end{tabular}
\hfill
   \begin{tabular}{rl}
      \multicolumn{2}{l}{\textbf{PLUSQUAMPERFEKT}}\\
   & waren + \textbf{getreten}\\
  \end{tabular}
  
  \vspace{10pt}



\begin{tabular}{lll}
\multicolumn{3}{l}{\textbf{MIT PRÄPOSITIONEN}}\\
& ---\\
\end{tabular}
\hfill
\begin{tabular}{lll}
\multicolumn{3}{l}{\textbf{TRENNBARE VERBEN}}\\
& ---\\
\end{tabular}

  \vspace{-5pt}

\end{verbbox}

%trinken (A1) – to drink
\begin{verbbox-acc}\addcontentsline{toc}{subsubsection}{trinken (to drink)}
    {\Large \bf trinken} (to drink) \hfill Type: \textbf{irregular, + Akkusativ} \hfill Auxiliary Verb: \textbf{haben} 
    
     \vspace{5pt}

    \hrule
    
    \vspace{5pt}

  \begin{tabular}{rl}
     \multicolumn{2}{l}{\textbf{PRÄSENS} }\\
    ich & \textbf{trinke}\\
    du & \textbf{trinkst}\\
    er/sie/es & \textbf{trinkt}\\
    wir & \textbf{trinken}\\
    ihr & \textbf{trinkt}\\
    sie/Sie & \textbf{trinken}\\
  \end{tabular}
  \hfill
     \begin{tabular}{rl}
    \multicolumn{2}{l}{\textbf{PRÄTERITUM} }\\
    ich & \textbf{trank}\\
    du & \textbf{trankest}\\
    er/sie/es & \textbf{trank}\\
    wir & \textbf{tranken}\\
    ihr & \textbf{trankt}\\
    sie/Sie & \textbf{tranken}\\
  \end{tabular}
  \hfill
  \begin{tabular}{rl}
     \multicolumn{2}{l}{\textbf{IMPERATIV} }\\
   & \textbf{trink(e) (du)!}\\
   & \textbf{trinken wir!}\\
   & \textbf{trinkt (ihr)!}\\
   & \textbf{trinken Sie!}\\
   & \vspace{10pt}
  \end{tabular}

    \vspace{10pt}

 \begin{tabular}{rl}
     \multicolumn{2}{l}{\textbf{PERFEKT}}\\
   & haben + \textbf{getrunken}\\
  \end{tabular}
\hfill
   \begin{tabular}{rl}
   
    \multicolumn{2}{l}{\textbf{FUTURE I} }\\
   & werden + \textbf{trinken}\\
 
  \end{tabular}
\hfill
   \begin{tabular}{rl}
      \multicolumn{2}{l}{\textbf{PLUSQUAMPERFEKT}}\\
   & hatten + \textbf{getrunken}\\
  \end{tabular}
  
  \vspace{-5pt}


\end{verbbox-acc}
%treffen (A2) – to meet
\begin{verbbox-acc}\addcontentsline{toc}{subsubsection}{treffen (to meet)}
    {\Large \bf treffen} (to meet) \hfill Type: \textbf{irregular, + Akkusativ} \hfill Auxiliary Verb: \textbf{haben} 
    
     \vspace{5pt}

    \hrule
    
    \vspace{5pt}

  \begin{tabular}{rl}
     \multicolumn{2}{l}{\textbf{PRÄSENS} }\\
    ich & \textbf{treffe}\\
    du & \textbf{triffst}\\
    er/sie/es & \textbf{trifft}\\
    wir & \textbf{treffen}\\
    ihr & \textbf{trefft}\\
    sie/Sie & \textbf{treffen}\\
  \end{tabular}
  \hfill
     \begin{tabular}{rl}
    \multicolumn{2}{l}{\textbf{PRÄTERITUM} }\\
    ich & \textbf{traf}\\
    du & \textbf{trafst}\\
    er/sie/es & \textbf{traf}\\
    wir & \textbf{trafen}\\
    ihr & \textbf{traft}\\
    sie/Sie & \textbf{trafen}\\
  \end{tabular}
  \hfill
  \begin{tabular}{rl}
     \multicolumn{2}{l}{\textbf{IMPERATIV} }\\
   & \textbf{triff (du)!}\\
   & \textbf{treffen wir!}\\
   & \textbf{trefft (ihr)!}\\
   & \textbf{treffen Sie!}\\
   & \vspace{10pt}
  \end{tabular}

    \vspace{10pt}

 \begin{tabular}{rl}
     \multicolumn{2}{l}{\textbf{PERFEKT}}\\
  &  haben + \textbf{getroffen}\\
  \end{tabular}
\hfill
   \begin{tabular}{rl}
   
    \multicolumn{2}{l}{\textbf{FUTURE I} }\\
  &  werden + \textbf{treffen}\\
 
  \end{tabular}
\hfill
   \begin{tabular}{rl}
      \multicolumn{2}{l}{\textbf{PLUSQUAMPERFEKT}}\\
   & hatten + \textbf{getroffen}\\
  \end{tabular}
  
  \vspace{-5pt}


\end{verbbox-acc}

%trösten (to comfort)
\begin{verbbox}\addcontentsline{toc}{subsubsection}{trösten (to comfort)}
{\Large \bf trösten} (to comfort) \hfill Type: \textbf{regular} \hfill Auxiliary Verb: \textbf{haben}

\vspace{5pt}
\hrule
\vspace{5pt}

\begin{tabular}{rl}
\multicolumn{2}{l}{\textbf{PRÄSENS}}\\
ich & \textbf{tröste}\\
du & \textbf{tröstest}\\
er/sie/es & \textbf{tröstet}\\
wir & \textbf{trösten}\\
ihr & \textbf{tröstet}\\
sie/Sie & \textbf{trösten}\\
\end{tabular}
\hfill
\begin{tabular}{rl}
\multicolumn{2}{l}{\textbf{PRÄTERITUM}}\\
ich & \textbf{tröstete}\\
du & \textbf{tröstetest}\\
er/sie/es & \textbf{tröstete}\\
wir & \textbf{trösteten}\\
ihr & \textbf{tröstetet}\\
sie/Sie & \textbf{trösteten}\\
\end{tabular}
\hfill
\begin{tabular}{rl}
\multicolumn{2}{l}{\textbf{IMPERATIV}}\\
& \textbf{tröste (du)!}\\
& \textbf{trösten wir!}\\
& \textbf{tröstet (ihr)!}\\
& \textbf{trösten Sie!}\\
\end{tabular}

\vspace{10pt}

\begin{tabular}{rl}
\multicolumn{2}{l}{\textbf{PERFEKT}}\\
& haben + \textbf{getröstet}\\
\end{tabular}
\hfill
\begin{tabular}{rl}
\multicolumn{2}{l}{\textbf{FUTURE I}}\\
& werden + \textbf{trösten}\\
\end{tabular}
\hfill
\begin{tabular}{rl}
\multicolumn{2}{l}{\textbf{PLUSQUAMPERFEKT}}\\
& hatten + \textbf{getröstet}\\
\end{tabular}

\vspace{10pt}

\begin{tabular}{lll}
\multicolumn{3}{l}{\textbf{MIT PRÄPOSITIONEN}}\\
& \textbf{trösten mit} + Dat & (to comfort with)\\
& \textbf{trösten über} + Akk & (to console about)\\
\end{tabular}
\hfill
\begin{tabular}{lll}
\multicolumn{3}{l}{\textbf{TRENNBARE VERBEN}}\\
\vspace{20pt}
\end{tabular}
\vspace{-5pt}
\end{verbbox}

%tun (to do)
\begin{verbbox-acc}\addcontentsline{toc}{subsubsection}{tun (to do)}
{\Large \bf tun} (to do) \hfill Type: \textbf{irregular, + Akkusativ} \hfill Auxiliary Verb: \textbf{haben}

\vspace{5pt}
\hrule
\vspace{5pt}

\begin{tabular}{rl}
\multicolumn{2}{l}{\textbf{PRÄSENS}}\\
ich & \textbf{tue}\\
du & \textbf{tust}\\
er/sie/es & \textbf{tut}\\
wir & \textbf{tun}\\
ihr & \textbf{tut}\\
sie/Sie & \textbf{tun}\\
\end{tabular}
\hfill
\begin{tabular}{rl}
\multicolumn{2}{l}{\textbf{PRÄTERITUM}}\\
ich & \textbf{tat}\\
du & \textbf{tatest}\\
er/sie/es & \textbf{tat}\\
wir & \textbf{taten}\\
ihr & \textbf{tatet}\\
sie/Sie & \textbf{taten}\\
\end{tabular}
\hfill
\begin{tabular}{rl}
\multicolumn{2}{l}{\textbf{IMPERATIV}}\\
& \textbf{tu(e) (du)!}\\
& \textbf{tun wir!}\\
& \textbf{tut (ihr)!}\\
& \textbf{tun Sie!}\\
\end{tabular}

\vspace{10pt}

\begin{tabular}{rl}
\multicolumn{2}{l}{\textbf{PERFEKT}}\\
& haben + \textbf{getan}\\
\end{tabular}
\hfill
\begin{tabular}{rl}
\multicolumn{2}{l}{\textbf{FUTURE I}}\\
& werden + \textbf{tun}\\
\end{tabular}
\hfill
\begin{tabular}{rl}
\multicolumn{2}{l}{\textbf{PLUSQUAMPERFEKT}}\\
& hatten + \textbf{getan}\\
\end{tabular}

\vspace{10pt}

\begin{tabular}{lll}
\multicolumn{3}{l}{\textbf{MIT PRÄPOSITIONEN}}\\
& \textbf{tun für} + Akk & (to do for)\\
& \textbf{tun gegen} + Akk & (to do against / about)\\
\vspace{30pt}
\end{tabular}
\hfill
\begin{tabular}{lll}
\multicolumn{3}{l}{\textbf{TRENNBARE VERBEN}}\\
& \textbf{mit$|$tun} & (to participate / be involved)\\
& \textbf{ab$|$tun} & (to dismiss / remove)\\
& \textbf{weg$|$tun} & (to put away / remove)\\
& \textbf{hinzu$|$tun} & (to add)\\
& \textbf{weh$|$tun} & (to hurt / cause pain)\\
\end{tabular}

\vspace{-5pt}
\end{verbbox-acc}

\subsection{U Verbs}

%üben (to practice)
\begin{verbbox}\addcontentsline{toc}{subsubsection}{üben (to practice)}
    {\Large \bf üben} (to practice) \hfill Type: \textbf{regular} \hfill Auxiliary Verb: \textbf{haben} 
    
     \vspace{5pt}

    \hrule
    
    \vspace{5pt}

  \begin{tabular}{rl}
     \multicolumn{2}{l}{\textbf{PRÄSENS} }\\
    ich & \textbf{übe}\\
    du & \textbf{übst}\\
    er/sie/es & \textbf{übt}\\
    wir & \textbf{üben}\\
    ihr & \textbf{übt}\\
    sie/Sie & \textbf{üben}\\
  \end{tabular}
  \hfill
     \begin{tabular}{rl}
    \multicolumn{2}{l}{\textbf{PRÄTERITUM} }\\
    ich & \textbf{übte}\\
    du & \textbf{übtest}\\
    er/sie/es & \textbf{übte}\\
    wir & \textbf{übten}\\
    ihr & \textbf{übtet}\\
    sie/Sie & \textbf{übten}\\
  \end{tabular}
  \hfill
  \begin{tabular}{rl}
     \multicolumn{2}{l}{\textbf{IMPERATIV} }\\
   & \textbf{üb(e) (du)!}\\
   & \textbf{üben wir!}\\
   & \textbf{übt (ihr)!}\\
   & \textbf{üben Sie!}\\
   & \vspace{10pt}
  \end{tabular}

    \vspace{10pt}

 \begin{tabular}{rl}
     \multicolumn{2}{l}{\textbf{PERFEKT}}\\
   & haben + \textbf{geübt}\\
  \end{tabular}
\hfill
   \begin{tabular}{rl}
    \multicolumn{2}{l}{\textbf{FUTURE I} }\\
   & werde + \textbf{üben}\\
  \end{tabular}
\hfill
   \begin{tabular}{rl}
      \multicolumn{2}{l}{\textbf{PLUSQUAMPERFEKT}}\\
   & hatten + \textbf{geübt}\\
  \end{tabular}

   \vspace{10pt}

   \begin{tabular}{lll}
      \multicolumn{3}{l}{\textbf{MIT PRÄPOSITIONEN}}\\
& \textbf{üben für} + Akk & (to practice for)\\
& \textbf{üben mit} + Dat & (to practice with)\\
\end{tabular}
  \hfill
   \begin{tabular}{lll}
      \multicolumn{3}{l}{\textbf{TRENNBARE VERBEN}}\\
 & \textbf{ein$|$üben} & (to rehearse / practice)\\
 & \textbf{aus$|$üben} & (to exercise / carry out)\\
\vspace{10pt}
  \end{tabular}

  \vspace{-5pt}

\end{verbbox}

%überlegen (B1) – to think about
\begin{verbbox-acc}\addcontentsline{toc}{subsubsection}{überlegen (to think about)}
    {\Large \bf überlegen} (to think about) \hfill Type: \textbf{irregular, + Akkusativ} \hfill Auxiliary Verb: \textbf{haben} 
    
     \vspace{5pt}

    \hrule
    
    \vspace{5pt}

  \begin{tabular}{rl}
     \multicolumn{2}{l}{\textbf{PRÄSENS} }\\
    ich & \textbf{überlege}\\
    du & \textbf{überlegst}\\
    er/sie/es & \textbf{überlegt}\\
    wir & \textbf{überlegen}\\
    ihr & \textbf{überlegt}\\
    sie/Sie & \textbf{überlegen}\\
  \end{tabular}
  \hfill
     \begin{tabular}{rl}
    \multicolumn{2}{l}{\textbf{PRÄTERITUM} }\\
    ich & \textbf{überlegte}\\
    du & \textbf{überlegtest}\\
    er/sie/es & \textbf{überlegte}\\
    wir & \textbf{überlegten}\\
    ihr & \textbf{überlegtet}\\
    sie/Sie & \textbf{überlegten}\\
  \end{tabular}
  \hfill
  \begin{tabular}{rl}
     \multicolumn{2}{l}{\textbf{IMPERATIV} }\\
   & \textbf{überleg(e) (du)!}\\
   & \textbf{überlegen wir!}\\
   & \textbf{überlegt (ihr)!}\\
   & \textbf{überlegen Sie!}\\
   & \vspace{10pt}
  \end{tabular}

    \vspace{10pt}

 \begin{tabular}{rl}
     \multicolumn{2}{l}{\textbf{PERFEKT}}\\
   & haben + \textbf{überlegt}\\
  \end{tabular}
\hfill
   \begin{tabular}{rl}
   
    \multicolumn{2}{l}{\textbf{FUTURE I} }\\
  &  werden + \textbf{überlegen}\\
 
  \end{tabular}
\hfill
   \begin{tabular}{rl}
      \multicolumn{2}{l}{\textbf{PLUSQUAMPERFEKT}}\\
  &  hatten + \textbf{überlegt}\\
  \end{tabular}
  
  \vspace{-5pt}


\end{verbbox-acc}

%übernehmen (to take over / to assume)
\begin{verbbox}\addcontentsline{toc}{subsubsection}{übernehmen (to take over / to assume)}
    {\Large \bf übernehmen} (to take over / to assume) \hfill Type: \textbf{irregular} \hfill Auxiliary Verb: \textbf{haben} 
    
     \vspace{5pt}

    \hrule
    
    \vspace{5pt}

  \begin{tabular}{rl}
     \multicolumn{2}{l}{\textbf{PRÄSENS} }\\
    ich & \textbf{übernehme}\\
    du & \textbf{übernimmst}\\
    er/sie/es & \textbf{übernimmt}\\
    wir & \textbf{übernehmen}\\
    ihr & \textbf{übernehmt}\\
    sie/Sie & \textbf{übernehmen}\\
  \end{tabular}
  \hfill
     \begin{tabular}{rl}
    \multicolumn{2}{l}{\textbf{PRÄTERITUM} }\\
    ich & \textbf{übernahm}\\
    du & \textbf{übernahmst}\\
    er/sie/es & \textbf{übernahm}\\
    wir & \textbf{übernahmen}\\
    ihr & \textbf{übernahmt}\\
    sie/Sie & \textbf{übernahmen}\\
  \end{tabular}
  \hfill
  \begin{tabular}{rl}
     \multicolumn{2}{l}{\textbf{IMPERATIV} }\\
   & \textbf{übernimm (du)!}\\
   & \textbf{übernehmen wir!}\\
   & \textbf{übernehmt (ihr)!}\\
   & \textbf{übernehmen Sie!}\\
   & \vspace{10pt}
  \end{tabular}

    \vspace{10pt}

 \begin{tabular}{rl}
     \multicolumn{2}{l}{\textbf{PERFEKT}}\\
   & haben + \textbf{übernommen}\\
  \end{tabular}
\hfill
   \begin{tabular}{rl}
    \multicolumn{2}{l}{\textbf{FUTURE I} }\\
   & werden + \textbf{übernehmen}\\
  \end{tabular}
\hfill
   \begin{tabular}{rl}
      \multicolumn{2}{l}{\textbf{PLUSQUAMPERFEKT}}\\
   & hatten + \textbf{übernommen}\\
  \end{tabular}

   \vspace{10pt}

   \begin{tabular}{lll}
      \multicolumn{3}{l}{\textbf{MIT PRÄPOSITIONEN}}\\
& \textbf{übernehmen von} + Dat & (to take over from)\\
& \textbf{übernehmen für} + Akk & (to assume responsibility for)\\
\end{tabular}
  \hfill
   \begin{tabular}{lll}
      \multicolumn{3}{l}{\textbf{TRENNBARE VERBEN}}\\
 & \textbf{---} & (none common)\\
\vspace{10pt}
  \end{tabular}

  \vspace{-5pt}

\end{verbbox}

%überweisen (to transfer money / refer)
\begin{verbbox}\addcontentsline{toc}{subsubsection}{überweisen (to transfer money / refer)}
{\Large \bf überweisen} (to transfer money / refer) \hfill Type: \textbf{regular} \hfill Auxiliary Verb: \textbf{haben}

\vspace{5pt}
\hrule
\vspace{5pt}

\begin{tabular}{rl}
\multicolumn{2}{l}{\textbf{PRÄSENS}}\\
ich & \textbf{überweise}\\
du & \textbf{überweist}\\
er/sie/es & \textbf{überweist}\\
wir & \textbf{überweisen}\\
ihr & \textbf{überweist}\\
sie/Sie & \textbf{überweisen}\\
\end{tabular}
\hfill
\begin{tabular}{rl}
\multicolumn{2}{l}{\textbf{PRÄTERITUM}}\\
ich & \textbf{überwies}\\
du & \textbf{überwiesest}\\
er/sie/es & \textbf{überwies}\\
wir & \textbf{überwiesen}\\
ihr & \textbf{überwieset}\\
sie/Sie & \textbf{überwiesen}\\
\end{tabular}
\hfill
\begin{tabular}{rl}
\multicolumn{2}{l}{\textbf{IMPERATIV}}\\
& \textbf{überweise (du)!}\\
& \textbf{überweisen wir!}\\
& \textbf{überweist (ihr)!}\\
& \textbf{überweisen Sie!}\\
\end{tabular}

\vspace{10pt}

\begin{tabular}{rl}
\multicolumn{2}{l}{\textbf{PERFEKT}}\\
& haben + \textbf{überwiesen}\\
\end{tabular}
\hfill
\begin{tabular}{rl}
\multicolumn{2}{l}{\textbf{FUTURE I}}\\
& werden + \textbf{überweisen}\\
\end{tabular}
\hfill
\begin{tabular}{rl}
\multicolumn{2}{l}{\textbf{PLUSQUAMPERFEKT}}\\
& hatten + \textbf{überwiesen}\\
\end{tabular}

\vspace{10pt}

\begin{tabular}{lll}
\multicolumn{3}{l}{\textbf{MIT PRÄPOSITIONEN}}\\
& \textbf{überweisen an} + Akk & (to transfer to / refer to)\\
\end{tabular}
\hfill
\begin{tabular}{lll}
\multicolumn{3}{l}{\textbf{TRENNBARE VERBEN}}\\
\vspace{10pt}\\
\end{tabular}
\vspace{-5pt}
\end{verbbox}

%umarmen (to hug)
\begin{verbbox-acc}\addcontentsline{toc}{subsubsection}{umarmen (to hug)}
{\Large \bf umarmen} (to hug) \hfill Type: \textbf{regular, + Akkusativ} \hfill Auxiliary Verb: \textbf{haben}

\vspace{5pt}
\hrule
\vspace{5pt}

\begin{tabular}{rl}
\multicolumn{2}{l}{\textbf{PRÄSENS}}\\
ich & \textbf{umarme}\\
du & \textbf{umarmst}\\
er/sie/es & \textbf{umarmt}\\
wir & \textbf{umarmen}\\
ihr & \textbf{umarmt}\\
sie/Sie & \textbf{umarmen}\\
\end{tabular}
\hfill
\begin{tabular}{rl}
\multicolumn{2}{l}{\textbf{PRÄTERITUM}}\\
ich & \textbf{umarmte}\\
du & \textbf{umarmtest}\\
er/sie/es & \textbf{umarmte}\\
wir & \textbf{umarmten}\\
ihr & \textbf{umarmtet}\\
sie/Sie & \textbf{umarmten}\\
\end{tabular}
\hfill
\begin{tabular}{rl}
\multicolumn{2}{l}{\textbf{IMPERATIV}}\\
& \textbf{umarme (du)!}\\
& \textbf{umarmen wir!}\\
& \textbf{umarmt (ihr)!}\\
& \textbf{umarmen Sie!}\\
\end{tabular}

\vspace{10pt}

\begin{tabular}{rl}
\multicolumn{2}{l}{\textbf{PERFEKT}}\\
& haben + \textbf{umarmt}\\
\end{tabular}
\hfill
\begin{tabular}{rl}
\multicolumn{2}{l}{\textbf{FUTURE I}}\\
& werden + \textbf{umarmen}\\
\end{tabular}
\hfill
\begin{tabular}{rl}
\multicolumn{2}{l}{\textbf{PLUSQUAMPERFEKT}}\\
& hatten + \textbf{umarmt}\\
\end{tabular}

\vspace{10pt}

\begin{tabular}{lll}
\multicolumn{3}{l}{\textbf{MIT PRÄPOSITIONEN}}\\
& ---\\
\end{tabular}
\hfill
\begin{tabular}{lll}
\multicolumn{3}{l}{\textbf{TRENNBARE VERBEN}}\\
& ---\\
\end{tabular}

\vspace{-5pt}
\end{verbbox-acc}

%unterhalten (to entertain)
\begin{verbbox-acc}\addcontentsline{toc}{subsubsection}{unterhalten (to entertain)}
{\Large \bf unterhalten} (to entertain) \hfill Type: \textbf{irregular, + Akkusativ} \hfill Auxiliary Verb: \textbf{haben}

\vspace{5pt}
\hrule
\vspace{5pt}

\begin{tabular}{rl}
\multicolumn{2}{l}{\textbf{PRÄSENS}}\\
ich & \textbf{unterhalte}\\
du & \textbf{unterhältst}\\
er/sie/es & \textbf{unterhält}\\
wir & \textbf{unterhalten}\\
ihr & \textbf{unterhaltet}\\
sie/Sie & \textbf{unterhalten}\\
\end{tabular}
\hfill
\begin{tabular}{rl}
\multicolumn{2}{l}{\textbf{PRÄTERITUM}}\\
ich & \textbf{unterhielt}\\
du & \textbf{unterhieltest}\\
er/sie/es & \textbf{unterhielt}\\
wir & \textbf{unterhielten}\\
ihr & \textbf{unterhieltet}\\
sie/Sie & \textbf{unterhielten}\\
\end{tabular}
\hfill
\begin{tabular}{rl}
\multicolumn{2}{l}{\textbf{IMPERATIV}}\\
& \textbf{unterhalte (du)!}\\
& \textbf{unterhalten wir!}\\
& \textbf{unterhaltet (ihr)!}\\
& \textbf{unterhalten Sie!}\\
\end{tabular}

\vspace{10pt}

\begin{tabular}{rl}
\multicolumn{2}{l}{\textbf{PERFEKT}}\\
& haben + \textbf{unterhalten}\\
\end{tabular}
\hfill
\begin{tabular}{rl}
\multicolumn{2}{l}{\textbf{FUTURE I}}\\
& werden + \textbf{unterhalten}\\
\end{tabular}
\hfill
\begin{tabular}{rl}
\multicolumn{2}{l}{\textbf{PLUSQUAMPERFEKT}}\\
& hatten + \textbf{unterhalten}\\
\end{tabular}

\vspace{10pt}

\begin{tabular}{lll}
\multicolumn{3}{l}{\textbf{MIT PRÄPOSITIONEN}}\\
& ---\\
\end{tabular}
\hfill
\begin{tabular}{lll}
\multicolumn{3}{l}{\textbf{TRENNBARE VERBEN}}\\
& ---\\
\end{tabular}

\vspace{-5pt}
\end{verbbox-acc}

%unternehmen (to undertake / to do)
\begin{verbbox}\addcontentsline{toc}{subsubsection}{unternehmen (to undertake / to do)}
    {\Large \bf unternehmen} (to undertake / to do) \hfill Type: \textbf{irregular} \hfill Auxiliary Verb: \textbf{haben} 
    
     \vspace{5pt}

    \hrule
    
    \vspace{5pt}

  \begin{tabular}{rl}
     \multicolumn{2}{l}{\textbf{PRÄSENS} }\\
    ich & \textbf{unternehme}\\
    du & \textbf{unternimmst}\\
    er/sie/es & \textbf{unternimmt}\\
    wir & \textbf{unternehmen}\\
    ihr & \textbf{unternehmt}\\
    sie/Sie & \textbf{unternehmen}\\
  \end{tabular}
  \hfill
     \begin{tabular}{rl}
    \multicolumn{2}{l}{\textbf{PRÄTERITUM} }\\
    ich & \textbf{unternahm}\\
    du & \textbf{unternahmst}\\
    er/sie/es & \textbf{unternahm}\\
    wir & \textbf{unternahmen}\\
    ihr & \textbf{unternahmt}\\
    sie/Sie & \textbf{unternahmen}\\
  \end{tabular}
  \hfill
  \begin{tabular}{rl}
     \multicolumn{2}{l}{\textbf{IMPERATIV} }\\
   & \textbf{unternimm (du)!}\\
   & \textbf{unternehmen wir!}\\
   & \textbf{unternehmt (ihr)!}\\
   & \textbf{unternehmen Sie!}\\
   & \vspace{10pt}
  \end{tabular}

    \vspace{10pt}

 \begin{tabular}{rl}
     \multicolumn{2}{l}{\textbf{PERFEKT}}\\
   & haben + \textbf{unternommen}\\
  \end{tabular}
\hfill
   \begin{tabular}{rl}
    \multicolumn{2}{l}{\textbf{FUTURE I} }\\
   & werden + \textbf{unternehmen}\\
  \end{tabular}
\hfill
   \begin{tabular}{rl}
      \multicolumn{2}{l}{\textbf{PLUSQUAMPERFEKT}}\\
   & hatten + \textbf{unternommen}\\
  \end{tabular}

   \vspace{10pt}

   \begin{tabular}{lll}
      \multicolumn{3}{l}{\textbf{MIT PRÄPOSITIONEN}}\\
& \textbf{unternehmen gegen} + Akk & (to undertake against)\\
& \textbf{unternehmen mit} + Dat & (to do with)\\
\end{tabular}
  \hfill
   \begin{tabular}{lll}
      \multicolumn{3}{l}{\textbf{TRENNBARE VERBEN}}\\
 & \textbf{---} & (none common)\\
\vspace{10pt}
  \end{tabular}

  \vspace{-5pt}

\end{verbbox}

%unterscheiden (to distinguish)
\begin{verbbox}\addcontentsline{toc}{subsubsection}{unterscheiden (to distinguish)}
{\Large \bf unterscheiden} (to distinguish) \hfill Type: \textbf{irregular} \hfill Auxiliary Verb: \textbf{haben} 

\vspace{5pt}
\hrule
\vspace{5pt}

\begin{tabular}{rl}
\multicolumn{2}{l}{\textbf{PRÄSENS}}\\
ich & \textbf{unterscheide}\\
du & \textbf{unterscheidest}\\
er/sie/es & \textbf{unterscheidet}\\
wir & \textbf{unterscheiden}\\
ihr & \textbf{unterscheidet}\\
sie/Sie & \textbf{unterscheiden}\\
\end{tabular}
\hfill
\begin{tabular}{rl}
\multicolumn{2}{l}{\textbf{PRÄTERITUM}}\\
ich & \textbf{unterschied}\\
du & \textbf{unterschiedst}\\
er/sie/es & \textbf{unterschied}\\
wir & \textbf{unterschieden}\\
ihr & \textbf{unterschiedet}\\
sie/Sie & \textbf{unterschieden}\\
\end{tabular}
\hfill
\begin{tabular}{rl}
\multicolumn{2}{l}{\textbf{IMPERATIV}}\\
& \textbf{unterscheide (du)!}\\
& \textbf{unterscheiden wir!}\\
& \textbf{unterscheidet (ihr)!}\\
& \textbf{unterscheiden Sie!}\\
\end{tabular}

\vspace{10pt}

\begin{tabular}{rl}
\multicolumn{2}{l}{\textbf{PERFEKT}}\\
& haben + \textbf{unterschieden}\\
\end{tabular}
\hfill
\begin{tabular}{rl}
\multicolumn{2}{l}{\textbf{FUTURE I}}\\
& werden + \textbf{unterscheiden}\\
\end{tabular}
\hfill
\begin{tabular}{rl}
\multicolumn{2}{l}{\textbf{PLUSQUAMPERFEKT}}\\
& hatten + \textbf{unterschieden}\\
\end{tabular}

\vspace{10pt}

\begin{tabular}{lll}
\multicolumn{3}{l}{\textbf{MIT PRÄPOSITIONEN}}\\
& \textbf{unterscheiden zwischen} + Dat & (to distinguish between)\\
\end{tabular}
\hfill
\begin{tabular}{lll}
\multicolumn{3}{l}{\textbf{TRENNBARE VERBEN}}\\
& ---\\
\end{tabular}

\vspace{-5pt}
\end{verbbox}

%unterschreiben - to sign something
\begin{verbbox-acc}\addcontentsline{toc}{subsubsection}{unterschreiben (to sign something)}
    {\Large \bf unterschreiben} (to sign something) \hfill Type: \textbf{irregular, + Akkusativ} \hfill Auxiliary Verb: \textbf{haben} 
    
     \vspace{5pt}

    \hrule
    
    \vspace{5pt}

  \begin{tabular}{rl}
     \multicolumn{2}{l}{\textbf{PRÄSENS} }\\
    ich & \textbf{unterschreibe}\\
    du & \textbf{unterschreibst}\\
    er/sie/es & \textbf{unterschreibt}\\
    wir & \textbf{unterschreiben}\\
    ihr & \textbf{unterschreibt}\\
    sie/Sie & \textbf{unterschreiben}\\
  \end{tabular}
  \hfill
     \begin{tabular}{rl}
    \multicolumn{2}{l}{\textbf{PRÄTERITUM} }\\
    ich & \textbf{unterschrieb}\\
    du & \textbf{unterschriebst}\\
    er/sie/es & \textbf{unterschreib}\\
    wir & \textbf{unterschrieben}\\
    ihr & \textbf{unterschriebt}\\
    sie/Sie & \textbf{unterschrieben}\\
  \end{tabular}
  \hfill
  \begin{tabular}{rl}
     \multicolumn{2}{l}{\textbf{IMPERATIV} }\\
   & \textbf{unterschreib(e) (du)!}\\
   & \textbf{unterschreiben wir!}\\
   & \textbf{unterschreibt (ihr)!}\\
   & \textbf{unterschreiben Sie!}\\
   & \vspace{10pt}
  \end{tabular}

    \vspace{10pt}

 \begin{tabular}{rl}
     \multicolumn{2}{l}{\textbf{PERFEKT}}\\
   & haben + \textbf{unterschrieben}\\
  \end{tabular}
\hfill
   \begin{tabular}{rl}
   
    \multicolumn{2}{l}{\textbf{FUTURE I} }\\
   & werden + \textbf{unterschreiben}\\
 
  \end{tabular}
\hfill
   \begin{tabular}{rl}
      \multicolumn{2}{l}{\textbf{PLUSQUAMPERFEKT}}\\
   & hatten + \textbf{unterschrieben}\\
  \end{tabular}
  
  \vspace{-5pt}


\end{verbbox-acc}

%unterstützen (to support)
\begin{verbbox}\addcontentsline{toc}{subsubsection}{unterstützen (to support)}
    {\Large \bf unterstützen} (to support) \hfill Type: \textbf{regular} \hfill Auxiliary Verb: \textbf{haben} 
    
     \vspace{5pt}

    \hrule
    
    \vspace{5pt}

  \begin{tabular}{rl}
     \multicolumn{2}{l}{\textbf{PRÄSENS} }\\
    ich & \textbf{unterstütze}\\
    du & \textbf{unterstützt}\\
    er/sie/es & \textbf{unterstützt}\\
    wir & \textbf{unterstützen}\\
    ihr & \textbf{unterstützt}\\
    sie/Sie & \textbf{unterstützen}\\
  \end{tabular}
  \hfill
     \begin{tabular}{rl}
    \multicolumn{2}{l}{\textbf{PRÄTERITUM} }\\
    ich & \textbf{unterstützte}\\
    du & \textbf{unterstütztest}\\
    er/sie/es & \textbf{unterstützte}\\
    wir & \textbf{unterstützten}\\
    ihr & \textbf{unterstütztet}\\
    sie/Sie & \textbf{unterstützten}\\
  \end{tabular}
  \hfill
  \begin{tabular}{rl}
     \multicolumn{2}{l}{\textbf{IMPERATIV} }\\
   & \textbf{unterstütz(e) (du)!}\\
   & \textbf{unterstützen wir!}\\
   & \textbf{unterstützt (ihr)!}\\
   & \textbf{unterstützen Sie!}\\
   & \vspace{10pt}
  \end{tabular}

    \vspace{10pt}

 \begin{tabular}{rl}
     \multicolumn{2}{l}{\textbf{PERFEKT}}\\
   & haben + \textbf{unterstützt}\\
  \end{tabular}
\hfill
   \begin{tabular}{rl}
   
    \multicolumn{2}{l}{\textbf{FUTURE I} }\\
   & werden + \textbf{unterstützen}\\
 
  \end{tabular}
\hfill
   \begin{tabular}{rl}
      \multicolumn{2}{l}{\textbf{PLUSQUAMPERFEKT}}\\
   & hatten + \textbf{unterstützt}\\
  \end{tabular}

   \vspace{10pt}

   \begin{tabular}{lll}
      \multicolumn{3}{l}{\textbf{MIT PRÄPOSITIONEN}}\\
& \textbf{unterstützen bei} + Dat & (to support with / help with)\\
& \textbf{unterstützen mit} + Dat & (to support with / by means of)\\
\end{tabular}
  \hfill
   \begin{tabular}{lll}
      \multicolumn{3}{l}{\textbf{TRENNBARE VERBEN}}\\
 & \textbf{---} & (none common)\\
\vspace{10pt}
  \end{tabular}


  \vspace{-5pt}


\end{verbbox}

%untersuchen (B1) – to examine
\begin{verbbox}\addcontentsline{toc}{subsubsection}{untersuchen (to examine)}
    {\Large \bf untersuchen} (to examine) \hfill Type: \textbf{irregular} \hfill Auxiliary Verb: \textbf{haben} 
    
     \vspace{5pt}

    \hrule
    
    \vspace{5pt}

  \begin{tabular}{rl}
     \multicolumn{2}{l}{\textbf{PRÄSENS} }\\
    ich & \textbf{untersuche}\\
    du & \textbf{untersuchst}\\
    er/sie/es & \textbf{untersucht}\\
    wir & \textbf{untersuchen}\\
    ihr & \textbf{untersucht}\\
    sie/Sie & \textbf{untersuchen}\\
  \end{tabular}
  \hfill
     \begin{tabular}{rl}
    \multicolumn{2}{l}{\textbf{PRÄTERITUM} }\\
    ich & \textbf{untersuchte}\\
    du & \textbf{untersuchtest}\\
    er/sie/es & \textbf{untersuchte}\\
    wir & \textbf{untersuchten}\\
    ihr & \textbf{untersuchtet}\\
    sie/Sie & \textbf{untersuchten}\\
  \end{tabular}
  \hfill
  \begin{tabular}{rl}
     \multicolumn{2}{l}{\textbf{IMPERATIV} }\\
   & \textbf{untersuch(e) (du)!}\\
   & \textbf{untersuchen wir!}\\
   & \textbf{untersucht (ihr)!}\\
   & \textbf{untersuchen Sie!}\\
   & \vspace{10pt}
  \end{tabular}

    \vspace{10pt}

 \begin{tabular}{rl}
     \multicolumn{2}{l}{\textbf{PERFEKT}}\\
   & haben + \textbf{untersucht}\\
  \end{tabular}
\hfill
   \begin{tabular}{rl}
   
    \multicolumn{2}{l}{\textbf{FUTURE I} }\\
   & werden + \textbf{untersuchen}\\
 
  \end{tabular}
\hfill
   \begin{tabular}{rl}
      \multicolumn{2}{l}{\textbf{PLUSQUAMPERFEKT}}\\
   & hatten + \textbf{untersucht}\\
  \end{tabular}
  
  \vspace{10pt}



\begin{tabular}{lll}
\multicolumn{3}{l}{\textbf{MIT PRÄPOSITIONEN}}\\
& ---\\
\end{tabular}
\hfill
\begin{tabular}{lll}
\multicolumn{3}{l}{\textbf{TRENNBARE VERBEN}}\\
& ---\\
\end{tabular}

  \vspace{-5pt}

\end{verbbox}

\subsection{V Verbs}

%verändern (B1) – to change
\begin{verbbox-acc}\addcontentsline{toc}{subsubsection}{verändern (to change)}
    {\Large \bf verändern} (to change) \hfill Type: \textbf{irregular, + Akkusativ} \hfill Auxiliary Verb: \textbf{haben} 
    
     \vspace{5pt}

    \hrule
    
    \vspace{5pt}

  \begin{tabular}{rl}
     \multicolumn{2}{l}{\textbf{PRÄSENS} }\\
    ich & \textbf{verändere}\\
    du & \textbf{veränderst}\\
    er/sie/es & \textbf{verändert}\\
    wir & \textbf{verändern}\\
    ihr & \textbf{verändert}\\
    sie/Sie & \textbf{verändern}\\
  \end{tabular}
  \hfill
     \begin{tabular}{rl}
    \multicolumn{2}{l}{\textbf{PRÄTERITUM} }\\
    ich & \textbf{veränderte}\\
    du & \textbf{verändertest}\\
    er/sie/es & \textbf{veränderte}\\
    wir & \textbf{veränderten}\\
    ihr & \textbf{verändertet}\\
    sie/Sie & \textbf{veränderten}\\
  \end{tabular}
  \hfill
  \begin{tabular}{rl}
     \multicolumn{2}{l}{\textbf{IMPERATIV} }\\
   & \textbf{veränder(e) (du)!}\\
   & \textbf{verändern wir!}\\
   & \textbf{verändert (ihr)!}\\
   & \textbf{verändern Sie!}\\
   & \vspace{10pt}
  \end{tabular}

    \vspace{10pt}

 \begin{tabular}{rl}
     \multicolumn{2}{l}{\textbf{PERFEKT}}\\
  &  haben + \textbf{verändert}\\
  \end{tabular}
\hfill
   \begin{tabular}{rl}
   
    \multicolumn{2}{l}{\textbf{FUTURE I} }\\
  &  werden + \textbf{veränderen}\\
 
  \end{tabular}
\hfill
   \begin{tabular}{rl}
      \multicolumn{2}{l}{\textbf{PLUSQUAMPERFEKT}}\\
   & hatten + \textbf{verändert}\\
  \end{tabular}
  
  \vspace{-5pt}


\end{verbbox-acc}

%verbinden (to connect / link)
\begin{verbbox}\addcontentsline{toc}{subsubsection}{verbinden (to connect / link)}
{\Large \bf verbinden} (to connect / link) \hfill Type: \textbf{regular} \hfill Auxiliary Verb: \textbf{haben}

\vspace{5pt}
\hrule
\vspace{5pt}

\begin{tabular}{rl}
\multicolumn{2}{l}{\textbf{PRÄSENS}}\\
ich & \textbf{verbinde}\\
du & \textbf{verbindest}\\
er/sie/es & \textbf{verbindet}\\
wir & \textbf{verbinden}\\
ihr & \textbf{verbindet}\\
sie/Sie & \textbf{verbinden}\\
\end{tabular}
\hfill
\begin{tabular}{rl}
\multicolumn{2}{l}{\textbf{PRÄTERITUM}}\\
ich & \textbf{verband}\\
du & \textbf{verbandst}\\
er/sie/es & \textbf{verband}\\
wir & \textbf{verbanden}\\
ihr & \textbf{verbandet}\\
sie/Sie & \textbf{verbanden}\\
\end{tabular}
\hfill
\begin{tabular}{rl}
\multicolumn{2}{l}{\textbf{IMPERATIV}}\\
& \textbf{verbinde (du)!}\\
& \textbf{verbinden wir!}\\
& \textbf{verbindet (ihr)!}\\
& \textbf{verbinden Sie!}\\
\end{tabular}

\vspace{10pt}

\begin{tabular}{rl}
\multicolumn{2}{l}{\textbf{PERFEKT}}\\
& haben + \textbf{verbunden}\\
\end{tabular}
\hfill
\begin{tabular}{rl}
\multicolumn{2}{l}{\textbf{FUTURE I}}\\
& werden + \textbf{verbinden}\\
\end{tabular}
\hfill
\begin{tabular}{rl}
\multicolumn{2}{l}{\textbf{PLUSQUAMPERFEKT}}\\
& hatten + \textbf{verbunden}\\
\end{tabular}

\vspace{10pt}

\begin{tabular}{lll}
\multicolumn{3}{l}{\textbf{MIT PRÄPOSITIONEN}}\\
& \textbf{verbinden mit} + Dat & (to connect with)\\
& \textbf{verbinden zu} + Dat & (to link to / form into)\\
\end{tabular}
\hfill
\begin{tabular}{lll}
\multicolumn{3}{l}{\textbf{TRENNBARE VERBEN}}\\
\vspace{20pt}
\end{tabular}
\vspace{-5pt}
\end{verbbox}

%verbreiten (to spread / distribute)
\begin{verbbox}\addcontentsline{toc}{subsubsection}{verbreiten (to spread / distribute)}
{\Large \bf verbreiten} (to spread / distribute) \hfill Type: \textbf{regular} \hfill Auxiliary Verb: \textbf{haben}

\vspace{5pt}
\hrule
\vspace{5pt}

\begin{tabular}{rl}
\multicolumn{2}{l}{\textbf{PRÄSENS}}\\
ich & \textbf{verbreite}\\
du & \textbf{verbreitest}\\
er/sie/es & \textbf{verbreitet}\\
wir & \textbf{verbreiten}\\
ihr & \textbf{verbreitet}\\
sie/Sie & \textbf{verbreiten}\\
\end{tabular}
\hfill
\begin{tabular}{rl}
\multicolumn{2}{l}{\textbf{PRÄTERITUM}}\\
ich & \textbf{verbreitete}\\
du & \textbf{verbreitetest}\\
er/sie/es & \textbf{verbreitete}\\
wir & \textbf{verbreiteten}\\
ihr & \textbf{verbreitetet}\\
sie/Sie & \textbf{verbreiteten}\\
\end{tabular}
\hfill
\begin{tabular}{rl}
\multicolumn{2}{l}{\textbf{IMPERATIV}}\\
& \textbf{verbreite (du)!}\\
& \textbf{verbreiten wir!}\\
& \textbf{verbreitet (ihr)!}\\
& \textbf{verbreiten Sie!}\\
\end{tabular}

\vspace{10pt}

\begin{tabular}{rl}
\multicolumn{2}{l}{\textbf{PERFEKT}}\\
& haben + \textbf{verbreitet}\\
\end{tabular}
\hfill
\begin{tabular}{rl}
\multicolumn{2}{l}{\textbf{FUTURE I}}\\
& werden + \textbf{verbreiten}\\
\end{tabular}
\hfill
\begin{tabular}{rl}
\multicolumn{2}{l}{\textbf{PLUSQUAMPERFEKT}}\\
& hatten + \textbf{verbreitet}\\
\end{tabular}

\vspace{10pt}

\begin{tabular}{lll}
\multicolumn{3}{l}{\textbf{MIT PRÄPOSITIONEN}}\\
& \textbf{verbreiten unter} + Dat & (to spread among)\\
& \textbf{verbreiten über} + Akk & (to spread about)\\
& \textbf{verbreiten durch} + Akk & (to spread through)\\
\end{tabular}
\hfill
\begin{tabular}{lll}
\multicolumn{3}{l}{\textbf{TRENNBARE VERBEN}}\\
\vspace{30pt}\end{tabular}
\vspace{-5pt}
\end{verbbox}

%verbringen (to spend time)
\begin{verbbox-acc}\addcontentsline{toc}{subsubsection}{verbringen (to spend time)}
{\Large \bf verbringen} (to spend time) \hfill Type: \textbf{irregular, + Akkusativ} \hfill Auxiliary Verb: \textbf{haben}

\vspace{5pt}
\hrule
\vspace{5pt}

\begin{tabular}{rl}
\multicolumn{2}{l}{\textbf{PRÄSENS}}\\
ich & \textbf{verbringe}\\
du & \textbf{verbringst}\\
er/sie/es & \textbf{verbringt}\\
wir & \textbf{verbringen}\\
ihr & \textbf{verbringt}\\
sie/Sie & \textbf{verbringen}\\
\end{tabular}
\hfill
\begin{tabular}{rl}
\multicolumn{2}{l}{\textbf{PRÄTERITUM}}\\
ich & \textbf{verbrachte}\\
du & \textbf{verbrachtest}\\
er/sie/es & \textbf{verbrachte}\\
wir & \textbf{verbrachten}\\
ihr & \textbf{verbrachtet}\\
sie/Sie & \textbf{verbrachten}\\
\end{tabular}
\hfill
\begin{tabular}{rl}
\multicolumn{2}{l}{\textbf{IMPERATIV}}\\
& \textbf{verbring(e) (du)!}\\
& \textbf{verbringen wir!}\\
& \textbf{verbringt (ihr)!}\\
& \textbf{verbringen Sie!}\\
\end{tabular}

\vspace{10pt}

\begin{tabular}{rl}
\multicolumn{2}{l}{\textbf{PERFEKT}}\\
& haben + \textbf{verbracht}\\
\end{tabular}
\hfill
\begin{tabular}{rl}
\multicolumn{2}{l}{\textbf{FUTURE I}}\\
& werden + \textbf{verbringen}\\
\end{tabular}
\hfill
\begin{tabular}{rl}
\multicolumn{2}{l}{\textbf{PLUSQUAMPERFEKT}}\\
& hatten + \textbf{verbracht}\\
\end{tabular}

\vspace{10pt}

\begin{tabular}{lll}
\multicolumn{3}{l}{\textbf{MIT PRÄPOSITIONEN}}\\
& ---\\
\end{tabular}
\hfill
\begin{tabular}{lll}
\multicolumn{3}{l}{\textbf{TRENNBARE VERBEN}}\\
& ---\\
\end{tabular}

\vspace{-5pt}
\end{verbbox-acc}

%vereinbaren (to arrange / agree on)
\begin{verbbox}\addcontentsline{toc}{subsubsection}{vereinbaren (to arrange / agree on)}
{\Large \bf vereinbaren} (to arrange / agree on) \hfill Type: \textbf{regular} \hfill Auxiliary Verb: \textbf{haben}

\vspace{5pt}
\hrule
\vspace{5pt}

\begin{tabular}{rl}
\multicolumn{2}{l}{\textbf{PRÄSENS}}\\
ich & \textbf{vereinbare}\\
du & \textbf{vereinbarst}\\
er/sie/es & \textbf{vereinbart}\\
wir & \textbf{vereinbaren}\\
ihr & \textbf{vereinbart}\\
sie/Sie & \textbf{vereinbaren}\\
\end{tabular}
\hfill
\begin{tabular}{rl}
\multicolumn{2}{l}{\textbf{PRÄTERITUM}}\\
ich & \textbf{vereinbarte}\\
du & \textbf{vereinbartest}\\
er/sie/es & \textbf{vereinbarte}\\
wir & \textbf{vereinbarten}\\
ihr & \textbf{vereinbartet}\\
sie/Sie & \textbf{vereinbarten}\\
\end{tabular}
\hfill
\begin{tabular}{rl}
\multicolumn{2}{l}{\textbf{IMPERATIV}}\\
& \textbf{vereinbare (du)!}\\
& \textbf{vereinbaren wir!}\\
& \textbf{vereinbart (ihr)!}\\
& \textbf{vereinbaren Sie!}\\
\end{tabular}

\vspace{10pt}

\begin{tabular}{rl}
\multicolumn{2}{l}{\textbf{PERFEKT}}\\
& haben + \textbf{vereinbart}\\
\end{tabular}
\hfill
\begin{tabular}{rl}
\multicolumn{2}{l}{\textbf{FUTURE I}}\\
& werden + \textbf{vereinbaren}\\
\end{tabular}
\hfill
\begin{tabular}{rl}
\multicolumn{2}{l}{\textbf{PLUSQUAMPERFEKT}}\\
& hatten + \textbf{vereinbart}\\
\end{tabular}

\vspace{10pt}

\begin{tabular}{lll}
\multicolumn{3}{l}{\textbf{MIT PRÄPOSITIONEN}}\\
& \textbf{vereinbaren mit} + Dat & (to arrange with / agree with someone)\\
& \textbf{vereinbaren über} + Akk & (to agree on / about something)\\
\end{tabular}
\hfill
\begin{tabular}{lll}
\multicolumn{3}{l}{\textbf{TRENNBARE VERBEN}}\\
\vspace{20pt}\end{tabular}
\vspace{-5pt}
\end{verbbox}



%vergleichen (A2) – to compare
\begin{verbbox-dat}\addcontentsline{toc}{subsubsection}{vergleichen (to compare)}
    {\Large \bf vergleichen} (to compare) \hfill Type: \textbf{irregular, + Accuative + Dativ} \hfill Auxiliary Verb: \textbf{haben} 
    
     \vspace{5pt}

    \hrule
    
    \vspace{5pt}

  \begin{tabular}{rl}
     \multicolumn{2}{l}{\textbf{PRÄSENS} }\\
    ich & \textbf{vergleiche}\\
    du & \textbf{vergleichst}\\
    er/sie/es & \textbf{vergleicht}\\
    wir & \textbf{vergleichen}\\
    ihr & \textbf{vergleicht}\\
    sie/Sie & \textbf{vergleichen}\\
  \end{tabular}
  \hfill
     \begin{tabular}{rl}
    \multicolumn{2}{l}{\textbf{PRÄTERITUM} }\\
    ich & \textbf{verglich}\\
    du & \textbf{verglichst}\\
    er/sie/es & \textbf{verglich}\\
    wir & \textbf{verglichen}\\
    ihr & \textbf{verglicht}\\
    sie/Sie & \textbf{verglichen}\\
  \end{tabular}
  \hfill
  \begin{tabular}{rl}
     \multicolumn{2}{l}{\textbf{IMPERATIV} }\\
   & \textbf{vergleich(e) (du)!}\\
   & \textbf{vergleichen wir!}\\
   & \textbf{vergleicht (ihr)!}\\
   & \textbf{vergleichen Sie!}\\
   & \vspace{10pt}
  \end{tabular}

    \vspace{10pt}

 \begin{tabular}{rl}
     \multicolumn{2}{l}{\textbf{PERFEKT}}\\
  &  haben + \textbf{verglichen}\\
  \end{tabular}
\hfill
   \begin{tabular}{rl}
   
    \multicolumn{2}{l}{\textbf{FUTURE I} }\\
  &  werden + \textbf{vergleichen}\\
 
  \end{tabular}
\hfill
   \begin{tabular}{rl}
      \multicolumn{2}{l}{\textbf{PLUSQUAMPERFEKT}}\\
  &  hatten + \textbf{verglichen}\\
  \end{tabular}
  
  \vspace{-5pt}


\end{verbbox-dat}

%vergessen (A1) – to forget
\begin{verbbox-acc}\addcontentsline{toc}{subsubsection}{vergessen (to forget)}
    {\Large \bf vergessen} (to forget) \hfill Type: \textbf{irregular, + Akkusativ} \hfill Auxiliary Verb: \textbf{haben} 
    
     \vspace{5pt}

    \hrule
    
    \vspace{5pt}

  \begin{tabular}{rl}
     \multicolumn{2}{l}{\textbf{PRÄSENS} }\\
    ich & \textbf{vergesse}\\
    du & \textbf{vergisst}\\
    er/sie/es & \textbf{vergisst}\\
    wir & \textbf{vergessen}\\
    ihr & \textbf{vergesst}\\
    sie/Sie & \textbf{vergessen}\\
  \end{tabular}
  \hfill
     \begin{tabular}{rl}
    \multicolumn{2}{l}{\textbf{PRÄTERITUM} }\\
    ich & \textbf{vergaß}\\
    du & \textbf{vergaßest}\\
    er/sie/es & \textbf{vergaß}\\
    wir & \textbf{vergaßen}\\
    ihr & \textbf{vergaßt}\\
    sie/Sie & \textbf{vergaßen}\\
  \end{tabular}
  \hfill
  \begin{tabular}{rl}
     \multicolumn{2}{l}{\textbf{IMPERATIV} }\\
   & \textbf{gehör(e) (du)!}\\
   & \textbf{gehören wir!}\\
   & \textbf{gehört (ihr)!}\\
   & \textbf{gehören Sie!}\\
   & \vspace{10pt}
  \end{tabular}

    \vspace{10pt}

 \begin{tabular}{rl}
     \multicolumn{2}{l}{\textbf{PERFEKT}}\\
  &  haben + \textbf{vergessen}\\
  \end{tabular}
\hfill
   \begin{tabular}{rl}
   
    \multicolumn{2}{l}{\textbf{FUTURE I} }\\
  &  werden + \textbf{vergessen}\\
 
  \end{tabular}
\hfill
   \begin{tabular}{rl}
      \multicolumn{2}{l}{\textbf{PLUSQUAMPERFEKT}}\\
  &  hatten + \textbf{vergessen}\\
  \end{tabular}
  
  \vspace{-5pt}


\end{verbbox-acc}

%verhindern (to prevent)
\begin{verbbox}\addcontentsline{toc}{subsubsection}{verhindern (to prevent)}
    {\Large \bf verhindern} (to prevent) \hfill Type: \textbf{regular, inseparable} \hfill Auxiliary Verb: \textbf{haben} 
    
     \vspace{5pt}

    \hrule
    
    \vspace{5pt}

  \begin{tabular}{rl}
     \multicolumn{2}{l}{\textbf{PRÄSENS} }\\
    ich & \textbf{verhindere}\\
    du & \textbf{verhinderst}\\
    er/sie/es & \textbf{verhindert}\\
    wir & \textbf{verhindern}\\
    ihr & \textbf{verhindert}\\
    sie/Sie & \textbf{verhindern}\\
  \end{tabular}
  \hfill
     \begin{tabular}{rl}
    \multicolumn{2}{l}{\textbf{PRÄTERITUM} }\\
    ich & \textbf{verhinderte}\\
    du & \textbf{verhindertest}\\
    er/sie/es & \textbf{verhinderte}\\
    wir & \textbf{verhinderten}\\
    ihr & \textbf{verhindertet}\\
    sie/Sie & \textbf{verhinderten}\\
  \end{tabular}
  \hfill
  \begin{tabular}{rl}
     \multicolumn{2}{l}{\textbf{IMPERATIV} }\\
   & \textbf{verhindere (du)!}\\
   & \textbf{verhindern wir!}\\
   & \textbf{verhindert (ihr)!}\\
   & \textbf{verhindern Sie!}\\
   & \vspace{10pt}
  \end{tabular}

    \vspace{10pt}

 \begin{tabular}{rl}
     \multicolumn{2}{l}{\textbf{PERFEKT}}\\
   & haben + \textbf{verhindert}\\
  \end{tabular}
\hfill
   \begin{tabular}{rl}
    \multicolumn{2}{l}{\textbf{FUTURE I} }\\
   & werden + \textbf{verhindern}\\
  \end{tabular}
\hfill
   \begin{tabular}{rl}
      \multicolumn{2}{l}{\textbf{PLUSQUAMPERFEKT}}\\
   & hatten + \textbf{verhindert}\\
  \end{tabular}

   \vspace{10pt}

   \begin{tabular}{lll}
      \multicolumn{3}{l}{\textbf{MIT PRÄPOSITIONEN}}\\
& \textbf{verhindern an} + Dat & (to prevent from)\\
& \textbf{verhindern durch} + Akk & (to prevent by means of)\\
\end{tabular}
  \hfill
   \begin{tabular}{lll}
      \multicolumn{3}{l}{\textbf{TRENNBARE VERBEN}}\\
 & \textbf{---} & (none common)\\
\vspace{10pt}
  \end{tabular}

  \vspace{-5pt}

\end{verbbox}



%verleihen (to lend / award)
\begin{verbbox}\addcontentsline{toc}{subsubsection}{verleihen (to lend / award)}
    {\Large \bf verleihen} (to lend / award) \hfill Type: \textbf{regular} \hfill Auxiliary Verb: \textbf{haben} 
    
     \vspace{5pt}

    \hrule
    
    \vspace{5pt}

  \begin{tabular}{rl}
     \multicolumn{2}{l}{\textbf{PRÄSENS} }\\
    ich & \textbf{verleihe}\\
    du & \textbf{verleihst}\\
    er/sie/es & \textbf{verleiht}\\
    wir & \textbf{verleihen}\\
    ihr & \textbf{verleiht}\\
    sie/Sie & \textbf{verleihen}\\
  \end{tabular}
  \hfill
  \begin{tabular}{rl}
    \multicolumn{2}{l}{\textbf{PRÄTERITUM} }\\
    ich & \textbf{verlieh}\\
    du & \textbf{verliehst}\\
    er/sie/es & \textbf{verlieh}\\
    wir & \textbf{verliehen}\\
    ihr & \textbf{verlieht}\\
    sie/Sie & \textbf{verliehen}\\
  \end{tabular}
  \hfill
  \begin{tabular}{rl}
     \multicolumn{2}{l}{\textbf{IMPERATIV} }\\
   & \textbf{verleihe (du)!}\\
   & \textbf{verleihen wir!}\\
   & \textbf{verleiht (ihr)!}\\
   & \textbf{verleihen Sie!}\\
   & \vspace{10pt}
  \end{tabular}

 \vspace{10pt}

 \begin{tabular}{rl}
     \multicolumn{2}{l}{\textbf{PERFEKT}}\\
   & haben + \textbf{verliehen}\\
  \end{tabular}
\hfill
\begin{tabular}{rl}
    \multicolumn{2}{l}{\textbf{FUTURE I}}\\
   & werden + \textbf{verleihen}\\
  \end{tabular}
\hfill
\begin{tabular}{rl}
    \multicolumn{2}{l}{\textbf{PLUSQUAMPERFEKT}}\\
   & hatten + \textbf{verliehen}\\
  \end{tabular}

 \vspace{10pt}

\begin{tabular}{lll}
      \multicolumn{3}{l}{\textbf{MIT PRÄPOSITIONEN}}\\
& \textbf{verleihen an} + Akk & (to lend to)\\
& \textbf{verleihen für} + Akk & (to award for)\\
\end{tabular}
  \hfill
\begin{tabular}{lll}
      \multicolumn{3}{l}{\textbf{TRENNBARE VERBEN}}\\
 & \textbf{ver$|$leihen} & (to lend out / award)\\
\vspace{10pt}
  \end{tabular}


\vspace{-5pt}
\end{verbbox}

%verletzen (to injure / to hurt)
\begin{verbbox}\addcontentsline{toc}{subsubsection}{verletzen (to injure / to hurt)}
    {\Large \bf verletzen} (to injure / to hurt) \hfill Type: \textbf{regular} \hfill Auxiliary Verb: \textbf{haben} 
    
     \vspace{5pt}

    \hrule
    
    \vspace{5pt}

  \begin{tabular}{rl}
     \multicolumn{2}{l}{\textbf{PRÄSENS} }\\
    ich & \textbf{verletze}\\
    du & \textbf{verletzt}\\
    er/sie/es & \textbf{verletzt}\\
    wir & \textbf{verletzen}\\
    ihr & \textbf{verletzt}\\
    sie/Sie & \textbf{verletzen}\\
  \end{tabular}
  \hfill
     \begin{tabular}{rl}
    \multicolumn{2}{l}{\textbf{PRÄTERITUM} }\\
    ich & \textbf{verletzte}\\
    du & \textbf{verletztest}\\
    er/sie/es & \textbf{verletzte}\\
    wir & \textbf{verletzten}\\
    ihr & \textbf{verletztet}\\
    sie/Sie & \textbf{verletzten}\\
  \end{tabular}
  \hfill
  \begin{tabular}{rl}
     \multicolumn{2}{l}{\textbf{IMPERATIV} }\\
   & \textbf{verletz(e) (du)!}\\
   & \textbf{verletzen wir!}\\
   & \textbf{verletzt (ihr)!}\\
   & \textbf{verletzen Sie!}\\
   & \vspace{10pt}
  \end{tabular}

    \vspace{10pt}

 \begin{tabular}{rl}
     \multicolumn{2}{l}{\textbf{PERFEKT}}\\
   & haben + \textbf{verletzt}\\
  \end{tabular}
\hfill
   \begin{tabular}{rl}
    \multicolumn{2}{l}{\textbf{FUTURE I} }\\
   & werde + \textbf{verletzen}\\
  \end{tabular}
\hfill
   \begin{tabular}{rl}
      \multicolumn{2}{l}{\textbf{PLUSQUAMPERFEKT}}\\
   & hatten + \textbf{verletzt}\\
  \end{tabular}

   \vspace{10pt}

   \begin{tabular}{lll}
      \multicolumn{3}{l}{\textbf{MIT PRÄPOSITIONEN}}\\
& \textbf{sich verletzen an} + Dat & (to injure oneself on)\\
& \textbf{verletzen mit} + Dat & (to injure with)\\
\end{tabular}
  \hfill
   \begin{tabular}{lll}
      \multicolumn{3}{l}{\textbf{TRENNBARE VERBEN}}\\
 & \textbf{---} & (none common)\\
\vspace{10pt}
  \end{tabular}

  \vspace{-5pt}

\end{verbbox}

%verlieren (to lose)
\begin{verbbox}\addcontentsline{toc}{subsubsection}{verlieren (to lose)}
    {\Large \bf verlieren} (to lose) \hfill Type: \textbf{irregular} \hfill Auxiliary Verb: \textbf{haben} 
    
     \vspace{5pt}

    \hrule
    
    \vspace{5pt}

  \begin{tabular}{rl}
     \multicolumn{2}{l}{\textbf{PRÄSENS} }\\
    ich & \textbf{verliere}\\
    du & \textbf{verlierst}\\
    er/sie/es & \textbf{verliert}\\
    wir & \textbf{verlieren}\\
    ihr & \textbf{verliert}\\
    sie/Sie & \textbf{verlieren}\\
  \end{tabular}
  \hfill
     \begin{tabular}{rl}
    \multicolumn{2}{l}{\textbf{PRÄTERITUM} }\\
    ich & \textbf{verlor}\\
    du & \textbf{verlorst}\\
    er/sie/es & \textbf{verlor}\\
    wir & \textbf{verloren}\\
    ihr & \textbf{verlort}\\
    sie/Sie & \textbf{verloren}\\
  \end{tabular}
  \hfill
  \begin{tabular}{rl}
     \multicolumn{2}{l}{\textbf{IMPERATIV} }\\
   & \textbf{verlier(e) (du)!}\\
   & \textbf{verlieren wir!}\\
   & \textbf{verliert (ihr)!}\\
   & \textbf{verlieren Sie!}\\
   & \vspace{10pt}
  \end{tabular}

    \vspace{10pt}

 \begin{tabular}{rl}
     \multicolumn{2}{l}{\textbf{PERFEKT}}\\
   & haben + \textbf{verloren}\\
  \end{tabular}
\hfill
   \begin{tabular}{rl}
    \multicolumn{2}{l}{\textbf{FUTURE I} }\\
   & werde + \textbf{verlieren}\\
  \end{tabular}
\hfill
   \begin{tabular}{rl}
      \multicolumn{2}{l}{\textbf{PLUSQUAMPERFEKT}}\\
   & hatten + \textbf{verloren}\\
  \end{tabular}

   \vspace{10pt}

   \begin{tabular}{lll}
      \multicolumn{3}{l}{\textbf{MIT PRÄPOSITIONEN}}\\
& \textbf{verlieren an} + Akk & (to lose to)\\
& \textbf{verlieren durch} + Akk & (to lose through / because of)\\
\end{tabular}
  \hfill
   \begin{tabular}{lll}
      \multicolumn{3}{l}{\textbf{TRENNBARE VERBEN}}\\
 & \textbf{---} & (none common)\\
\vspace{10pt}
  \end{tabular}

  \vspace{-5pt}

\end{verbbox}

%vermeiden - to avoid
\begin{verbbox}\addcontentsline{toc}{subsubsection}{vermeiden (to avoid)}
{\Large \bf vermeiden} (to avoid) \hfill Type: \textbf{irregular} \hfill Auxiliary Verb: \textbf{haben}

\vspace{5pt}
\hrule
\vspace{5pt}

\begin{tabular}{rl}
\multicolumn{2}{l}{\textbf{PRÄSENS}}\\
ich & \textbf{vermeide}\\
du & \textbf{vermeidest}\\
er/sie/es & \textbf{vermeidet}\\
wir & \textbf{vermeiden}\\
ihr & \textbf{vermeidet}\\
sie/Sie & \textbf{vermeiden}\\
\end{tabular}
\hfill
\begin{tabular}{rl}
\multicolumn{2}{l}{\textbf{PRÄTERITUM}}\\
ich & \textbf{vermied}\\
du & \textbf{vermiedest}\\
er/sie/es & \textbf{vermied}\\
wir & \textbf{vermieden}\\
ihr & \textbf{vermiedet}\\
sie/Sie & \textbf{vermieden}\\
\end{tabular}
\hfill
\begin{tabular}{rl}
\multicolumn{2}{l}{\textbf{IMPERATIV}}\\
& \textbf{vermeide (du)!}\\
& \textbf{vermeiden wir!}\\
& \textbf{vermeidet (ihr)!}\\
& \textbf{vermeiden Sie!}\\
\end{tabular}

\vspace{10pt}

\begin{tabular}{rl}
\multicolumn{2}{l}{\textbf{PERFEKT}}\\
& haben + \textbf{vermieden}\\
\end{tabular}
\hfill
\begin{tabular}{rl}
\multicolumn{2}{l}{\textbf{FUTURE I}}\\
& werden + \textbf{vermeiden}\\
\end{tabular}
\hfill
\begin{tabular}{rl}
\multicolumn{2}{l}{\textbf{PLUSQUAMPERFEKT}}\\
& hatten + \textbf{vermieden}\\
\end{tabular}

\vspace{10pt}

\begin{tabular}{lll}
\multicolumn{3}{l}{\textbf{MIT PRÄPOSITIONEN}}\\
& --- & \\
\end{tabular}
\hfill
\begin{tabular}{lll}
\multicolumn{3}{l}{\textbf{TRENNBARE VERBEN}}\\
& --- &\\
\end{tabular}
\vspace{-5pt}
\end{verbbox}

%vermissen (to miss someone)
\begin{verbbox}\addcontentsline{toc}{subsubsection}{vermissen (to miss someone)}
{\Large \bf vermissen} (to miss someone) \hfill Type: \textbf{regular} \hfill Auxiliary Verb: \textbf{haben}

\vspace{5pt}
\hrule
\vspace{5pt}

\begin{tabular}{rl}
\multicolumn{2}{l}{\textbf{PRÄSENS}}\\
ich & \textbf{vermisse}\\
du & \textbf{vermisst}\\
er/sie/es & \textbf{vermisst}\\
wir & \textbf{vermissen}\\
ihr & \textbf{vermisst}\\
sie/Sie & \textbf{vermissen}\\
\end{tabular}
\hfill
\begin{tabular}{rl}
\multicolumn{2}{l}{\textbf{PRÄTERITUM}}\\
ich & \textbf{vermisste}\\
du & \textbf{vermisstest}\\
er/sie/es & \textbf{vermisste}\\
wir & \textbf{vermissten}\\
ihr & \textbf{vermisstet}\\
sie/Sie & \textbf{vermissten}\\
\end{tabular}
\hfill
\begin{tabular}{rl}
\multicolumn{2}{l}{\textbf{IMPERATIV}}\\
& \textbf{vermisse (du)!}\\
& \textbf{vermissen wir!}\\
& \textbf{vermisst (ihr)!}\\
& \textbf{vermissen Sie!}\\
\end{tabular}

\vspace{10pt}

\begin{tabular}{rl}
\multicolumn{2}{l}{\textbf{PERFEKT}}\\
& haben + \textbf{vermisst}\\
\end{tabular}
\hfill
\begin{tabular}{rl}
\multicolumn{2}{l}{\textbf{FUTURE I}}\\
& werden + \textbf{vermissen}\\
\end{tabular}
\hfill
\begin{tabular}{rl}
\multicolumn{2}{l}{\textbf{PLUSQUAMPERFEKT}}\\
& hatten + \textbf{vermisst}\\
\end{tabular}

\vspace{10pt}

\begin{tabular}{lll}
\multicolumn{3}{l}{\textbf{MIT PRÄPOSITIONEN}}\\
& \textbf{vermissen} + Akk & (to miss someone/something)\\
\end{tabular}
\hfill
\begin{tabular}{lll}
\multicolumn{3}{l}{\textbf{TRENNBARE VERBEN}}\\
& --- & \\
\end{tabular}
\vspace{-5pt}
\end{verbbox}

%verputzen
\begin{verbbox}\addcontentsline{toc}{subsubsection}{verputzen (to plaster / render)}
    {\Large \bf verputzen} (to plaster / render) \hfill Type: \textbf{regular} \hfill Auxiliary Verb: \textbf{haben} 
    
     \vspace{5pt}

    \hrule
    
    \vspace{5pt}

  \begin{tabular}{rl}
     \multicolumn{2}{l}{\textbf{PRÄSENS} }\\
    ich & \textbf{verputze}\\
    du & \textbf{verputzt}\\
    er/sie/es & \textbf{verputzt}\\
    wir & \textbf{verputzen}\\
    ihr & \textbf{verputzt}\\
    sie/Sie & \textbf{verputzen}\\
  \end{tabular}
  \hfill
  \begin{tabular}{rl}
    \multicolumn{2}{l}{\textbf{PRÄTERITUM} }\\
    ich & \textbf{verputzte}\\
    du & \textbf{verputztest}\\
    er/sie/es & \textbf{verputzte}\\
    wir & \textbf{verputzten}\\
    ihr & \textbf{verputztet}\\
    sie/Sie & \textbf{verputzten}\\
  \end{tabular}
  \hfill
  \begin{tabular}{rl}
     \multicolumn{2}{l}{\textbf{IMPERATIV} }\\
   & \textbf{verputze (du)!}\\
   & \textbf{verputzen wir!}\\
   & \textbf{verputzt (ihr)!}\\
   & \textbf{verputzen Sie!}\\
   & \vspace{10pt}
  \end{tabular}

 \vspace{10pt}

 \begin{tabular}{rl}
     \multicolumn{2}{l}{\textbf{PERFEKT}}\\
   & haben + \textbf{verputzt}\\
  \end{tabular}
\hfill
\begin{tabular}{rl}
    \multicolumn{2}{l}{\textbf{FUTURE I}}\\
   & werden + \textbf{verputzen}\\
  \end{tabular}
\hfill
\begin{tabular}{rl}
    \multicolumn{2}{l}{\textbf{PLUSQUAMPERFEKT}}\\
   & hatten + \textbf{verputzt}\\
  \end{tabular}

 \vspace{10pt}

\begin{tabular}{lll}
      \multicolumn{3}{l}{\textbf{MIT PRÄPOSITIONEN}}\\
& \textbf{verputzen mit} + Dat & (to plaster with)\\
\end{tabular}
  \hfill
\begin{tabular}{lll}
      \multicolumn{3}{l}{\textbf{TRENNBARE VERBEN}}\\
 & \textbf{ein$|$verputzen} & (to plaster in / apply)\\
  \end{tabular}


\vspace{-5pt}
\end{verbbox}

%verschreiben (to prescribe)
\begin{verbbox}\addcontentsline{toc}{subsubsection}{verschreiben (to prescribe)}
{\Large \bf verschreiben} (to prescribe) \hfill Type: \textbf{irregular} \hfill Auxiliary Verb: \textbf{haben / sich + haben}

\vspace{5pt}
\hrule
\vspace{5pt}

\begin{tabular}{rl}
\multicolumn{2}{l}{\textbf{PRÄSENS}}\\
ich & \textbf{verschreibe}\\
du & \textbf{verschreibst}\\
er/sie/es & \textbf{verschreibt}\\
wir & \textbf{verschreiben}\\
ihr & \textbf{verschreibt}\\
sie/Sie & \textbf{verschreiben}\\
\end{tabular}
\hfill
\begin{tabular}{rl}
\multicolumn{2}{l}{\textbf{PRÄTERITUM}}\\
ich & \textbf{verschrieb}\\
du & \textbf{verschriebst}\\
er/sie/es & \textbf{verschrieb}\\
wir & \textbf{verschrieben}\\
ihr & \textbf{verschriebt}\\
sie/Sie & \textbf{verschrieben}\\
\end{tabular}
\hfill
\begin{tabular}{rl}
\multicolumn{2}{l}{\textbf{IMPERATIV}}\\
& \textbf{verschreibe (du)!}\\
& \textbf{verschreiben wir!}\\
& \textbf{verschreibt (ihr)!}\\
& \textbf{verschreiben Sie!}\\
\end{tabular}

\vspace{10pt}

\begin{tabular}{rl}
\multicolumn{2}{l}{\textbf{PERFEKT}}\\
& haben + \textbf{verschrieben}\\
\end{tabular}
\hfill
\begin{tabular}{rl}
\multicolumn{2}{l}{\textbf{FUTURE I}}\\
& werden + \textbf{verschreiben}\\
\end{tabular}
\hfill
\begin{tabular}{rl}
\multicolumn{2}{l}{\textbf{PLUSQUAMPERFEKT}}\\
& hatten + \textbf{verschrieben}\\
\end{tabular}

\vspace{10pt}

\begin{tabular}{lll}
\multicolumn{3}{l}{\textbf{MIT PRÄPOSITIONEN}}\\
\end{tabular}
\hfill
\begin{tabular}{lll}
\multicolumn{3}{l}{\textbf{TRENNBARE VERBEN}}\\
\end{tabular}
\vspace{-5pt}
\end{verbbox}

%verschwenden (to waste)
\begin{verbbox}\addcontentsline{toc}{subsubsection}{verschwenden (to waste)}
    {\Large \bf verschwenden} (to waste) \hfill Type: \textbf{regular} \hfill Auxiliary Verb: \textbf{haben} 
    
     \vspace{5pt}

    \hrule
    
    \vspace{5pt}

  \begin{tabular}{rl}
     \multicolumn{2}{l}{\textbf{PRÄSENS} }\\
    ich & \textbf{verschwende}\\
    du & \textbf{verschwendest}\\
    er/sie/es & \textbf{verschwendet}\\
    wir & \textbf{verschwenden}\\
    ihr & \textbf{verschwendet}\\
    sie/Sie & \textbf{verschwenden}\\
  \end{tabular}
  \hfill
     \begin{tabular}{rl}
    \multicolumn{2}{l}{\textbf{PRÄTERITUM} }\\
    ich & \textbf{verschwendete}\\
    du & \textbf{verschwendetest}\\
    er/sie/es & \textbf{verschwendete}\\
    wir & \textbf{verschwendeten}\\
    ihr & \textbf{verschwendetet}\\
    sie/Sie & \textbf{verschwendeten}\\
  \end{tabular}
  \hfill
  \begin{tabular}{rl}
     \multicolumn{2}{l}{\textbf{IMPERATIV} }\\
   & \textbf{verschwende (du)!}\\
   & \textbf{verschwenden wir!}\\
   & \textbf{verschwendet (ihr)!}\\
   & \textbf{verschwenden Sie!}\\
   & \vspace{10pt}
  \end{tabular}

    \vspace{10pt}

 \begin{tabular}{rl}
     \multicolumn{2}{l}{\textbf{PERFEKT}}\\
   & haben + \textbf{verschwendet}\\
  \end{tabular}
\hfill
   \begin{tabular}{rl}
    \multicolumn{2}{l}{\textbf{FUTURE I} }\\
   & werde + \textbf{verschwenden}\\
  \end{tabular}
\hfill
   \begin{tabular}{rl}
      \multicolumn{2}{l}{\textbf{PLUSQUAMPERFEKT}}\\
   & hatten + \textbf{verschwendet}\\
  \end{tabular}

   \vspace{10pt}

   \begin{tabular}{lll}
      \multicolumn{3}{l}{\textbf{MIT PRÄPOSITIONEN}}\\
& \textbf{verschwenden an} + Akk & (to waste on)\\
& \textbf{verschwenden für} + Akk & (to waste for)\\
\end{tabular}
  \hfill
   \begin{tabular}{lll}
      \multicolumn{3}{l}{\textbf{TRENNBARE VERBEN}}\\
 & \textbf{---} & (none common)\\
\vspace{10pt}
  \end{tabular}

  \vspace{-5pt}

\end{verbbox}

%verschwinden (to disappear)
\begin{verbbox}\addcontentsline{toc}{subsubsection}{verschwinden (to disappear)}
    {\Large \bf verschwinden} (to disappear) \hfill Type: \textbf{irregular} \hfill Auxiliary Verb: \textbf{sein} 
    
     \vspace{5pt}

    \hrule
    
    \vspace{5pt}

  \begin{tabular}{rl}
     \multicolumn{2}{l}{\textbf{PRÄSENS} }\\
    ich & \textbf{verschwinde}\\
    du & \textbf{verschwindest}\\
    er/sie/es & \textbf{verschwindet}\\
    wir & \textbf{verschwinden}\\
    ihr & \textbf{verschwindet}\\
    sie/Sie & \textbf{verschwinden}\\
  \end{tabular}
  \hfill
     \begin{tabular}{rl}
    \multicolumn{2}{l}{\textbf{PRÄTERITUM} }\\
    ich & \textbf{verschwand}\\
    du & \textbf{verschwandest}\\
    er/sie/es & \textbf{verschwand}\\
    wir & \textbf{verschwanden}\\
    ihr & \textbf{verschwandet}\\
    sie/Sie & \textbf{verschwanden}\\
  \end{tabular}
  \hfill
  \begin{tabular}{rl}
     \multicolumn{2}{l}{\textbf{IMPERATIV} }\\
   & \textbf{verschwind(e) (du)!}\\
   & \textbf{verschwinden wir!}\\
   & \textbf{verschwindet (ihr)!}\\
   & \textbf{verschwinden Sie!}\\
   & \vspace{10pt}
  \end{tabular}

    \vspace{10pt}

 \begin{tabular}{rl}
     \multicolumn{2}{l}{\textbf{PERFEKT}}\\
   & sein + \textbf{verschwunden}\\
  \end{tabular}
\hfill
   \begin{tabular}{rl}
    \multicolumn{2}{l}{\textbf{FUTURE I} }\\
   & werde + \textbf{verschwinden}\\
  \end{tabular}
\hfill
   \begin{tabular}{rl}
      \multicolumn{2}{l}{\textbf{PLUSQUAMPERFEKT}}\\
   & waren + \textbf{verschwunden}\\
  \end{tabular}

   \vspace{10pt}

   \begin{tabular}{lll}
      \multicolumn{3}{l}{\textbf{MIT PRÄPOSITIONEN}}\\
& \textbf{verschwinden aus} + Dat & (to disappear from)\\
& \textbf{verschwinden in} + Dat & (to disappear in / into)\\
\end{tabular}
  \hfill
   \begin{tabular}{lll}
      \multicolumn{3}{l}{\textbf{TRENNBARE VERBEN}}\\
 & \textbf{---} & (none common)\\
\vspace{10pt}
  \end{tabular}

  \vspace{-5pt}

\end{verbbox}

%verstehen (A1) – to understand
\begin{verbbox-acc}\addcontentsline{toc}{subsubsection}{verstehen (to understand)}
    {\Large \bf verstehen} (to understand) \hfill Type: \textbf{irregular, + Akkusativ} \hfill Auxiliary Verb: \textbf{haben} 
    
     \vspace{5pt}

    \hrule
    
    \vspace{5pt}

  \begin{tabular}{rl}
     \multicolumn{2}{l}{\textbf{PRÄSENS} }\\
    ich & \textbf{verstehe}\\
    du & \textbf{verstehst}\\
    er/sie/es & \textbf{versteht}\\
    wir & \textbf{verstehen}\\
    ihr & \textbf{versteht}\\
    sie/Sie & \textbf{verstehen}\\
  \end{tabular}
  \hfill
     \begin{tabular}{rl}
    \multicolumn{2}{l}{\textbf{PRÄTERITUM} }\\
    ich & \textbf{verstand}\\
    du & \textbf{verstandest}\\
    er/sie/es & \textbf{verstand}\\
    wir & \textbf{verstanden}\\
    ihr & \textbf{verstandet}\\
    sie/Sie & \textbf{verstanden}\\
  \end{tabular}
  \hfill
  \begin{tabular}{rl}
     \multicolumn{2}{l}{\textbf{IMPERATIV} }\\
   & \textbf{versteh(e) (du)!}\\
   & \textbf{verstehen wir!}\\
   & \textbf{versteht (ihr)!}\\
   & \textbf{verstehen Sie!}\\
   & \vspace{10pt}
  \end{tabular}

    \vspace{10pt}

 \begin{tabular}{rl}
     \multicolumn{2}{l}{\textbf{PERFEKT}}\\
  &  haben + \textbf{verstanden}\\
  \end{tabular}
\hfill
   \begin{tabular}{rl}
   
    \multicolumn{2}{l}{\textbf{FUTURE I} }\\
  &  werden + \textbf{verstehen}\\
 
  \end{tabular}
\hfill
   \begin{tabular}{rl}
      \multicolumn{2}{l}{\textbf{PLUSQUAMPERFEKT}}\\
  &  hatten + \textbf{verstanden}\\
  \end{tabular}
  
  \vspace{-5pt}


\end{verbbox-acc}

%versuchen (A2) – to try
\begin{verbbox-acc}\addcontentsline{toc}{subsubsection}{versuchen (to try)}
    {\Large \bf versuchen} (to try) \hfill Type: \textbf{irregular, + Akkusativ} \hfill Auxiliary Verb: \textbf{haben} 
    
     \vspace{5pt}

    \hrule
    
    \vspace{5pt}

  \begin{tabular}{rl}
     \multicolumn{2}{l}{\textbf{PRÄSENS} }\\
    ich & \textbf{versuche}\\
    du & \textbf{versuchst}\\
    er/sie/es & \textbf{versucht}\\
    wir & \textbf{versuchen}\\
    ihr & \textbf{versucht}\\
    sie/Sie & \textbf{versuchen}\\
  \end{tabular}
  \hfill
     \begin{tabular}{rl}
    \multicolumn{2}{l}{\textbf{PRÄTERITUM} }\\
    ich & \textbf{versuchte}\\
    du & \textbf{versuchtest}\\
    er/sie/es & \textbf{versuchte}\\
    wir & \textbf{versuchten}\\
    ihr & \textbf{versuchtet}\\
    sie/Sie & \textbf{versuchten}\\
  \end{tabular}
  \hfill
  \begin{tabular}{rl}
     \multicolumn{2}{l}{\textbf{IMPERATIV} }\\
   & \textbf{versuch(e) (du)!}\\
   & \textbf{versuchen wir!}\\
   & \textbf{versucht (ihr)!}\\
   & \textbf{versuchen Sie!}\\
   & \vspace{10pt}
  \end{tabular}

    \vspace{10pt}

 \begin{tabular}{rl}
     \multicolumn{2}{l}{\textbf{PERFEKT}}\\
  &  haben + \textbf{versucht}\\
  \end{tabular}
\hfill
   \begin{tabular}{rl}
   
    \multicolumn{2}{l}{\textbf{FUTURE I} }\\
  &  werden + \textbf{versuchen}\\
 
  \end{tabular}
\hfill
   \begin{tabular}{rl}
      \multicolumn{2}{l}{\textbf{PLUSQUAMPERFEKT}}\\
   & hatten + \textbf{versucht}\\
  \end{tabular}
  
  \vspace{-5pt}


\end{verbbox-acc}

%vertiefen (to deepen)
\begin{verbbox-acc}\addcontentsline{toc}{subsubsection}{vertiefen (to deepen)}
{\Large \bf vertiefen} (to deepen) \hfill Type: \textbf{regular, + Akkusativ} \hfill Auxiliary Verb: \textbf{haben}

\vspace{5pt}
\hrule
\vspace{5pt}

\begin{tabular}{rl}
\multicolumn{2}{l}{\textbf{PRÄSENS}}\\
ich & \textbf{vertiefe}\\
du & \textbf{vertiefst}\\
er/sie/es & \textbf{vertieft}\\
wir & \textbf{vertiefen}\\
ihr & \textbf{vertieft}\\
sie/Sie & \textbf{vertiefen}\\
\end{tabular}
\hfill
\begin{tabular}{rl}
\multicolumn{2}{l}{\textbf{PRÄTERITUM}}\\
ich & \textbf{vertiefte}\\
du & \textbf{vertieftest}\\
er/sie/es & \textbf{vertiefte}\\
wir & \textbf{vertieften}\\
ihr & \textbf{vertieftet}\\
sie/Sie & \textbf{vertieften}\\
\end{tabular}
\hfill
\begin{tabular}{rl}
\multicolumn{2}{l}{\textbf{IMPERATIV}}\\
& \textbf{vertiefe (du)!}\\
& \textbf{vertiefen wir!}\\
& \textbf{vertieft (ihr)!}\\
& \textbf{vertiefen Sie!}\\
\end{tabular}

\vspace{10pt}

\begin{tabular}{rl}
\multicolumn{2}{l}{\textbf{PERFEKT}}\\
& haben + \textbf{vertieft}\\
\end{tabular}
\hfill
\begin{tabular}{rl}
\multicolumn{2}{l}{\textbf{FUTURE I}}\\
& werden + \textbf{vertiefen}\\
\end{tabular}
\hfill
\begin{tabular}{rl}
\multicolumn{2}{l}{\textbf{PLUSQUAMPERFEKT}}\\
& hatten + \textbf{vertieft}\\
\end{tabular}

\vspace{10pt}

\begin{tabular}{lll}
\multicolumn{3}{l}{\textbf{MIT PRÄPOSITIONEN}}\\
& \textbf{vertiefen in} + Akk & (to deepen into / immerse in)\\
\end{tabular}
\hfill
\begin{tabular}{lll}
\multicolumn{3}{l}{\textbf{TRENNBARE VERBEN}}\\
& ---\\
\end{tabular}

\vspace{-5pt}
\end{verbbox-acc}

%vertrauen (to trust)
\begin{verbbox}\addcontentsline{toc}{subsubsection}{vertrauen (to trust)}
{\Large \bf vertrauen} (to trust) \hfill Type: \textbf{regular} \hfill Auxiliary Verb: \textbf{haben}

\vspace{5pt}
\hrule
\vspace{5pt}

\begin{tabular}{rl}
\multicolumn{2}{l}{\textbf{PRÄSENS}}\\
ich & \textbf{vertraue}\\
du & \textbf{vertraust}\\
er/sie/es & \textbf{vertraut}\\
wir & \textbf{vertrauen}\\
ihr & \textbf{vertraut}\\
sie/Sie & \textbf{vertrauen}\\
\end{tabular}
\hfill
\begin{tabular}{rl}
\multicolumn{2}{l}{\textbf{PRÄTERITUM}}\\
ich & \textbf{vertraute}\\
du & \textbf{vertrautest}\\
er/sie/es & \textbf{vertraute}\\
wir & \textbf{vertrauten}\\
ihr & \textbf{vertrautet}\\
sie/Sie & \textbf{vertrauten}\\
\end{tabular}
\hfill
\begin{tabular}{rl}
\multicolumn{2}{l}{\textbf{IMPERATIV}}\\
& \textbf{vertraue (du)!}\\
& \textbf{vertrauen wir!}\\
& \textbf{vertraut (ihr)!}\\
& \textbf{vertrauen Sie!}\\
\end{tabular}

\vspace{10pt}

\begin{tabular}{rl}
\multicolumn{2}{l}{\textbf{PERFEKT}}\\
& haben + \textbf{vertraut}\\
\end{tabular}
\hfill
\begin{tabular}{rl}
\multicolumn{2}{l}{\textbf{FUTURE I}}\\
& werden + \textbf{vertrauen}\\
\end{tabular}
\hfill
\begin{tabular}{rl}
\multicolumn{2}{l}{\textbf{PLUSQUAMPERFEKT}}\\
& hatten + \textbf{vertraut}\\
\end{tabular}

\vspace{10pt}

\begin{tabular}{lll}
\multicolumn{3}{l}{\textbf{MIT PRÄPOSITIONEN}}\\
& \textbf{vertrauen auf} + Akk & (to rely on)\\
& \textbf{vertrauen in} + Akk & (to trust in)\\
& \textbf{vertrauen} + Dat & (to trust someone)\\
\end{tabular}
\hfill
\begin{tabular}{lll}
\multicolumn{3}{l}{\textbf{TRENNBARE VERBEN}}\\
\vspace{30pt}\\
\end{tabular}
\vspace{-5pt}
\end{verbbox}

%verzeihen (to forgive)
\begin{verbbox-dat}\addcontentsline{toc}{subsubsection}{verzeihen (to forgive)}
    {\Large \bf verzeihen} (to forgive) \hfill Type: \textbf{irregular, + Dativ} \hfill Auxiliary Verb: \textbf{haben} 
    
     \vspace{5pt}

    \hrule
    
    \vspace{5pt}

  \begin{tabular}{rl}
     \multicolumn{2}{l}{\textbf{PRÄSENS} }\\
    ich & \textbf{verzeihe}\\
    du & \textbf{verzeihst}\\
    er/sie/es & \textbf{verzeiht}\\
    wir & \textbf{verzeihen}\\
    ihr & \textbf{verzeiht}\\
    sie/Sie & \textbf{verzeihen}\\
  \end{tabular}
  \hfill
     \begin{tabular}{rl}
    \multicolumn{2}{l}{\textbf{PRÄTERITUM} }\\
    ich & \textbf{verzieh}\\
    du & \textbf{verziehst}\\
    er/sie/es & \textbf{verzieh}\\
    wir & \textbf{verziehen}\\
    ihr & \textbf{verzieht}\\
    sie/Sie & \textbf{verziehen}\\
  \end{tabular}
  \hfill
  \begin{tabular}{rl}
     \multicolumn{2}{l}{\textbf{IMPERATIV} }\\
   & \textbf{verzeih(e) (du)!}\\
   & \textbf{verzeihen wir!}\\
   & \textbf{verzeiht (ihr)!}\\
   & \textbf{verzeihen Sie!}\\
   & \vspace{10pt}
  \end{tabular}

    \vspace{10pt}

 \begin{tabular}{rl}
     \multicolumn{2}{l}{\textbf{PERFEKT}}\\
   & haben + \textbf{verziehen}\\
  \end{tabular}
\hfill
   \begin{tabular}{rl}
   
    \multicolumn{2}{l}{\textbf{FUTURE I} }\\
   & werden + \textbf{verzeihen}\\
 
  \end{tabular}
\hfill
   \begin{tabular}{rl}
      \multicolumn{2}{l}{\textbf{PLUSQUAMPERFEKT}}\\
   & hatten + \textbf{verziehen}\\
  \end{tabular}

   \vspace{10pt}

   \begin{tabular}{lll}
      \multicolumn{3}{l}{\textbf{MIT PRÄPOSITIONEN}}\\
 & -- & --\\
\end{tabular}
  \hfill
   \begin{tabular}{lll}
      \multicolumn{3}{l}{\textbf{TRENNBARE VERBEN}}\\
 & -- & --\\
  \end{tabular}


  \vspace{-5pt}


\end{verbbox-dat}

%verzichten - (to do without / renounce)
\begin{verbbox}\addcontentsline{toc}{subsubsection}{verzichten (to do without / renounce)}
{\Large \bf verzichten} (to do without) \hfill Type: \textbf{regular} \hfill Auxiliary Verb: \textbf{haben}

\vspace{5pt}
\hrule
\vspace{5pt}

\begin{tabular}{rl}
\multicolumn{2}{l}{\textbf{PRÄSENS}}\\
ich & \textbf{verzichte}\\
du & \textbf{verzichtest}\\
er/sie/es & \textbf{verzichtet}\\
wir & \textbf{verzichten}\\
ihr & \textbf{verzichtet}\\
sie/Sie & \textbf{verzichten}\\
\end{tabular}
\hfill
\begin{tabular}{rl}
\multicolumn{2}{l}{\textbf{PRÄTERITUM}}\\
ich & \textbf{verzichtete}\\
du & \textbf{verzichtetest}\\
er/sie/es & \textbf{verzichtete}\\
wir & \textbf{verzichteten}\\
ihr & \textbf{verzichtetet}\\
sie/Sie & \textbf{verzichteten}\\
\end{tabular}
\hfill
\begin{tabular}{rl}
\multicolumn{2}{l}{\textbf{IMPERATIV}}\\
& \textbf{verzichte (du)!}\\
& \textbf{verzichten wir!}\\
& \textbf{verzichtet (ihr)!}\\
& \textbf{verzichten Sie!}\\
\end{tabular}

\vspace{10pt}

\begin{tabular}{rl}
\multicolumn{2}{l}{\textbf{PERFEKT}}\\
& haben + \textbf{verzichtet}\\
\end{tabular}
\hfill
\begin{tabular}{rl}
\multicolumn{2}{l}{\textbf{FUTURE I}}\\
& werden + \textbf{verzichten}\\
\end{tabular}
\hfill
\begin{tabular}{rl}
\multicolumn{2}{l}{\textbf{PLUSQUAMPERFEKT}}\\
& hatten + \textbf{verzichtet}\\
\end{tabular}

\vspace{10pt}

\begin{tabular}{lll}
\multicolumn{3}{l}{\textbf{MIT PRÄPOSITIONEN}}\\
& \textbf{verzichten auf} + Akk & (to give up / do without)\\
\end{tabular}
\hfill
\begin{tabular}{lll}
\multicolumn{3}{l}{\textbf{TRENNBARE VERBEN}}\\
& --- & \\
\end{tabular}
\vspace{-5pt}
\end{verbbox}

%verzieren (to decorate / adorn)
\begin{verbbox-acc}\addcontentsline{toc}{subsubsection}{verzieren (to decorate / adorn)}
{\Large \bf verzieren} (to decorate / adorn) \hfill Type: \textbf{regular, + Akkusativ} \hfill Auxiliary Verb: \textbf{haben}

\vspace{5pt}
\hrule
\vspace{5pt}

\begin{tabular}{rl}
\multicolumn{2}{l}{\textbf{PRÄSENS}}\\
ich & \textbf{verziere}\\
du & \textbf{verzierst}\\
er/sie/es & \textbf{verziert}\\
wir & \textbf{verzieren}\\
ihr & \textbf{verziert}\\
sie/Sie & \textbf{verzieren}\\
\end{tabular}
\hfill
\begin{tabular}{rl}
\multicolumn{2}{l}{\textbf{PRÄTERITUM}}\\
ich & \textbf{verzierte}\\
du & \textbf{verziertest}\\
er/sie/es & \textbf{verzierte}\\
wir & \textbf{verzierten}\\
ihr & \textbf{verziertet}\\
sie/Sie & \textbf{verzierten}\\
\end{tabular}
\hfill
\begin{tabular}{rl}
\multicolumn{2}{l}{\textbf{IMPERATIV}}\\
& \textbf{verziere (du)!}\\
& \textbf{verzieren wir!}\\
& \textbf{verziert (ihr)!}\\
& \textbf{verzieren Sie!}\\
\end{tabular}

\vspace{10pt}

\begin{tabular}{rl}
\multicolumn{2}{l}{\textbf{PERFEKT}}\\
& haben + \textbf{verziert}\\
\end{tabular}
\hfill
\begin{tabular}{rl}
\multicolumn{2}{l}{\textbf{FUTURE I}}\\
& werden + \textbf{verzieren}\\
\end{tabular}
\hfill
\begin{tabular}{rl}
\multicolumn{2}{l}{\textbf{PLUSQUAMPERFEKT}}\\
& hatten + \textbf{verziert}\\
\end{tabular}

\vspace{10pt}

\begin{tabular}{lll}
\multicolumn{3}{l}{\textbf{MIT PRÄPOSITIONEN}}\\
& \textbf{verzieren mit} + Dat & (to decorate with)\\
\end{tabular}
\hfill
\begin{tabular}{lll}
\multicolumn{3}{l}{\textbf{TRENNBARE VERBEN}}\\
& ---\\
\end{tabular}

\vspace{-5pt}
\end{verbbox-acc}

\subsection{W Verbs}

%wählen (A2) – to choose, vote
\begin{verbbox-acc}\addcontentsline{toc}{subsubsection}{wählen (to choose / vote)}
    {\Large \bf wählen} (to choose / vote) \hfill Type: \textbf{regular} \hfill Auxiliary Verb: \textbf{haben} 
    
     \vspace{5pt}

    \hrule
    
    \vspace{5pt}

  \begin{tabular}{rl}
     \multicolumn{2}{l}{\textbf{PRÄSENS} }\\
    ich & \textbf{wähle}\\
    du & \textbf{wählst}\\
    er/sie/es & \textbf{wählt}\\
    wir & \textbf{wählen}\\
    ihr & \textbf{wählt}\\
    sie/Sie & \textbf{wählen}\\
  \end{tabular}
  \hfill
  \begin{tabular}{rl}
    \multicolumn{2}{l}{\textbf{PRÄTERITUM} }\\
    ich & \textbf{wählte}\\
    du & \textbf{wähltest}\\
    er/sie/es & \textbf{wählte}\\
    wir & \textbf{wählten}\\
    ihr & \textbf{wähltet}\\
    sie/Sie & \textbf{wählten}\\
  \end{tabular}
  \hfill
  \begin{tabular}{rl}
     \multicolumn{2}{l}{\textbf{IMPERATIV} }\\
   & \textbf{wähle (du)!}\\
   & \textbf{wählen wir!}\\
   & \textbf{wählt (ihr)!}\\
   & \textbf{wählen Sie!}\\
   & \vspace{10pt}
  \end{tabular}

 \vspace{10pt}

 \begin{tabular}{rl}
     \multicolumn{2}{l}{\textbf{PERFEKT}}\\
   & haben + \textbf{gewählt}\\
  \end{tabular}
\hfill
\begin{tabular}{rl}
    \multicolumn{2}{l}{\textbf{FUTURE I}}\\
   & werden + \textbf{wählen}\\
  \end{tabular}
\hfill
\begin{tabular}{rl}
    \multicolumn{2}{l}{\textbf{PLUSQUAMPERFEKT}}\\
   & hatten + \textbf{gewählt}\\
  \end{tabular}

 \vspace{10pt}

\begin{tabular}{lll}
      \multicolumn{3}{l}{\textbf{MIT PRÄPOSITIONEN}}\\
& \textbf{wählen für} + Akk & (to vote for / choose for)\\
& \textbf{wählen zwischen} + Dat & (to choose between)\\
& \textbf{wählen zu} + Dat & (to elect as / choose as)\\
\end{tabular}
  \hfill
\begin{tabular}{lll}
      \multicolumn{3}{l}{\textbf{TRENNBARE VERBEN}}\\
 & \textbf{aus$|$wählen} & (to select / pick out)\\
\vspace{20pt}
  \end{tabular}
  
  \vspace{-5pt}

\end{verbbox-acc}

%warten (A1) – to wait
\begin{verbbox}\addcontentsline{toc}{subsubsection}{warten (to wait)}
    {\Large \bf warten} (to wait) \hfill Type: \textbf{regular} \hfill Auxiliary Verb: \textbf{haben} 
    
     \vspace{5pt}

    \hrule
    
    \vspace{5pt}

  \begin{tabular}{rl}
     \multicolumn{2}{l}{\textbf{PRÄSENS} }\\
    ich & \textbf{warte}\\
    du & \textbf{wartest}\\
    er/sie/es & \textbf{wartet}\\
    wir & \textbf{warten}\\
    ihr & \textbf{wartet}\\
    sie/Sie & \textbf{warten}\\
  \end{tabular}
  \hfill
     \begin{tabular}{rl}
    \multicolumn{2}{l}{\textbf{PRÄTERITUM} }\\
    ich & \textbf{wartete}\\
    du & \textbf{wartetest}\\
    er/sie/es & \textbf{wartete}\\
    wir & \textbf{warteten}\\
    ihr & \textbf{wartetet}\\
    sie/Sie & \textbf{warteten}\\
  \end{tabular}
  \hfill
  \begin{tabular}{rl}
     \multicolumn{2}{l}{\textbf{IMPERATIV} }\\
   & \textbf{wart(e) (du)!}\\
   & \textbf{warten wir!}\\
   & \textbf{wartet (ihr)!}\\
   & \textbf{warten Sie!}\\
   & \vspace{10pt}
  \end{tabular}

    \vspace{10pt}

 \begin{tabular}{rl}
     \multicolumn{2}{l}{\textbf{PERFEKT}}\\
  &  haben + \textbf{gewartet}\\
  \end{tabular}
\hfill
   \begin{tabular}{rl}
   
    \multicolumn{2}{l}{\textbf{FUTURE I} }\\
  &  werden + \textbf{warten}\\
 
  \end{tabular}
\hfill
   \begin{tabular}{rl}
      \multicolumn{2}{l}{\textbf{PLUSQUAMPERFEKT}}\\
  &  hatten + \textbf{gewartet}\\
  \end{tabular}
  
  \vspace{10pt}



\begin{tabular}{lll}
\multicolumn{3}{l}{\textbf{MIT PRÄPOSITIONEN}}\\
& ---\\
\end{tabular}
\hfill
\begin{tabular}{lll}
\multicolumn{3}{l}{\textbf{TRENNBARE VERBEN}}\\
& ---\\
\end{tabular}

  \vspace{-5pt}

\end{verbbox}

%waschen (A1) – to wash
\begin{verbbox-acc}\addcontentsline{toc}{subsubsection}{waschen (to wash)}
    {\Large \bf waschen} (to wash) \hfill Type: \textbf{regular, + Akkusativ} \hfill Auxiliary Verb: \textbf{haben} 
    
     \vspace{5pt}

    \hrule
    
    \vspace{5pt}

  \begin{tabular}{rl}
     \multicolumn{2}{l}{\textbf{PRÄSENS} }\\
    ich & \textbf{wasche}\\
    du & \textbf{wäschst}\\
    er/sie/es & \textbf{wäscht}\\
    wir & \textbf{waschen}\\
    ihr & \textbf{wascht}\\
    sie/Sie & \textbf{waschen}\\
  \end{tabular}
  \hfill
     \begin{tabular}{rl}
    \multicolumn{2}{l}{\textbf{PRÄTERITUM} }\\
    ich & \textbf{wusch}\\
    du & \textbf{wuschst}\\
    er/sie/es & \textbf{wusch}\\
    wir & \textbf{wuschen}\\
    ihr & \textbf{wuscht}\\
    sie/Sie & \textbf{wuschen}\\
  \end{tabular}
  \hfill
  \begin{tabular}{rl}
     \multicolumn{2}{l}{\textbf{IMPERATIV} }\\
   & \textbf{wasch(e) (du)!}\\
   & \textbf{waschen wir!}\\
   & \textbf{wascht (ihr)!}\\
   & \textbf{waschen Sie!}\\
   & \vspace{10pt}
  \end{tabular}

    \vspace{10pt}

 \begin{tabular}{rl}
     \multicolumn{2}{l}{\textbf{PERFEKT}}\\
  &  haben + \textbf{gewaschen}\\
  \end{tabular}
\hfill
   \begin{tabular}{rl}
   
    \multicolumn{2}{l}{\textbf{FUTURE I} }\\
  &  werden + \textbf{waschen}\\
 
  \end{tabular}
\hfill
   \begin{tabular}{rl}
      \multicolumn{2}{l}{\textbf{PLUSQUAMPERFEKT}}\\
  &  hatten + \textbf{gewaschen}\\
  \end{tabular}
  
  \vspace{-5pt}


\end{verbbox-acc}

%wechseln (to change / switch)
\begin{verbbox}\addcontentsline{toc}{subsubsection}{wechseln (to change / switch)}
{\Large \bf wechseln} (to change / switch) \hfill Type: \textbf{regular} \hfill Auxiliary Verb: \textbf{haben / sein (rarely)} 

\vspace{5pt}
\hrule
\vspace{5pt}

\begin{tabular}{rl}
\multicolumn{2}{l}{\textbf{PRÄSENS}}\\
ich & \textbf{wechsle}\\
du & \textbf{wechselst}\\
er/sie/es & \textbf{wechselt}\\
wir & \textbf{wechseln}\\
ihr & \textbf{wechselt}\\
sie/Sie & \textbf{wechseln}\\
\end{tabular}
\hfill
\begin{tabular}{rl}
\multicolumn{2}{l}{\textbf{PRÄTERITUM}}\\
ich & \textbf{wechselte}\\
du & \textbf{wechseltest}\\
er/sie/es & \textbf{wechselte}\\
wir & \textbf{wechselten}\\
ihr & \textbf{wechseltet}\\
sie/Sie & \textbf{wechselten}\\
\end{tabular}
\hfill
\begin{tabular}{rl}
\multicolumn{2}{l}{\textbf{IMPERATIV}}\\
& \textbf{wechsle (du)!}\\
& \textbf{wechseln wir!}\\
& \textbf{wechselt (ihr)!}\\
& \textbf{wechseln Sie!}\\
\end{tabular}

\vspace{10pt}

\begin{tabular}{rl}
\multicolumn{2}{l}{\textbf{PERFEKT}}\\
& haben + \textbf{gewechselt}\\
\end{tabular}
\hfill
\begin{tabular}{rl}
\multicolumn{2}{l}{\textbf{FUTURE I}}\\
& werden + \textbf{wechseln}\\
\end{tabular}
\hfill
\begin{tabular}{rl}
\multicolumn{2}{l}{\textbf{PLUSQUAMPERFEKT}}\\
& hatten + \textbf{gewechselt}\\
\end{tabular}

\vspace{10pt}

\begin{tabular}{lll}
\multicolumn{3}{l}{\textbf{MIT PRÄPOSITIONEN}}\\
& \textbf{wechseln zu} + Dat & (to switch to)\\
& \textbf{wechseln mit} + Dat & (to exchange with)\\
\end{tabular}
\hfill
\begin{tabular}{lll}
\multicolumn{3}{l}{\textbf{TRENNBARE VERBEN}}\\
& \textbf{um$|$wechseln} & (to change / swap around)\\
\vspace{10pt}
\end{tabular}
\vspace{-5pt}
\end{verbbox}

%weinen (to cry)
\begin{verbbox}\addcontentsline{toc}{subsubsection}{weinen (to cry)}
{\Large \bf weinen} (to cry) \hfill Type: \textbf{regular} \hfill Auxiliary Verb: \textbf{haben}

\vspace{5pt}
\hrule
\vspace{5pt}

\begin{tabular}{rl}
\multicolumn{2}{l}{\textbf{PRÄSENS}}\\
ich & \textbf{weine}\\
du & \textbf{weinst}\\
er/sie/es & \textbf{weint}\\
wir & \textbf{weinen}\\
ihr & \textbf{weint}\\
sie/Sie & \textbf{weinen}\\
\end{tabular}
\hfill
\begin{tabular}{rl}
\multicolumn{2}{l}{\textbf{PRÄTERITUM}}\\
ich & \textbf{weinte}\\
du & \textbf{weintest}\\
er/sie/es & \textbf{weinte}\\
wir & \textbf{weinten}\\
ihr & \textbf{weintet}\\
sie/Sie & \textbf{weinten}\\
\end{tabular}
\hfill
\begin{tabular}{rl}
\multicolumn{2}{l}{\textbf{IMPERATIV}}\\
& \textbf{wein(e) (du)!}\\
& \textbf{weinen wir!}\\
& \textbf{weint (ihr)!}\\
& \textbf{weinen Sie!}\\
\end{tabular}

\vspace{10pt}

\begin{tabular}{rl}
\multicolumn{2}{l}{\textbf{PERFEKT}}\\
& haben + \textbf{geweint}\\
\end{tabular}
\hfill
\begin{tabular}{rl}
\multicolumn{2}{l}{\textbf{FUTURE I}}\\
& werden + \textbf{weinen}\\
\end{tabular}
\hfill
\begin{tabular}{rl}
\multicolumn{2}{l}{\textbf{PLUSQUAMPERFEKT}}\\
& hatten + \textbf{geweint}\\
\end{tabular}

\vspace{10pt}

\begin{tabular}{lll}
\multicolumn{3}{l}{\textbf{MIT PRÄPOSITIONEN}}\\
& \textbf{weinen über} + Akk & (to cry about)\\
\end{tabular}
\hfill
\begin{tabular}{lll}
\multicolumn{3}{l}{\textbf{TRENNBARE VERBEN}}\\
& ---\\
\end{tabular}
\vspace{-5pt}
\end{verbbox}

%werfen (to throw)
\begin{verbbox}\addcontentsline{toc}{subsubsection}{werfen (to throw)}
    {\Large \bf werfen} (to throw) \hfill Type: \textbf{irregular} \hfill Auxiliary Verb: \textbf{haben} 
    
     \vspace{5pt}

    \hrule
    
    \vspace{5pt}

  \begin{tabular}{rl}
     \multicolumn{2}{l}{\textbf{PRÄSENS} }\\
    ich & \textbf{werfe}\\
    du & \textbf{wirfst}\\
    er/sie/es & \textbf{wirft}\\
    wir & \textbf{werfen}\\
    ihr & \textbf{werft}\\
    sie/Sie & \textbf{werfen}\\
  \end{tabular}
  \hfill
     \begin{tabular}{rl}
    \multicolumn{2}{l}{\textbf{PRÄTERITUM} }\\
    ich & \textbf{warf}\\
    du & \textbf{warfst}\\
    er/sie/es & \textbf{warf}\\
    wir & \textbf{warfen}\\
    ihr & \textbf{warft}\\
    sie/Sie & \textbf{warfen}\\
  \end{tabular}
  \hfill
  \begin{tabular}{rl}
     \multicolumn{2}{l}{\textbf{IMPERATIV} }\\
   & \textbf{wirf (du)!}\\
   & \textbf{werfen wir!}\\
   & \textbf{werft (ihr)!}\\
   & \textbf{werfen Sie!}\\
   & \vspace{10pt}
  \end{tabular}

    \vspace{10pt}

 \begin{tabular}{rl}
     \multicolumn{2}{l}{\textbf{PERFEKT}}\\
   & haben + \textbf{geworfen}\\
  \end{tabular}
\hfill
   \begin{tabular}{rl}
    \multicolumn{2}{l}{\textbf{FUTURE I} }\\
   & werde + \textbf{werfen}\\
  \end{tabular}
\hfill
   \begin{tabular}{rl}
      \multicolumn{2}{l}{\textbf{PLUSQUAMPERFEKT}}\\
   & hatten + \textbf{geworfen}\\
  \end{tabular}

   \vspace{10pt}

   \begin{tabular}{lll}
      \multicolumn{3}{l}{\textbf{MIT PRÄPOSITIONEN}}\\
& \textbf{werfen auf} + Akk & (to throw onto / at)\\
& \textbf{werfen in} + Akk & (to throw into)\\
\end{tabular}
  \hfill
   \begin{tabular}{lll}
      \multicolumn{3}{l}{\textbf{TRENNBARE VERBEN}}\\
 & \textbf{weg$|$werfen} & (to throw away)\\
 & \textbf{ein$|$werfen} & (to throw in / insert)\\
\vspace{10pt}
  \end{tabular}

  \vspace{-5pt}

\end{verbbox}


%widersprechen (to contradict)
\begin{verbbox-dat}\addcontentsline{toc}{subsubsection}{widersprechen (to contradict)}
    {\Large \bf widersprechen} (to contradict) \hfill Type: \textbf{irregular, + Dativ} \hfill Auxiliary Verb: \textbf{haben} 
    
     \vspace{5pt}

    \hrule
    
    \vspace{5pt}

  \begin{tabular}{rl}
     \multicolumn{2}{l}{\textbf{PRÄSENS} }\\
    ich & \textbf{widerspreche}\\
    du & \textbf{widersprichst}\\
    er/sie/es & \textbf{widerspricht}\\
    wir & \textbf{widersprechen}\\
    ihr & \textbf{widersprecht}\\
    sie/Sie & \textbf{widersprechen}\\
  \end{tabular}
  \hfill
     \begin{tabular}{rl}
    \multicolumn{2}{l}{\textbf{PRÄTERITUM} }\\
    ich & \textbf{widersprach}\\
    du & \textbf{widersprachst}\\
    er/sie/es & \textbf{widersprach}\\
    wir & \textbf{widersprachen}\\
    ihr & \textbf{widerspracht}\\
    sie/Sie & \textbf{widersprachen}\\
  \end{tabular}
  \hfill
  \begin{tabular}{rl}
     \multicolumn{2}{l}{\textbf{IMPERATIV} }\\
   & \textbf{widersprich (du)!}\\
   & \textbf{widersprechen wir!}\\
   & \textbf{widersprecht (ihr)!}\\
   & \textbf{widersprechen Sie!}\\
   & \vspace{10pt}
  \end{tabular}

    \vspace{10pt}

 \begin{tabular}{rl}
     \multicolumn{2}{l}{\textbf{PERFEKT}}\\
   & haben + \textbf{widersprochen}\\
  \end{tabular}
\hfill
   \begin{tabular}{rl}
   
    \multicolumn{2}{l}{\textbf{FUTURE I} }\\
   & werden + \textbf{widersprechen}\\
 
  \end{tabular}
\hfill
   \begin{tabular}{rl}
      \multicolumn{2}{l}{\textbf{PLUSQUAMPERFEKT}}\\
   & hatten + \textbf{widersprochen}\\
  \end{tabular}

   \vspace{10pt}

   \begin{tabular}{lll}
      \multicolumn{3}{l}{\textbf{MIT PRÄPOSITIONEN}}\\
 & -- & --\\
\end{tabular}
  \hfill
   \begin{tabular}{lll}
      \multicolumn{3}{l}{\textbf{TRENNBARE VERBEN}}\\
 & -- & --\\
  \end{tabular}


  \vspace{-5pt}


\end{verbbox-dat}

%wirken - to appear
\begin{verbbox}\addcontentsline{toc}{subsubsection}{wirken (to appear)}
    {\Large \bf wirken} (to appear) \hfill Type: \textbf{regular} \hfill Auxiliary Verb: \textbf{haben} 
    
     \vspace{5pt}

    \hrule
    
    \vspace{5pt}

  \begin{tabular}{rl}
     \multicolumn{2}{l}{\textbf{PRÄSENS} }\\
    ich & \textbf{wirke}\\
    du & \textbf{wirkst}\\
    er/sie/es & \textbf{wirkt}\\
    wir & \textbf{wirken}\\
    ihr & \textbf{wirkt}\\
    sie/Sie & \textbf{wirken}\\
  \end{tabular}
  \hfill
     \begin{tabular}{rl}
    \multicolumn{2}{l}{\textbf{PRÄTERITUM} }\\
    ich & \textbf{wirkte}\\
    du & \textbf{wirktest}\\
    er/sie/es & \textbf{wirkte}\\
    wir & \textbf{wirkten}\\
    ihr & \textbf{wirktet}\\
    sie/Sie & \textbf{wirkten}\\
  \end{tabular}
  \hfill
  \begin{tabular}{rl}
     \multicolumn{2}{l}{\textbf{IMPERATIV} }\\
   & \textbf{wirk(e) (du)!}\\
   & \textbf{wirken wir!}\\
   & \textbf{wirkt (ihr)!}\\
   & \textbf{wirken Sie!}\\
   & \vspace{10pt}
  \end{tabular}

    \vspace{10pt}

 \begin{tabular}{rl}
     \multicolumn{2}{l}{\textbf{PERFEKT}}\\
  &  haben + \textbf{gewirkt}\\
  \end{tabular}
\hfill
   \begin{tabular}{rl}
   
    \multicolumn{2}{l}{\textbf{FUTURE I} }\\
  &  werden + \textbf{wirken}\\
 
  \end{tabular}
\hfill
   \begin{tabular}{rl}
      \multicolumn{2}{l}{\textbf{PLUSQUAMPERFEKT}}\\
  &  hatten + \textbf{gewirkt}\\
  \end{tabular}

   \vspace{10pt}

   \begin{tabular}{lll}
      \multicolumn{3}{l}{\textbf{MIT PRÄPOSITIONEN}}\\
& \textbf{wirken auf} + Akk & (to influence)\\
& \textbf{wirken gegen} + Akk & (to work against)\\
&\textbf{wirken als} + Nom & (to function as)\\
&\textbf{wirken wie} + Nom & (to seem like)\\
&\textbf{wirken durch} + Akk & (to have an effect through)\\
&\textbf{wirken mit} + Dat & (to contribute with)
  \end{tabular}
  \hfill
   \begin{tabular}{lll}
      \multicolumn{3}{l}{\textbf{TRENNBARE VERBEN}}\\
 & \textbf{mit$|$wirken} & (to participate to)\\
\vspace{40pt}
  \end{tabular}


  \vspace{-5pt}


\end{verbbox}


%wissen (A1) – to know (a fact)
\begin{verbbox}\addcontentsline{toc}{subsubsection}{wissen (to know)}
    {\Large \bf wissen} (to know) \hfill Type: \textbf{irregular} \hfill Auxiliary Verb: \textbf{haben} 
    
     \vspace{5pt}

    \hrule
    
    \vspace{5pt}

  \begin{tabular}{rl}
     \multicolumn{2}{l}{\textbf{PRÄSENS} }\\
    ich & \textbf{weiß}\\
    du & \textbf{weißt}\\
    er/sie/es & \textbf{weiß}\\
    wir & \textbf{wissen}\\
    ihr & \textbf{wisst}\\
    sie/Sie & \textbf{wissen}\\
  \end{tabular}
  \hfill
     \begin{tabular}{rl}
    \multicolumn{2}{l}{\textbf{PRÄTERITUM} }\\
    ich & \textbf{wusste}\\
    du & \textbf{wusstest}\\
    er/sie/es & \textbf{wusste}\\
    wir & \textbf{wussten}\\
    ihr & \textbf{wusstet}\\
    sie/Sie & \textbf{wussten}\\
  \end{tabular}
  \hfill
  \begin{tabular}{rl}
     \multicolumn{2}{l}{\textbf{IMPERATIV} }\\
   & \textbf{wisse (du)!}\\
   & \textbf{wissen wir!}\\
   & \textbf{wisst (ihr)!}\\
   & \textbf{wissen Sie!}\\
   & \vspace{10pt}
  \end{tabular}

    \vspace{10pt}

 \begin{tabular}{rl}
     \multicolumn{2}{l}{\textbf{PERFEKT}}\\
  &  haben + \textbf{gewusst}\\
  \end{tabular}
\hfill
   \begin{tabular}{rl}
   
    \multicolumn{2}{l}{\textbf{FUTURE I} }\\
  &  werden + \textbf{wissen}\\
 
  \end{tabular}
\hfill
   \begin{tabular}{rl}
      \multicolumn{2}{l}{\textbf{PLUSQUAMPERFEKT}}\\
   & hatten + \textbf{gewusst}\\
  \end{tabular}
  
  \vspace{10pt}



\begin{tabular}{lll}
\multicolumn{3}{l}{\textbf{MIT PRÄPOSITIONEN}}\\
& ---\\
\end{tabular}
\hfill
\begin{tabular}{lll}
\multicolumn{3}{l}{\textbf{TRENNBARE VERBEN}}\\
& ---\\
\end{tabular}

  \vspace{-5pt}

\end{verbbox}

%wohnen (A1) – to live / reside
\begin{verbbox}\addcontentsline{toc}{subsubsection}{wohnen (to live)}
    {\Large \bf wohnen} (to live) \hfill Type: \textbf{regular} \hfill Auxiliary Verb: \textbf{haben} 
    
     \vspace{5pt}

    \hrule
    
    \vspace{5pt}

  \begin{tabular}{rl}
     \multicolumn{2}{l}{\textbf{PRÄSENS} }\\
    ich & \textbf{wohne}\\
    du & \textbf{wohnst}\\
    er/sie/es & \textbf{wohnt}\\
    wir & \textbf{wohnen}\\
    ihr & \textbf{wohnt}\\
    sie/Sie & \textbf{wohnen}\\
  \end{tabular}
  \hfill
     \begin{tabular}{rl}
    \multicolumn{2}{l}{\textbf{PRÄTERITUM} }\\
    ich & \textbf{wohnte}\\
    du & \textbf{wohntest}\\
    er/sie/es & \textbf{wohnte}\\
    wir & \textbf{wohnten}\\
    ihr & \textbf{wohntet}\\
    sie/Sie & \textbf{wohnten}\\
  \end{tabular}
  \hfill
  \begin{tabular}{rl}
     \multicolumn{2}{l}{\textbf{IMPERATIV} }\\
   & \textbf{wohn(e) (du)!}\\
   & \textbf{wohnen wir!}\\
   & \textbf{wohnt (ihr)!}\\
   & \textbf{wohnen Sie!}\\
   & \vspace{10pt}
  \end{tabular}

    \vspace{10pt}

 \begin{tabular}{rl}
     \multicolumn{2}{l}{\textbf{PERFEKT}}\\
   & haben + \textbf{gewohnt}\\
  \end{tabular}
\hfill
   \begin{tabular}{rl}
   
    \multicolumn{2}{l}{\textbf{FUTURE I} }\\
  &  werden + \textbf{wohnen}\\
 
  \end{tabular}
\hfill
   \begin{tabular}{rl}
      \multicolumn{2}{l}{\textbf{PLUSQUAMPERFEKT}}\\
   & hatten + \textbf{gewohnt}\\
  \end{tabular}
  
  \vspace{10pt}



\begin{tabular}{lll}
\multicolumn{3}{l}{\textbf{MIT PRÄPOSITIONEN}}\\
& ---\\
\end{tabular}
\hfill
\begin{tabular}{lll}
\multicolumn{3}{l}{\textbf{TRENNBARE VERBEN}}\\
& ---\\
\end{tabular}

  \vspace{-5pt}

\end{verbbox}

\subsection{X Verbs}

\subsection{Y Verbs}

\subsection{Z Verbs}

%zahlen (to pay)
\begin{verbbox}\addcontentsline{toc}{subsubsection}{zahlen (to pay)}
{\Large \bf zahlen} (to pay) \hfill Type: \textbf{regular} \hfill Auxiliary Verb: \textbf{haben}

\vspace{5pt}
\hrule
\vspace{5pt}

\begin{tabular}{rl}
\multicolumn{2}{l}{\textbf{PRÄSENS}}\\
ich & \textbf{zahle}\\
du & \textbf{zahlst}\\
er/sie/es & \textbf{zahlt}\\
wir & \textbf{zahlen}\\
ihr & \textbf{zahlt}\\
sie/Sie & \textbf{zahlen}\\
\end{tabular}
\hfill
\begin{tabular}{rl}
\multicolumn{2}{l}{\textbf{PRÄTERITUM}}\\
ich & \textbf{zahlte}\\
du & \textbf{zahltest}\\
er/sie/es & \textbf{zahlte}\\
wir & \textbf{zahlten}\\
ihr & \textbf{zahltet}\\
sie/Sie & \textbf{zahlten}\\
\end{tabular}
\hfill
\begin{tabular}{rl}
\multicolumn{2}{l}{\textbf{IMPERATIV}}\\
& \textbf{zahle (du)!}\\
& \textbf{zahlen wir!}\\
& \textbf{zahlt (ihr)!}\\
& \textbf{zahlen Sie!}\\
\end{tabular}

\vspace{10pt}

\begin{tabular}{rl}
\multicolumn{2}{l}{\textbf{PERFEKT}}\\
& haben + \textbf{gezahlt}\\
\end{tabular}
\hfill
\begin{tabular}{rl}
\multicolumn{2}{l}{\textbf{FUTURE I}}\\
& werden + \textbf{zahlen}\\
\end{tabular}
\hfill
\begin{tabular}{rl}
\multicolumn{2}{l}{\textbf{PLUSQUAMPERFEKT}}\\
& hatten + \textbf{gezahlt}\\
\end{tabular}

\vspace{10pt}

\begin{tabular}{lll}
\multicolumn{3}{l}{\textbf{MIT PRÄPOSITIONEN}}\\
& \textbf{zahlen für} + Akk & (to pay for)\\
& \textbf{zahlen an} + Akk & (to pay to)\\
\end{tabular}
\hfill
\begin{tabular}{lll}
\multicolumn{3}{l}{\textbf{TRENNBARE VERBEN}}\\
\vspace{20pt}
\end{tabular}
\vspace{-5pt}
\end{verbbox}



%zeigen (A1) – to show
\begin{verbbox-acc}\addcontentsline{toc}{subsubsection}{zeigen (to show / point out)}
    {\Large \bf zeigen} (to show / point out) \hfill Type: \textbf{regular} \hfill Auxiliary Verb: \textbf{haben} 
    
     \vspace{5pt}

    \hrule
    
    \vspace{5pt}

  \begin{tabular}{rl}
     \multicolumn{2}{l}{\textbf{PRÄSENS} }\\
    ich & \textbf{zeige}\\
    du & \textbf{zeigst}\\
    er/sie/es & \textbf{zeigt}\\
    wir & \textbf{zeigen}\\
    ihr & \textbf{zeigt}\\
    sie/Sie & \textbf{zeigen}\\
  \end{tabular}
  \hfill
  \begin{tabular}{rl}
    \multicolumn{2}{l}{\textbf{PRÄTERITUM} }\\
    ich & \textbf{zeigte}\\
    du & \textbf{zeigtest}\\
    er/sie/es & \textbf{zeigte}\\
    wir & \textbf{zeigten}\\
    ihr & \textbf{zeigtet}\\
    sie/Sie & \textbf{zeigten}\\
  \end{tabular}
  \hfill
  \begin{tabular}{rl}
     \multicolumn{2}{l}{\textbf{IMPERATIV} }\\
   & \textbf{zeige (du)!}\\
   & \textbf{zeigen wir!}\\
   & \textbf{zeigt (ihr)!}\\
   & \textbf{zeigen Sie!}\\
   & \vspace{10pt}
  \end{tabular}

 \vspace{10pt}

 \begin{tabular}{rl}
     \multicolumn{2}{l}{\textbf{PERFEKT}}\\
   & haben + \textbf{gezeigt}\\
  \end{tabular}
\hfill
\begin{tabular}{rl}
    \multicolumn{2}{l}{\textbf{FUTURE I}}\\
   & werden + \textbf{zeigen}\\
  \end{tabular}
\hfill
\begin{tabular}{rl}
    \multicolumn{2}{l}{\textbf{PLUSQUAMPERFEKT}}\\
   & hatten + \textbf{gezeigt}\\
  \end{tabular}

 \vspace{10pt}

\begin{tabular}{lll}
      \multicolumn{3}{l}{\textbf{MIT PRÄPOSITIONEN}}\\
& \textbf{zeigen auf} + Akk & (to point to / indicate)\\
& \textbf{zeigen für} + Akk & (to show for / demonstrate)\\
& \textbf{zeigen mit} + Dat & (to show with)\\
& \textbf{zeigen von} + Dat & (to show of / about)\\
\end{tabular}
  \hfill
\begin{tabular}{lll}
      \multicolumn{3}{l}{\textbf{TRENNBARE VERBEN}}\\
 & \textbf{vor$|$zeigen} & (to demonstrate / exhibit)\\
 & \textbf{auf$|$zeigen} & (to point out / indicate)\\
 & \textbf{nach$|$zeigen} & (to show afterward / demonstrate later)\\
\vspace{10pt}
  \end{tabular}
  
  \vspace{-5pt}

\end{verbbox-acc}

%zerstören (to destroy)
\begin{verbbox}\addcontentsline{toc}{subsubsection}{zerstören (to destroy)}
    {\Large \bf zerstören} (to destroy) \hfill Type: \textbf{regular, inseparable} \hfill Auxiliary Verb: \textbf{haben} 
    
     \vspace{5pt}

    \hrule
    
    \vspace{5pt}

  \begin{tabular}{rl}
     \multicolumn{2}{l}{\textbf{PRÄSENS} }\\
    ich & \textbf{zerstöre}\\
    du & \textbf{zerstörst}\\
    er/sie/es & \textbf{zerstört}\\
    wir & \textbf{zerstören}\\
    ihr & \textbf{zerstört}\\
    sie/Sie & \textbf{zerstören}\\
  \end{tabular}
  \hfill
     \begin{tabular}{rl}
    \multicolumn{2}{l}{\textbf{PRÄTERITUM} }\\
    ich & \textbf{zerstörte}\\
    du & \textbf{zerstörtest}\\
    er/sie/es & \textbf{zerstörte}\\
    wir & \textbf{zerstörten}\\
    ihr & \textbf{zerstörtet}\\
    sie/Sie & \textbf{zerstörten}\\
  \end{tabular}
  \hfill
  \begin{tabular}{rl}
     \multicolumn{2}{l}{\textbf{IMPERATIV} }\\
   & \textbf{zerstöre (du)!}\\
   & \textbf{zerstören wir!}\\
   & \textbf{zerstört (ihr)!}\\
   & \textbf{zerstören Sie!}\\
   & \vspace{10pt}
  \end{tabular}

    \vspace{10pt}

 \begin{tabular}{rl}
     \multicolumn{2}{l}{\textbf{PERFEKT}}\\
   & haben + \textbf{zerstört}\\
  \end{tabular}
\hfill
   \begin{tabular}{rl}
    \multicolumn{2}{l}{\textbf{FUTURE I} }\\
   & werden + \textbf{zerstören}\\
  \end{tabular}
\hfill
   \begin{tabular}{rl}
      \multicolumn{2}{l}{\textbf{PLUSQUAMPERFEKT}}\\
   & hatten + \textbf{zerstört}\\
  \end{tabular}

   \vspace{10pt}

   \begin{tabular}{lll}
      \multicolumn{3}{l}{\textbf{MIT PRÄPOSITIONEN}}\\
& \textbf{zerstören durch} + Akk & (to destroy through / by)\\
& \textbf{zerstören mit} + Dat & (to destroy with)\\
\end{tabular}
  \hfill
   \begin{tabular}{lll}
      \multicolumn{3}{l}{\textbf{TRENNBARE VERBEN}}\\
 & \textbf{---} & (none common)\\
\vspace{10pt}
  \end{tabular}

  \vspace{-5pt}

\end{verbbox}

%ziehen (A2) – to pull
\begin{verbbox}\addcontentsline{toc}{subsubsection}{ziehen (to pull)}
    {\Large \bf ziehen} (to pull) \hfill Type: \textbf{irregular} \hfill Auxiliary Verb: \textbf{haben} 
    
     \vspace{5pt}

    \hrule
    
    \vspace{5pt}

  \begin{tabular}{rl}
     \multicolumn{2}{l}{\textbf{PRÄSENS} }\\
    ich & \textbf{ziehe}\\
    du & \textbf{ziehst}\\
    er/sie/es & \textbf{zieht}\\
    wir & \textbf{ziehen}\\
    ihr & \textbf{zieht}\\
    sie/Sie & \textbf{ziehen}\\
  \end{tabular}
  \hfill
     \begin{tabular}{rl}
    \multicolumn{2}{l}{\textbf{PRÄTERITUM} }\\
    ich & \textbf{zog}\\
    du & \textbf{zogst}\\
    er/sie/es & \textbf{zog}\\
    wir & \textbf{zogen}\\
    ihr & \textbf{zogt}\\
    sie/Sie & \textbf{zogen}\\
  \end{tabular}
  \hfill
  \begin{tabular}{rl}
     \multicolumn{2}{l}{\textbf{IMPERATIV} }\\
   & \textbf{zieh(e) (du)!}\\
   & \textbf{ziehen wir!}\\
   & \textbf{zieht (ihr)!}\\
   & \textbf{ziehen Sie!}\\
   & \vspace{10pt}
  \end{tabular}

    \vspace{10pt}

 \begin{tabular}{rl}
     \multicolumn{2}{l}{\textbf{PERFEKT}}\\
  &  haben + \textbf{gezogen}\\
  \end{tabular}
\hfill
   \begin{tabular}{rl}
   
    \multicolumn{2}{l}{\textbf{FUTURE I} }\\
  &  werden + \textbf{ziehen}\\
 
  \end{tabular}
\hfill
   \begin{tabular}{rl}
      \multicolumn{2}{l}{\textbf{PLUSQUAMPERFEKT}}\\
  &  hatten + \textbf{gezogen}
  \end{tabular}
  


  \vspace{10pt}

   \begin{tabular}{lll}
      \multicolumn{3}{l}{\textbf{MIT PRÄPOSITIONEN}}\\
& \textbf{ziehen an} + Dat & (to pull on / at)\\
& \textbf{ziehen nach} + Dat & (to move to / after)\\
& \textbf{ziehen in} + Akk & (to move into)\\
& \textbf{ziehen aus} + Dat & (to move out of)\\
& \textbf{ziehen für} + Akk & (to pull for / draw for)\\
\vspace{10pt}
  \end{tabular}
  \hfill
   \begin{tabular}{lll}
      \multicolumn{3}{l}{\textbf{TRENNBARE VERBEN}}\\
 & \textbf{ein$|$ziehen} & (to move in / enlist / draw in)\\
 & \textbf{aus$|$ziehen} & (to move out / undress / extract)\\
 & \textbf{um$|$ziehen} & (to move / relocate / change clothes)\\
 & \textbf{zurück$|$ziehen} & (to pull back / withdraw)\\
 & \textbf{ab$|$ziehen} & (to subtract / withdraw / remove)\\
 & \textbf{auf$|$ziehen} & (to wind up / draw up / open)\\
 & \textbf{durch$|$ziehen} & (to pull through / carry out / complete)
  \end{tabular}


  \vspace{-5pt}


\end{verbbox}









achten
bemühen
gießen
untershützen
sich gewöhnen
klapen
kämpfen
sich erholen
sich erkundigen
protestieren
sich konzentrieren
schimpfen
riecehn
sich schützen
sterben
streiten 
sich verabschieden
sich wenden
zeifeln



































\begin{comment}



\begin{verbbox}\addcontentsline{toc}{subsubsection}{wirken (to appear)}
    {\Large \bf wirken} (to appear) \hfill Type: \textbf{regular} \hfill Auxiliary Verb: \textbf{haben} 
    
     \vspace{5pt}

    \hrule
    
    \vspace{5pt}

  \begin{tabular}{rl}
     \multicolumn{2}{l}{\textbf{PRÄSENS} }\\
    ich & \textbf{wirke}\\
    du & \textbf{wirkst}\\
    er/sie/es & \textbf{wirkt}\\
    wir & \textbf{wirken}\\
    ihr & \textbf{wirkt}\\
    sie/Sie & \textbf{wirken}\\
  \end{tabular}
  \hfill
     \begin{tabular}{rl}
    \multicolumn{2}{l}{\textbf{PRÄTERITUM} }\\
    ich & \textbf{wirkte}\\
    du & \textbf{wirktest}\\
    er/sie/es & \textbf{wirkte}\\
    wir & \textbf{wirkten}\\
    ihr & \textbf{wirktet}\\
    sie/Sie & \textbf{wirkten}\\
  \end{tabular}
  \hfill
  \begin{tabular}{rl}
     \multicolumn{2}{l}{\textbf{IMPERATIV} }\\
   & \textbf{wirk(e) (du)!}\\
   & \textbf{wirken wir!}\\
   & \textbf{wirkt (ihr)!}\\
   & \textbf{wirken Sie!}\\
   & \vspace{10pt}
  \end{tabular}

    \vspace{10pt}

 \begin{tabular}{rl}
     \multicolumn{2}{l}{\textbf{PERFEKT}}\\
   & haben + \textbf{gewirkt}\\
  \end{tabular}
\hfill
   \begin{tabular}{rl}
   
    \multicolumn{2}{l}{\textbf{FUTURE I} }\\
   & werden + \textbf{wirken}\\
 
  \end{tabular}
\hfill
   \begin{tabular}{rl}
      \multicolumn{2}{l}{\textbf{PLUSQUAMPERFEKT}}\\
   & hatten + \textbf{gewirkt}\\
  \end{tabular}

   \vspace{10pt}

   \begin{tabular}{lll}
      \multicolumn{3}{l}{\textbf{MIT PRÄPOSITIONEN}}\\
& \textbf{wirken auf} + Akk & (to influence)\\
& \textbf{wirken gegen} + Akk & (to work against)\\
\end{tabular}
  \hfill
   \begin{tabular}{lll}
      \multicolumn{3}{l}{\textbf{TRENNBARE VERBEN}}\\
 & \textbf{mit$|$wirken} & (to participate to)\\
\vspace{10pt}
  \end{tabular}


  \vspace{-5pt}


\end{verbbox}
\end{comment}


\end{document}